\documentclass[11pt,a4paper]{report}
\usepackage{fontspec,amsmath,amssymb,amsthm,microtype,xcolor}
\setmainfont{lmroman10-regular.otf}[BoldFont=lmroman10-bold.otf,ItalicFont=lmroman10-italic.otf,BoldItalicFont=lmroman10-bolditalic.otf]
\usepackage[most]{tcolorbox}
\usepackage{mathrsfs}
\providecommand{\lt}{<}
\providecommand{\gt}{>}
\setmonofont{DejaVuSansMono.ttf}[BoldFont=DejaVuSansMono-Bold.ttf,ItalicFont=DejaVuSansMono-Oblique.ttf,BoldItalicFont=DejaVuSansMono-BoldOblique.ttf,Scale=0.86]
\newfontfamily\leanfallback{STIXTwoMath-Regular.otf}[Scale=0.86]
\definecolor{leanrule}{gray}{0.62}
\definecolor{leanback}{gray}{0.962}
\definecolor{leantitle}{gray}{0.22}
\definecolor{leanmuted}{gray}{0.38}
\newtcolorbox{leanbox}[1]{enhanced,breakable,colback=leanback,colframe=leanrule,
  boxrule=0.4pt,arc=0.6mm,left=2.4mm,right=2.4mm,top=1.8mm,bottom=1.8mm,
  before skip=0.8em,after skip=1em,fonttitle=\small\bfseries,coltitle=leantitle,
  colbacktitle=leanback,titlerule=0.3pt,title={#1}}
\newlength{\leancharwidth}
\newenvironment{leancode}{\par\smallskip\begingroup\ttfamily\footnotesize\raggedright
  \settowidth{\leancharwidth}{M}\setlength{\parskip}{0pt}\setlength{\parindent}{0pt}
  \frenchspacing}{\par\endgroup\smallskip}
\newcommand{\leanline}[2]{\hangindent=\dimexpr#1\leancharwidth+4\leancharwidth\relax
  \hangafter=1\noindent\hspace*{#1\leancharwidth}#2\par}
\newcommand{\leansource}[1]{\par{\footnotesize\color{leanmuted}#1}\par}
\newcommand{\leanlabel}[1]{\par\smallskip\noindent{\small\bfseries #1}\par}
\usepackage[margin=25mm]{geometry}
\usepackage[unicode,hidelinks]{hyperref}
\hypersetup{pdfauthor={LangevinLean Project},bookmarksdepth=2}
\setlength{\emergencystretch}{3em}
\setlength{\parskip}{0.4em}
\setlength{\parindent}{0pt}
\setcounter{tocdepth}{1}
\makeatletter
\renewcommand*\l@section{\@dottedtocline{1}{1.5em}{3.8em}}
\renewcommand*\@pnumwidth{2.5em}
\renewcommand*\@tocrmarg{3.5em}
\makeatother
\setcounter{secnumdepth}{2}
\allowdisplaybreaks[2]
\raggedbottom
\title{\texorpdfstring{Langevin Monte Carlo with Average Smoothness}{Langevin Monte Carlo with Average Smoothness}\\[0.6em]\large Mathematical proofs}
\author{LangevinLean Project}
\date{Draft report LL-09\\[0.4em]\small Content revision 1c93063f2624}
\begin{document}
\pagenumbering{roman}
\maketitle
\chapter*{About this document}
\addcontentsline{toc}{chapter}{About this document}
Chapter 1 states the assumptions and the main results. Chapter 2 proves them, following the argument from the hypotheses to the conclusion in order. Chapter 3 proves the supporting lemmas, each before its first use. The proofs are written for readers with an undergraduate background in analysis, linear algebra and probability; definitions are recalled where they are first used, and each inequality is justified.

\section*{Mathematics and its Lean formalization}
Every result is followed by a grey box headed \emph{Lean}. The box reproduces the exact statement of the corresponding Lean declaration, read from the source file and line range named in the box, and then explains how the Lean proof follows the mathematical argument: which hypothesis encodes which assumption, and which formal step carries out which step of the proof. The mathematics can be read without the boxes; the boxes connect each claim to the statement that Lean has checked.

\section*{Checking the proofs}
The results are verified by Lean 4 with the Mathlib library, at the versions pinned in the repository \url{https://github.com/odakyildiz/LangevinLean}. After installing the Lean toolchain manager elan and cloning the repository, run from its root:
\begin{leancode}
\leanline{0}{lake\ exe\ cache\ get}
\leanline{0}{lake\ build\ AverageLMC}
\leanline{0}{lake\ env\ lean\ AverageLMC/Check.lean}
\end{leancode}
The second command checks every proof in the library. The third prints the axioms on which the main theorems depend; they are only the standard axioms of propositional extensionality, choice and quotient soundness.

Prepared by the LangevinLean Project with AI assistance, these notes form part of an ongoing effort to formalize Langevin algorithms comprehensively. This is draft work and may contain mistakes.

\section*{Literature and scope}
\begin{enumerate}
\item Arnak S. Dalalyan and Avetik Karagulyan (2026). \emph{Improved Guarantees for Langevin Monte Carlo with Average Smoothness}. arXiv preprint arXiv:2605.31413, version 1 (29 May 2026). \url{https://arxiv.org/abs/2605.31413v1}.
\end{enumerate}
Results from the earlier reports \cite{LangevinLean2026ULASC,LangevinLean2026ULALSI,LangevinLean2026IPLA,LangevinLean2026Proximal} are used below. Each application cites the section containing the required result.

\tableofcontents
\clearpage
\pagenumbering{arabic}
\chapter{Assumptions and main results}
\section{\texorpdfstring{Theorem 1: the actual one-step average-smoothness bound}{Theorem 1: the actual one-step average-smoothness bound}}


Let \(f:\mathbb R^p\to\mathbb R\) be differentiable,
\(m\)-strongly convex with \(M_\infty\)-Lipschitz gradient, where
\(m>0\). The coordinate constants are the literal Taylor suprema in the
paper, and \(M_{\rm av}=p^{-1}\sum_j M_j\). For the actual independent
Gaussian step \(X^+=X-h\nabla f(X)+\sqrt{2h}\,Z\), every \(h\ge0\) satisfies


\[W_2^2(\nu^+,\pi)\le
\frac{1-mh}{1+mh}W_2^2(\nu,\pi)+
\frac{(M_{\rm av}+m)h^2p}{1+mh}+
\frac{h^2(hM_\infty-1)}{1+mh}\mathbb E_\nu\|\nabla f\|^2.\]


The energy comparison, Gaussian smoothing, Brownian orthogonality and
transport infimum are proved from these hypotheses. This real-valued local
formula is stated for finite quadratic transport, so every signed energy term
is defined. The nonnegative global estimates below also cover infinite
initial transport.

\begin{leanbox}{Lean: \texttt{lmc\_\allowbreak{}one\_\allowbreak{}step}}
\leansource{\texttt{lmc\_\allowbreak{}one\_\allowbreak{}step} \quad AverageLMC/\allowbreak{}GaussianLMC.lean, lines 19–50}
\begin{leancode}
\leanline{0}{theorem\ lmc\_one\_step\ \{p\ :\ ℕ\}\ (hp\ :\ 0\ <\ p)\ (f\ :\ State\ p\ →\ ℝ)}
\leanline{4}{(m\ h\ :\ ℝ)\ (L\ :\ ℝ≥0)\ (hm\ :\ 0\ <\ m)\ (hh\ :\ 0\ ≤\ h)}
\leanline{4}{(hf\ :\ Differentiable\ ℝ\ f)\ (hsc\ :\ StrongConvexOn\ univ\ m\ f)}
\leanline{4}{(hLip\ :\ LipschitzWith\ L\ (gradient\ f))}
\leanline{4}{(μ\ :\ Measure\ (State\ p))\ [IsProbabilityMeasure\ μ]}
\leanline{4}{(hfinite\ :\ wasserstein2\ μ\ (gibbsMeasure\ f)\ ≠\ ⊤)\ :}
\leanline{4}{wasserstein2Real\ (ulaLawStep\ f\ h\ μ)\ (gibbsMeasure\ f)\ \textasciicircum{}\ 2\ ≤}
\leanline{6}{(1\ -\ m\ *\ h)\ /\ (1\ +\ m\ *\ h)\ *\ wasserstein2Real\ μ\ (gibbsMeasure\ f)\ \textasciicircum{}\ 2\ +}
\leanline{6}{((averageSmoothness\ f\ +\ m)\ *\ h\ \textasciicircum{}\ 2\ *\ p)\ /\ (1\ +\ m\ *\ h)\ +}
\leanline{6}{h\ \textasciicircum{}\ 2\ *\ (h\ *\ (L\ :\ ℝ)\ -\ 1)\ /\ (1\ +\ m\ *\ h)\ *}
\leanline{8}{(∫\ x,\ ‖gradient\ f\ x‖\ \textasciicircum{}\ 2\ ∂μ)}
\end{leancode}
\end{leanbox}

{\small\color{gray}Supporting proof: Section~\ref{proof-module-AverageLMC-GaussianLMC-lean}.\par}
\section{\texorpdfstring{Theorem 1: constant-step convergence}{Theorem 1: constant-step convergence}}


For finite \(W_2(\nu_0,\pi)\) and \(0\le h\le1/M_\infty\), the score correction
is nonpositive. Iterating the actual Gaussian transition gives


\[W_2^2(\nu_k,\pi)\le e^{-2mkh}W_2^2(\nu_0,\pi)
+\frac{(M_{\rm av}+m)hp}{2m}.\]


The proof propagates finite moments and discharges the recurrence for each
actual marginal law. An extended-nonnegative version proves the same bound
for every probability initialization, without adding a moment hypothesis.

\begin{leanbox}{Lean: \texttt{lmc\_\allowbreak{}constant\_\allowbreak{}step}}
\leansource{\texttt{lmc\_\allowbreak{}constant\_\allowbreak{}step} \quad AverageLMC/\allowbreak{}Theorem.lean, lines 84–111}
\begin{leancode}
\leanline{0}{theorem\ lmc\_constant\_step\ \{p\ :\ ℕ\}\ (hp\ :\ 0\ <\ p)\ (f\ :\ State\ p\ →\ ℝ)}
\leanline{4}{(m\ h\ :\ ℝ)\ (L\ :\ ℝ≥0)\ (hm\ :\ 0\ <\ m)\ (hh\ :\ 0\ ≤\ h)}
\leanline{4}{(hf\ :\ Differentiable\ ℝ\ f)\ (hsc\ :\ StrongConvexOn\ univ\ m\ f)}
\leanline{4}{(hLip\ :\ LipschitzWith\ L\ (gradient\ f))\ (hstep\ :\ h\ *\ (L\ :\ ℝ)\ ≤\ 1)}
\leanline{4}{(μ\ :\ Measure\ (State\ p))\ [IsProbabilityMeasure\ μ]}
\leanline{4}{(hfinite\ :\ wasserstein2\ μ\ (gibbsMeasure\ f)\ ≠\ ⊤)\ (k\ :\ ℕ)\ :}
\leanline{4}{wasserstein2Real\ ((ulaLawStep\ f\ h)\textasciicircum{}[k]\ μ)\ (gibbsMeasure\ f)\ \textasciicircum{}\ 2\ ≤}
\leanline{6}{Real.exp\ (-2\ *\ m\ *\ (k\ :\ ℝ)\ *\ h)\ *\ wasserstein2Real\ μ\ (gibbsMeasure\ f)\ \textasciicircum{}\ 2\ +}
\leanline{8}{(averageSmoothness\ f\ +\ m)\ *\ h\ *\ p\ /\ (2\ *\ m)}
\end{leancode}
\end{leanbox}

{\small\color{gray}Supporting proof: Section~\ref{proof-module-AverageLMC-Theorem-lean}.\par}
\section{\texorpdfstring{Proposition 1: the exact decreasing-step schedule}{Proposition 1: the exact decreasing-step schedule}}


Write \(\kappa_\infty=M_\infty/m\),
\(\kappa_{\rm av}=M_{\rm av}/m\), and suppose
\(W_2^2(\nu_0,\pi)\le\varepsilon_0^2p/m\). Set


\[k_0=\left\lceil\frac{\kappa_\infty}{2}
\log_+\frac{2\kappa_\infty\varepsilon_0^2}{\kappa_{\rm av}+1}\right\rceil,
\qquad h_k=\frac1{M_\infty+m(k-k_0)_+}.\]


The actual varying-step chain obeys, for every \(k\ge k_0\),


\[W_2^2(\nu_k,\pi)\le
\frac{\kappa_{\rm av}+1}{\kappa_\infty+k-k_0}\frac p m.\]


The literal schedule, positive-part burn-in and iteration ceiling are
checked, including zero initial error. The paired statement bounds the
real-valued distance, so finiteness of \(W_2(\nu_0,\pi)\) is a separate
hypothesis there; a companion version states the initial bound in
\([0,\infty]\) and derives this finiteness from it.

\begin{leanbox}{Lean: \texttt{lmc\_\allowbreak{}variable\_\allowbreak{}step}}
\leansource{\texttt{lmc\_\allowbreak{}variable\_\allowbreak{}step} \quad AverageLMC/\allowbreak{}Theorem.lean, lines 168–227}
\begin{leancode}
\leanline{0}{theorem\ lmc\_variable\_step\ \{p\ :\ ℕ\}\ (hp\ :\ 0\ <\ p)\ (f\ :\ State\ p\ →\ ℝ)}
\leanline{4}{(m\ ε₀\ :\ ℝ)\ (L\ :\ ℝ≥0)\ (hm\ :\ 0\ <\ m)}
\leanline{4}{(hf\ :\ Differentiable\ ℝ\ f)\ (hsc\ :\ StrongConvexOn\ univ\ m\ f)}
\leanline{4}{(hLip\ :\ LipschitzWith\ L\ (gradient\ f))}
\leanline{4}{(μ\ :\ Measure\ (State\ p))\ [IsProbabilityMeasure\ μ]}
\leanline{4}{(hfinite\ :\ wasserstein2\ μ\ (gibbsMeasure\ f)\ ≠\ ⊤)}
\leanline{4}{(hinit\ :\ wasserstein2Real\ μ\ (gibbsMeasure\ f)\ \textasciicircum{}\ 2\ ≤\ ε₀\ \textasciicircum{}\ 2\ *\ p\ /\ m)}
\leanline{4}{(k\ :\ ℕ)\ (hk\ :\ warmup\ ((L\ :\ ℝ)\ /\ m)\ (averageSmoothness\ f\ /\ m\ +\ 1)\ ε₀\ ≤\ k)\ :}
\leanline{4}{let\ κ\ :=\ (L\ :\ ℝ)\ /\ m}
\leanline{4}{let\ a\ :=\ averageSmoothness\ f\ /\ m\ +\ 1}
\leanline{4}{let\ k₀\ :=\ warmup\ κ\ a\ ε₀}
\leanline{4}{wasserstein2Real\ (lmcLaws\ f\ (lmcVariableSchedule\ (L\ :\ ℝ)\ m\ k₀)\ μ\ k)}
\leanline{6}{(gibbsMeasure\ f)\ \textasciicircum{}\ 2\ ≤}
\leanline{6}{a\ /\ (κ\ +\ ((k\ -\ k₀\ :\ ℕ)\ :\ ℝ))\ *\ p\ /\ m}
\end{leancode}
\end{leanbox}

{\small\color{gray}Supporting proof: Section~\ref{proof-module-AverageLMC-Theorem-lean}.\par}
\section{\texorpdfstring{Theorem 2 (th:3): LMC with the exact Laplacian condition}{Theorem 2 (th:3): LMC with the exact Laplacian condition}}


Let \(f\in C^3(\mathbb R^p)\) be \(m\)-strongly
convex with \(M_\infty\)-Lipschitz gradient, \(m>0\), and assume exactly


\[\|\nabla\Delta f(x)\|^2\le M_\Delta^2p\qquad(x\in\mathbb R^p),\quad M_\Delta\ge0.\]


For the actual Gaussian LMC chain with finite \(W_2(\nu_0,\pi)\) and
\(0<h<1/M_\infty\),


\[W_2(\nu_k,\pi)\le(1-mh)^kW_2(\nu_0,\pi)
+\frac{(2M_\infty\sqrt{M_{\rm av}}+M_\Delta)h\sqrt p}{m}.\]


The proof derives a weighted Dynkin identity for the constructed stationary
Brownian diffusion from fresh Gaussian Euler steps. Compact \(C^2\) approximations
of its score preserve precisely the source \(C^3\) assumption; the Gibbs score
moment and gradient-generator bound retain the exact average smoothness.
Clipping and the transport infimum give the displayed constant without an
optimal-coupling premise. An extended version covers arbitrary probability
initializations. No Hessian Lipschitz condition or fourth derivative is used.

\begin{leanbox}{Lean: \texttt{lmc\_\allowbreak{}laplacian}}
\leansource{\texttt{lmc\_\allowbreak{}laplacian} \quad AverageLMC/\allowbreak{}LaplacianTheorem.lean, lines 45–84}
\begin{leancode}
\leanline{0}{theorem\ lmc\_laplacian\ \{p\ :\ ℕ\}\ (hp\ :\ 0\ <\ p)\ (f\ :\ State\ p\ →\ ℝ)}
\leanline{4}{(m\ h\ :\ ℝ)\ (L\ :\ ℝ≥0)\ (hm\ :\ 0\ <\ m)\ (hh\ :\ 0\ <\ h)}
\leanline{4}{(hsc\ :\ StrongConvexOn\ univ\ m\ f)\ (hLip\ :\ LipschitzWith\ L\ (gradient\ f))}
\leanline{4}{(hsmall\ :\ h\ *\ (L\ :\ ℝ)\ <\ 1)\ (Mdelta\ :\ ℝ)\ (hMdelta\ :\ 0\ ≤\ Mdelta)}
\leanline{4}{(hδ\ :\ LaplacianSmooth\ f\ Mdelta)\ (μ\ :\ Measure\ (State\ p))\ [IsProbabilityMeasure\ μ]}
\leanline{4}{(hfinite\ :\ wasserstein2\ μ\ (gibbsMeasure\ f)\ ≠\ ⊤)\ (k\ :\ ℕ)\ :}
\leanline{4}{wasserstein2Real\ ((ulaLawStep\ f\ h)\textasciicircum{}[k]\ μ)\ (gibbsMeasure\ f)\ ≤}
\leanline{6}{(1\ -\ m\ *\ h)\ \textasciicircum{}\ k\ *\ wasserstein2Real\ μ\ (gibbsMeasure\ f)\ +}
\leanline{6}{(2\ *\ (L\ :\ ℝ)\ *\ Real.sqrt\ (averageSmoothness\ f)\ +\ Mdelta)\ *\ h\ *\ Real.sqrt\ p\ /\ m}
\end{leancode}
\end{leanbox}

{\small\color{gray}Supporting proof: Section~\ref{proof-module-AverageLMC-LaplacianTheorem-lean}.\par}
\section{\texorpdfstring{The SGLD main theorem with exact component RMS smoothness}{The SGLD main theorem with exact component RMS smoothness}}


Let \(f=\sum_{i=1}^n f_i\), each component have a Lipschitz
gradient, and the full potential be \(n\bar m\)-strongly convex. With its
unique minimizer \(\theta_*\), the actual update is


\[X^+=X-hn(\nabla f_I(X)-\nabla f_I(\theta_*))+\sqrt{2h}\,Z,\]


where the uniform index and standard Gaussian are independent of each
other and the initialization. The component seminorms define exactly
\(M_{\rm rms}^2=n^{-1}\sum_i\|\nabla f_i\|_{\rm Lip}^2\), and
\(\bar M_{\rm av}=M_{\rm av}(f)/n\), as in the source. For
\(0\le h\le\bar m/(3nM_{\rm rms}^2)\),


\[W_2^2(\nu_k,\pi)\le(1+n\bar m h)^{-k}W_2^2(\nu_0,\pi)
+\left(\bar M_{\rm av}+2\bar m+
\frac{2M_{\rm rms}^2}{\bar m}\right)\frac{hp}{\bar m}.\]


The estimator is proved unbiased, its variance is computed exactly, and
disintegration constructs the required coupling. No convexity assumption on
the individual components or additional derivative order is used. The
bound is stated in \([0,\infty]\), so it covers every probability
initialization.

\begin{leanbox}{Lean: \texttt{sgld\_\allowbreak{}constant\_\allowbreak{}step\_\allowbreak{}paper}}
\leansource{\texttt{sgld\_\allowbreak{}constant\_\allowbreak{}step\_\allowbreak{}paper} \quad AverageLMC/\allowbreak{}SGLDTheorem.lean, lines 239–256}
\begin{leancode}
\leanline{0}{theorem\ sgld\_constant\_step\_paper\ \{p\ n\ :\ ℕ\}\ (hp\ :\ 0\ <\ p)\ (hn\ :\ 0\ <\ n)}
\leanline{4}{(f\ :\ Fin\ n\ →\ State\ p\ →\ ℝ)\ (hf\ :\ ∀\ i,\ Differentiable\ ℝ\ (f\ i))}
\leanline{4}{(L\ :\ Fin\ n\ →\ ℝ≥0)\ (hLip\ :\ ∀\ i,\ LipschitzWith\ (L\ i)\ (gradient\ (f\ i)))}
\leanline{4}{(mbar\ h\ :\ ℝ)\ (hmbar\ :\ 0\ <\ mbar)\ (hh\ :\ 0\ ≤\ h)}
\leanline{4}{(hsc\ :\ StrongConvexOn\ univ\ ((n\ :\ ℝ)\ *\ mbar)\ (componentPotential\ f))}
\leanline{4}{(hstep\ :\ h\ ≤\ mbar\ /\ (3\ *\ n\ *\ rmsComponentSmoothness\ f\ \textasciicircum{}\ 2))}
\leanline{4}{(anchor\ :\ State\ p)\ (hanchor\ :\ ∀\ x,\ componentPotential\ f\ anchor\ ≤\ componentPotential\ f\ x)}
\leanline{4}{(μ\ :\ Measure\ (State\ p))\ [IsProbabilityMeasure\ μ]\ (k\ :\ ℕ)\ :}
\leanline{4}{wasserstein2\ (sgldLaws\ f\ anchor\ h\ μ\ k)\ (gibbsMeasure\ (componentPotential\ f))\ \textasciicircum{}\ 2\ ≤}
\leanline{6}{ENNReal.ofReal\ ((1\ +\ (n\ :\ ℝ)\ *\ mbar\ *\ h)⁻¹\ \textasciicircum{}\ k)\ *}
\leanline{8}{wasserstein2\ μ\ (gibbsMeasure\ (componentPotential\ f))\ \textasciicircum{}\ 2\ +}
\leanline{6}{ENNReal.ofReal\ ((averageSmoothness\ (fun\ x\ ↦\ componentPotential\ f\ x\ /\ n)\ +}
\leanline{8}{2\ *\ mbar\ +\ 2\ *\ rmsComponentSmoothness\ f\ \textasciicircum{}\ 2\ /\ mbar)\ *\ h\ *\ p\ /\ mbar)}
\end{leancode}
\end{leanbox}

{\small\color{gray}Supporting proof: Section~\ref{proof-module-AverageLMC-SGLDTheorem-lean}.\par}
\chapter{Mathematical argument}
\section{\texorpdfstring{Gaussian smoothing and the energy comparison}{Gaussian smoothing and the energy comparison}}



\textbf{Objects and hypotheses.} Let \(p\ge1\), \(m>0\), and let \(f:\mathbb R^p\to\mathbb R\) be differentiable, \(m\)-strongly convex and have an \(M_\infty\)-Lipschitz gradient. Thus

\[f(y)\ge f(x)+\langle\nabla f(x),y-x\rangle+\frac m2\|y-x\|^2,\qquad\|\nabla f(x)-\nabla f(y)\|\le M_\infty\|x-y\|.\]

The Gibbs probability law is \(\pi(dx)=Z^{-1}e^{-f(x)}dx\), where strong convexity makes \(0<Z<\infty\) and gives finite second moments. A coupling of two probability laws is a joint law with those marginals, and \(W_2\) is the infimum of the root-mean-square displacement over couplings. Work first with \(W_2(\nu,\pi)<\infty\). Throughout, a norm of a random vector means \(\|X\|_{L^2}=(\mathbb E\|X\|_2^2)^{1/2}\); its inner product is \(\mathbb E\langle X,Y\rangle\). The gradient bound \(\|\nabla f(x)\|\le\|\nabla f(0)\|+M_\infty\|x\|\) makes all score second moments finite for the states considered. The two Taylor bounds give \(|f(x)|\le|f(0)|+\|\nabla f(0)\|\|x\|+M_\infty\|x\|^2/2\), justifying the expected potential differences.




The
supremum defining \(M_j\) gives
\(f(x+ae_j)\le f(x)+a\,\partial_jf(x)+\tfrac12M_ja^2\) for all \(x\) and
\(a\). Insert the coordinates of a standard Gaussian vector \(\xi\) one at a
time, with \(a=\sqrt{2s}\,\xi_j\). Each \(\xi_j\) is centered and independent
of the coordinates already inserted, so its linear term has mean zero, while
\(\mathbb E a^2=2s\). Summing over \(j\),


\[\mathbb E f(x+\sqrt{2s}\,\xi)-f(x)\le s\sum_{j=1}^pM_j=pM_{\rm av}s
\qquad(s\ge0).\]


Couple \(\vartheta\sim\nu\) with \(Y_0\sim\pi\) at finite cost. Let
\(Y_t=Y_0-\int_0^t\nabla f(Y_s)\,ds+B_t\), \(B_t=\sqrt2W_t\), be the
stationary Langevin diffusion driven by a Brownian motion independent of
\((\vartheta,Y_0)\). Put \(A=\vartheta-h\nabla f(\vartheta)\), so that
\(A+B_h\sim\nu^+\) and \(Y_t\sim\pi\). Write \(w=W_2^2(\nu^+,\pi)\) and use
\(L^2\) norms. The common noise cancels:


\[\|A+B_h-Y_h\|^2=\|A-Y_0\|^2+2\int_0^h\langle A+B_t-Y_t,\nabla f(Y_t)\rangle\,dt.\]


Strong convexity bounds the integrand by
\(2\mathbb E[f(A+B_t)-f(Y_t)]-m\|A+B_t-Y_t\|^2\). Stationarity and smoothing at
the independent centre \(A\) give
\(\mathbb E[f(A+B_t)-f(Y_t)]\le\mathbb E[f(A)-f(Y_0)]+pM_{\rm av}t\). The
future increment \(B_h-B_t\) is orthogonal to \(A+B_t-Y_t\) and has squared
norm \(2(h-t)p\). Since \(w\le\|A+B_h-Y_t\|^2\) for every \(t\), this gives
\(\|A+B_t-Y_t\|^2\ge w-2(h-t)p\). Integrating over \([0,h]\) yields


\[(1+mh)\,w\le\|A-Y_0\|^2+2h\,\mathbb E[f(A)-f(Y_0)]+(M_{\rm av}+m)h^2p.\]\begin{leanbox}{Lean: \texttt{integral\_\allowbreak{}stdGaussian\_\allowbreak{}smoothing}, \texttt{gaussian\_\allowbreak{}energy\_\allowbreak{}bound}}
\leansource{\texttt{integral\_\allowbreak{}stdGaussian\_\allowbreak{}smoothing} \quad AverageLMC/\allowbreak{}Smoothing.lean, lines 154–175}
\begin{leancode}
\leanline{0}{theorem\ integral\_stdGaussian\_smoothing\ \{p\ :\ ℕ\}\ (hp\ :\ 0\ <\ p)}
\leanline{4}{(f\ :\ State\ p\ →\ ℝ)\ (M\ :\ ℝ≥0)\ (hf\ :\ Differentiable\ ℝ\ f)}
\leanline{4}{(hLip\ :\ LipschitzWith\ M\ (gradient\ f))\ (x\ :\ State\ p)\ (s\ :\ ℝ)\ (hs\ :\ 0\ ≤\ s)\ :}
\leanline{4}{(∫\ z,\ f\ (x\ +\ Real.sqrt\ (2\ *\ s)\ •\ z)\ -\ f\ x\ ∂stdGaussian\ (State\ p))\ ≤}
\leanline{6}{averageSmoothness\ f\ *\ (p\ :\ ℝ)\ *\ s}
\end{leancode}
\leansource{\texttt{gaussian\_\allowbreak{}energy\_\allowbreak{}bound} \quad AverageLMC/\allowbreak{}GaussianPerturbation.lean, lines 58–106}
\begin{leancode}
\leanline{0}{theorem\ gaussian\_energy\_bound\ (f\ :\ E\ →\ ℝ)\ (m\ h\ d\ a\ b\ w\ :\ ℝ)\ (L\ :\ ℝ≥0)}
\leanline{4}{(hf\ :\ Differentiable\ ℝ\ f)\ (hsc\ :\ StrongConvexOn\ univ\ m\ f)}
\leanline{4}{(hLip\ :\ LipschitzWith\ L\ (gradient\ f))\ (hm\ :\ 0\ ≤\ m)\ (hh\ :\ 0\ ≤\ h)}
\leanline{4}{(X\ :\ Lp\ E\ 2\ P)\ (Y\ B\ :\ ℝ\ →\ Lp\ E\ 2\ P)\ (hY\ :\ Continuous\ Y)\ (hB\ :\ Continuous\ B)}
\leanline{4}{(hYe\ :\ ∀\ t\ ∈\ Icc\ 0\ h,\ Y\ t\ -\ Y\ 0\ =}
\leanline{6}{-(∫\ s\ in\ 0..t,\ ULA.lpDrift\ P\ (gradient\ f)\ hLip\ (Y\ s))\ +\ B\ t)}
\leanline{4}{(horth\ :\ ∀\ t\ ∈\ Icc\ 0\ h,\ {\leanfallback ⟪}X\ +\ B\ t\ -\ Y\ t,\ B\ h\ -\ B\ t{\leanfallback ⟫}\ =\ 0)}
\leanline{4}{(hnoise\ :\ ∀\ t\ ∈\ Icc\ 0\ h,\ ‖B\ h\ -\ B\ t‖\ \textasciicircum{}\ 2\ =\ 2\ *\ (h\ -\ t)\ *\ d)}
\leanline{4}{(hcost\ :\ ∀\ t\ ∈\ Icc\ 0\ h,\ w\ ≤\ ‖X\ +\ B\ h\ -\ Y\ t‖\ \textasciicircum{}\ 2)}
\leanline{4}{(henergy\ :\ ∀\ t\ ∈\ Icc\ 0\ h,}
\leanline{6}{(∫\ ω,\ f\ ((X\ +\ B\ t)\ ω)\ ∂P)\ -\ (∫\ ω,\ f\ (Y\ t\ ω)\ ∂P)\ ≤\ a\ +\ b\ *\ t)\ :}
\leanline{4}{(1\ +\ m\ *\ h)\ *\ w\ ≤\ ‖X\ -\ Y\ 0‖\ \textasciicircum{}\ 2\ +\ 2\ *\ h\ *\ a\ +\ h\ \textasciicircum{}\ 2\ *\ b\ +\ m\ *\ h\ \textasciicircum{}\ 2\ *\ d}
\end{leancode}
\leanlabel{How the Lean proof follows the argument}
\texttt{integral\_\allowbreak{}std\allowbreak{}Gaussian\_\allowbreak{}smoothing} specializes the independent-coordinate smoothing theorem to the standard Gaussian law. In its finite-coordinate induction, centering removes the linear term and the unit second moment evaluates the quadratic term. \texttt{gaussian\_\allowbreak{}energy\_\allowbreak{}bound} inserts this smoothing inequality into the common-noise stationary comparison; the Brownian orthogonality identity contributes the extra curvature-times-dimension term. All these are integral statements with integrability discharged before linearity is used.
\end{leanbox}
\section{\texorpdfstring{The one-step Wasserstein bound}{The one-step Wasserstein bound}}


With
\(G=\mathbb E\|\nabla f(\vartheta)\|^2\), strong convexity gives the
gradient-step inequality


\[\|A-Y_0\|^2\le(1-mh)\|\vartheta-Y_0\|^2+2h\,\mathbb E[f(Y_0)-f(\vartheta)]+h^2G,\]


and the \(M_\infty\)-Lipschitz gradient gives the descent inequality
\(2\mathbb E[f(A)-f(\vartheta)]\le-h(2-hM_\infty)G\). Adding both to the energy
comparison cancels every potential value:


\[(1+mh)\,W_2^2(\nu^+,\pi)\le(1-mh)\|\vartheta-Y_0\|^2+(M_{\rm av}+m)h^2p
+h^2(hM_\infty-1)G.\]


The output law does not depend on the initial coupling. Taking the
infimum over finite-cost couplings, which is valid also when \(1-mh<0\),
and dividing by \(1+mh\) gives Theorem 1's one-step bound for \(h>0\). At
\(h=0\) the transition is the identity.




\textbf{The cancellation explicitly.} Strong convexity at \(\vartheta\), with comparison point \(Y_0\), gives

\[-2h\mathbb E\langle\nabla f(\vartheta),\vartheta-Y_0\rangle\le2h\mathbb E[f(Y_0)-f(\vartheta)]-mh\|\vartheta-Y_0\|_{L^2}^2.\]

Expanding \(\|\vartheta-h\nabla f(\vartheta)-Y_0\|_{L^2}^2\) produces the gradient-step bound. For the descent bound apply the upper Taylor model at \(\vartheta\) with displacement \(-h\nabla f(\vartheta)\): its linear term is \(-h\|\nabla f(\vartheta)\|^2\) and its quadratic term is \(M_\infty h^2\|\nabla f(\vartheta)\|^2/2\). Adding the two inequalities to the energy comparison cancels \(f(Y_0)\), \(f(\vartheta)\) and \(f(A)\) exactly. The remaining score coefficient is \(h^2-2h^2+h^3M_\infty=h^2(hM_\infty-1)\). Every denominator is positive because \(m>0\) and \(h\ge0\).


\begin{leanbox}{Lean: \texttt{lp\_\allowbreak{}gradient\_\allowbreak{}step}, \texttt{lp\_\allowbreak{}gradient\_\allowbreak{}descent}, \texttt{lmc\_\allowbreak{}one\_\allowbreak{}step}}
\leansource{\texttt{lp\_\allowbreak{}gradient\_\allowbreak{}step} \quad AverageLMC/\allowbreak{}ExpectedGradientStep.lean, lines 15–33}
\begin{leancode}
\leanline{0}{theorem\ lp\_gradient\_step\ (f\ :\ E\ →\ ℝ)\ (m\ h\ :\ ℝ)\ (L\ :\ ℝ≥0)}
\leanline{4}{(hf\ :\ Differentiable\ ℝ\ f)\ (hsc\ :\ StrongConvexOn\ Set.univ\ m\ f)}
\leanline{4}{(hLip\ :\ LipschitzWith\ L\ (gradient\ f))\ (hh\ :\ 0\ ≤\ h)\ (X\ Y\ :\ Lp\ E\ 2\ P)\ :}
\leanline{4}{‖X\ -\ h\ •\ ULA.lpDrift\ P\ (gradient\ f)\ hLip\ X\ -\ Y‖\ \textasciicircum{}\ 2\ ≤}
\leanline{6}{(1\ -\ m\ *\ h)\ *\ ‖X\ -\ Y‖\ \textasciicircum{}\ 2\ +}
\leanline{6}{2\ *\ h\ *\ ((∫\ ω,\ f\ (Y\ ω)\ ∂P)\ -\ ∫\ ω,\ f\ (X\ ω)\ ∂P)\ +}
\leanline{6}{h\ \textasciicircum{}\ 2\ *\ ‖ULA.lpDrift\ P\ (gradient\ f)\ hLip\ X‖\ \textasciicircum{}\ 2}
\end{leancode}
\leansource{\texttt{lp\_\allowbreak{}gradient\_\allowbreak{}descent} \quad AverageLMC/\allowbreak{}ExpectedGradientStep.lean, lines 37–63}
\begin{leancode}
\leanline{0}{theorem\ lp\_gradient\_descent\ (f\ :\ E\ →\ ℝ)\ (h\ :\ ℝ)\ (L\ :\ ℝ≥0)}
\leanline{4}{(hf\ :\ Differentiable\ ℝ\ f)\ (hLip\ :\ LipschitzWith\ L\ (gradient\ f))}
\leanline{4}{(hh\ :\ 0\ ≤\ h)\ (X\ :\ Lp\ E\ 2\ P)\ :}
\leanline{4}{2\ *\ ((∫\ ω,\ f\ ((X\ -\ h\ •\ ULA.lpDrift\ P\ (gradient\ f)\ hLip\ X)\ ω)\ ∂P)\ -}
\leanline{6}{∫\ ω,\ f\ (X\ ω)\ ∂P)\ ≤}
\leanline{6}{-h\ *\ (2\ -\ h\ *\ (L\ :\ ℝ))\ *\ ‖ULA.lpDrift\ P\ (gradient\ f)\ hLip\ X‖\ \textasciicircum{}\ 2}
\end{leancode}
\leansource{\texttt{lmc\_\allowbreak{}one\_\allowbreak{}step} \quad AverageLMC/\allowbreak{}GaussianLMC.lean, lines 19–50}
\begin{leancode}
\leanline{0}{theorem\ lmc\_one\_step\ \{p\ :\ ℕ\}\ (hp\ :\ 0\ <\ p)\ (f\ :\ State\ p\ →\ ℝ)}
\leanline{4}{(m\ h\ :\ ℝ)\ (L\ :\ ℝ≥0)\ (hm\ :\ 0\ <\ m)\ (hh\ :\ 0\ ≤\ h)}
\leanline{4}{(hf\ :\ Differentiable\ ℝ\ f)\ (hsc\ :\ StrongConvexOn\ univ\ m\ f)}
\leanline{4}{(hLip\ :\ LipschitzWith\ L\ (gradient\ f))}
\leanline{4}{(μ\ :\ Measure\ (State\ p))\ [IsProbabilityMeasure\ μ]}
\leanline{4}{(hfinite\ :\ wasserstein2\ μ\ (gibbsMeasure\ f)\ ≠\ ⊤)\ :}
\leanline{4}{wasserstein2Real\ (ulaLawStep\ f\ h\ μ)\ (gibbsMeasure\ f)\ \textasciicircum{}\ 2\ ≤}
\leanline{6}{(1\ -\ m\ *\ h)\ /\ (1\ +\ m\ *\ h)\ *\ wasserstein2Real\ μ\ (gibbsMeasure\ f)\ \textasciicircum{}\ 2\ +}
\leanline{6}{((averageSmoothness\ f\ +\ m)\ *\ h\ \textasciicircum{}\ 2\ *\ p)\ /\ (1\ +\ m\ *\ h)\ +}
\leanline{6}{h\ \textasciicircum{}\ 2\ *\ (h\ *\ (L\ :\ ℝ)\ -\ 1)\ /\ (1\ +\ m\ *\ h)\ *}
\leanline{8}{(∫\ x,\ ‖gradient\ f\ x‖\ \textasciicircum{}\ 2\ ∂μ)}
\end{leancode}
\leanlabel{How the Lean proof follows the argument}
\texttt{lp\_\allowbreak{}gradient\_\allowbreak{}step} and \texttt{lp\_\allowbreak{}gradient\_\allowbreak{}descent} are the two expected deterministic inequalities whose potential terms cancel. \texttt{lmc\_\allowbreak{}one\_\allowbreak{}step} first separates the zero step. In the positive branch, \texttt{hsm} is the exact coordinate smoothing estimate, \texttt{houtput} proves the output transport finite, and \texttt{wasserstein2\_\allowbreak{}sq\_\allowbreak{}bound\_\allowbreak{}of\_\allowbreak{}finite\_\allowbreak{}couplings} passes the common bound to the infimum even when its input-cost coefficient has either sign. The final field normalization divides by the strictly positive denominator.
\end{leanbox}
\section{\texorpdfstring{Constant-step convergence}{Constant-step convergence}}


For
\(0\le h\le1/M_\infty\) the score term is nonpositive and is dropped. The
coordinate bounds give \(m\le M_{\rm av}\le M_\infty\), so \(t=mh\in[0,1]\).
Finite second moments propagate from \(\nu_0\), so
\(r_k=W_2^2(\nu_k,\pi)\) satisfies, at every step,


\[r_{k+1}\le q\,r_k+\frac{(M_{\rm av}+m)h^2p}{1+t},\qquad q=\frac{1-t}{1+t}.\]


Since \(1-q=2t/(1+t)\), the fixed point is
\(c=(M_{\rm av}+m)hp/(2m)\), and induction gives
\(r_k\le q^kr_0+(1-q^k)c\le q^kr_0+c\). Finally \(q\le e^{-2t}\) on
\([0,1]\), so


\[W_2^2(\nu_k,\pi)\le e^{-2mkh}W_2^2(\nu_0,\pi)+\frac{(M_{\rm av}+m)hp}{2m}.\]


If \(W_2(\nu_0,\pi)=\infty\), the same bound holds trivially in
\([0,\infty]\).




\textbf{Induction explicitly.} If \(h=0\), the chain is constant and the assertion follows by equality. If \(h>0\), write \(F=(M_{\rm av}+m)h^2p/(1+mh)\). Since \(0\le q<1\), the identity \(F=(1-q)c\) holds with \(c=(M_{\rm av}+m)hp/(2m)\). At \(k=0\) the proposed sharper bound \(r_k\le q^kr_0+(1-q^k)c\) is equality. If it holds at \(k\), multiply it by \(q\ge0\), add \((1-q)c\), and obtain \(q^{k+1}r_0+(1-q^{k+1})c\). This proves the induction. For \(0\le t<1\), differentiation gives

\[\frac{d}{dt}\left(\log\frac{1-t}{1+t}+2t\right)=-\frac2{1-t^2}+2=-\frac{2t^2}{1-t^2}\le0.\]

The expression vanishes at zero, so \(q\le e^{-2t}\). At \(t=1\), \(q=0\) and the same inequality is immediate. Raising this nonnegative comparison to the natural power \(k\) yields the displayed exponential rate.


\begin{leanbox}{Lean: \texttt{contraction\_\allowbreak{}le\_\allowbreak{}exp}, \texttt{constant\_\allowbreak{}step\_\allowbreak{}normalized}, \texttt{lmc\_\allowbreak{}constant\_\allowbreak{}step}}
\leansource{\texttt{contraction\_\allowbreak{}le\_\allowbreak{}exp} \quad AverageLMC/\allowbreak{}Recurrence.lean, lines 21–34}
\begin{leancode}
\leanline{0}{theorem\ contraction\_le\_exp\ \{t\ :\ ℝ\}\ (ht\ :\ 0\ ≤\ t)\ (ht1\ :\ t\ ≤\ 1)\ :}
\leanline{4}{(1\ -\ t)\ /\ (1\ +\ t)\ ≤\ Real.exp\ (-2\ *\ t)}
\end{leancode}
\leansource{\texttt{constant\_\allowbreak{}step\_\allowbreak{}normalized} \quad AverageLMC/\allowbreak{}Recurrence.lean, lines 78–83}
\begin{leancode}
\leanline{0}{theorem\ constant\_step\_normalized\ (r\ :\ ℕ\ →\ ℝ)\ (a\ t\ :\ ℝ)}
\leanline{4}{(ha\ :\ 0\ ≤\ a)\ (ht\ :\ 0\ ≤\ t)\ (ht1\ :\ t\ ≤\ 1)\ (hr0\ :\ 0\ ≤\ r\ 0)}
\leanline{4}{(hstep\ :\ ∀\ j,\ r\ (j\ +\ 1)\ ≤}
\leanline{6}{(1\ -\ t)\ /\ (1\ +\ t)\ *\ r\ j\ +\ a\ *\ t\ \textasciicircum{}\ 2\ /\ (1\ +\ t))\ (n\ :\ ℕ)\ :}
\leanline{4}{r\ n\ ≤\ Real.exp\ (-2\ *\ t\ *\ (n\ :\ ℝ))\ *\ r\ 0\ +\ a\ *\ t\ /\ 2}
\end{leancode}
\leansource{\texttt{lmc\_\allowbreak{}constant\_\allowbreak{}step} \quad AverageLMC/\allowbreak{}Theorem.lean, lines 84–111}
\begin{leancode}
\leanline{0}{theorem\ lmc\_constant\_step\ \{p\ :\ ℕ\}\ (hp\ :\ 0\ <\ p)\ (f\ :\ State\ p\ →\ ℝ)}
\leanline{4}{(m\ h\ :\ ℝ)\ (L\ :\ ℝ≥0)\ (hm\ :\ 0\ <\ m)\ (hh\ :\ 0\ ≤\ h)}
\leanline{4}{(hf\ :\ Differentiable\ ℝ\ f)\ (hsc\ :\ StrongConvexOn\ univ\ m\ f)}
\leanline{4}{(hLip\ :\ LipschitzWith\ L\ (gradient\ f))\ (hstep\ :\ h\ *\ (L\ :\ ℝ)\ ≤\ 1)}
\leanline{4}{(μ\ :\ Measure\ (State\ p))\ [IsProbabilityMeasure\ μ]}
\leanline{4}{(hfinite\ :\ wasserstein2\ μ\ (gibbsMeasure\ f)\ ≠\ ⊤)\ (k\ :\ ℕ)\ :}
\leanline{4}{wasserstein2Real\ ((ulaLawStep\ f\ h)\textasciicircum{}[k]\ μ)\ (gibbsMeasure\ f)\ \textasciicircum{}\ 2\ ≤}
\leanline{6}{Real.exp\ (-2\ *\ m\ *\ (k\ :\ ℝ)\ *\ h)\ *\ wasserstein2Real\ μ\ (gibbsMeasure\ f)\ \textasciicircum{}\ 2\ +}
\leanline{8}{(averageSmoothness\ f\ +\ m)\ *\ h\ *\ p\ /\ (2\ *\ m)}
\end{leancode}
\leanlabel{How the Lean proof follows the argument}
\texttt{lmc\_\allowbreak{}one\_\allowbreak{}step\_\allowbreak{}stable} first proves that the score correction is nonpositive. \texttt{lmc\_\allowbreak{}constant\_\allowbreak{}step} defines the actual squared-distance sequence as \texttt{r}; \texttt{lmc\allowbreak{}Laws\_\allowbreak{}finite} allows the one-step theorem to be applied at every index. The block \texttt{hs} is the scalar recurrence after normalization. \texttt{constant\_\allowbreak{}step\_\allowbreak{}normalized} iterates it and uses \texttt{contraction\_\allowbreak{}le\_\allowbreak{}exp}. Finally \texttt{lmc\allowbreak{}Laws\_\allowbreak{}const} identifies the recursively constructed marginal laws with iteration of the Gaussian transition.
\end{leanbox}
\section{\texorpdfstring{The decreasing-step schedule}{The decreasing-step schedule}}


Normalize
\(r_k=mW_2^2(\nu_k,\pi)/p\), \(a=\kappa_{\rm av}+1\) and
\(t_k=mh_k=1/(\kappa_\infty+(k-k_0)_+)\). Since \(h_kM_\infty\le1\), the
stable one-step bound reads


\[r_{k+1}\le\frac{1-t_k}{1+t_k}\,r_k+\frac{a\,t_k^2}{1+t_k}.\]


Before \(k_0\) the step is constant, \(t_k=1/\kappa_\infty\le1\), so the
constant-step estimate gives
\(r_{k_0}\le e^{-2k_0/\kappa_\infty}\varepsilon_0^2+a/(2\kappa_\infty)\). The
ceiling defining \(k_0\) makes
\(e^{-2k_0/\kappa_\infty}\varepsilon_0^2\le a/(2\kappa_\infty)\), hence
\(r_{k_0}\le a/\kappa_\infty\). After the warmup, \(r\le a/x\) with \(x\ge1\)
implies


\[\frac{1-1/x}{1+1/x}\,r+\frac{a/x^2}{1+1/x}
\le\frac{x-1}{x+1}\cdot\frac ax+\frac a{x(x+1)}=\frac a{x+1}.\]


Induction with \(x=\kappa_\infty+k-k_0\) gives
\(r_k\le a/(\kappa_\infty+k-k_0)\) for \(k\ge k_0\). Multiplying by \(p/m\)
is Proposition 1. When \(\varepsilon_0=0\), \(k_0=0\) and the induction starts
at once.




\textbf{The warmup ceiling.} For \(\varepsilon_0\ne0\), set \(R=2\kappa_\infty\varepsilon_0^2/a\) and \(\log_+R=\max(\log R,0)\). If \(R\le1\), the initial error already satisfies \(\varepsilon_0^2\le a/(2\kappa_\infty)\), and the warmup may be zero. If \(R>1\), the ceiling ensures \(2k_0/\kappa_\infty\ge\log R\); exponentiating its negative gives \(e^{-2k_0/\kappa_\infty}\varepsilon_0^2\le a/(2\kappa_\infty)\). In both cases the constant-step bias contributes the other half of \(a/\kappa_\infty\). After warmup \(x=\kappa_\infty+k-k_0\ge1\), so the factor \((x-1)/(x+1)\) multiplying the induction hypothesis is nonnegative. The rational calculation above therefore preserves the inequality. Positive dimension and positive \(m\) justify the final multiplication by \(p/m\). When \(\varepsilon_0=0\), the logarithmic warmup is defined to be zero and \(r_0=0\le a/\kappa_\infty\), giving the same induction directly.


\begin{leanbox}{Lean: \texttt{warmup\_\allowbreak{}decay}, \texttt{harmonic\_\allowbreak{}step}, \texttt{lmc\_\allowbreak{}variable\_\allowbreak{}step}}
\leansource{\texttt{warmup\_\allowbreak{}decay} \quad AverageLMC/\allowbreak{}Complexity.lean, lines 39–56}
\begin{leancode}
\leanline{0}{theorem\ warmup\_decay\ \{κ\ a\ ε₀\ :\ ℝ\}\ (hκ\ :\ 0\ <\ κ)\ (ha\ :\ 0\ <\ a)}
\leanline{4}{(hε₀\ :\ 0\ <\ ε₀)\ :}
\leanline{4}{Real.exp\ (-2\ *\ (warmup\ κ\ a\ ε₀\ :\ ℝ)\ /\ κ)\ *\ ε₀\ \textasciicircum{}\ 2\ ≤\ a\ /\ (2\ *\ κ)}
\end{leancode}
\leansource{\texttt{harmonic\_\allowbreak{}step} \quad AverageLMC/\allowbreak{}Recurrence.lean, lines 86–100}
\begin{leancode}
\leanline{0}{theorem\ harmonic\_step\ \{a\ x\ r\ :\ ℝ\}\ (hx\ :\ 1\ ≤\ x)}
\leanline{4}{(hr\ :\ r\ ≤\ a\ /\ x)\ :}
\leanline{4}{(1\ -\ 1\ /\ x)\ /\ (1\ +\ 1\ /\ x)\ *\ r\ +\ a\ *\ (1\ /\ x)\ \textasciicircum{}\ 2\ /\ (1\ +\ 1\ /\ x)}
\leanline{6}{≤\ a\ /\ (x\ +\ 1)}
\end{leancode}
\leansource{\texttt{lmc\_\allowbreak{}variable\_\allowbreak{}step} \quad AverageLMC/\allowbreak{}Theorem.lean, lines 168–227}
\begin{leancode}
\leanline{0}{theorem\ lmc\_variable\_step\ \{p\ :\ ℕ\}\ (hp\ :\ 0\ <\ p)\ (f\ :\ State\ p\ →\ ℝ)}
\leanline{4}{(m\ ε₀\ :\ ℝ)\ (L\ :\ ℝ≥0)\ (hm\ :\ 0\ <\ m)}
\leanline{4}{(hf\ :\ Differentiable\ ℝ\ f)\ (hsc\ :\ StrongConvexOn\ univ\ m\ f)}
\leanline{4}{(hLip\ :\ LipschitzWith\ L\ (gradient\ f))}
\leanline{4}{(μ\ :\ Measure\ (State\ p))\ [IsProbabilityMeasure\ μ]}
\leanline{4}{(hfinite\ :\ wasserstein2\ μ\ (gibbsMeasure\ f)\ ≠\ ⊤)}
\leanline{4}{(hinit\ :\ wasserstein2Real\ μ\ (gibbsMeasure\ f)\ \textasciicircum{}\ 2\ ≤\ ε₀\ \textasciicircum{}\ 2\ *\ p\ /\ m)}
\leanline{4}{(k\ :\ ℕ)\ (hk\ :\ warmup\ ((L\ :\ ℝ)\ /\ m)\ (averageSmoothness\ f\ /\ m\ +\ 1)\ ε₀\ ≤\ k)\ :}
\leanline{4}{let\ κ\ :=\ (L\ :\ ℝ)\ /\ m}
\leanline{4}{let\ a\ :=\ averageSmoothness\ f\ /\ m\ +\ 1}
\leanline{4}{let\ k₀\ :=\ warmup\ κ\ a\ ε₀}
\leanline{4}{wasserstein2Real\ (lmcLaws\ f\ (lmcVariableSchedule\ (L\ :\ ℝ)\ m\ k₀)\ μ\ k)}
\leanline{6}{(gibbsMeasure\ f)\ \textasciicircum{}\ 2\ ≤}
\leanline{6}{a\ /\ (κ\ +\ ((k\ -\ k₀\ :\ ℕ)\ :\ ℝ))\ *\ p\ /\ m}
\end{leancode}
\leanlabel{How the Lean proof follows the argument}
\texttt{warmup\_\allowbreak{}decay} checks the initial exponential allocation, including the logarithmic ceiling. \texttt{harmonic\_\allowbreak{}step} proves the single rational induction step after warmup. \texttt{lmc\_\allowbreak{}variable\_\allowbreak{}step} verifies admissibility of every actual step in \texttt{hs}, normalizes the initial cost in \texttt{hi}, and derives the normalized recurrence in \texttt{hrec}. \texttt{variable\_\allowbreak{}step\_\allowbreak{}normalized\_\allowbreak{}all} handles zero initial error explicitly before applying the scalar warmup-and-tail estimate.
\end{leanbox}
\section{\texorpdfstring{The Laplacian refinement}{The Laplacian refinement}}


Use the coupling and
stationary diffusion of the first step, with \(L^2\) inner products. Put
\(\Psi=(\vartheta-h\nabla f(\vartheta))-(Y_0-h\nabla f(Y_0))\) and
\(R=\int_0^h(\nabla f(Y_t)-\nabla f(Y_0))\,dt\). The output error is
\(A+B_h-Y_h=\Psi+R\), and for \(hM_\infty\le1\) the gradient step contracts:
\(\|\Psi\|\le(1-mh)\|\vartheta-Y_0\|\).




The diffusion generator maps \(\nabla f\) to
\(\nabla\Delta f-\nabla^2f\,\nabla f\). Because \(\Psi\) is independent of the
Brownian motion, a weighted Dynkin identity, stationarity,
\(\|\nabla\Delta f\|\le M_\Delta\sqrt p\) and the Gibbs score moment
\(\mathbb E_\pi\|\nabla f\|^2\le pM_{\rm av}\) give


\[\langle\Psi,\nabla f(Y_t)-\nabla f(Y_0)\rangle\le t\,\|\Psi\|\,G_\Delta,
\qquad G_\Delta=M_\infty\sqrt{M_{\rm av}p}+M_\Delta\sqrt p.\]


The Dynkin identity is derived for compactly supported \(C^2\) tests from
fresh Gaussian Euler steps. Cutoffs of the score components and radial
clipping of \(\Psi\) extend it to this pairing. The integral equation,
Brownian moments and \(\mathbb E_\pi\|\nabla f\|^2\le pM_\infty\) give
\(\|R\|\le\tfrac32M_\infty h^{3/2}\sqrt p\). Expanding \(\|\Psi+R\|^2\) and
taking the coupling infimum, \(e_k=W_2(\nu_k,\pi)\) and \(q=1-mh\) satisfy


\[e_{k+1}^2\le q^2e_k^2+q\,h^2G_\Delta\,e_k+\bigl(\tfrac32M_\infty h^{3/2}\sqrt p\bigr)^2.\]


A recurrence \(e_{k+1}^2\le q^2e_k^2+2qCe_k+B^2\) with
\(C\le(1-q)a\) and \(B^2\le(1-q^2)b^2\) gives \(e_k\le q^ke_0+a+b\). Take
\(a=hG_\Delta/(2m)\) and \(b=3hM_\infty\sqrt{M_{\rm av}p}/(2m)\); the second
condition uses \(m\le M_{\rm av}\). Then


\[a+b=\frac{(4M_\infty\sqrt{M_{\rm av}}+M_\Delta)h\sqrt p}{2m}
\le\frac{(2M_\infty\sqrt{M_{\rm av}}+M_\Delta)h\sqrt p}{m},\]


which is Theorem 2 for \(0<h<1/M_\infty\).




\textbf{Bounds used by the recurrence.} Here the additional assumptions are \(f\in C^3\), \(M_\Delta\ge0\) and \(\|\nabla\Delta f(x)\|^2\le M_\Delta^2p\) for every \(x\), with \(0<hM_\infty<1\). The derivative of the gradient has operator norm at most \(M_\infty\). Thus Cauchy–Schwarz and the score moment bound give

\[\|\nabla\Delta f(Y_t)-\nabla^2f(Y_t)\nabla f(Y_t)\|_{L^2}\le(M_\Delta+M_\infty\sqrt{M_{\rm av}})\sqrt p=G_\Delta.\]

Integrating the pairing bound over time gives \(2\langle\Psi,R\rangle\le h^2G_\Delta\|\Psi\|_{L^2}\). For the total defect, the integral equation, Minkowski and the stationary score moment give \(\|Y_t-Y_0\|_{L^2}\le t\sqrt{M_\infty p}+\sqrt{2pt}\). Hence

\[\begin{aligned}\|R\|_{L^2}&\le M_\infty\int_0^h\left(t\sqrt{M_\infty p}+\sqrt{2pt}\right)dt\\&=M_\infty\left(\frac{h^2}2\sqrt{M_\infty p}+\frac{2\sqrt2}3h^{3/2}\sqrt p\right)\\&\le\left(\frac12+\frac{2\sqrt2}3\right)M_\infty h^{3/2}\sqrt p\le\frac32M_\infty h^{3/2}\sqrt p.\end{aligned}\]

The penultimate inequality uses \(hM_\infty\le1\); the last uses \(\sqrt2\le3/2\). For the stated recurrence floor, \(C=h^2G_\Delta/2=(1-q)a\). Also \(B^2\le mh\,b^2\) follows from \(m\le M_{\rm av}\), and \(mh\le1-q^2=mh(2-mh)\). These verify both scalar recurrence hypotheses. The final enlargement of the nonnegative \(M_\Delta\) term gives precisely the printed coefficient.


\begin{leanbox}{Lean: \texttt{lp\_\allowbreak{}stationary\_\allowbreak{}gradient\_\allowbreak{}pairing\_\allowbreak{}bound}, \texttt{lp\_\allowbreak{}laplacian\_\allowbreak{}one\_\allowbreak{}step\_\allowbreak{}bound}, \texttt{laplacian\_\allowbreak{}recurrence\_\allowbreak{}bound\_\allowbreak{}of\_\allowbreak{}global\_\allowbreak{}defect}}
\leansource{\texttt{lp\_\allowbreak{}stationary\_\allowbreak{}gradient\_\allowbreak{}pairing\_\allowbreak{}bound} \quad AverageLMC/\allowbreak{}LaplacianPairing.lean, lines 50–210}
\begin{leancode}
\leanline{0}{theorem\ lp\_stationary\_gradient\_pairing\_bound\ \{p\ :\ ℕ\}\ (hp\ :\ 0\ <\ p)}
\leanline{4}{(f\ :\ State\ p\ →\ ℝ)\ (m\ :\ ℝ)\ (L\ :\ ℝ≥0)\ (hm\ :\ 0\ <\ m)}
\leanline{4}{(hsc\ :\ StrongConvexOn\ univ\ m\ f)\ (hLip\ :\ LipschitzWith\ L\ (gradient\ f))}
\leanline{4}{(Mdelta\ :\ ℝ)\ (hMdelta\ :\ 0\ ≤\ Mdelta)\ (hδ\ :\ LaplacianSmooth\ f\ Mdelta)}
\leanline{4}{(initial\ :\ Lp\ (State\ p)\ 2\ P)}
\leanline{4}{(B\ :\ Fin\ p\ →\ ℝ≥0\ →\ Ω\ →\ ℝ)\ (hB\ :\ ∀\ i,\ IsPreBrownianReal\ (B\ i)\ P)}
\leanline{4}{(hscalar\ :\ ∀\ r\ :\ ℝ,\ 0\ <\ r\ →\ iIndepFun\ (fun\ q\ :\ ℕ\ ×\ Fin\ p\ ↦\ fun\ a\ ↦}
\leanline{6}{ULA.Hilbert.brownianEulerNoise\ P\ B\ hB\ r\ q.1\ a\ q.2)\ P)}
\leanline{4}{(Ψ\ :\ Lp\ (State\ p)\ 2\ P)}
\leanline{4}{(hind\ :\ ∀\ r\ :\ ℝ,\ 0\ <\ r\ →\ IndepFun\ (fun\ a\ ↦\ (Ψ\ a,\ initial\ a))}
\leanline{6}{(fun\ a\ k\ ↦\ ULA.Hilbert.brownianEulerNoise\ P\ B\ hB\ r\ k\ a)\ P)}
\leanline{4}{(Y\ :\ ℝ\ →\ Lp\ (State\ p)\ 2\ P)\ (hY\ :\ Continuous\ Y)}
\leanline{4}{(T\ :\ ℝ)\ (hT\ :\ 0\ ≤\ T)}
\leanline{4}{(hYe\ :\ ∀\ t\ ∈\ Icc\ 0\ T,\ Y\ t\ -\ initial\ =}
\leanline{6}{-(∫\ s\ in\ 0..t,\ ULA.lpDrift\ P\ (gradient\ f)\ hLip\ (Y\ s))\ +}
\leanline{8}{ULA.brownianLangevinNoise\ P\ B\ hB\ 0\ t)}
\leanline{4}{(hstationary\ :\ ∀\ t\ ∈\ Icc\ 0\ T,\ P.map\ (Y\ t)\ =\ ULA.gibbsMeasure\ f)\ :}
\leanline{4}{inner\ ℝ\ Ψ\ (ULA.lpDrift\ P\ (gradient\ f)\ hLip\ (Y\ T)\ -}
\leanline{6}{ULA.lpDrift\ P\ (gradient\ f)\ hLip\ initial)\ ≤}
\leanline{8}{T\ *\ ‖Ψ‖\ *\ ((L\ :\ ℝ)\ *\ Real.sqrt\ (averageSmoothness\ f\ *\ (p\ :\ ℝ))\ +}
\leanline{10}{Mdelta\ *\ Real.sqrt\ (p\ :\ ℝ))}
\end{leancode}
\leansource{\texttt{lp\_\allowbreak{}laplacian\_\allowbreak{}one\_\allowbreak{}step\_\allowbreak{}bound} \quad AverageLMC/\allowbreak{}LaplacianOneStep.lean, lines 21–88}
\begin{leancode}
\leanline{0}{theorem\ lp\_laplacian\_one\_step\_bound}
\leanline{4}{(f\ :\ E\ →\ ℝ)\ (m\ h\ d\ Rgen\ :\ ℝ)\ (M\ :\ ℝ≥0)}
\leanline{4}{(hf\ :\ Differentiable\ ℝ\ f)\ (hsc\ :\ StrongConvexOn\ univ\ m\ f)}
\leanline{4}{(hLip\ :\ LipschitzWith\ M\ (gradient\ f))\ (hm\ :\ 0\ <\ m)\ (hmM\ :\ m\ ≤\ (M\ :\ ℝ))}
\leanline{4}{(hh\ :\ 0\ ≤\ h)\ (hd\ :\ 0\ ≤\ d)\ (hsmall\ :\ h\ *\ (M\ :\ ℝ)\ ≤\ 1)\ (hRgen\ :\ 0\ ≤\ Rgen)}
\leanline{4}{(X\ :\ Lp\ E\ 2\ P)\ (Y\ B\ :\ ℝ\ →\ Lp\ E\ 2\ P)\ (hY\ :\ Continuous\ Y)}
\leanline{4}{(hG\ :\ ∀\ t\ ∈\ Icc\ 0\ h,\ ‖ULA.lpDrift\ P\ (gradient\ f)\ hLip\ (Y\ t)‖\ ≤}
\leanline{6}{Real.sqrt\ ((M\ :\ ℝ)\ *\ d))}
\leanline{4}{(hB\ :\ ∀\ t\ ∈\ Icc\ 0\ h,\ ‖B\ t‖\ ≤\ Real.sqrt\ (2\ *\ d)\ *\ Real.sqrt\ t)}
\leanline{4}{(hYe\ :\ ∀\ t\ ∈\ Icc\ 0\ h,\ Y\ t\ -\ Y\ 0\ =}
\leanline{6}{-(∫\ s\ in\ 0..t,\ ULA.lpDrift\ P\ (gradient\ f)\ hLip\ (Y\ s))\ +\ B\ t)}
\leanline{4}{(hpair\ :\ ∀\ t\ ∈\ Icc\ 0\ h,}
\leanline{6}{{\leanfallback ⟪}(X\ -\ h\ •\ ULA.lpDrift\ P\ (gradient\ f)\ hLip\ X)\ -}
\leanline{10}{(Y\ 0\ -\ h\ •\ ULA.lpDrift\ P\ (gradient\ f)\ hLip\ (Y\ 0)),}
\leanline{8}{ULA.lpDrift\ P\ (gradient\ f)\ hLip\ (Y\ t)\ -}
\leanline{10}{ULA.lpDrift\ P\ (gradient\ f)\ hLip\ (Y\ 0){\leanfallback ⟫}\ ≤}
\leanline{8}{t\ *\ ‖(X\ -\ h\ •\ ULA.lpDrift\ P\ (gradient\ f)\ hLip\ X)\ -}
\leanline{10}{(Y\ 0\ -\ h\ •\ ULA.lpDrift\ P\ (gradient\ f)\ hLip\ (Y\ 0))‖\ *\ Rgen)\ :}
\leanline{4}{‖X\ -\ h\ •\ ULA.lpDrift\ P\ (gradient\ f)\ hLip\ X\ +\ B\ h\ -\ Y\ h‖\ \textasciicircum{}\ 2\ ≤}
\leanline{6}{(1\ -\ m\ *\ h)\ \textasciicircum{}\ 2\ *\ ‖X\ -\ Y\ 0‖\ \textasciicircum{}\ 2\ +}
\leanline{6}{h\ \textasciicircum{}\ 2\ *\ (1\ -\ m\ *\ h)\ *\ ‖X\ -\ Y\ 0‖\ *\ Rgen\ +}
\leanline{6}{((3\ /\ 2\ :\ ℝ)\ *\ (M\ :\ ℝ)\ *\ h\ *\ Real.sqrt\ h\ *\ Real.sqrt\ d)\ \textasciicircum{}\ 2}
\end{leancode}
\leansource{\texttt{laplacian\_\allowbreak{}recurrence\_\allowbreak{}bound\_\allowbreak{}of\_\allowbreak{}global\_\allowbreak{}defect} \quad AverageLMC/\allowbreak{}LaplacianLocalError.lean, lines 57–116}
\begin{leancode}
\leanline{0}{theorem\ laplacian\_recurrence\_bound\_of\_global\_defect}
\leanline{4}{(e\ :\ ℕ\ →\ ℝ)\ (m\ M\ h\ A\ δ\ S\ :\ ℝ)}
\leanline{4}{(he\ :\ ∀\ k,\ 0\ ≤\ e\ k)\ (hm\ :\ 0\ <\ m)\ (hmM\ :\ m\ ≤\ M)}
\leanline{4}{(hh\ :\ 0\ ≤\ h)\ (hhM\ :\ h\ *\ M\ ≤\ 1)}
\leanline{4}{(hA\ :\ 0\ ≤\ A)\ (hmA\ :\ m\ ≤\ A\ \textasciicircum{}\ 2)\ (hδ\ :\ 0\ ≤\ δ)\ (hS\ :\ 0\ ≤\ S)}
\leanline{4}{(hstep\ :\ ∀\ k,\ (e\ (k\ +\ 1))\ \textasciicircum{}\ 2\ ≤}
\leanline{6}{(1\ -\ m\ *\ h)\ \textasciicircum{}\ 2\ *\ (e\ k)\ \textasciicircum{}\ 2\ +}
\leanline{8}{(1\ -\ m\ *\ h)\ *\ h\ \textasciicircum{}\ 2\ *\ (M\ *\ A\ +\ δ)\ *\ S\ *\ e\ k\ +}
\leanline{10}{((3\ /\ 2\ :\ ℝ)\ *\ M\ *\ h\ *\ Real.sqrt\ h\ *\ S)\ \textasciicircum{}\ 2)\ (n\ :\ ℕ)\ :}
\leanline{4}{e\ n\ ≤\ (1\ -\ m\ *\ h)\ \textasciicircum{}\ n\ *\ e\ 0\ +\ (2\ *\ M\ *\ A\ +\ δ)\ *\ h\ *\ S\ /\ m}
\end{leancode}
\leanlabel{How the Lean proof follows the argument}
\texttt{lp\_\allowbreak{}stationary\_\allowbreak{}gradient\_\allowbreak{}pairing\_\allowbreak{}bound} bounds the cross term with the actual stationary diffusion and the stated gradient-of-Laplacian assumption. \texttt{lp\_\allowbreak{}laplacian\_\allowbreak{}one\_\allowbreak{}step\_\allowbreak{}bound} combines that pairing with the local integral remainder and gradient-step contraction. \texttt{laplacian\_\allowbreak{}recurrence\_\allowbreak{}bound\_\allowbreak{}of\_\allowbreak{}global\_\allowbreak{}defect} iterates the resulting quadratic-plus-linear scalar estimate. The final theorem \texttt{lmc\_\allowbreak{}laplacian} obtains the required finite-distance recurrence for the actual Gaussian laws and simplifies its constant using the exact Gibbs score moment bound.
\end{leanbox}
\section{\texorpdfstring{Control-variate SGLD}{Control-variate SGLD}}


Let
\(g_I=n(\nabla f_I(\vartheta)-\nabla f_I(\theta_*))\). Since
\(\nabla f(\theta_*)=0\), its mean over the uniform index is
\(\nabla f(\vartheta)\). The SGLD output is the LMC output plus
\(-h(g_I-\nabla f(\vartheta))\), which has conditional mean zero given
\((\vartheta,\xi)\). Lifting any coupling of the LMC output with \(\pi\) to the
index by disintegration, the cross term vanishes, so


\[W_2^2(\nu^+,\pi)\le W_2^2(\nu^+_{\rm LMC},\pi)
+h^2\bigl(\mathbb E\|g_I\|^2-\mathbb E\|\nabla f(\vartheta)\|^2\bigr).\]


Write \(m=n\bar m\) and let \(M\) be a Lipschitz constant of \(\nabla f\).
The one-step LMC bound turns the score coefficient into
\(h^2(hM-1)/(1+mh)-h^2\le h^2(hM-2)/(1+mh)\), which is Proposition 3's first
claim. The component seminorms give
\(\mathbb E_I\|g_I\|^2\le n^2M_{\rm rms}^2\|\vartheta-\theta_*\|^2\), and the
Gibbs bound \(\mathbb E_\pi\|Y-\theta_*\|^2\le p/m\) gives
\(\mathbb E\|\vartheta-\theta_*\|^2\le2W_2^2(\nu,\pi)+2p/m\). Take
\(M=R=nM_{\rm rms}\). Because \(m\le R\), the step condition
\(3R^2h\le m\) implies \(hR\le1/3\), so the score term is nonpositive and


\[r_{k+1}\le\Bigl(\frac{1-mh}{1+mh}+2R^2h^2\Bigr)r_k
+\Bigl(\frac{M_{\rm av}(f)+m}{1+mh}+\frac{2R^2}m\Bigr)h^2p.\]


The step condition makes the contraction at most \((1+mh)^{-1}\) and the
forcing at most \((M_{\rm av}(f)+2m+2R^2/m)h^2p/(1+mh)\). The fixed point is
\((M_{\rm av}(f)+2m+2R^2/m)hp/m\). Substituting \(m=n\bar m\),
\(R=nM_{\rm rms}\) and \(M_{\rm av}(f)=n\bar M_{\rm av}\) gives the SGLD
bound; an infinite initial distance makes it trivial.




\textbf{Why the variance adds.} Conditional on \((\vartheta,\xi)\), the index \(I\) is uniform and its centered error \(\eta=g_I-\nabla f(\vartheta)\) has mean zero. Expanding a lifted coupled square gives \(\mathbb E\|D-h\eta\|^2=\mathbb E\|D\|^2+h^2\mathbb E\|\eta\|^2\), because the conditional expectation of \(\langle D,\eta\rangle\) is zero. Also \(\mathbb E_I\|\eta\|^2=\mathbb E_I\|g_I\|^2-\|\nabla f(\vartheta)\|^2\). If \(L_i\) is the actual Lipschitz seminorm of \(\nabla f_i\), then

\[\mathbb E_I\|g_I\|^2=n\sum_i\|\nabla f_i(\vartheta)-\nabla f_i(\theta_*)\|^2\le n\sum_iL_i^2\|\vartheta-\theta_*\|^2=n^2M_{\rm rms}^2\|\vartheta-\theta_*\|^2.\]

For any coupling with Gibbs, \(\|\vartheta-\theta_*\|^2\le2\|\vartheta-Y\|^2+2\|Y-\theta_*\|^2\). Integrate, use the Gibbs second moment, and pass costs to their infimum to obtain the displayed second-moment bound. With \(R=nM_{\rm rms}\), the small-step condition \(3R^2h\le m\) and \(m\le R\) imply \(mh\le1/3\). The forcing comparison uses \(2R^2h\le m\): after multiplication by \(1+mh\), its extra term is at most \(mh\) in the required coefficient. The recurrence fixed point is the forcing divided by \(1-(1+mh)^{-1}=mh/(1+mh)\). At \(h=0\) the transition is the identity and this division is unnecessary.


\begin{leanbox}{Lean: \texttt{sgld\_\allowbreak{}variance\_\allowbreak{}addition}, \texttt{controlVariateEstimator\_\allowbreak{}second\_\allowbreak{}moment\_\allowbreak{}le}, \texttt{sgld\_\allowbreak{}one\_\allowbreak{}step\_\allowbreak{}rms} and 1 more}
\leansource{\texttt{sgld\_\allowbreak{}variance\_\allowbreak{}addition} \quad AverageLMC/\allowbreak{}SGLDLaw.lean, lines 98–141}
\begin{leancode}
\leanline{0}{theorem\ sgld\_variance\_addition\ \{p\ n\ :\ ℕ\}\ (hn\ :\ 0\ <\ n)}
\leanline{4}{(f\ :\ Fin\ n\ →\ State\ p\ →\ ℝ)\ (hf\ :\ ∀\ i,\ Differentiable\ ℝ\ (f\ i))}
\leanline{4}{(L\ :\ Fin\ n\ →\ ℝ≥0)\ (hLip\ :\ ∀\ i,\ LipschitzWith\ (L\ i)\ (gradient\ (f\ i)))}
\leanline{4}{(M\ :\ ℝ≥0)\ (hfull\ :\ LipschitzWith\ M\ (gradient\ (componentPotential\ f)))}
\leanline{4}{(anchor\ :\ State\ p)\ (hstar\ :\ gradient\ (componentPotential\ f)\ anchor\ =\ 0)}
\leanline{4}{(h\ :\ ℝ)\ (μ\ π\ :\ Measure\ (State\ p))\ [IsProbabilityMeasure\ μ]\ [IsProbabilityMeasure\ π]}
\leanline{4}{(hμ\ :\ MemLp\ (fun\ x\ :\ State\ p\ ↦\ x)\ 2\ μ)\ (hπ\ :\ MemLp\ (fun\ x\ :\ State\ p\ ↦\ x)\ 2\ π)\ :}
\leanline{4}{ULA.wasserstein2Real\ (sgldLawStep\ f\ anchor\ h\ μ)\ π\ \textasciicircum{}\ 2\ ≤}
\leanline{6}{ULA.wasserstein2Real\ (ULA.ulaLawStep\ (componentPotential\ f)\ h\ μ)\ π\ \textasciicircum{}\ 2\ +}
\leanline{8}{h\ \textasciicircum{}\ 2\ *\ ∫\ q\ :\ (State\ p\ ×\ State\ p)\ ×\ Fin\ n,}
\leanline{10}{‖controlVariateError\ f\ anchor\ q.1.1\ q.2‖\ \textasciicircum{}\ 2}
\leanline{12}{∂(μ.prod\ (stdGaussian\ (State\ p))).prod\ (uniformIndexLaw\ n)}
\end{leancode}
\leansource{\texttt{controlVariateEstimator\_\allowbreak{}second\_\allowbreak{}moment\_\allowbreak{}le} \quad AverageLMC/\allowbreak{}SGLDEstimator.lean, lines 152–167}
\begin{leancode}
\leanline{0}{theorem\ controlVariateEstimator\_second\_moment\_le\ \{p\ n\ :\ ℕ\}\ (hp\ :\ 0\ <\ p)\ (hn\ :\ 0\ <\ n)}
\leanline{4}{(f\ :\ Fin\ n\ →\ State\ p\ →\ ℝ)\ (L\ :\ Fin\ n\ →\ ℝ≥0)}
\leanline{4}{(hLip\ :\ ∀\ i,\ LipschitzWith\ (L\ i)\ (gradient\ (f\ i)))\ (xstar\ x\ :\ State\ p)\ :}
\leanline{4}{(∫\ i,\ ‖controlVariateEstimator\ f\ xstar\ x\ i‖\ \textasciicircum{}\ 2\ ∂uniformIndexLaw\ n)\ ≤}
\leanline{6}{(n\ :\ ℝ)\ \textasciicircum{}\ 2\ *\ rmsComponentSmoothness\ f\ \textasciicircum{}\ 2\ *\ ‖x\ -\ xstar‖\ \textasciicircum{}\ 2}
\end{leancode}
\leansource{\texttt{sgld\_\allowbreak{}one\_\allowbreak{}step\_\allowbreak{}rms} \quad AverageLMC/\allowbreak{}SGLDTheorem.lean, lines 56–103}
\begin{leancode}
\leanline{0}{theorem\ sgld\_one\_step\_rms\ \{p\ n\ :\ ℕ\}\ (hp\ :\ 0\ <\ p)\ (hn\ :\ 0\ <\ n)}
\leanline{4}{(f\ :\ Fin\ n\ →\ State\ p\ →\ ℝ)\ (hf\ :\ ∀\ i,\ Differentiable\ ℝ\ (f\ i))}
\leanline{4}{(L\ :\ Fin\ n\ →\ ℝ≥0)\ (hLip\ :\ ∀\ i,\ LipschitzWith\ (L\ i)\ (gradient\ (f\ i)))}
\leanline{4}{(m\ h\ :\ ℝ)\ (M\ :\ ℝ≥0)\ (hm\ :\ 0\ <\ m)\ (hh\ :\ 0\ ≤\ h)}
\leanline{4}{(hsc\ :\ StrongConvexOn\ univ\ m\ (componentPotential\ f))}
\leanline{4}{(hfull\ :\ LipschitzWith\ M\ (gradient\ (componentPotential\ f)))}
\leanline{4}{(hstep\ :\ h\ *\ (M\ :\ ℝ)\ ≤\ 2)}
\leanline{4}{(anchor\ :\ State\ p)\ (hanchor\ :\ ∀\ x,\ componentPotential\ f\ anchor\ ≤\ componentPotential\ f\ x)}
\leanline{4}{(hstar\ :\ gradient\ (componentPotential\ f)\ anchor\ =\ 0)}
\leanline{4}{(μ\ :\ Measure\ (State\ p))\ [IsProbabilityMeasure\ μ]}
\leanline{4}{(hfinite\ :\ wasserstein2\ μ\ (gibbsMeasure\ (componentPotential\ f))\ ≠\ ⊤)\ :}
\leanline{4}{let\ R\ :=\ (n\ :\ ℝ)\ *\ rmsComponentSmoothness\ f}
\leanline{4}{wasserstein2Real\ (sgldLawStep\ f\ anchor\ h\ μ)\ (gibbsMeasure\ (componentPotential\ f))\ \textasciicircum{}\ 2\ ≤}
\leanline{6}{((1\ -\ m\ *\ h)\ /\ (1\ +\ m\ *\ h)\ +\ 2\ *\ R\ \textasciicircum{}\ 2\ *\ h\ \textasciicircum{}\ 2)\ *}
\leanline{8}{wasserstein2Real\ μ\ (gibbsMeasure\ (componentPotential\ f))\ \textasciicircum{}\ 2\ +}
\leanline{6}{((averageSmoothness\ (componentPotential\ f)\ +\ m)\ /\ (1\ +\ m\ *\ h)\ +\ 2\ *\ R\ \textasciicircum{}\ 2\ /\ m)\ *}
\leanline{8}{h\ \textasciicircum{}\ 2\ *\ p}
\end{leancode}
\leansource{\texttt{sgld\_\allowbreak{}recurrence\_\allowbreak{}bound} \quad AverageLMC/\allowbreak{}SGLDRecurrence.lean, lines 52–80}
\begin{leancode}
\leanline{0}{theorem\ sgld\_recurrence\_bound\ (r\ :\ ℕ\ →\ ℝ)\ (m\ R\ A\ h\ p\ :\ ℝ)}
\leanline{4}{(hm\ :\ 0\ <\ m)\ (hh\ :\ 0\ ≤\ h)\ (hp\ :\ 0\ ≤\ p)\ (hA\ :\ 0\ ≤\ A)}
\leanline{4}{(hR\ :\ m\ \textasciicircum{}\ 2\ ≤\ R\ \textasciicircum{}\ 2)\ (hstep\ :\ 3\ *\ R\ \textasciicircum{}\ 2\ *\ h\ ≤\ m)}
\leanline{4}{(hr\ :\ ∀\ n,\ 0\ ≤\ r\ n)}
\leanline{4}{(hrec\ :\ ∀\ n,\ r\ (n\ +\ 1)\ ≤}
\leanline{6}{((1\ -\ m\ *\ h)\ /\ (1\ +\ m\ *\ h)\ +\ 2\ *\ R\ \textasciicircum{}\ 2\ *\ h\ \textasciicircum{}\ 2)\ *\ r\ n\ +}
\leanline{6}{((A\ +\ m)\ /\ (1\ +\ m\ *\ h)\ +\ 2\ *\ R\ \textasciicircum{}\ 2\ /\ m)\ *\ h\ \textasciicircum{}\ 2\ *\ p)}
\leanline{4}{(k\ :\ ℕ)\ :}
\leanline{4}{r\ k\ ≤\ (1\ +\ m\ *\ h)⁻¹\ \textasciicircum{}\ k\ *\ r\ 0\ +}
\leanline{6}{(A\ +\ 2\ *\ m\ +\ 2\ *\ R\ \textasciicircum{}\ 2\ /\ m)\ *\ h\ *\ p\ /\ m}
\end{leancode}
\leanlabel{How the Lean proof follows the argument}
\texttt{sgld\_\allowbreak{}variance\_\allowbreak{}addition} uses a lifted coupling to remove the centered estimator cross term. \texttt{control\allowbreak{}Variate\allowbreak{}Estimator\_\allowbreak{}second\_\allowbreak{}moment\_\allowbreak{}le} bounds the estimator using the actual component Lipschitz seminorms. \texttt{sgld\_\allowbreak{}one\_\allowbreak{}step\_\allowbreak{}rms} combines this variance bound with the one-step LMC comparison. \texttt{sgld\_\allowbreak{}recurrence\_\allowbreak{}bound} uses the small-step condition to bound the contraction and forcing, then sums the affine recurrence. The extended theorem separates infinite initial transport before converting any distance to a real number.
\end{leanbox}
\chapter{Supporting proofs}
The following sections give the definitions and prove the lemmas used in the convergence arguments.

\section{\texorpdfstring{Coordinate smoothness constants of the potential}{Coordinate smoothness constants of the potential}}
\label{proof-module-AverageLMC-Assumptions-lean}


This section fixes the objects in terms of which every later estimate is stated. The state space is \(\mathbb R^p\) with the Euclidean inner product \(\langle x,y\rangle=\sum_{j=1}^px_jy_j\), the norm \(\|x\|=\langle x,x\rangle^{1/2}\) and the canonical basis \(e_1,\dots,e_p\). The potential is a function \(f:\mathbb R^p\to\mathbb R\) with gradient \(\nabla f\) and partial derivatives \(\partial_jf(x)=\langle\nabla f(x),e_j\rangle\). Following the notation section of the paper, the coordinate smoothness constants, their average and the global smoothness constant are


\[\begin{aligned}
M_j&=\sup_{x\in\mathbb R^p,\ a\ne0}
\frac{2\bigl(f(x+ae_j)-f(x)-a\,\partial_jf(x)\bigr)}{a^2},\\
M_{\rm av}&=\frac1p\sum_{j=1}^pM_j,\\
M_\infty&=\sup_{x\ne y}
\frac{\|\nabla f(x)-\nabla f(y)\|}{\|x-y\|}.
\end{aligned}\]


All three are computed from the potential itself; they are not additional parameters. For a differentiable \(f\) whose gradient is \(M\)-Lipschitz we prove that both suprema are finite, that \(\nabla f\) is \(M_\infty\)-Lipschitz, and that the coordinate quadratic upper model \(f(x+ae_j)\le f(x)+a\,\partial_jf(x)+\tfrac12M_ja^2\) holds. If in addition \(f\) is \(m\)-strongly convex, then \(m\le M_j\le M_\infty\) for every \(j\) and \(m\le M_{\rm av}\le M_\infty\). No second derivative is used for these facts.




The section also defines the Hessian matrix \(\nabla^2f(x)\), the Laplacian \(\Delta f(x)\), the third-order condition of the paper's Section 4, \(\sup_x\|\nabla\Delta f(x)\|^2\le M_\Delta^2p\) for a three times continuously differentiable potential, and the root-mean-square constant \(M_{\rm rms}\) for a finite sum of potentials, and it proves \(\operatorname{tr}\nabla^2f=\Delta f\). The coordinate upper model is the input of the Gaussian smoothing lemma, the comparisons between \(m\), \(M_{\rm av}\) and \(M_\infty\) are used when the one-step bounds are iterated, and the remaining definitions enter the Laplacian-smooth theorem and the stochastic-gradient theorem.



Earlier results used here: \cite[Section~3.9]{LangevinLean2026ULASC}; \cite[Section~3.2]{LangevinLean2026ULASC}.
\subsection{\texorpdfstring{Coordinates and coordinate Taylor quotients}{Coordinates and coordinate Taylor quotients}}
\label{report-step-AverageLMC-Assumptions-lean-1}


\textbf{Definitions.} Let \(p\in\mathbb N\). Write \(e_j\in\mathbb R^p\) for the vector whose \(j\)th coordinate is \(1\) and whose other coordinates are \(0\). For a function \(f:\mathbb R^p\to\mathbb R\) the gradient \(\nabla f(x)\) is the vector representing the derivative, \(Df(x)v=\langle\nabla f(x),v\rangle\) for all \(v\), and \(\partial_jf(x)\) is its \(j\)th coordinate. For \(x\in\mathbb R^p\) and a real number \(a\ne0\) the coordinate Taylor quotient is


\[Q_j(x,a)=\frac{2\bigl(f(x+ae_j)-f(x)-a\,\partial_jf(x)\bigr)}{a^2},\]


that is, twice the error of the first-order approximation of \(f\) along the \(j\)th coordinate line, divided by the squared displacement. Let \(\mathcal Q_j=\{Q_j(x,a):x\in\mathbb R^p,\ a\in\mathbb R,\ a\ne0\}\). The coordinate smoothness constant is \(M_j=\sup\mathcal Q_j\), and for \(p\ge1\) the average smoothness constant is \(M_{\rm av}=\frac1p\sum_{j=1}^pM_j\). The supremum of a nonempty set of real numbers is a real number exactly when the set is bounded above; Step 3 shows that \(\mathcal Q_j\) is bounded above whenever \(\nabla f\) is Lipschitz, and \(M_j\) is used only under that hypothesis.




\textbf{Claim.} For every \(f\), every \(j\in\{1,\dots,p\}\) and every \(x\in\mathbb R^p\):



\begin{enumerate}


\item the set \(\mathcal Q_j\) is nonempty;

\item \(\|e_j\|=1\);

\item \(\langle x,e_j\rangle=x_j\). In particular \(\langle\nabla f(x),e_j\rangle=\partial_jf(x)\).

\end{enumerate}



\textbf{Proof.} (1) The number \(Q_j(0,1)\) is obtained from the admissible choice \(x=0\), \(a=1\ne0\), so it belongs to \(\mathcal Q_j\). (2) The only nonzero coordinate of \(e_j\) equals \(1\), so \(\|e_j\|^2=\sum_i(e_j)_i^2=1\) and \(\|e_j\|=1\). (3) By definition of the inner product, \(\langle x,e_j\rangle=\sum_ix_i(e_j)_i=x_j\), since \((e_j)_i=0\) for \(i\ne j\) and \((e_j)_j=1\). Applying this with \(x\) replaced by \(\nabla f(x)\) gives the last assertion. \(\square\)




Identities (2) and (3) are used to evaluate the inner product and the norm of a coordinate displacement \(ae_j\): \(\langle v,ae_j\rangle=a\,v_j\) and \(\|ae_j\|^2=a^2\).

\begin{leanbox}{Lean: \texttt{State}, \texttt{coordinateVector}, \texttt{coordinateTaylorQuotient} and 6 more}
\leansource{\texttt{State} \quad AverageLMC/\allowbreak{}Assumptions.lean, lines 26–26}
\begin{leancode}
\leanline{0}{abbrev\ State\ (p\ :\ ℕ)\ :=\ EuclideanSpace\ ℝ\ (Fin\ p)}
\end{leancode}
\leansource{\texttt{coordinateVector} \quad AverageLMC/\allowbreak{}Assumptions.lean, lines 29–29}
\begin{leancode}
\leanline{0}{def\ coordinateVector\ \{p\ :\ ℕ\}\ (j\ :\ Fin\ p)\ :\ State\ p\ :=\ EuclideanSpace.single\ j\ 1}
\end{leancode}
\leansource{\texttt{coordinateTaylorQuotient} \quad AverageLMC/\allowbreak{}Assumptions.lean, lines 32–34}
\begin{leancode}
\leanline{0}{def\ coordinateTaylorQuotient\ \{p\ :\ ℕ\}\ (f\ :\ State\ p\ →\ ℝ)\ (j\ :\ Fin\ p)}
\leanline{4}{(x\ :\ State\ p)\ (a\ :\ ℝ)\ :\ ℝ\ :=}
\leanline{2}{2\ *\ (f\ (x\ +\ a\ •\ coordinateVector\ j)\ -\ f\ x\ -\ a\ *\ (gradient\ f\ x)\ j)\ /\ a\ \textasciicircum{}\ 2}
\end{leancode}
\leansource{\texttt{coordinateTaylorQuotients} \quad AverageLMC/\allowbreak{}Assumptions.lean, lines 37–38}
\begin{leancode}
\leanline{0}{def\ coordinateTaylorQuotients\ \{p\ :\ ℕ\}\ (f\ :\ State\ p\ →\ ℝ)\ (j\ :\ Fin\ p)\ :\ Set\ ℝ\ :=}
\leanline{2}{\{q\ |\ ∃\ x\ a,\ a\ ≠\ 0\ ∧\ q\ =\ coordinateTaylorQuotient\ f\ j\ x\ a\}}
\end{leancode}
\leansource{\texttt{coordinateSmoothness} \quad AverageLMC/\allowbreak{}Assumptions.lean, lines 42–43}
\begin{leancode}
\leanline{0}{def\ coordinateSmoothness\ \{p\ :\ ℕ\}\ (f\ :\ State\ p\ →\ ℝ)\ (j\ :\ Fin\ p)\ :\ ℝ\ :=}
\leanline{2}{sSup\ (coordinateTaylorQuotients\ f\ j)}
\end{leancode}
\leansource{\texttt{averageSmoothness} \quad AverageLMC/\allowbreak{}Assumptions.lean, lines 46–47}
\begin{leancode}
\leanline{0}{def\ averageSmoothness\ \{p\ :\ ℕ\}\ (f\ :\ State\ p\ →\ ℝ)\ :\ ℝ\ :=}
\leanline{2}{(∑\ j,\ coordinateSmoothness\ f\ j)\ /\ (p\ :\ ℝ)}
\end{leancode}
\leansource{\texttt{coordinateTaylorQuotients\_\allowbreak{}nonempty} \quad AverageLMC/\allowbreak{}Assumptions.lean, lines 72–74}
\begin{leancode}
\leanline{0}{theorem\ coordinateTaylorQuotients\_nonempty\ \{p\ :\ ℕ\}\ (f\ :\ State\ p\ →\ ℝ)\ (j\ :\ Fin\ p)\ :}
\leanline{4}{(coordinateTaylorQuotients\ f\ j).Nonempty}
\end{leancode}
\leansource{\texttt{norm\_\allowbreak{}coordinateVector} \quad AverageLMC/\allowbreak{}Assumptions.lean, lines 76–78}
\begin{leancode}
\leanline{0}{@[simp]\ theorem\ norm\_coordinateVector\ \{p\ :\ ℕ\}\ (j\ :\ Fin\ p)\ :}
\leanline{4}{‖coordinateVector\ j‖\ =\ 1}
\end{leancode}
\leansource{\texttt{inner\_\allowbreak{}coordinateVector} \quad AverageLMC/\allowbreak{}Assumptions.lean, lines 80–82}
\begin{leancode}
\leanline{0}{@[simp]\ theorem\ inner\_coordinateVector\ \{p\ :\ ℕ\}\ (x\ :\ State\ p)\ (j\ :\ Fin\ p)\ :}
\leanline{4}{inner\ ℝ\ x\ (coordinateVector\ j)\ =\ x\ j}
\end{leancode}
\leanlabel{How the Lean proof follows the argument}
The state space is \texttt{State p}, an abbreviation for \texttt{Euclidean\allowbreak{}Space \ensuremath{\mathbb R} (Fin p)}: real functions on the index type \texttt{Fin p} \(=\{0,\dots,p-1\}\), with the Euclidean inner product and norm. Coordinates are therefore indexed from \(0\) to \(p-1\) in the source and from \(1\) to \(p\) in the text. The definitions match the mathematics as follows.



\begin{enumerate}


\item \texttt{coordinate\allowbreak{}Vector j} is \texttt{Euclidean\allowbreak{}Space.\allowbreak{}single j 1}, the vector \(e_j\).

\item \texttt{coordinate\allowbreak{}Taylor\allowbreak{}Quotient f j x a} is the expression \(2(f(x+ae_j)-f(x)-a\,\partial_jf(x))/a^2\). Here \texttt{gradient f x} is the Mathlib gradient, the vector associated with the Fréchet derivative by the inner product, and \texttt{(gradient f x) j} is its coordinate \(\partial_jf(x)\). (Mathlib assigns the derivative \(0\) at points where \(f\) is not differentiable; all results that use the quotient assume differentiability.)

\item \texttt{coordinate\allowbreak{}Taylor\allowbreak{}Quotients f j} is the set \(\mathcal Q_j\), written as \texttt{\{q | \ensuremath{\exists} x a, a ≠ 0 ∧ q = coordinate\allowbreak{}Taylor\allowbreak{}Quotient f j x a\}}.

\item \texttt{coordinate\allowbreak{}Smoothness f j} is \texttt{s\allowbreak{}Sup} of this set. Mathlib's \texttt{s\allowbreak{}Sup} on \(\mathbb R\) is the least upper bound for nonempty sets that are bounded above and is \(0\) by convention otherwise; the later lemmas establish boundedness before using the value, so the definition agrees with \(M_j\) wherever it is used.

\item \texttt{average\allowbreak{}Smoothness f} is \texttt{(∑ j, coordinate\allowbreak{}Smoothness f j) /\allowbreak{} (p : \ensuremath{\mathbb R})}, the average \(M_{\rm av}\); \texttt{(p : \ensuremath{\mathbb R})} is the dimension cast to a real number.

\end{enumerate}



\texttt{coordinate\allowbreak{}Taylor\allowbreak{}Quotients\_\allowbreak{}nonempty} is proved by the explicit element \texttt{\ensuremath{\langle}coordinate\allowbreak{}Taylor\allowbreak{}Quotient f j 0 1, 0, 1, one\_\allowbreak{}ne\_\allowbreak{}zero, rfl\ensuremath{\rangle}}: the quotient \(Q_j(0,1)\), the witnesses \(x=0\), \(a=1\), the proof \texttt{one\_\allowbreak{}ne\_\allowbreak{}zero} of \(1\ne0\), and \texttt{rfl} for the defining equality. \texttt{norm\_\allowbreak{}coordinate\allowbreak{}Vector} and \texttt{inner\_\allowbreak{}coordinate\allowbreak{}Vector} are claims (2) and (3); both are closed by \texttt{simp} after unfolding \texttt{coordinate\allowbreak{}Vector}, the second with the Mathlib identity \texttt{Euclidean\allowbreak{}Space.\allowbreak{}inner\_\allowbreak{}single\_\allowbreak{}right}. They carry the \texttt{@[simp]} attribute, so later simplifications rewrite \(\|e_j\|\) to \(1\) and \(\langle x,e_j\rangle\) to \(x_j\) automatically.
\end{leanbox}
\subsection{\texorpdfstring{The Lipschitz seminorm of the gradient}{The Lipschitz seminorm of the gradient}}
\label{report-step-AverageLMC-Assumptions-lean-2}


\textbf{Definition.} For \(f:\mathbb R^p\to\mathbb R\) let \(\mathcal G=\bigl\{\|\nabla f(x)-\nabla f(y)\|/\|x-y\|:x,y\in\mathbb R^p,\ x\ne y\bigr\}\) and \(M_\infty=\sup\mathcal G\), the Lipschitz seminorm of the gradient. A number \(M\ge0\) is a Lipschitz constant of \(\nabla f\) if \(\|\nabla f(x)-\nabla f(y)\|\le M\|x-y\|\) for all \(x,y\).




\textbf{Claim.} Let \(p\ge1\), and let \(\nabla f\) be Lipschitz with some constant \(M\ge0\). Then



\begin{enumerate}


\item \(M_\infty\le M\);

\item \(M_\infty\ge0\);

\item \(\nabla f\) is Lipschitz with constant \(M_\infty\): \(\|\nabla f(x)-\nabla f(y)\|\le M_\infty\|x-y\|\) for all \(x,y\in\mathbb R^p\).

\end{enumerate}



Thus \(M_\infty\) is the smallest Lipschitz constant of \(\nabla f\), and later estimates may be stated with this exact constant.




\textbf{Proof.} \emph{The set \(\mathcal G\) is nonempty.} Since \(p\ge1\), the vector \(e_1\) exists, and \(e_1\ne0\) because \(\|e_1\|=1\ne0=\|0\|\). Hence the quotient at the pair \((0,e_1)\) belongs to \(\mathcal G\).




\emph{\(M\) is an upper bound of \(\mathcal G\).} If \(x\ne y\), then \(\|x-y\|>0\), and dividing the Lipschitz inequality \(\|\nabla f(x)-\nabla f(y)\|\le M\|x-y\|\) by \(\|x-y\|\) gives \(\|\nabla f(x)-\nabla f(y)\|/\|x-y\|\le M\).




(1) A nonempty set of reals with an upper bound has a least upper bound, and the least upper bound is at most every upper bound; hence \(M_\infty\le M\).




(2) Every element of \(\mathcal G\) is a quotient of two nonnegative numbers, hence nonnegative. The element at \((0,e_1)\) is therefore \(\ge0\), and it is \(\le M_\infty\) because \(M_\infty\) is an upper bound of \(\mathcal G\).




(3) If \(x=y\), both sides of the inequality are \(0\). If \(x\ne y\), the quotient at \((x,y)\) belongs to \(\mathcal G\), so \(\|\nabla f(x)-\nabla f(y)\|/\|x-y\|\le M_\infty\); multiplying by \(\|x-y\|>0\) gives the inequality. \(\square\)




The hypothesis \(p\ge1\) is used only to produce an element of \(\mathcal G\): for \(p=0\) the space is \(\{0\}\), \(\mathcal G\) is empty and its supremum is not defined.

\begin{leanbox}{Lean: \texttt{globalSmoothness}, \texttt{globalSmoothness\_\allowbreak{}le\_\allowbreak{}bound}, \texttt{globalSmoothness\_\allowbreak{}nonneg} and 1 more}
\leansource{\texttt{globalSmoothness} \quad AverageLMC/\allowbreak{}Assumptions.lean, lines 50–52}
\begin{leancode}
\leanline{0}{def\ globalSmoothness\ \{p\ :\ ℕ\}\ (f\ :\ State\ p\ →\ ℝ)\ :\ ℝ\ :=}
\leanline{2}{sSup\ \{q\ :\ ℝ\ |\ ∃\ x\ y\ :\ State\ p,\ x\ ≠\ y\ ∧}
\leanline{4}{q\ =\ ‖gradient\ f\ x\ -\ gradient\ f\ y‖\ /\ ‖x\ -\ y‖\}}
\end{leancode}
\leansource{\texttt{globalSmoothness\_\allowbreak{}le\_\allowbreak{}bound} \quad AverageLMC/\allowbreak{}Assumptions.lean, lines 86–100}
\begin{leancode}
\leanline{0}{theorem\ globalSmoothness\_le\_bound\ \{p\ :\ ℕ\}\ (hp\ :\ 0\ <\ p)}
\leanline{4}{(f\ :\ State\ p\ →\ ℝ)\ (M\ :\ ℝ≥0)\ (hLip\ :\ LipschitzWith\ M\ (gradient\ f))\ :}
\leanline{4}{globalSmoothness\ f\ ≤\ (M\ :\ ℝ)}
\end{leancode}
\leansource{\texttt{globalSmoothness\_\allowbreak{}nonneg} \quad AverageLMC/\allowbreak{}Assumptions.lean, lines 102–120}
\begin{leancode}
\leanline{0}{theorem\ globalSmoothness\_nonneg\ \{p\ :\ ℕ\}\ (hp\ :\ 0\ <\ p)}
\leanline{4}{(f\ :\ State\ p\ →\ ℝ)\ (M\ :\ ℝ≥0)\ (hLip\ :\ LipschitzWith\ M\ (gradient\ f))\ :}
\leanline{4}{0\ ≤\ globalSmoothness\ f}
\end{leancode}
\leansource{\texttt{gradient\_\allowbreak{}lipschitz\_\allowbreak{}globalSmoothness} \quad AverageLMC/\allowbreak{}Assumptions.lean, lines 124–140}
\begin{leancode}
\leanline{0}{theorem\ gradient\_lipschitz\_globalSmoothness\ \{p\ :\ ℕ\}\ (hp\ :\ 0\ <\ p)}
\leanline{4}{(f\ :\ State\ p\ →\ ℝ)\ (M\ :\ ℝ≥0)\ (hLip\ :\ LipschitzWith\ M\ (gradient\ f))\ :}
\leanline{4}{LipschitzWith\ ⟨globalSmoothness\ f,\ globalSmoothness\_nonneg\ hp\ f\ M\ hLip⟩}
\leanline{6}{(gradient\ f)}
\end{leancode}
\leanlabel{How the Lean proof follows the argument}
\texttt{global\allowbreak{}Smoothness f} is \texttt{s\allowbreak{}Sup} of the set \texttt{\{q : \ensuremath{\mathbb R} | \ensuremath{\exists} x y : State p, x ≠ y ∧ q = \ensuremath{\Vert}gradient f x -\allowbreak{} gradient f y\ensuremath{\Vert} /\allowbreak{} \ensuremath{\Vert}x -\allowbreak{} y\ensuremath{\Vert}\}}, that is \(M_\infty=\sup\mathcal G\). The three theorems take \texttt{hp : 0 < p} and a Lipschitz hypothesis \texttt{h\allowbreak{}Lip : Lipschitz\allowbreak{}With M (gradient f)}, where \texttt{M : \ensuremath{\mathbb R}\ensuremath{\ge}0} is a nonnegative real; \texttt{h\allowbreak{}Lip.\allowbreak{}norm\_\allowbreak{}sub\_\allowbreak{}le x y} is the inequality \(\|\nabla f(x)-\nabla f(y)\|\le M\|x-y\|\).



\begin{enumerate}


\item In \texttt{global\allowbreak{}Smoothness\_\allowbreak{}le\_\allowbreak{}bound}, \texttt{j : Fin p := \ensuremath{\langle}0, hp\ensuremath{\rangle}} is the first coordinate and \texttt{he} proves \(0\ne e_1\) by applying the norm to a hypothetical equality (\texttt{congr\allowbreak{}Arg norm h}) and simplifying. \texttt{cs\allowbreak{}Sup\_\allowbreak{}le} reduces claim (1) to nonemptiness, witnessed by the pair \((0,e_1)\), and to the bound \(q\le M\) for each element. For an element with \(x\ne y\), \texttt{div\_\allowbreak{}le\_\allowbreak{}iff\ensuremath{{}_{0}}} with the positivity proof \texttt{norm\_\allowbreak{}pos\_\allowbreak{}iff.\allowbreak{}mpr (sub\_\allowbreak{}ne\_\allowbreak{}zero.\allowbreak{}mpr hxy)} of \(\|x-y\|>0\) turns the quotient bound into the Lipschitz inequality.

\item \texttt{global\allowbreak{}Smoothness\_\allowbreak{}nonneg} first proves \texttt{hb}, that the set is bounded above by \(M\), then chains \texttt{div\_\allowbreak{}nonneg} (the quotient at \((0,e_1)\) is nonnegative) with \texttt{le\_\allowbreak{}cs\allowbreak{}Sup hb} (that quotient is at most the supremum). This is claim (2).

\item \texttt{gradient\_\allowbreak{}lipschitz\_\allowbreak{}global\allowbreak{}Smoothness} concludes \texttt{Lipschitz\allowbreak{}With \ensuremath{\langle}global\allowbreak{}Smoothness f, global\allowbreak{}Smoothness\_\allowbreak{}nonneg hp f M h\allowbreak{}Lip\ensuremath{\rangle} (gradient f)}: the constant \(M_\infty\) is packaged as an element of \texttt{\ensuremath{\mathbb R}\ensuremath{\ge}0} together with the proof of its nonnegativity. \texttt{Lipschitz\allowbreak{}With.\allowbreak{}of\_\allowbreak{}dist\_\allowbreak{}le\_\allowbreak{}mul} reduces the goal to \(\operatorname{dist}(\nabla f(x),\nabla f(y))\le M_\infty\operatorname{dist}(x,y)\), \texttt{dist\_\allowbreak{}eq\_\allowbreak{}norm} rewrites distances as norms, and \texttt{by\_\allowbreak{}cases hxy : x = y} separates the trivial case from the case \(x\ne y\), where \texttt{(div\_\allowbreak{}le\_\allowbreak{}iff\ensuremath{{}_{0}} …).\allowbreak{}mp} applied to \texttt{le\_\allowbreak{}cs\allowbreak{}Sup hb \ensuremath{\langle}x, y, hxy, rfl\ensuremath{\rangle}} gives claim (3).

\end{enumerate}
\end{leanbox}
\subsection{\texorpdfstring{The coordinate quadratic upper model}{The coordinate quadratic upper model}}
\label{report-step-AverageLMC-Assumptions-lean-3}


\textbf{Claim.} Let \(f:\mathbb R^p\to\mathbb R\) be differentiable and let \(\nabla f\) be Lipschitz with constant \(M\ge0\). Then for every \(j\):



\begin{enumerate}


\item \(Q_j(x,a)\le M\) for every \(x\in\mathbb R^p\) and every \(a\ne0\);

\item \(\mathcal Q_j\) is bounded above, so \(M_j=\sup\mathcal Q_j\) is a real number, and \(M_j\le M\);

\item for every \(x\in\mathbb R^p\) and every \(a\in\mathbb R\), including \(a=0\),
\[f(x+ae_j)\le f(x)+a\,\partial_jf(x)+\frac{M_j}{2}\,a^2.\qquad(1)\]

\end{enumerate}



\textbf{Proof.} We use the descent inequality: for all \(x,y\in\mathbb R^p\),


\[f(y)\le f(x)+\langle\nabla f(x),y-x\rangle+\frac M2\|y-x\|^2.\qquad(2)\]


To recall its proof, put \(d=y-x\) and \(g(t)=f(x+td)\) for \(t\in[0,1]\). By the chain rule \(g'(t)=\langle\nabla f(x+td),d\rangle\), which is continuous in \(t\) because \(\nabla f\) is Lipschitz. By the fundamental theorem of calculus, the Cauchy–Schwarz inequality and the Lipschitz bound,


\[f(y)-f(x)-\langle\nabla f(x),d\rangle=\int_0^1\langle\nabla f(x+td)-\nabla f(x),d\rangle\,dt\le\int_0^1Mt\|d\|^2\,dt=\frac M2\|d\|^2.\]


(1) Fix \(x\) and \(a\ne0\), and apply (2) with \(y=x+ae_j\). Then \(y-x=ae_j\), \(\langle\nabla f(x),ae_j\rangle=a\,\partial_jf(x)\) and \(\|ae_j\|^2=a^2\|e_j\|^2=a^2\) by Step 1. Hence


\[f(x+ae_j)-f(x)-a\,\partial_jf(x)\le\frac M2\,a^2.\]


Multiplying by \(2\) and dividing by \(a^2>0\) gives \(Q_j(x,a)\le M\).




(2) By (1), \(M\) is an upper bound of \(\mathcal Q_j\). The set is nonempty by Step 1, so its supremum \(M_j\) exists and is at most the upper bound \(M\).




(3) If \(a=0\), both sides of (1) equal \(f(x)\). If \(a\ne0\), then \(Q_j(x,a)\in\mathcal Q_j\), so \(Q_j(x,a)\le M_j\) by definition of the supremum (which exists by (2)). Multiplying this inequality by \(a^2/2>0\) gives \(f(x+ae_j)-f(x)-a\,\partial_jf(x)\le\frac{M_j}2a^2\), which is (1). \(\square\)




Inequality (1) has the coefficient \(M_j/2\) in place of the coefficient \(M/2\) given by (2) along the line \(x+\mathbb Re_j\); the two can differ substantially when the potential is much stiffer in some directions than along the coordinate axes. Only differentiability and a Lipschitz gradient are used; the existence of second derivatives is not needed.

\begin{leanbox}{Lean: \texttt{coordinateTaylorQuotient\_\allowbreak{}le}, \texttt{coordinateTaylorQuotients\_\allowbreak{}bddAbove}, \texttt{coordinateSmoothness\_\allowbreak{}le\_\allowbreak{}global\_\allowbreak{}bound} and 1 more}
\leansource{\texttt{coordinateTaylorQuotient\_\allowbreak{}le} \quad AverageLMC/\allowbreak{}Assumptions.lean, lines 142–153}
\begin{leancode}
\leanline{0}{theorem\ coordinateTaylorQuotient\_le\ \{p\ :\ ℕ\}\ (f\ :\ State\ p\ →\ ℝ)\ (M\ :\ ℝ≥0)}
\leanline{4}{(hdiff\ :\ Differentiable\ ℝ\ f)\ (hLip\ :\ LipschitzWith\ M\ (gradient\ f))}
\leanline{4}{(j\ :\ Fin\ p)\ (x\ :\ State\ p)\ (a\ :\ ℝ)\ (ha\ :\ a\ ≠\ 0)\ :}
\leanline{4}{coordinateTaylorQuotient\ f\ j\ x\ a\ ≤\ (M\ :\ ℝ)}
\end{leancode}
\leansource{\texttt{coordinateTaylorQuotients\_\allowbreak{}bddAbove} \quad AverageLMC/\allowbreak{}Assumptions.lean, lines 155–160}
\begin{leancode}
\leanline{0}{theorem\ coordinateTaylorQuotients\_bddAbove\ \{p\ :\ ℕ\}\ (f\ :\ State\ p\ →\ ℝ)\ (M\ :\ ℝ≥0)}
\leanline{4}{(hdiff\ :\ Differentiable\ ℝ\ f)\ (hLip\ :\ LipschitzWith\ M\ (gradient\ f))\ (j\ :\ Fin\ p)\ :}
\leanline{4}{BddAbove\ (coordinateTaylorQuotients\ f\ j)}
\end{leancode}
\leansource{\texttt{coordinateSmoothness\_\allowbreak{}le\_\allowbreak{}global\_\allowbreak{}bound} \quad AverageLMC/\allowbreak{}Assumptions.lean, lines 162–167}
\begin{leancode}
\leanline{0}{theorem\ coordinateSmoothness\_le\_global\_bound\ \{p\ :\ ℕ\}\ (f\ :\ State\ p\ →\ ℝ)\ (M\ :\ ℝ≥0)}
\leanline{4}{(hdiff\ :\ Differentiable\ ℝ\ f)\ (hLip\ :\ LipschitzWith\ M\ (gradient\ f))\ (j\ :\ Fin\ p)\ :}
\leanline{4}{coordinateSmoothness\ f\ j\ ≤\ (M\ :\ ℝ)}
\end{leancode}
\leansource{\texttt{coordinate\_\allowbreak{}upper\_\allowbreak{}bound} \quad AverageLMC/\allowbreak{}Assumptions.lean, lines 171–184}
\begin{leancode}
\leanline{0}{theorem\ coordinate\_upper\_bound\ \{p\ :\ ℕ\}\ (f\ :\ State\ p\ →\ ℝ)\ (M\ :\ ℝ≥0)}
\leanline{4}{(hdiff\ :\ Differentiable\ ℝ\ f)\ (hLip\ :\ LipschitzWith\ M\ (gradient\ f))}
\leanline{4}{(j\ :\ Fin\ p)\ (x\ :\ State\ p)\ (a\ :\ ℝ)\ :}
\leanline{4}{f\ (x\ +\ a\ •\ coordinateVector\ j)\ ≤\ f\ x\ +\ a\ *\ (gradient\ f\ x)\ j\ +}
\leanline{6}{coordinateSmoothness\ f\ j\ /\ 2\ *\ a\ \textasciicircum{}\ 2}
\end{leancode}
\leanlabel{How the Lean proof follows the argument}
\begin{enumerate}


\item \texttt{coordinate\allowbreak{}Taylor\allowbreak{}Quotient\_\allowbreak{}le} takes \texttt{hdiff : Differentiable \ensuremath{\mathbb R} f}, \texttt{h\allowbreak{}Lip} and \texttt{ha : a ≠ 0}. The descent inequality (2) is \texttt{ULA.\allowbreak{}Analysis.\allowbreak{}smooth\_\allowbreak{}upper\_\allowbreak{}bound}, proved in the unadjusted Langevin library and applied here with \texttt{y = x + a \ensuremath{\bullet} coordinate\allowbreak{}Vector j} to obtain \texttt{hu}. The \texttt{simp only} call with \texttt{add\_\allowbreak{}sub\_\allowbreak{}cancel\_\allowbreak{}left}, \texttt{inner\_\allowbreak{}smul\_\allowbreak{}right}, \texttt{inner\_\allowbreak{}coordinate\allowbreak{}Vector}, \texttt{norm\_\allowbreak{}smul}, \texttt{Real.\allowbreak{}norm\_\allowbreak{}eq\_\allowbreak{}abs}, \texttt{norm\_\allowbreak{}coordinate\allowbreak{}Vector} and \texttt{mul\_\allowbreak{}one} rewrites \(y-x\) as \(ae_j\), the inner product as \(a\,\partial_jf(x)\) and the norm as \(|a|\); \texttt{sq\_\allowbreak{}abs} replaces \(|a|^2\) by \(a^2\). After unfolding the quotient, \texttt{div\_\allowbreak{}le\_\allowbreak{}iff\ensuremath{{}_{0}} (sq\_\allowbreak{}pos\_\allowbreak{}of\_\allowbreak{}ne\_\allowbreak{}zero ha)} turns the division by \(a^2>0\) into a multiplication and \texttt{nlinarith} concludes.

\item \texttt{coordinate\allowbreak{}Taylor\allowbreak{}Quotients\_\allowbreak{}bdd\allowbreak{}Above} exhibits \(M\) as an upper bound (the anonymous constructor \texttt{\ensuremath{\langle}(M : \ensuremath{\mathbb R}), ?\_\allowbreak{}\ensuremath{\rangle}}) and checks it on each element with the previous lemma. \texttt{coordinate\allowbreak{}Smoothness\_\allowbreak{}le\_\allowbreak{}global\_\allowbreak{}bound} applies \texttt{cs\allowbreak{}Sup\_\allowbreak{}le} with the nonemptiness lemma of Step 1; this is claim (2).

\item \texttt{coordinate\_\allowbreak{}upper\_\allowbreak{}bound} is (1). After \texttt{by\_\allowbreak{}cases ha : a = 0}, the case \(a=0\) is closed by \texttt{simp}. Otherwise \texttt{hs} is \texttt{le\_\allowbreak{}cs\allowbreak{}Sup} applied to the bounded set and to the membership of \(Q_j(x,a)\); \texttt{change} restates it as \(Q_j(x,a)\le M_j\), \texttt{div\_\allowbreak{}le\_\allowbreak{}iff\ensuremath{{}_{0}} (sq\_\allowbreak{}pos\_\allowbreak{}of\_\allowbreak{}ne\_\allowbreak{}zero ha)} clears the denominator, and \texttt{nlinarith} rearranges.

\end{enumerate}
\end{leanbox}


Earlier results used here: \cite[Section~3.9.1]{LangevinLean2026ULASC}.
\subsection{\texorpdfstring{Comparison of strong convexity, average smoothness and global smoothness}{Comparison of strong convexity, average smoothness and global smoothness}}
\label{report-step-AverageLMC-Assumptions-lean-4}


\textbf{Claim.} Let \(f:\mathbb R^p\to\mathbb R\) be differentiable with \(M\)-Lipschitz gradient, \(M\ge0\).



\begin{enumerate}


\item If \(f\) is \(m\)-strongly convex for some \(m\in\mathbb R\), then \(m\le M_j\) for every \(j\).

\item If \(p\ge1\), then \(M_{\rm av}\le M\).

\item If \(p\ge1\), then \(M_j\le M_\infty\) for every \(j\), and \(M_{\rm av}\le M_\infty\).

\item If \(p\ge1\) and \(f\) is \(m\)-strongly convex, then \(m\le M_{\rm av}\).

\end{enumerate}



Consequently, for \(p\ge1\) and an \(m\)-strongly convex \(f\), \(m\le M_{\rm av}\le M_\infty\le M\).




Here \(f\) is \(m\)-strongly convex if \(f((1-t)x+ty)\le(1-t)f(x)+tf(y)-\frac m2t(1-t)\|x-y\|^2\) for all \(x,y\in\mathbb R^p\) and \(t\in[0,1]\).




\textbf{Proof.} (1) For a differentiable \(m\)-strongly convex function the supporting-plane inequality holds:


\[f(y)\ge f(x)+\langle\nabla f(x),y-x\rangle+\frac m2\|y-x\|^2\qquad\text{for all }x,y.\qquad(3)\]


To recall why, note that \(g(v)=f(v)-\frac m2\|v\|^2\) is convex: inserting the identity \(\|(1-t)x+ty\|^2=(1-t)\|x\|^2+t\|y\|^2-t(1-t)\|x-y\|^2\) into the definition of strong convexity gives exactly the convexity inequality for \(g\). A differentiable convex function lies above its tangent planes, \(g(y)\ge g(x)+Dg(x)(y-x)\), and \(Dg(x)v=\langle\nabla f(x),v\rangle-m\langle x,v\rangle\). Since \(\|y\|^2-\|x\|^2-2\langle x,y-x\rangle=\|y-x\|^2\), this tangent-plane inequality is (3).




Apply (3) with \(x=0\) and \(y=e_j\), and apply the coordinate upper model (1) of Step 3 with \(x=0\) and \(a=1\). Using \(\langle\nabla f(0),e_j\rangle=\partial_jf(0)\) and \(\|e_j\|=1\),


\[f(0)+\partial_jf(0)+\frac m2\le f(e_j)\le f(0)+\partial_jf(0)+\frac{M_j}2.\]


Cancelling \(f(0)+\partial_jf(0)\) and multiplying by \(2\) gives \(m\le M_j\).




(2) By Step 3, \(M_j\le M\) for each of the \(p\) indices. Summing, \(\sum_jM_j\le pM\), and dividing by \(p>0\) gives \(M_{\rm av}\le M\).




(3) By Step 2, \(M_\infty\ge0\) and \(\nabla f\) is \(M_\infty\)-Lipschitz. Applying Step 3 and part (2) with the Lipschitz constant \(M_\infty\) in place of \(M\) gives \(M_j\le M_\infty\) and \(M_{\rm av}\le M_\infty\).




(4) Summing (1) over \(j\) gives \(pm\le\sum_jM_j\), and dividing by \(p>0\) gives \(m\le M_{\rm av}\). The final chain combines (4), (3) and Step 2. \(\square\)




The constants \(M_j\), \(M_{\rm av}\) and \(M_\infty\) in these inequalities are the suprema defined from \(f\), not arbitrary upper bounds; this is what allows the later theorems to be stated with the constants of the paper. The inequality \(m\le M_{\rm av}\) is used, for example, to guarantee \(mh\le1\) when \(hM_\infty\le1\).

\begin{leanbox}{Lean: \texttt{strongConvex\_\allowbreak{}le\_\allowbreak{}coordinateSmoothness}, \texttt{averageSmoothness\_\allowbreak{}le\_\allowbreak{}global\_\allowbreak{}bound}, \texttt{coordinateSmoothness\_\allowbreak{}le\_\allowbreak{}globalSmoothness} and 2 more}
\leansource{\texttt{strongConvex\_\allowbreak{}le\_\allowbreak{}coordinateSmoothness} \quad AverageLMC/\allowbreak{}Assumptions.lean, lines 186–195}
\begin{leancode}
\leanline{0}{theorem\ strongConvex\_le\_coordinateSmoothness\ \{p\ :\ ℕ\}\ (f\ :\ State\ p\ →\ ℝ)}
\leanline{4}{(m\ :\ ℝ)\ (M\ :\ ℝ≥0)\ (hdiff\ :\ Differentiable\ ℝ\ f)}
\leanline{4}{(hsc\ :\ StrongConvexOn\ Set.univ\ m\ f)\ (hLip\ :\ LipschitzWith\ M\ (gradient\ f))}
\leanline{4}{(j\ :\ Fin\ p)\ :\ m\ ≤\ coordinateSmoothness\ f\ j}
\end{leancode}
\leansource{\texttt{averageSmoothness\_\allowbreak{}le\_\allowbreak{}global\_\allowbreak{}bound} \quad AverageLMC/\allowbreak{}Assumptions.lean, lines 197–203}
\begin{leancode}
\leanline{0}{theorem\ averageSmoothness\_le\_global\_bound\ \{p\ :\ ℕ\}\ (hp\ :\ 0\ <\ p)}
\leanline{4}{(f\ :\ State\ p\ →\ ℝ)\ (M\ :\ ℝ≥0)\ (hdiff\ :\ Differentiable\ ℝ\ f)}
\leanline{4}{(hLip\ :\ LipschitzWith\ M\ (gradient\ f))\ :\ averageSmoothness\ f\ ≤\ (M\ :\ ℝ)}
\end{leancode}
\leansource{\texttt{coordinateSmoothness\_\allowbreak{}le\_\allowbreak{}globalSmoothness} \quad AverageLMC/\allowbreak{}Assumptions.lean, lines 206–212}
\begin{leancode}
\leanline{0}{theorem\ coordinateSmoothness\_le\_globalSmoothness\ \{p\ :\ ℕ\}\ (hp\ :\ 0\ <\ p)}
\leanline{4}{(f\ :\ State\ p\ →\ ℝ)\ (M\ :\ ℝ≥0)\ (hdiff\ :\ Differentiable\ ℝ\ f)}
\leanline{4}{(hLip\ :\ LipschitzWith\ M\ (gradient\ f))\ (j\ :\ Fin\ p)\ :}
\leanline{4}{coordinateSmoothness\ f\ j\ ≤\ globalSmoothness\ f}
\end{leancode}
\leansource{\texttt{averageSmoothness\_\allowbreak{}le\_\allowbreak{}globalSmoothness} \quad AverageLMC/\allowbreak{}Assumptions.lean, lines 216–222}
\begin{leancode}
\leanline{0}{theorem\ averageSmoothness\_le\_globalSmoothness\ \{p\ :\ ℕ\}\ (hp\ :\ 0\ <\ p)}
\leanline{4}{(f\ :\ State\ p\ →\ ℝ)\ (M\ :\ ℝ≥0)\ (hdiff\ :\ Differentiable\ ℝ\ f)}
\leanline{4}{(hLip\ :\ LipschitzWith\ M\ (gradient\ f))\ :}
\leanline{4}{averageSmoothness\ f\ ≤\ globalSmoothness\ f}
\end{leancode}
\leansource{\texttt{strongConvex\_\allowbreak{}le\_\allowbreak{}averageSmoothness} \quad AverageLMC/\allowbreak{}Assumptions.lean, lines 224–231}
\begin{leancode}
\leanline{0}{theorem\ strongConvex\_le\_averageSmoothness\ \{p\ :\ ℕ\}\ (hp\ :\ 0\ <\ p)}
\leanline{4}{(f\ :\ State\ p\ →\ ℝ)\ (m\ :\ ℝ)\ (M\ :\ ℝ≥0)\ (hdiff\ :\ Differentiable\ ℝ\ f)}
\leanline{4}{(hsc\ :\ StrongConvexOn\ Set.univ\ m\ f)\ (hLip\ :\ LipschitzWith\ M\ (gradient\ f))\ :}
\leanline{4}{m\ ≤\ averageSmoothness\ f}
\end{leancode}
\leanlabel{How the Lean proof follows the argument}
Strong convexity is the hypothesis \texttt{hsc : Strong\allowbreak{}Convex\allowbreak{}On Set.\allowbreak{}univ m f}. Mathlib defines it as uniform convexity with modulus \(r\mapsto\frac m2r^2\) on the whole space: \(f(ax+by)\le af(x)+bf(y)-ab\frac m2\|x-y\|^2\) whenever \(a,b\ge0\) and \(a+b=1\), which is the definition used above.



\begin{enumerate}


\item \texttt{strong\allowbreak{}Convex\_\allowbreak{}le\_\allowbreak{}coordinate\allowbreak{}Smoothness}: \texttt{hl} is (3) at \(x=0\), \(y=e_j\), supplied by \texttt{ULA.\allowbreak{}strong\_\allowbreak{}first\_\allowbreak{}order hdiff hsc (0 : State p) (coordinate\allowbreak{}Vector j)} from the unadjusted Langevin library; \texttt{hu} is \texttt{coordinate\_\allowbreak{}upper\_\allowbreak{}bound} at \(x=0\), \(a=1\). Two \texttt{simp only} calls evaluate \(\|e_j\|=1\), \(1^2=1\) and \(\langle\nabla f(0),e_j\rangle=\partial_jf(0)\), and \texttt{linarith} combines the two inequalities.

\item \texttt{average\allowbreak{}Smoothness\_\allowbreak{}le\_\allowbreak{}global\_\allowbreak{}bound} unfolds the average, uses \texttt{div\_\allowbreak{}le\_\allowbreak{}iff\ensuremath{{}_{0}}} with \(p>0\) (the cast \texttt{(by exact\_\allowbreak{}mod\_\allowbreak{}cast hp : (0 : \ensuremath{\mathbb R}) < p)}) to reduce to \(\sum_jM_j\le M\,p\), and obtains this from \texttt{Finset.\allowbreak{}sum\_\allowbreak{}le\_\allowbreak{}sum} applied to the bounds \texttt{coordinate\allowbreak{}Smoothness\_\allowbreak{}le\_\allowbreak{}global\_\allowbreak{}bound}; \texttt{simpa [mul\_\allowbreak{}comm]} identifies the sum of the constant \(M\) over the \(p\) indices with \(pM\).

\item \texttt{coordinate\allowbreak{}Smoothness\_\allowbreak{}le\_\allowbreak{}global\allowbreak{}Smoothness} and \texttt{average\allowbreak{}Smoothness\_\allowbreak{}le\_\allowbreak{}global\allowbreak{}Smoothness} are the previous bounds instantiated with the Lipschitz constant \texttt{\ensuremath{\langle}global\allowbreak{}Smoothness f, global\allowbreak{}Smoothness\_\allowbreak{}nonneg hp f M h\allowbreak{}Lip\ensuremath{\rangle}} and the Lipschitz proof \texttt{gradient\_\allowbreak{}lipschitz\_\allowbreak{}global\allowbreak{}Smoothness hp f M h\allowbreak{}Lip} of Step 2.

\item \texttt{strong\allowbreak{}Convex\_\allowbreak{}le\_\allowbreak{}average\allowbreak{}Smoothness} uses \texttt{le\_\allowbreak{}div\_\allowbreak{}iff\ensuremath{{}_{0}}} and \texttt{Finset.\allowbreak{}sum\_\allowbreak{}le\_\allowbreak{}sum} over the bounds of part (1).

\end{enumerate}
\end{leanbox}


Earlier results used here: \cite[Section~3.2.2]{LangevinLean2026ULASC}.
\subsection{\texorpdfstring{Hessian, Laplacian and the constants of the third-order and finite-sum settings}{Hessian, Laplacian and the constants of the third-order and finite-sum settings}}
\label{report-step-AverageLMC-Assumptions-lean-5}


\textbf{Definitions.} Let \(D^2f(x)\) denote the second derivative of \(f\) at \(x\), a bilinear form on \(\mathbb R^p\) (as usual for derivatives defined at every point, it is taken to be \(0\) where it does not exist).



\begin{itemize}


\item The \emph{Hessian matrix} is \(\nabla^2f(x)=\bigl(D^2f(x)(e_i,e_j)\bigr)_{i,j=1}^p\).

\item The \emph{Laplacian} is the trace of \(D^2f(x)\) with respect to the inner product: \(\Delta f(x)=\sum_{i=1}^pD^2f(x)(v_i,v_i)\) for an orthonormal basis \(v_1,\dots,v_p\). The sum does not depend on the orthonormal basis, because a change of orthonormal basis is given by an orthogonal matrix \(O\) and \(\sum_iD^2f(x)(Ov_i,Ov_i)=\operatorname{tr}(O^\top AO)=\operatorname{tr}A\) for the matrix \(A\) of \(D^2f(x)\).

\item For \(M_\Delta\in\mathbb R\), the potential satisfies the \emph{Laplacian smoothness condition} if \(f\) is three times continuously differentiable and \(\|\nabla\Delta f(x)\|^2\le M_\Delta^2\,p\) for every \(x\in\mathbb R^p\). No Lipschitz condition on the Hessian is part of this condition.

\item For a finite family \(f_1,\dots,f_n:\mathbb R^p\to\mathbb R\), let \(M_\infty(f_i)\) be the gradient seminorm of Step 2 for \(f_i\). The root-mean-square component smoothness is \(M_{\rm rms}=\bigl(\frac1n\sum_{i=1}^nM_\infty(f_i)^2\bigr)^{1/2}\).

\end{itemize}



\textbf{Claim.} For every \(f:\mathbb R^p\to\mathbb R\) and every \(x\in\mathbb R^p\), \(\operatorname{tr}\nabla^2f(x)=\Delta f(x)\).




\textbf{Proof.} The canonical basis \(e_1,\dots,e_p\) is orthonormal, so by the definition of the Laplacian with \(v_i=e_i\),


\[\Delta f(x)=\sum_{i=1}^pD^2f(x)(e_i,e_i)=\sum_{i=1}^p\bigl(\nabla^2f(x)\bigr)_{ii}=\operatorname{tr}\nabla^2f(x).\qquad\square\]\begin{leanbox}{Lean: \texttt{hessian}, \texttt{laplacian}, \texttt{LaplacianSmooth} and 2 more}
\leansource{\texttt{hessian} \quad AverageLMC/\allowbreak{}Assumptions.lean, lines 55–57}
\begin{leancode}
\leanline{0}{def\ hessian\ \{p\ :\ ℕ\}\ (f\ :\ State\ p\ →\ ℝ)\ (x\ :\ State\ p)\ :}
\leanline{4}{Matrix\ (Fin\ p)\ (Fin\ p)\ ℝ\ :=}
\leanline{2}{fun\ i\ j\ ↦\ iteratedFDeriv\ ℝ\ 2\ f\ x\ ![coordinateVector\ i,\ coordinateVector\ j]}
\end{leancode}
\leansource{\texttt{laplacian} \quad AverageLMC/\allowbreak{}Assumptions.lean, lines 60–61}
\begin{leancode}
\leanline{0}{def\ laplacian\ \{p\ :\ ℕ\}\ (f\ :\ State\ p\ →\ ℝ)\ :\ State\ p\ →\ ℝ\ :=}
\leanline{2}{Laplacian.laplacian\ f}
\end{leancode}
\leansource{\texttt{LaplacianSmooth} \quad AverageLMC/\allowbreak{}Assumptions.lean, lines 65–66}
\begin{leancode}
\leanline{0}{def\ LaplacianSmooth\ \{p\ :\ ℕ\}\ (f\ :\ State\ p\ →\ ℝ)\ (Mdelta\ :\ ℝ)\ :\ Prop\ :=}
\leanline{2}{ContDiff\ ℝ\ 3\ f\ ∧\ ∀\ x,\ ‖gradient\ (laplacian\ f)\ x‖\ \textasciicircum{}\ 2\ ≤\ Mdelta\ \textasciicircum{}\ 2\ *\ (p\ :\ ℝ)}
\end{leancode}
\leansource{\texttt{rmsComponentSmoothness} \quad AverageLMC/\allowbreak{}Assumptions.lean, lines 69–70}
\begin{leancode}
\leanline{0}{def\ rmsComponentSmoothness\ \{p\ n\ :\ ℕ\}\ (f\ :\ Fin\ n\ →\ State\ p\ →\ ℝ)\ :\ ℝ\ :=}
\leanline{2}{Real.sqrt\ ((∑\ i,\ globalSmoothness\ (f\ i)\ \textasciicircum{}\ 2)\ /\ (n\ :\ ℝ))}
\end{leancode}
\leansource{\texttt{trace\_\allowbreak{}hessian} \quad AverageLMC/\allowbreak{}Assumptions.lean, lines 234–238}
\begin{leancode}
\leanline{0}{theorem\ trace\_hessian\ \{p\ :\ ℕ\}\ (f\ :\ State\ p\ →\ ℝ)\ (x\ :\ State\ p)\ :}
\leanline{4}{Matrix.trace\ (hessian\ f\ x)\ =\ laplacian\ f\ x}
\end{leancode}
\leanlabel{How the Lean proof follows the argument}
\begin{enumerate}


\item \texttt{hessian f x} is the matrix \texttt{fun i j \ensuremath{\mapsto} iterated\allowbreak{}FDeriv \ensuremath{\mathbb R} 2 f x ![coordinate\allowbreak{}Vector i, coordinate\allowbreak{}Vector j]}. Here \texttt{iterated\allowbreak{}FDeriv \ensuremath{\mathbb R} 2 f x} is the second Fréchet derivative as a multilinear map of two arguments, and \texttt{![u, v]} is the list of its two arguments.

\item \texttt{laplacian f} is Mathlib's \texttt{Laplacian.\allowbreak{}laplacian f} (notation \(\Delta\)), defined without coordinates by contracting the second derivative with the canonical covariant tensor of the inner product; this is the basis-independent trace above.

\item \texttt{Laplacian\allowbreak{}Smooth f Mdelta} is the conjunction \texttt{Cont\allowbreak{}Diff \ensuremath{\mathbb R} 3 f ∧ \ensuremath{\forall} x, \ensuremath{\Vert}gradient (laplacian f) x\ensuremath{\Vert} \textasciicircum{} 2 \ensuremath{\le} Mdelta \textasciicircum{} 2 * (p : \ensuremath{\mathbb R})}; \texttt{Cont\allowbreak{}Diff \ensuremath{\mathbb R} 3 f} means three times continuously differentiable, and the second conjunct is the paper's condition (lip:2).

\item \texttt{rms\allowbreak{}Component\allowbreak{}Smoothness f} takes the family as \texttt{f : Fin n \ensuremath{\to} State p \ensuremath{\to} \ensuremath{\mathbb R}} and returns \texttt{Real.\allowbreak{}sqrt ((∑ i, global\allowbreak{}Smoothness (f i) \textasciicircum{} 2) /\allowbreak{} (n : \ensuremath{\mathbb R}))}, which is \(M_{\rm rms}\).

\item \texttt{trace\_\allowbreak{}hessian} rewrites the Laplacian with the Mathlib identity \texttt{Inner\allowbreak{}Product\allowbreak{}Space.\allowbreak{}laplacian\_\allowbreak{}eq\_\allowbreak{}iterated\allowbreak{}FDeriv\_\allowbreak{}orthonormal\allowbreak{}Basis}, which expresses \(\Delta f(x)\) as \(\sum_iD^2f(x)(v_i,v_i)\) for any orthonormal basis, applied to the standard basis \texttt{Euclidean\allowbreak{}Space.\allowbreak{}basis\allowbreak{}Fun (Fin p) \ensuremath{\mathbb R}}. Unfolding \texttt{Matrix.\allowbreak{}trace}, \texttt{hessian} and \texttt{coordinate\allowbreak{}Vector} and evaluating the basis vectors with \texttt{Euclidean\allowbreak{}Space.\allowbreak{}basis\allowbreak{}Fun\_\allowbreak{}apply}, \texttt{simp} identifies the two sums.

\end{enumerate}
\end{leanbox}
\section{\texorpdfstring{The gradient-step and descent inequalities}{The gradient-step and descent inequalities}}
\label{proof-module-AverageLMC-GradientStep-lean}


This section proves the two deterministic inequalities used in the proof of Theorem 1. They concern one gradient step \(\theta\mapsto\theta-h\nabla f(\theta)\) with step \(h\ge0\) for a differentiable potential \(f\) on a real Hilbert space \(H\) (in the applications \(H=\mathbb R^p\)). The first, the gradient-step lemma (eq:basic-gd) of the paper, uses \(m\)-strong convexity:


\[\|\theta-h\nabla f(\theta)-y\|^2\le(1-mh)\|\theta-y\|^2+2h\bigl(f(y)-f(\theta)\bigr)+h^2\|\nabla f(\theta)\|^2\qquad(\theta,y\in H).\]


The second, the deterministic half (eq:exp-upper) of the smoothing lemma, uses an \(M\)-Lipschitz gradient:


\[2\bigl(f(\theta-h\nabla f(\theta))-f(\theta)\bigr)\le-h(2-hM)\|\nabla f(\theta)\|^2.\]


In the proof of Theorem 1 the first inequality is applied with \(\theta=\vartheta\) and \(y=Y\), where \(Y\sim\pi\) is coupled with the current iterate \(\vartheta\); the second controls the potential after the gradient step. When the two are added, the potential values cancel and the coefficient \(h^2(hM-1)\) of \(\|\nabla f(\vartheta)\|^2\) remains.

\subsection{\texorpdfstring{The gradient-step lemma}{The gradient-step lemma}}
\label{report-step-AverageLMC-GradientStep-lean-1}


\textbf{Claim.} Let \(H\) be a real Hilbert space, let \(f:H\to\mathbb R\) be differentiable and \(m\)-strongly convex for some \(m\in\mathbb R\), and let \(h\ge0\). Then for all \(\theta,y\in H\),


\[\|\theta-h\nabla f(\theta)-y\|^2\le(1-mh)\|\theta-y\|^2+2h\bigl(f(y)-f(\theta)\bigr)+h^2\|\nabla f(\theta)\|^2.\]


\textbf{Proof.} Write \(g=\nabla f(\theta)\) and \(d=\theta-y\). Since \(\theta-hg-y=d-hg\), expanding the square of the norm with \(\|u-v\|^2=\|u\|^2-2\langle u,v\rangle+\|v\|^2\) and \(\|hg\|=|h|\,\|g\|=h\|g\|\) gives


\[\|\theta-hg-y\|^2=\|d\|^2-2h\langle d,g\rangle+h^2\|g\|^2.\qquad(1)\]


The supporting-plane inequality of a differentiable \(m\)-strongly convex function (inequality (3) of the previous section, valid in every real Hilbert space) at the base point \(\theta\) and the point \(y\) reads


\[f(\theta)+\langle g,y-\theta\rangle+\frac m2\|y-\theta\|^2\le f(y).\]


Since \(y-\theta=-d\), we have \(\langle g,y-\theta\rangle=-\langle d,g\rangle\) and \(\|y-\theta\|=\|d\|\), so the inequality becomes


\[-\langle d,g\rangle\le f(y)-f(\theta)-\frac m2\|d\|^2.\qquad(2)\]


Multiplying (2) by the nonnegative number \(2h\) preserves the inequality: \(-2h\langle d,g\rangle\le2h(f(y)-f(\theta))-mh\|d\|^2\). Substituting this bound into (1),


\[\|\theta-hg-y\|^2\le\|d\|^2-mh\|d\|^2+2h\bigl(f(y)-f(\theta)\bigr)+h^2\|g\|^2,\]


which is the claim, since \(\|d\|^2-mh\|d\|^2=(1-mh)\|\theta-y\|^2\). \(\square\)




The hypothesis \(h\ge0\) is used twice: to write \(\|hg\|=h\|g\|\) and to multiply (2) by \(2h\). For \(h=0\) the claim reduces to \(\|\theta-y\|^2\le\|\theta-y\|^2\). No sign condition on \(m\) and no smoothness of \(\nabla f\) are needed.

\begin{leanbox}{Lean: \texttt{gradient\_\allowbreak{}step}}
\leansource{\texttt{gradient\_\allowbreak{}step} \quad AverageLMC/\allowbreak{}GradientStep.lean, lines 24–39}
\begin{leancode}
\leanline{0}{theorem\ gradient\_step\ (f\ :\ H\ →\ ℝ)\ (m\ h\ :\ ℝ)\ (hdiff\ :\ Differentiable\ ℝ\ f)}
\leanline{4}{(hsc\ :\ StrongConvexOn\ Set.univ\ m\ f)\ (hh\ :\ 0\ ≤\ h)\ (θ\ y\ :\ H)\ :}
\leanline{4}{‖θ\ -\ h\ •\ gradient\ f\ θ\ -\ y‖\ \textasciicircum{}\ 2\ ≤}
\leanline{6}{(1\ -\ m\ *\ h)\ *\ ‖θ\ -\ y‖\ \textasciicircum{}\ 2\ +\ 2\ *\ h\ *\ (f\ y\ -\ f\ θ)\ +}
\leanline{8}{h\ \textasciicircum{}\ 2\ *\ ‖gradient\ f\ θ‖\ \textasciicircum{}\ 2}
\end{leancode}
\leanlabel{How the Lean proof follows the argument}
The ambient space is a type \texttt{H} with the instances \texttt{[Normed\allowbreak{}Add\allowbreak{}Comm\allowbreak{}Group H] [Inner\allowbreak{}Product\allowbreak{}Space \ensuremath{\mathbb R} H] [Complete\allowbreak{}Space H]}, that is, a real Hilbert space; \texttt{gradient f \ensuremath{\theta}} is the vector representing the derivative of \texttt{f} at \texttt{\ensuremath{\theta}} and \texttt{h \ensuremath{\bullet} v} is scalar multiplication. The hypotheses of \texttt{gradient\_\allowbreak{}step} are \texttt{hdiff : Differentiable \ensuremath{\mathbb R} f}, \texttt{hsc : Strong\allowbreak{}Convex\allowbreak{}On Set.\allowbreak{}univ m f} and \texttt{hh : 0 \ensuremath{\le} h}.



\begin{enumerate}


\item \texttt{hfirst} is the supporting-plane inequality, the library lemma \texttt{ULA.\allowbreak{}strong\_\allowbreak{}first\_\allowbreak{}order hdiff hsc \ensuremath{\theta} y}.

\item \texttt{hnorm} (\(\|y-\theta\|=\|\theta-y\|\), by \texttt{norm\_\allowbreak{}sub\_\allowbreak{}rev}) and \texttt{hinner} (\(\langle g,y-\theta\rangle=-\langle\theta-y,g\rangle\), by writing \(y-\theta=-(\theta-y)\) with \texttt{module}, then \texttt{inner\_\allowbreak{}neg\_\allowbreak{}right} and \texttt{real\_\allowbreak{}inner\_\allowbreak{}comm}) rewrite \texttt{hfirst} into (2).

\item \texttt{hscaled} multiplies (2) by \(2h\ge0\) with \texttt{mul\_\allowbreak{}le\_\allowbreak{}mul\_\allowbreak{}of\_\allowbreak{}nonneg\_\allowbreak{}left}; the nonnegativity of \(2h\) is checked by \texttt{positivity}.

\item The goal is rewritten into the expansion (1): the vector identity \(\theta-hg-y=(\theta-y)-hg\) (by \texttt{module}), \texttt{norm\_\allowbreak{}sub\_\allowbreak{}sq\_\allowbreak{}real} for the square of a difference, \texttt{real\_\allowbreak{}inner\_\allowbreak{}smul\_\allowbreak{}right} to extract \(h\) from the inner product, and \texttt{norm\_\allowbreak{}smul}, \texttt{Real.\allowbreak{}norm\_\allowbreak{}eq\_\allowbreak{}abs}, \texttt{abs\_\allowbreak{}of\_\allowbreak{}nonneg hh} for \(\|hg\|=h\|g\|\). \texttt{nlinarith} then combines the rewritten goal with \texttt{hscaled}.

\end{enumerate}
\end{leanbox}


Earlier results used here: \cite[Section~3.2.2]{LangevinLean2026ULASC}.
\subsection{\texorpdfstring{Descent of the potential along a gradient step}{Descent of the potential along a gradient step}}
\label{report-step-AverageLMC-GradientStep-lean-2}


\textbf{Claim.} Let \(H\) be a real Hilbert space, let \(f:H\to\mathbb R\) be differentiable with \(M\)-Lipschitz gradient for some \(M\ge0\), and let \(\theta\in H\) and \(h\ge0\). Then


\[2\bigl(f(\theta-h\nabla f(\theta))-f(\theta)\bigr)\le-h(2-hM)\|\nabla f(\theta)\|^2.\qquad(3)\]


If moreover \(hM\le2\), then \(f(\theta-h\nabla f(\theta))\le f(\theta)\).




\textbf{Proof.} Write \(g=\nabla f(\theta)\). The descent inequality \(f(z)\le f(\theta)+\langle g,z-\theta\rangle+\frac M2\|z-\theta\|^2\), proved in the previous section from the fundamental theorem of calculus along the segment from \(\theta\) to \(z\), applies to \(z=\theta-hg\). There \(z-\theta=-hg\), so \(\langle g,z-\theta\rangle=-h\langle g,g\rangle=-h\|g\|^2\) and \(\|z-\theta\|^2=|-h|^2\|g\|^2=h^2\|g\|^2\). Hence


\[f(\theta-hg)\le f(\theta)-h\|g\|^2+\frac M2h^2\|g\|^2.\]


Subtracting \(f(\theta)\) and multiplying by \(2\),


\[2\bigl(f(\theta-hg)-f(\theta)\bigr)\le-2h\|g\|^2+Mh^2\|g\|^2=-h(2-hM)\|g\|^2,\]


which is (3). If \(hM\le2\), the three factors \(-h\le0\), \(2-hM\ge0\) and \(\|g\|^2\ge0\) make the right-hand side of (3) a product of a nonpositive and two nonnegative numbers, hence \(\le0\). Therefore \(2(f(\theta-hg)-f(\theta))\le0\), which is the second assertion. \(\square\)




In Theorem 1 the step satisfies \(hM_\infty\le1\), so the second assertion applies: the gradient step does not increase the potential. Only (3) itself enters the one-step bound, where its right-hand side combines with the term \(h^2\|\nabla f(\vartheta)\|^2\) of the gradient-step lemma.

\begin{leanbox}{Lean: \texttt{gradient\_\allowbreak{}descent}, \texttt{gradient\_\allowbreak{}descent\_\allowbreak{}nonincreasing}}
\leansource{\texttt{gradient\_\allowbreak{}descent} \quad AverageLMC/\allowbreak{}GradientStep.lean, lines 43–51}
\begin{leancode}
\leanline{0}{theorem\ gradient\_descent\ (f\ :\ H\ →\ ℝ)\ (M\ :\ ℝ≥0)\ (hdiff\ :\ Differentiable\ ℝ\ f)}
\leanline{4}{(hLip\ :\ LipschitzWith\ M\ (gradient\ f))\ (θ\ :\ H)\ (h\ :\ ℝ)\ (hh\ :\ 0\ ≤\ h)\ :}
\leanline{4}{2\ *\ (f\ (θ\ -\ h\ •\ gradient\ f\ θ)\ -\ f\ θ)\ ≤}
\leanline{6}{-h\ *\ (2\ -\ h\ *\ (M\ :\ ℝ))\ *\ ‖gradient\ f\ θ‖\ \textasciicircum{}\ 2}
\end{leancode}
\leansource{\texttt{gradient\_\allowbreak{}descent\_\allowbreak{}nonincreasing} \quad AverageLMC/\allowbreak{}GradientStep.lean, lines 54–63}
\begin{leancode}
\leanline{0}{theorem\ gradient\_descent\_nonincreasing\ (f\ :\ H\ →\ ℝ)\ (M\ :\ ℝ≥0)}
\leanline{4}{(hdiff\ :\ Differentiable\ ℝ\ f)\ (hLip\ :\ LipschitzWith\ M\ (gradient\ f))}
\leanline{4}{(θ\ :\ H)\ (h\ :\ ℝ)\ (hh\ :\ 0\ ≤\ h)\ (hstep\ :\ h\ *\ (M\ :\ ℝ)\ ≤\ 2)\ :}
\leanline{4}{f\ (θ\ -\ h\ •\ gradient\ f\ θ)\ ≤\ f\ θ}
\end{leancode}
\leanlabel{How the Lean proof follows the argument}
The Lipschitz hypothesis is \texttt{h\allowbreak{}Lip : Lipschitz\allowbreak{}With M (gradient f)} with \texttt{M : \ensuremath{\mathbb R}\ensuremath{\ge}0}, a nonnegative real; \texttt{(M : \ensuremath{\mathbb R})} is its coercion to \(\mathbb R\).



\begin{enumerate}


\item In \texttt{gradient\_\allowbreak{}descent}, \texttt{hs} is the descent inequality \texttt{ULA.\allowbreak{}Analysis.\allowbreak{}smooth\_\allowbreak{}upper\_\allowbreak{}bound} applied at \texttt{\ensuremath{\theta}} and \texttt{\ensuremath{\theta} -\allowbreak{} h \ensuremath{\bullet} gradient f \ensuremath{\theta}}. \texttt{heq} computes the displacement \(-hg\) (by \texttt{module}), and the rewrites \texttt{real\_\allowbreak{}inner\_\allowbreak{}smul\_\allowbreak{}right}, \texttt{real\_\allowbreak{}inner\_\allowbreak{}self\_\allowbreak{}eq\_\allowbreak{}norm\_\allowbreak{}sq}, \texttt{norm\_\allowbreak{}smul}, \texttt{Real.\allowbreak{}norm\_\allowbreak{}eq\_\allowbreak{}abs}, \texttt{abs\_\allowbreak{}neg} and \texttt{abs\_\allowbreak{}of\_\allowbreak{}nonneg hh} evaluate \(\langle g,-hg\rangle=-h\|g\|^2\) and \(\|-hg\|=h\|g\|\) inside \texttt{hs}. \texttt{nlinarith} multiplies by \(2\) and rearranges, giving (3).

\item \texttt{gradient\_\allowbreak{}descent\_\allowbreak{}nonincreasing} adds \texttt{hstep : h * (M : \ensuremath{\mathbb R}) \ensuremath{\le} 2}. \texttt{hneg} shows that the right-hand side of (3) is \(\le0\) by two applications of \texttt{mul\_\allowbreak{}nonpos\_\allowbreak{}of\_\allowbreak{}nonpos\_\allowbreak{}of\_\allowbreak{}nonneg}, with \texttt{neg\_\allowbreak{}nonpos.\allowbreak{}mpr hh} for \(-h\le0\), \texttt{sub\_\allowbreak{}nonneg.\allowbreak{}mpr hstep} for \(2-hM\ge0\) and \texttt{sq\_\allowbreak{}nonneg} for \(\|g\|^2\ge0\); \texttt{linarith} combines \texttt{hs} (inequality (3)) with \texttt{hneg}.

\end{enumerate}
\end{leanbox}


Earlier results used here: \cite[Section~3.9.1]{LangevinLean2026ULASC}.
\section{\texorpdfstring{Quadratic integrability and exact energy of the synchronous process}{Quadratic integrability and exact energy of the synchronous process}}
\label{proof-module-AverageLMC-PathwiseEnergy-lean}


All expectations and derivatives in the stochastic proof are justified on the actual space of square-integrable random variables. A smooth potential need not be globally Lipschitz: its quadratic growth is the correct condition here. The common additive noise cancels before the squared error is differentiated.



Writing \(\langle U,V\rangle_{L^2}=\mathbb E\langle U,V\rangle\), strong convexity and the common-noise integral equation give


\[2\langle X-Y,\nabla f(Y)\rangle_{L^2}\le 2\bigl(\mathbb Ef(X)-\mathbb Ef(Y)\bigr)-m\|X-Y\|_{L^2}^{2}.\]
\[\|X+B_h-Y_h\|_{L^2}^{2}=\|X-Y_0\|_{L^2}^{2}+2\int_0^h\langle X+B_t-Y_t,\nabla f(Y_t)\rangle_{L^2}\,dt.\]

Earlier results used here: \cite[Section~3.15]{LangevinLean2026Proximal}.
\subsection{\texorpdfstring{Smooth potentials are integrable on \(L^{2}\)}{Smooth potentials are integrable on L²}}
\label{report-step-AverageLMC-PathwiseEnergy-lean-1}


Apply the quadratic upper model both to \(f\) and to \(-f\) at zero. These two bounds control the absolute value of \(f(x)\) by the sum of \(|f(0)|\), the product of the initial gradient norm and \(\|x\|\), and \(L\|x\|^{2}/2\). A square-integrable random vector has integrable norm and squared norm, so comparison proves integrability of \(f\) evaluated at its \(L^{2}\) representative. The second theorem transfers the result from the quotient-space representative to any supplied random vector with finite second moment by its almost-everywhere equality with the canonical square-integrable representative.

\begin{leanbox}{Lean: \texttt{integrable\_\allowbreak{}smooth\_\allowbreak{}compLp}, \texttt{integrable\_\allowbreak{}smooth\_\allowbreak{}comp}}
\leansource{\texttt{integrable\_\allowbreak{}smooth\_\allowbreak{}compLp} \quad AverageLMC/\allowbreak{}PathwiseEnergy.lean, lines 18–44}
\begin{leancode}
\leanline{0}{theorem\ integrable\_smooth\_compLp\ (f\ :\ E\ →\ ℝ)\ (L\ :\ ℝ≥0)}
\leanline{4}{(hf\ :\ Differentiable\ ℝ\ f)\ (hLip\ :\ LipschitzWith\ L\ (gradient\ f))}
\leanline{4}{(X\ :\ Lp\ E\ 2\ P)\ :\ Integrable\ (fun\ ω\ ↦\ f\ (X\ ω))\ P}
\end{leancode}
\leansource{\texttt{integrable\_\allowbreak{}smooth\_\allowbreak{}comp} \quad AverageLMC/\allowbreak{}PathwiseEnergy.lean, lines 47–52}
\begin{leancode}
\leanline{0}{theorem\ integrable\_smooth\_comp\ (f\ :\ E\ →\ ℝ)\ (L\ :\ ℝ≥0)}
\leanline{4}{(hf\ :\ Differentiable\ ℝ\ f)\ (hLip\ :\ LipschitzWith\ L\ (gradient\ f))}
\leanline{4}{(X\ :\ Ω\ →\ E)\ (hX\ :\ MemLp\ X\ 2\ P)\ :\ Integrable\ (fun\ ω\ ↦\ f\ (X\ ω))\ P}
\end{leancode}
\leanlabel{How the Lean proof follows the argument}
\texttt{integrable\_\allowbreak{}smooth\_\allowbreak{}comp\allowbreak{}Lp} gives the formal counterpart of this step. The main intermediate facts are \texttt{hu}, \texttt{hneg}, \texttt{h\allowbreak{}Ln}, \texttt{hl}, \texttt{hgrowth}, \texttt{hn} in that order. Named proof applications include \texttt{h\allowbreak{}Lip.\allowbreak{}neg}; the adjacent proof shows their hypotheses and order of use.




\texttt{integrable\_\allowbreak{}smooth\_\allowbreak{}comp} gives the formal counterpart of this step.
\end{leanbox}


Earlier results used here: \cite[Section~3.9.1]{LangevinLean2026ULASC}; \cite[Section~3.11.1]{LangevinLean2026IPLA}.
\subsection{\texorpdfstring{Integrate the strong supporting plane}{Integrate the strong supporting plane}}
\label{report-step-AverageLMC-PathwiseEnergy-lean-2}


At each outcome, strong convexity gives \(f(X)-f(Y)\) at least the inner product of \(\nabla f(Y)\) with \(X-Y\) plus \(m\|X-Y\|^{2}/2\). Every scalar term is integrable by quadratic growth, and the inner product of two \(L^{2}\) vectors is integrable by the \(L^{2}\) Cauchy–Schwarz inequality. Integration, the definition of the \(L^{2}\) inner product, and the exact squared-norm integral identity give the supporting-plane inequality on \(L^{2}\) itself. The lifted vector field is the actual gradient applied pointwise, rather than a supplied random drift.

\begin{leanbox}{Lean: \texttt{lp\_\allowbreak{}strong\_\allowbreak{}gradient\_\allowbreak{}support}}
\leansource{\texttt{lp\_\allowbreak{}strong\_\allowbreak{}gradient\_\allowbreak{}support} \quad AverageLMC/\allowbreak{}PathwiseEnergy.lean, lines 55–85}
\begin{leancode}
\leanline{0}{theorem\ lp\_strong\_gradient\_support\ (f\ :\ E\ →\ ℝ)\ (m\ :\ ℝ)\ (L\ :\ ℝ≥0)}
\leanline{4}{(hf\ :\ Differentiable\ ℝ\ f)\ (hsc\ :\ StrongConvexOn\ univ\ m\ f)}
\leanline{4}{(hLip\ :\ LipschitzWith\ L\ (gradient\ f))\ (X\ Y\ :\ Lp\ E\ 2\ P)\ :}
\leanline{4}{2\ *\ {\leanfallback ⟪}X\ -\ Y,\ ULA.lpDrift\ P\ (gradient\ f)\ hLip\ Y{\leanfallback ⟫}\ ≤}
\leanline{6}{2\ *\ ((∫\ ω,\ f\ (X\ ω)\ ∂P)\ -\ ∫\ ω,\ f\ (Y\ ω)\ ∂P)\ -\ m\ *\ ‖X\ -\ Y‖\ \textasciicircum{}\ 2}
\end{leancode}
\leanlabel{How the Lean proof follows the argument}
\texttt{lp\_\allowbreak{}strong\_\allowbreak{}gradient\_\allowbreak{}support} gives the formal counterpart of this step. The main intermediate facts are \texttt{hi\allowbreak{}X}, \texttt{hi\allowbreak{}Y}, \texttt{hn}, \texttt{hright}, \texttt{hi}, \texttt{heq} in that order. Named proof applications include \texttt{hi}; the adjacent proof shows their hypotheses and order of use.
\end{leanbox}


Earlier results used here: \cite[Section~3.13.2]{LangevinLean2026ULASC}; \cite[Section~3.13.3]{LangevinLean2026ULASC}; \cite[Section~3.2.2]{LangevinLean2026ULASC}; \cite[Section~3.11.1]{LangevinLean2026IPLA}; \cite[Section~3.15.2]{LangevinLean2026Proximal}.
\subsection{\texorpdfstring{Differentiate after the common noise cancels}{Differentiate after the common noise cancels}}
\label{report-step-AverageLMC-PathwiseEnergy-lean-3}


For a continuous vector field \(g\), the primitive \(A+\int_0^t g(s)\,ds\) has derivative \(g(t)\). The derivative of its squared Hilbert norm is twice its inner product with \(g(t)\); the fundamental theorem of calculus proves the displayed exact energy identity. In the Langevin integral equation, \(X+B(t)-Y(t)\) equals \(X-Y(0)+\int_0^t\nabla f(Y(s))\,ds\) because the same \(B\) drives both processes. Substitute this expression into the Hilbert identity. No derivative of a Brownian path and no stochastic isometry is assumed.

\begin{leanbox}{Lean: \texttt{integral\_\allowbreak{}curve\_\allowbreak{}norm\_\allowbreak{}sq}, \texttt{langevin\_\allowbreak{}energy\_\allowbreak{}identity}}
\leansource{\texttt{integral\_\allowbreak{}curve\_\allowbreak{}norm\_\allowbreak{}sq} \quad AverageLMC/\allowbreak{}PathwiseEnergy.lean, lines 89–103}
\begin{leancode}
\leanline{0}{theorem\ integral\_curve\_norm\_sq\ (g\ :\ ℝ\ →\ E)\ (hg\ :\ Continuous\ g)\ (A\ :\ E)\ (h\ :\ ℝ)\ :}
\leanline{4}{‖A\ +\ ∫\ s\ in\ 0..h,\ g\ s‖\ \textasciicircum{}\ 2\ =\ ‖A‖\ \textasciicircum{}\ 2\ +}
\leanline{6}{∫\ t\ in\ 0..h,\ 2\ *\ {\leanfallback ⟪}A\ +\ ∫\ s\ in\ 0..t,\ g\ s,\ g\ t{\leanfallback ⟫}}
\end{leancode}
\leansource{\texttt{langevin\_\allowbreak{}energy\_\allowbreak{}identity} \quad AverageLMC/\allowbreak{}PathwiseEnergy.lean, lines 106–127}
\begin{leancode}
\leanline{0}{theorem\ langevin\_energy\_identity\ (f\ :\ E\ →\ ℝ)\ (L\ :\ ℝ≥0)}
\leanline{4}{(hLip\ :\ LipschitzWith\ L\ (gradient\ f))\ (X\ :\ Lp\ E\ 2\ P)}
\leanline{4}{(Y\ B\ :\ ℝ\ →\ Lp\ E\ 2\ P)\ (hY\ :\ Continuous\ Y)\ (h\ :\ ℝ)\ (hh\ :\ 0\ ≤\ h)}
\leanline{4}{(hYe\ :\ ∀\ t\ ∈\ Icc\ 0\ h,\ Y\ t\ -\ Y\ 0\ =}
\leanline{6}{-(∫\ s\ in\ 0..t,\ ULA.lpDrift\ P\ (gradient\ f)\ hLip\ (Y\ s))\ +\ B\ t)\ :}
\leanline{4}{‖X\ +\ B\ h\ -\ Y\ h‖\ \textasciicircum{}\ 2\ =\ ‖X\ -\ Y\ 0‖\ \textasciicircum{}\ 2\ +}
\leanline{6}{∫\ t\ in\ 0..h,\ 2\ *\ {\leanfallback ⟪}X\ +\ B\ t\ -\ Y\ t,\ ULA.lpDrift\ P\ (gradient\ f)\ hLip\ (Y\ t){\leanfallback ⟫}}
\end{leancode}
\leanlabel{How the Lean proof follows the argument}
\texttt{integral\_\allowbreak{}curve\_\allowbreak{}norm\_\allowbreak{}sq} gives the formal counterpart of this step. The main intermediate facts are \texttt{G}, \texttt{hd}, \texttt{hn}, \texttt{h\allowbreak{}G}, \texttt{he} in that order. Its remaining equalities are discharged by the simplifications displayed in the proof.




\texttt{langevin\_\allowbreak{}energy\_\allowbreak{}identity} gives the formal counterpart of this step. The main intermediate facts are \texttt{g}, \texttt{hg}, \texttt{he} in that order. Named proof applications include \texttt{interval\allowbreak{}Integral.\allowbreak{}integral\_\allowbreak{}congr}; the adjacent proof shows their hypotheses and order of use.
\end{leanbox}


Earlier results used here: \cite[Section~3.13.2]{LangevinLean2026ULASC}; \cite[Section~3.13.5]{LangevinLean2026ULASC}.
\section{\texorpdfstring{Random-center smoothing and the integrated Gaussian energy bound}{Random-center smoothing and the integrated Gaussian energy bound}}
\label{proof-module-AverageLMC-GaussianPerturbation-lean}


This module carries out the expectation and time-integration steps used in the improved one-step estimate. Independence promotes a deterministic smoothing estimate to a random center; strong convexity and orthogonality then give the exact Gaussian energy constant.



The energy theorem uses an affine upper bound for the expected potential difference, the exact Brownian increment moment, orthogonality, and a common lower bound \(w\) for \(\|X+B_h-Y_t\|_{L^2}^2\). Its key input and conclusion are


\[\mathbb Ef(X+B_t)-\mathbb Ef(Y_t)\le a+bt\qquad(0\le t\le h).\]
\[(1+mh)w\le\|X-Y_0\|_{L^2}^{2}+2ha+h^2b+mh^2d.\]

Earlier results used here: \cite[Section~3.18]{LangevinLean2026Proximal}; \cite[Section~3.45]{LangevinLean2026ULASC}.
\subsection{\texorpdfstring{Promote smoothing to an independent random center}{Promote smoothing to an independent random center}}
\label{report-step-AverageLMC-GaussianPerturbation-lean-1}


For independent \(L^{2}\) vectors \(X\) and \(Z\), their joint pushforward measure is the product of their marginal laws. The function \(F(x,z)=f(x+cz)-f(x)\) is continuous and integrable under this joint law because both arguments have finite second moment and \(f\) has quadratic growth. The pushforward identity and Fubini turn its expectation into the iterated integral. Integrate the pointwise smoothing bound over the law of \(X\), which is a probability measure, to retain the same constant. Finally use almost-everywhere representatives to identify the pointwise sum with the sum in \(L^{2}\).

\begin{leanbox}{Lean: \texttt{independent\_\allowbreak{}smoothing}}
\leansource{\texttt{independent\_\allowbreak{}smoothing} \quad AverageLMC/\allowbreak{}GaussianPerturbation.lean, lines 21–53}
\begin{leancode}
\leanline{0}{theorem\ independent\_smoothing\ (f\ :\ E\ →\ ℝ)\ (L\ :\ ℝ≥0)}
\leanline{4}{(hf\ :\ Differentiable\ ℝ\ f)\ (hLip\ :\ LipschitzWith\ L\ (gradient\ f))}
\leanline{4}{(X\ Z\ :\ Lp\ E\ 2\ P)\ (η\ :\ Measure\ E)\ (hZ\ :\ P.map\ Z\ =\ η)}
\leanline{4}{(hind\ :\ IndepFun\ X\ Z\ P)\ (c\ C\ :\ ℝ)}
\leanline{4}{(hsm\ :\ ∀\ x,\ (∫\ z,\ f\ (x\ +\ c\ •\ z)\ -\ f\ x\ ∂η)\ ≤\ C)\ :}
\leanline{4}{(∫\ ω,\ f\ ((X\ +\ c\ •\ Z)\ ω)\ ∂P)\ -\ (∫\ ω,\ f\ (X\ ω)\ ∂P)\ ≤\ C}
\end{leancode}
\leanlabel{How the Lean proof follows the argument}
\texttt{independent\_\allowbreak{}smoothing} gives the formal counterpart of this step. The main intermediate facts are \texttt{F}, \texttt{h\allowbreak{}Fm}, \texttt{hjoint}, \texttt{hpair}, \texttt{hcomp}, \texttt{h\allowbreak{}Fint} in that order. Named proof applications include \texttt{hsm}, \texttt{integral\_\allowbreak{}congr\_\allowbreak{}ae}, \texttt{heq}; the adjacent proof shows their hypotheses and order of use.
\end{leanbox}
\subsection{\texorpdfstring{Integrate strong convexity with the future-noise correction}{Integrate strong convexity with the future-noise correction}}
\label{report-step-AverageLMC-GaussianPerturbation-lean-2}


Let \(w\) be a common lower bound for the squared coupling costs between \(X+B(h)\) and the stationary comparison \(Y(t)\). Orthogonality expands the latter cost exactly as \(\|X+B(t)-Y(t)\|^{2}+2(h-t)d\). Integrate the strong supporting-plane bound, using an affine upper bound \(a+b\) \(t\) for the potential expectation difference. The integrand is bounded by \(2a+2mhd-mw+(2b-2md)t\). Its explicit integral is \(2ha+mh^{2}d-mhw+h^{2}b\). Combine this with the exact synchronous energy identity and the endpoint coupling bound to obtain \((1+mh)w\le \|X-Y(0)\|^{2}+2ha+h^{2}b+mh^{2}d\). Every input involving noise or stationarity is discharged for the actual Gaussian step by the canonical construction in the following modules.

\begin{leanbox}{Lean: \texttt{gaussian\_\allowbreak{}energy\_\allowbreak{}bound}}
\leansource{\texttt{gaussian\_\allowbreak{}energy\_\allowbreak{}bound} \quad AverageLMC/\allowbreak{}GaussianPerturbation.lean, lines 58–106}
\begin{leancode}
\leanline{0}{theorem\ gaussian\_energy\_bound\ (f\ :\ E\ →\ ℝ)\ (m\ h\ d\ a\ b\ w\ :\ ℝ)\ (L\ :\ ℝ≥0)}
\leanline{4}{(hf\ :\ Differentiable\ ℝ\ f)\ (hsc\ :\ StrongConvexOn\ univ\ m\ f)}
\leanline{4}{(hLip\ :\ LipschitzWith\ L\ (gradient\ f))\ (hm\ :\ 0\ ≤\ m)\ (hh\ :\ 0\ ≤\ h)}
\leanline{4}{(X\ :\ Lp\ E\ 2\ P)\ (Y\ B\ :\ ℝ\ →\ Lp\ E\ 2\ P)\ (hY\ :\ Continuous\ Y)\ (hB\ :\ Continuous\ B)}
\leanline{4}{(hYe\ :\ ∀\ t\ ∈\ Icc\ 0\ h,\ Y\ t\ -\ Y\ 0\ =}
\leanline{6}{-(∫\ s\ in\ 0..t,\ ULA.lpDrift\ P\ (gradient\ f)\ hLip\ (Y\ s))\ +\ B\ t)}
\leanline{4}{(horth\ :\ ∀\ t\ ∈\ Icc\ 0\ h,\ {\leanfallback ⟪}X\ +\ B\ t\ -\ Y\ t,\ B\ h\ -\ B\ t{\leanfallback ⟫}\ =\ 0)}
\leanline{4}{(hnoise\ :\ ∀\ t\ ∈\ Icc\ 0\ h,\ ‖B\ h\ -\ B\ t‖\ \textasciicircum{}\ 2\ =\ 2\ *\ (h\ -\ t)\ *\ d)}
\leanline{4}{(hcost\ :\ ∀\ t\ ∈\ Icc\ 0\ h,\ w\ ≤\ ‖X\ +\ B\ h\ -\ Y\ t‖\ \textasciicircum{}\ 2)}
\leanline{4}{(henergy\ :\ ∀\ t\ ∈\ Icc\ 0\ h,}
\leanline{6}{(∫\ ω,\ f\ ((X\ +\ B\ t)\ ω)\ ∂P)\ -\ (∫\ ω,\ f\ (Y\ t\ ω)\ ∂P)\ ≤\ a\ +\ b\ *\ t)\ :}
\leanline{4}{(1\ +\ m\ *\ h)\ *\ w\ ≤\ ‖X\ -\ Y\ 0‖\ \textasciicircum{}\ 2\ +\ 2\ *\ h\ *\ a\ +\ h\ \textasciicircum{}\ 2\ *\ b\ +\ m\ *\ h\ \textasciicircum{}\ 2\ *\ d}
\end{leancode}
\leanlabel{How the Lean proof follows the argument}
\texttt{gaussian\_\allowbreak{}energy\_\allowbreak{}bound} gives the formal counterpart of this step. The main intermediate facts are \texttt{g}, \texttt{hg}, \texttt{hi}, \texttt{hpoint}, \texttt{hr}, \texttt{hmono} in that order. Its remaining equalities are discharged by the simplifications displayed in the proof.
\end{leanbox}


Earlier results used here: \cite[Section~3.13.2]{LangevinLean2026ULASC}; \cite[Section~3.13.5]{LangevinLean2026ULASC}.
\section{\texorpdfstring{Exact future Brownian moments and derived orthogonality}{Exact future Brownian moments and derived orthogonality}}
\label{proof-module-AverageLMC-GaussianBrownian-lean}


The improved energy estimate uses the actual Brownian process underlying the stationary Langevin solution. Its future increments have exact moments, and their orthogonality to the coupled past is derived from independence and adaptedness.



For the scaled noise \(B_t=\sqrt2W_t\), genuine independence and the constructed adapted integral solution imply, for \(0\le t\le h\),


\[\mathbb E(B_h-B_t)=0,\qquad\|B_h-B_t\|_{L^2}^{2}=2(h-t)p.\]
\[\langle X+B_t-Y_t,\,B_h-B_t\rangle_{L^2}=0.\]\subsection{\texorpdfstring{Center and measure future Brownian noise}{Center and measure future Brownian noise}}
\label{report-step-AverageLMC-GaussianBrownian-lean-1}


The Langevin additive noise is \(\sqrt 2\) times a Brownian increment. Its expectation vanishes at every time, so the expectation of a future difference also vanishes. Algebraically, the difference between the noises at \(h\) and \(t\) is the additive noise over the segment starting at \(t\) with length \(h-t\). The existing exact Brownian \(L^{2}\) norm formula therefore gives its squared norm as \(2(h-t)d\). This uses the true elapsed-time Gaussian variance and covers the endpoint \(t=h\).

\begin{leanbox}{Lean: \texttt{brownian\_\allowbreak{}future\_\allowbreak{}mean\_\allowbreak{}zero}, \texttt{brownian\_\allowbreak{}future\_\allowbreak{}norm\_\allowbreak{}sq}}
\leansource{\texttt{brownian\_\allowbreak{}future\_\allowbreak{}mean\_\allowbreak{}zero} \quad AverageLMC/\allowbreak{}GaussianBrownian.lean, lines 15–25}
\begin{leancode}
\leanline{0}{theorem\ brownian\_future\_mean\_zero}
\leanline{4}{(B\ :\ ι\ →\ ℝ≥0\ →\ Ω\ →\ ℝ)\ (hB\ :\ ∀\ i,\ IsPreBrownianReal\ (B\ i)\ P)\ (t\ h\ :\ ℝ)\ :}
\leanline{4}{(∫\ ω,\ (ULA.brownianLangevinNoise\ P\ B\ hB\ 0\ h\ -}
\leanline{6}{ULA.brownianLangevinNoise\ P\ B\ hB\ 0\ t\ :\ Lp\ (EuclideanSpace\ ℝ\ ι)\ 2\ P)\ ω\ ∂P)\ =\ 0}
\end{leancode}
\leansource{\texttt{brownian\_\allowbreak{}future\_\allowbreak{}norm\_\allowbreak{}sq} \quad AverageLMC/\allowbreak{}GaussianBrownian.lean, lines 28–43}
\begin{leancode}
\leanline{0}{theorem\ brownian\_future\_norm\_sq}
\leanline{4}{(B\ :\ ι\ →\ ℝ≥0\ →\ Ω\ →\ ℝ)\ (hB\ :\ ∀\ i,\ IsPreBrownianReal\ (B\ i)\ P)}
\leanline{4}{(t\ h\ :\ ℝ)\ (ht\ :\ 0\ ≤\ t)\ (hth\ :\ t\ ≤\ h)\ :}
\leanline{4}{‖ULA.brownianLangevinNoise\ P\ B\ hB\ 0\ h\ -}
\leanline{6}{ULA.brownianLangevinNoise\ P\ B\ hB\ 0\ t‖\ \textasciicircum{}\ 2\ =}
\leanline{6}{2\ *\ (h\ -\ t)\ *\ (Fintype.card\ ι\ :\ ℝ)}
\end{leancode}
\leanlabel{How the Lean proof follows the argument}
\texttt{brownian\_\allowbreak{}future\_\allowbreak{}mean\_\allowbreak{}zero} gives the formal counterpart of this step.




\texttt{brownian\_\allowbreak{}future\_\allowbreak{}norm\_\allowbreak{}sq} gives the formal counterpart of this step. The main intermediate facts are \texttt{he} in that order. Its remaining equalities are discharged by the simplifications displayed in the proof.
\end{leanbox}


Earlier results used here: \cite[Section~3.25.1]{LangevinLean2026ULASC}; \cite[Section~3.25.4]{LangevinLean2026ULASC}; \cite[Section~3.11.1]{LangevinLean2026IPLA}; \cite[Section~3.18.4]{LangevinLean2026Proximal}.
\subsection{\texorpdfstring{Derive the orthogonality of the actual error}{Derive the orthogonality of the actual error}}
\label{report-step-AverageLMC-GaussianBrownian-lean-2}


The complete initial state and stopped Brownian past are independent of the future Brownian increment. Projection onto the initial-state and current-noise coordinates gives zero \(L^{2}\) inner product with each of them because the future increment is centered. For \(Y(t)\), the natural-past measurability of the actual integral solution has already been proved by the invariant-subspace construction; applying it to the identity vector field gives the same zero inner product. Subtract and add the three identities to prove that \(X+B(t)-Y(t)\) is orthogonal to \(B(h)-B(t)\). No stochastic-error orthogonality condition is assumed in the final LMC theorem.

\begin{leanbox}{Lean: \texttt{gaussian\_\allowbreak{}error\_\allowbreak{}future\_\allowbreak{}inner\_\allowbreak{}zero}}
\leansource{\texttt{gaussian\_\allowbreak{}error\_\allowbreak{}future\_\allowbreak{}inner\_\allowbreak{}zero} \quad AverageLMC/\allowbreak{}GaussianBrownian.lean, lines 47–82}
\begin{leancode}
\leanline{0}{theorem\ gaussian\_error\_future\_inner\_zero}
\leanline{4}{(f\ :\ EuclideanSpace\ ℝ\ ι\ →\ ℝ)\ (L\ :\ ℝ≥0)}
\leanline{4}{(hLip\ :\ LipschitzWith\ L\ (gradient\ f))}
\leanline{4}{(X\ Z\ :\ Lp\ (EuclideanSpace\ ℝ\ ι)\ 2\ P)}
\leanline{4}{(B\ :\ ι\ →\ ℝ≥0\ →\ Ω\ →\ ℝ)\ (hB\ :\ ∀\ i,\ IsPreBrownianReal\ (B\ i)\ P)}
\leanline{4}{(hmeas\ :\ ∀\ i\ s,\ Measurable\ (B\ i\ s))}
\leanline{4}{(hpaths\ :\ iIndepFun\ (fun\ i\ ω\ s\ ↦\ B\ i\ s\ ω)\ P)}
\leanline{4}{(hX\ :\ IndepFun\ X\ (fun\ ω\ (i\ :\ ι)\ (s\ :\ ℝ≥0)\ ↦\ B\ i\ s\ ω)\ P)}
\leanline{4}{(hZ\ :\ IndepFun\ Z\ (fun\ ω\ (i\ :\ ι)\ (s\ :\ ℝ≥0)\ ↦\ B\ i\ s\ ω)\ P)}
\leanline{4}{(Y\ :\ ℝ\ →\ Lp\ (EuclideanSpace\ ℝ\ ι)\ 2\ P)\ (hY\ :\ Continuous\ Y)}
\leanline{4}{(t\ h\ :\ ℝ)\ (ht\ :\ 0\ ≤\ t)\ (hth\ :\ t\ ≤\ h)}
\leanline{4}{(hYe\ :\ ∀\ s\ ∈\ Icc\ 0\ h,\ Y\ s\ -\ Z\ =}
\leanline{6}{-(∫\ r\ in\ 0..s,\ ULA.lpDrift\ P\ (gradient\ f)\ hLip\ (Y\ r))\ +}
\leanline{8}{ULA.brownianLangevinNoise\ P\ B\ hB\ 0\ s)\ :}
\leanline{4}{{\leanfallback ⟪}X\ +\ ULA.brownianLangevinNoise\ P\ B\ hB\ 0\ t\ -\ Y\ t,}
\leanline{6}{ULA.brownianLangevinNoise\ P\ B\ hB\ 0\ h\ -}
\leanline{8}{ULA.brownianLangevinNoise\ P\ B\ hB\ 0\ t{\leanfallback ⟫}\ =\ 0}
\end{leancode}
\leanlabel{How the Lean proof follows the argument}
\texttt{gaussian\_\allowbreak{}error\_\allowbreak{}future\_\allowbreak{}inner\_\allowbreak{}zero} gives the formal counterpart of this step. The main intermediate facts are \texttt{N}, \texttt{hind}, \texttt{h\allowbreak{}IX}, \texttt{h\allowbreak{}IB}, \texttt{hmean}, \texttt{hx} in that order. Named proof applications include \texttt{hind.\allowbreak{}comp}, \texttt{Lp.\allowbreak{}ext}, \texttt{ULA.\allowbreak{}lp\allowbreak{}Drift\_\allowbreak{}coe}; the adjacent proof shows their hypotheses and order of use.
\end{leanbox}


Earlier results used here: \cite[Section~3.13.2]{LangevinLean2026ULASC}; \cite[Section~3.13.3]{LangevinLean2026ULASC}; \cite[Section~3.25.1]{LangevinLean2026ULASC}; \cite[Section~3.18.3]{LangevinLean2026Proximal}; \cite[Section~3.18.4]{LangevinLean2026Proximal}; \cite[Section~3.18.5]{LangevinLean2026Proximal}.
\section{\texorpdfstring{Expected potential and gradient-step energies on \(L^{2}\)}{Expected potential and gradient-step energies on L²}}
\label{proof-module-AverageLMC-ExpectedGradientStep-lean}


The deterministic gradient-step and descent inequalities extend to the actual Hilbert space of square-integrable random variables. Potential integrability is proved from quadratic growth, and the lifted drift is identified almost everywhere with the pointwise gradient. These estimates retain the full gradient energy, including its sign-sensitive correction in the sampler one-step bound.

\subsection{\texorpdfstring{Integrate the strong gradient-step inequality}{Integrate the strong gradient-step inequality}}
\label{report-step-AverageLMC-ExpectedGradientStep-lean-1}


Apply the integrated strong supporting-plane inequality to \(Y\) and \(X\). Reverse the displacement \(Y\)-\(X\) to \(X\)-\(Y\) and multiply by nonnegative \(h\). Expand the squared \(L^{2}\) distance after subtracting \(h\) times the lifted gradient. The mixed inner product is bounded by the potential difference and strong convexity, giving

\[\|X-h\nabla f(X)-Y\|_{L^2}^2\le(1-mh)\|X-Y\|_{L^2}^2+2h(\mathbb Ef(Y)-\mathbb Ef(X))+h^2\|\nabla f(X)\|_{L^2}^2.\]

The norms are genuine \(L^{2}\) norms and every expectation is finite.

\begin{leanbox}{Lean: \texttt{lp\_\allowbreak{}gradient\_\allowbreak{}step}}
\leansource{\texttt{lp\_\allowbreak{}gradient\_\allowbreak{}step} \quad AverageLMC/\allowbreak{}ExpectedGradientStep.lean, lines 15–33}
\begin{leancode}
\leanline{0}{theorem\ lp\_gradient\_step\ (f\ :\ E\ →\ ℝ)\ (m\ h\ :\ ℝ)\ (L\ :\ ℝ≥0)}
\leanline{4}{(hf\ :\ Differentiable\ ℝ\ f)\ (hsc\ :\ StrongConvexOn\ Set.univ\ m\ f)}
\leanline{4}{(hLip\ :\ LipschitzWith\ L\ (gradient\ f))\ (hh\ :\ 0\ ≤\ h)\ (X\ Y\ :\ Lp\ E\ 2\ P)\ :}
\leanline{4}{‖X\ -\ h\ •\ ULA.lpDrift\ P\ (gradient\ f)\ hLip\ X\ -\ Y‖\ \textasciicircum{}\ 2\ ≤}
\leanline{6}{(1\ -\ m\ *\ h)\ *\ ‖X\ -\ Y‖\ \textasciicircum{}\ 2\ +}
\leanline{6}{2\ *\ h\ *\ ((∫\ ω,\ f\ (Y\ ω)\ ∂P)\ -\ ∫\ ω,\ f\ (X\ ω)\ ∂P)\ +}
\leanline{6}{h\ \textasciicircum{}\ 2\ *\ ‖ULA.lpDrift\ P\ (gradient\ f)\ hLip\ X‖\ \textasciicircum{}\ 2}
\end{leancode}
\leanlabel{How the Lean proof follows the argument}
\texttt{lp\_\allowbreak{}gradient\_\allowbreak{}step} gives the formal counterpart of this step. The main intermediate facts are \texttt{hs}, \texttt{hn}, \texttt{hi}, \texttt{hscaled} in that order. The final step is the algebraic normalization or ordered arithmetic displayed in the proof.
\end{leanbox}


Earlier results used here: \cite[Section~3.13.2]{LangevinLean2026ULASC}.
\subsection{\texorpdfstring{Integrate deterministic descent without dropping the score term}{Integrate deterministic descent without dropping the score term}}
\label{report-step-AverageLMC-ExpectedGradientStep-lean-2}


The pointwise deterministic descent theorem holds outside a null set where the lifted state arithmetic and lifted gradient agree with their ordinary representatives. Both potential terms are integrable by smooth quadratic growth, and the squared score is integrable by the Lipschitz drift construction. Integrate the pointwise estimate, use linearity, and identify the integral of the squared lifted gradient with its squared \(L^{2}\) norm. The result is

\[2(\mathbb Ef(X-h\nabla f(X))-\mathbb Ef(X))\le-h(2-hM)\|\nabla f(X)\|_{L^2}^2.\]\begin{leanbox}{Lean: \texttt{lp\_\allowbreak{}gradient\_\allowbreak{}descent}}
\leansource{\texttt{lp\_\allowbreak{}gradient\_\allowbreak{}descent} \quad AverageLMC/\allowbreak{}ExpectedGradientStep.lean, lines 37–63}
\begin{leancode}
\leanline{0}{theorem\ lp\_gradient\_descent\ (f\ :\ E\ →\ ℝ)\ (h\ :\ ℝ)\ (L\ :\ ℝ≥0)}
\leanline{4}{(hf\ :\ Differentiable\ ℝ\ f)\ (hLip\ :\ LipschitzWith\ L\ (gradient\ f))}
\leanline{4}{(hh\ :\ 0\ ≤\ h)\ (X\ :\ Lp\ E\ 2\ P)\ :}
\leanline{4}{2\ *\ ((∫\ ω,\ f\ ((X\ -\ h\ •\ ULA.lpDrift\ P\ (gradient\ f)\ hLip\ X)\ ω)\ ∂P)\ -}
\leanline{6}{∫\ ω,\ f\ (X\ ω)\ ∂P)\ ≤}
\leanline{6}{-h\ *\ (2\ -\ h\ *\ (L\ :\ ℝ))\ *\ ‖ULA.lpDrift\ P\ (gradient\ f)\ hLip\ X‖\ \textasciicircum{}\ 2}
\end{leancode}
\leanlabel{How the Lean proof follows the argument}
\texttt{lp\_\allowbreak{}gradient\_\allowbreak{}descent} gives the formal counterpart of this step. The main intermediate facts are \texttt{G}, \texttt{hi\allowbreak{}X}, \texttt{hi\allowbreak{}Y}, \texttt{hi\allowbreak{}G}, \texttt{hpoint}, \texttt{hb} in that order. Named proof applications include \texttt{gradient\_\allowbreak{}descent}, \texttt{hb}; the adjacent proof shows their hypotheses and order of use.
\end{leanbox}


Earlier results used here: \cite[Section~3.13.2]{LangevinLean2026ULASC}; \cite[Section~3.13.3]{LangevinLean2026ULASC}; \cite[Section~3.11.1]{LangevinLean2026IPLA}; \cite[Section~3.15.2]{LangevinLean2026Proximal}.
\section{\texorpdfstring{Couplings, the transport infimum and lifts of output couplings}{Couplings, the transport infimum and lifts of output couplings}}
\label{proof-module-AverageLMC-SquaredTransport-lean}


This section contains two general facts about quadratic transport that are used when a one-step bound is proved by constructing couplings. Let \(E\) be a normed space with its Borel sets (in the applications \(E=\mathbb R^p\)). For measures \(\mu,\pi\) on \(E\), a \emph{coupling} of \(\mu\) and \(\pi\) is a measure \(\gamma\) on \(E\times E\) whose first and second marginals are \(\mu\) and \(\pi\). Its cost is \(C(\gamma)=\bigl(\int\|x-y\|^2\,\gamma(dx,dy)\bigr)^{1/2}\in[0,\infty]\), and the Wasserstein distance is \(\mathsf W_2(\mu,\pi)=\inf_\gamma C(\gamma)\in[0,\infty]\), the infimum over all couplings. No optimal coupling is assumed to exist.




The first result passes a bound from couplings to the distance: if \(\mathsf W_2(\mu,\pi)<\infty\) and \(r\le a\,C(\gamma)^2+b\) for every coupling \(\gamma\) of finite cost, then \(r\le a\,\mathsf W_2(\mu,\pi)^2+b\), for every real coefficient \(a\). The second result lifts a coupling of an output law to the input space: if \(F:S\to E\) is measurable, \(P\) is a probability measure on \(S\) and \(\gamma\) couples the law \(P\circ F^{-1}\) of \(F\) with \(\pi\), then there is a probability measure \(R\) on \(S\times E\) with first marginal \(P\) under which the pair \((F(s),y)\) has law \(\gamma\). In the stochastic-gradient argument \(s\) is the input of one step, \(F(s)\) the full-gradient output, and the lift allows an independent random index to be adjoined to an arbitrary coupling of the full-gradient output with the target.



Earlier results used here: \cite[Section~3.22]{LangevinLean2026ULASC}.
\subsection{\texorpdfstring{From bounds over couplings to a bound on the squared distance}{From bounds over couplings to a bound on the squared distance}}
\label{report-step-AverageLMC-SquaredTransport-lean-1}


\textbf{Claim.} Let \(\mu,\pi\) be measures on \(E\) with \(\mathsf W_2(\mu,\pi)<\infty\), and write \(w=\mathsf W_2(\mu,\pi)\in[0,\infty)\). Let \(r,a,b\in\mathbb R\) and suppose that


\[r\le a\,C(\gamma)^2+b\qquad\text{for every coupling }\gamma\text{ of }\mu\text{ and }\pi\text{ with }C(\gamma)<\infty.\qquad(1)\]


Then \(r\le a\,w^2+b\). In particular, if \(\mu',\pi'\) are measures with \(\mathsf W_2(\mu',\pi')<\infty\) and \(\mathsf W_2(\mu',\pi')^2\le a\,C(\gamma)^2+b\) for every finite-cost coupling \(\gamma\) of \(\mu\) and \(\pi\), then \(\mathsf W_2(\mu',\pi')^2\le a\,\mathsf W_2(\mu,\pi)^2+b\).




\textbf{Proof.} \emph{Case \(a\ge0\).} Let \(\varepsilon>0\). Since \(w<w+\varepsilon\) and \(w\) is the infimum of the costs, some coupling \(\gamma\) satisfies \(C(\gamma)<w+\varepsilon\); in particular \(C(\gamma)<\infty\). Squaring the inequality between the nonnegative numbers \(C(\gamma)\) and \(w+\varepsilon\) gives \(C(\gamma)^2\le(w+\varepsilon)^2\), and multiplying by \(a\ge0\) and using (1),


\[r\le a\,C(\gamma)^2+b\le a\,(w+\varepsilon)^2+b.\]


This holds for every \(\varepsilon>0\). The function \(\varepsilon\mapsto a(w+\varepsilon)^2+b\) is continuous, so letting \(\varepsilon\downarrow0\) gives \(r\le a\,w^2+b\).




\emph{Case \(a<0\).} Since \(w<\infty\), not every coupling has infinite cost, for otherwise the infimum would be \(\infty\); hence there is a coupling \(\gamma\) with \(C(\gamma)<\infty\). By the definition of the infimum, \(w\le C(\gamma)\), hence \(w^2\le C(\gamma)^2\), and multiplying by \(a<0\) reverses this inequality: \(a\,C(\gamma)^2\le a\,w^2\). With (1), \(r\le a\,C(\gamma)^2+b\le a\,w^2+b\).




The second statement is the first with \(r=\mathsf W_2(\mu',\pi')^2\). \(\square\)




The finiteness assumption on \(\mathsf W_2(\mu,\pi)\) is needed in both cases: the approximation \(C(\gamma)<w+\varepsilon\) and the existence of a finite-cost coupling require a finite infimum. The case \(a<0\) occurs in Theorem 1 when \(1-mh<0\).

\begin{leanbox}{Lean: \texttt{real\_\allowbreak{}bound\_\allowbreak{}of\_\allowbreak{}finite\_\allowbreak{}couplings}, \texttt{wasserstein2\_\allowbreak{}sq\_\allowbreak{}bound\_\allowbreak{}of\_\allowbreak{}finite\_\allowbreak{}couplings}}
\leansource{\texttt{real\_\allowbreak{}bound\_\allowbreak{}of\_\allowbreak{}finite\_\allowbreak{}couplings} \quad AverageLMC/\allowbreak{}SquaredTransport.lean, lines 25–64}
\begin{leancode}
\leanline{0}{theorem\ real\_bound\_of\_finite\_couplings}
\leanline{4}{(μ\ π\ :\ Measure\ E)\ (r\ a\ b\ :\ ℝ)}
\leanline{4}{(hfinite\ :\ ULA.wasserstein2\ μ\ π\ ≠\ ⊤)}
\leanline{4}{(hfamily\ :\ ∀\ c\ :\ ULA.Coupling\ μ\ π,\ ULA.couplingCost\ c.joint\ ≠\ ⊤\ →}
\leanline{6}{r\ ≤\ a\ *\ (ULA.couplingCost\ c.joint).toReal\ \textasciicircum{}\ 2\ +\ b)\ :}
\leanline{4}{r\ ≤\ a\ *\ ULA.wasserstein2Real\ μ\ π\ \textasciicircum{}\ 2\ +\ b}
\end{leancode}
\leansource{\texttt{wasserstein2\_\allowbreak{}sq\_\allowbreak{}bound\_\allowbreak{}of\_\allowbreak{}finite\_\allowbreak{}couplings} \quad AverageLMC/\allowbreak{}SquaredTransport.lean, lines 68–77}
\begin{leancode}
\leanline{0}{theorem\ wasserstein2\_sq\_bound\_of\_finite\_couplings}
\leanline{4}{(μ\ π\ μ'\ π'\ :\ Measure\ E)\ (a\ b\ :\ ℝ)}
\leanline{4}{(hfinite\ :\ ULA.wasserstein2\ μ\ π\ ≠\ ⊤)}
\leanline{4}{(\_houtput\ :\ ULA.wasserstein2\ μ'\ π'\ ≠\ ⊤)}
\leanline{4}{(hfamily\ :\ ∀\ c\ :\ ULA.Coupling\ μ\ π,\ ULA.couplingCost\ c.joint\ ≠\ ⊤\ →}
\leanline{6}{ULA.wasserstein2Real\ μ'\ π'\ \textasciicircum{}\ 2\ ≤\ a\ *\ (ULA.couplingCost\ c.joint).toReal\ \textasciicircum{}\ 2\ +\ b)\ :}
\leanline{4}{ULA.wasserstein2Real\ μ'\ π'\ \textasciicircum{}\ 2\ ≤\ a\ *\ ULA.wasserstein2Real\ μ\ π\ \textasciicircum{}\ 2\ +\ b}
\end{leancode}
\leanlabel{How the Lean proof follows the argument}
The cost is \texttt{ULA.\allowbreak{}coupling\allowbreak{}Cost}, the \(L^2\) seminorm \texttt{e\allowbreak{}Lp\allowbreak{}Norm (fun p \ensuremath{\mapsto} p.\allowbreak{}1 -\allowbreak{} p.\allowbreak{}2) 2 \ensuremath{\gamma}} with values in \texttt{\ensuremath{\mathbb R}\ensuremath{\ge}0\ensuremath{\infty}} \(=[0,\infty]\), and \texttt{ULA.\allowbreak{}wasserstein2 \ensuremath{\mu} \ensuremath{\pi}} is the infimum \texttt{{\leanfallback ⨅} c : Coupling \ensuremath{\mu} \ensuremath{\pi}, coupling\allowbreak{}Cost c.\allowbreak{}joint} over the type of couplings (a joint measure \texttt{c.\allowbreak{}joint} with the two marginal identities). \texttt{ULA.\allowbreak{}wasserstein2\allowbreak{}Real} is its conversion \texttt{to\allowbreak{}Real} to a real number, which agrees with \(\mathsf W_2\) when the latter is finite. Hypothesis (1) is \texttt{hfamily}, and finiteness is \texttt{hfinite : ULA.\allowbreak{}wasserstein2 \ensuremath{\mu} \ensuremath{\pi} ≠ \ensuremath{\top}}.



\begin{enumerate}


\item In \texttt{real\_\allowbreak{}bound\_\allowbreak{}of\_\allowbreak{}finite\_\allowbreak{}couplings}, \texttt{by\_\allowbreak{}cases ha : 0 \ensuremath{\le} a} separates the two cases. For \(a\ge0\), \texttt{heps} proves the bound with \(w+\varepsilon\): \texttt{hi} places \(\mathsf W_2\) strictly below \texttt{ENNReal.\allowbreak{}of\allowbreak{}Real (w + \ensuremath{\varepsilon})} (using \texttt{ENNReal.\allowbreak{}of\allowbreak{}Real\_\allowbreak{}to\allowbreak{}Real} and \texttt{ENNReal.\allowbreak{}of\allowbreak{}Real\_\allowbreak{}lt\_\allowbreak{}of\allowbreak{}Real\_\allowbreak{}iff}), \texttt{i\allowbreak{}Inf\_\allowbreak{}lt\_\allowbreak{}iff} extracts a coupling \texttt{c} with smaller cost, \texttt{hreal} converts this to real numbers with \texttt{ENNReal.\allowbreak{}to\allowbreak{}Real\_\allowbreak{}mono} and \texttt{ENNReal.\allowbreak{}to\allowbreak{}Real\_\allowbreak{}of\allowbreak{}Real}, and \texttt{sq\_\allowbreak{}le\_\allowbreak{}sq\ensuremath{{}_{0}}} squares it. \texttt{hlim} is the continuity of \(\varepsilon\mapsto a(w+\varepsilon)^2+b\) at \(0\) along \(\varepsilon>0\) (the filter \texttt{nhds\allowbreak{}Within 0 (Ioi 0)}), and \texttt{le\_\allowbreak{}of\_\allowbreak{}tendsto\_\allowbreak{}of\_\allowbreak{}tendsto} passes to the limit.

\item For \(a<0\), \texttt{ULA.\allowbreak{}exists\_\allowbreak{}finite\allowbreak{}Cost\_\allowbreak{}coupling} gives a finite-cost coupling, \texttt{ULA.\allowbreak{}wasserstein2\_\allowbreak{}le\_\allowbreak{}coupling\allowbreak{}Cost} gives \(w\le C(\gamma)\), and \texttt{mul\_\allowbreak{}le\_\allowbreak{}mul\_\allowbreak{}of\_\allowbreak{}nonpos\_\allowbreak{}left} multiplies the squared inequality by \(a\le0\).

\item \texttt{wasserstein2\_\allowbreak{}sq\_\allowbreak{}bound\_\allowbreak{}of\_\allowbreak{}finite\_\allowbreak{}couplings} applies the first theorem with \(r=\mathsf W_2(\mu',\pi')^2\). Its hypothesis \texttt{\_\allowbreak{}houtput} (finiteness of \(\mathsf W_2(\mu',\pi')\)) is not needed by the proof; it records that the real value of the output distance is the true distance.

\end{enumerate}
\end{leanbox}


Earlier results used here: \cite[Section~3.17.1]{LangevinLean2026ULASC}; \cite[Section~3.17.2]{LangevinLean2026ULASC}; \cite[Section~3.17.3]{LangevinLean2026ULASC}; \cite[Section~3.17.4]{LangevinLean2026ULASC}; \cite[Section~3.17.5]{LangevinLean2026ULASC}; \cite[Section~3.18.2]{LangevinLean2026ULASC}.
\subsection{\texorpdfstring{Lifting an output coupling to the input space}{Lifting an output coupling to the input space}}
\label{report-step-AverageLMC-SquaredTransport-lean-2}


\textbf{Setting.} Let \((S,\mathcal S)\) be a measurable space with a probability measure \(P\), let \(E\) be a standard Borel space (for instance \(\mathbb R^p\)), let \(F:S\to E\) be measurable, and let \(\gamma\) be a coupling of the law \(P\circ F^{-1}\) of \(F\) and a measure \(\pi\) on \(E\). Since the first marginal of \(\gamma\) is a probability measure, \(\gamma\) is a probability measure. Because \(E\) is standard Borel, \(\gamma\) admits a disintegration with respect to its first coordinate: a Markov kernel \(K\) from \(E\) to \(E\) (each \(K(x,\cdot)\) is a probability measure and \(x\mapsto K(x,B)\) is measurable) such that


\[\begin{aligned}
\gamma(C)&=\int_E K(x,C_x)\,(P\circ F^{-1})(dx),\\
C_x&=\{y:(x,y)\in C\},\qquad C\subseteq E\times E\text{ measurable}.
\qquad(2)
\end{aligned}\]


Define the measure \(R\) on \(S\times E\) by


\[R(G)=\int_S K\bigl(F(s),G_s\bigr)\,P(ds),\qquad G_s=\{y:(s,y)\in G\},\]


that is, \(s\) is drawn from \(P\) and then \(y\) from \(K(F(s),\cdot)\).




\textbf{Claim.} (1) \(R\) is a probability measure. (2) The first marginal of \(R\) is \(P\). (3) The image of \(R\) under \((s,y)\mapsto(F(s),y)\) is \(\gamma\). (4) The second marginal of \(R\) is \(\pi\).




\textbf{Proof.} (1) and (2): for measurable \(A\subseteq S\), \(R(A\times E)=\int_AK(F(s),E)\,P(ds)=P(A)\), because \(K(F(s),E)=1\). Taking \(A=S\) gives \(R(S\times E)=1\).




(3) Let \(C\subseteq E\times E\) be measurable and \(G=\{(s,y):(F(s),y)\in C\}\). Then \(G_s=C_{F(s)}\), so


\[R(G)=\int_SK\bigl(F(s),C_{F(s)}\bigr)\,P(ds)=\int_EK(x,C_x)\,(P\circ F^{-1})(dx)=\gamma(C),\]


where the middle equality is the change-of-variables formula for the image measure \(P\circ F^{-1}\) applied to the measurable function \(x\mapsto K(x,C_x)\), and the last equality is (2).




(4) The second coordinate of \((F(s),y)\) is \(y\). Taking the image of the measure in (3) under the second projection, the second marginal of \(R\) equals the second marginal of \(\gamma\), which is \(\pi\). \(\square\)




In probabilistic terms: on the probability space \((S\times E,R)\), the coordinate \(s\) has law \(P\), the coordinate \(y\) has law \(\pi\), and \((F(s),y)\) has the prescribed joint law \(\gamma\). Every coupling of the output law with the target is therefore realized jointly with the input of the step.

\begin{leanbox}{Lean: \texttt{liftedOutputCoupling}, \texttt{liftedOutputCoupling\_\allowbreak{}probability}, \texttt{liftedOutputCoupling\_\allowbreak{}fst} and 2 more}
\leansource{\texttt{liftedOutputCoupling} \quad AverageLMC/\allowbreak{}SquaredTransport.lean, lines 82–86}
\begin{leancode}
\leanline{0}{def\ liftedOutputCoupling\ (P\ :\ Measure\ S)\ [IsProbabilityMeasure\ P]}
\leanline{4}{(F\ :\ S\ →\ E)\ (hF\ :\ Measurable\ F)\ \{π\ :\ Measure\ E\}}
\leanline{4}{(c\ :\ ULA.Coupling\ (P.map\ F)\ π)\ :\ Measure\ (S\ ×\ E)\ :=}
\leanline{2}{have\ :=\ c.isProbabilityMeasure}
\leanline{2}{P\ ⊗ₘ\ (c.joint.condKernel.comap\ F\ hF)}
\end{leancode}
\leansource{\texttt{liftedOutputCoupling\_\allowbreak{}probability} \quad AverageLMC/\allowbreak{}SquaredTransport.lean, lines 88–93}
\begin{leancode}
\leanline{0}{instance\ liftedOutputCoupling\_probability\ (P\ :\ Measure\ S)\ [IsProbabilityMeasure\ P]}
\leanline{4}{(F\ :\ S\ →\ E)\ (hF\ :\ Measurable\ F)\ \{π\ :\ Measure\ E\}}
\leanline{4}{(c\ :\ ULA.Coupling\ (P.map\ F)\ π)\ :}
\leanline{4}{IsProbabilityMeasure\ (liftedOutputCoupling\ P\ F\ hF\ c)\ :=\ by}
\leanline{2}{unfold\ liftedOutputCoupling}
\leanline{2}{infer\_instance}
\end{leancode}
\leansource{\texttt{liftedOutputCoupling\_\allowbreak{}fst} \quad AverageLMC/\allowbreak{}SquaredTransport.lean, lines 95–100}
\begin{leancode}
\leanline{0}{theorem\ liftedOutputCoupling\_fst\ (P\ :\ Measure\ S)\ [IsProbabilityMeasure\ P]}
\leanline{4}{(F\ :\ S\ →\ E)\ (hF\ :\ Measurable\ F)\ \{π\ :\ Measure\ E\}}
\leanline{4}{(c\ :\ ULA.Coupling\ (P.map\ F)\ π)\ :}
\leanline{4}{(liftedOutputCoupling\ P\ F\ hF\ c).map\ Prod.fst\ =\ P}
\end{leancode}
\leansource{\texttt{liftedOutputCoupling\_\allowbreak{}output} \quad AverageLMC/\allowbreak{}SquaredTransport.lean, lines 104–119}
\begin{leancode}
\leanline{0}{theorem\ liftedOutputCoupling\_output\ (P\ :\ Measure\ S)\ [IsProbabilityMeasure\ P]}
\leanline{4}{(F\ :\ S\ →\ E)\ (hF\ :\ Measurable\ F)\ \{π\ :\ Measure\ E\}}
\leanline{4}{(c\ :\ ULA.Coupling\ (P.map\ F)\ π)\ :}
\leanline{4}{(liftedOutputCoupling\ P\ F\ hF\ c).map\ (fun\ q\ ↦\ (F\ q.1,\ q.2))\ =\ c.joint}
\end{leancode}
\leansource{\texttt{liftedOutputCoupling\_\allowbreak{}snd} \quad AverageLMC/\allowbreak{}SquaredTransport.lean, lines 121–128}
\begin{leancode}
\leanline{0}{theorem\ liftedOutputCoupling\_snd\ (P\ :\ Measure\ S)\ [IsProbabilityMeasure\ P]}
\leanline{4}{(F\ :\ S\ →\ E)\ (hF\ :\ Measurable\ F)\ \{π\ :\ Measure\ E\}}
\leanline{4}{(c\ :\ ULA.Coupling\ (P.map\ F)\ π)\ :}
\leanline{4}{(liftedOutputCoupling\ P\ F\ hF\ c).map\ Prod.snd\ =\ π}
\end{leancode}
\leanlabel{How the Lean proof follows the argument}
\begin{enumerate}


\item \texttt{lifted\allowbreak{}Output\allowbreak{}Coupling P F h\allowbreak{}F c} is \texttt{P ⊗ₘ (c.\allowbreak{}joint.\allowbreak{}cond\allowbreak{}Kernel.\allowbreak{}comap F h\allowbreak{}F)}. Here \texttt{c.\allowbreak{}joint.\allowbreak{}cond\allowbreak{}Kernel} is Mathlib's conditional kernel of a measure on a product whose second factor is standard Borel (the kernel \(K\) of (2)), \texttt{.\allowbreak{}comap F h\allowbreak{}F} precomposes it with the measurable map \(F\) (the kernel \(s\mapsto K(F(s),\cdot)\)), and \texttt{⊗ₘ} is the composition-product of a measure with a kernel, which is the definition of \(R\). The line \texttt{have := c.\allowbreak{}is\allowbreak{}Probability\allowbreak{}Measure} supplies the fact that the coupling is a probability measure, required for the conditional kernel.

\item \texttt{lifted\allowbreak{}Output\allowbreak{}Coupling\_\allowbreak{}probability} is an \texttt{instance}, so that later statements may use \(R\) as a probability measure; \texttt{infer\_\allowbreak{}instance} finds it because a composition-product of a probability measure with a Markov kernel is a probability measure. This is claim (1).

\item \texttt{lifted\allowbreak{}Output\allowbreak{}Coupling\_\allowbreak{}fst} is claim (2), the Mathlib identity \texttt{Measure.\allowbreak{}fst\_\allowbreak{}comp\allowbreak{}Prod}.

\item \texttt{lifted\allowbreak{}Output\allowbreak{}Coupling\_\allowbreak{}output} is claim (3). \texttt{hfst} records that the first marginal of the coupling is \texttt{P.\allowbreak{}map F}, and \texttt{hdis := c.\allowbreak{}joint.\allowbreak{}disintegrate c.\allowbreak{}joint.\allowbreak{}cond\allowbreak{}Kernel} is the disintegration identity (2), rewritten with \texttt{hfst}. After \texttt{ext A h\allowbreak{}A} the two measures are compared on a measurable set: \texttt{Measure.\allowbreak{}map\_\allowbreak{}apply}, \texttt{Measure.\allowbreak{}comp\allowbreak{}Prod\_\allowbreak{}apply} (twice) and \texttt{lintegral\_\allowbreak{}map} (the change of variables under \(F\)) reduce both sides to the same integral, closed by \texttt{rfl}.

\item \texttt{lifted\allowbreak{}Output\allowbreak{}Coupling\_\allowbreak{}snd} is claim (4): it maps claim (3) by \texttt{Prod.\allowbreak{}snd} (\texttt{congr\allowbreak{}Arg}), composes the two maps with \texttt{Measure.\allowbreak{}map\_\allowbreak{}map}, and uses the marginal identity \texttt{c.\allowbreak{}snd}.

\end{enumerate}
\end{leanbox}


Earlier results used here: \cite[Section~3.17.1]{LangevinLean2026ULASC}.
\section{\texorpdfstring{Construct the stationary coupling of one actual Gaussian LMC step}{Construct the stationary coupling of one actual Gaussian LMC step}}
\label{proof-module-AverageLMC-GaussianCanonical-lean}


The theorem starts from an arbitrary finite-cost initial coupling and the original smooth strongly convex potential. It constructs the independent Brownian extension, the actual integral solution, its stationary Gibbs law, and the exact Gaussian output law. The only local auxiliary input is the coordinate smoothing theorem proved separately.



Let \(\gamma\) be any finite-cost coupling of \(\nu\) and \(\pi\), with \(C_\gamma^2=\int\|x-y\|^2\,d\gamma\). The canonical independent Brownian extension realizes the actual update


\[\vartheta^+=\vartheta-h\nabla f(\vartheta)+\sqrt{2h}\,\xi,\qquad\xi\sim N(0,I_p),\quad\xi\perp\!\!\!\perp\vartheta.\]
\[(1+mh)W_2^2(\nu^+,\pi)\le(1-mh)C_\gamma^2+(M_{\rm av}+m)h^2p+h^2(hL-1)\int\|\nabla f(x)\|^2\,d\nu(x).\]\subsection{\texorpdfstring{Construct, identify and estimate the actual transition}{Construct, identify and estimate the actual transition}}
\label{report-step-AverageLMC-GaussianCanonical-lean-1}


Extend the initial coupling by the explicit canonical product Brownian probability space. Its initial algorithm state \(X\) has law \(\mu\) and its reference state \(Z\) has the true Gibbs law. The Lipschitz drift constructs a continuous \(L^{2}\) Langevin integral solution; stationarity follows from the already proved Gibbs-stationarity theorem. Independent scalar Brownian paths provide independent standard Gaussian innovations at every positive time. The algorithm center \(A=X-h\nabla f(X)\) is a measurable function of the initial state, so it remains independent of those innovations. The endpoint \(A+B(h)\) has exactly the ordinary Gaussian LMC transition law. Local coordinate smoothing at \(A\) gives the affine potential bound needed by the Gaussian energy theorem. Integrated strong convexity controls \(\|A-Z\|^{2}\), and integrated deterministic descent controls \(\mathbb E f(A)-\mathbb E f(X)\). These potential terms cancel. Identifying the lifted gradient norm with the score moment under \(\mu\) leaves precisely \((1-mh)\) times the initial coupling cost squared, \((M_{\mathrm{av}}+m)h^{2}d\), and \(h^2(hL-1)\) times the score second moment. Output transport finiteness follows from the constructed \(L^{2}\) coupling before any real-valued distance is used.

\begin{leanbox}{Lean: \texttt{lmc\_\allowbreak{}coupling\_\allowbreak{}step}}
\leansource{\texttt{lmc\_\allowbreak{}coupling\_\allowbreak{}step} \quad AverageLMC/\allowbreak{}GaussianCanonical.lean, lines 19–195}
\begin{leancode}
\leanline{0}{theorem\ lmc\_coupling\_step\ (f\ :\ EuclideanSpace\ ℝ\ ι\ →\ ℝ)\ (m\ h\ Mav\ :\ ℝ)\ (L\ :\ ℝ≥0)}
\leanline{4}{(hm\ :\ 0\ <\ m)\ (hh\ :\ 0\ <\ h)\ (hf\ :\ Differentiable\ ℝ\ f)}
\leanline{4}{(hsc\ :\ StrongConvexOn\ univ\ m\ f)\ (hLip\ :\ LipschitzWith\ L\ (gradient\ f))}
\leanline{4}{(hsm\ :\ ∀\ x\ s,\ 0\ ≤\ s\ →\ (∫\ z,\ f\ (x\ +\ Real.sqrt\ (2\ *\ s)\ •\ z)\ -\ f\ x}
\leanline{6}{∂stdGaussian\ (EuclideanSpace\ ℝ\ ι))\ ≤\ Mav\ *\ (Fintype.card\ ι\ :\ ℝ)\ *\ s)}
\leanline{4}{(μ\ :\ Measure\ (EuclideanSpace\ ℝ\ ι))\ [IsProbabilityMeasure\ μ]}
\leanline{4}{(c\ :\ Coupling\ μ\ (gibbsMeasure\ f))\ (hc\ :\ couplingCost\ c.joint\ ≠\ ⊤)\ :}
\leanline{4}{wasserstein2\ (ulaLawStep\ f\ h\ μ)\ (gibbsMeasure\ f)\ ≠\ ⊤\ ∧}
\leanline{4}{(1\ +\ m\ *\ h)\ *\ wasserstein2Real\ (ulaLawStep\ f\ h\ μ)\ (gibbsMeasure\ f)\ \textasciicircum{}\ 2\ ≤}
\leanline{6}{(1\ -\ m\ *\ h)\ *\ (couplingCost\ c.joint).toReal\ \textasciicircum{}\ 2\ +}
\leanline{6}{(Mav\ +\ m)\ *\ h\ \textasciicircum{}\ 2\ *\ (Fintype.card\ ι\ :\ ℝ)\ +}
\leanline{6}{h\ \textasciicircum{}\ 2\ *\ (h\ *\ (L\ :\ ℝ)\ -\ 1)\ *\ (∫\ x,\ ‖gradient\ f\ x‖\ \textasciicircum{}\ 2\ ∂μ)}
\end{leancode}
\leanlabel{How the Lean proof follows the argument}
\texttt{lmc\_\allowbreak{}coupling\_\allowbreak{}step} gives the formal counterpart of this step. The main intermediate facts are \texttt{Q}, \texttt{P}, \texttt{h\ensuremath{\pi}}, \texttt{X}, \texttt{Z}, \texttt{W} in that order. Named proof applications include \texttt{lp\allowbreak{}Ula\allowbreak{}Update\_\allowbreak{}coe}, \texttt{ne\_\allowbreak{}top\_\allowbreak{}of\_\allowbreak{}le\_\allowbreak{}ne\_\allowbreak{}top}, \texttt{h\allowbreak{}Ye}; the adjacent proof shows their hypotheses and order of use.
\end{leanbox}


Earlier results used here: \cite[Section~3.13.12]{LangevinLean2026ULASC}; \cite[Section~3.13.2]{LangevinLean2026ULASC}; \cite[Section~3.13.3]{LangevinLean2026ULASC}; \cite[Section~3.13.9]{LangevinLean2026ULASC}; \cite[Section~3.17.1]{LangevinLean2026ULASC}; \cite[Section~3.17.12]{LangevinLean2026ULASC}; \cite[Section~3.17.2]{LangevinLean2026ULASC}; \cite[Section~3.17.3]{LangevinLean2026ULASC}; \cite[Section~3.17.4]{LangevinLean2026ULASC}; \cite[Section~3.22.17]{LangevinLean2026ULASC}; \cite[Section~3.22.18]{LangevinLean2026ULASC}; \cite[Section~3.22.19]{LangevinLean2026ULASC}; \cite[Section~3.22.2]{LangevinLean2026ULASC}; \cite[Section~3.22.20]{LangevinLean2026ULASC}; \cite[Section~3.22.21]{LangevinLean2026ULASC}; \cite[Section~3.22.22]{LangevinLean2026ULASC}; \cite[Section~3.22.23]{LangevinLean2026ULASC}; \cite[Section~3.22.3]{LangevinLean2026ULASC}; \cite[Section~3.22.7]{LangevinLean2026ULASC}; \cite[Section~3.22.8]{LangevinLean2026ULASC}; \cite[Section~3.23.1]{LangevinLean2026ULASC}; \cite[Section~3.23.2]{LangevinLean2026ULASC}; \cite[Section~3.25.1]{LangevinLean2026ULASC}; \cite[Section~3.25.2]{LangevinLean2026ULASC}; \cite[Section~3.25.6]{LangevinLean2026ULASC}; \cite[Section~3.26.3]{LangevinLean2026ULASC}; \cite[Section~3.29.7]{LangevinLean2026ULASC}; \cite[Section~3.29.8]{LangevinLean2026ULASC}; \cite[Section~3.32.12]{LangevinLean2026ULASC}; \cite[Section~3.32.13]{LangevinLean2026ULASC}; \cite[Section~3.32.15]{LangevinLean2026ULASC}; \cite[Section~3.32.16]{LangevinLean2026ULASC}; \cite[Section~3.32.17]{LangevinLean2026ULASC}; \cite[Section~3.32.19]{LangevinLean2026ULASC}; \cite[Section~3.36.7]{LangevinLean2026ULASC}; \cite[Section~3.42.6]{LangevinLean2026ULASC}; \cite[Section~3.42.7]{LangevinLean2026ULASC}; \cite[Section~3.8.11]{LangevinLean2026ULASC}; \cite[Section~3.8.16]{LangevinLean2026ULASC}; \cite[Section~3.8.17]{LangevinLean2026ULASC}; \cite[Section~3.8.19]{LangevinLean2026ULASC}; \cite[Section~3.8.20]{LangevinLean2026ULASC}; \cite[Section~3.8.21]{LangevinLean2026ULASC}; \cite[Section~3.8.22]{LangevinLean2026ULASC}.
\section{\texorpdfstring{Smoothing by arbitrary independent centered coordinates}{Smoothing by arbitrary independent centered coordinates}}
\label{proof-module-AverageLMC-Smoothing-lean}


This module proves the full expectation conclusion of the manuscript smoothing lemma. The coordinates may have any distributions with zero mean and unit second moment; Gaussianity is used only in the final specialization. The bound uses the exact coordinate Taylor suprema and their arithmetic average. All potential, score, and product integrability follows from the supplied coordinate second moments and the gradient Lipschitz assumption.



Earlier results used here: \cite[Section~3.13]{LangevinLean2026ULASC}; \cite[Section~3.4]{LangevinLean2026ULASC}.
\subsection{\texorpdfstring{Build square-integrable partial coordinate sums}{Build square-integrable partial coordinate sums}}
\label{report-step-AverageLMC-Smoothing-lean-1}


For a finite set \(S\) of coordinate indices, define \(Z_S\) as the sum of \(\xi_j\) \(e_j\) over \(j\) in \(S\). Each scalar coordinate is square integrable, and the map \(a\) to \(ae_j\) is continuous linear, so every vector summand and every finite partial sum is square integrable. The actual gradient evaluated at an \(L^{2}\) state is also \(L^{2}\): apply the existing Lipschitz drift lift to the canonical \(L^{2}\) representative, then transfer back along its almost-everywhere equality. Adding the deterministic center and scaling preserve these moments.

\begin{leanbox}{Lean: \texttt{memLp\_\allowbreak{}gradient\_\allowbreak{}comp}, \texttt{coordinatePartial}, \texttt{memLp\_\allowbreak{}coordinatePartial}}
\leansource{\texttt{memLp\_\allowbreak{}gradient\_\allowbreak{}comp} \quad AverageLMC/\allowbreak{}Smoothing.lean, lines 28–35}
\begin{leancode}
\leanline{0}{private\ theorem\ memLp\_gradient\_comp\ \{p\ :\ ℕ\}\ (μ\ :\ Measure\ Ω)\ [IsProbabilityMeasure\ μ]}
\leanline{4}{(f\ :\ State\ p\ →\ ℝ)\ (M\ :\ ℝ≥0)\ (hLip\ :\ LipschitzWith\ M\ (gradient\ f))}
\leanline{4}{(X\ :\ Ω\ →\ State\ p)\ (hX\ :\ MemLp\ X\ 2\ μ)\ :}
\leanline{4}{MemLp\ (fun\ ω\ ↦\ gradient\ f\ (X\ ω))\ 2\ μ}
\end{leancode}
\leansource{\texttt{coordinatePartial} \quad AverageLMC/\allowbreak{}Smoothing.lean, lines 37–39}
\begin{leancode}
\leanline{0}{private\ def\ coordinatePartial\ \{p\ :\ ℕ\}\ (ξ\ :\ Fin\ p\ →\ Ω\ →\ ℝ)}
\leanline{4}{(s\ :\ Finset\ (Fin\ p))\ (ω\ :\ Ω)\ :\ State\ p\ :=}
\leanline{2}{∑\ j\ ∈\ s,\ ξ\ j\ ω\ •\ coordinateVector\ j}
\end{leancode}
\leansource{\texttt{memLp\_\allowbreak{}coordinatePartial} \quad AverageLMC/\allowbreak{}Smoothing.lean, lines 41–46}
\begin{leancode}
\leanline{0}{private\ theorem\ memLp\_coordinatePartial\ \{p\ :\ ℕ\}\ (μ\ :\ Measure\ Ω)}
\leanline{4}{(ξ\ :\ Fin\ p\ →\ Ω\ →\ ℝ)\ (hξ\ :\ ∀\ j,\ MemLp\ (ξ\ j)\ 2\ μ)\ (s\ :\ Finset\ (Fin\ p))\ :}
\leanline{4}{MemLp\ (coordinatePartial\ ξ\ s)\ 2\ μ}
\end{leancode}
\leanlabel{How the Lean proof follows the argument}
\texttt{mem\allowbreak{}Lp\_\allowbreak{}gradient\_\allowbreak{}comp} gives the formal counterpart of this step. The main intermediate facts are \texttt{h} in that order. Named proof applications include \texttt{h.\allowbreak{}ae\_\allowbreak{}eq}; the adjacent proof shows their hypotheses and order of use.




\texttt{coordinate\allowbreak{}Partial} names the mathematical object defined above. Its defining expression appears in the adjacent statement; subsequent identities may unfold this definition directly.




\texttt{mem\allowbreak{}Lp\_\allowbreak{}coordinate\allowbreak{}Partial} gives the formal counterpart of this step. Named proof applications include \texttt{mem\allowbreak{}Lp\_\allowbreak{}finset\allowbreak{}Sum}; the adjacent proof shows their hypotheses and order of use.
\end{leanbox}


Earlier results used here: \cite[Section~3.13.4]{LangevinLean2026ULASC}.
\subsection{\texorpdfstring{Insert one independent coordinate and eliminate its linear term}{Insert one independent coordinate and eliminate its linear term}}
\label{report-step-AverageLMC-Smoothing-lean-2}


Induct on the finite coordinate set. When inserting \(i\) outside \(S\), apply the exact coordinate model at \(x+c\) \(Z_S\) with displacement \(c\) \(\xi_i\). It gives

\[f(x+cZ_{S\cup\{i\}})\le f(x+cZ_S)+c\xi_i\partial_i f(x+cZ_S)+\tfrac12 M_i c^2\xi_i^2.\]

Independence of the scalar coordinates survives the maps \(a\) to \(ae_j\). Consequently their partial vector sum is independent of the new coordinate. Composing with the actual gradient coordinate proves independence of \(\xi_i\) and the displayed derivative term. Their product is integrable by the \(L^{2}\) Holder inequality, and its integral factors; the zero first moment makes it vanish. The quadratic term integrates to \(M_i\) \(c^2/2\). Adding this to the inductive bound proves the finite-set estimate, with no assumed stochastic error inequality.




\textbf{Induction in detail.} At the empty set the perturbation is zero and the expected potential is \(f(x)\). Suppose the bound holds for \(S\), and choose \(i\notin S\). Set \(G_i=\partial_i f(x+cZ_S)\). Global gradient Lipschitz continuity and the finite second moment of \(Z_S\) imply \(G_i\in L^2\); hence Cauchy–Schwarz makes \(\xi_iG_i\) integrable. The variable \(G_i\) is a measurable function of the coordinates in \(S\), so independence gives \(\mathbb E[\xi_iG_i]=(\mathbb E\xi_i)(\mathbb EG_i)=0\). The coordinate model then gives the previous expectation plus \(M_ic^2/2\), since \(\mathbb E\xi_i^2=1\). The induction hypothesis and \(\sum_{j\in S\cup\{i\}}M_j=\sum_{j\in S}M_j+M_i\) complete the step. This holds for every real \(c\), including \(c=0\), and for arbitrary independent centered unit-second-moment coordinates with measurable representatives. Gaussianity is used only in the later specialization. \(\square\)

\begin{leanbox}{Lean: \texttt{integral\_\allowbreak{}coordinatePartial\_\allowbreak{}le}}
\leansource{\texttt{integral\_\allowbreak{}coordinatePartial\_\allowbreak{}le} \quad AverageLMC/\allowbreak{}Smoothing.lean, lines 50–117}
\begin{leancode}
\leanline{0}{theorem\ integral\_coordinatePartial\_le\ \{p\ :\ ℕ\}\ (μ\ :\ Measure\ Ω)\ [IsProbabilityMeasure\ μ]}
\leanline{4}{(f\ :\ State\ p\ →\ ℝ)\ (M\ :\ ℝ≥0)\ (hf\ :\ Differentiable\ ℝ\ f)}
\leanline{4}{(hLip\ :\ LipschitzWith\ M\ (gradient\ f))}
\leanline{4}{(ξ\ :\ Fin\ p\ →\ Ω\ →\ ℝ)\ (hξm\ :\ ∀\ j,\ Measurable\ (ξ\ j))}
\leanline{4}{(hξ\ :\ ∀\ j,\ MemLp\ (ξ\ j)\ 2\ μ)\ (hind\ :\ iIndepFun\ ξ\ μ)}
\leanline{4}{(hmean\ :\ ∀\ j,\ ∫\ ω,\ ξ\ j\ ω\ ∂μ\ =\ 0)}
\leanline{4}{(hvar\ :\ ∀\ j,\ ∫\ ω,\ ξ\ j\ ω\ \textasciicircum{}\ 2\ ∂μ\ =\ 1)}
\leanline{4}{(x\ :\ State\ p)\ (c\ :\ ℝ)\ (s\ :\ Finset\ (Fin\ p))\ :}
\leanline{4}{(∫\ ω,\ f\ (x\ +\ c\ •\ coordinatePartial\ ξ\ s\ ω)\ ∂μ)\ ≤}
\leanline{6}{f\ x\ +\ c\ \textasciicircum{}\ 2\ /\ 2\ *\ ∑\ j\ ∈\ s,\ coordinateSmoothness\ f\ j}
\end{leancode}
\leanlabel{How the Lean proof follows the argument}
\texttt{integral\_\allowbreak{}coordinate\allowbreak{}Partial\_\allowbreak{}le} derives the partial-state moments \texttt{h\allowbreak{}X} and potential integrability \texttt{h\allowbreak{}F} before the finite-set induction. In the insert branch \texttt{G} is the new score coordinate and \texttt{h\allowbreak{}G} proves its square integrability. \texttt{h\allowbreak{}IG} obtains its independence from the new coordinate by composing independence of the finite partial sum with the score function. \texttt{hprod} supplies product integrability, and \texttt{hzero} factors the integral and uses centering. The coordinate Taylor model is then integrated, its squared-coordinate integral is replaced by the unit-moment hypothesis, and the induction hypothesis supplies the sum over the old coordinates.
\end{leanbox}
\subsection{\texorpdfstring{Sum every coordinate and recover the exact average-smoothness coefficient}{Sum every coordinate and recover the exact average-smoothness coefficient}}
\label{report-step-AverageLMC-Smoothing-lean-3}


For the full index set, the coordinate sum is exactly the Euclidean vector with coordinates \(\xi_j\). Choose \(c\) equal to the square root of \(2s\), whose square is \(2s\) because \(s\) is nonnegative. Smooth potentials have quadratic growth, so the perturbed potential and its difference from the deterministic center are integrable. The finite-set theorem yields

\[\mathbb E[f(x+\sqrt{2s}\,\xi)-f(x)]\le s\sum_j M_j=M_{\rm av}ps.\]

The final equality uses the literal definition \(M_{\mathrm{av}}=p^{-1}\sum_jM_j\) and the positive dimension.

\begin{leanbox}{Lean: \texttt{integral\_\allowbreak{}coordinate\_\allowbreak{}smoothing}}
\leansource{\texttt{integral\_\allowbreak{}coordinate\_\allowbreak{}smoothing} \quad AverageLMC/\allowbreak{}Smoothing.lean, lines 121–151}
\begin{leancode}
\leanline{0}{theorem\ integral\_coordinate\_smoothing\ \{p\ :\ ℕ\}\ (hp\ :\ 0\ <\ p)}
\leanline{4}{(μ\ :\ Measure\ Ω)\ [IsProbabilityMeasure\ μ]}
\leanline{4}{(f\ :\ State\ p\ →\ ℝ)\ (M\ :\ ℝ≥0)\ (hf\ :\ Differentiable\ ℝ\ f)}
\leanline{4}{(hLip\ :\ LipschitzWith\ M\ (gradient\ f))}
\leanline{4}{(ξ\ :\ Fin\ p\ →\ Ω\ →\ ℝ)\ (hξm\ :\ ∀\ j,\ Measurable\ (ξ\ j))}
\leanline{4}{(hξ\ :\ ∀\ j,\ MemLp\ (ξ\ j)\ 2\ μ)\ (hind\ :\ iIndepFun\ ξ\ μ)}
\leanline{4}{(hmean\ :\ ∀\ j,\ ∫\ ω,\ ξ\ j\ ω\ ∂μ\ =\ 0)}
\leanline{4}{(hvar\ :\ ∀\ j,\ ∫\ ω,\ ξ\ j\ ω\ \textasciicircum{}\ 2\ ∂μ\ =\ 1)}
\leanline{4}{(x\ :\ State\ p)\ (s\ :\ ℝ)\ (hs\ :\ 0\ ≤\ s)\ :}
\leanline{4}{(∫\ ω,\ f\ (x\ +\ Real.sqrt\ (2\ *\ s)\ •\ WithLp.toLp\ 2\ (fun\ j\ ↦\ ξ\ j\ ω))\ -\ f\ x\ ∂μ)\ ≤}
\leanline{6}{averageSmoothness\ f\ *\ (p\ :\ ℝ)\ *\ s}
\end{leancode}
\leanlabel{How the Lean proof follows the argument}
\texttt{integral\_\allowbreak{}coordinate\_\allowbreak{}smoothing} gives the formal counterpart of this step. The main intermediate facts are \texttt{hi}, \texttt{heq}, \texttt{h\allowbreak{}X}, \texttt{h\allowbreak{}F}, \texttt{hsum} in that order. Its remaining equalities are discharged by the simplifications displayed in the proof.
\end{leanbox}
\subsection{\texorpdfstring{Specialize to the actual standard Gaussian measure}{Specialize to the actual standard Gaussian measure}}
\label{report-step-AverageLMC-Smoothing-lean-4}


Take the product of \(n=p\) real Gaussian laws with mean zero and variance one. Each coordinate evaluation preserves its marginal Gaussian measure; the product coordinates are mutually independent, square integrable, centered, and have second moment one. Apply the arbitrary-coordinate theorem on this product. Its pushforward is identified by the Euclidean coordinate map with the actual standard Gaussian law, so the change-of-measure integral is precisely the Gaussian smoothing integral used by the LMC transition.

\begin{leanbox}{Lean: \texttt{integral\_\allowbreak{}stdGaussian\_\allowbreak{}smoothing}}
\leansource{\texttt{integral\_\allowbreak{}stdGaussian\_\allowbreak{}smoothing} \quad AverageLMC/\allowbreak{}Smoothing.lean, lines 154–175}
\begin{leancode}
\leanline{0}{theorem\ integral\_stdGaussian\_smoothing\ \{p\ :\ ℕ\}\ (hp\ :\ 0\ <\ p)}
\leanline{4}{(f\ :\ State\ p\ →\ ℝ)\ (M\ :\ ℝ≥0)\ (hf\ :\ Differentiable\ ℝ\ f)}
\leanline{4}{(hLip\ :\ LipschitzWith\ M\ (gradient\ f))\ (x\ :\ State\ p)\ (s\ :\ ℝ)\ (hs\ :\ 0\ ≤\ s)\ :}
\leanline{4}{(∫\ z,\ f\ (x\ +\ Real.sqrt\ (2\ *\ s)\ •\ z)\ -\ f\ x\ ∂stdGaussian\ (State\ p))\ ≤}
\leanline{6}{averageSmoothness\ f\ *\ (p\ :\ ℝ)\ *\ s}
\end{leancode}
\leanlabel{How the Lean proof follows the argument}
\texttt{integral\_\allowbreak{}std\allowbreak{}Gaussian\_\allowbreak{}smoothing} gives the formal counterpart of this step. The main intermediate facts are \texttt{\ensuremath{\mu}}, \texttt{hl}, \texttt{h\ensuremath{\xi}}, \texttt{hmean}, \texttt{hvar}, \texttt{h} in that order.
\end{leanbox}


Earlier results used here: \cite[Section~3.4.1]{LangevinLean2026ULASC}.
\section{\texorpdfstring{The exact averaged-smoothness one-step LMC theorem}{The exact averaged-smoothness one-step LMC theorem}}
\label{proof-module-AverageLMC-GaussianLMC-lean}


This is Theorem 1, claim 1, for the actual Gaussian Markov transition and the literal average coordinate Taylor supremum. It assumes no coupling recurrence, stationary-reference bundle, or desired convergence estimate.



For the actual independent Gaussian step, \(M_{\rm av}\) is the paper's average of the exact coordinate constants, and \(L\) is the global gradient Lipschitz constant. Theorem 1's local inequality is


\[W_2^2(\nu^+,\pi)\le\frac{1-mh}{1+mh}W_2^2(\nu,\pi)+\frac{(M_{\rm av}+m)h^2p}{1+mh}+\frac{h^2(hL-1)}{1+mh}\int\|\nabla f(x)\|^2\,d\nu(x),\qquad h\ge0.\]\subsection{\texorpdfstring{Include the zero step exactly}{Include the zero step exactly}}
\label{report-step-AverageLMC-GaussianLMC-lean-1}


At \(h=0\), the deterministic displacement and Gaussian coefficient both vanish. The transition is the projection onto the first coordinate of \(\mu\) times the standard Gaussian probability law; its pushforward is \(\mu\). This proves the zero-step case without a limiting argument.

\begin{leanbox}{Lean: \texttt{ulaLawStep\_\allowbreak{}zero}}
\leansource{\texttt{ulaLawStep\_\allowbreak{}zero} \quad AverageLMC/\allowbreak{}GaussianLMC.lean, lines 12–14}
\begin{leancode}
\leanline{0}{@[simp]\ theorem\ ulaLawStep\_zero\ \{p\ :\ ℕ\}\ (f\ :\ State\ p\ →\ ℝ)\ (μ\ :\ Measure\ (State\ p))\ :}
\leanline{4}{ulaLawStep\ f\ 0\ μ\ =\ μ}
\end{leancode}
\leanlabel{How the Lean proof follows the argument}
\texttt{ula\allowbreak{}Law\allowbreak{}Step\_\allowbreak{}zero} simplifies \texttt{ula\allowbreak{}Law\allowbreak{}Step} and \texttt{mala\allowbreak{}Proposal} at zero. \texttt{Measure.\allowbreak{}map\_\allowbreak{}fst\_\allowbreak{}prod} identifies the first-coordinate pushforward of the input-times-Gaussian product with the input measure.
\end{leanbox}


Earlier results used here: \cite[Section~3.34.1]{LangevinLean2026ULASC}; \cite[Section~3.36.7]{LangevinLean2026ULASC}.
\subsection{\texorpdfstring{Use exact coordinate smoothing and take the transport infimum}{Use exact coordinate smoothing and take the transport infimum}}
\label{report-step-AverageLMC-GaussianLMC-lean-2}


Instantiate the canonical coupling theorem with the proved standard-Gaussian coordinate smoothing bound and \(M_{\mathrm{av}}\) equal to the average of the literal coordinate Taylor suprema. Every finite initial coupling satisfies the same quadratic estimate. The squared-transport infimum theorem passes this estimate to the true initial Wasserstein distance; it also handles the possible negative coefficient \(1-mh\) when \(h\) is large. Divide by the strictly positive \(1+mh\) and collect the terms to obtain the paper’s exact formula for every \(h\ge 0\). The finite-distance hypothesis expresses the finite-second-moment interpretation of this displayed real-valued inequality.




\textbf{Hypotheses and algebra.} Here \(p\ge1\), \(m>0\), \(h\ge0\), \(f\) is differentiable and \(m\)-strongly convex, its gradient is \(L\)-Lipschitz with \(L\ge0\), and \(\mu\) is a probability law at finite quadratic transport from Gibbs. No upper bound on \(h\) is used. For a coupled initial pair \((X,Y)\), put \(A=X-h\nabla f(X)\) and \(G=\mathbb E\|\nabla f(X)\|^2\). The energy estimate and deterministic gradient-step estimates give

\[\begin{aligned}(1+mh)W_2^2(\nu^+,\pi)&\le\mathbb E\|X-Y\|^2-2h\mathbb E\langle\nabla f(X),X-Y\rangle+h^2G\\&\quad+2h\mathbb E[f(A)-f(Y)]+(M_{\rm av}+m)h^2p\\&\le(1-mh)\mathbb E\|X-Y\|^2+2h\mathbb E[f(A)-f(X)]+h^2G+(M_{\rm av}+m)h^2p\\&\le(1-mh)\mathbb E\|X-Y\|^2+h^2(hL-1)G+(M_{\rm av}+m)h^2p.\end{aligned}\]

The second line uses strong convexity at \(X\); the last uses the Lipschitz-gradient descent bound \(f(A)-f(X)\le-h(1-hL/2)\|\nabla f(X)\|^2\). The output law and \(G\) depend only on \(\mu\). Letting finite initial costs tend to their infimum proves the same polynomial inequality in the exact initial transport cost by continuity; this limit argument does not require \(1-mh\ge0\). Division by \(1+mh>0\) proves the claim. \(\square\)

\begin{leanbox}{Lean: \texttt{lmc\_\allowbreak{}one\_\allowbreak{}step}}
\leansource{\texttt{lmc\_\allowbreak{}one\_\allowbreak{}step} \quad AverageLMC/\allowbreak{}GaussianLMC.lean, lines 19–50}
\begin{leancode}
\leanline{0}{theorem\ lmc\_one\_step\ \{p\ :\ ℕ\}\ (hp\ :\ 0\ <\ p)\ (f\ :\ State\ p\ →\ ℝ)}
\leanline{4}{(m\ h\ :\ ℝ)\ (L\ :\ ℝ≥0)\ (hm\ :\ 0\ <\ m)\ (hh\ :\ 0\ ≤\ h)}
\leanline{4}{(hf\ :\ Differentiable\ ℝ\ f)\ (hsc\ :\ StrongConvexOn\ univ\ m\ f)}
\leanline{4}{(hLip\ :\ LipschitzWith\ L\ (gradient\ f))}
\leanline{4}{(μ\ :\ Measure\ (State\ p))\ [IsProbabilityMeasure\ μ]}
\leanline{4}{(hfinite\ :\ wasserstein2\ μ\ (gibbsMeasure\ f)\ ≠\ ⊤)\ :}
\leanline{4}{wasserstein2Real\ (ulaLawStep\ f\ h\ μ)\ (gibbsMeasure\ f)\ \textasciicircum{}\ 2\ ≤}
\leanline{6}{(1\ -\ m\ *\ h)\ /\ (1\ +\ m\ *\ h)\ *\ wasserstein2Real\ μ\ (gibbsMeasure\ f)\ \textasciicircum{}\ 2\ +}
\leanline{6}{((averageSmoothness\ f\ +\ m)\ *\ h\ \textasciicircum{}\ 2\ *\ p)\ /\ (1\ +\ m\ *\ h)\ +}
\leanline{6}{h\ \textasciicircum{}\ 2\ *\ (h\ *\ (L\ :\ ℝ)\ -\ 1)\ /\ (1\ +\ m\ *\ h)\ *}
\leanline{8}{(∫\ x,\ ‖gradient\ f\ x‖\ \textasciicircum{}\ 2\ ∂μ)}
\end{leancode}
\leanlabel{How the Lean proof follows the argument}
\texttt{lp\_\allowbreak{}gradient\_\allowbreak{}step} and \texttt{lp\_\allowbreak{}gradient\_\allowbreak{}descent} are the two expected deterministic inequalities whose potential terms cancel. \texttt{lmc\_\allowbreak{}one\_\allowbreak{}step} first separates the zero step. In the positive branch, \texttt{hsm} is the exact coordinate smoothing estimate, \texttt{houtput} proves the output transport finite, and \texttt{wasserstein2\_\allowbreak{}sq\_\allowbreak{}bound\_\allowbreak{}of\_\allowbreak{}finite\_\allowbreak{}couplings} passes the common bound to the infimum even when its input-cost coefficient has either sign. The final field normalization divides by the strictly positive denominator.
\end{leanbox}


Earlier results used here: \cite[Section~3.17.3]{LangevinLean2026ULASC}; \cite[Section~3.17.4]{LangevinLean2026ULASC}; \cite[Section~3.18.2]{LangevinLean2026ULASC}; \cite[Section~3.26.3]{LangevinLean2026ULASC}; \cite[Section~3.36.7]{LangevinLean2026ULASC}.
\section{\texorpdfstring{Scalar recurrences for the normalized error}{Scalar recurrences for the normalized error}}
\label{proof-module-AverageLMC-Recurrence-lean}


This section solves the scalar recurrences that turn a one-step Wasserstein bound into a bound after many steps. No probability enters: the statements concern real sequences that are assumed to satisfy a one-step inequality. In the applications the sequence is the normalized error \(r_k=\frac mp\mathsf W_2^2(\nu_k,\pi)\) of the Langevin Monte Carlo chain, \(t=mh\) is the dimensionless step, and \(a=\kappa_{\rm av}+1\) with \(\kappa_{\rm av}=M_{\rm av}/m\). Multiplying the one-step bound of Theorem 1 by \(m/p\) and dropping its nonpositive last term gives


\[r_{k+1}\le\frac{1-t}{1+t}\,r_k+\frac{a\,t^2}{1+t},\]


because \(\frac mp\cdot\frac{(M_{\rm av}+m)h^2p}{1+mh}=\frac{(\kappa_{\rm av}+1)(mh)^2}{1+mh}\). We prove the elementary inequality \(\frac{1-t}{1+t}\le e^{-2t}\) on \([0,1]\), the iteration of an affine inequality whose fixed point is known, the resulting constant-step bound \(r_n\le e^{-2tn}r_0+\frac{at}2\) (the scalar content of the second claim of Theorem 1), and the induction step of the decreasing-step schedule of Proposition 1, in which \(t_k=1/(\kappa+k)\) and \(r_k\le a/(\kappa+k)\).



Earlier results used here: \cite[Section~3.15]{LangevinLean2026ULASC}.
\subsection{\texorpdfstring{The contraction factor is at most an exponential}{The contraction factor is at most an exponential}}
\label{report-step-AverageLMC-Recurrence-lean-1}


\textbf{Claim.} For every real \(t\) with \(0\le t\le1\),


\[\frac{1-t}{1+t}\le e^{-2t}.\]


\textbf{Proof.} If \(t=1\), the left-hand side is \(0\) and the right-hand side \(e^{-2}\) is positive. Let \(0\le t<1\), so that \(1+t>0\) and \(1-t>0\). Consider \(\varphi(s)=\log(1+s)-\log(1-s)-2s\) on \([0,1)\). Then \(\varphi(0)=0\) and


\[\varphi'(s)=\frac1{1+s}+\frac1{1-s}-2=\frac{2}{1-s^2}-2=\frac{2s^2}{1-s^2}\ge0,\]


so \(\varphi\) is nondecreasing by the mean value theorem and \(\varphi(t)\ge0\), that is,


\[2t\le\log\frac{1+t}{1-t}.\qquad(1)\]


(Equivalently, \(t\) is the first term of the series \(\frac12\log\frac{1+t}{1-t}=\sum_{i\ge0}\frac{t^{2i+1}}{2i+1}\), whose terms are nonnegative.) Since \(\frac{1-t}{1+t}=\bigl(\frac{1+t}{1-t}\bigr)^{-1}\) and \(\log(x^{-1})=-\log x\) for \(x>0\), (1) gives \(\log\frac{1-t}{1+t}\le-2t\). The exponential is increasing and \(\frac{1-t}{1+t}>0\), so exponentiating gives the claim. \(\square\)




In the applications \(t=mh\), and \(t\le1\) follows from \(m\le M_\infty\) and \(hM_\infty\le1\).

\begin{leanbox}{Lean: \texttt{contraction\_\allowbreak{}le\_\allowbreak{}exp}}
\leansource{\texttt{contraction\_\allowbreak{}le\_\allowbreak{}exp} \quad AverageLMC/\allowbreak{}Recurrence.lean, lines 21–34}
\begin{leancode}
\leanline{0}{theorem\ contraction\_le\_exp\ \{t\ :\ ℝ\}\ (ht\ :\ 0\ ≤\ t)\ (ht1\ :\ t\ ≤\ 1)\ :}
\leanline{4}{(1\ -\ t)\ /\ (1\ +\ t)\ ≤\ Real.exp\ (-2\ *\ t)}
\end{leancode}
\leanlabel{How the Lean proof follows the argument}
\texttt{contraction\_\allowbreak{}le\_\allowbreak{}exp} takes \texttt{ht : 0 \ensuremath{\le} t} and \texttt{ht1 : t \ensuremath{\le} 1}. \texttt{rcases eq\_\allowbreak{}or\_\allowbreak{}lt\_\allowbreak{}of\_\allowbreak{}le ht1} splits off the case \(t=1\), closed by \texttt{simpa} with the positivity of \(e^{-2}\) (\texttt{Real.\allowbreak{}exp\_\allowbreak{}pos}). For \(t<1\), \texttt{hl} is the Mathlib inequality \texttt{Real.\allowbreak{}sum\_\allowbreak{}range\_\allowbreak{}le\_\allowbreak{}log\_\allowbreak{}div ht h 1}, which states \(\sum_{i<n}t^{2i+1}/(2i+1)\le\frac12\log\frac{1+t}{1-t}\); with \(n=1\) and after \texttt{simp only} it is (1). \texttt{hi} writes the contraction as the inverse of \((1+t)/(1-t)\). The goal is moved to logarithms with \texttt{Real.\allowbreak{}log\_\allowbreak{}le\_\allowbreak{}iff\_\allowbreak{}le\_\allowbreak{}exp} (valid because the quotient is positive, \texttt{div\_\allowbreak{}pos hm hp}), \texttt{Real.\allowbreak{}log\_\allowbreak{}inv} turns the logarithm of the inverse into \(-\log\frac{1+t}{1-t}\), and \texttt{linarith} concludes from \texttt{hl}.
\end{leanbox}
\subsection{\texorpdfstring{Iterating an affine inequality over a finite horizon}{Iterating an affine inequality over a finite horizon}}
\label{report-step-AverageLMC-Recurrence-lean-2}


\textbf{Claim.} Let \((r_k)_{k\ge0}\) be real numbers, let \(q\ge0\), \(c\ge0\) and \(n\in\mathbb N\), and suppose that


\[r_{j+1}\le q\,r_j+(1-q)\,c\qquad\text{for every }j<n.\]


Then \(r_n\le q^nr_0+c\).




\textbf{Proof.} We first prove by induction on \(k\) that


\[r_k\le q^kr_0+(1-q^k)\,c\qquad\text{for every }k\le n.\qquad(2)\]


\emph{Base case \(k=0\).} Since \(q^0=1\), the right-hand side is \(r_0\).




\emph{Inductive step.} Let \(k+1\le n\) and assume (2) for \(k\) (which is allowed since \(k\le n\)). Because \(q\ge0\), multiplying the induction hypothesis by \(q\) preserves it: \(q\,r_k\le q^{k+1}r_0+(q-q^{k+1})c\). Since \(k<n\), the hypothesis applies with \(j=k\):


\[r_{k+1}\le q\,r_k+(1-q)c\le q^{k+1}r_0+(q-q^{k+1}+1-q)\,c=q^{k+1}r_0+(1-q^{k+1})\,c.\]


This proves (2). For \(k=n\), (2) gives \(r_n\le q^nr_0+c-q^nc\le q^nr_0+c\), because \(q^nc\ge0\) as a product of the nonnegative numbers \(q^n\) and \(c\). \(\square\)




The constant \(c\) is the fixed point of the map \(x\mapsto qx+(1-q)c\). The hypothesis is required only for the first \(n\) indices, which is what allows the bound to be applied to the warmup phase of a schedule that changes after \(n\) steps. No condition \(q\le1\) is needed.

\begin{leanbox}{Lean: \texttt{affine\_\allowbreak{}floor\_\allowbreak{}up\_\allowbreak{}to}}
\leansource{\texttt{affine\_\allowbreak{}floor\_\allowbreak{}up\_\allowbreak{}to} \quad AverageLMC/\allowbreak{}Recurrence.lean, lines 37–52}
\begin{leancode}
\leanline{0}{theorem\ affine\_floor\_up\_to\ (r\ :\ ℕ\ →\ ℝ)\ (q\ c\ :\ ℝ)\ (hq\ :\ 0\ ≤\ q)}
\leanline{4}{(hc\ :\ 0\ ≤\ c)\ (n\ :\ ℕ)}
\leanline{4}{(hstep\ :\ ∀\ j\ <\ n,\ r\ (j\ +\ 1)\ ≤\ q\ *\ r\ j\ +\ (1\ -\ q)\ *\ c)\ :}
\leanline{4}{r\ n\ ≤\ q\ \textasciicircum{}\ n\ *\ r\ 0\ +\ c}
\end{leancode}
\leanlabel{How the Lean proof follows the argument}
In \texttt{affine\_\allowbreak{}floor\_\allowbreak{}up\_\allowbreak{}to} the sequence is \texttt{r : \ensuremath{\mathbb N} \ensuremath{\to} \ensuremath{\mathbb R}} and the finite-horizon hypothesis is \texttt{hstep : \ensuremath{\forall} j < n, r (j + 1) \ensuremath{\le} q * r j + (1 -\allowbreak{} q) * c}. The strengthened statement (2) is \texttt{hstrong}, proved by \texttt{induction k with}: the \texttt{zero} case is \texttt{simp}; in the \texttt{succ} case \texttt{hprev} is the induction hypothesis (applicable because \texttt{Nat.\allowbreak{}le\_\allowbreak{}of\_\allowbreak{}succ\_\allowbreak{}le hk} gives \(k\le n\)), \texttt{hmul} multiplies it by \(q\ge0\) with \texttt{mul\_\allowbreak{}le\_\allowbreak{}mul\_\allowbreak{}of\_\allowbreak{}nonneg\_\allowbreak{}left}, \texttt{hs} is \texttt{hstep k} (with \texttt{Nat.\allowbreak{}lt\_\allowbreak{}of\_\allowbreak{}succ\_\allowbreak{}le hk} for \(k<n\)), \texttt{rw [pow\_\allowbreak{}succ]} writes \(q^{k+1}=q^kq\), and \texttt{nlinarith} performs the displayed computation. The final \texttt{nlinarith} uses \(q^nc\ge0\), supplied as \texttt{mul\_\allowbreak{}nonneg (pow\_\allowbreak{}nonneg hq n) hc}.
\end{leanbox}
\subsection{\texorpdfstring{The constant-step bound}{The constant-step bound}}
\label{report-step-AverageLMC-Recurrence-lean-3}


\textbf{Claim.} Let \((r_k)_{k\ge0}\) be real numbers and \(a,t\in\mathbb R\) with \(a\ge0\), \(0\le t\le1\) and \(r_0\ge0\). Let \(n\in\mathbb N\) and suppose that


\[r_{j+1}\le\frac{1-t}{1+t}\,r_j+\frac{a\,t^2}{1+t}\qquad\text{for every }j<n.\qquad(3)\]


Then


\[r_n\le e^{-2tn}\,r_0+\frac{a\,t}2.\qquad(4)\]


In particular, if (3) holds for every \(j\in\mathbb N\), then (4) holds for every \(n\in\mathbb N\).




\textbf{Proof.} Put \(q=\frac{1-t}{1+t}\) and \(c=\frac{at}2\). Since \(0\le t\le1\), we have \(1+t>0\) and \(1-t\ge0\), hence \(q\ge0\); also \(c\ge0\) because \(a,t\ge0\). The forcing term in (3) equals \((1-q)c\):


\[(1-q)\,c=\frac{(1+t)-(1-t)}{1+t}\cdot\frac{at}2=\frac{2t}{1+t}\cdot\frac{at}2=\frac{a\,t^2}{1+t}.\]


Thus (3) is the hypothesis of Step 2 with these \(q\) and \(c\), and Step 2 gives \(r_n\le q^nr_0+\frac{at}2\). By Step 1, \(0\le q\le e^{-2t}\); raising to the power \(n\) preserves inequalities between nonnegative numbers, so \(q^n\le(e^{-2t})^n=e^{-2tn}\). Multiplying by \(r_0\ge0\) gives \(q^nr_0\le e^{-2tn}r_0\), which proves (4). The second assertion applies the first to each \(n\), the hypothesis for all \(j\) containing the hypothesis for \(j<n\). \(\square\)




With \(r_k=\frac mp\mathsf W_2^2(\nu_k,\pi)\), \(t=mh\) and \(a=\kappa_{\rm av}+1\), multiplying (4) by \(p/m\) gives \(\mathsf W_2^2(\nu_n,\pi)\le e^{-2mnh}\mathsf W_2^2(\nu_0,\pi)+\frac{(M_{\rm av}+m)hp}{2m}\), since \(\frac pm\cdot\frac{(\kappa_{\rm av}+1)mh}2=\frac{(M_{\rm av}+m)hp}{2m}\). This is the second claim of Theorem 1.

\begin{leanbox}{Lean: \texttt{constant\_\allowbreak{}step\_\allowbreak{}normalized\_\allowbreak{}up\_\allowbreak{}to}, \texttt{constant\_\allowbreak{}step\_\allowbreak{}normalized}}
\leansource{\texttt{constant\_\allowbreak{}step\_\allowbreak{}normalized\_\allowbreak{}up\_\allowbreak{}to} \quad AverageLMC/\allowbreak{}Recurrence.lean, lines 56–75}
\begin{leancode}
\leanline{0}{theorem\ constant\_step\_normalized\_up\_to\ (r\ :\ ℕ\ →\ ℝ)\ (a\ t\ :\ ℝ)}
\leanline{4}{(ha\ :\ 0\ ≤\ a)\ (ht\ :\ 0\ ≤\ t)\ (ht1\ :\ t\ ≤\ 1)\ (hr0\ :\ 0\ ≤\ r\ 0)\ (n\ :\ ℕ)}
\leanline{4}{(hstep\ :\ ∀\ j\ <\ n,\ r\ (j\ +\ 1)\ ≤}
\leanline{6}{(1\ -\ t)\ /\ (1\ +\ t)\ *\ r\ j\ +\ a\ *\ t\ \textasciicircum{}\ 2\ /\ (1\ +\ t))\ :}
\leanline{4}{r\ n\ ≤\ Real.exp\ (-2\ *\ t\ *\ (n\ :\ ℝ))\ *\ r\ 0\ +\ a\ *\ t\ /\ 2}
\end{leancode}
\leansource{\texttt{constant\_\allowbreak{}step\_\allowbreak{}normalized} \quad AverageLMC/\allowbreak{}Recurrence.lean, lines 78–83}
\begin{leancode}
\leanline{0}{theorem\ constant\_step\_normalized\ (r\ :\ ℕ\ →\ ℝ)\ (a\ t\ :\ ℝ)}
\leanline{4}{(ha\ :\ 0\ ≤\ a)\ (ht\ :\ 0\ ≤\ t)\ (ht1\ :\ t\ ≤\ 1)\ (hr0\ :\ 0\ ≤\ r\ 0)}
\leanline{4}{(hstep\ :\ ∀\ j,\ r\ (j\ +\ 1)\ ≤}
\leanline{6}{(1\ -\ t)\ /\ (1\ +\ t)\ *\ r\ j\ +\ a\ *\ t\ \textasciicircum{}\ 2\ /\ (1\ +\ t))\ (n\ :\ ℕ)\ :}
\leanline{4}{r\ n\ ≤\ Real.exp\ (-2\ *\ t\ *\ (n\ :\ ℝ))\ *\ r\ 0\ +\ a\ *\ t\ /\ 2}
\end{leancode}
\leanlabel{How the Lean proof follows the argument}
\begin{enumerate}


\item \texttt{constant\_\allowbreak{}step\_\allowbreak{}normalized\_\allowbreak{}up\_\allowbreak{}to} has the hypotheses \texttt{ha : 0 \ensuremath{\le} a}, \texttt{ht : 0 \ensuremath{\le} t}, \texttt{ht1 : t \ensuremath{\le} 1}, \texttt{hr0 : 0 \ensuremath{\le} r 0} and the finite-horizon recurrence \texttt{hstep}, which is (3). \texttt{hq} is \(q\ge0\) (\texttt{div\_\allowbreak{}nonneg}), and \texttt{hfix} is the identity \(\frac{at^2}{1+t}=(1-q)\frac{at}2\), proved by \texttt{field\_\allowbreak{}simp} and \texttt{ring}.

\item \texttt{h} applies \texttt{affine\_\allowbreak{}floor\_\allowbreak{}up\_\allowbreak{}to} (Step 2) with \(c=at/2\), after \texttt{simpa only [hfix]} rewrites the forcing term of \texttt{hstep}.

\item \texttt{hp} raises the inequality of \texttt{contraction\_\allowbreak{}le\_\allowbreak{}exp} to the \(n\)th power with \texttt{pow\_\allowbreak{}le\_\allowbreak{}pow\_\allowbreak{}left\ensuremath{{}_{0}}}, and \texttt{heq} identifies \((e^{-2t})^n\) with \(e^{-2tn}\) through \texttt{Real.\allowbreak{}exp\_\allowbreak{}nat\_\allowbreak{}mul}. The last line multiplies by \(r_0\ge0\) (\texttt{mul\_\allowbreak{}le\_\allowbreak{}mul\_\allowbreak{}of\_\allowbreak{}nonneg\_\allowbreak{}right hp hr0}) and chains with \texttt{h}.

\item \texttt{constant\_\allowbreak{}step\_\allowbreak{}normalized} is the version for a recurrence valid at every index: it calls the previous theorem with \texttt{fun j \_\allowbreak{} \ensuremath{\mapsto} hstep j}.

\end{enumerate}
\end{leanbox}
\subsection{\texorpdfstring{The induction step of the decreasing-step schedule}{The induction step of the decreasing-step schedule}}
\label{report-step-AverageLMC-Recurrence-lean-4}


\textbf{Claim.} (1) Let \(a,x,r\in\mathbb R\) with \(x\ge1\) and \(r\le a/x\). Then


\[\frac{1-1/x}{1+1/x}\,r+\frac{a\,(1/x)^2}{1+1/x}\le\frac a{x+1}.\]


(2) Let \((r_n)_{n\ge0}\) be real numbers and \(a,\kappa\in\mathbb R\) with \(a\ge0\), \(\kappa\ge1\) and \(r_0\le a/\kappa\). Suppose that for every \(n\in\mathbb N\), with \(t_n=1/(\kappa+n)\),


\[r_{n+1}\le\frac{1-t_n}{1+t_n}\,r_n+\frac{a\,t_n^2}{1+t_n}.\qquad(5)\]


Then \(r_n\le\frac a{\kappa+n}\) for every \(n\in\mathbb N\).




\textbf{Proof.} (1) Since \(x\ge1>0\), we have \(1/x\le1\) and \(1+1/x>0\), so the coefficient \(q=\frac{1-1/x}{1+1/x}\) is nonnegative. Multiplying \(r\le a/x\) by \(q\ge0\),


\[q\,r+\frac{a(1/x)^2}{1+1/x}\le q\,\frac ax+\frac{a(1/x)^2}{1+1/x}.\]


Multiplying numerators and denominators by \(x\) gives \(q=\frac{x-1}{x+1}\) and \(\frac{(1/x)^2}{1+1/x}=\frac1{x(x+1)}\), so the right-hand side equals


\[\frac{a(x-1)}{x(x+1)}+\frac a{x(x+1)}=\frac{a\,x}{x(x+1)}=\frac a{x+1}.\]


(2) Induction on \(n\). For \(n=0\) the claim is the hypothesis \(r_0\le a/\kappa\). Assume \(r_n\le a/(\kappa+n)\) and put \(x=\kappa+n\). Then \(x\ge1\) because \(\kappa\ge1\) and \(n\ge0\), and \(t_n=1/x\). By (5) and part (1),


\[r_{n+1}\le\frac{1-1/x}{1+1/x}\,r_n+\frac{a(1/x)^2}{1+1/x}\le\frac a{x+1}=\frac a{\kappa+(n+1)}.\qquad\square\]


Part (1) uses no sign condition on \(a\), and the hypothesis \(a\ge0\) of part (2) is not used in its proof. In Proposition 1, \(x=\kappa_\infty+k-k_0\) and \(t_n=mh_k\) is the decreasing step after the warmup phase; part (1) states that the bound \(a/x\) is propagated to \(a/(x+1)\) by one step of the schedule.

\begin{leanbox}{Lean: \texttt{harmonic\_\allowbreak{}step}, \texttt{variable\_\allowbreak{}step\_\allowbreak{}tail}}
\leansource{\texttt{harmonic\_\allowbreak{}step} \quad AverageLMC/\allowbreak{}Recurrence.lean, lines 86–100}
\begin{leancode}
\leanline{0}{theorem\ harmonic\_step\ \{a\ x\ r\ :\ ℝ\}\ (hx\ :\ 1\ ≤\ x)}
\leanline{4}{(hr\ :\ r\ ≤\ a\ /\ x)\ :}
\leanline{4}{(1\ -\ 1\ /\ x)\ /\ (1\ +\ 1\ /\ x)\ *\ r\ +\ a\ *\ (1\ /\ x)\ \textasciicircum{}\ 2\ /\ (1\ +\ 1\ /\ x)}
\leanline{6}{≤\ a\ /\ (x\ +\ 1)}
\end{leancode}
\leansource{\texttt{variable\_\allowbreak{}step\_\allowbreak{}tail} \quad AverageLMC/\allowbreak{}Recurrence.lean, lines 104–117}
\begin{leancode}
\leanline{0}{theorem\ variable\_step\_tail\ (r\ :\ ℕ\ →\ ℝ)\ (a\ κ\ :\ ℝ)}
\leanline{4}{(\_ha\ :\ 0\ ≤\ a)\ (hκ\ :\ 1\ ≤\ κ)\ (hinit\ :\ r\ 0\ ≤\ a\ /\ κ)}
\leanline{4}{(hstep\ :\ ∀\ n,\ r\ (n\ +\ 1)\ ≤}
\leanline{6}{(1\ -\ 1\ /\ (κ\ +\ (n\ :\ ℝ)))\ /\ (1\ +\ 1\ /\ (κ\ +\ (n\ :\ ℝ)))\ *\ r\ n\ +}
\leanline{6}{a\ *\ (1\ /\ (κ\ +\ (n\ :\ ℝ)))\ \textasciicircum{}\ 2\ /\ (1\ +\ 1\ /\ (κ\ +\ (n\ :\ ℝ))))}
\leanline{4}{(n\ :\ ℕ)\ :\ r\ n\ ≤\ a\ /\ (κ\ +\ (n\ :\ ℝ))}
\end{leancode}
\leanlabel{How the Lean proof follows the argument}
\begin{enumerate}


\item \texttt{harmonic\_\allowbreak{}step} has the hypotheses \texttt{hx : 1 \ensuremath{\le} x} and \texttt{hr : r \ensuremath{\le} a /\allowbreak{} x}. The auxiliary facts are \(x>0\) (\texttt{hx0}), \(1+1/x>0\) (\texttt{hd}), \(1/x\le1\) (\texttt{hdiv}, from \texttt{div\_\allowbreak{}le\_\allowbreak{}one}) and \(q\ge0\) (\texttt{hq}). The \texttt{calc} block has two lines: the first multiplies \texttt{hr} by \(q\) with \texttt{mul\_\allowbreak{}le\_\allowbreak{}mul\_\allowbreak{}of\_\allowbreak{}nonneg\_\allowbreak{}left} and adds the unchanged forcing term (\texttt{add\_\allowbreak{}le\_\allowbreak{}add … le\_\allowbreak{}rfl}); the second is the algebraic identity, proved by \texttt{field\_\allowbreak{}simp} and \texttt{ring}.

\item \texttt{variable\_\allowbreak{}step\_\allowbreak{}tail} has the hypotheses \texttt{\_\allowbreak{}ha : 0 \ensuremath{\le} a} (unused, as the leading underscore records), \texttt{h\ensuremath{\kappa} : 1 \ensuremath{\le} \ensuremath{\kappa}}, \texttt{hinit : r 0 \ensuremath{\le} a /\allowbreak{} \ensuremath{\kappa}}, and \texttt{hstep}, which is (5) written with \(t_n\) spelled out as \texttt{1 /\allowbreak{} (\ensuremath{\kappa} + (n : \ensuremath{\mathbb R}))}. The proof is \texttt{induction n with}: the \texttt{zero} case is \texttt{simpa using hinit}; in the \texttt{succ} case \texttt{hx} proves \(\kappa+n\ge1\) from \texttt{Nat.\allowbreak{}cast\_\allowbreak{}nonneg}, \texttt{(hstep n).\allowbreak{}trans (harmonic\_\allowbreak{}step hx ih)} chains (5) with part (1), and \texttt{simpa only [Nat.\allowbreak{}cast\_\allowbreak{}add, Nat.\allowbreak{}cast\_\allowbreak{}one, add\_\allowbreak{}assoc]} rewrites \(\kappa+n+1\) as \(\kappa+(n+1)\).

\end{enumerate}
\end{leanbox}
\section{\texorpdfstring{The decreasing-step schedule and its iteration count}{The decreasing-step schedule and its iteration count}}
\label{proof-module-AverageLMC-Complexity-lean}


This section carries out the scalar part of Proposition 1, the decreasing-step version of Langevin Monte Carlo. Throughout, \(\kappa=\kappa_\infty=M_\infty/m\), \(a=\kappa_{\rm av}+1\), and \(r_k=\frac mp\mathsf W_2^2(\nu_k,\pi)\) is the normalized error, so that the assumption \(\mathsf W_2^2(\nu_0,\pi)\le\varepsilon_0^2p/m\) of the proposition reads \(r_0\le\varepsilon_0^2\). With \(\log_+x=\max(\log x,0)\), the warmup length and the step sizes of the proposition are


\[\begin{aligned}
k_0&=\Bigl\lceil\frac\kappa2\log_+\Bigl(\frac{2\kappa\varepsilon_0^2}{a}\Bigr)\Bigr\rceil,\\
h_k&=\frac1{M_\infty+m(k-k_0)_+},\\
t_k&=mh_k=\frac1{\kappa+(k-k_0)_+}.
\end{aligned}\]


Assuming that the normalized one-step recurrence \(r_{k+1}\le\frac{1-t_k}{1+t_k}r_k+\frac{a\,t_k^2}{1+t_k}\) holds at every step, we prove \(r_k\le a/(\kappa+k-k_0)\) for all \(k\ge k_0\), which after multiplication by \(p/m\) is the error bound of the proposition, and we prove that \(r_k\le\varepsilon^2\) as soon as \(k\ge k_\varepsilon=k_0+\lceil(a/\varepsilon^2-\kappa)_+\rceil\). The recurrence itself is a consequence of the one-step bound of Theorem 1, which is proved elsewhere; the present section uses it as a hypothesis.

\subsection{\texorpdfstring{The positive logarithm, the warmup length and the schedule}{The positive logarithm, the warmup length and the schedule}}
\label{report-step-AverageLMC-Complexity-lean-1}


\textbf{Definitions.} For \(x>0\) let \(\log_+x=\max(\log x,0)\). For real \(y\), \(\lceil y\rceil\) denotes the smallest natural number \(n\) with \(y\le n\) (so \(\lceil y\rceil=0\) when \(y\le0\)), and for natural numbers \((k-k_0)_+=\max(k-k_0,0)\). Given \(\kappa,a,\varepsilon_0\in\mathbb R\), define


\[k_0=k_0(\kappa,a,\varepsilon_0)=\Bigl\lceil\frac\kappa2\log_+\Bigl(\frac{2\kappa\varepsilon_0^2}{a}\Bigr)\Bigr\rceil,\qquad t_k=\frac1{\kappa+(k-k_0)_+}\quad(k\in\mathbb N).\]


With \(\kappa=\kappa_\infty\) and \(a=\kappa_{\rm av}+1\), \(k_0\) is the warmup length of Proposition 1, and \(t_k=mh_k\) because \(\frac m{M_\infty+m(k-k_0)_+}=\frac1{\kappa_\infty+(k-k_0)_+}\). For \(k\le k_0\) the step is constant, \(t_k=1/\kappa\); afterwards \(t_k=1/(\kappa+k-k_0)\) decreases harmonically.




\textbf{Claim.} If \(\kappa>0\), then \(0<t_k\le1/\kappa\) for every \(k_0\) and every \(k\in\mathbb N\).




\textbf{Proof.} The denominator \(\kappa+(k-k_0)_+\) is at least \(\kappa>0\), because \((k-k_0)_+\ge0\). Hence \(t_k>0\), and since \(x\mapsto1/x\) is decreasing on \((0,\infty)\), \(t_k\le1/\kappa\). \(\square\)




In terms of the step size, \(t_k\le1/\kappa_\infty\) is \(h_kM_\infty\le1\): every step of the schedule satisfies the step condition of Theorem 1 under which the last term of the one-step bound is nonpositive.

\begin{leanbox}{Lean: \texttt{positiveLog}, \texttt{warmup}, \texttt{variableStep} and 1 more}
\leansource{\texttt{positiveLog} \quad AverageLMC/\allowbreak{}Complexity.lean, lines 17–17}
\begin{leancode}
\leanline{0}{noncomputable\ def\ positiveLog\ (x\ :\ ℝ)\ :\ ℝ\ :=\ max\ (Real.log\ x)\ 0}
\end{leancode}
\leansource{\texttt{warmup} \quad AverageLMC/\allowbreak{}Complexity.lean, lines 20–21}
\begin{leancode}
\leanline{0}{noncomputable\ def\ warmup\ (κ\ a\ ε₀\ :\ ℝ)\ :\ ℕ\ :=}
\leanline{2}{⌈κ\ /\ 2\ *\ positiveLog\ (2\ *\ κ\ *\ ε₀\ \textasciicircum{}\ 2\ /\ a)⌉₊}
\end{leancode}
\leansource{\texttt{variableStep} \quad AverageLMC/\allowbreak{}Complexity.lean, lines 24–25}
\begin{leancode}
\leanline{0}{noncomputable\ def\ variableStep\ (κ\ :\ ℝ)\ (k₀\ k\ :\ ℕ)\ :\ ℝ\ :=}
\leanline{2}{1\ /\ (κ\ +\ ((k\ -\ k₀\ :\ ℕ)\ :\ ℝ))}
\end{leancode}
\leansource{\texttt{variableStep\_\allowbreak{}le} \quad AverageLMC/\allowbreak{}Complexity.lean, lines 28–35}
\begin{leancode}
\leanline{0}{theorem\ variableStep\_le\ \{κ\ :\ ℝ\}\ (hκ\ :\ 0\ <\ κ)\ (k₀\ k\ :\ ℕ)\ :}
\leanline{4}{0\ <\ variableStep\ κ\ k₀\ k\ ∧\ variableStep\ κ\ k₀\ k\ ≤\ 1\ /\ κ}
\end{leancode}
\leanlabel{How the Lean proof follows the argument}
\begin{enumerate}


\item \texttt{positive\allowbreak{}Log x} is \texttt{max (Real.\allowbreak{}log x) 0}. Mathlib's \texttt{Real.\allowbreak{}log} is defined on all of \(\mathbb R\) (with \(\log0=0\) and \(\log x=\log|x|\)); only positive arguments occur below.

\item \texttt{warmup \ensuremath{\kappa} a \ensuremath{\varepsilon}\ensuremath{{}_{0}}} is \texttt{⌈\ensuremath{\kappa} /\allowbreak{} 2 * positive\allowbreak{}Log (2 * \ensuremath{\kappa} * \ensuremath{\varepsilon}\ensuremath{{}_{0}} \textasciicircum{} 2 /\allowbreak{} a)⌉\ensuremath{{}_{+}}}, where \texttt{⌈·⌉\ensuremath{{}_{+}}} is the natural-number ceiling \texttt{Nat.\allowbreak{}ceil}, equal to \(0\) on nonpositive arguments. This is \(k_0\).

\item \texttt{variable\allowbreak{}Step \ensuremath{\kappa} k\ensuremath{{}_{0}} k} is \texttt{1 /\allowbreak{} (\ensuremath{\kappa} + ((k -\allowbreak{} k\ensuremath{{}_{0}} : \ensuremath{\mathbb N}) : \ensuremath{\mathbb R}))}. The subtraction \texttt{k -\allowbreak{} k\ensuremath{{}_{0}}} is performed in \(\mathbb N\), where it is truncated at \(0\); this is \((k-k_0)_+\), which is then cast to \(\mathbb R\).

\item \texttt{variable\allowbreak{}Step\_\allowbreak{}le} proves the two inequalities with \texttt{constructor}: positivity of the denominator follows from \texttt{Nat.\allowbreak{}cast\_\allowbreak{}nonneg} and is closed by \texttt{positivity}; the upper bound is \texttt{one\_\allowbreak{}div\_\allowbreak{}le\_\allowbreak{}one\_\allowbreak{}div\_\allowbreak{}of\_\allowbreak{}le h\ensuremath{\kappa}}, the antitonicity of \(x\mapsto1/x\).

\end{enumerate}
\end{leanbox}
\subsection{\texorpdfstring{The warmup length reduces the initial error}{The warmup length reduces the initial error}}
\label{report-step-AverageLMC-Complexity-lean-2}


\textbf{Claim.} Let \(\kappa>0\), \(a>0\) and \(\varepsilon_0>0\), and let \(k_0=k_0(\kappa,a,\varepsilon_0)\). Then


\[e^{-2k_0/\kappa}\,\varepsilon_0^2\le\frac a{2\kappa}.\]


\textbf{Proof.} Let \(x=2\kappa\varepsilon_0^2/a\), which is positive. By the definition of the ceiling, \(k_0\ge\frac\kappa2\log_+x\), and \(\log_+x\ge\log x\) by the definition of the maximum. Hence \(k_0\ge\frac\kappa2\log x\), and multiplying by \(-2/\kappa<0\) reverses the inequality:


\[-\frac{2k_0}\kappa\le-\log x.\]


The exponential is increasing and \(\varepsilon_0^2\ge0\), so


\[e^{-2k_0/\kappa}\varepsilon_0^2\le e^{-\log x}\varepsilon_0^2=\frac{\varepsilon_0^2}{x}=\frac{a\,\varepsilon_0^2}{2\kappa\varepsilon_0^2}=\frac a{2\kappa}.\qquad\square\]


When \(x\le1\) the positive logarithm vanishes and \(k_0=0\); the argument covers this case, since \(\log x\le0=\log_+x\). The inequality says that after \(k_0\) constant steps the contribution of the initial error, \(e^{-2k_0/\kappa}r_0\le e^{-2k_0/\kappa}\varepsilon_0^2\), is at most the persistent error \(a/(2\kappa)\) of the constant-step bound with \(t=1/\kappa\).

\begin{leanbox}{Lean: \texttt{warmup\_\allowbreak{}decay}}
\leansource{\texttt{warmup\_\allowbreak{}decay} \quad AverageLMC/\allowbreak{}Complexity.lean, lines 39–56}
\begin{leancode}
\leanline{0}{theorem\ warmup\_decay\ \{κ\ a\ ε₀\ :\ ℝ\}\ (hκ\ :\ 0\ <\ κ)\ (ha\ :\ 0\ <\ a)}
\leanline{4}{(hε₀\ :\ 0\ <\ ε₀)\ :}
\leanline{4}{Real.exp\ (-2\ *\ (warmup\ κ\ a\ ε₀\ :\ ℝ)\ /\ κ)\ *\ ε₀\ \textasciicircum{}\ 2\ ≤\ a\ /\ (2\ *\ κ)}
\end{leancode}
\leanlabel{How the Lean proof follows the argument}
\texttt{warmup\_\allowbreak{}decay} has the hypotheses \texttt{h\ensuremath{\kappa} : 0 < \ensuremath{\kappa}}, \texttt{ha : 0 < a} and \texttt{h\ensuremath{\varepsilon}\ensuremath{{}_{0}} : 0 < \ensuremath{\varepsilon}\ensuremath{{}_{0}}}. \texttt{hr} is \(x>0\) (\texttt{positivity}), \texttt{hc} is the ceiling inequality \texttt{Nat.\allowbreak{}le\_\allowbreak{}ceil}, \texttt{hl} is \(\log x\le\log_+x\) (\texttt{le\_\allowbreak{}max\_\allowbreak{}left}), and \texttt{hexp} is \(-2k_0/\kappa\le-\log x\), obtained with \texttt{div\_\allowbreak{}le\_\allowbreak{}iff\ensuremath{{}_{0}} h\ensuremath{\kappa}} and \texttt{nlinarith}. The \texttt{calc} block first applies the monotonicity of the exponential (\texttt{Real.\allowbreak{}exp\_\allowbreak{}le\_\allowbreak{}exp.\allowbreak{}mpr hexp}) multiplied by \(\varepsilon_0^2\ge0\), then evaluates \(e^{-\log x}\varepsilon_0^2\) with \texttt{Real.\allowbreak{}exp\_\allowbreak{}neg}, \texttt{Real.\allowbreak{}exp\_\allowbreak{}log hr} and \texttt{field\_\allowbreak{}simp}.
\end{leanbox}
\subsection{\texorpdfstring{The error bound after the warmup}{The error bound after the warmup}}
\label{report-step-AverageLMC-Complexity-lean-3}


\textbf{Claim.} Let \(a>0\), \(\kappa\ge1\), \(\varepsilon_0>0\), let \(k_0=k_0(\kappa,a,\varepsilon_0)\) and let \(t_k\) be the schedule of Step 1. Let \((r_k)_{k\ge0}\) be real numbers with \(0\le r_0\le\varepsilon_0^2\) such that


\[r_{k+1}\le\frac{1-t_k}{1+t_k}\,r_k+\frac{a\,t_k^2}{1+t_k}\qquad\text{for every }k\in\mathbb N.\qquad(1)\]


Then for every \(k\ge k_0\),


\[r_k\le\frac a{\kappa+(k-k_0)}.\]


\textbf{Proof.} \emph{Warmup phase.} For \(j<k_0\) we have \((j-k_0)_+=0\), so \(t_j=1/\kappa\), and (1) is the constant-step recurrence with \(t=1/\kappa\). Since \(\kappa\ge1\), \(0<1/\kappa\le1\); moreover \(a\ge0\) and \(r_0\ge0\). The finite-horizon constant-step bound of the previous section, with horizon \(n=k_0\), gives


\[r_{k_0}\le e^{-2k_0/\kappa}\,r_0+\frac a{2\kappa},\]


because \(e^{-2(1/\kappa)k_0}=e^{-2k_0/\kappa}\) and \(\frac a2\cdot\frac1\kappa=\frac a{2\kappa}\). Multiplying \(r_0\le\varepsilon_0^2\) by \(e^{-2k_0/\kappa}\ge0\) and using Step 2,


\[e^{-2k_0/\kappa}\,r_0\le e^{-2k_0/\kappa}\,\varepsilon_0^2\le\frac a{2\kappa},\qquad\text{hence}\qquad r_{k_0}\le\frac a{2\kappa}+\frac a{2\kappa}=\frac a\kappa.\]


\emph{Decreasing phase.} Put \(s_n=r_{k_0+n}\) for \(n\in\mathbb N\). Then \(s_0\le a/\kappa\), and since \((k_0+n-k_0)_+=n\), the step at index \(k_0+n\) is \(t_{k_0+n}=1/(\kappa+n)\); thus (1) at \(k=k_0+n\) is exactly the recurrence (5) of the previous section for \((s_n)\). Its induction (part (2) of the last step there, with \(a\ge0\) and \(\kappa\ge1\)) gives \(s_n\le a/(\kappa+n)\) for every \(n\). For \(k\ge k_0\), take \(n=k-k_0\), so that \(k_0+n=k\). \(\square\)




Multiplying the conclusion by \(p/m\) gives \(\mathsf W_2^2(\nu_k,\pi)\le\frac{\kappa_{\rm av}+1}{\kappa_\infty+k-k_0}\cdot\frac pm\), the error bound of Proposition 1. The hypothesis \(\kappa\ge1\) is used twice: for \(t=1/\kappa\le1\) in the warmup and for \(x=\kappa+n\ge1\) in the decreasing phase; in the application it is \(M_\infty\ge m\).

\begin{leanbox}{Lean: \texttt{variable\_\allowbreak{}step\_\allowbreak{}normalized}}
\leansource{\texttt{variable\_\allowbreak{}step\_\allowbreak{}normalized} \quad AverageLMC/\allowbreak{}Complexity.lean, lines 59–96}
\begin{leancode}
\leanline{0}{theorem\ variable\_step\_normalized\ (r\ :\ ℕ\ →\ ℝ)\ (a\ κ\ ε₀\ :\ ℝ)}
\leanline{4}{(ha\ :\ 0\ <\ a)\ (hκ\ :\ 1\ ≤\ κ)\ (hε₀\ :\ 0\ <\ ε₀)}
\leanline{4}{(hr0\ :\ 0\ ≤\ r\ 0)\ (hinit\ :\ r\ 0\ ≤\ ε₀\ \textasciicircum{}\ 2)}
\leanline{4}{(hstep\ :\ ∀\ k,\ r\ (k\ +\ 1)\ ≤}
\leanline{6}{(1\ -\ variableStep\ κ\ (warmup\ κ\ a\ ε₀)\ k)\ /}
\leanline{8}{(1\ +\ variableStep\ κ\ (warmup\ κ\ a\ ε₀)\ k)\ *\ r\ k\ +}
\leanline{6}{a\ *\ variableStep\ κ\ (warmup\ κ\ a\ ε₀)\ k\ \textasciicircum{}\ 2\ /}
\leanline{8}{(1\ +\ variableStep\ κ\ (warmup\ κ\ a\ ε₀)\ k))}
\leanline{4}{(k\ :\ ℕ)\ (hk\ :\ warmup\ κ\ a\ ε₀\ ≤\ k)\ :}
\leanline{4}{r\ k\ ≤\ a\ /\ (κ\ +\ ((k\ -\ warmup\ κ\ a\ ε₀\ :\ ℕ)\ :\ ℝ))}
\end{leancode}
\leanlabel{How the Lean proof follows the argument}
\texttt{variable\_\allowbreak{}step\_\allowbreak{}normalized} states (1) as \texttt{hstep} with \texttt{variable\allowbreak{}Step \ensuremath{\kappa} (warmup \ensuremath{\kappa} a \ensuremath{\varepsilon}\ensuremath{{}_{0}}) k} in place of \(t_k\), and the conclusion for \texttt{k} with \texttt{hk : warmup \ensuremath{\kappa} a \ensuremath{\varepsilon}\ensuremath{{}_{0}} \ensuremath{\le} k}.



\begin{enumerate}


\item \texttt{hfirst} applies \texttt{constant\_\allowbreak{}step\_\allowbreak{}normalized\_\allowbreak{}up\_\allowbreak{}to} with \(t=1/\kappa\) and horizon \texttt{k\ensuremath{{}_{0}}}. Its finite-horizon hypothesis is checked for \(j<k_0\) by computing \texttt{hj0 : j -\allowbreak{} k\ensuremath{{}_{0}} = 0} (\texttt{Nat.\allowbreak{}sub\_\allowbreak{}eq\_\allowbreak{}zero\_\allowbreak{}of\_\allowbreak{}le}) and simplifying \texttt{hstep j} with it. \texttt{heq} rewrites \(-2\cdot\frac1\kappa\cdot k_0\) as \(-2k_0/\kappa\).

\item \texttt{hdecay} is \texttt{warmup\_\allowbreak{}decay}, \texttt{hmul} multiplies \texttt{hinit} by the nonnegative exponential (\texttt{Real.\allowbreak{}exp\_\allowbreak{}nonneg}), and \texttt{hwarm}, the bound \(r_{k_0}\le a/\kappa\), follows by \texttt{linarith} after rewriting \(a\cdot\frac1\kappa/2=\frac a{2\kappa}\) (\texttt{heq\ensuremath{{}_{2}}}) and \(\frac a{2\kappa}+\frac a{2\kappa}=\frac a\kappa\) (\texttt{hadd}).

\item \texttt{htail} applies \texttt{variable\_\allowbreak{}step\_\allowbreak{}tail} to the shifted sequence \texttt{fun n \ensuremath{\mapsto} r (k\ensuremath{{}_{0}} + n)}; the index identity \texttt{heq\ensuremath{{}_{3}}}, \((k_0+n)-k_0=n\) in \(\mathbb N\), is proved by \texttt{omega}, and \texttt{hstep (k\ensuremath{{}_{0}} + n)} rewritten with it is the shifted recurrence. Finally \texttt{Nat.\allowbreak{}add\_\allowbreak{}sub\_\allowbreak{}of\_\allowbreak{}le hk} gives \(k_0+(k-k_0)=k\).

\end{enumerate}
\end{leanbox}
\subsection{\texorpdfstring{The iteration count reaches the prescribed accuracy}{The iteration count reaches the prescribed accuracy}}
\label{report-step-AverageLMC-Complexity-lean-4}


\textbf{Definition.} For \(\kappa,a,\varepsilon_0,\varepsilon\in\mathbb R\), the iteration count of Proposition 1 is


\[k_\varepsilon=k_0(\kappa,a,\varepsilon_0)+\Bigl\lceil\max\Bigl(\frac a{\varepsilon^2}-\kappa,\,0\Bigr)\Bigr\rceil.\]


\textbf{Claim.} (1) Let \(\kappa>0\), \(a>0\) and \(\varepsilon>0\). If \(k\in\mathbb N\) satisfies \(k\ge k_\varepsilon\), then \(k\ge k_0\) and


\[\frac a{\kappa+(k-k_0)}\le\varepsilon^2.\]


(2) Under the hypotheses of Step 3 and \(\varepsilon>0\), \(r_k\le\varepsilon^2\) for every \(k\ge k_\varepsilon\).




\textbf{Proof.} (1) Write \(N=\lceil\max(a/\varepsilon^2-\kappa,0)\rceil\in\mathbb N\), so \(k\ge k_0+N\). Then \(k\ge k_0\), and \(k-k_0\ge N\) as natural numbers. By the definition of the ceiling and of the maximum,


\[k-k_0\ge N\ge\max\Bigl(\frac a{\varepsilon^2}-\kappa,0\Bigr)\ge\frac a{\varepsilon^2}-\kappa,\qquad\text{so}\qquad\kappa+(k-k_0)\ge\frac a{\varepsilon^2}.\]


The number \(\kappa+(k-k_0)\) is positive because \(\kappa>0\). Multiplying the last inequality by \(\varepsilon^2>0\) gives \(a\le\varepsilon^2(\kappa+(k-k_0))\), and dividing by \(\kappa+(k-k_0)>0\) gives the claim.




(2) Since \(\kappa\ge1>0\), part (1) gives \(k\ge k_0\) and \(a/(\kappa+(k-k_0))\le\varepsilon^2\). Step 3 gives \(r_k\le a/(\kappa+(k-k_0))\), and the two inequalities combine. \(\square\)




In the original variables, \(r_k\le\varepsilon^2\) is \(\mathsf W_2^2(\nu_k,\pi)\le\varepsilon^2p/m\), and \(k_\varepsilon=k_0+\lceil((\kappa_{\rm av}+1)/\varepsilon^2-\kappa_\infty)_+\rceil\) is the iteration count of Proposition 1. The hypothesis \(a>0\) of part (1) is not used in its proof.

\begin{leanbox}{Lean: \texttt{variableIterations}, \texttt{variable\_\allowbreak{}iterations\_\allowbreak{}envelope}, \texttt{variable\_\allowbreak{}step\_\allowbreak{}accuracy}}
\leansource{\texttt{variableIterations} \quad AverageLMC/\allowbreak{}Complexity.lean, lines 99–100}
\begin{leancode}
\leanline{0}{noncomputable\ def\ variableIterations\ (κ\ a\ ε₀\ ε\ :\ ℝ)\ :\ ℕ\ :=}
\leanline{2}{warmup\ κ\ a\ ε₀\ +\ ⌈max\ (a\ /\ ε\ \textasciicircum{}\ 2\ -\ κ)\ 0⌉₊}
\end{leancode}
\leansource{\texttt{variable\_\allowbreak{}iterations\_\allowbreak{}envelope} \quad AverageLMC/\allowbreak{}Complexity.lean, lines 104–123}
\begin{leancode}
\leanline{0}{theorem\ variable\_iterations\_envelope\ \{κ\ a\ ε₀\ ε\ :\ ℝ\}}
\leanline{4}{(hκ\ :\ 0\ <\ κ)\ (\_ha\ :\ 0\ <\ a)\ (hε\ :\ 0\ <\ ε)}
\leanline{4}{\{k\ :\ ℕ\}\ (hk\ :\ variableIterations\ κ\ a\ ε₀\ ε\ ≤\ k)\ :}
\leanline{4}{warmup\ κ\ a\ ε₀\ ≤\ k\ ∧}
\leanline{6}{a\ /\ (κ\ +\ ((k\ -\ warmup\ κ\ a\ ε₀\ :\ ℕ)\ :\ ℝ))\ ≤\ ε\ \textasciicircum{}\ 2}
\end{leancode}
\leansource{\texttt{variable\_\allowbreak{}step\_\allowbreak{}accuracy} \quad AverageLMC/\allowbreak{}Complexity.lean, lines 126–136}
\begin{leancode}
\leanline{0}{theorem\ variable\_step\_accuracy\ (r\ :\ ℕ\ →\ ℝ)\ (a\ κ\ ε₀\ ε\ :\ ℝ)}
\leanline{4}{(ha\ :\ 0\ <\ a)\ (hκ\ :\ 1\ ≤\ κ)\ (hε₀\ :\ 0\ <\ ε₀)\ (hε\ :\ 0\ <\ ε)}
\leanline{4}{(hr0\ :\ 0\ ≤\ r\ 0)\ (hinit\ :\ r\ 0\ ≤\ ε₀\ \textasciicircum{}\ 2)}
\leanline{4}{(hstep\ :\ ∀\ k,\ r\ (k\ +\ 1)\ ≤}
\leanline{6}{(1\ -\ variableStep\ κ\ (warmup\ κ\ a\ ε₀)\ k)\ /}
\leanline{8}{(1\ +\ variableStep\ κ\ (warmup\ κ\ a\ ε₀)\ k)\ *\ r\ k\ +}
\leanline{6}{a\ *\ variableStep\ κ\ (warmup\ κ\ a\ ε₀)\ k\ \textasciicircum{}\ 2\ /}
\leanline{8}{(1\ +\ variableStep\ κ\ (warmup\ κ\ a\ ε₀)\ k))}
\leanline{4}{(k\ :\ ℕ)\ (hk\ :\ variableIterations\ κ\ a\ ε₀\ ε\ ≤\ k)\ :\ r\ k\ ≤\ ε\ \textasciicircum{}\ 2}
\end{leancode}
\leanlabel{How the Lean proof follows the argument}
\begin{enumerate}


\item \texttt{variable\allowbreak{}Iterations \ensuremath{\kappa} a \ensuremath{\varepsilon}\ensuremath{{}_{0}} \ensuremath{\varepsilon}} is \texttt{warmup \ensuremath{\kappa} a \ensuremath{\varepsilon}\ensuremath{{}_{0}} + ⌈max (a /\allowbreak{} \ensuremath{\varepsilon} \textasciicircum{} 2 -\allowbreak{} \ensuremath{\kappa}) 0⌉\ensuremath{{}_{+}}}, the count \(k_\varepsilon\).

\item \texttt{variable\_\allowbreak{}iterations\_\allowbreak{}envelope} has the hypotheses \texttt{h\ensuremath{\kappa} : 0 < \ensuremath{\kappa}}, \texttt{\_\allowbreak{}ha : 0 < a} (unused) and \texttt{h\ensuremath{\varepsilon} : 0 < \ensuremath{\varepsilon}}, and returns the conjunction of \(k\ge k_0\) and the envelope bound. \texttt{hk\ensuremath{{}_{0}}} and \texttt{hnat} are the two natural-number inequalities, obtained by \texttt{omega} after unfolding the count; \texttt{hc} passes to real numbers through \texttt{Nat.\allowbreak{}le\_\allowbreak{}ceil} and \texttt{exact\_\allowbreak{}mod\_\allowbreak{}cast}; \texttt{hden} is the positivity of \(\kappa+(k-k_0)\). The goal is reduced by \texttt{div\_\allowbreak{}le\_\allowbreak{}iff\ensuremath{{}_{0}} hden} to \(a\le\varepsilon^2(\kappa+(k-k_0))\), which \texttt{nlinarith} derives from \texttt{hq} (\(a/\varepsilon^2\le\kappa+(k-k_0)\)) cleared of its denominator by \texttt{div\_\allowbreak{}le\_\allowbreak{}iff\ensuremath{{}_{0}} he}.

\item \texttt{variable\_\allowbreak{}step\_\allowbreak{}accuracy} chains \texttt{variable\_\allowbreak{}step\_\allowbreak{}normalized}, applied with \texttt{he.\allowbreak{}1} (\(k\ge k_0\)), and the envelope bound \texttt{he.\allowbreak{}2}.

\end{enumerate}
\end{leanbox}
\section{\texorpdfstring{Affine transport bounds for arbitrary initial distributions}{Affine transport bounds for arbitrary initial distributions}}
\label{proof-module-AverageLMC-ExtendedTransport-lean}


The main theorems are stated for every initial probability distribution, including distributions at infinite Wasserstein distance from the target. For this purpose distances are taken in the extended half-line \([0,\infty]\), with the conventions \(x+\infty=\infty\) and \(c\cdot\infty=\infty\) for \(c>0\). This section shows that an affine bound \(\mathsf W_2^2(\mu',\pi')\le a\,\mathsf W_2^2(\mu,\pi)+b\), proved under the assumption \(\mathsf W_2(\mu,\pi)<\infty\), holds in \([0,\infty]\) without that assumption, provided \(a>0\) and \(b\ge0\): if the initial distance is infinite, the right-hand side is infinite. No moment assumption on the initial distribution is therefore needed in the final statements.

\subsection{\texorpdfstring{Extending a finite-distance bound to infinite initial distance}{Extending a finite-distance bound to infinite initial distance}}
\label{report-step-AverageLMC-ExtendedTransport-lean-1}


\textbf{Claim.} Let \(\mu,\pi,\mu',\pi'\) be measures on a normed space \(E\), and let \(a>0\) and \(b\ge0\). Assume that, whenever \(\mathsf W_2(\mu,\pi)<\infty\), also \(\mathsf W_2(\mu',\pi')<\infty\) and


\[\mathsf W_2(\mu',\pi')^2\le a\,\mathsf W_2(\mu,\pi)^2+b\qquad\text{(an inequality between real numbers).}\]


Then, in \([0,\infty]\), \(\mathsf W_2(\mu',\pi')^2\le a\,\mathsf W_2(\mu,\pi)^2+b\) holds without any finiteness assumption.




\textbf{Proof.} If \(\mathsf W_2(\mu,\pi)=\infty\), then \(\mathsf W_2(\mu,\pi)^2=\infty\) and \(a\cdot\infty=\infty\) because \(a>0\), so the right-hand side is \(\infty\) and the inequality holds. If \(\mathsf W_2(\mu,\pi)<\infty\), then by assumption \(\mathsf W_2(\mu',\pi')<\infty\), so both distances are real numbers and the real inequality holds. The inclusion of \([0,\infty)\) into \([0,\infty]\) preserves order, sums and products of nonnegative numbers; since \(a\,\mathsf W_2(\mu,\pi)^2\ge0\) and \(b\ge0\), the real inequality is also an inequality in \([0,\infty]\). \(\square\)




Both sign conditions are needed. If \(a=0\), then \(0\cdot\infty=0\) in \([0,\infty]\) and an infinite initial distance would not make the right-hand side infinite. If \(b<0\), the right-hand side could not be interpreted in \([0,\infty]\).

\begin{leanbox}{Lean: \texttt{squared\_\allowbreak{}transport\_\allowbreak{}extended}}
\leansource{\texttt{squared\_\allowbreak{}transport\_\allowbreak{}extended} \quad AverageLMC/\allowbreak{}ExtendedTransport.lean, lines 12–26}
\begin{leancode}
\leanline{0}{theorem\ squared\_transport\_extended\ \{E\ :\ Type*\}\ [NormedAddCommGroup\ E]\ [MeasurableSpace\ E]}
\leanline{4}{(μ\ π\ μ'\ π'\ :\ Measure\ E)\ (a\ b\ :\ ℝ)\ (ha\ :\ 0\ <\ a)\ (hb\ :\ 0\ ≤\ b)}
\leanline{4}{(houtput\ :\ ULA.wasserstein2\ μ\ π\ ≠\ ⊤\ →\ ULA.wasserstein2\ μ'\ π'\ ≠\ ⊤)}
\leanline{4}{(hbound\ :\ ULA.wasserstein2\ μ\ π\ ≠\ ⊤\ →}
\leanline{6}{ULA.wasserstein2Real\ μ'\ π'\ \textasciicircum{}\ 2\ ≤\ a\ *\ ULA.wasserstein2Real\ μ\ π\ \textasciicircum{}\ 2\ +\ b)\ :}
\leanline{4}{ULA.wasserstein2\ μ'\ π'\ \textasciicircum{}\ 2\ ≤}
\leanline{6}{ENNReal.ofReal\ a\ *\ ULA.wasserstein2\ μ\ π\ \textasciicircum{}\ 2\ +\ ENNReal.ofReal\ b}
\end{leancode}
\leanlabel{How the Lean proof follows the argument}
\texttt{squared\_\allowbreak{}transport\_\allowbreak{}extended} states the conclusion in \texttt{\ensuremath{\mathbb R}\ensuremath{\ge}0\ensuremath{\infty}}: \texttt{ULA.\allowbreak{}wasserstein2 \ensuremath{\mu}' \ensuremath{\pi}' \textasciicircum{} 2 \ensuremath{\le} ENNReal.\allowbreak{}of\allowbreak{}Real a * ULA.\allowbreak{}wasserstein2 \ensuremath{\mu} \ensuremath{\pi} \textasciicircum{} 2 + ENNReal.\allowbreak{}of\allowbreak{}Real b}, where \texttt{ENNReal.\allowbreak{}of\allowbreak{}Real} embeds a real number into \([0,\infty]\) (negative numbers are sent to \(0\)). The hypotheses are \texttt{ha : 0 < a}, \texttt{hb : 0 \ensuremath{\le} b}, \texttt{houtput} (finiteness of the output distance when the input distance is finite) and \texttt{hbound} (the real inequality under the same condition, written with \texttt{ULA.\allowbreak{}wasserstein2\allowbreak{}Real}).



\begin{enumerate}


\item \texttt{by\_\allowbreak{}cases hi : ULA.\allowbreak{}wasserstein2 \ensuremath{\mu} \ensuremath{\pi} = \ensuremath{\top}}. In the infinite case, \texttt{simp} evaluates the right-hand side to \(\top=\infty\), using \texttt{ENNReal.\allowbreak{}of\allowbreak{}Real\_\allowbreak{}pos.\allowbreak{}mpr ha} to know that the coefficient is nonzero.

\item In the finite case, \texttt{h} applies the monotone map \texttt{ENNReal.\allowbreak{}of\allowbreak{}Real} to \texttt{hbound hi}; \texttt{ENNReal.\allowbreak{}of\allowbreak{}Real\_\allowbreak{}add} and \texttt{ENNReal.\allowbreak{}of\allowbreak{}Real\_\allowbreak{}mul} split the image of the right-hand side (using the nonnegativity of both summands), and \texttt{simpa only} with \texttt{ENNReal.\allowbreak{}of\allowbreak{}Real\_\allowbreak{}to\allowbreak{}Real} for the two finite distances (\texttt{houtput hi} and \texttt{hi}) identifies the images of the real distances with the extended ones.

\end{enumerate}
\end{leanbox}


Earlier results used here: \cite[Section~3.17.3]{LangevinLean2026ULASC}; \cite[Section~3.17.4]{LangevinLean2026ULASC}.
\section{\texorpdfstring{Actual-law constant and varying-step averaged LMC bounds}{Actual-law constant and varying-step averaged LMC bounds}}
\label{proof-module-AverageLMC-Theorem-lean}


The scalar recurrences are now instantiated by the actual Gaussian algorithm. Its marginal probability laws and finite moments are derived at every step. The exact constant-step exponential bound and the complete variable-step proposition, including its warmup and accuracy ceiling, follow with the paper’s constants.



The actual Gaussian LMC laws satisfy the constant-step bound, including arbitrary initial probability laws in extended transport, when \(0\le h\le1/L\):


\[W_2^2(\nu_k,\pi)\le e^{-2mkh}W_2^2(\nu_0,\pi)+\frac{(M_{\rm av}+m)hp}{2m}.\]


For Proposition 1 put \(\kappa=L/m\), \(a=M_{\rm av}/m+1\), and assume \(W_2^2(\nu_0,\pi)\le\varepsilon_0^2p/m\). With \(\log_+u=\max\{\log u,0\}\), the exact schedule and error envelope are


\[\begin{aligned}k_0&=\left\lceil\frac\kappa2\log_+\!\left(\frac{2\kappa\varepsilon_0^2}{a}\right)\right\rceil,\qquad h_k=\frac1{L+m(k-k_0)_+},\\ W_2^2(\nu_k,\pi)&\le\frac{a}{\kappa+k-k_0}\frac pm\qquad(k\ge k_0).\end{aligned}\]\subsection{\texorpdfstring{Define actual marginal laws and propagate finite moments}{Define actual marginal laws and propagate finite moments}}
\label{report-step-AverageLMC-Theorem-lean-1}


The recursive law at step \(k+1\) is the independent Gaussian Euler transition at the deterministic step \(h(k)\), applied to the current law. The zero and successor identities expose this recursion. A pushforward of a product of probability laws remains a probability law, proving normalization at every \(k\). For each nonnegative step, the already constructed \(L^{2}\) Brownian realization gives finite output transport whenever initial transport is finite; the zero-step identity covers \(h=0\). Induction gives finite moments for every law of an arbitrary nonnegative deterministic schedule.

\begin{leanbox}{Lean: \texttt{lmcLaws}, \texttt{lmcLaws\_\allowbreak{}zero}, \texttt{lmcLaws\_\allowbreak{}succ} and 3 more}
\leansource{\texttt{lmcLaws} \quad AverageLMC/\allowbreak{}Theorem.lean, lines 13–16}
\begin{leancode}
\leanline{0}{def\ lmcLaws\ \{p\ :\ ℕ\}\ (f\ :\ State\ p\ →\ ℝ)\ (steps\ :\ ℕ\ →\ ℝ)}
\leanline{4}{(μ\ :\ Measure\ (State\ p))\ :\ ℕ\ →\ Measure\ (State\ p)}
\leanline{2}{|\ 0\ =>\ μ}
\leanline{2}{|\ k\ +\ 1\ =>\ ulaLawStep\ f\ (steps\ k)\ (lmcLaws\ f\ steps\ μ\ k)}
\end{leancode}
\leansource{\texttt{lmcLaws\_\allowbreak{}zero} \quad AverageLMC/\allowbreak{}Theorem.lean, lines 18–19}
\begin{leancode}
\leanline{0}{@[simp]\ theorem\ lmcLaws\_zero\ \{p\ :\ ℕ\}\ (f\ :\ State\ p\ →\ ℝ)\ (steps\ :\ ℕ\ →\ ℝ)}
\leanline{4}{(μ\ :\ Measure\ (State\ p))\ :\ lmcLaws\ f\ steps\ μ\ 0\ =\ μ}
\end{leancode}
\leansource{\texttt{lmcLaws\_\allowbreak{}succ} \quad AverageLMC/\allowbreak{}Theorem.lean, lines 21–23}
\begin{leancode}
\leanline{0}{@[simp]\ theorem\ lmcLaws\_succ\ \{p\ :\ ℕ\}\ (f\ :\ State\ p\ →\ ℝ)\ (steps\ :\ ℕ\ →\ ℝ)}
\leanline{4}{(μ\ :\ Measure\ (State\ p))\ (k\ :\ ℕ)\ :}
\leanline{4}{lmcLaws\ f\ steps\ μ\ (k\ +\ 1)\ =\ ulaLawStep\ f\ (steps\ k)\ (lmcLaws\ f\ steps\ μ\ k)}
\end{leancode}
\leansource{\texttt{lmcLaws\_\allowbreak{}isProbability} \quad AverageLMC/\allowbreak{}Theorem.lean, lines 25–32}
\begin{leancode}
\leanline{0}{instance\ lmcLaws\_isProbability\ \{p\ :\ ℕ\}\ (f\ :\ State\ p\ →\ ℝ)\ (steps\ :\ ℕ\ →\ ℝ)}
\leanline{4}{(μ\ :\ Measure\ (State\ p))\ [IsProbabilityMeasure\ μ]\ (k\ :\ ℕ)\ :}
\leanline{4}{IsProbabilityMeasure\ (lmcLaws\ f\ steps\ μ\ k)\ :=\ by}
\leanline{2}{induction\ k\ with}
\leanline{2}{|\ zero\ =>\ exact\ inferInstanceAs\ (IsProbabilityMeasure\ μ)}
\leanline{2}{|\ succ\ k\ ih\ =>}
\leanline{4}{change\ IsProbabilityMeasure\ ((lmcLaws\ f\ steps\ μ\ k).prod\ (stdGaussian\ (State\ p))\ |>.map\ \_)}
\leanline{4}{infer\_instance}
\end{leancode}
\leansource{\texttt{lmc\_\allowbreak{}step\_\allowbreak{}finite} \quad AverageLMC/\allowbreak{}Theorem.lean, lines 35–43}
\begin{leancode}
\leanline{0}{theorem\ lmc\_step\_finite\ \{p\ :\ ℕ\}\ (f\ :\ State\ p\ →\ ℝ)\ (m\ h\ :\ ℝ)\ (L\ :\ ℝ≥0)}
\leanline{4}{(hm\ :\ 0\ <\ m)\ (hh\ :\ 0\ ≤\ h)\ (hf\ :\ Differentiable\ ℝ\ f)}
\leanline{4}{(hsc\ :\ StrongConvexOn\ univ\ m\ f)\ (hLip\ :\ LipschitzWith\ L\ (gradient\ f))}
\leanline{4}{(μ\ :\ Measure\ (State\ p))\ [IsProbabilityMeasure\ μ]}
\leanline{4}{(hfinite\ :\ wasserstein2\ μ\ (gibbsMeasure\ f)\ ≠\ ⊤)\ :}
\leanline{4}{wasserstein2\ (ulaLawStep\ f\ h\ μ)\ (gibbsMeasure\ f)\ ≠\ ⊤}
\end{leancode}
\leansource{\texttt{lmcLaws\_\allowbreak{}finite} \quad AverageLMC/\allowbreak{}Theorem.lean, lines 46–54}
\begin{leancode}
\leanline{0}{theorem\ lmcLaws\_finite\ \{p\ :\ ℕ\}\ (f\ :\ State\ p\ →\ ℝ)\ (m\ :\ ℝ)\ (L\ :\ ℝ≥0)}
\leanline{4}{(hm\ :\ 0\ <\ m)\ (hf\ :\ Differentiable\ ℝ\ f)\ (hsc\ :\ StrongConvexOn\ univ\ m\ f)}
\leanline{4}{(hLip\ :\ LipschitzWith\ L\ (gradient\ f))\ (steps\ :\ ℕ\ →\ ℝ)\ (hsteps\ :\ ∀\ k,\ 0\ ≤\ steps\ k)}
\leanline{4}{(μ\ :\ Measure\ (State\ p))\ [IsProbabilityMeasure\ μ]}
\leanline{4}{(hfinite\ :\ wasserstein2\ μ\ (gibbsMeasure\ f)\ ≠\ ⊤)\ (k\ :\ ℕ)\ :}
\leanline{4}{wasserstein2\ (lmcLaws\ f\ steps\ μ\ k)\ (gibbsMeasure\ f)\ ≠\ ⊤}
\end{leancode}
\leanlabel{How the Lean proof follows the argument}
\texttt{lmc\allowbreak{}Laws} names the mathematical object defined above. Its defining expression appears in the adjacent statement; subsequent identities may unfold this definition directly.




\texttt{lmc\allowbreak{}Laws\_\allowbreak{}zero} is a definitional equality: \texttt{rfl} closes it because both sides reduce to the same expression.




\texttt{lmc\allowbreak{}Laws\_\allowbreak{}succ} is a definitional equality: \texttt{rfl} closes it because both sides reduce to the same expression.




\texttt{lmc\allowbreak{}Laws\_\allowbreak{}is\allowbreak{}Probability} names the mathematical object defined above. Its defining expression appears in the adjacent statement; subsequent identities may unfold this definition directly.




\texttt{lmc\_\allowbreak{}step\_\allowbreak{}finite} gives the formal counterpart of this step. Its remaining equalities are discharged by the simplifications displayed in the proof.




\texttt{lmc\allowbreak{}Laws\_\allowbreak{}finite} proceeds by \texttt{induction k with}. The base cases establish the initial claims and the inductive branches use the induction hypotheses.
\end{leanbox}


Earlier results used here: \cite[Section~3.17.3]{LangevinLean2026ULASC}; \cite[Section~3.26.3]{LangevinLean2026ULASC}; \cite[Section~3.36.7]{LangevinLean2026ULASC}; \cite[Section~3.45.6]{LangevinLean2026ULASC}.
\subsection{\texorpdfstring{Remove the nonpositive score correction and iterate a constant step}{Remove the nonpositive score correction and iterate a constant step}}
\label{report-step-AverageLMC-Theorem-lean-2}


When \(hL\le1\), the score correction coefficient is nonpositive and the score second moment is nonnegative, so the exact one-step theorem implies the two-term affine recurrence. A constant schedule agrees by induction with iteration of the actual Gaussian transition. Strong convexity and coordinate smoothness imply \(m\le M_{\mathrm{av}}\le L\), hence \(0\le mh\le 1\). Apply the scalar recurrence theorem with \(t=mh\) and forcing coefficient \((M_{\mathrm{av}}+m)p/m^2\). Its exponential contraction becomes \(\exp(-2mkh)\), and its fixed point becomes exactly \((M_{\mathrm{av}}+m)hp/(2m)\). This proves Theorem 1’s global claim for actual LMC laws.

\begin{leanbox}{Lean: \texttt{lmc\_\allowbreak{}one\_\allowbreak{}step\_\allowbreak{}stable}, \texttt{lmcLaws\_\allowbreak{}const}, \texttt{lmc\_\allowbreak{}constant\_\allowbreak{}step}}
\leansource{\texttt{lmc\_\allowbreak{}one\_\allowbreak{}step\_\allowbreak{}stable} \quad AverageLMC/\allowbreak{}Theorem.lean, lines 57–73}
\begin{leancode}
\leanline{0}{theorem\ lmc\_one\_step\_stable\ \{p\ :\ ℕ\}\ (hp\ :\ 0\ <\ p)\ (f\ :\ State\ p\ →\ ℝ)}
\leanline{4}{(m\ h\ :\ ℝ)\ (L\ :\ ℝ≥0)\ (hm\ :\ 0\ <\ m)\ (hh\ :\ 0\ ≤\ h)}
\leanline{4}{(hf\ :\ Differentiable\ ℝ\ f)\ (hsc\ :\ StrongConvexOn\ univ\ m\ f)}
\leanline{4}{(hLip\ :\ LipschitzWith\ L\ (gradient\ f))\ (hstep\ :\ h\ *\ (L\ :\ ℝ)\ ≤\ 1)}
\leanline{4}{(μ\ :\ Measure\ (State\ p))\ [IsProbabilityMeasure\ μ]}
\leanline{4}{(hfinite\ :\ wasserstein2\ μ\ (gibbsMeasure\ f)\ ≠\ ⊤)\ :}
\leanline{4}{wasserstein2Real\ (ulaLawStep\ f\ h\ μ)\ (gibbsMeasure\ f)\ \textasciicircum{}\ 2\ ≤}
\leanline{6}{(1\ -\ m\ *\ h)\ /\ (1\ +\ m\ *\ h)\ *\ wasserstein2Real\ μ\ (gibbsMeasure\ f)\ \textasciicircum{}\ 2\ +}
\leanline{6}{((averageSmoothness\ f\ +\ m)\ *\ h\ \textasciicircum{}\ 2\ *\ p)\ /\ (1\ +\ m\ *\ h)}
\end{leancode}
\leansource{\texttt{lmcLaws\_\allowbreak{}const} \quad AverageLMC/\allowbreak{}Theorem.lean, lines 76–81}
\begin{leancode}
\leanline{0}{theorem\ lmcLaws\_const\ \{p\ :\ ℕ\}\ (f\ :\ State\ p\ →\ ℝ)\ (h\ :\ ℝ)}
\leanline{4}{(μ\ :\ Measure\ (State\ p))\ (k\ :\ ℕ)\ :}
\leanline{4}{lmcLaws\ f\ (fun\ \_\ ↦\ h)\ μ\ k\ =\ (ulaLawStep\ f\ h)\textasciicircum{}[k]\ μ}
\end{leancode}
\leansource{\texttt{lmc\_\allowbreak{}constant\_\allowbreak{}step} \quad AverageLMC/\allowbreak{}Theorem.lean, lines 84–111}
\begin{leancode}
\leanline{0}{theorem\ lmc\_constant\_step\ \{p\ :\ ℕ\}\ (hp\ :\ 0\ <\ p)\ (f\ :\ State\ p\ →\ ℝ)}
\leanline{4}{(m\ h\ :\ ℝ)\ (L\ :\ ℝ≥0)\ (hm\ :\ 0\ <\ m)\ (hh\ :\ 0\ ≤\ h)}
\leanline{4}{(hf\ :\ Differentiable\ ℝ\ f)\ (hsc\ :\ StrongConvexOn\ univ\ m\ f)}
\leanline{4}{(hLip\ :\ LipschitzWith\ L\ (gradient\ f))\ (hstep\ :\ h\ *\ (L\ :\ ℝ)\ ≤\ 1)}
\leanline{4}{(μ\ :\ Measure\ (State\ p))\ [IsProbabilityMeasure\ μ]}
\leanline{4}{(hfinite\ :\ wasserstein2\ μ\ (gibbsMeasure\ f)\ ≠\ ⊤)\ (k\ :\ ℕ)\ :}
\leanline{4}{wasserstein2Real\ ((ulaLawStep\ f\ h)\textasciicircum{}[k]\ μ)\ (gibbsMeasure\ f)\ \textasciicircum{}\ 2\ ≤}
\leanline{6}{Real.exp\ (-2\ *\ m\ *\ (k\ :\ ℝ)\ *\ h)\ *\ wasserstein2Real\ μ\ (gibbsMeasure\ f)\ \textasciicircum{}\ 2\ +}
\leanline{8}{(averageSmoothness\ f\ +\ m)\ *\ h\ *\ p\ /\ (2\ *\ m)}
\end{leancode}
\leanlabel{How the Lean proof follows the argument}
\texttt{lmc\_\allowbreak{}one\_\allowbreak{}step\_\allowbreak{}stable} first proves that the score correction is nonpositive. \texttt{lmc\_\allowbreak{}constant\_\allowbreak{}step} defines the actual squared-distance sequence as \texttt{r}; \texttt{lmc\allowbreak{}Laws\_\allowbreak{}finite} allows the one-step theorem to be applied at every index. The block \texttt{hs} is the scalar recurrence after normalization. \texttt{constant\_\allowbreak{}step\_\allowbreak{}normalized} iterates it and uses \texttt{contraction\_\allowbreak{}le\_\allowbreak{}exp}. Finally \texttt{lmc\allowbreak{}Laws\_\allowbreak{}const} identifies the recursively constructed marginal laws with iteration of the Gaussian transition.
\end{leanbox}


Earlier results used here: \cite[Section~3.17.3]{LangevinLean2026ULASC}; \cite[Section~3.17.4]{LangevinLean2026ULASC}; \cite[Section~3.26.3]{LangevinLean2026ULASC}; \cite[Section~3.36.7]{LangevinLean2026ULASC}.
\subsection{\texorpdfstring{Normalize the variable schedule without changing its constants}{Normalize the variable schedule without changing its constants}}
\label{report-step-AverageLMC-Theorem-lean-3}


Multiplying the proved stable one-step inequality by \(m/p\) gives the dimensionless recurrence for \(r(k)=m\) \(W_{2}^{2}/p\) with \(a=M_{\mathrm{av}}/m+1\). The schedule is initially written as the dimensionless schedule divided by \(m\); elementary division shows it is exactly \(1/(L+m(k-k_{0})_{+})\), as displayed in Proposition 1. The scalar warmup proof is extended to zero initial error: then \(k_{0}=0\) and the harmonic tail applies immediately. For nonzero initial error, its absolute value has the same square and the same warmup, so the existing positive-parameter scalar theorem applies without adding a sign restriction.

\begin{leanbox}{Lean: \texttt{lmc\_\allowbreak{}normalized\_\allowbreak{}step}, \texttt{lmcVariableSchedule}, \texttt{lmcVariableSchedule\_\allowbreak{}eq} and 1 more}
\leansource{\texttt{lmc\_\allowbreak{}normalized\_\allowbreak{}step} \quad AverageLMC/\allowbreak{}Theorem.lean, lines 114–128}
\begin{leancode}
\leanline{0}{theorem\ lmc\_normalized\_step\ \{p\ :\ ℕ\}\ (hp\ :\ 0\ <\ p)\ (f\ :\ State\ p\ →\ ℝ)}
\leanline{4}{(m\ h\ :\ ℝ)\ (L\ :\ ℝ≥0)\ (hm\ :\ 0\ <\ m)\ (hh\ :\ 0\ ≤\ h)}
\leanline{4}{(hf\ :\ Differentiable\ ℝ\ f)\ (hsc\ :\ StrongConvexOn\ univ\ m\ f)}
\leanline{4}{(hLip\ :\ LipschitzWith\ L\ (gradient\ f))\ (hstep\ :\ h\ *\ (L\ :\ ℝ)\ ≤\ 1)}
\leanline{4}{(μ\ :\ Measure\ (State\ p))\ [IsProbabilityMeasure\ μ]}
\leanline{4}{(hfinite\ :\ wasserstein2\ μ\ (gibbsMeasure\ f)\ ≠\ ⊤)\ :}
\leanline{4}{m\ /\ p\ *\ wasserstein2Real\ (ulaLawStep\ f\ h\ μ)\ (gibbsMeasure\ f)\ \textasciicircum{}\ 2\ ≤}
\leanline{6}{(1\ -\ m\ *\ h)\ /\ (1\ +\ m\ *\ h)\ *}
\leanline{8}{(m\ /\ p\ *\ wasserstein2Real\ μ\ (gibbsMeasure\ f)\ \textasciicircum{}\ 2)\ +}
\leanline{6}{(averageSmoothness\ f\ /\ m\ +\ 1)\ *\ (m\ *\ h)\ \textasciicircum{}\ 2\ /\ (1\ +\ m\ *\ h)}
\end{leancode}
\leansource{\texttt{lmcVariableSchedule} \quad AverageLMC/\allowbreak{}Theorem.lean, lines 131–132}
\begin{leancode}
\leanline{0}{def\ lmcVariableSchedule\ (M\ m\ :\ ℝ)\ (k₀\ k\ :\ ℕ)\ :\ ℝ\ :=}
\leanline{2}{variableStep\ (M\ /\ m)\ k₀\ k\ /\ m}
\end{leancode}
\leansource{\texttt{lmcVariableSchedule\_\allowbreak{}eq} \quad AverageLMC/\allowbreak{}Theorem.lean, lines 135–141}
\begin{leancode}
\leanline{0}{theorem\ lmcVariableSchedule\_eq\ \{M\ m\ :\ ℝ\}\ (hm\ :\ 0\ <\ m)\ (hM\ :\ 0\ <\ M)}
\leanline{4}{(k₀\ k\ :\ ℕ)\ :}
\leanline{4}{lmcVariableSchedule\ M\ m\ k₀\ k\ =\ 1\ /\ (M\ +\ m\ *\ ((k\ -\ k₀\ :\ ℕ)\ :\ ℝ))}
\end{leancode}
\leansource{\texttt{variable\_\allowbreak{}step\_\allowbreak{}normalized\_\allowbreak{}all} \quad AverageLMC/\allowbreak{}Theorem.lean, lines 144–164}
\begin{leancode}
\leanline{0}{private\ theorem\ variable\_step\_normalized\_all\ (r\ :\ ℕ\ →\ ℝ)\ (a\ κ\ ε₀\ :\ ℝ)}
\leanline{4}{(ha\ :\ 0\ <\ a)\ (hκ\ :\ 1\ ≤\ κ)\ (hr0\ :\ 0\ ≤\ r\ 0)\ (hinit\ :\ r\ 0\ ≤\ ε₀\ \textasciicircum{}\ 2)}
\leanline{4}{(hstep\ :\ ∀\ k,\ r\ (k\ +\ 1)\ ≤}
\leanline{6}{(1\ -\ variableStep\ κ\ (warmup\ κ\ a\ ε₀)\ k)\ /}
\leanline{8}{(1\ +\ variableStep\ κ\ (warmup\ κ\ a\ ε₀)\ k)\ *\ r\ k\ +}
\leanline{6}{a\ *\ variableStep\ κ\ (warmup\ κ\ a\ ε₀)\ k\ \textasciicircum{}\ 2\ /}
\leanline{8}{(1\ +\ variableStep\ κ\ (warmup\ κ\ a\ ε₀)\ k))}
\leanline{4}{(k\ :\ ℕ)\ (hk\ :\ warmup\ κ\ a\ ε₀\ ≤\ k)\ :}
\leanline{4}{r\ k\ ≤\ a\ /\ (κ\ +\ ((k\ -\ warmup\ κ\ a\ ε₀\ :\ ℕ)\ :\ ℝ))}
\end{leancode}
\leanlabel{How the Lean proof follows the argument}
\texttt{lmc\_\allowbreak{}normalized\_\allowbreak{}step} multiplies the stable one-step inequality by \(m/p>0\) and simplifies the fractions. \texttt{lmc\allowbreak{}Variable\allowbreak{}Schedule} divides the scalar normalized step by \(m\); \texttt{lmc\allowbreak{}Variable\allowbreak{}Schedule\_\allowbreak{}eq} proves its equality to the displayed reciprocal schedule by clearing positive denominators. \texttt{variable\_\allowbreak{}step\_\allowbreak{}normalized\_\allowbreak{}all} splits \texttt{h\ensuremath{\varepsilon}}, the zero-initial-error case. There the warmup is zero and \texttt{variable\_\allowbreak{}step\_\allowbreak{}tail} applies directly. Otherwise it uses the unchanged square of the absolute initial-error parameter to invoke \texttt{variable\_\allowbreak{}step\_\allowbreak{}normalized}.
\end{leanbox}


Earlier results used here: \cite[Section~3.17.3]{LangevinLean2026ULASC}; \cite[Section~3.17.4]{LangevinLean2026ULASC}; \cite[Section~3.26.3]{LangevinLean2026ULASC}; \cite[Section~3.36.7]{LangevinLean2026ULASC}.
\subsection{\texorpdfstring{Prove the entire variable-step proposition for the actual chain}{Prove the entire variable-step proposition for the actual chain}}
\label{report-step-AverageLMC-Theorem-lean-4}


The coordinate comparison gives \(\kappa =L/m\ge 1\) and \(a>0\). The chosen schedule is positive and satisfies \(h(k)L\le 1\) at every step. Use the actual-law normalized estimate and the derived finite moments to discharge every recurrence premise. The warmup and harmonic-tail theorem then gives \(r(k)\le a/(\kappa +k-k_{0})\) for all \(k\ge k_{0}\). Multiply back by \(p/m\) to recover the exact paper envelope. Finally, the explicit positive-part ceiling defining the sufficient iteration count makes this envelope at most \(\varepsilon ^{2}\); multiplying by \(p/m\) proves the stated sufficient iteration threshold for the actual algorithm.

\begin{leanbox}{Lean: \texttt{lmc\_\allowbreak{}variable\_\allowbreak{}step}, \texttt{lmc\_\allowbreak{}variable\_\allowbreak{}step\_\allowbreak{}accuracy}}
\leansource{\texttt{lmc\_\allowbreak{}variable\_\allowbreak{}step} \quad AverageLMC/\allowbreak{}Theorem.lean, lines 168–227}
\begin{leancode}
\leanline{0}{theorem\ lmc\_variable\_step\ \{p\ :\ ℕ\}\ (hp\ :\ 0\ <\ p)\ (f\ :\ State\ p\ →\ ℝ)}
\leanline{4}{(m\ ε₀\ :\ ℝ)\ (L\ :\ ℝ≥0)\ (hm\ :\ 0\ <\ m)}
\leanline{4}{(hf\ :\ Differentiable\ ℝ\ f)\ (hsc\ :\ StrongConvexOn\ univ\ m\ f)}
\leanline{4}{(hLip\ :\ LipschitzWith\ L\ (gradient\ f))}
\leanline{4}{(μ\ :\ Measure\ (State\ p))\ [IsProbabilityMeasure\ μ]}
\leanline{4}{(hfinite\ :\ wasserstein2\ μ\ (gibbsMeasure\ f)\ ≠\ ⊤)}
\leanline{4}{(hinit\ :\ wasserstein2Real\ μ\ (gibbsMeasure\ f)\ \textasciicircum{}\ 2\ ≤\ ε₀\ \textasciicircum{}\ 2\ *\ p\ /\ m)}
\leanline{4}{(k\ :\ ℕ)\ (hk\ :\ warmup\ ((L\ :\ ℝ)\ /\ m)\ (averageSmoothness\ f\ /\ m\ +\ 1)\ ε₀\ ≤\ k)\ :}
\leanline{4}{let\ κ\ :=\ (L\ :\ ℝ)\ /\ m}
\leanline{4}{let\ a\ :=\ averageSmoothness\ f\ /\ m\ +\ 1}
\leanline{4}{let\ k₀\ :=\ warmup\ κ\ a\ ε₀}
\leanline{4}{wasserstein2Real\ (lmcLaws\ f\ (lmcVariableSchedule\ (L\ :\ ℝ)\ m\ k₀)\ μ\ k)}
\leanline{6}{(gibbsMeasure\ f)\ \textasciicircum{}\ 2\ ≤}
\leanline{6}{a\ /\ (κ\ +\ ((k\ -\ k₀\ :\ ℕ)\ :\ ℝ))\ *\ p\ /\ m}
\end{leancode}
\leansource{\texttt{lmc\_\allowbreak{}variable\_\allowbreak{}step\_\allowbreak{}accuracy} \quad AverageLMC/\allowbreak{}Theorem.lean, lines 231–252}
\begin{leancode}
\leanline{0}{theorem\ lmc\_variable\_step\_accuracy\ \{p\ :\ ℕ\}\ (hp\ :\ 0\ <\ p)\ (f\ :\ State\ p\ →\ ℝ)}
\leanline{4}{(m\ ε₀\ ε\ :\ ℝ)\ (L\ :\ ℝ≥0)\ (hm\ :\ 0\ <\ m)\ (hε\ :\ 0\ <\ ε)}
\leanline{4}{(hf\ :\ Differentiable\ ℝ\ f)\ (hsc\ :\ StrongConvexOn\ univ\ m\ f)}
\leanline{4}{(hLip\ :\ LipschitzWith\ L\ (gradient\ f))}
\leanline{4}{(μ\ :\ Measure\ (State\ p))\ [IsProbabilityMeasure\ μ]}
\leanline{4}{(hfinite\ :\ wasserstein2\ μ\ (gibbsMeasure\ f)\ ≠\ ⊤)}
\leanline{4}{(hinit\ :\ wasserstein2Real\ μ\ (gibbsMeasure\ f)\ \textasciicircum{}\ 2\ ≤\ ε₀\ \textasciicircum{}\ 2\ *\ p\ /\ m)}
\leanline{4}{(k\ :\ ℕ)\ (hk\ :\ variableIterations\ ((L\ :\ ℝ)\ /\ m)\ (averageSmoothness\ f\ /\ m\ +\ 1)\ ε₀\ ε\ ≤\ k)\ :}
\leanline{4}{let\ κ\ :=\ (L\ :\ ℝ)\ /\ m}
\leanline{4}{let\ a\ :=\ averageSmoothness\ f\ /\ m\ +\ 1}
\leanline{4}{let\ k₀\ :=\ warmup\ κ\ a\ ε₀}
\leanline{4}{wasserstein2Real\ (lmcLaws\ f\ (lmcVariableSchedule\ (L\ :\ ℝ)\ m\ k₀)\ μ\ k)}
\leanline{6}{(gibbsMeasure\ f)\ \textasciicircum{}\ 2\ ≤\ ε\ \textasciicircum{}\ 2\ *\ p\ /\ m}
\end{leancode}
\leanlabel{How the Lean proof follows the argument}
\texttt{warmup\_\allowbreak{}decay} checks the initial exponential allocation, including the logarithmic ceiling. \texttt{harmonic\_\allowbreak{}step} proves the single rational induction step after warmup. \texttt{lmc\_\allowbreak{}variable\_\allowbreak{}step} verifies admissibility of every actual step in \texttt{hs}, normalizes the initial cost in \texttt{hi}, and derives the normalized recurrence in \texttt{hrec}. \texttt{variable\_\allowbreak{}step\_\allowbreak{}normalized\_\allowbreak{}all} handles zero initial error explicitly before applying the scalar warmup-and-tail estimate.
\end{leanbox}


Earlier results used here: \cite[Section~3.17.3]{LangevinLean2026ULASC}; \cite[Section~3.17.4]{LangevinLean2026ULASC}; \cite[Section~3.26.3]{LangevinLean2026ULASC}.
\subsection{\texorpdfstring{Remove separate initial-moment assumptions from the global statements}{Remove separate initial-moment assumptions from the global statements}}
\label{report-step-AverageLMC-Theorem-lean-5}


The exponential contraction factor is strictly positive. If the initial extended Wasserstein distance is infinite, the extended global constant-step inequality has infinite right side and follows immediately. If it is finite, the proved real-valued estimate and finite-output transport convert back to the extended nonnegative inequality. For the variable-step proposition, its stated finite initial squared-distance bound already forces finite initial distance; taking real values gives exactly the initial-error inequality used above. Applying the actual-law variable envelope and accuracy theorems therefore requires no independent finite-moment premise.

\begin{leanbox}{Lean: \texttt{lmc\_\allowbreak{}constant\_\allowbreak{}step\_\allowbreak{}extended}, \texttt{finite\_\allowbreak{}initial\_\allowbreak{}error\_\allowbreak{}of\_\allowbreak{}extended}, \texttt{lmc\_\allowbreak{}variable\_\allowbreak{}step\_\allowbreak{}extended\_\allowbreak{}initial} and 1 more}
\leansource{\texttt{lmc\_\allowbreak{}constant\_\allowbreak{}step\_\allowbreak{}extended} \quad AverageLMC/\allowbreak{}Theorem.lean, lines 256–275}
\begin{leancode}
\leanline{0}{theorem\ lmc\_constant\_step\_extended\ \{p\ :\ ℕ\}\ (hp\ :\ 0\ <\ p)\ (f\ :\ State\ p\ →\ ℝ)}
\leanline{4}{(m\ h\ :\ ℝ)\ (L\ :\ ℝ≥0)\ (hm\ :\ 0\ <\ m)\ (hh\ :\ 0\ ≤\ h)}
\leanline{4}{(hf\ :\ Differentiable\ ℝ\ f)\ (hsc\ :\ StrongConvexOn\ univ\ m\ f)}
\leanline{4}{(hLip\ :\ LipschitzWith\ L\ (gradient\ f))\ (hstep\ :\ h\ *\ (L\ :\ ℝ)\ ≤\ 1)}
\leanline{4}{(μ\ :\ Measure\ (State\ p))\ [IsProbabilityMeasure\ μ]\ (k\ :\ ℕ)\ :}
\leanline{4}{wasserstein2\ ((ulaLawStep\ f\ h)\textasciicircum{}[k]\ μ)\ (gibbsMeasure\ f)\ \textasciicircum{}\ 2\ ≤}
\leanline{6}{ENNReal.ofReal\ (Real.exp\ (-2\ *\ m\ *\ (k\ :\ ℝ)\ *\ h))\ *}
\leanline{8}{wasserstein2\ μ\ (gibbsMeasure\ f)\ \textasciicircum{}\ 2\ +}
\leanline{6}{ENNReal.ofReal\ ((averageSmoothness\ f\ +\ m)\ *\ h\ *\ p\ /\ (2\ *\ m))}
\end{leancode}
\leansource{\texttt{finite\_\allowbreak{}initial\_\allowbreak{}error\_\allowbreak{}of\_\allowbreak{}extended} \quad AverageLMC/\allowbreak{}Theorem.lean, lines 279–289}
\begin{leancode}
\leanline{0}{theorem\ finite\_initial\_error\_of\_extended\ \{p\ :\ ℕ\}\ (μ\ π\ :\ Measure\ (State\ p))}
\leanline{4}{(m\ ε₀\ :\ ℝ)\ (hm\ :\ 0\ <\ m)}
\leanline{4}{(hinit\ :\ wasserstein2\ μ\ π\ \textasciicircum{}\ 2\ ≤\ ENNReal.ofReal\ (ε₀\ \textasciicircum{}\ 2\ *\ p\ /\ m))\ :}
\leanline{4}{wasserstein2\ μ\ π\ ≠\ ⊤\ ∧\ wasserstein2Real\ μ\ π\ \textasciicircum{}\ 2\ ≤\ ε₀\ \textasciicircum{}\ 2\ *\ p\ /\ m}
\end{leancode}
\leansource{\texttt{lmc\_\allowbreak{}variable\_\allowbreak{}step\_\allowbreak{}extended\_\allowbreak{}initial} \quad AverageLMC/\allowbreak{}Theorem.lean, lines 293–306}
\begin{leancode}
\leanline{0}{theorem\ lmc\_variable\_step\_extended\_initial\ \{p\ :\ ℕ\}\ (hp\ :\ 0\ <\ p)\ (f\ :\ State\ p\ →\ ℝ)}
\leanline{4}{(m\ ε₀\ :\ ℝ)\ (L\ :\ ℝ≥0)\ (hm\ :\ 0\ <\ m)}
\leanline{4}{(hf\ :\ Differentiable\ ℝ\ f)\ (hsc\ :\ StrongConvexOn\ univ\ m\ f)}
\leanline{4}{(hLip\ :\ LipschitzWith\ L\ (gradient\ f))}
\leanline{4}{(μ\ :\ Measure\ (State\ p))\ [IsProbabilityMeasure\ μ]}
\leanline{4}{(hinit\ :\ wasserstein2\ μ\ (gibbsMeasure\ f)\ \textasciicircum{}\ 2\ ≤\ ENNReal.ofReal\ (ε₀\ \textasciicircum{}\ 2\ *\ p\ /\ m))}
\leanline{4}{(k\ :\ ℕ)\ (hk\ :\ warmup\ ((L\ :\ ℝ)\ /\ m)\ (averageSmoothness\ f\ /\ m\ +\ 1)\ ε₀\ ≤\ k)\ :}
\leanline{4}{let\ κ\ :=\ (L\ :\ ℝ)\ /\ m}
\leanline{4}{let\ a\ :=\ averageSmoothness\ f\ /\ m\ +\ 1}
\leanline{4}{let\ k₀\ :=\ warmup\ κ\ a\ ε₀}
\leanline{4}{wasserstein2Real\ (lmcLaws\ f\ (lmcVariableSchedule\ (L\ :\ ℝ)\ m\ k₀)\ μ\ k)}
\leanline{6}{(gibbsMeasure\ f)\ \textasciicircum{}\ 2\ ≤\ a\ /\ (κ\ +\ ((k\ -\ k₀\ :\ ℕ)\ :\ ℝ))\ *\ p\ /\ m}
\end{leancode}
\leansource{\texttt{lmc\_\allowbreak{}variable\_\allowbreak{}step\_\allowbreak{}accuracy\_\allowbreak{}extended\_\allowbreak{}initial} \quad AverageLMC/\allowbreak{}Theorem.lean, lines 310–323}
\begin{leancode}
\leanline{0}{theorem\ lmc\_variable\_step\_accuracy\_extended\_initial\ \{p\ :\ ℕ\}\ (hp\ :\ 0\ <\ p)\ (f\ :\ State\ p\ →\ ℝ)}
\leanline{4}{(m\ ε₀\ ε\ :\ ℝ)\ (L\ :\ ℝ≥0)\ (hm\ :\ 0\ <\ m)\ (hε\ :\ 0\ <\ ε)}
\leanline{4}{(hf\ :\ Differentiable\ ℝ\ f)\ (hsc\ :\ StrongConvexOn\ univ\ m\ f)}
\leanline{4}{(hLip\ :\ LipschitzWith\ L\ (gradient\ f))}
\leanline{4}{(μ\ :\ Measure\ (State\ p))\ [IsProbabilityMeasure\ μ]}
\leanline{4}{(hinit\ :\ wasserstein2\ μ\ (gibbsMeasure\ f)\ \textasciicircum{}\ 2\ ≤\ ENNReal.ofReal\ (ε₀\ \textasciicircum{}\ 2\ *\ p\ /\ m))}
\leanline{4}{(k\ :\ ℕ)\ (hk\ :\ variableIterations\ ((L\ :\ ℝ)\ /\ m)\ (averageSmoothness\ f\ /\ m\ +\ 1)\ ε₀\ ε\ ≤\ k)\ :}
\leanline{4}{let\ κ\ :=\ (L\ :\ ℝ)\ /\ m}
\leanline{4}{let\ a\ :=\ averageSmoothness\ f\ /\ m\ +\ 1}
\leanline{4}{let\ k₀\ :=\ warmup\ κ\ a\ ε₀}
\leanline{4}{wasserstein2Real\ (lmcLaws\ f\ (lmcVariableSchedule\ (L\ :\ ℝ)\ m\ k₀)\ μ\ k)}
\leanline{6}{(gibbsMeasure\ f)\ \textasciicircum{}\ 2\ ≤\ ε\ \textasciicircum{}\ 2\ *\ p\ /\ m}
\end{leancode}
\leanlabel{How the Lean proof follows the argument}
\texttt{lmc\_\allowbreak{}constant\_\allowbreak{}step\_\allowbreak{}extended} splits finite and infinite initial extended distance, converting the finite real theorem back with \texttt{ENNReal.\allowbreak{}of\allowbreak{}Real}. \texttt{finite\_\allowbreak{}initial\_\allowbreak{}error\_\allowbreak{}of\_\allowbreak{}extended} derives finiteness from the stated finite squared-distance bound. \texttt{lmc\_\allowbreak{}variable\_\allowbreak{}step\_\allowbreak{}extended\_\allowbreak{}initial} and \texttt{lmc\_\allowbreak{}variable\_\allowbreak{}step\_\allowbreak{}accuracy\_\allowbreak{}extended\_\allowbreak{}initial} then use that derived fact as the initial finiteness premise of their real-valued counterparts. Thus the extended statements do not interpret an infinite distance through \texttt{ENNReal.\allowbreak{}to\allowbreak{}Real}.
\end{leanbox}


Earlier results used here: \cite[Section~3.17.3]{LangevinLean2026ULASC}; \cite[Section~3.17.4]{LangevinLean2026ULASC}; \cite[Section~3.26.3]{LangevinLean2026ULASC}; \cite[Section~3.36.7]{LangevinLean2026ULASC}.
\section{\texorpdfstring{The scalar quadratic recurrence and its invariant floor}{The scalar quadratic recurrence and its invariant floor}}
\label{proof-module-AverageLMC-LaplacianRecurrence-lean}


This module isolates arithmetic used after the actual LMC one-step estimate is proved. Its recurrence premise is an intermediate numerical interface, not a replacement for the stochastic theorem.



The invariant-floor calculation retains the quadratic cross term:


\[\begin{gathered}
e_{k+1}^2\le q^2e_k^2+2qCe_k+B^2,\qquad
C\le(1-q)F,\\
2qCF+B^2\le(1-q^2)F^2
\quad\Longrightarrow\quad e_n\le q^ne_0+F.
\end{gathered}\]


With the coefficients of the Laplacian theorem (Theorem 2, source label th:3), the certified floor is


\[e_n\le(1-mh)^ne_0+\frac{(2MA+\delta)hS}{m},
\qquad A=\sqrt{M_{\mathrm{av}}},\quad S=\sqrt p.\]\subsection{\texorpdfstring{Keep the quadratic cross term}{Keep the quadratic cross term}}
\label{report-step-AverageLMC-LaplacianRecurrence-lean-1}


Induct on the iteration number with the comparison \(q^k\) \(e_{0}+F\). The inequalities \(C\le (1-q)F\) and \(2qCF+B^{2}\le (1-q^{2})F^{2}\) make this comparison invariant. Nonnegativity allows comparison of squares without replacing the local error by a triangle bound.

\begin{leanbox}{Lean: \texttt{quadratic\_\allowbreak{}recurrence\_\allowbreak{}floor}}
\leansource{\texttt{quadratic\_\allowbreak{}recurrence\_\allowbreak{}floor} \quad AverageLMC/\allowbreak{}LaplacianRecurrence.lean, lines 19–42}
\begin{leancode}
\leanline{0}{theorem\ quadratic\_recurrence\_floor}
\leanline{4}{(e\ :\ ℕ\ →\ ℝ)\ (q\ C\ B\ F\ :\ ℝ)}
\leanline{4}{(he\ :\ ∀\ k,\ 0\ ≤\ e\ k)\ (hq\ :\ 0\ ≤\ q)\ (\_hq1\ :\ q\ ≤\ 1)}
\leanline{4}{(hC\ :\ 0\ ≤\ C)\ (hF\ :\ 0\ ≤\ F)}
\leanline{4}{(hCF\ :\ C\ ≤\ (1\ -\ q)\ *\ F)}
\leanline{4}{(hfloor\ :\ 2\ *\ q\ *\ C\ *\ F\ +\ B\ \textasciicircum{}\ 2\ ≤\ (1\ -\ q\ \textasciicircum{}\ 2)\ *\ F\ \textasciicircum{}\ 2)}
\leanline{4}{(hstep\ :\ ∀\ k,\ (e\ (k\ +\ 1))\ \textasciicircum{}\ 2\ ≤}
\leanline{6}{q\ \textasciicircum{}\ 2\ *\ (e\ k)\ \textasciicircum{}\ 2\ +\ 2\ *\ q\ *\ C\ *\ e\ k\ +\ B\ \textasciicircum{}\ 2)\ (n\ :\ ℕ)\ :}
\leanline{4}{e\ n\ ≤\ q\ \textasciicircum{}\ n\ *\ e\ 0\ +\ F}
\end{leancode}
\leanlabel{How the Lean proof follows the argument}
\texttt{quadratic\_\allowbreak{}recurrence\_\allowbreak{}floor} proceeds by \texttt{induction n with}. The base cases establish the initial claims and the inductive branches use the induction hypotheses. Named proof applications include \texttt{hbound}; the adjacent proof shows their hypotheses and order of use.
\end{leanbox}
\subsection{\texorpdfstring{Split the floor into two contributions}{Split the floor into two contributions}}
\label{report-step-AverageLMC-LaplacianRecurrence-lean-2}


Choose a drift floor \(a\) and a variance floor \(b\). From \(C\le (1-q)a\) and \(B^{2}\le (1-q^{2})b^{2}\), expand the square of \(a+b\) and verify the invariant-floor hypotheses. This retains the order \(h\) accumulation of the quadratic defect.

\begin{leanbox}{Lean: \texttt{quadratic\_\allowbreak{}recurrence\_\allowbreak{}split\_\allowbreak{}floor}}
\leansource{\texttt{quadratic\_\allowbreak{}recurrence\_\allowbreak{}split\_\allowbreak{}floor} \quad AverageLMC/\allowbreak{}LaplacianRecurrence.lean, lines 46–66}
\begin{leancode}
\leanline{0}{theorem\ quadratic\_recurrence\_split\_floor}
\leanline{4}{(e\ :\ ℕ\ →\ ℝ)\ (q\ C\ B\ a\ b\ :\ ℝ)}
\leanline{4}{(he\ :\ ∀\ k,\ 0\ ≤\ e\ k)\ (hq\ :\ 0\ ≤\ q)\ (hq1\ :\ q\ ≤\ 1)}
\leanline{4}{(hC\ :\ 0\ ≤\ C)\ (ha\ :\ 0\ ≤\ a)\ (hb\ :\ 0\ ≤\ b)}
\leanline{4}{(hCa\ :\ C\ ≤\ (1\ -\ q)\ *\ a)}
\leanline{4}{(hBb\ :\ B\ \textasciicircum{}\ 2\ ≤\ (1\ -\ q\ \textasciicircum{}\ 2)\ *\ b\ \textasciicircum{}\ 2)}
\leanline{4}{(hstep\ :\ ∀\ k,\ (e\ (k\ +\ 1))\ \textasciicircum{}\ 2\ ≤}
\leanline{6}{q\ \textasciicircum{}\ 2\ *\ (e\ k)\ \textasciicircum{}\ 2\ +\ 2\ *\ q\ *\ C\ *\ e\ k\ +\ B\ \textasciicircum{}\ 2)\ (n\ :\ ℕ)\ :}
\leanline{4}{e\ n\ ≤\ q\ \textasciicircum{}\ n\ *\ e\ 0\ +\ a\ +\ b}
\end{leancode}
\leanlabel{How the Lean proof follows the argument}
\texttt{quadratic\_\allowbreak{}recurrence\_\allowbreak{}split\_\allowbreak{}floor} gives the formal counterpart of this step. The main intermediate facts are \texttt{hq0}, \texttt{h\allowbreak{}CF}, \texttt{hcross}, \texttt{hrest}, \texttt{hfloor} in that order. Named proof applications include \texttt{quadratic\_\allowbreak{}recurrence\_\allowbreak{}floor}; the adjacent proof shows their hypotheses and order of use.
\end{leanbox}
\subsection{\texorpdfstring{Recover the manuscript coefficient}{Recover the manuscript coefficient}}
\label{report-step-AverageLMC-LaplacianRecurrence-lean-3}


Set \(q=1-mh\), \(C=h^{2}(MA+\delta )S/2\) and use the manuscript-style local-defect coefficient. The conditions \(m\le M\) and \(hM\le1\) control all square roots. Explicit floors add to \((2MA+\delta)hS/m\), giving the claimed iteration bound with \(A=\sqrt{M_{\mathrm{av}}}\) and \(S=\sqrt p\).

\begin{leanbox}{Lean: \texttt{laplacian\_\allowbreak{}quadratic\_\allowbreak{}recurrence\_\allowbreak{}bound}}
\leansource{\texttt{laplacian\_\allowbreak{}quadratic\_\allowbreak{}recurrence\_\allowbreak{}bound} \quad AverageLMC/\allowbreak{}LaplacianRecurrence.lean, lines 71–146}
\begin{leancode}
\leanline{0}{theorem\ laplacian\_quadratic\_recurrence\_bound}
\leanline{4}{(e\ :\ ℕ\ →\ ℝ)\ (m\ M\ h\ A\ δ\ S\ :\ ℝ)}
\leanline{4}{(he\ :\ ∀\ k,\ 0\ ≤\ e\ k)\ (hm\ :\ 0\ <\ m)\ (hmM\ :\ m\ ≤\ M)}
\leanline{4}{(hh\ :\ 0\ ≤\ h)\ (hhM\ :\ h\ *\ M\ ≤\ 1)}
\leanline{4}{(hA\ :\ 0\ ≤\ A)\ (hδ\ :\ 0\ ≤\ δ)\ (hS\ :\ 0\ ≤\ S)}
\leanline{4}{(hstep\ :\ ∀\ k,\ (e\ (k\ +\ 1))\ \textasciicircum{}\ 2\ ≤}
\leanline{6}{(1\ -\ m\ *\ h)\ \textasciicircum{}\ 2\ *\ (e\ k)\ \textasciicircum{}\ 2\ +}
\leanline{8}{(1\ -\ m\ *\ h)\ *\ h\ \textasciicircum{}\ 2\ *\ (M\ *\ A\ +\ δ)\ *\ S\ *\ e\ k\ +}
\leanline{10}{(h\ \textasciicircum{}\ 2\ /\ 2\ *\ δ\ *\ S\ +}
\leanline{12}{(3\ /\ 2\ :\ ℝ)\ *\ h\ *\ Real.sqrt\ h\ *\ Real.sqrt\ M\ *\ A\ *\ S)\ \textasciicircum{}\ 2)\ (n\ :\ ℕ)\ :}
\leanline{4}{e\ n\ ≤\ (1\ -\ m\ *\ h)\ \textasciicircum{}\ n\ *\ e\ 0\ +\ (2\ *\ M\ *\ A\ +\ δ)\ *\ h\ *\ S\ /\ m}
\end{leancode}
\leanlabel{How the Lean proof follows the argument}
\texttt{laplacian\_\allowbreak{}quadratic\_\allowbreak{}recurrence\_\allowbreak{}bound} gives the formal counterpart of this step. The main intermediate facts are \texttt{h\allowbreak{}M}, \texttt{hmh}, \texttt{q}, \texttt{C}, \texttt{B}, \texttt{a} in that order. Named proof applications include \texttt{h\allowbreak{}Bb0.\allowbreak{}trans}, \texttt{mul\_\allowbreak{}le\_\allowbreak{}mul\_\allowbreak{}of\_\allowbreak{}nonneg\_\allowbreak{}right}, \texttt{hbound}; the adjacent proof shows their hypotheses and order of use.
\end{leanbox}
\section{\texorpdfstring{A sufficient total-defect bound with the exact final coefficient}{A sufficient total-defect bound with the exact final coefficient}}
\label{proof-module-AverageLMC-LaplacianLocalError-lean}


The improved rate needs generator cancellation in the cross term, while an ordinary Brownian increment estimate suffices for the total squared local defect. This avoids adding stochastic-integral assumptions to the paper.



For the actual integral solution and its Brownian increment bounds, \(hM\le1\) gives


\[\left\|\int_0^h\bigl(g(Y_t)-g(Y_0)\bigr)\,dt\right\|
\le\frac32Mh\sqrt h\,\sqrt d.\]


The final scalar argument uses \(q=1-mh\), \(m\le A^2\), and the precise implication


\[\begin{gathered}
e_{k+1}^2\le q^2e_k^2+qh^2(MA+\delta)S e_k
+\left(\frac32Mh\sqrt h\,S\right)^2\\
\Longrightarrow\quad e_n\le q^ne_0+\frac{(2MA+\delta)hS}{m}.
\end{gathered}\]

Earlier results used here: \cite[Section~3.11]{LangevinLean2026ULASC}.
\subsection{\texorpdfstring{Control the total local defect}{Control the total local defect}}
\label{report-step-AverageLMC-LaplacianLocalError-lean-1}


Apply the already proved integral-equation defect bound. A score norm at most \(\sqrt{Md}\), Brownian increment norm at most \(\sqrt{2d}\sqrt t\), and \(hM\le1\) give a total defect at most \((3/2)Mh\sqrt h\sqrt d\). The numerical estimate follows from \(\sqrt 2\le 3/2\).

\begin{leanbox}{Lean: \texttt{localDefect\_\allowbreak{}three\_\allowbreak{}halves\_\allowbreak{}bound}}
\leansource{\texttt{localDefect\_\allowbreak{}three\_\allowbreak{}halves\_\allowbreak{}bound} \quad AverageLMC/\allowbreak{}LaplacianLocalError.lean, lines 25–52}
\begin{leancode}
\leanline{0}{theorem\ localDefect\_three\_halves\_bound}
\leanline{4}{(g\ :\ H\ →\ H)\ (Y\ B\ :\ ℝ\ →\ H)\ (M\ :\ ℝ≥0)\ (h\ d\ :\ ℝ)}
\leanline{4}{(hh\ :\ 0\ ≤\ h)\ (hd\ :\ 0\ ≤\ d)\ (hsmall\ :\ h\ *\ (M\ :\ ℝ)\ ≤\ 1)}
\leanline{4}{(hLip\ :\ LipschitzWith\ M\ g)}
\leanline{4}{(hG\ :\ ∀\ t\ ∈\ Set.Icc\ 0\ h,\ ‖g\ (Y\ t)‖\ ≤\ Real.sqrt\ ((M\ :\ ℝ)\ *\ d))}
\leanline{4}{(hB\ :\ ∀\ t\ ∈\ Set.Icc\ 0\ h,\ ‖B\ t‖\ ≤\ Real.sqrt\ (2\ *\ d)\ *\ Real.sqrt\ t)}
\leanline{4}{(heq\ :\ ∀\ t\ ∈\ Set.Icc\ 0\ h,}
\leanline{6}{Y\ t\ -\ Y\ 0\ =\ -(∫\ s\ in\ 0..t,\ g\ (Y\ s))\ +\ B\ t)\ :}
\leanline{4}{‖ULA.localDefect\ g\ Y\ h‖\ ≤}
\leanline{6}{(3\ /\ 2\ :\ ℝ)\ *\ (M\ :\ ℝ)\ *\ h\ *\ Real.sqrt\ h\ *\ Real.sqrt\ d}
\end{leancode}
\leanlabel{How the Lean proof follows the argument}
\texttt{local\allowbreak{}Defect\_\allowbreak{}three\_\allowbreak{}halves\_\allowbreak{}bound} gives the formal counterpart of this step. The main intermediate facts are \texttt{hroot}, \texttt{htwo}, \texttt{hnum}, \texttt{hb}, \texttt{hscale}, \texttt{hscale2} in that order. The final step is the algebraic normalization or ordered arithmetic displayed in the proof.
\end{leanbox}


Earlier results used here: \cite[Section~3.11.1]{LangevinLean2026ULASC}; \cite[Section~3.11.7]{LangevinLean2026ULASC}.
\subsection{\texorpdfstring{Preserve the exact order \(h\) floor}{Preserve the exact order h floor}}
\label{report-step-AverageLMC-LaplacianLocalError-lean-2}


For a recurrence with cross coefficient \((1-mh)h^{2}(MA+\delta )S\) and the preceding total defect, use \(m\le A^{2}\) to absorb its variance floor. The explicit floors \(a=h(MA+\delta )S/(2m)\) and \(b=3hMAS/(2m)\) sum to at most the manuscript's \((2MA+\delta)hS/m\). The final law-level theorem later proves the recurrence rather than assuming it.

\begin{leanbox}{Lean: \texttt{laplacian\_\allowbreak{}recurrence\_\allowbreak{}bound\_\allowbreak{}of\_\allowbreak{}global\_\allowbreak{}defect}}
\leansource{\texttt{laplacian\_\allowbreak{}recurrence\_\allowbreak{}bound\_\allowbreak{}of\_\allowbreak{}global\_\allowbreak{}defect} \quad AverageLMC/\allowbreak{}LaplacianLocalError.lean, lines 57–116}
\begin{leancode}
\leanline{0}{theorem\ laplacian\_recurrence\_bound\_of\_global\_defect}
\leanline{4}{(e\ :\ ℕ\ →\ ℝ)\ (m\ M\ h\ A\ δ\ S\ :\ ℝ)}
\leanline{4}{(he\ :\ ∀\ k,\ 0\ ≤\ e\ k)\ (hm\ :\ 0\ <\ m)\ (hmM\ :\ m\ ≤\ M)}
\leanline{4}{(hh\ :\ 0\ ≤\ h)\ (hhM\ :\ h\ *\ M\ ≤\ 1)}
\leanline{4}{(hA\ :\ 0\ ≤\ A)\ (hmA\ :\ m\ ≤\ A\ \textasciicircum{}\ 2)\ (hδ\ :\ 0\ ≤\ δ)\ (hS\ :\ 0\ ≤\ S)}
\leanline{4}{(hstep\ :\ ∀\ k,\ (e\ (k\ +\ 1))\ \textasciicircum{}\ 2\ ≤}
\leanline{6}{(1\ -\ m\ *\ h)\ \textasciicircum{}\ 2\ *\ (e\ k)\ \textasciicircum{}\ 2\ +}
\leanline{8}{(1\ -\ m\ *\ h)\ *\ h\ \textasciicircum{}\ 2\ *\ (M\ *\ A\ +\ δ)\ *\ S\ *\ e\ k\ +}
\leanline{10}{((3\ /\ 2\ :\ ℝ)\ *\ M\ *\ h\ *\ Real.sqrt\ h\ *\ S)\ \textasciicircum{}\ 2)\ (n\ :\ ℕ)\ :}
\leanline{4}{e\ n\ ≤\ (1\ -\ m\ *\ h)\ \textasciicircum{}\ n\ *\ e\ 0\ +\ (2\ *\ M\ *\ A\ +\ δ)\ *\ h\ *\ S\ /\ m}
\end{leancode}
\leanlabel{How the Lean proof follows the argument}
\texttt{laplacian\_\allowbreak{}recurrence\_\allowbreak{}bound\_\allowbreak{}of\_\allowbreak{}global\_\allowbreak{}defect} gives the formal counterpart of this step. The main intermediate facts are \texttt{h\allowbreak{}M}, \texttt{hmh}, \texttt{q}, \texttt{C}, \texttt{B}, \texttt{a} in that order. Named proof applications include \texttt{mul\_\allowbreak{}le\_\allowbreak{}mul\_\allowbreak{}of\_\allowbreak{}nonneg\_\allowbreak{}right}, \texttt{div\_\allowbreak{}le\_\allowbreak{}div\_\allowbreak{}of\_\allowbreak{}nonneg\_\allowbreak{}right}, \texttt{hbound.\allowbreak{}trans}; the adjacent proof shows their hypotheses and order of use.
\end{leanbox}
\section{\texorpdfstring{Exact average-smoothness Gibbs score and Laplacian moments}{Exact average-smoothness Gibbs score and Laplacian moments}}
\label{proof-module-AverageLMC-ScoreMoment-lean}


The first-order part proves \(\mathbb E_\pi\|\nabla f\|^2\le pM_{\rm av}\) without assuming any Hessian or integration-by-parts identity. The second part proves Lemma 7 for \(C^{2}\) convex potentials with finite Gibbs normalization: the true score energy equals the \(L^{1}\) Laplacian integral, bounded by the exact Laplacian supremum and then the exact average coordinate smoothness. Strong convexity supplies normalization automatically in the main application.



Earlier results used here: \cite[Section~3.27]{LangevinLean2026ULASC}; \cite[Section~3.46]{LangevinLean2026ULASC}.
\subsection{\texorpdfstring{Obtain the exact coordinate score MGF by translation}{Obtain the exact coordinate score MGF by translation}}
\label{report-step-AverageLMC-ScoreMoment-lean-1}


The coordinate quadratic model implies

\[e^{-f(x)}e^{-a\partial_j f(x)}\le e^{M_j a^2/2}e^{-f(x+ae_j)}.\]

Integrate the nonnegative functions against Lebesgue measure. Translation invariance leaves the partition function unchanged, and finite positive normalization cancels it. Replacing a by minus \(t\) produces the Gaussian MGF bound for the actual \(j\)th score coordinate, with proxy \(M_j\). The same finite exponential integral bound supplies exponential integrability, so this is an actual sub-Gaussian statement rather than a formal integral comparison.

\begin{leanbox}{Lean: \texttt{coordinate\_\allowbreak{}score\_\allowbreak{}lmgf\_\allowbreak{}le}, \texttt{coordinate\_\allowbreak{}score\_\allowbreak{}subgaussian}}
\leansource{\texttt{coordinate\_\allowbreak{}score\_\allowbreak{}lmgf\_\allowbreak{}le} \quad AverageLMC/\allowbreak{}ScoreMoment.lean, lines 24–64}
\begin{leancode}
\leanline{0}{theorem\ coordinate\_score\_lmgf\_le\ \{p\ :\ ℕ\}}
\leanline{4}{(f\ :\ State\ p\ →\ ℝ)\ (M\ :\ ℝ≥0)\ (hdiff\ :\ Differentiable\ ℝ\ f)}
\leanline{4}{(hLip\ :\ LipschitzWith\ M\ (gradient\ f))\ (hZ\ :\ ULA.gibbsPartition\ f\ ≠\ ⊤)}
\leanline{4}{(j\ :\ Fin\ p)\ (a\ :\ ℝ)\ :}
\leanline{4}{(∫⁻\ x,\ ENNReal.ofReal\ (Real.exp\ (-a\ *\ (gradient\ f\ x)\ j))}
\leanline{6}{∂ULA.gibbsMeasure\ f)\ ≤}
\leanline{6}{ENNReal.ofReal\ (Real.exp\ (coordinateSmoothness\ f\ j\ /\ 2\ *\ a\ \textasciicircum{}\ 2))}
\end{leancode}
\leansource{\texttt{coordinate\_\allowbreak{}score\_\allowbreak{}subgaussian} \quad AverageLMC/\allowbreak{}ScoreMoment.lean, lines 66–91}
\begin{leancode}
\leanline{0}{theorem\ coordinate\_score\_subgaussian\ \{p\ :\ ℕ\}}
\leanline{4}{(f\ :\ State\ p\ →\ ℝ)\ (M\ :\ ℝ≥0)\ (hdiff\ :\ Differentiable\ ℝ\ f)}
\leanline{4}{(hLip\ :\ LipschitzWith\ M\ (gradient\ f))\ (hZ\ :\ ULA.gibbsPartition\ f\ ≠\ ⊤)}
\leanline{4}{(j\ :\ Fin\ p)\ (hj\ :\ 0\ ≤\ coordinateSmoothness\ f\ j)\ :}
\leanline{4}{HasSubgaussianMGF\ (fun\ x\ ↦\ (gradient\ f\ x)\ j)}
\leanline{6}{⟨coordinateSmoothness\ f\ j,\ hj⟩\ (ULA.gibbsMeasure\ f)}
\end{leancode}
\leanlabel{How the Lean proof follows the argument}
\texttt{coordinate\_\allowbreak{}score\_\allowbreak{}lmgf\_\allowbreak{}le} gives the formal counterpart of this step. The main intermediate facts are \texttt{hw}, \texttt{hc}, \texttt{hg}, \texttt{hpoint}, \texttt{ht} in that order. Named proof applications include \texttt{ENNReal.\allowbreak{}of\allowbreak{}Real\_\allowbreak{}le\_\allowbreak{}of\allowbreak{}Real}, \texttt{Real.\allowbreak{}exp\_\allowbreak{}le\_\allowbreak{}exp.\allowbreak{}mpr}; the adjacent proof shows their hypotheses and order of use.




\texttt{coordinate\_\allowbreak{}score\_\allowbreak{}subgaussian} gives the formal counterpart of this step. The main intermediate facts are \texttt{he}, \texttt{hi}, \texttt{h} in that order. Its remaining equalities are discharged by the simplifications displayed in the proof.
\end{leanbox}


Earlier results used here: \cite[Section~3.26.1]{LangevinLean2026ULASC}; \cite[Section~3.26.2]{LangevinLean2026ULASC}; \cite[Section~3.26.3]{LangevinLean2026ULASC}; \cite[Section~3.26.9]{LangevinLean2026ULASC}.
\subsection{\texorpdfstring{Sum sharp second moments without second differentiability}{Sum sharp second moments without second differentiability}}
\label{report-step-AverageLMC-ScoreMoment-lean-2}


Strong convexity implies \(m\) at most every coordinate Taylor supremum, so each \(M_j\) is nonnegative. Apply the established sub-Gaussian second-moment theorem to the coordinate score: its uncentered second moment is at most \(M_j\), and square integrability is derived at the same time. The squared Euclidean norm is the sum of the squared coordinates. Integral linearity and the finite sum therefore give

\[\mathbb E_\pi\|\nabla f\|^2\le\sum_j M_j=pM_{\rm av}.\]

Only differentiability, a Lipschitz gradient, and positive strong convexity are used in this sharper score bound.

\begin{leanbox}{Lean: \texttt{integral\_\allowbreak{}norm\_\allowbreak{}sq\_\allowbreak{}gradient\_\allowbreak{}gibbs\_\allowbreak{}le}}
\leansource{\texttt{integral\_\allowbreak{}norm\_\allowbreak{}sq\_\allowbreak{}gradient\_\allowbreak{}gibbs\_\allowbreak{}le} \quad AverageLMC/\allowbreak{}ScoreMoment.lean, lines 95–127}
\begin{leancode}
\leanline{0}{theorem\ integral\_norm\_sq\_gradient\_gibbs\_le\ \{p\ :\ ℕ\}\ (hp\ :\ 0\ <\ p)}
\leanline{4}{(f\ :\ State\ p\ →\ ℝ)\ (m\ :\ ℝ)\ (M\ :\ ℝ≥0)\ (hm\ :\ 0\ <\ m)}
\leanline{4}{(hdiff\ :\ Differentiable\ ℝ\ f)\ (hsc\ :\ StrongConvexOn\ Set.univ\ m\ f)}
\leanline{4}{(hLip\ :\ LipschitzWith\ M\ (gradient\ f))\ :}
\leanline{4}{(∫\ x,\ ‖gradient\ f\ x‖\ \textasciicircum{}\ 2\ ∂ULA.gibbsMeasure\ f)\ ≤}
\leanline{6}{(p\ :\ ℝ)\ *\ averageSmoothness\ f}
\end{leancode}
\leanlabel{How the Lean proof follows the argument}
\texttt{integral\_\allowbreak{}norm\_\allowbreak{}sq\_\allowbreak{}gradient\_\allowbreak{}gibbs\_\allowbreak{}le} gives the formal counterpart of this step. The main intermediate facts are \texttt{h\allowbreak{}Z}, \texttt{hj}, \texttt{hi}, \texttt{hb} in that order. Named proof applications include \texttt{ENNReal.\allowbreak{}inv\_\allowbreak{}mul\_\allowbreak{}cancel}, \texttt{integral\_\allowbreak{}congr\_\allowbreak{}ae}, \texttt{ae\_\allowbreak{}of\_\allowbreak{}all}; the adjacent proof shows their hypotheses and order of use.
\end{leanbox}


Earlier results used here: \cite[Section~3.26.10]{LangevinLean2026ULASC}; \cite[Section~3.26.3]{LangevinLean2026ULASC}; \cite[Section~3.26.9]{LangevinLean2026ULASC}; \cite[Section~3.27.5]{LangevinLean2026ULASC}.
\subsection{\texorpdfstring{Identify coordinate derivatives and control all weighted integrals}{Identify coordinate derivatives and control all weighted integrals}}
\label{report-step-AverageLMC-ScoreMoment-lean-3}


A normalized Gibbs integral is integrable exactly when its exponential-weighted Lebesgue integrand is integrable, since the partition function is finite and positive. The coordinate score equals the directional Frechet derivative in \(e_j\). Differentiating that expression identifies its derivative in \(e_j\) with the actual diagonal Hessian entry. \(C^{2}\) regularity makes this coordinate score \(C^{1}\). Projecting the global gradient Lipschitz inequality bounds the coordinate score with the same Lipschitz constant, and hence bounds its derivative. These facts justify the weighted integration by parts used next.

\begin{leanbox}{Lean: \texttt{integrable\_\allowbreak{}gibbs\_\allowbreak{}weighted}, \texttt{fderiv\_\allowbreak{}coordinate\_\allowbreak{}score}, \texttt{coordinate\_\allowbreak{}score\_\allowbreak{}contDiff} and 1 more}
\leansource{\texttt{integrable\_\allowbreak{}gibbs\_\allowbreak{}weighted} \quad AverageLMC/\allowbreak{}ScoreMoment.lean, lines 129–140}
\begin{leancode}
\leanline{0}{private\ theorem\ integrable\_gibbs\_weighted\ \{p\ :\ ℕ\}}
\leanline{4}{(f\ :\ State\ p\ →\ ℝ)\ (hf\ :\ Differentiable\ ℝ\ f)\ (hZ\ :\ ULA.gibbsPartition\ f\ ≠\ ⊤)}
\leanline{4}{(g\ :\ State\ p\ →\ ℝ)\ :\ Integrable\ g\ (ULA.gibbsMeasure\ f)\ ↔}
\leanline{6}{Integrable\ (fun\ x\ ↦\ Real.exp\ (-f\ x)\ *\ g\ x)\ volume}
\end{leancode}
\leansource{\texttt{fderiv\_\allowbreak{}coordinate\_\allowbreak{}score} \quad AverageLMC/\allowbreak{}ScoreMoment.lean, lines 144–153}
\begin{leancode}
\leanline{0}{theorem\ fderiv\_coordinate\_score\ \{p\ :\ ℕ\}\ (f\ :\ State\ p\ →\ ℝ)}
\leanline{4}{(hf\ :\ ContDiff\ ℝ\ 2\ f)\ (x\ :\ State\ p)\ (j\ :\ Fin\ p)\ :}
\leanline{4}{fderiv\ ℝ\ (fun\ y\ ↦\ (gradient\ f\ y)\ j)\ x\ (coordinateVector\ j)\ =\ hessian\ f\ x\ j\ j}
\end{leancode}
\leansource{\texttt{coordinate\_\allowbreak{}score\_\allowbreak{}contDiff} \quad AverageLMC/\allowbreak{}ScoreMoment.lean, lines 155–163}
\begin{leancode}
\leanline{0}{private\ theorem\ coordinate\_score\_contDiff\ \{p\ :\ ℕ\}\ (f\ :\ State\ p\ →\ ℝ)}
\leanline{4}{(hf\ :\ ContDiff\ ℝ\ 2\ f)\ (j\ :\ Fin\ p)\ :}
\leanline{4}{ContDiff\ ℝ\ 1\ (fun\ y\ ↦\ (gradient\ f\ y)\ j)}
\end{leancode}
\leansource{\texttt{coordinate\_\allowbreak{}score\_\allowbreak{}lipschitz} \quad AverageLMC/\allowbreak{}ScoreMoment.lean, lines 165–172}
\begin{leancode}
\leanline{0}{private\ theorem\ coordinate\_score\_lipschitz\ \{p\ :\ ℕ\}\ (f\ :\ State\ p\ →\ ℝ)\ (M\ :\ ℝ≥0)}
\leanline{4}{(hLip\ :\ LipschitzWith\ M\ (gradient\ f))\ (j\ :\ Fin\ p)\ :}
\leanline{4}{LipschitzWith\ M\ (fun\ x\ ↦\ (gradient\ f\ x)\ j)}
\end{leancode}
\leanlabel{How the Lean proof follows the argument}
\texttt{integrable\_\allowbreak{}gibbs\_\allowbreak{}weighted} gives the formal counterpart of this step. The main intermediate facts are \texttt{hw} in that order. Its remaining equalities are discharged by the simplifications displayed in the proof.




\texttt{fderiv\_\allowbreak{}coordinate\_\allowbreak{}score} gives the formal counterpart of this step. The main intermediate facts are \texttt{heq} in that order. Its remaining equalities are discharged by the simplifications displayed in the proof.




\texttt{coordinate\_\allowbreak{}score\_\allowbreak{}cont\allowbreak{}Diff} gives the formal counterpart of this step. The main intermediate facts are \texttt{heq} in that order. Its remaining equalities are discharged by the simplifications displayed in the proof.




\texttt{coordinate\_\allowbreak{}score\_\allowbreak{}lipschitz} gives the formal counterpart of this step. Named proof applications include \texttt{Lipschitz\allowbreak{}With.\allowbreak{}of\_\allowbreak{}dist\_\allowbreak{}le\_\allowbreak{}mul}; the adjacent proof shows their hypotheses and order of use.
\end{leanbox}


Earlier results used here: \cite[Section~3.26.1]{LangevinLean2026ULASC}; \cite[Section~3.26.2]{LangevinLean2026ULASC}; \cite[Section~3.26.3]{LangevinLean2026ULASC}; \cite[Section~3.26.9]{LangevinLean2026ULASC}.
\subsection{\texorpdfstring{Prove the score-Hessian identity by global integration by parts}{Prove the score-Hessian identity by global integration by parts}}
\label{report-step-AverageLMC-ScoreMoment-lean-4}


Set \(\psi\) equal to the \(j\)th score coordinate. Its first and second moments are finite by the first-order Gibbs score theorem. Its directional derivative is bounded by \(M\), so that derivative is Gibbs integrable as well. The three unnormalized functions needed for global integration by parts are therefore integrable: the derivative of \(\psi\) times \(e^{-f}\), \(\psi\) times the derivative of \(e^{-f}\), and \(\psi e^{-f}\). The chain rule gives \(D_je^{-f}=-e^{-f}\psi\). Integration by parts cancels its minus sign and proves

\[\int(\partial_j f)^2\,d\pi=\int\partial_{jj}f\,d\pi.\]

No decay or boundary term is supplied as a hypothesis. Summing the finite coordinate identities and using \(\operatorname{tr}(D^2f)=\Delta f\) proves the full score-Laplacian identity. The partition-based statement works whenever normalization is finite; strong convexity supplies normalization automatically.

\begin{leanbox}{Lean: \texttt{integral\_\allowbreak{}coordinate\_\allowbreak{}score\_\allowbreak{}sq\_\allowbreak{}eq\_\allowbreak{}hessian}, \texttt{integral\_\allowbreak{}norm\_\allowbreak{}sq\_\allowbreak{}gradient\_\allowbreak{}eq\_\allowbreak{}laplacian\_\allowbreak{}of\_\allowbreak{}partition}, \texttt{integral\_\allowbreak{}norm\_\allowbreak{}sq\_\allowbreak{}gradient\_\allowbreak{}eq\_\allowbreak{}laplacian}}
\leansource{\texttt{integral\_\allowbreak{}coordinate\_\allowbreak{}score\_\allowbreak{}sq\_\allowbreak{}eq\_\allowbreak{}hessian} \quad AverageLMC/\allowbreak{}ScoreMoment.lean, lines 176–240}
\begin{leancode}
\leanline{0}{theorem\ integral\_coordinate\_score\_sq\_eq\_hessian\ \{p\ :\ ℕ\}}
\leanline{4}{(f\ :\ State\ p\ →\ ℝ)\ (M\ :\ ℝ≥0)\ (hf\ :\ ContDiff\ ℝ\ 2\ f)}
\leanline{4}{(hLip\ :\ LipschitzWith\ M\ (gradient\ f))\ (hZ\ :\ ULA.gibbsPartition\ f\ ≠\ ⊤)}
\leanline{4}{(j\ :\ Fin\ p)\ :}
\leanline{4}{(∫\ x,\ (gradient\ f\ x)\ j\ \textasciicircum{}\ 2\ ∂ULA.gibbsMeasure\ f)\ =}
\leanline{6}{∫\ x,\ hessian\ f\ x\ j\ j\ ∂ULA.gibbsMeasure\ f}
\end{leancode}
\leansource{\texttt{integral\_\allowbreak{}norm\_\allowbreak{}sq\_\allowbreak{}gradient\_\allowbreak{}eq\_\allowbreak{}laplacian\_\allowbreak{}of\_\allowbreak{}partition} \quad AverageLMC/\allowbreak{}ScoreMoment.lean, lines 244–282}
\begin{leancode}
\leanline{0}{theorem\ integral\_norm\_sq\_gradient\_eq\_laplacian\_of\_partition\ \{p\ :\ ℕ\}}
\leanline{4}{(f\ :\ State\ p\ →\ ℝ)\ (M\ :\ ℝ≥0)\ (hf\ :\ ContDiff\ ℝ\ 2\ f)}
\leanline{4}{(hLip\ :\ LipschitzWith\ M\ (gradient\ f))\ (hZ\ :\ ULA.gibbsPartition\ f\ ≠\ ⊤)\ :}
\leanline{4}{(∫\ x,\ ‖gradient\ f\ x‖\ \textasciicircum{}\ 2\ ∂ULA.gibbsMeasure\ f)\ =}
\leanline{6}{∫\ x,\ laplacian\ f\ x\ ∂ULA.gibbsMeasure\ f}
\end{leancode}
\leansource{\texttt{integral\_\allowbreak{}norm\_\allowbreak{}sq\_\allowbreak{}gradient\_\allowbreak{}eq\_\allowbreak{}laplacian} \quad AverageLMC/\allowbreak{}ScoreMoment.lean, lines 285–292}
\begin{leancode}
\leanline{0}{theorem\ integral\_norm\_sq\_gradient\_eq\_laplacian\ \{p\ :\ ℕ\}}
\leanline{4}{(f\ :\ State\ p\ →\ ℝ)\ (m\ :\ ℝ)\ (M\ :\ ℝ≥0)\ (hm\ :\ 0\ <\ m)}
\leanline{4}{(hf\ :\ ContDiff\ ℝ\ 2\ f)\ (hsc\ :\ StrongConvexOn\ Set.univ\ m\ f)}
\leanline{4}{(hLip\ :\ LipschitzWith\ M\ (gradient\ f))\ :}
\leanline{4}{(∫\ x,\ ‖gradient\ f\ x‖\ \textasciicircum{}\ 2\ ∂ULA.gibbsMeasure\ f)\ =}
\leanline{6}{∫\ x,\ laplacian\ f\ x\ ∂ULA.gibbsMeasure\ f}
\end{leancode}
\leanlabel{How the Lean proof follows the argument}
\texttt{integral\_\allowbreak{}coordinate\_\allowbreak{}score\_\allowbreak{}sq\_\allowbreak{}eq\_\allowbreak{}hessian} gives the formal counterpart of this step. The main intermediate facts are \texttt{hdiff}, \texttt{\ensuremath{\psi}}, \texttt{h\ensuremath{\psi}}, \texttt{h\ensuremath{\psi}\allowbreak{}Lp}, \texttt{hi\ensuremath{\psi}}, \texttt{hi\ensuremath{\psi}2} in that order. Named proof applications include \texttt{ae\_\allowbreak{}of\_\allowbreak{}all}, \texttt{hibp}, \texttt{congr\allowbreak{}Arg}; the adjacent proof shows their hypotheses and order of use.




\texttt{integral\_\allowbreak{}norm\_\allowbreak{}sq\_\allowbreak{}gradient\_\allowbreak{}eq\_\allowbreak{}laplacian\_\allowbreak{}of\_\allowbreak{}partition} gives the formal counterpart of this step. The main intermediate facts are \texttt{hdiff}, \texttt{hscore}, \texttt{h\allowbreak{}H} in that order. Named proof applications include \texttt{integral\_\allowbreak{}congr\_\allowbreak{}ae}, \texttt{ae\_\allowbreak{}of\_\allowbreak{}all}, \texttt{integral\_\allowbreak{}coordinate\_\allowbreak{}score\_\allowbreak{}sq\_\allowbreak{}eq\_\allowbreak{}hessian}; the adjacent proof shows their hypotheses and order of use.




\texttt{integral\_\allowbreak{}norm\_\allowbreak{}sq\_\allowbreak{}gradient\_\allowbreak{}eq\_\allowbreak{}laplacian} gives the formal counterpart of this step.
\end{leanbox}


Earlier results used here: \cite[Section~3.26.10]{LangevinLean2026ULASC}; \cite[Section~3.26.2]{LangevinLean2026ULASC}; \cite[Section~3.26.3]{LangevinLean2026ULASC}; \cite[Section~3.26.9]{LangevinLean2026ULASC}; \cite[Section~3.27.7]{LangevinLean2026ULASC}; \cite[Section~3.46.3]{LangevinLean2026ULASC}.
\subsection{\texorpdfstring{Differentiate the coordinate quadratic model to bound each Hessian diagonal}{Differentiate the coordinate quadratic model to bound each Hessian diagonal}}
\label{report-step-AverageLMC-ScoreMoment-lean-5}


Apply the coordinate upper model at \(x\) with displacement \(ae_j\) and at \(x+ae_j\) with displacement minus \(ae_j\). Addition gives \(a\) times the difference of the coordinate scores at most \(M_j\) \(a^2\). For \(a>0\), divide by \(a^2\) and take the right derivative at zero. Its actual derivative is the diagonal Hessian entry, so \(\partial_{jj}f(x)\) is at most \(M_j\). Taking the trace proves the pointwise Laplacian bound \(p\) \(M_{\rm av}\). Conversely strong gradient monotonicity along the same coordinate line bounds its difference quotient below by \(m\). Taking the same derivative gives \(m\) at most every diagonal Hessian entry. In the convex case \(m=0\), their sum is nonnegative.

\begin{leanbox}{Lean: \texttt{hessian\_\allowbreak{}diagonal\_\allowbreak{}le\_\allowbreak{}coordinateSmoothness}, \texttt{laplacian\_\allowbreak{}le\_\allowbreak{}averageSmoothness}, \texttt{strongConvex\_\allowbreak{}le\_\allowbreak{}hessian\_\allowbreak{}diagonal} and 1 more}
\leansource{\texttt{hessian\_\allowbreak{}diagonal\_\allowbreak{}le\_\allowbreak{}coordinateSmoothness} \quad AverageLMC/\allowbreak{}ScoreMoment.lean, lines 296–324}
\begin{leancode}
\leanline{0}{theorem\ hessian\_diagonal\_le\_coordinateSmoothness\ \{p\ :\ ℕ\}}
\leanline{4}{(f\ :\ State\ p\ →\ ℝ)\ (M\ :\ ℝ≥0)\ (hf\ :\ ContDiff\ ℝ\ 2\ f)}
\leanline{4}{(hLip\ :\ LipschitzWith\ M\ (gradient\ f))\ (x\ :\ State\ p)\ (j\ :\ Fin\ p)\ :}
\leanline{4}{hessian\ f\ x\ j\ j\ ≤\ coordinateSmoothness\ f\ j}
\end{leancode}
\leansource{\texttt{laplacian\_\allowbreak{}le\_\allowbreak{}averageSmoothness} \quad AverageLMC/\allowbreak{}ScoreMoment.lean, lines 326–338}
\begin{leancode}
\leanline{0}{theorem\ laplacian\_le\_averageSmoothness\ \{p\ :\ ℕ\}\ (hp\ :\ 0\ <\ p)}
\leanline{4}{(f\ :\ State\ p\ →\ ℝ)\ (M\ :\ ℝ≥0)\ (hf\ :\ ContDiff\ ℝ\ 2\ f)}
\leanline{4}{(hLip\ :\ LipschitzWith\ M\ (gradient\ f))\ (x\ :\ State\ p)\ :}
\leanline{4}{laplacian\ f\ x\ ≤\ (p\ :\ ℝ)\ *\ averageSmoothness\ f}
\end{leancode}
\leansource{\texttt{strongConvex\_\allowbreak{}le\_\allowbreak{}hessian\_\allowbreak{}diagonal} \quad AverageLMC/\allowbreak{}ScoreMoment.lean, lines 341–369}
\begin{leancode}
\leanline{0}{theorem\ strongConvex\_le\_hessian\_diagonal\ \{p\ :\ ℕ\}}
\leanline{4}{(f\ :\ State\ p\ →\ ℝ)\ (m\ :\ ℝ)\ (hf\ :\ ContDiff\ ℝ\ 2\ f)}
\leanline{4}{(hsc\ :\ StrongConvexOn\ Set.univ\ m\ f)\ (x\ :\ State\ p)\ (j\ :\ Fin\ p)\ :}
\leanline{4}{m\ ≤\ hessian\ f\ x\ j\ j}
\end{leancode}
\leansource{\texttt{laplacian\_\allowbreak{}nonneg\_\allowbreak{}of\_\allowbreak{}convex} \quad AverageLMC/\allowbreak{}ScoreMoment.lean, lines 372–379}
\begin{leancode}
\leanline{0}{theorem\ laplacian\_nonneg\_of\_convex\ \{p\ :\ ℕ\}}
\leanline{4}{(f\ :\ State\ p\ →\ ℝ)\ (hf\ :\ ContDiff\ ℝ\ 2\ f)\ (hconv\ :\ ConvexOn\ ℝ\ Set.univ\ f)}
\leanline{4}{(x\ :\ State\ p)\ :\ 0\ ≤\ laplacian\ f\ x}
\end{leancode}
\leanlabel{How the Lean proof follows the argument}
\texttt{hessian\_\allowbreak{}diagonal\_\allowbreak{}le\_\allowbreak{}coordinate\allowbreak{}Smoothness} gives the formal counterpart of this step. The main intermediate facts are \texttt{hdiff}, \texttt{h\ensuremath{\psi}}, \texttt{hline}, \texttt{hd}, \texttt{hb}, \texttt{hlim} in that order. Named proof applications include \texttt{le\_\allowbreak{}of\_\allowbreak{}tendsto\_\allowbreak{}of\_\allowbreak{}tendsto}, \texttt{hb}; the adjacent proof shows their hypotheses and order of use.




\texttt{laplacian\_\allowbreak{}le\_\allowbreak{}average\allowbreak{}Smoothness} gives the formal counterpart of this step. The main intermediate facts are \texttt{h} in that order. The final step is the algebraic normalization or ordered arithmetic displayed in the proof.




\texttt{strong\allowbreak{}Convex\_\allowbreak{}le\_\allowbreak{}hessian\_\allowbreak{}diagonal} gives the formal counterpart of this step. The main intermediate facts are \texttt{hdiff}, \texttt{h\ensuremath{\psi}}, \texttt{hline}, \texttt{hd}, \texttt{hb}, \texttt{hlim} in that order. Named proof applications include \texttt{le\_\allowbreak{}of\_\allowbreak{}tendsto\_\allowbreak{}of\_\allowbreak{}tendsto}, \texttt{hb}; the adjacent proof shows their hypotheses and order of use.




\texttt{laplacian\_\allowbreak{}nonneg\_\allowbreak{}of\_\allowbreak{}convex} gives the formal counterpart of this step. Named proof applications include \texttt{Finset.\allowbreak{}sum\_\allowbreak{}nonneg}, \texttt{strong\allowbreak{}Convex\_\allowbreak{}le\_\allowbreak{}hessian\_\allowbreak{}diagonal}; the adjacent proof shows their hypotheses and order of use.
\end{leanbox}


Earlier results used here: \cite[Section~3.2.3]{LangevinLean2026ULASC}.
\subsection{\texorpdfstring{Use the genuine Laplacian supremum and complete Lemma 7}{Use the genuine Laplacian supremum and complete Lemma 7}}
\label{report-step-AverageLMC-ScoreMoment-lean-6}


Define \(M_{1,\Delta}\) as the supremum of the range of the true Laplacian, divided by \(p\). Its range is nonempty, and the pointwise bound \(p\) \(M_{\rm av}\) makes it bounded above. The defining supremum therefore bounds every Laplacian value and is itself at most \(p\) \(M_{\rm av}\). Dividing by positive \(p\) gives \(M_{1,\Delta}\) at most \(M_{\rm av}\). \(C^{2}\) regularity also makes the Laplacian continuous, as a finite sum of continuous score derivatives. Its nonnegativity and finite pointwise bound prove actual integrability. Under the normalized Gibbs law, the score identity and these pointwise comparisons give exactly

\[\mathbb E_\pi\|\nabla f\|^2=\mathbb E_\pi|\Delta f|\le pM_{1,\Delta}\le pM_{\rm av}.\]\begin{leanbox}{Lean: \texttt{averageLaplacianSmoothness}, \texttt{averageLaplacianSmoothness\_\allowbreak{}le\_\allowbreak{}averageSmoothness}, \texttt{laplacian\_\allowbreak{}le\_\allowbreak{}averageLaplacianSmoothness} and 2 more}
\leansource{\texttt{averageLaplacianSmoothness} \quad AverageLMC/\allowbreak{}ScoreMoment.lean, lines 383–384}
\begin{leancode}
\leanline{0}{def\ averageLaplacianSmoothness\ \{p\ :\ ℕ\}\ (f\ :\ State\ p\ →\ ℝ)\ :\ ℝ\ :=}
\leanline{2}{sSup\ (Set.range\ (laplacian\ f))\ /\ (p\ :\ ℝ)}
\end{leancode}
\leansource{\texttt{averageLaplacianSmoothness\_\allowbreak{}le\_\allowbreak{}averageSmoothness} \quad AverageLMC/\allowbreak{}ScoreMoment.lean, lines 386–394}
\begin{leancode}
\leanline{0}{theorem\ averageLaplacianSmoothness\_le\_averageSmoothness\ \{p\ :\ ℕ\}\ (hp\ :\ 0\ <\ p)}
\leanline{4}{(f\ :\ State\ p\ →\ ℝ)\ (M\ :\ ℝ≥0)\ (hf\ :\ ContDiff\ ℝ\ 2\ f)}
\leanline{4}{(hLip\ :\ LipschitzWith\ M\ (gradient\ f))\ :}
\leanline{4}{averageLaplacianSmoothness\ f\ ≤\ averageSmoothness\ f}
\end{leancode}
\leansource{\texttt{laplacian\_\allowbreak{}le\_\allowbreak{}averageLaplacianSmoothness} \quad AverageLMC/\allowbreak{}ScoreMoment.lean, lines 397–408}
\begin{leancode}
\leanline{0}{theorem\ laplacian\_le\_averageLaplacianSmoothness\ \{p\ :\ ℕ\}\ (hp\ :\ 0\ <\ p)}
\leanline{4}{(f\ :\ State\ p\ →\ ℝ)\ (M\ :\ ℝ≥0)\ (hf\ :\ ContDiff\ ℝ\ 2\ f)}
\leanline{4}{(hLip\ :\ LipschitzWith\ M\ (gradient\ f))\ (x\ :\ State\ p)\ :}
\leanline{4}{laplacian\ f\ x\ ≤\ (p\ :\ ℝ)\ *\ averageLaplacianSmoothness\ f}
\end{leancode}
\leansource{\texttt{continuous\_\allowbreak{}laplacian} \quad AverageLMC/\allowbreak{}ScoreMoment.lean, lines 410–419}
\begin{leancode}
\leanline{0}{theorem\ continuous\_laplacian\ \{p\ :\ ℕ\}\ (f\ :\ State\ p\ →\ ℝ)\ (hf\ :\ ContDiff\ ℝ\ 2\ f)\ :}
\leanline{4}{Continuous\ (laplacian\ f)}
\end{leancode}
\leansource{\texttt{score\_\allowbreak{}laplacian\_\allowbreak{}identity\_\allowbreak{}and\_\allowbreak{}bounds} \quad AverageLMC/\allowbreak{}ScoreMoment.lean, lines 424–454}
\begin{leancode}
\leanline{0}{theorem\ score\_laplacian\_identity\_and\_bounds\ \{p\ :\ ℕ\}\ (hp\ :\ 0\ <\ p)}
\leanline{4}{(f\ :\ State\ p\ →\ ℝ)\ (M\ :\ ℝ≥0)\ (hf\ :\ ContDiff\ ℝ\ 2\ f)}
\leanline{4}{(hconv\ :\ ConvexOn\ ℝ\ Set.univ\ f)\ (hLip\ :\ LipschitzWith\ M\ (gradient\ f))}
\leanline{4}{(hZ\ :\ ULA.gibbsPartition\ f\ ≠\ ⊤)\ :}
\leanline{4}{(∫\ x,\ ‖gradient\ f\ x‖\ \textasciicircum{}\ 2\ ∂ULA.gibbsMeasure\ f)\ =}
\leanline{8}{(∫\ x,\ |laplacian\ f\ x|\ ∂ULA.gibbsMeasure\ f)\ ∧}
\leanline{6}{(∫\ x,\ |laplacian\ f\ x|\ ∂ULA.gibbsMeasure\ f)\ ≤}
\leanline{8}{(p\ :\ ℝ)\ *\ averageLaplacianSmoothness\ f\ ∧}
\leanline{6}{(p\ :\ ℝ)\ *\ averageLaplacianSmoothness\ f\ ≤\ (p\ :\ ℝ)\ *\ averageSmoothness\ f}
\end{leancode}
\leanlabel{How the Lean proof follows the argument}
\texttt{average\allowbreak{}Laplacian\allowbreak{}Smoothness} names the mathematical object defined above. Its defining expression appears in the adjacent statement; subsequent identities may unfold this definition directly.




\texttt{average\allowbreak{}Laplacian\allowbreak{}Smoothness\_\allowbreak{}le\_\allowbreak{}average\allowbreak{}Smoothness} gives the formal counterpart of this step. Named proof applications include \texttt{cs\allowbreak{}Sup\_\allowbreak{}le}, \texttt{laplacian\_\allowbreak{}le\_\allowbreak{}average\allowbreak{}Smoothness}; the adjacent proof shows their hypotheses and order of use.




\texttt{laplacian\_\allowbreak{}le\_\allowbreak{}average\allowbreak{}Laplacian\allowbreak{}Smoothness} gives the formal counterpart of this step. The main intermediate facts are \texttt{hb}, \texttt{h} in that order. Named proof applications include \texttt{laplacian\_\allowbreak{}le\_\allowbreak{}average\allowbreak{}Smoothness}; the adjacent proof shows their hypotheses and order of use.




\texttt{continuous\_\allowbreak{}laplacian} gives the formal counterpart of this step. The main intermediate facts are \texttt{heq} in that order. Named proof applications include \texttt{continuous\_\allowbreak{}finset\allowbreak{}Sum}; the adjacent proof shows their hypotheses and order of use.




\texttt{score\_\allowbreak{}laplacian\_\allowbreak{}identity\_\allowbreak{}and\_\allowbreak{}bounds} gives the formal counterpart of this step. The main intermediate facts are \texttt{hdiff}, \texttt{hn}, \texttt{habs}, \texttt{hi} in that order. Named proof applications include \texttt{ENNReal.\allowbreak{}inv\_\allowbreak{}mul\_\allowbreak{}cancel}, \texttt{ae\_\allowbreak{}of\_\allowbreak{}all}, \texttt{laplacian\_\allowbreak{}le\_\allowbreak{}average\allowbreak{}Smoothness}; the adjacent proof shows their hypotheses and order of use.
\end{leanbox}


Earlier results used here: \cite[Section~3.26.2]{LangevinLean2026ULASC}; \cite[Section~3.26.3]{LangevinLean2026ULASC}; \cite[Section~3.26.9]{LangevinLean2026ULASC}.
\section{\texorpdfstring{Commute the actual gradient and Laplacian under \(C^{3}\) regularity}{Commute the actual gradient and Laplacian under C³ regularity}}
\label{proof-module-AverageLMC-LaplacianGradient-lean}


The third-order analysis needs the identity \(\Delta \nabla f=\nabla \Delta f\) for the true derivatives. The proof uses Schwarz symmetry twice, and requires three continuous derivatives rather than analyticity or a Lipschitz Hessian.



Three continuous derivatives suffice for the directional Schwarz identities and the vector Laplacian-gradient commutation used by the diffusion generator:


\[D_uD_uD_vf=D_vD_uD_uf\qquad\bigl(f\in C^3(\mathbb R^p)\bigr).\]
\[\Delta(\nabla f)=\nabla(\Delta f).\]\subsection{\texorpdfstring{Fixed directional derivatives and Schwarz symmetry}{Fixed directional derivatives and Schwarz symmetry}}
\label{report-step-AverageLMC-LaplacianGradient-lean-1}


A directional derivative is the actual Fréchet derivative applied to a fixed vector. A continuously differentiable derivative map loses one order of regularity, and evaluation at a fixed vector preserves the remaining regularity. Differentiating that evaluation identifies a second directional derivative with the actual second Fréchet derivative applied to the two directions. Schwarz symmetry for \(C^{2}\) functions then commutes the two directional derivatives as functions of the base point. These statements are proved for arbitrary real normed spaces, before choosing Euclidean coordinates.

\begin{leanbox}{Lean: \texttt{directionalDerivative}, \texttt{contDiff\_\allowbreak{}directionalDerivative}, \texttt{directionalDerivative\_\allowbreak{}two} and 1 more}
\leansource{\texttt{directionalDerivative} \quad AverageLMC/\allowbreak{}LaplacianGradient.lean, lines 13–14}
\begin{leancode}
\leanline{0}{def\ directionalDerivative\ \{E\ :\ Type*\}\ [NormedAddCommGroup\ E]\ [NormedSpace\ ℝ\ E]}
\leanline{4}{(f\ :\ E\ →\ ℝ)\ (v\ :\ E)\ :\ E\ →\ ℝ\ :=\ fun\ x\ ↦\ fderiv\ ℝ\ f\ x\ v}
\end{leancode}
\leansource{\texttt{contDiff\_\allowbreak{}directionalDerivative} \quad AverageLMC/\allowbreak{}LaplacianGradient.lean, lines 18–21}
\begin{leancode}
\leanline{0}{theorem\ contDiff\_directionalDerivative\ \{n\ :\ ℕ\}\ (f\ :\ E\ →\ ℝ)}
\leanline{4}{(hf\ :\ ContDiff\ ℝ\ (n\ +\ 1)\ f)\ (v\ :\ E)\ :}
\leanline{4}{ContDiff\ ℝ\ n\ (directionalDerivative\ f\ v)}
\end{leancode}
\leansource{\texttt{directionalDerivative\_\allowbreak{}two} \quad AverageLMC/\allowbreak{}LaplacianGradient.lean, lines 23–31}
\begin{leancode}
\leanline{0}{theorem\ directionalDerivative\_two\ (f\ :\ E\ →\ ℝ)\ (hf\ :\ ContDiff\ ℝ\ 2\ f)}
\leanline{4}{(u\ v\ x\ :\ E)\ :}
\leanline{4}{directionalDerivative\ (directionalDerivative\ f\ v)\ u\ x\ =}
\leanline{6}{fderiv\ ℝ\ (fderiv\ ℝ\ f)\ x\ u\ v}
\end{leancode}
\leansource{\texttt{directionalDerivative\_\allowbreak{}commute} \quad AverageLMC/\allowbreak{}LaplacianGradient.lean, lines 34–39}
\begin{leancode}
\leanline{0}{theorem\ directionalDerivative\_commute\ (f\ :\ E\ →\ ℝ)\ (hf\ :\ ContDiff\ ℝ\ 2\ f)\ (u\ v\ :\ E)\ :}
\leanline{4}{directionalDerivative\ (directionalDerivative\ f\ v)\ u\ =}
\leanline{6}{directionalDerivative\ (directionalDerivative\ f\ u)\ v}
\end{leancode}
\leanlabel{How the Lean proof follows the argument}
\texttt{directional\allowbreak{}Derivative} names the mathematical object defined above. Its defining expression appears in the adjacent statement; subsequent identities may unfold this definition directly.




\texttt{cont\allowbreak{}Diff\_\allowbreak{}directional\allowbreak{}Derivative} gives the formal counterpart of this step.




\texttt{directional\allowbreak{}Derivative\_\allowbreak{}two} gives the formal counterpart of this step. Its remaining equalities are discharged by the simplifications displayed in the proof.




\texttt{directional\allowbreak{}Derivative\_\allowbreak{}commute} gives the formal counterpart of this step.
\end{leanbox}
\subsection{\texorpdfstring{Commute the needed three derivatives without analyticity}{Commute the needed three derivatives without analyticity}}
\label{report-step-AverageLMC-LaplacianGradient-lean-2}


For a \(C^{3}\) potential, first commute the two inner directional derivatives using \(C^{2}\) symmetry of \(f\). Then commute the outer pair using \(C^{2}\) symmetry of one first directional derivative. This gives \(D_{u}D_{u}D_{v}f=D_{v}D_{u}D_{u}f\). The argument establishes exactly the third-order permutation needed for the Laplacian and avoids any appeal to an analytic-only iterated-derivative permutation theorem.

\begin{leanbox}{Lean: \texttt{directionalDerivative\_\allowbreak{}three}}
\leansource{\texttt{directionalDerivative\_\allowbreak{}three} \quad AverageLMC/\allowbreak{}LaplacianGradient.lean, lines 42–48}
\begin{leancode}
\leanline{0}{theorem\ directionalDerivative\_three\ (f\ :\ E\ →\ ℝ)\ (hf\ :\ ContDiff\ ℝ\ 3\ f)}
\leanline{4}{(u\ v\ :\ E)\ :}
\leanline{4}{directionalDerivative\ (directionalDerivative\ (directionalDerivative\ f\ v)\ u)\ u\ =}
\leanline{6}{directionalDerivative\ (directionalDerivative\ (directionalDerivative\ f\ u)\ u)\ v}
\end{leancode}
\leanlabel{How the Lean proof follows the argument}
\texttt{directional\allowbreak{}Derivative\_\allowbreak{}three} gives the formal counterpart of this step. Named proof applications include \texttt{directional\allowbreak{}Derivative\_\allowbreak{}commute}; the adjacent proof shows their hypotheses and order of use.
\end{leanbox}
\subsection{\texorpdfstring{Expand the canonical trace and differentiate the finite sum}{Expand the canonical trace and differentiate the finite sum}}
\label{report-step-AverageLMC-LaplacianGradient-lean-3}


The canonical orthonormal-basis formula for the Laplacian expresses it as the sum of second directional derivatives along the coordinate vectors. The \(j\)th coordinate of the actual gradient equals the directional derivative along \(e_j\) by the gradient–Fréchet-derivative identity. A finite sum of differentiable scalar functions can be differentiated term by term. Finally, the gradient is \(C^{2}\) because it is the continuous Riesz inverse applied to the \(C^{2}\) derivative map of a \(C^{3}\) potential.

\begin{leanbox}{Lean: \texttt{laplacian\_\allowbreak{}eq\_\allowbreak{}directionalDerivative}, \texttt{gradient\_\allowbreak{}coordinate\_\allowbreak{}directionalDerivative}, \texttt{directionalDerivative\_\allowbreak{}sum} and 1 more}
\leansource{\texttt{laplacian\_\allowbreak{}eq\_\allowbreak{}directionalDerivative} \quad AverageLMC/\allowbreak{}LaplacianGradient.lean, lines 52–62}
\begin{leancode}
\leanline{0}{theorem\ laplacian\_eq\_directionalDerivative\ \{p\ :\ ℕ\}\ (f\ :\ State\ p\ →\ ℝ)}
\leanline{4}{(hf\ :\ ContDiff\ ℝ\ 2\ f)\ :}
\leanline{4}{laplacian\ f\ =\ fun\ x\ ↦\ ∑\ i\ :\ Fin\ p,}
\leanline{6}{directionalDerivative\ (directionalDerivative\ f\ (coordinateVector\ i))\ (coordinateVector\ i)\ x}
\end{leancode}
\leansource{\texttt{gradient\_\allowbreak{}coordinate\_\allowbreak{}directionalDerivative} \quad AverageLMC/\allowbreak{}LaplacianGradient.lean, lines 65–68}
\begin{leancode}
\leanline{0}{theorem\ gradient\_coordinate\_directionalDerivative\ \{p\ :\ ℕ\}\ (f\ :\ State\ p\ →\ ℝ)\ (j\ :\ Fin\ p)\ :}
\leanline{4}{(fun\ x\ ↦\ (gradient\ f\ x)\ j)\ =\ directionalDerivative\ f\ (coordinateVector\ j)}
\end{leancode}
\leansource{\texttt{directionalDerivative\_\allowbreak{}sum} \quad AverageLMC/\allowbreak{}LaplacianGradient.lean, lines 72–78}
\begin{leancode}
\leanline{0}{theorem\ directionalDerivative\_sum\ \{ι\ :\ Type*\}\ [Fintype\ ι]}
\leanline{4}{(g\ :\ ι\ →\ E\ →\ ℝ)\ (hg\ :\ ∀\ i,\ Differentiable\ ℝ\ (g\ i))\ (v\ x\ :\ E)\ :}
\leanline{4}{directionalDerivative\ (fun\ y\ ↦\ ∑\ i,\ g\ i\ y)\ v\ x\ =}
\leanline{6}{∑\ i,\ directionalDerivative\ (g\ i)\ v\ x}
\end{leancode}
\leansource{\texttt{contDiff\_\allowbreak{}gradient\_\allowbreak{}of\_\allowbreak{}three} \quad AverageLMC/\allowbreak{}LaplacianGradient.lean, lines 81–85}
\begin{leancode}
\leanline{0}{theorem\ contDiff\_gradient\_of\_three\ \{p\ :\ ℕ\}\ (f\ :\ State\ p\ →\ ℝ)\ (hf\ :\ ContDiff\ ℝ\ 3\ f)\ :}
\leanline{4}{ContDiff\ ℝ\ 2\ (gradient\ f)}
\end{leancode}
\leanlabel{How the Lean proof follows the argument}
\texttt{laplacian\_\allowbreak{}eq\_\allowbreak{}directional\allowbreak{}Derivative} gives the formal counterpart of this step. Named proof applications include \texttt{Finset.\allowbreak{}sum\_\allowbreak{}congr}; the adjacent proof shows their hypotheses and order of use.




\texttt{gradient\_\allowbreak{}coordinate\_\allowbreak{}directional\allowbreak{}Derivative} gives the formal counterpart of this step. Its remaining equalities are discharged by the simplifications displayed in the proof.




\texttt{directional\allowbreak{}Derivative\_\allowbreak{}sum} gives the formal counterpart of this step. Its remaining equalities are discharged by the simplifications displayed in the proof.




\texttt{cont\allowbreak{}Diff\_\allowbreak{}gradient\_\allowbreak{}of\_\allowbreak{}three} gives the formal counterpart of this step.
\end{leanbox}
\subsection{\texorpdfstring{Identify the vector Laplacian coordinate by coordinate}{Identify the vector Laplacian coordinate by coordinate}}
\label{report-step-AverageLMC-LaplacianGradient-lean-4}


Continuous linear projection onto coordinate \(j\) commutes with the vector Laplacian. Its scalar observable is \(D_jf\), whose Laplacian is the sum of \(D_iD_iD_j f\). The \(j\)th coordinate of \(\nabla \Delta f\) is \(D_j\) of the finite Laplacian sum, namely the sum of \(D_jD_iD_i f\). Apply the proved third-derivative commutation term by term, then use equality of all Euclidean coordinates to obtain the full vector identity.

\begin{leanbox}{Lean: \texttt{laplacian\_\allowbreak{}gradient\_\allowbreak{}eq\_\allowbreak{}gradient\_\allowbreak{}laplacian}}
\leansource{\texttt{laplacian\_\allowbreak{}gradient\_\allowbreak{}eq\_\allowbreak{}gradient\_\allowbreak{}laplacian} \quad AverageLMC/\allowbreak{}LaplacianGradient.lean, lines 89–112}
\begin{leancode}
\leanline{0}{theorem\ laplacian\_gradient\_eq\_gradient\_laplacian\ \{p\ :\ ℕ\}}
\leanline{4}{(f\ :\ State\ p\ →\ ℝ)\ (hf\ :\ ContDiff\ ℝ\ 3\ f)\ (x\ :\ State\ p)\ :}
\leanline{4}{Laplacian.laplacian\ (gradient\ f)\ x\ =\ gradient\ (laplacian\ f)\ x}
\end{leancode}
\leanlabel{How the Lean proof follows the argument}
\texttt{laplacian\_\allowbreak{}gradient\_\allowbreak{}eq\_\allowbreak{}gradient\_\allowbreak{}laplacian} gives the formal counterpart of this step. The main intermediate facts are \texttt{hgrad}, \texttt{hp} in that order. Named proof applications include \texttt{congr\allowbreak{}Fun}; the adjacent proof shows their hypotheses and order of use.
\end{leanbox}
\section{\texorpdfstring{The actual gradient generator and its derived \(L^{2}\) bound}{The actual gradient generator and its derived L² bound}}
\label{proof-module-AverageLMC-LaplacianGenerator-lean}


The vector generator in the third-order Langevin argument is represented by true derivatives of the same potential. Its continuity, finite second moments and quantitative \(L^{2}\) bound follow from the paper’s \(C^{3}\), gradient-Lipschitz and squared Laplacian-gradient assumptions.



The actual vector generator of the gradient and its sharp stationary \(L^2\) bound are


\[G_f(x)=\nabla\Delta f(x)-D(\nabla f)(x)[\nabla f(x)]=\Delta(\nabla f)(x)-D(\nabla f)(x)[\nabla f(x)].\]
\[Y\sim\pi\quad\Longrightarrow\quad\|G_f(Y)\|_{L^2}\le L\sqrt{M_{\rm av}p}+M_\Delta\sqrt p.\]\subsection{\texorpdfstring{Define and identify the true generator}{Define and identify the true generator}}
\label{report-step-AverageLMC-LaplacianGenerator-lean-1}


The vector field is \(\nabla \Delta f-D(\nabla f)\nabla f\). The directional-derivative trace formula makes \(\Delta f\) continuously differentiable for \(C^{3}\) \(f\). The Riesz representation of its derivative gives continuity of \(\nabla \Delta f\); evaluation of the continuous Hessian on the continuous gradient gives continuity of the second term. The proved vector commutation identity identifies this field exactly with \(\Delta (\nabla f)-D(\nabla f)\nabla f\), the Langevin generator applied to \(\nabla f\).

\begin{leanbox}{Lean: \texttt{gradientGenerator}, \texttt{contDiff\_\allowbreak{}laplacian\_\allowbreak{}one}, \texttt{continuous\_\allowbreak{}gradientGenerator} and 1 more}
\leansource{\texttt{gradientGenerator} \quad AverageLMC/\allowbreak{}LaplacianGenerator.lean, lines 10–11}
\begin{leancode}
\leanline{0}{def\ gradientGenerator\ \{p\ :\ ℕ\}\ (f\ :\ State\ p\ →\ ℝ)\ (x\ :\ State\ p)\ :\ State\ p\ :=}
\leanline{2}{gradient\ (laplacian\ f)\ x\ -\ fderiv\ ℝ\ (gradient\ f)\ x\ (gradient\ f\ x)}
\end{leancode}
\leansource{\texttt{contDiff\_\allowbreak{}laplacian\_\allowbreak{}one} \quad AverageLMC/\allowbreak{}LaplacianGenerator.lean, lines 14–18}
\begin{leancode}
\leanline{0}{theorem\ contDiff\_laplacian\_one\ \{p\ :\ ℕ\}\ (f\ :\ State\ p\ →\ ℝ)\ (hf\ :\ ContDiff\ ℝ\ 3\ f)\ :}
\leanline{4}{ContDiff\ ℝ\ 1\ (laplacian\ f)}
\end{leancode}
\leansource{\texttt{continuous\_\allowbreak{}gradientGenerator} \quad AverageLMC/\allowbreak{}LaplacianGenerator.lean, lines 21–29}
\begin{leancode}
\leanline{0}{theorem\ continuous\_gradientGenerator\ \{p\ :\ ℕ\}\ (f\ :\ State\ p\ →\ ℝ)\ (hf\ :\ ContDiff\ ℝ\ 3\ f)\ :}
\leanline{4}{Continuous\ (gradientGenerator\ f)}
\end{leancode}
\leansource{\texttt{gradientGenerator\_\allowbreak{}eq} \quad AverageLMC/\allowbreak{}LaplacianGenerator.lean, lines 33–37}
\begin{leancode}
\leanline{0}{theorem\ gradientGenerator\_eq\ \{p\ :\ ℕ\}\ (f\ :\ State\ p\ →\ ℝ)\ (hf\ :\ ContDiff\ ℝ\ 3\ f)\ (x\ :\ State\ p)\ :}
\leanline{4}{gradientGenerator\ f\ x\ =\ Laplacian.laplacian\ (gradient\ f)\ x\ -}
\leanline{6}{fderiv\ ℝ\ (gradient\ f)\ x\ (gradient\ f\ x)}
\end{leancode}
\leanlabel{How the Lean proof follows the argument}
\texttt{gradient\allowbreak{}Generator} names the mathematical object defined above. Its defining expression appears in the adjacent statement; subsequent identities may unfold this definition directly.




\texttt{cont\allowbreak{}Diff\_\allowbreak{}laplacian\_\allowbreak{}one} gives the formal counterpart of this step. Named proof applications include \texttt{Cont\allowbreak{}Diff.\allowbreak{}sum}; the adjacent proof shows their hypotheses and order of use.




\texttt{continuous\_\allowbreak{}gradient\allowbreak{}Generator} gives the formal counterpart of this step. The main intermediate facts are \texttt{hgrad}, \texttt{h\allowbreak{}Lap}, \texttt{h\allowbreak{}G} in that order. Named proof applications include \texttt{h\allowbreak{}G.\allowbreak{}sub}; the adjacent proof shows their hypotheses and order of use.




\texttt{gradient\allowbreak{}Generator\_\allowbreak{}eq} gives the formal counterpart of this step.
\end{leanbox}
\subsection{\texorpdfstring{Bound the two actual derivative terms pointwise}{Bound the two actual derivative terms pointwise}}
\label{report-step-AverageLMC-LaplacianGenerator-lean-2}


A globally \(L\)-Lipschitz gradient has Fréchet derivative operator norm at most \(L\) at every point. Its value on \(\nabla f\) is therefore bounded by \(L\|\nabla f\|\). Separately, the paper’s hypothesis \(\|\nabla \Delta f\|^{2}\le M_\Delta ^{2}p\) implies \(\|\nabla \Delta f\|\le M_\Delta \sqrt p\) because both norms and \(M_\Delta\) are nonnegative. Neither bound asks for the Hessian to be Lipschitz.

\begin{leanbox}{Lean: \texttt{hessian\_\allowbreak{}gradient\_\allowbreak{}norm\_\allowbreak{}le}, \texttt{gradient\_\allowbreak{}laplacian\_\allowbreak{}norm\_\allowbreak{}le}}
\leansource{\texttt{hessian\_\allowbreak{}gradient\_\allowbreak{}norm\_\allowbreak{}le} \quad AverageLMC/\allowbreak{}LaplacianGenerator.lean, lines 41–45}
\begin{leancode}
\leanline{0}{theorem\ hessian\_gradient\_norm\_le\ \{p\ :\ ℕ\}\ (f\ :\ State\ p\ →\ ℝ)\ (L\ :\ ℝ≥0)}
\leanline{4}{(hLip\ :\ LipschitzWith\ L\ (gradient\ f))\ (x\ :\ State\ p)\ :}
\leanline{4}{‖fderiv\ ℝ\ (gradient\ f)\ x\ (gradient\ f\ x)‖\ ≤\ (L\ :\ ℝ)\ *\ ‖gradient\ f\ x‖}
\end{leancode}
\leansource{\texttt{gradient\_\allowbreak{}laplacian\_\allowbreak{}norm\_\allowbreak{}le} \quad AverageLMC/\allowbreak{}LaplacianGenerator.lean, lines 49–56}
\begin{leancode}
\leanline{0}{theorem\ gradient\_laplacian\_norm\_le\ \{p\ :\ ℕ\}\ (f\ :\ State\ p\ →\ ℝ)\ (Mdelta\ :\ ℝ)}
\leanline{4}{(hMdelta\ :\ 0\ ≤\ Mdelta)\ (hδ\ :\ LaplacianSmooth\ f\ Mdelta)\ (x\ :\ State\ p)\ :}
\leanline{4}{‖gradient\ (laplacian\ f)\ x‖\ ≤\ Mdelta\ *\ Real.sqrt\ (p\ :\ ℝ)}
\end{leancode}
\leanlabel{How the Lean proof follows the argument}
\texttt{hessian\_\allowbreak{}gradient\_\allowbreak{}norm\_\allowbreak{}le} gives the formal counterpart of this step.




\texttt{gradient\_\allowbreak{}laplacian\_\allowbreak{}norm\_\allowbreak{}le} gives the formal counterpart of this step.
\end{leanbox}
\subsection{\texorpdfstring{Construct finite second moments for the generator}{Construct finite second moments for the generator}}
\label{report-step-AverageLMC-LaplacianGenerator-lean-3}


The Laplacian-gradient term is a continuous function bounded uniformly by \(M_\Delta \sqrt p\), hence belongs to \(L^{2}\) on every probability space. The Hessian-gradient term is continuous and bounded by \(L\) times the actual gradient; the latter has finite second moment because it is a Lipschitz field applied to an \(L^{2}\) state. Subtraction constructs the finite second moment of the actual generator itself. These are conclusions of the potential assumptions, not moment premises supplied to a local-error theorem.

\begin{leanbox}{Lean: \texttt{laplacian\_\allowbreak{}gradient\_\allowbreak{}memLp}, \texttt{hessian\_\allowbreak{}gradient\_\allowbreak{}memLp}, \texttt{gradientGenerator\_\allowbreak{}memLp}}
\leansource{\texttt{laplacian\_\allowbreak{}gradient\_\allowbreak{}memLp} \quad AverageLMC/\allowbreak{}LaplacianGenerator.lean, lines 61–71}
\begin{leancode}
\leanline{0}{theorem\ laplacian\_gradient\_memLp\ \{p\ :\ ℕ\}\ (f\ :\ State\ p\ →\ ℝ)\ (Mdelta\ :\ ℝ)}
\leanline{4}{(hMdelta\ :\ 0\ ≤\ Mdelta)\ (hδ\ :\ LaplacianSmooth\ f\ Mdelta)\ (Y\ :\ Lp\ (State\ p)\ 2\ P)\ :}
\leanline{4}{MemLp\ (fun\ ω\ ↦\ gradient\ (laplacian\ f)\ (Y\ ω))\ 2\ P}
\end{leancode}
\leansource{\texttt{hessian\_\allowbreak{}gradient\_\allowbreak{}memLp} \quad AverageLMC/\allowbreak{}LaplacianGenerator.lean, lines 74–82}
\begin{leancode}
\leanline{0}{theorem\ hessian\_gradient\_memLp\ \{p\ :\ ℕ\}\ (f\ :\ State\ p\ →\ ℝ)\ (L\ :\ ℝ≥0)}
\leanline{4}{(hf\ :\ ContDiff\ ℝ\ 3\ f)\ (hLip\ :\ LipschitzWith\ L\ (gradient\ f))\ (Y\ :\ Lp\ (State\ p)\ 2\ P)\ :}
\leanline{4}{MemLp\ (fun\ ω\ ↦\ fderiv\ ℝ\ (gradient\ f)\ (Y\ ω)\ (gradient\ f\ (Y\ ω)))\ 2\ P}
\end{leancode}
\leansource{\texttt{gradientGenerator\_\allowbreak{}memLp} \quad AverageLMC/\allowbreak{}LaplacianGenerator.lean, lines 86–91}
\begin{leancode}
\leanline{0}{theorem\ gradientGenerator\_memLp\ \{p\ :\ ℕ\}\ (f\ :\ State\ p\ →\ ℝ)\ (L\ :\ ℝ≥0)}
\leanline{4}{(hLip\ :\ LipschitzWith\ L\ (gradient\ f))\ (Mdelta\ :\ ℝ)\ (hMdelta\ :\ 0\ ≤\ Mdelta)}
\leanline{4}{(hδ\ :\ LaplacianSmooth\ f\ Mdelta)\ (Y\ :\ Lp\ (State\ p)\ 2\ P)\ :}
\leanline{4}{MemLp\ (fun\ ω\ ↦\ gradientGenerator\ f\ (Y\ ω))\ 2\ P}
\end{leancode}
\leanlabel{How the Lean proof follows the argument}
\texttt{laplacian\_\allowbreak{}gradient\_\allowbreak{}mem\allowbreak{}Lp} gives the formal counterpart of this step. The main intermediate facts are \texttt{h\allowbreak{}Lap}, \texttt{hc} in that order. Named proof applications include \texttt{ae\_\allowbreak{}of\_\allowbreak{}all}; the adjacent proof shows their hypotheses and order of use.




\texttt{hessian\_\allowbreak{}gradient\_\allowbreak{}mem\allowbreak{}Lp} gives the formal counterpart of this step. The main intermediate facts are \texttt{hgrad}, \texttt{hc} in that order.




\texttt{gradient\allowbreak{}Generator\_\allowbreak{}mem\allowbreak{}Lp} gives the formal counterpart of this step.
\end{leanbox}


Earlier results used here: \cite[Section~3.13.4]{LangevinLean2026ULASC}.
\subsection{\texorpdfstring{Derive the sharp stationary generator constant}{Derive the sharp stationary generator constant}}
\label{report-step-AverageLMC-LaplacianGenerator-lean-4}


Represent each generator term in \(L^{2}\) and use the exact almost-everywhere pointwise identities. The first representative has norm at most \(M_\Delta \sqrt p\); the second has norm at most \(L\) times the lifted score norm. The triangle inequality therefore bounds the actual generator norm by \(L\|\nabla f(Y)\|_{2}+M_\Delta \sqrt p\). When \(Y\) has the true Gibbs law, pushforward integration identifies its score second moment with the Gibbs integral. The sharp coordinate-score theorem bounds that integral by \(M_{\mathrm{av}}p\). Taking its nonnegative square root yields the required stationary constant \(L\sqrt{M_{\mathrm{av}}p}+M_\Delta\sqrt p\).

\begin{leanbox}{Lean: \texttt{lp\_\allowbreak{}gradientGenerator\_\allowbreak{}norm\_\allowbreak{}le}, \texttt{lp\_\allowbreak{}gradientGenerator\_\allowbreak{}gibbs\_\allowbreak{}norm\_\allowbreak{}le}}
\leansource{\texttt{lp\_\allowbreak{}gradientGenerator\_\allowbreak{}norm\_\allowbreak{}le} \quad AverageLMC/\allowbreak{}LaplacianGenerator.lean, lines 95–120}
\begin{leancode}
\leanline{0}{theorem\ lp\_gradientGenerator\_norm\_le\ \{p\ :\ ℕ\}\ (f\ :\ State\ p\ →\ ℝ)\ (L\ :\ ℝ≥0)}
\leanline{4}{(hLip\ :\ LipschitzWith\ L\ (gradient\ f))\ (Mdelta\ :\ ℝ)\ (hMdelta\ :\ 0\ ≤\ Mdelta)}
\leanline{4}{(hδ\ :\ LaplacianSmooth\ f\ Mdelta)\ (Y\ :\ Lp\ (State\ p)\ 2\ P)\ :}
\leanline{4}{‖(gradientGenerator\_memLp\ P\ f\ L\ hLip\ Mdelta\ hMdelta\ hδ\ Y).toLp}
\leanline{6}{(fun\ ω\ ↦\ gradientGenerator\ f\ (Y\ ω))‖\ ≤}
\leanline{6}{(L\ :\ ℝ)\ *\ ‖ULA.lpDrift\ P\ (gradient\ f)\ hLip\ Y‖\ +\ Mdelta\ *\ Real.sqrt\ (p\ :\ ℝ)}
\end{leancode}
\leansource{\texttt{lp\_\allowbreak{}gradientGenerator\_\allowbreak{}gibbs\_\allowbreak{}norm\_\allowbreak{}le} \quad AverageLMC/\allowbreak{}LaplacianGenerator.lean, lines 124–144}
\begin{leancode}
\leanline{0}{theorem\ lp\_gradientGenerator\_gibbs\_norm\_le\ \{p\ :\ ℕ\}\ (hp\ :\ 0\ <\ p)}
\leanline{4}{(f\ :\ State\ p\ →\ ℝ)\ (m\ :\ ℝ)\ (L\ :\ ℝ≥0)\ (hm\ :\ 0\ <\ m)}
\leanline{4}{(hsc\ :\ StrongConvexOn\ univ\ m\ f)\ (hLip\ :\ LipschitzWith\ L\ (gradient\ f))}
\leanline{4}{(Mdelta\ :\ ℝ)\ (hMdelta\ :\ 0\ ≤\ Mdelta)\ (hδ\ :\ LaplacianSmooth\ f\ Mdelta)}
\leanline{4}{(Y\ :\ Lp\ (State\ p)\ 2\ P)\ (hlaw\ :\ P.map\ Y\ =\ ULA.gibbsMeasure\ f)\ :}
\leanline{4}{‖(gradientGenerator\_memLp\ P\ f\ L\ hLip\ Mdelta\ hMdelta\ hδ\ Y).toLp}
\leanline{6}{(fun\ ω\ ↦\ gradientGenerator\ f\ (Y\ ω))‖\ ≤}
\leanline{6}{(L\ :\ ℝ)\ *\ Real.sqrt\ (averageSmoothness\ f\ *\ (p\ :\ ℝ))\ +\ Mdelta\ *\ Real.sqrt\ (p\ :\ ℝ)}
\end{leancode}
\leanlabel{How the Lean proof follows the argument}
\texttt{lp\_\allowbreak{}gradient\allowbreak{}Generator\_\allowbreak{}norm\_\allowbreak{}le} gives the formal counterpart of this step. The main intermediate facts are \texttt{D}, \texttt{H}, \texttt{h\allowbreak{}D}, \texttt{h\allowbreak{}H}, \texttt{heq}, \texttt{h\allowbreak{}Dbound} in that order. Named proof applications include \texttt{hb}, \texttt{Lp.\allowbreak{}norm\_\allowbreak{}le\_\allowbreak{}mul\_\allowbreak{}norm\_\allowbreak{}of\_\allowbreak{}ae\_\allowbreak{}le\_\allowbreak{}mul}, \texttt{hessian\_\allowbreak{}gradient\_\allowbreak{}norm\_\allowbreak{}le}; the adjacent proof shows their hypotheses and order of use.




\texttt{lp\_\allowbreak{}gradient\allowbreak{}Generator\_\allowbreak{}gibbs\_\allowbreak{}norm\_\allowbreak{}le} gives the formal counterpart of this step. The main intermediate facts are \texttt{hf}, \texttt{hs}, \texttt{hscore} in that order. Named proof applications include \texttt{hs}; the adjacent proof shows their hypotheses and order of use.
\end{leanbox}


Earlier results used here: \cite[Section~3.13.2]{LangevinLean2026ULASC}; \cite[Section~3.13.3]{LangevinLean2026ULASC}; \cite[Section~3.13.9]{LangevinLean2026ULASC}; \cite[Section~3.26.3]{LangevinLean2026ULASC}.
\section{\texorpdfstring{The third-order one-step energy estimate for the actual Euler coupling}{The third-order one-step energy estimate for the actual Euler coupling}}
\label{proof-module-AverageLMC-LaplacianOneStep-lean}


The proof retains the cancellation in the defect cross term while bounding the complete defect by the existing first-order Brownian estimate. All gradient contraction and energy algebra concern the actual lifted gradient and actual integral solution.



In the \(L^2\) integral equation write \(g=\nabla f\), \(q=1-mh\), and \(e=\|X-Y_0\|_{L^2}\). The local defect and deterministic Euler difference give the exact decomposition


\[\begin{aligned}\Psi&=(X-hg(X))-(Y_0-hg(Y_0)),\\D&=\int_0^h\bigl(g(Y_t)-g(Y_0)\bigr)\,dt,\\X-hg(X)+B_h-Y_h&=\Psi+D.\end{aligned}\]


The proved pairing bound with constant \(R_{\rm gen}\), together with the score and Brownian moment bounds in dimension parameter \(d\), yields


\[\|X-hg(X)+B_h-Y_h\|_{L^2}^{2}\le q^2e^2+h^2qeR_{\rm gen}+\left(\frac32Lh\sqrt h\sqrt d\right)^2.\]

Earlier results used here: \cite[Section~3.13]{LangevinLean2026ULASC}.
\subsection{\texorpdfstring{Integrate the derived pairing bound and retain the exact cross coefficient}{Integrate the derived pairing bound and retain the exact cross coefficient}}
\label{report-step-AverageLMC-LaplacianOneStep-lean-1}


Let \(\Psi\) be the difference between the deterministic gradient steps of \(X\) and \(Y(0)\), and let \(D\) be the exact drift integral omitted by Euler. The integral equation gives the exact output error \(\Psi +D\). On the admissible range \(hM\le1\) and \(m\le M\), the actual smooth strongly convex gradient step contracts by \(q=1-mh\ge 0\), so \(\|\Psi \|\le q\|X-Y(0)\|\). The first-order total-defect theorem gives \(\|D\|\le (3/2)Mh\sqrt h\sqrt d\). The separately proved weak-Dynkin pairing inequality bounds the pairing with the instantaneous gradient increment by \(t\|\Psi \|R_{\mathrm{gen}}\). Continuous linear integration then bounds \(\langle \Psi ,D\rangle\) by \(h^{2}\|\Psi \|R_{\mathrm{gen}}/2\). Expanding \(\|\Psi +D\|^{2}\) and substituting the three bounds yields \(q^{2}\) times the original error squared, the exact cross term \(h^{2}q\|X-Y(0)\|R_{\mathrm{gen}}\), and the squared total-defect bound. The pairing premise is an intermediate analytic consequence of the actual stationary diffusion, discharged by the Dynkin and cutoff construction.

\begin{leanbox}{Lean: \texttt{lp\_\allowbreak{}laplacian\_\allowbreak{}one\_\allowbreak{}step\_\allowbreak{}bound}}
\leansource{\texttt{lp\_\allowbreak{}laplacian\_\allowbreak{}one\_\allowbreak{}step\_\allowbreak{}bound} \quad AverageLMC/\allowbreak{}LaplacianOneStep.lean, lines 21–88}
\begin{leancode}
\leanline{0}{theorem\ lp\_laplacian\_one\_step\_bound}
\leanline{4}{(f\ :\ E\ →\ ℝ)\ (m\ h\ d\ Rgen\ :\ ℝ)\ (M\ :\ ℝ≥0)}
\leanline{4}{(hf\ :\ Differentiable\ ℝ\ f)\ (hsc\ :\ StrongConvexOn\ univ\ m\ f)}
\leanline{4}{(hLip\ :\ LipschitzWith\ M\ (gradient\ f))\ (hm\ :\ 0\ <\ m)\ (hmM\ :\ m\ ≤\ (M\ :\ ℝ))}
\leanline{4}{(hh\ :\ 0\ ≤\ h)\ (hd\ :\ 0\ ≤\ d)\ (hsmall\ :\ h\ *\ (M\ :\ ℝ)\ ≤\ 1)\ (hRgen\ :\ 0\ ≤\ Rgen)}
\leanline{4}{(X\ :\ Lp\ E\ 2\ P)\ (Y\ B\ :\ ℝ\ →\ Lp\ E\ 2\ P)\ (hY\ :\ Continuous\ Y)}
\leanline{4}{(hG\ :\ ∀\ t\ ∈\ Icc\ 0\ h,\ ‖ULA.lpDrift\ P\ (gradient\ f)\ hLip\ (Y\ t)‖\ ≤}
\leanline{6}{Real.sqrt\ ((M\ :\ ℝ)\ *\ d))}
\leanline{4}{(hB\ :\ ∀\ t\ ∈\ Icc\ 0\ h,\ ‖B\ t‖\ ≤\ Real.sqrt\ (2\ *\ d)\ *\ Real.sqrt\ t)}
\leanline{4}{(hYe\ :\ ∀\ t\ ∈\ Icc\ 0\ h,\ Y\ t\ -\ Y\ 0\ =}
\leanline{6}{-(∫\ s\ in\ 0..t,\ ULA.lpDrift\ P\ (gradient\ f)\ hLip\ (Y\ s))\ +\ B\ t)}
\leanline{4}{(hpair\ :\ ∀\ t\ ∈\ Icc\ 0\ h,}
\leanline{6}{{\leanfallback ⟪}(X\ -\ h\ •\ ULA.lpDrift\ P\ (gradient\ f)\ hLip\ X)\ -}
\leanline{10}{(Y\ 0\ -\ h\ •\ ULA.lpDrift\ P\ (gradient\ f)\ hLip\ (Y\ 0)),}
\leanline{8}{ULA.lpDrift\ P\ (gradient\ f)\ hLip\ (Y\ t)\ -}
\leanline{10}{ULA.lpDrift\ P\ (gradient\ f)\ hLip\ (Y\ 0){\leanfallback ⟫}\ ≤}
\leanline{8}{t\ *\ ‖(X\ -\ h\ •\ ULA.lpDrift\ P\ (gradient\ f)\ hLip\ X)\ -}
\leanline{10}{(Y\ 0\ -\ h\ •\ ULA.lpDrift\ P\ (gradient\ f)\ hLip\ (Y\ 0))‖\ *\ Rgen)\ :}
\leanline{4}{‖X\ -\ h\ •\ ULA.lpDrift\ P\ (gradient\ f)\ hLip\ X\ +\ B\ h\ -\ Y\ h‖\ \textasciicircum{}\ 2\ ≤}
\leanline{6}{(1\ -\ m\ *\ h)\ \textasciicircum{}\ 2\ *\ ‖X\ -\ Y\ 0‖\ \textasciicircum{}\ 2\ +}
\leanline{6}{h\ \textasciicircum{}\ 2\ *\ (1\ -\ m\ *\ h)\ *\ ‖X\ -\ Y\ 0‖\ *\ Rgen\ +}
\leanline{6}{((3\ /\ 2\ :\ ℝ)\ *\ (M\ :\ ℝ)\ *\ h\ *\ Real.sqrt\ h\ *\ Real.sqrt\ d)\ \textasciicircum{}\ 2}
\end{leancode}
\leanlabel{How the Lean proof follows the argument}
\texttt{lp\_\allowbreak{}laplacian\_\allowbreak{}one\_\allowbreak{}step\_\allowbreak{}bound} gives the formal counterpart of this step. The main intermediate facts are \texttt{g}, \texttt{\ensuremath{\Psi}}, \texttt{D}, \texttt{hg}, \texttt{hgc}, \texttt{he} in that order. Named proof applications include \texttt{hs}; the adjacent proof shows their hypotheses and order of use.
\end{leanbox}


Earlier results used here: \cite[Section~3.11.1]{LangevinLean2026ULASC}; \cite[Section~3.11.3]{LangevinLean2026ULASC}; \cite[Section~3.13.2]{LangevinLean2026ULASC}; \cite[Section~3.13.5]{LangevinLean2026ULASC}; \cite[Section~3.13.6]{LangevinLean2026ULASC}.
\section{\texorpdfstring{Uniform Euler mesh convergence in the actual \(L^{2}\) space}{Uniform Euler mesh convergence in the actual L² space}}
\label{proof-module-AverageLMC-LaplacianEulerLimit-lean}


The existing additive integral equation is used directly in a Banach space. Matching-noise Euler convergence is strengthened from its terminal value to every mesh point, giving the uniform drift moments required to sum weak errors.



Let \(E_k^{(n)}\) be the matching-noise Euler chain with mesh \(h_n=T/(n+1)\). The integral equation implies simultaneous mesh convergence:


\[\max_{0\le k\le n+1}\left\|E_k^{(n)}-Y(kh_n)\right\|
\longrightarrow0.\]


For every uniformly continuous observation \(F\) into a Banach space,


\[\sum_{k=0}^{n}h_nF\bigl(E_k^{(n)}\bigr)
\longrightarrow\int_0^T F(Y(t))\,dt.\]

Earlier results used here: \cite[Section~3.16]{LangevinLean2026ULASC}.
\subsection{\texorpdfstring{Estimate every grid point with one drift modulus}{Estimate every grid point with one drift modulus}}
\label{report-step-AverageLMC-LaplacianEulerLimit-lean-1}


Restrict the existing Euler drift-modulus estimate to the shorter interval ending at \(kh\). Monotonicity of \(kh\) and the exponential factor gives the same upper bound \(T\varepsilon e^{LT}\) at every point \(k\le n\).

\begin{leanbox}{Lean: \texttt{additiveEulerChain\_\allowbreak{}grid\_\allowbreak{}error\_\allowbreak{}le\_\allowbreak{}of\_\allowbreak{}drift\_\allowbreak{}modulus}}
\leansource{\texttt{additiveEulerChain\_\allowbreak{}grid\_\allowbreak{}error\_\allowbreak{}le\_\allowbreak{}of\_\allowbreak{}drift\_\allowbreak{}modulus} \quad AverageLMC/\allowbreak{}LaplacianEulerLimit.lean, lines 23–46}
\begin{leancode}
\leanline{0}{theorem\ additiveEulerChain\_grid\_error\_le\_of\_drift\_modulus}
\leanline{4}{(g\ :\ H\ →\ H)\ (L\ :\ ℝ≥0)\ (hLip\ :\ LipschitzWith\ L\ g)}
\leanline{4}{(x\ :\ H)\ (B\ Y\ :\ ℝ\ →\ H)\ (hB0\ :\ B\ 0\ =\ 0)\ (hY\ :\ Continuous\ Y)}
\leanline{4}{(T\ h\ ε\ :\ ℝ)\ (hT\ :\ 0\ ≤\ T)\ (hh\ :\ 0\ ≤\ h)\ (hε\ :\ 0\ ≤\ ε)\ (n\ :\ ℕ)}
\leanline{4}{(hpartition\ :\ (n\ :\ ℝ)\ *\ h\ =\ T)}
\leanline{4}{(hYe\ :\ ∀\ t\ ∈\ Icc\ 0\ T,\ Y\ t\ -\ x\ =\ -(∫\ s\ in\ 0..t,\ g\ (Y\ s))\ +\ B\ t)}
\leanline{4}{(hmod\ :\ ∀\ s\ ∈\ Icc\ 0\ T,\ ∀\ t\ ∈\ Icc\ 0\ T,\ dist\ s\ t\ ≤\ h\ →}
\leanline{6}{‖g\ (Y\ s)\ -\ g\ (Y\ t)‖\ ≤\ ε)\ (k\ :\ ℕ)\ (hk\ :\ k\ ≤\ n)\ :}
\leanline{4}{‖ULA.additiveEulerChain\ g\ h\ x\ B\ k\ -\ Y\ ((k\ :\ ℝ)\ *\ h)‖\ ≤}
\leanline{6}{T\ *\ ε\ *\ Real.exp\ ((L\ :\ ℝ)\ *\ T)}
\end{leancode}
\leanlabel{How the Lean proof follows the argument}
\texttt{additive\allowbreak{}Euler\allowbreak{}Chain\_\allowbreak{}grid\_\allowbreak{}error\_\allowbreak{}le\_\allowbreak{}of\_\allowbreak{}drift\_\allowbreak{}modulus} gives the formal counterpart of this step. The main intermediate facts are \texttt{hkh}, \texttt{hk\allowbreak{}T}, \texttt{hb} in that order. Named proof applications include \texttt{mul\_\allowbreak{}le\_\allowbreak{}mul\_\allowbreak{}of\_\allowbreak{}nonneg\_\allowbreak{}right}, \texttt{hb.\allowbreak{}trans}, \texttt{mul\_\allowbreak{}le\_\allowbreak{}mul}; the adjacent proof shows their hypotheses and order of use.
\end{leanbox}


Earlier results used here: \cite[Section~3.16.1]{LangevinLean2026ULASC}; \cite[Section~3.16.7]{LangevinLean2026ULASC}.
\subsection{\texorpdfstring{Let the compact-time modulus shrink}{Let the compact-time modulus shrink}}
\label{report-step-AverageLMC-LaplacianEulerLimit-lean-2}


The drift along the continuous integral solution is uniformly continuous on [0,\(T\)]. Choose a drift tolerance from the requested path error, then refine the mesh below its continuity radius. The preceding estimate proves simultaneous convergence of all Euler grid points.

\begin{leanbox}{Lean: \texttt{additiveEulerChain\_\allowbreak{}grid\_\allowbreak{}uniform\_\allowbreak{}tendsto}}
\leansource{\texttt{additiveEulerChain\_\allowbreak{}grid\_\allowbreak{}uniform\_\allowbreak{}tendsto} \quad AverageLMC/\allowbreak{}LaplacianEulerLimit.lean, lines 51–82}
\begin{leancode}
\leanline{0}{theorem\ additiveEulerChain\_grid\_uniform\_tendsto}
\leanline{4}{(g\ :\ H\ →\ H)\ (L\ :\ ℝ≥0)\ (hLip\ :\ LipschitzWith\ L\ g)}
\leanline{4}{(x\ :\ H)\ (B\ Y\ :\ ℝ\ →\ H)\ (hB0\ :\ B\ 0\ =\ 0)\ (hY\ :\ Continuous\ Y)}
\leanline{4}{(T\ :\ ℝ)\ (hT\ :\ 0\ ≤\ T)}
\leanline{4}{(hYe\ :\ ∀\ t\ ∈\ Icc\ 0\ T,\ Y\ t\ -\ x\ =\ -(∫\ s\ in\ 0..t,\ g\ (Y\ s))\ +\ B\ t)\ :}
\leanline{4}{∀\ ε\ >\ 0,\ ∃\ N\ :\ ℕ,\ ∀\ n\ ≥\ N,\ ∀\ k\ ≤\ n\ +\ 1,}
\leanline{6}{‖ULA.additiveEulerChain\ g\ (T\ /\ ((n\ +\ 1\ :\ ℕ)\ :\ ℝ))\ x\ B\ k\ -}
\leanline{8}{Y\ ((k\ :\ ℝ)\ *\ (T\ /\ ((n\ +\ 1\ :\ ℕ)\ :\ ℝ)))‖\ <\ ε}
\end{leancode}
\leanlabel{How the Lean proof follows the argument}
\texttt{additive\allowbreak{}Euler\allowbreak{}Chain\_\allowbreak{}grid\_\allowbreak{}uniform\_\allowbreak{}tendsto} gives the formal counterpart of this step. The main intermediate facts are \texttt{C}, \texttt{\ensuremath{\eta}}, \texttt{h\allowbreak{}C}, \texttt{h\ensuremath{\eta}}, \texttt{hscale}, \texttt{hsmall} in that order. Named proof applications include \texttt{hosc}, \texttt{hb.\allowbreak{}trans\_\allowbreak{}lt}; the adjacent proof shows their hypotheses and order of use.
\end{leanbox}


Earlier results used here: \cite[Section~3.16.1]{LangevinLean2026ULASC}.
\subsection{\texorpdfstring{Derive uniform state and score moments}{Derive uniform state and score moments}}
\label{report-step-AverageLMC-LaplacianEulerLimit-lean-3}


The true continuous path has bounded norm on the compact time interval. Eventually all Euler grid points are within distance one of this path. Lipschitz growth then bounds every drift value by \(L\) times this common state radius plus \(\|g(0)\|\). In the \(L^{2}\) application these are actual uniform second-moment bounds.

\begin{leanbox}{Lean: \texttt{additiveEulerChain\_\allowbreak{}grid\_\allowbreak{}eventually\_\allowbreak{}bounded}, \texttt{additiveEulerChain\_\allowbreak{}drift\_\allowbreak{}grid\_\allowbreak{}eventually\_\allowbreak{}bounded}}
\leansource{\texttt{additiveEulerChain\_\allowbreak{}grid\_\allowbreak{}eventually\_\allowbreak{}bounded} \quad AverageLMC/\allowbreak{}LaplacianEulerLimit.lean, lines 87–110}
\begin{leancode}
\leanline{0}{theorem\ additiveEulerChain\_grid\_eventually\_bounded}
\leanline{4}{(g\ :\ H\ →\ H)\ (L\ :\ ℝ≥0)\ (hLip\ :\ LipschitzWith\ L\ g)}
\leanline{4}{(x\ :\ H)\ (B\ Y\ :\ ℝ\ →\ H)\ (hB0\ :\ B\ 0\ =\ 0)\ (hY\ :\ Continuous\ Y)}
\leanline{4}{(T\ :\ ℝ)\ (hT\ :\ 0\ ≤\ T)}
\leanline{4}{(hYe\ :\ ∀\ t\ ∈\ Icc\ 0\ T,\ Y\ t\ -\ x\ =\ -(∫\ s\ in\ 0..t,\ g\ (Y\ s))\ +\ B\ t)\ :}
\leanline{4}{∃\ R\ :\ ℝ,\ 0\ ≤\ R\ ∧\ ∃\ N\ :\ ℕ,\ ∀\ n\ ≥\ N,\ ∀\ k\ ≤\ n\ +\ 1,}
\leanline{6}{‖ULA.additiveEulerChain\ g\ (T\ /\ ((n\ +\ 1\ :\ ℕ)\ :\ ℝ))\ x\ B\ k‖\ ≤\ R}
\end{leancode}
\leansource{\texttt{additiveEulerChain\_\allowbreak{}drift\_\allowbreak{}grid\_\allowbreak{}eventually\_\allowbreak{}bounded} \quad AverageLMC/\allowbreak{}LaplacianEulerLimit.lean, lines 115–130}
\begin{leancode}
\leanline{0}{theorem\ additiveEulerChain\_drift\_grid\_eventually\_bounded}
\leanline{4}{(g\ :\ H\ →\ H)\ (L\ :\ ℝ≥0)\ (hLip\ :\ LipschitzWith\ L\ g)}
\leanline{4}{(x\ :\ H)\ (B\ Y\ :\ ℝ\ →\ H)\ (hB0\ :\ B\ 0\ =\ 0)\ (hY\ :\ Continuous\ Y)}
\leanline{4}{(T\ :\ ℝ)\ (hT\ :\ 0\ ≤\ T)}
\leanline{4}{(hYe\ :\ ∀\ t\ ∈\ Icc\ 0\ T,\ Y\ t\ -\ x\ =\ -(∫\ s\ in\ 0..t,\ g\ (Y\ s))\ +\ B\ t)\ :}
\leanline{4}{∃\ R\ :\ ℝ,\ 0\ ≤\ R\ ∧\ ∃\ N\ :\ ℕ,\ ∀\ n\ ≥\ N,\ ∀\ k\ ≤\ n\ +\ 1,}
\leanline{6}{‖g\ (ULA.additiveEulerChain\ g\ (T\ /\ ((n\ +\ 1\ :\ ℕ)\ :\ ℝ))\ x\ B\ k)‖\ ≤\ R}
\end{leancode}
\leanlabel{How the Lean proof follows the argument}
\texttt{additive\allowbreak{}Euler\allowbreak{}Chain\_\allowbreak{}grid\_\allowbreak{}eventually\_\allowbreak{}bounded} gives the formal counterpart of this step. The main intermediate facts are \texttt{h}, \texttt{hh}, \texttt{hgrid}, \texttt{hbound}, \texttt{he}, \texttt{hy} in that order. The final step is the algebraic normalization or ordered arithmetic displayed in the proof.




\texttt{additive\allowbreak{}Euler\allowbreak{}Chain\_\allowbreak{}drift\_\allowbreak{}grid\_\allowbreak{}eventually\_\allowbreak{}bounded} gives the formal counterpart of this step. The main intermediate facts are \texttt{z}, \texttt{hz}, \texttt{hg}, \texttt{hn} in that order. Named proof applications include \texttt{h\allowbreak{}Lip.\allowbreak{}norm\_\allowbreak{}sub\_\allowbreak{}le}, \texttt{hn.\allowbreak{}trans}; the adjacent proof shows their hypotheses and order of use.
\end{leanbox}


Earlier results used here: \cite[Section~3.16.1]{LangevinLean2026ULASC}.
\subsection{\texorpdfstring{Prove the uniform left-sum estimate}{Prove the uniform left-sum estimate}}
\label{report-step-AverageLMC-LaplacianEulerLimit-lean-4}


On each partition cell, integrate the difference between the left-endpoint value and the actual continuous integrand. A modulus \(\varepsilon\) bounds its norm by \(h\varepsilon\). The adjacent-interval identity and the norm of a finite sum give total error at most \(T\varepsilon\). Compact-time uniform continuity then proves convergence of the uniform Riemann sums.

\begin{leanbox}{Lean: \texttt{uniform\_\allowbreak{}left\_\allowbreak{}sum\_\allowbreak{}error\_\allowbreak{}le\_\allowbreak{}of\_\allowbreak{}modulus}, \texttt{uniform\_\allowbreak{}left\_\allowbreak{}sum\_\allowbreak{}tendsto}}
\leansource{\texttt{uniform\_\allowbreak{}left\_\allowbreak{}sum\_\allowbreak{}error\_\allowbreak{}le\_\allowbreak{}of\_\allowbreak{}modulus} \quad AverageLMC/\allowbreak{}LaplacianEulerLimit.lean, lines 133–172}
\begin{leancode}
\leanline{0}{theorem\ uniform\_left\_sum\_error\_le\_of\_modulus}
\leanline{4}{(f\ :\ ℝ\ →\ H)\ (hf\ :\ Continuous\ f)\ (T\ h\ ε\ :\ ℝ)}
\leanline{4}{(hh\ :\ 0\ ≤\ h)\ (hε\ :\ 0\ ≤\ ε)\ (n\ :\ ℕ)\ (hpartition\ :\ (n\ :\ ℝ)\ *\ h\ =\ T)}
\leanline{4}{(hmod\ :\ ∀\ s\ ∈\ Icc\ 0\ T,\ ∀\ t\ ∈\ Icc\ 0\ T,\ dist\ s\ t\ ≤\ h\ →\ ‖f\ s\ -\ f\ t‖\ ≤\ ε)\ :}
\leanline{4}{‖(∑\ k\ ∈\ Finset.range\ n,\ h\ •\ f\ ((k\ :\ ℝ)\ *\ h))\ -\ ∫\ t\ in\ 0..T,\ f\ t‖\ ≤\ T\ *\ ε}
\end{leancode}
\leansource{\texttt{uniform\_\allowbreak{}left\_\allowbreak{}sum\_\allowbreak{}tendsto} \quad AverageLMC/\allowbreak{}LaplacianEulerLimit.lean, lines 175–197}
\begin{leancode}
\leanline{0}{theorem\ uniform\_left\_sum\_tendsto}
\leanline{4}{(f\ :\ ℝ\ →\ H)\ (hf\ :\ Continuous\ f)\ (T\ :\ ℝ)\ (hT\ :\ 0\ ≤\ T)\ :}
\leanline{4}{Tendsto\ (fun\ n\ :\ ℕ\ ↦\ ∑\ k\ ∈\ Finset.range\ (n\ +\ 1),}
\leanline{6}{(T\ /\ ((n\ +\ 1\ :\ ℕ)\ :\ ℝ))\ •\ f\ ((k\ :\ ℝ)\ *\ (T\ /\ ((n\ +\ 1\ :\ ℕ)\ :\ ℝ))))}
\leanline{6}{atTop\ ({\leanfallback 𝓝}\ (∫\ t\ in\ 0..T,\ f\ t))}
\end{leancode}
\leanlabel{How the Lean proof follows the argument}
\texttt{uniform\_\allowbreak{}left\_\allowbreak{}sum\_\allowbreak{}error\_\allowbreak{}le\_\allowbreak{}of\_\allowbreak{}modulus} gives the formal counterpart of this step. The main intermediate facts are \texttt{hcell}, \texttt{hi} in that order. Named proof applications include \texttt{mul\_\allowbreak{}le\_\allowbreak{}mul\_\allowbreak{}of\_\allowbreak{}nonneg\_\allowbreak{}right}, \texttt{hmod}; the adjacent proof shows their hypotheses and order of use.




\texttt{uniform\_\allowbreak{}left\_\allowbreak{}sum\_\allowbreak{}tendsto} gives the formal counterpart of this step. The main intermediate facts are \texttt{\ensuremath{\eta}}, \texttt{h\ensuremath{\eta}}, \texttt{hscale}, \texttt{hsmall}, \texttt{huc}, \texttt{hmesh} in that order. Named proof applications include \texttt{Metric.\allowbreak{}tendsto\_\allowbreak{}at\allowbreak{}Top.\allowbreak{}mpr}, \texttt{hosc}; the adjacent proof shows their hypotheses and order of use.
\end{leanbox}
\subsection{\texorpdfstring{Pass observations of Euler states to time integrals}{Pass observations of Euler states to time integrals}}
\label{report-step-AverageLMC-LaplacianEulerLimit-lean-5}


Uniform continuity of an observation sends the simultaneous Euler grid error to one common observation error. Sum this error over the partition, and combine it with the continuous path's left Riemann-sum limit. No differentiability of the Brownian path is required.

\begin{leanbox}{Lean: \texttt{additiveEulerChain\_\allowbreak{}test\_\allowbreak{}sum\_\allowbreak{}tendsto}}
\leansource{\texttt{additiveEulerChain\_\allowbreak{}test\_\allowbreak{}sum\_\allowbreak{}tendsto} \quad AverageLMC/\allowbreak{}LaplacianEulerLimit.lean, lines 201–257}
\begin{leancode}
\leanline{0}{theorem\ additiveEulerChain\_test\_sum\_tendsto}
\leanline{4}{\{G\ :\ Type*\}\ [NormedAddCommGroup\ G]\ [NormedSpace\ ℝ\ G]\ [CompleteSpace\ G]}
\leanline{4}{(g\ :\ H\ →\ H)\ (L\ :\ ℝ≥0)\ (hLip\ :\ LipschitzWith\ L\ g)}
\leanline{4}{(x\ :\ H)\ (B\ Y\ :\ ℝ\ →\ H)\ (hB0\ :\ B\ 0\ =\ 0)\ (hY\ :\ Continuous\ Y)}
\leanline{4}{(T\ :\ ℝ)\ (hT\ :\ 0\ ≤\ T)}
\leanline{4}{(hYe\ :\ ∀\ t\ ∈\ Icc\ 0\ T,\ Y\ t\ -\ x\ =\ -(∫\ s\ in\ 0..t,\ g\ (Y\ s))\ +\ B\ t)}
\leanline{4}{(F\ :\ H\ →\ G)\ (hF\ :\ UniformContinuous\ F)\ :}
\leanline{4}{Tendsto\ (fun\ n\ :\ ℕ\ ↦\ ∑\ k\ ∈\ Finset.range\ (n\ +\ 1),}
\leanline{6}{(T\ /\ ((n\ +\ 1\ :\ ℕ)\ :\ ℝ))\ •}
\leanline{8}{F\ (ULA.additiveEulerChain\ g\ (T\ /\ ((n\ +\ 1\ :\ ℕ)\ :\ ℝ))\ x\ B\ k))}
\leanline{6}{atTop\ ({\leanfallback 𝓝}\ (∫\ t\ in\ 0..T,\ F\ (Y\ t)))}
\end{leancode}
\leanlabel{How the Lean proof follows the argument}
\texttt{additive\allowbreak{}Euler\allowbreak{}Chain\_\allowbreak{}test\_\allowbreak{}sum\_\allowbreak{}tendsto} gives the formal counterpart of this step. The main intermediate facts are \texttt{\ensuremath{\eta}}, \texttt{h\ensuremath{\eta}}, \texttt{hscale}, \texttt{hsmall}, \texttt{h}, \texttt{A} in that order. Named proof applications include \texttt{mul\_\allowbreak{}le\_\allowbreak{}mul\_\allowbreak{}of\_\allowbreak{}nonneg\_\allowbreak{}left}, \texttt{h\allowbreak{}N\ensuremath{{}_{1}}}, \texttt{norm\_\allowbreak{}add\_\allowbreak{}le}; the adjacent proof shows their hypotheses and order of use.
\end{leanbox}


Earlier results used here: \cite[Section~3.16.1]{LangevinLean2026ULASC}.
\section{\texorpdfstring{From summed Gaussian residuals to the exact Dynkin identity}{From summed Gaussian residuals to the exact Dynkin identity}}
\label{proof-module-AverageLMC-LaplacianDynkinLimit-lean}


This is the deterministic limit argument connecting the proved Gaussian weak steps to a continuous integral solution. It does not assume a stochastic calculus identity or a Wasserstein error recurrence.



For \(h_n=T/(n+1)\), the residual hypothesis is arbitrarily small at first order:


\[\forall\varepsilon>0\ \exists K\ge0:\quad
|r_n|\le C\bigl(\varepsilon+K\sqrt{h_n}\bigr)
\ \text{eventually},\qquad r_n\longrightarrow0.\]


Matching-noise Euler convergence and uniform continuity of the two observations then give the exact identity


\[\operatorname{test}(Y(T))-\operatorname{test}(Y(0))
=\int_0^T\operatorname{generator}(Y(t))\,dt.\]\subsection{\texorpdfstring{Force the accumulated weak residual to zero}{Force the accumulated weak residual to zero}}
\label{report-step-AverageLMC-LaplacianDynkinLimit-lean-1}


Suppose every positive \(\varepsilon\) permits an eventual bound \(C[\varepsilon +K\sqrt{T/(n+1)}]\), with \(K\) allowed to depend on \(\varepsilon\). First fix \(\varepsilon\) small relative to the desired error; then refine the mesh so its square-root term is equally small. This proves that the actual scalar residual converges to zero.

\begin{leanbox}{Lean: \texttt{weakEuler\_\allowbreak{}residual\_\allowbreak{}tendsto\_\allowbreak{}zero}}
\leansource{\texttt{weakEuler\_\allowbreak{}residual\_\allowbreak{}tendsto\_\allowbreak{}zero} \quad AverageLMC/\allowbreak{}LaplacianDynkinLimit.lean, lines 24–46}
\begin{leancode}
\leanline{0}{theorem\ weakEuler\_residual\_tendsto\_zero}
\leanline{4}{(r\ :\ ℕ\ →\ ℝ)\ (T\ C\ :\ ℝ)\ (hT\ :\ 0\ ≤\ T)\ (hC\ :\ 0\ ≤\ C)}
\leanline{4}{(hres\ :\ ∀\ ε\ >\ 0,\ ∃\ K\ :\ ℝ,\ 0\ ≤\ K\ ∧\ ∃\ N\ :\ ℕ,\ ∀\ n\ ≥\ N,}
\leanline{6}{|r\ n|\ ≤\ C\ *\ (ε\ +\ K\ *\ Real.sqrt\ (T\ /\ ((n\ +\ 1\ :\ ℕ)\ :\ ℝ))))\ :}
\leanline{4}{Tendsto\ r\ atTop\ ({\leanfallback 𝓝}\ 0)}
\end{leancode}
\leanlabel{How the Lean proof follows the argument}
\texttt{weak\allowbreak{}Euler\_\allowbreak{}residual\_\allowbreak{}tendsto\_\allowbreak{}zero} gives the formal counterpart of this step. The main intermediate facts are \texttt{\ensuremath{\eta}}, \texttt{h\ensuremath{\eta}}, \texttt{hscale}, \texttt{hmesh}, \texttt{hroot}, \texttt{hbound} in that order. Named proof applications include \texttt{Metric.\allowbreak{}tendsto\_\allowbreak{}at\allowbreak{}Top.\allowbreak{}mpr}; the adjacent proof shows their hypotheses and order of use.
\end{leanbox}
\subsection{\texorpdfstring{Identify the unique limiting residual}{Identify the unique limiting residual}}
\label{report-step-AverageLMC-LaplacianDynkinLimit-lean-2}


Uniform continuity sends the terminal matching-noise Euler state to the terminal test observation. Uniform mesh convergence sends its generator sum to the time integral. Their difference therefore converges to the proposed Dynkin residual, while the summed weak remainder converges to zero. Uniqueness of real limits gives the exact identity.

\begin{leanbox}{Lean: \texttt{weighted\_\allowbreak{}dynkin\_\allowbreak{}of\_\allowbreak{}weak\_\allowbreak{}euler\_\allowbreak{}consistency}}
\leansource{\texttt{weighted\_\allowbreak{}dynkin\_\allowbreak{}of\_\allowbreak{}weak\_\allowbreak{}euler\_\allowbreak{}consistency} \quad AverageLMC/\allowbreak{}LaplacianDynkinLimit.lean, lines 53–80}
\begin{leancode}
\leanline{0}{theorem\ weighted\_dynkin\_of\_weak\_euler\_consistency}
\leanline{4}{(g\ :\ H\ →\ H)\ (L\ :\ ℝ≥0)\ (hLip\ :\ LipschitzWith\ L\ g)}
\leanline{4}{(x\ :\ H)\ (B\ Y\ :\ ℝ\ →\ H)\ (hB0\ :\ B\ 0\ =\ 0)\ (hY\ :\ Continuous\ Y)}
\leanline{4}{(T\ :\ ℝ)\ (hT\ :\ 0\ ≤\ T)}
\leanline{4}{(hYe\ :\ ∀\ t\ ∈\ Icc\ 0\ T,\ Y\ t\ -\ x\ =\ -(∫\ s\ in\ 0..t,\ g\ (Y\ s))\ +\ B\ t)}
\leanline{4}{(test\ generator\ :\ H\ →\ ℝ)\ (htest\ :\ UniformContinuous\ test)}
\leanline{4}{(hgenerator\ :\ UniformContinuous\ generator)}
\leanline{4}{(C\ :\ ℝ)\ (hC\ :\ 0\ ≤\ C)}
\leanline{4}{(hres\ :\ ∀\ ε\ >\ 0,\ ∃\ K\ :\ ℝ,\ 0\ ≤\ K\ ∧\ ∃\ N\ :\ ℕ,\ ∀\ n\ ≥\ N,}
\leanline{6}{|test\ (ULA.additiveEulerChain\ g\ (T\ /\ ((n\ +\ 1\ :\ ℕ)\ :\ ℝ))\ x\ B\ (n\ +\ 1))\ -}
\leanline{8}{test\ x\ -\ ∑\ k\ ∈\ Finset.range\ (n\ +\ 1),\ (T\ /\ ((n\ +\ 1\ :\ ℕ)\ :\ ℝ))\ *}
\leanline{10}{generator\ (ULA.additiveEulerChain\ g\ (T\ /\ ((n\ +\ 1\ :\ ℕ)\ :\ ℝ))\ x\ B\ k)|\ ≤}
\leanline{8}{C\ *\ (ε\ +\ K\ *\ Real.sqrt\ (T\ /\ ((n\ +\ 1\ :\ ℕ)\ :\ ℝ))))\ :}
\leanline{4}{test\ (Y\ T)\ -\ test\ x\ =\ ∫\ t\ in\ 0..T,\ generator\ (Y\ t)}
\end{leancode}
\leanlabel{How the Lean proof follows the argument}
\texttt{weighted\_\allowbreak{}dynkin\_\allowbreak{}of\_\allowbreak{}weak\_\allowbreak{}euler\_\allowbreak{}consistency} gives the formal counterpart of this step. The main intermediate facts are \texttt{r}, \texttt{hzero}, \texttt{hendpoint}, \texttt{hsum}, \texttt{hlimit}, \texttt{he} in that order. Named proof applications include \texttt{sub\_\allowbreak{}eq\_\allowbreak{}zero.\allowbreak{}mp}; the adjacent proof shows their hypotheses and order of use.
\end{leanbox}


Earlier results used here: \cite[Section~3.16.1]{LangevinLean2026ULASC}; \cite[Section~3.16.8]{LangevinLean2026ULASC}.
\section{\texorpdfstring{Bounded weighted observations are uniformly continuous on \(L^{2}\)}{Bounded weighted observations are uniformly continuous on L²}}
\label{proof-module-AverageLMC-LaplacianWeightedTest-lean}


The observation is the actual probability integral \(E\)[\(w\) \(F(X)]\). Its continuity is proved from boundedness and a second-moment tail estimate, allowing weak Euler limits to be taken on the \(L^{2}\) state space.



The actual weighted observation is


\[\mathcal O_{w,F}(X)=\int w(a)F(X(a))\,P(da).\]


If \(|w|\le A\), \(|F|\le C\), and \(|F(x)-F(y)|\le\varepsilon\) whenever \(\|x-y\|\le\delta\), then


\[\left|\mathcal O_{w,F}(X)-\mathcal O_{w,F}(Y)\right|
\le A\left(\varepsilon+\frac{2C}{\delta^2}\|X-Y\|_{L^2(P)}^2\right).\]

Earlier results used here: \cite[Section~3.13]{LangevinLean2026ULASC}.
\subsection{\texorpdfstring{Define and integrate the weighted test}{Define and integrate the weighted test}}
\label{report-step-AverageLMC-LaplacianWeightedTest-lean-1}


Define the weighted observation \(wF(X)\) using the actual representative of \(X\) and the measure \(P\). Measurability of \(w\) and continuity of \(F\) give measurability of the product; the bounds \(|w|\le A\) almost everywhere and \(|F|\le C\) dominate it by the integrable constant \(AC\).

\begin{leanbox}{Lean: \texttt{lpWeightedTest}, \texttt{lpWeightedTest\_\allowbreak{}integrable}}
\leansource{\texttt{lpWeightedTest} \quad AverageLMC/\allowbreak{}LaplacianWeightedTest.lean, lines 28–29}
\begin{leancode}
\leanline{0}{def\ lpWeightedTest\ (w\ :\ Ω\ →\ ℝ)\ (F\ :\ H\ →\ ℝ)\ (X\ :\ Lp\ H\ 2\ P)\ :\ ℝ\ :=}
\leanline{2}{∫\ a,\ w\ a\ *\ F\ (X\ a)\ ∂P}
\end{leancode}
\leansource{\texttt{lpWeightedTest\_\allowbreak{}integrable} \quad AverageLMC/\allowbreak{}LaplacianWeightedTest.lean, lines 31–42}
\begin{leancode}
\leanline{0}{theorem\ lpWeightedTest\_integrable}
\leanline{4}{(w\ :\ Ω\ →\ ℝ)\ (hw\ :\ Measurable\ w)\ (A\ :\ ℝ)\ (hA\ :\ 0\ ≤\ A)}
\leanline{4}{(hwA\ :\ ∀ᵐ\ a\ ∂P,\ |w\ a|\ ≤\ A)}
\leanline{4}{(F\ :\ H\ →\ ℝ)\ (hF\ :\ Continuous\ F)\ (C\ :\ ℝ)\ (hC\ :\ 0\ ≤\ C)}
\leanline{4}{(hFC\ :\ ∀\ x,\ |F\ x|\ ≤\ C)\ (X\ :\ Lp\ H\ 2\ P)\ :}
\leanline{4}{Integrable\ (fun\ a\ ↦\ w\ a\ *\ F\ (X\ a))\ P}
\end{leancode}
\leanlabel{How the Lean proof follows the argument}
\texttt{lp\allowbreak{}Weighted\allowbreak{}Test} names the mathematical object defined above. Its defining expression appears in the adjacent statement; subsequent identities may unfold this definition directly.




\texttt{lp\allowbreak{}Weighted\allowbreak{}Test\_\allowbreak{}integrable} gives the formal counterpart of this step. Named proof applications include \texttt{mul\_\allowbreak{}le\_\allowbreak{}mul}; the adjacent proof shows their hypotheses and order of use.
\end{leanbox}
\subsection{\texorpdfstring{Make the continuity estimate explicit}{Make the continuity estimate explicit}}
\label{report-step-AverageLMC-LaplacianWeightedTest-lean-2}


If \(F\) oscillates by at most \(\varepsilon\) at distance at most \(\delta\), then everywhere \(|F(x)-F(y)|\le \varepsilon +(2C/\delta ^{2})\|x-y\|^{2}\). Integrate against the bounded weight and identify the squared \(L^{2}\) norm with the squared-distance integral. The resulting bound is \(A[\varepsilon +(2C/\delta ^{2})\|X-Y\|^{2}]\).

\begin{leanbox}{Lean: \texttt{lpWeightedTest\_\allowbreak{}difference\_\allowbreak{}le}}
\leansource{\texttt{lpWeightedTest\_\allowbreak{}difference\_\allowbreak{}le} \quad AverageLMC/\allowbreak{}LaplacianWeightedTest.lean, lines 46–100}
\begin{leancode}
\leanline{0}{theorem\ lpWeightedTest\_difference\_le}
\leanline{4}{(w\ :\ Ω\ →\ ℝ)\ (hw\ :\ Measurable\ w)\ (A\ :\ ℝ)\ (hA\ :\ 0\ ≤\ A)}
\leanline{4}{(hwA\ :\ ∀ᵐ\ a\ ∂P,\ |w\ a|\ ≤\ A)}
\leanline{4}{(F\ :\ H\ →\ ℝ)\ (hF\ :\ Continuous\ F)\ (C\ ε\ δ\ :\ ℝ)}
\leanline{4}{(hC\ :\ 0\ ≤\ C)\ (hε\ :\ 0\ ≤\ ε)\ (hδ\ :\ 0\ <\ δ)}
\leanline{4}{(hFC\ :\ ∀\ x,\ |F\ x|\ ≤\ C)}
\leanline{4}{(hmod\ :\ ∀\ x\ y,\ ‖x\ -\ y‖\ ≤\ δ\ →\ |F\ x\ -\ F\ y|\ ≤\ ε)}
\leanline{4}{(X\ Y\ :\ Lp\ H\ 2\ P)\ :}
\leanline{4}{|lpWeightedTest\ P\ w\ F\ X\ -\ lpWeightedTest\ P\ w\ F\ Y|\ ≤}
\leanline{6}{A\ *\ (ε\ +\ (2\ *\ C\ /\ δ\ \textasciicircum{}\ 2)\ *\ ‖X\ -\ Y‖\ \textasciicircum{}\ 2)}
\end{leancode}
\leanlabel{How the Lean proof follows the argument}
\texttt{lp\allowbreak{}Weighted\allowbreak{}Test\_\allowbreak{}difference\_\allowbreak{}le} separates the cases \texttt{hxy : \ensuremath{\Vert}x -\allowbreak{} y\ensuremath{\Vert} \ensuremath{\le} \ensuremath{\delta}} and its negation before combining them. The main intermediate facts are \texttt{\ensuremath{\beta}}, \texttt{h\ensuremath{\beta}}, \texttt{hp}, \texttt{hwi\allowbreak{}X}, \texttt{hwi\allowbreak{}Y}, \texttt{hdiff} in that order. Named proof applications include \texttt{mul\_\allowbreak{}le\_\allowbreak{}mul}, \texttt{ae\_\allowbreak{}of\_\allowbreak{}all}, \texttt{abs\_\allowbreak{}integral\_\allowbreak{}le\_\allowbreak{}integral\_\allowbreak{}abs.\allowbreak{}trans}; the adjacent proof shows their hypotheses and order of use.
\end{leanbox}
\subsection{\texorpdfstring{Transfer uniform continuity from state tests}{Transfer uniform continuity from state tests}}
\label{report-step-AverageLMC-LaplacianWeightedTest-lean-3}


Choose a small test oscillation first, then a sufficiently small \(L^{2}\) radius to control the quadratic tail term. This proves uniform continuity of \(X\mapsto\mathbb E[wF(X)]\) for every bounded uniformly continuous \(F\).

\begin{leanbox}{Lean: \texttt{lpWeightedTest\_\allowbreak{}uniformContinuous}}
\leansource{\texttt{lpWeightedTest\_\allowbreak{}uniformContinuous} \quad AverageLMC/\allowbreak{}LaplacianWeightedTest.lean, lines 104–135}
\begin{leancode}
\leanline{0}{theorem\ lpWeightedTest\_uniformContinuous}
\leanline{4}{(w\ :\ Ω\ →\ ℝ)\ (hw\ :\ Measurable\ w)\ (A\ :\ ℝ)\ (hA\ :\ 0\ ≤\ A)}
\leanline{4}{(hwA\ :\ ∀ᵐ\ a\ ∂P,\ |w\ a|\ ≤\ A)}
\leanline{4}{(F\ :\ H\ →\ ℝ)\ (hF\ :\ UniformContinuous\ F)\ (C\ :\ ℝ)\ (hC\ :\ 0\ ≤\ C)}
\leanline{4}{(hFC\ :\ ∀\ x,\ |F\ x|\ ≤\ C)\ :\ UniformContinuous\ (lpWeightedTest\ P\ w\ F)}
\end{leancode}
\leanlabel{How the Lean proof follows the argument}
\texttt{lp\allowbreak{}Weighted\allowbreak{}Test\_\allowbreak{}uniform\allowbreak{}Continuous} gives the formal counterpart of this step. The main intermediate facts are \texttt{\ensuremath{\eta}}, \texttt{h\ensuremath{\eta}}, \texttt{h\ensuremath{\eta}scale}, \texttt{\ensuremath{\beta}}, \texttt{h\ensuremath{\beta}}, \texttt{r} in that order. Named proof applications include \texttt{Metric.\allowbreak{}uniform\allowbreak{}Continuous\_\allowbreak{}iff.\allowbreak{}mpr}; the adjacent proof shows their hypotheses and order of use.
\end{leanbox}
\section{\texorpdfstring{Regularity of compact \(C^{2}\) Langevin tests and generators}{Regularity of compact C² Langevin tests and generators}}
\label{proof-module-AverageLMC-LaplacianCompactTest-lean}


The drift is the paper's Lipschitz gradient, and tests have only two continuous derivatives. Both the test and its actual Langevin generator yield bounded uniformly continuous observations needed in the Euler limit.



For every compactly supported \(C^2\) test and Lipschitz score, its actual Langevin generator remains continuous and compactly supported:


\[\mathcal L_f\phi=\Delta\phi-D\phi[\nabla f],\qquad
\phi\in C_c^2\ \Longrightarrow\ \mathcal L_f\phi\in C_c^0.\]


Consequently both scalar observations admit global finite bounds and uniform continuity, and hence


\[X\longmapsto\mathbb E[w\phi(X)],\qquad
X\longmapsto\mathbb E[w\mathcal L_f\phi(X)]
\quad\text{are uniformly continuous on }L^2(P)
\quad (w\in L^\infty(P)).\]

Earlier results used here: \cite[Section~3.46]{LangevinLean2026ULASC}.
\subsection{\texorpdfstring{Control the actual generator}{Control the actual generator}}
\label{report-step-AverageLMC-LaplacianCompactTest-lean-1}


Continuity of the Hessian makes the Laplacian continuous. The drift derivative term is continuous by continuity of the derivative and the Lipschitz gradient. Derivatives of compactly supported tests vanish off their support, so the generator also has compact support.

\begin{leanbox}{Lean: \texttt{continuous\_\allowbreak{}langevinGenerator\_\allowbreak{}compact\_\allowbreak{}test}, \texttt{hasCompactSupport\_\allowbreak{}langevinGenerator\_\allowbreak{}compact\_\allowbreak{}test}}
\leansource{\texttt{continuous\_\allowbreak{}langevinGenerator\_\allowbreak{}compact\_\allowbreak{}test} \quad AverageLMC/\allowbreak{}LaplacianCompactTest.lean, lines 23–27}
\begin{leancode}
\leanline{0}{theorem\ continuous\_langevinGenerator\_compact\_test}
\leanline{4}{(f\ :\ H\ →\ ℝ)\ (L\ :\ ℝ≥0)\ (hLip\ :\ LipschitzWith\ L\ (gradient\ f))}
\leanline{4}{(φ\ :\ H\ →\ ℝ)\ (hφ\ :\ ContDiff\ ℝ\ 2\ φ)\ :\ Continuous\ (ULA.langevinGenerator\ f\ φ)}
\end{leancode}
\leansource{\texttt{hasCompactSupport\_\allowbreak{}langevinGenerator\_\allowbreak{}compact\_\allowbreak{}test} \quad AverageLMC/\allowbreak{}LaplacianCompactTest.lean, lines 29–49}
\begin{leancode}
\leanline{0}{theorem\ hasCompactSupport\_langevinGenerator\_compact\_test}
\leanline{4}{(f\ φ\ :\ H\ →\ ℝ)\ (hc\ :\ HasCompactSupport\ φ)\ :}
\leanline{4}{HasCompactSupport\ (ULA.langevinGenerator\ f\ φ)}
\end{leancode}
\leanlabel{How the Lean proof follows the argument}
\texttt{continuous\_\allowbreak{}langevin\allowbreak{}Generator\_\allowbreak{}compact\_\allowbreak{}test} gives the formal counterpart of this step. Named proof applications include \texttt{continuous\_\allowbreak{}finset\allowbreak{}Sum}; the adjacent proof shows their hypotheses and order of use.




\texttt{has\allowbreak{}Compact\allowbreak{}Support\_\allowbreak{}langevin\allowbreak{}Generator\_\allowbreak{}compact\_\allowbreak{}test} proceeds by \texttt{induction s using Finset.\allowbreak{}induction\_\allowbreak{}on with}. The base cases establish the initial claims and the inductive branches use the induction hypotheses. The main intermediate facts are \texttt{hs} in that order. Named proof applications include \texttt{hs}; the adjacent proof shows their hypotheses and order of use.
\end{leanbox}


Earlier results used here: \cite[Section~3.46.1]{LangevinLean2026ULASC}; \cite[Section~3.46.2]{LangevinLean2026ULASC}; \cite[Section~3.46.8]{LangevinLean2026ULASC}; \cite[Section~3.46.9]{LangevinLean2026ULASC}.
\subsection{\texorpdfstring{Obtain global bounds and uniform continuity}{Obtain global bounds and uniform continuity}}
\label{report-step-AverageLMC-LaplacianCompactTest-lean-2}


A continuous compactly supported test is bounded and uniformly continuous. Apply the same facts to its generator, using the proved generator continuity and compact support. No extra regularity of the potential is introduced.

\begin{leanbox}{Lean: \texttt{compact\_\allowbreak{}test\_\allowbreak{}bounded\_\allowbreak{}uniformContinuous}, \texttt{compact\_\allowbreak{}test\_\allowbreak{}generator\_\allowbreak{}bounded\_\allowbreak{}uniformContinuous}}
\leansource{\texttt{compact\_\allowbreak{}test\_\allowbreak{}bounded\_\allowbreak{}uniformContinuous} \quad AverageLMC/\allowbreak{}LaplacianCompactTest.lean, lines 51–57}
\begin{leancode}
\leanline{0}{theorem\ compact\_test\_bounded\_uniformContinuous}
\leanline{4}{(φ\ :\ H\ →\ ℝ)\ (hφ\ :\ ContDiff\ ℝ\ 2\ φ)\ (hc\ :\ HasCompactSupport\ φ)\ :}
\leanline{4}{UniformContinuous\ φ\ ∧\ ∃\ C\ :\ ℝ,\ 0\ ≤\ C\ ∧\ ∀\ x,\ |φ\ x|\ ≤\ C}
\end{leancode}
\leansource{\texttt{compact\_\allowbreak{}test\_\allowbreak{}generator\_\allowbreak{}bounded\_\allowbreak{}uniformContinuous} \quad AverageLMC/\allowbreak{}LaplacianCompactTest.lean, lines 59–70}
\begin{leancode}
\leanline{0}{theorem\ compact\_test\_generator\_bounded\_uniformContinuous}
\leanline{4}{(f\ :\ H\ →\ ℝ)\ (L\ :\ ℝ≥0)\ (hLip\ :\ LipschitzWith\ L\ (gradient\ f))}
\leanline{4}{(φ\ :\ H\ →\ ℝ)\ (hφ\ :\ ContDiff\ ℝ\ 2\ φ)\ (hc\ :\ HasCompactSupport\ φ)\ :}
\leanline{4}{UniformContinuous\ (ULA.langevinGenerator\ f\ φ)\ ∧}
\leanline{6}{∃\ C\ :\ ℝ,\ 0\ ≤\ C\ ∧\ ∀\ x,\ |ULA.langevinGenerator\ f\ φ\ x|\ ≤\ C}
\end{leancode}
\leanlabel{How the Lean proof follows the argument}
\texttt{compact\_\allowbreak{}test\_\allowbreak{}bounded\_\allowbreak{}uniform\allowbreak{}Continuous} gives the formal counterpart of this step. Its remaining equalities are discharged by the simplifications displayed in the proof.




\texttt{compact\_\allowbreak{}test\_\allowbreak{}generator\_\allowbreak{}bounded\_\allowbreak{}uniform\allowbreak{}Continuous} gives the formal counterpart of this step. The main intermediate facts are \texttt{hcont}, \texttt{hcomp} in that order. Named proof applications include \texttt{h\allowbreak{}C}; the adjacent proof shows their hypotheses and order of use.
\end{leanbox}


Earlier results used here: \cite[Section~3.46.2]{LangevinLean2026ULASC}.
\subsection{\texorpdfstring{Build the \(L^{2}\) observation interface}{Build the L² observation interface}}
\label{report-step-AverageLMC-LaplacianCompactTest-lean-3}


Apply the weighted-test uniform continuity theorem to a compact \(C^{2}\) test and its generator. These are the precise observation functions appearing in the abstract matching-noise Euler limit.

\begin{leanbox}{Lean: \texttt{lpWeightedTest\_\allowbreak{}compact\_\allowbreak{}test\_\allowbreak{}uniformContinuous}}
\leansource{\texttt{lpWeightedTest\_\allowbreak{}compact\_\allowbreak{}test\_\allowbreak{}uniformContinuous} \quad AverageLMC/\allowbreak{}LaplacianCompactTest.lean, lines 74–86}
\begin{leancode}
\leanline{0}{theorem\ lpWeightedTest\_compact\_test\_uniformContinuous}
\leanline{4}{\{Ω\ :\ Type*\}\ [MeasurableSpace\ Ω]\ (P\ :\ Measure\ Ω)\ [IsProbabilityMeasure\ P]}
\leanline{4}{(w\ :\ Ω\ →\ ℝ)\ (hw\ :\ Measurable\ w)\ (A\ :\ ℝ)\ (hA\ :\ 0\ ≤\ A)}
\leanline{4}{(hwA\ :\ ∀ᵐ\ a\ ∂P,\ |w\ a|\ ≤\ A)}
\leanline{4}{(f\ :\ H\ →\ ℝ)\ (L\ :\ ℝ≥0)\ (hLip\ :\ LipschitzWith\ L\ (gradient\ f))}
\leanline{4}{(φ\ :\ H\ →\ ℝ)\ (hφ\ :\ ContDiff\ ℝ\ 2\ φ)\ (hc\ :\ HasCompactSupport\ φ)\ :}
\leanline{4}{UniformContinuous\ (lpWeightedTest\ P\ w\ φ)\ ∧}
\leanline{6}{UniformContinuous\ (lpWeightedTest\ P\ w\ (ULA.langevinGenerator\ f\ φ))}
\end{leancode}
\leanlabel{How the Lean proof follows the argument}
\texttt{lp\allowbreak{}Weighted\allowbreak{}Test\_\allowbreak{}compact\_\allowbreak{}test\_\allowbreak{}uniform\allowbreak{}Continuous} gives the formal counterpart of this step. Named proof applications include \texttt{\ensuremath{\langle}lp\allowbreak{}Weighted\allowbreak{}Test\_\allowbreak{}uniform\allowbreak{}Continuous}; the adjacent proof shows their hypotheses and order of use.
\end{leanbox}


Earlier results used here: \cite[Section~3.46.2]{LangevinLean2026ULASC}.
\section{\texorpdfstring{Exact Gaussian averaging of the quadratic Taylor polynomial}{Exact Gaussian averaging of the quadratic Taylor polynomial}}
\label{proof-module-AverageLMC-LaplacianGaussianTaylor-lean}


All expectations here use the actual standard Gaussian measure in a finite-dimensional real Hilbert space. The quadratic trace identity identifies the diffusion coefficient in the Langevin generator.



For an actual standard Gaussian vector and any continuous bilinear form,


\[\mathbb E[\langle u,Z\rangle\langle v,Z\rangle]=\langle u,v\rangle,
\qquad \mathbb E[B(Z,Z)]=\sum_i B(e_i,e_i),\qquad Z\sim N(0,I).\]


Writing the quadratic Taylor polynomial as \(\mathcal T_2\phi(x;w)=\phi(x)+D\phi(x)[w]+\tfrac12D^2\phi(x)[w,w]\), its Gaussian expectation is exactly


\[\mathbb E\,\mathcal T_2\phi\bigl(x;hv+\sqrt{2h}Z\bigr)
=\phi(x)+h\bigl(D\phi(x)[v]+\Delta\phi(x)\bigr)
+\frac{h^2}{2}D^2\phi(x)[v,v].\]

Earlier results used here: \cite[Section~3.55]{LangevinLean2026ULASC}; \cite[Section~3.37]{LangevinLean2026ULASC}.
\subsection{\texorpdfstring{Use the true covariance}{Use the true covariance}}
\label{report-step-AverageLMC-LaplacianGaussianTaylor-lean-1}


Apply the Gaussian covariance bilinear form to two linear observations. The standard Gaussian has mean zero and identity covariance, so the expectation of \ensuremath{\langle}\(u\),\(Z\)\ensuremath{\rangle}\ensuremath{\langle}\(v\),\(Z\)\ensuremath{\rangle} equals \ensuremath{\langle}\(u\),\(v\)\ensuremath{\rangle}. Square integrability justifies the product integral.

\begin{leanbox}{Lean: \texttt{integral\_\allowbreak{}gaussian\_\allowbreak{}inner\_\allowbreak{}mul}}
\leansource{\texttt{integral\_\allowbreak{}gaussian\_\allowbreak{}inner\_\allowbreak{}mul} \quad AverageLMC/\allowbreak{}LaplacianGaussianTaylor.lean, lines 27–34}
\begin{leancode}
\leanline{0}{theorem\ integral\_gaussian\_inner\_mul\ (u\ v\ :\ H)\ :}
\leanline{4}{(∫\ z\ :\ H,\ inner\ ℝ\ u\ z\ *\ inner\ ℝ\ v\ z\ ∂stdGaussian\ H)\ =\ inner\ ℝ\ u\ v}
\end{leancode}
\leanlabel{How the Lean proof follows the argument}
\texttt{integral\_\allowbreak{}gaussian\_\allowbreak{}inner\_\allowbreak{}mul} gives the formal counterpart of this step. The main intermediate facts are \texttt{hz}, \texttt{hc} in that order. Named proof applications include \texttt{hc.\allowbreak{}symm}; the adjacent proof shows their hypotheses and order of use.
\end{leanbox}
\subsection{\texorpdfstring{Turn a quadratic form into its trace}{Turn a quadratic form into its trace}}
\label{report-step-AverageLMC-LaplacianGaussianTaylor-lean-2}


Expand both arguments of a continuous bilinear form in the standard orthonormal basis. Interchange finite sums and integrals using Gaussian second moments. The covariance identity kills off-diagonal terms, leaving the sum of the diagonal values.

\begin{leanbox}{Lean: \texttt{integral\_\allowbreak{}gaussian\_\allowbreak{}bilinear}}
\leansource{\texttt{integral\_\allowbreak{}gaussian\_\allowbreak{}bilinear} \quad AverageLMC/\allowbreak{}LaplacianGaussianTaylor.lean, lines 37–73}
\begin{leancode}
\leanline{0}{theorem\ integral\_gaussian\_bilinear\ (B\ :\ H\ →L[ℝ]\ H\ →L[ℝ]\ ℝ)\ :}
\leanline{4}{(∫\ z\ :\ H,\ B\ z\ z\ ∂stdGaussian\ H)\ =}
\leanline{6}{∑\ i\ :\ Fin\ (Module.finrank\ ℝ\ H),}
\leanline{8}{B\ (stdOrthonormalBasis\ ℝ\ H\ i)\ (stdOrthonormalBasis\ ℝ\ H\ i)}
\end{leancode}
\leanlabel{How the Lean proof follows the argument}
\texttt{integral\_\allowbreak{}gaussian\_\allowbreak{}bilinear} gives the formal counterpart of this step. The main intermediate facts are \texttt{b}, \texttt{hz}, \texttt{hi}, \texttt{he}, \texttt{hefun} in that order. Named proof applications include \texttt{Finset.\allowbreak{}sum\_\allowbreak{}congr}; the adjacent proof shows their hypotheses and order of use.
\end{leanbox}
\subsection{\texorpdfstring{Average the entire Euler Taylor polynomial}{Average the entire Euler Taylor polynomial}}
\label{report-step-AverageLMC-LaplacianGaussianTaylor-lean-3}


Expand the increment \(hv+\sqrt{2h}Z\) in the linear and quadratic derivative terms. Gaussian centering removes its linear-noise and drift-noise terms; the trace is the actual Laplacian. The remaining deterministic quadratic drift is \(h^{2}D^{2}\varphi [v,v]/2\). This is an exact equality for the Taylor polynomial, prior to any remainder estimate.

\begin{leanbox}{Lean: \texttt{integral\_\allowbreak{}gaussian\_\allowbreak{}secondOrderExpansion}}
\leansource{\texttt{integral\_\allowbreak{}gaussian\_\allowbreak{}secondOrderExpansion} \quad AverageLMC/\allowbreak{}LaplacianGaussianTaylor.lean, lines 77–144}
\begin{leancode}
\leanline{0}{theorem\ integral\_gaussian\_secondOrderExpansion}
\leanline{4}{(φ\ :\ H\ →\ ℝ)\ (x\ v\ :\ H)\ (h\ :\ ℝ)\ (hh\ :\ 0\ ≤\ h)\ :}
\leanline{4}{(∫\ z\ :\ H,\ ULA.secondOrderExpansion\ φ\ x\ (h\ •\ v\ +\ Real.sqrt\ (2\ *\ h)\ •\ z)}
\leanline{6}{∂stdGaussian\ H)\ =}
\leanline{6}{φ\ x\ +\ h\ *\ (fderiv\ ℝ\ φ\ x\ v\ +\ Laplacian.laplacian\ φ\ x)\ +}
\leanline{8}{h\ \textasciicircum{}\ 2\ /\ 2\ *\ iteratedFDeriv\ ℝ\ 2\ φ\ x\ (fun\ \_\ ↦\ v)}
\end{leancode}
\leanlabel{How the Lean proof follows the argument}
\texttt{integral\_\allowbreak{}gaussian\_\allowbreak{}second\allowbreak{}Order\allowbreak{}Expansion} gives the formal counterpart of this step. The main intermediate facts are \texttt{D}, \texttt{B}, \texttt{s}, \texttt{hz}, \texttt{h\allowbreak{}D}, \texttt{h\allowbreak{}Bv} in that order. Named proof applications include \texttt{hquad.\allowbreak{}integrable}, \texttt{integral\_\allowbreak{}strong\allowbreak{}Dual\_\allowbreak{}std\allowbreak{}Gaussian}; the adjacent proof shows their hypotheses and order of use.
\end{leanbox}


Earlier results used here: \cite[Section~3.55.1]{LangevinLean2026ULASC}.
\section{\texorpdfstring{Uniform weak consistency for compact smooth tests}{Uniform weak consistency for compact smooth tests}}
\label{proof-module-AverageLMC-LaplacianWeakStep-lean}


This smooth-test development supplies the Gaussian calculus and drift-noise separation used by the later \(C^{2}\) extension. Smoothness is imposed only on auxiliary compact tests; the final potential keeps the paper's \(C^{3}\) hypothesis.



Drift and noise are separated by the deterministic estimate


\[\left|\phi(x+w+d)-\phi(x+w)-D\phi(x)[d]\right|
\le\frac L2\|d\|^2+L\|w\|\,\|d\|.\]


For every compact smooth test there is a constant \(K\ge0\), uniform in \(x,v\) and \(0\le h\le1\), such that


\[\left|\mathbb E\phi\bigl(x+hv+\sqrt{2h}Z\bigr)-\phi(x)
-h\bigl(D\phi(x)[v]+\Delta\phi(x)\bigr)\right|
\le Kh\sqrt h\,(1+\|v\|^2).\]

Earlier results used here: \cite[Section~3.9]{LangevinLean2026ULASC}; \cite[Section~3.46]{LangevinLean2026ULASC}.
\subsection{\texorpdfstring{Bound the test gradient globally}{Bound the test gradient globally}}
\label{report-step-AverageLMC-LaplacianWeakStep-lean-1}


Compact support bounds the second derivative everywhere. The derivative norm bound gives a Lipschitz derivative, and the Riesz isometry transfers this constant to the gradient.

\begin{leanbox}{Lean: \texttt{compactSmooth\_\allowbreak{}gradient\_\allowbreak{}lipschitz}}
\leansource{\texttt{compactSmooth\_\allowbreak{}gradient\_\allowbreak{}lipschitz} \quad AverageLMC/\allowbreak{}LaplacianWeakStep.lean, lines 28–42}
\begin{leancode}
\leanline{0}{theorem\ compactSmooth\_gradient\_lipschitz}
\leanline{4}{(φ\ :\ H\ →\ ℝ)\ (hφ\ :\ ContDiff\ ℝ\ ∞\ φ)\ (hc\ :\ HasCompactSupport\ φ)\ :}
\leanline{4}{∃\ L\ :\ ℝ≥0,\ LipschitzWith\ L\ (gradient\ φ)}
\end{leancode}
\leanlabel{How the Lean proof follows the argument}
\texttt{compact\allowbreak{}Smooth\_\allowbreak{}gradient\_\allowbreak{}lipschitz} gives the formal counterpart of this step. The main intermediate facts are \texttt{L}, \texttt{h\allowbreak{}Ddiff}, \texttt{h\allowbreak{}Dlip} in that order. Named proof applications include \texttt{lipschitz\allowbreak{}With\_\allowbreak{}of\_\allowbreak{}nnnorm\_\allowbreak{}fderiv\_\allowbreak{}le}, \texttt{hb}, \texttt{h\allowbreak{}Dlip.\allowbreak{}dist\_\allowbreak{}le\_\allowbreak{}mul}; the adjacent proof shows their hypotheses and order of use.
\end{leanbox}
\subsection{\texorpdfstring{Control a deterministic drift increment}{Control a deterministic drift increment}}
\label{report-step-AverageLMC-LaplacianWeakStep-lean-2}


Apply the smooth upper bound to \(\varphi\) and \(-\varphi\) to obtain an absolute first-order remainder bounded by \(L\|d\|^{2}/2\). Apply it at \(x+w\) and compare the two derivatives by gradient Lipschitz continuity, obtaining the additional \(L\|w\|\|d\|\) term. This separates a potentially large random drift from the small Gaussian increment.

\begin{leanbox}{Lean: \texttt{firstOrder\_\allowbreak{}remainder\_\allowbreak{}abs\_\allowbreak{}le}, \texttt{shifted\_\allowbreak{}firstOrder\_\allowbreak{}remainder\_\allowbreak{}abs\_\allowbreak{}le}}
\leansource{\texttt{firstOrder\_\allowbreak{}remainder\_\allowbreak{}abs\_\allowbreak{}le} \quad AverageLMC/\allowbreak{}LaplacianWeakStep.lean, lines 46–61}
\begin{leancode}
\leanline{0}{theorem\ firstOrder\_remainder\_abs\_le}
\leanline{4}{(φ\ :\ H\ →\ ℝ)\ (L\ :\ ℝ≥0)\ (hdiff\ :\ Differentiable\ ℝ\ φ)}
\leanline{4}{(hLip\ :\ LipschitzWith\ L\ (gradient\ φ))\ (x\ d\ :\ H)\ :}
\leanline{4}{|φ\ (x\ +\ d)\ -\ φ\ x\ -\ fderiv\ ℝ\ φ\ x\ d|\ ≤\ (L\ :\ ℝ)\ /\ 2\ *\ ‖d‖\ \textasciicircum{}\ 2}
\end{leancode}
\leansource{\texttt{shifted\_\allowbreak{}firstOrder\_\allowbreak{}remainder\_\allowbreak{}abs\_\allowbreak{}le} \quad AverageLMC/\allowbreak{}LaplacianWeakStep.lean, lines 65–81}
\begin{leancode}
\leanline{0}{theorem\ shifted\_firstOrder\_remainder\_abs\_le}
\leanline{4}{(φ\ :\ H\ →\ ℝ)\ (L\ :\ ℝ≥0)\ (hdiff\ :\ Differentiable\ ℝ\ φ)}
\leanline{4}{(hLip\ :\ LipschitzWith\ L\ (gradient\ φ))\ (x\ w\ d\ :\ H)\ :}
\leanline{4}{|φ\ (x\ +\ w\ +\ d)\ -\ φ\ (x\ +\ w)\ -\ fderiv\ ℝ\ φ\ x\ d|\ ≤}
\leanline{6}{(L\ :\ ℝ)\ /\ 2\ *\ ‖d‖\ \textasciicircum{}\ 2\ +\ (L\ :\ ℝ)\ *\ ‖w‖\ *\ ‖d‖}
\end{leancode}
\leanlabel{How the Lean proof follows the argument}
\texttt{first\allowbreak{}Order\_\allowbreak{}remainder\_\allowbreak{}abs\_\allowbreak{}le} gives the formal counterpart of this step. The main intermediate facts are \texttt{hneg}, \texttt{h\allowbreak{}Lip\allowbreak{}Neg}, \texttt{hu}, \texttt{hl} in that order. Named proof applications include \texttt{h\allowbreak{}Lip.\allowbreak{}neg}, \texttt{abs\_\allowbreak{}le.\allowbreak{}mpr}; the adjacent proof shows their hypotheses and order of use.




\texttt{shifted\_\allowbreak{}first\allowbreak{}Order\_\allowbreak{}remainder\_\allowbreak{}abs\_\allowbreak{}le} gives the formal counterpart of this step. The main intermediate facts are \texttt{hr}, \texttt{h\allowbreak{}D}, \texttt{heq} in that order. Named proof applications include \texttt{mul\_\allowbreak{}le\_\allowbreak{}mul\_\allowbreak{}of\_\allowbreak{}nonneg\_\allowbreak{}right}, \texttt{h\allowbreak{}Lip.\allowbreak{}norm\_\allowbreak{}sub\_\allowbreak{}le}; the adjacent proof shows their hypotheses and order of use.
\end{leanbox}


Earlier results used here: \cite[Section~3.9.1]{LangevinLean2026ULASC}.
\subsection{\texorpdfstring{Average the true Gaussian update}{Average the true Gaussian update}}
\label{report-step-AverageLMC-LaplacianWeakStep-lean-3}


Use the compact smooth Taylor remainder for the noise alone and the shifted first-order estimate for the drift. Gaussian first and third moments are finite. For \(0\le h\le 1\), absorb the deterministic \(h^{2}\) term into \(h\sqrt h\) and obtain one uniform constant \(K\) times \(h\sqrt h(1+\|v\|^{2})\), for arbitrary starting point and drift.

\begin{leanbox}{Lean: \texttt{compactSmooth\_\allowbreak{}gaussian\_\allowbreak{}weak\_\allowbreak{}step}}
\leansource{\texttt{compactSmooth\_\allowbreak{}gaussian\_\allowbreak{}weak\_\allowbreak{}step} \quad AverageLMC/\allowbreak{}LaplacianWeakStep.lean, lines 87–217}
\begin{leancode}
\leanline{0}{theorem\ compactSmooth\_gaussian\_weak\_step}
\leanline{4}{(φ\ :\ H\ →\ ℝ)\ (hφ\ :\ ContDiff\ ℝ\ ∞\ φ)\ (hc\ :\ HasCompactSupport\ φ)\ :}
\leanline{4}{∃\ K\ :\ ℝ,\ 0\ ≤\ K\ ∧\ ∀\ (x\ v\ :\ H)\ (h\ :\ ℝ),\ 0\ ≤\ h\ →\ h\ ≤\ 1\ →}
\leanline{6}{|(∫\ z\ :\ H,\ φ\ (x\ +\ h\ •\ v\ +\ Real.sqrt\ (2\ *\ h)\ •\ z)\ ∂stdGaussian\ H)\ -}
\leanline{8}{φ\ x\ -\ h\ *\ (fderiv\ ℝ\ φ\ x\ v\ +\ Laplacian.laplacian\ φ\ x)|\ ≤}
\leanline{10}{K\ *\ h\ *\ Real.sqrt\ h\ *\ (1\ +\ ‖v‖\ \textasciicircum{}\ 2)}
\end{leancode}
\leanlabel{How the Lean proof follows the argument}
\texttt{compact\allowbreak{}Smooth\_\allowbreak{}gaussian\_\allowbreak{}weak\_\allowbreak{}step} gives the formal counterpart of this step. The main intermediate facts are \texttt{h\ensuremath{\varphi}bound}, \texttt{hz1}, \texttt{hz3}, \texttt{G\ensuremath{{}_{1}}}, \texttt{G\ensuremath{{}_{3}}}, \texttt{K} in that order. Named proof applications include \texttt{hnoise.\allowbreak{}sub}, \texttt{abs\_\allowbreak{}integral\_\allowbreak{}le\_\allowbreak{}integral\_\allowbreak{}abs.\allowbreak{}trans}, \texttt{htotal.\allowbreak{}trans}; the adjacent proof shows their hypotheses and order of use.
\end{leanbox}


Earlier results used here: \cite[Section~3.55.1]{LangevinLean2026ULASC}; \cite[Section~3.55.3]{LangevinLean2026ULASC}.
\subsection{\texorpdfstring{Identify regularity of the diffusion term}{Identify regularity of the diffusion term}}
\label{report-step-AverageLMC-LaplacianWeakStep-lean-4}


Express the actual Laplacian as a finite sum of second derivative evaluations on an orthonormal basis. Continuity of the second derivative proves continuity of this sum.

\begin{leanbox}{Lean: \texttt{compactSmooth\_\allowbreak{}continuous\_\allowbreak{}laplacian}}
\leansource{\texttt{compactSmooth\_\allowbreak{}continuous\_\allowbreak{}laplacian} \quad AverageLMC/\allowbreak{}LaplacianWeakStep.lean, lines 221–234}
\begin{leancode}
\leanline{0}{theorem\ compactSmooth\_continuous\_laplacian}
\leanline{4}{(φ\ :\ H\ →\ ℝ)\ (hφ\ :\ ContDiff\ ℝ\ ∞\ φ)\ :\ Continuous\ (Laplacian.laplacian\ φ)}
\end{leancode}
\leanlabel{How the Lean proof follows the argument}
\texttt{compact\allowbreak{}Smooth\_\allowbreak{}continuous\_\allowbreak{}laplacian} gives the formal counterpart of this step. The main intermediate facts are \texttt{h\allowbreak{}DD}, \texttt{he} in that order. Named proof applications include \texttt{continuous\_\allowbreak{}finset\_\allowbreak{}sum}; the adjacent proof shows their hypotheses and order of use.
\end{leanbox}
\subsection{\texorpdfstring{Insert an actual bounded past weight}{Insert an actual bounded past weight}}
\label{report-step-AverageLMC-LaplacianWeakStep-lean-5}


Integrate the pointwise Gaussian weak residual against a measurable weight of absolute value at most \(A\). The drift needs only its true second moment, producing \(A\) \(K\) \(h\sqrt h(1+E\|v\|^{2})\). Product measurability and this integrable bound justify the averaging. The corresponding \(C^{2}\) estimate refines this bound to the regularity needed by the gradient cutoffs.

\begin{leanbox}{Lean: \texttt{compactSmooth\_\allowbreak{}gaussian\_\allowbreak{}weighted\_\allowbreak{}weak\_\allowbreak{}step}}
\leansource{\texttt{compactSmooth\_\allowbreak{}gaussian\_\allowbreak{}weighted\_\allowbreak{}weak\_\allowbreak{}step} \quad AverageLMC/\allowbreak{}LaplacianWeakStep.lean, lines 240–295}
\begin{leancode}
\leanline{0}{theorem\ compactSmooth\_gaussian\_weighted\_weak\_step}
\leanline{4}{\{Ω\ :\ Type*\}\ [MeasurableSpace\ Ω]\ (P\ :\ Measure\ Ω)\ [IsProbabilityMeasure\ P]}
\leanline{4}{(φ\ :\ H\ →\ ℝ)\ (hφ\ :\ ContDiff\ ℝ\ ∞\ φ)\ (hc\ :\ HasCompactSupport\ φ)}
\leanline{4}{(x\ v\ :\ Ω\ →\ H)\ (w\ :\ Ω\ →\ ℝ)\ (A\ :\ ℝ)}
\leanline{4}{(hx\ :\ Measurable\ x)\ (hv\ :\ Measurable\ v)\ (hw\ :\ Measurable\ w)}
\leanline{4}{(hv₂\ :\ MemLp\ v\ 2\ P)\ (hA\ :\ 0\ ≤\ A)\ (hwA\ :\ ∀\ a,\ |w\ a|\ ≤\ A)\ :}
\leanline{4}{∃\ K\ :\ ℝ,\ 0\ ≤\ K\ ∧\ ∀\ h\ :\ ℝ,\ 0\ ≤\ h\ →\ h\ ≤\ 1\ →}
\leanline{6}{|∫\ a,\ w\ a\ *\ ((∫\ z\ :\ H,}
\leanline{8}{φ\ (x\ a\ +\ h\ •\ v\ a\ +\ Real.sqrt\ (2\ *\ h)\ •\ z)\ ∂stdGaussian\ H)\ -}
\leanline{10}{φ\ (x\ a)\ -\ h\ *\ (fderiv\ ℝ\ φ\ (x\ a)\ (v\ a)\ +}
\leanline{12}{Laplacian.laplacian\ φ\ (x\ a)))\ ∂P|\ ≤}
\leanline{8}{A\ *\ K\ *\ h\ *\ Real.sqrt\ h\ *\ (1\ +\ ∫\ a,\ ‖v\ a‖\ \textasciicircum{}\ 2\ ∂P)}
\end{leancode}
\leanlabel{How the Lean proof follows the argument}
\texttt{compact\allowbreak{}Smooth\_\allowbreak{}gaussian\_\allowbreak{}weighted\_\allowbreak{}weak\_\allowbreak{}step} gives the formal counterpart of this step. The main intermediate facts are \texttt{r}, \texttt{h\ensuremath{\varphi}m}, \texttt{hjoint}, \texttt{hint}, \texttt{h\allowbreak{}D}, \texttt{h\allowbreak{}Lap} in that order. Named proof applications include \texttt{Measurable.\allowbreak{}aestrongly\allowbreak{}Measurable}, \texttt{abs\_\allowbreak{}integral\_\allowbreak{}le\_\allowbreak{}integral\_\allowbreak{}abs.\allowbreak{}trans}; the adjacent proof shows their hypotheses and order of use.
\end{leanbox}
\section{\texorpdfstring{Weak Gaussian Taylor expansion with exactly two derivatives}{Weak Gaussian Taylor expansion with exactly two derivatives}}
\label{proof-module-AverageLMC-CompactC2Taylor-lean}


The potential in the Laplacian theorem (Theorem 2, source th:3) is \(C^{3}\), so its gradient is only \(C^{2}\). These estimates establish the necessary compact-test weak expansion with precisely \(C^{2}\) regularity. A uniformly continuous compact Hessian supplies a freely small quadratic remainder plus a cubic bound; a third derivative of the test is never assumed. The Gaussian averaging requires only a second moment of the random drift.

\subsection{\texorpdfstring{Express the remainder using a genuine Hessian difference}{Express the remainder using a genuine Hessian difference}}
\label{report-step-AverageLMC-CompactC2Taylor-lean-1}


Restrict the test to the line \(x+t\) \(v\). The ordinary Taylor theorem with first-order polynomial and second derivative remainder gives a point \(u\) strictly between zero and one. The first derivative on this line is the Frechet derivative applied to \(v\); the second derivative is the actual second Frechet derivative applied twice to \(v\). Subtract the quadratic Taylor term at \(x\). The remaining error is exactly one half of the difference between the Hessians at \(x+u\) \(v\) and \(x\), evaluated twice on \(v\).

\begin{leanbox}{Lean: \texttt{secondOrderExpansion\_\allowbreak{}remainder\_\allowbreak{}eq\_\allowbreak{}second\_\allowbreak{}difference}}
\leansource{\texttt{secondOrderExpansion\_\allowbreak{}remainder\_\allowbreak{}eq\_\allowbreak{}second\_\allowbreak{}difference} \quad AverageLMC/\allowbreak{}CompactC2Taylor.lean, lines 17–56}
\begin{leancode}
\leanline{0}{theorem\ secondOrderExpansion\_remainder\_eq\_second\_difference}
\leanline{4}{(φ\ :\ H\ →\ ℝ)\ (hφ\ :\ ContDiff\ ℝ\ 2\ φ)\ (x\ v\ :\ H)\ :}
\leanline{4}{∃\ u\ ∈\ Ioo\ (0\ :\ ℝ)\ 1,}
\leanline{6}{φ\ (x\ +\ v)\ -\ ULA.secondOrderExpansion\ φ\ x\ v\ =}
\leanline{8}{(1\ /\ 2\ :\ ℝ)\ *\ ((iteratedFDeriv\ ℝ\ 2\ φ\ (x\ +\ u\ •\ v)\ -}
\leanline{10}{iteratedFDeriv\ ℝ\ 2\ φ\ x)\ (fun\ \_\ ↦\ v))}
\end{leancode}
\leanlabel{How the Lean proof follows the argument}
\texttt{second\allowbreak{}Order\allowbreak{}Expansion\_\allowbreak{}remainder\_\allowbreak{}eq\_\allowbreak{}second\_\allowbreak{}difference} gives the formal counterpart of this step. The main intermediate facts are \texttt{line}, \texttt{q}, \texttt{hline}, \texttt{hq}, \texttt{hline\allowbreak{}D}, \texttt{hline2} in that order. Named proof applications include \texttt{deriv\_\allowbreak{}const}; the adjacent proof shows their hypotheses and order of use.
\end{leanbox}


Earlier results used here: \cite[Section~3.55.1]{LangevinLean2026ULASC}.
\subsection{\texorpdfstring{Use uniform Hessian continuity without a third derivative}{Use uniform Hessian continuity without a third derivative}}
\label{report-step-AverageLMC-CompactC2Taylor-lean-2}


A compactly supported \(C^{2}\) test has a continuous compactly supported Hessian, hence a uniformly continuous and bounded Hessian. Given \(\eta\) positive, choose \(\delta\) so that Hessian differences at distance less than \(\delta\) are below \(2\eta\). If the displacement norm is below \(\delta\), the segment remainder is at most \(\eta\) times its squared norm. For larger displacements, the uniform Hessian bound gives a quadratic error, which is bounded by a fixed cubic multiple because the displacement norm is at least \(\delta\). Combining the cases proves

\[|R_2(x,v)|\le\eta\|v\|^2+C_\eta\|v\|^3.\]

The same bounded Hessian proves that the actual gradient of the compact \(C^{2}\) test is Lipschitz.

\begin{leanbox}{Lean: \texttt{compactC2\_\allowbreak{}secondOrderExpansion\_\allowbreak{}remainder}, \texttt{compactC2\_\allowbreak{}gradient\_\allowbreak{}lipschitz}}
\leansource{\texttt{compactC2\_\allowbreak{}secondOrderExpansion\_\allowbreak{}remainder} \quad AverageLMC/\allowbreak{}CompactC2Taylor.lean, lines 60–106}
\begin{leancode}
\leanline{0}{theorem\ compactC2\_secondOrderExpansion\_remainder}
\leanline{4}{(φ\ :\ H\ →\ ℝ)\ (hφ\ :\ ContDiff\ ℝ\ 2\ φ)\ (hc\ :\ HasCompactSupport\ φ)}
\leanline{4}{(η\ :\ ℝ)\ (hη\ :\ 0\ <\ η)\ :}
\leanline{4}{∃\ C\ :\ ℝ,\ 0\ ≤\ C\ ∧\ ∀\ x\ v,}
\leanline{6}{|φ\ (x\ +\ v)\ -\ ULA.secondOrderExpansion\ φ\ x\ v|\ ≤}
\leanline{8}{η\ *\ ‖v‖\ \textasciicircum{}\ 2\ +\ C\ *\ ‖v‖\ \textasciicircum{}\ 3}
\end{leancode}
\leansource{\texttt{compactC2\_\allowbreak{}gradient\_\allowbreak{}lipschitz} \quad AverageLMC/\allowbreak{}CompactC2Taylor.lean, lines 110–126}
\begin{leancode}
\leanline{0}{theorem\ compactC2\_gradient\_lipschitz}
\leanline{4}{(φ\ :\ H\ →\ ℝ)\ (hφ\ :\ ContDiff\ ℝ\ 2\ φ)\ (hc\ :\ HasCompactSupport\ φ)\ :}
\leanline{4}{∃\ L\ :\ ℝ≥0,\ LipschitzWith\ L\ (gradient\ φ)}
\end{leancode}
\leanlabel{How the Lean proof follows the argument}
\texttt{compact\allowbreak{}C2\_\allowbreak{}second\allowbreak{}Order\allowbreak{}Expansion\_\allowbreak{}remainder} separates the cases \texttt{hv : \ensuremath{\Vert}v\ensuremath{\Vert} < \ensuremath{\delta}} and its negation before combining them. The main intermediate facts are \texttt{h\allowbreak{}Hc}, \texttt{h\allowbreak{}Huc}, \texttt{B}, \texttt{h\allowbreak{}B}, \texttt{hbound}, \texttt{hop} in that order. Its remaining equalities are discharged by the simplifications displayed in the proof.




\texttt{compact\allowbreak{}C2\_\allowbreak{}gradient\_\allowbreak{}lipschitz} gives the formal counterpart of this step. The main intermediate facts are \texttt{C}, \texttt{L}, \texttt{h\allowbreak{}Ddiff}, \texttt{h\allowbreak{}Dlip} in that order. Named proof applications include \texttt{lipschitz\allowbreak{}With\_\allowbreak{}of\_\allowbreak{}nnnorm\_\allowbreak{}fderiv\_\allowbreak{}le}, \texttt{h\allowbreak{}Dlip.\allowbreak{}dist\_\allowbreak{}le\_\allowbreak{}mul}; the adjacent proof shows their hypotheses and order of use.
\end{leanbox}


Earlier results used here: \cite[Section~3.55.1]{LangevinLean2026ULASC}.
\subsection{\texorpdfstring{Separate the deterministic drift from Gaussian noise}{Separate the deterministic drift from Gaussian noise}}
\label{report-step-AverageLMC-CompactC2Taylor-lean-3}


Expand the test first at the Gaussian perturbation alone. Its linear term averages to zero and its quadratic term averages to \(h\) times the true Laplacian. Integrating the \(C^{2}\) remainder uses the exact Gaussian second moment and finite Gaussian third moment. Choose \(\eta\) small enough that the quadratic Gaussian contribution is at most \(\varepsilon\) \(h\). Insert the deterministic \(hv\) shift afterward using the first-order Lipschitz-gradient remainder. This contributes only \(h^2\) times the squared drift norm and a mixed \(h^{3/2}\) term. For \(h\) at most one, their sum is bounded by

\[h(\varepsilon+K_\varepsilon\sqrt h)(1+\|v\|^2).\]

The bound is uniform in the center, drift and admissible step size. No third moment of the drift is required.

\begin{leanbox}{Lean: \texttt{compactC2\_\allowbreak{}gaussian\_\allowbreak{}weak\_\allowbreak{}step}}
\leansource{\texttt{compactC2\_\allowbreak{}gaussian\_\allowbreak{}weak\_\allowbreak{}step} \quad AverageLMC/\allowbreak{}CompactC2Taylor.lean, lines 130–273}
\begin{leancode}
\leanline{0}{theorem\ compactC2\_gaussian\_weak\_step}
\leanline{4}{(φ\ :\ H\ →\ ℝ)\ (hφ\ :\ ContDiff\ ℝ\ 2\ φ)\ (hc\ :\ HasCompactSupport\ φ)}
\leanline{4}{(ε\ :\ ℝ)\ (hε\ :\ 0\ <\ ε)\ :}
\leanline{4}{∃\ K\ :\ ℝ,\ 0\ ≤\ K\ ∧\ ∀\ (x\ v\ :\ H)\ (h\ :\ ℝ),\ 0\ ≤\ h\ →\ h\ ≤\ 1\ →}
\leanline{6}{|(∫\ z\ :\ H,\ φ\ (x\ +\ h\ •\ v\ +\ Real.sqrt\ (2\ *\ h)\ •\ z)\ ∂stdGaussian\ H)\ -}
\leanline{8}{φ\ x\ -\ h\ *\ (fderiv\ ℝ\ φ\ x\ v\ +\ Laplacian.laplacian\ φ\ x)|\ ≤}
\leanline{10}{h\ *\ (ε\ +\ K\ *\ Real.sqrt\ h)\ *\ (1\ +\ ‖v‖\ \textasciicircum{}\ 2)}
\end{leancode}
\leanlabel{How the Lean proof follows the argument}
\texttt{compact\allowbreak{}C2\_\allowbreak{}gaussian\_\allowbreak{}weak\_\allowbreak{}step} gives the formal counterpart of this step. The main intermediate facts are \texttt{hz2}, \texttt{G\ensuremath{{}_{2}}}, \texttt{h\allowbreak{}G\ensuremath{{}_{2}}}, \texttt{\ensuremath{\eta}}, \texttt{h\ensuremath{\eta}}, \texttt{hz1} in that order. Named proof applications include \texttt{hnoise.\allowbreak{}sub}, \texttt{abs\_\allowbreak{}integral\_\allowbreak{}le\_\allowbreak{}integral\_\allowbreak{}abs.\allowbreak{}trans}, \texttt{htotal.\allowbreak{}trans}; the adjacent proof shows their hypotheses and order of use.
\end{leanbox}


Earlier results used here: \cite[Section~3.55.1]{LangevinLean2026ULASC}.
\subsection{\texorpdfstring{Lift the uniform bound through bounded past weights}{Lift the uniform bound through bounded past weights}}
\label{report-step-AverageLMC-CompactC2Taylor-lean-4}


The \(C^{2}\) Laplacian is continuous as a finite sum of continuous second derivatives. This makes the entire weak residual measurable for measurable random centers and drifts. If the past weight has absolute value at most \(A\) almost everywhere, multiply the pointwise residual bound by \(A\) and integrate. Square integrability of the random drift makes the majorant integrable. The resulting bound is \(Ah(\varepsilon+K\sqrt h)\) times one plus the drift second moment. The fixed-\(K\) version retains a supplied uniform pointwise constant unchanged, so the same constant can be used across every Euler mesh and mesh state.

\begin{leanbox}{Lean: \texttt{compactC2\_\allowbreak{}continuous\_\allowbreak{}laplacian}, \texttt{compactC2\_\allowbreak{}gaussian\_\allowbreak{}weighted\_\allowbreak{}weak\_\allowbreak{}step}, \texttt{gaussian\_\allowbreak{}weighted\_\allowbreak{}weak\_\allowbreak{}step\_\allowbreak{}of\_\allowbreak{}bound}}
\leansource{\texttt{compactC2\_\allowbreak{}continuous\_\allowbreak{}laplacian} \quad AverageLMC/\allowbreak{}CompactC2Taylor.lean, lines 276–289}
\begin{leancode}
\leanline{0}{theorem\ compactC2\_continuous\_laplacian}
\leanline{4}{(φ\ :\ H\ →\ ℝ)\ (hφ\ :\ ContDiff\ ℝ\ 2\ φ)\ :\ Continuous\ (Laplacian.laplacian\ φ)}
\end{leancode}
\leansource{\texttt{compactC2\_\allowbreak{}gaussian\_\allowbreak{}weighted\_\allowbreak{}weak\_\allowbreak{}step} \quad AverageLMC/\allowbreak{}CompactC2Taylor.lean, lines 294–351}
\begin{leancode}
\leanline{0}{theorem\ compactC2\_gaussian\_weighted\_weak\_step}
\leanline{4}{\{Ω\ :\ Type*\}\ [MeasurableSpace\ Ω]\ (P\ :\ Measure\ Ω)\ [IsProbabilityMeasure\ P]}
\leanline{4}{(φ\ :\ H\ →\ ℝ)\ (hφ\ :\ ContDiff\ ℝ\ 2\ φ)\ (hc\ :\ HasCompactSupport\ φ)}
\leanline{4}{(x\ v\ :\ Ω\ →\ H)\ (w\ :\ Ω\ →\ ℝ)\ (A\ :\ ℝ)}
\leanline{4}{(hx\ :\ Measurable\ x)\ (hv\ :\ Measurable\ v)\ (hw\ :\ Measurable\ w)}
\leanline{4}{(hv₂\ :\ MemLp\ v\ 2\ P)\ (hA\ :\ 0\ ≤\ A)\ (hwA\ :\ ∀ᵐ\ a\ ∂P,\ |w\ a|\ ≤\ A)}
\leanline{4}{(ε\ :\ ℝ)\ (hε\ :\ 0\ <\ ε)\ :}
\leanline{4}{∃\ K\ :\ ℝ,\ 0\ ≤\ K\ ∧\ ∀\ h\ :\ ℝ,\ 0\ ≤\ h\ →\ h\ ≤\ 1\ →}
\leanline{6}{|∫\ a,\ w\ a\ *\ ((∫\ z\ :\ H,}
\leanline{8}{φ\ (x\ a\ +\ h\ •\ v\ a\ +\ Real.sqrt\ (2\ *\ h)\ •\ z)\ ∂stdGaussian\ H)\ -}
\leanline{10}{φ\ (x\ a)\ -\ h\ *\ (fderiv\ ℝ\ φ\ (x\ a)\ (v\ a)\ +}
\leanline{12}{Laplacian.laplacian\ φ\ (x\ a)))\ ∂P|\ ≤}
\leanline{8}{A\ *\ h\ *\ (ε\ +\ K\ *\ Real.sqrt\ h)\ *\ (1\ +\ ∫\ a,\ ‖v\ a‖\ \textasciicircum{}\ 2\ ∂P)}
\end{leancode}
\leansource{\texttt{gaussian\_\allowbreak{}weighted\_\allowbreak{}weak\_\allowbreak{}step\_\allowbreak{}of\_\allowbreak{}bound} \quad AverageLMC/\allowbreak{}CompactC2Taylor.lean, lines 355–414}
\begin{leancode}
\leanline{0}{theorem\ gaussian\_weighted\_weak\_step\_of\_bound}
\leanline{4}{\{Ω\ :\ Type*\}\ [MeasurableSpace\ Ω]\ (P\ :\ Measure\ Ω)\ [IsProbabilityMeasure\ P]}
\leanline{4}{(φ\ :\ H\ →\ ℝ)\ (hφ\ :\ ContDiff\ ℝ\ 2\ φ)}
\leanline{4}{(x\ v\ :\ Ω\ →\ H)\ (w\ :\ Ω\ →\ ℝ)\ (A\ :\ ℝ)}
\leanline{4}{(hx\ :\ Measurable\ x)\ (hv\ :\ Measurable\ v)\ (hw\ :\ Measurable\ w)}
\leanline{4}{(hv₂\ :\ MemLp\ v\ 2\ P)\ (hA\ :\ 0\ ≤\ A)\ (hwA\ :\ ∀ᵐ\ a\ ∂P,\ |w\ a|\ ≤\ A)}
\leanline{4}{(ε\ :\ ℝ)\ (hε\ :\ 0\ <\ ε)\ (K\ :\ ℝ)\ (hK\ :\ 0\ ≤\ K)}
\leanline{4}{(hbound\ :\ ∀\ (x\ v\ :\ H)\ (h\ :\ ℝ),\ 0\ ≤\ h\ →\ h\ ≤\ 1\ →}
\leanline{6}{|(∫\ z\ :\ H,\ φ\ (x\ +\ h\ •\ v\ +\ Real.sqrt\ (2\ *\ h)\ •\ z)\ ∂stdGaussian\ H)\ -}
\leanline{8}{φ\ x\ -\ h\ *\ (fderiv\ ℝ\ φ\ x\ v\ +\ Laplacian.laplacian\ φ\ x)|\ ≤}
\leanline{10}{h\ *\ (ε\ +\ K\ *\ Real.sqrt\ h)\ *\ (1\ +\ ‖v‖\ \textasciicircum{}\ 2))\ :}
\leanline{4}{∀\ h\ :\ ℝ,\ 0\ ≤\ h\ →\ h\ ≤\ 1\ →}
\leanline{6}{|∫\ a,\ w\ a\ *\ ((∫\ z\ :\ H,}
\leanline{8}{φ\ (x\ a\ +\ h\ •\ v\ a\ +\ Real.sqrt\ (2\ *\ h)\ •\ z)\ ∂stdGaussian\ H)\ -}
\leanline{10}{φ\ (x\ a)\ -\ h\ *\ (fderiv\ ℝ\ φ\ (x\ a)\ (v\ a)\ +}
\leanline{12}{Laplacian.laplacian\ φ\ (x\ a)))\ ∂P|\ ≤}
\leanline{8}{A\ *\ h\ *\ (ε\ +\ K\ *\ Real.sqrt\ h)\ *\ (1\ +\ ∫\ a,\ ‖v\ a‖\ \textasciicircum{}\ 2\ ∂P)}
\end{leancode}
\leanlabel{How the Lean proof follows the argument}
\texttt{compact\allowbreak{}C2\_\allowbreak{}continuous\_\allowbreak{}laplacian} gives the formal counterpart of this step. The main intermediate facts are \texttt{h\allowbreak{}DD}, \texttt{he} in that order. Named proof applications include \texttt{continuous\_\allowbreak{}finset\_\allowbreak{}sum}; the adjacent proof shows their hypotheses and order of use.




\texttt{compact\allowbreak{}C2\_\allowbreak{}gaussian\_\allowbreak{}weighted\_\allowbreak{}weak\_\allowbreak{}step} gives the formal counterpart of this step. The main intermediate facts are \texttt{r}, \texttt{h\ensuremath{\varphi}m}, \texttt{hjoint}, \texttt{hint}, \texttt{h\allowbreak{}D}, \texttt{h\allowbreak{}Lap} in that order. Named proof applications include \texttt{Measurable.\allowbreak{}aestrongly\allowbreak{}Measurable}, \texttt{abs\_\allowbreak{}integral\_\allowbreak{}le\_\allowbreak{}integral\_\allowbreak{}abs.\allowbreak{}trans}; the adjacent proof shows their hypotheses and order of use.




\texttt{gaussian\_\allowbreak{}weighted\_\allowbreak{}weak\_\allowbreak{}step\_\allowbreak{}of\_\allowbreak{}bound} gives the formal counterpart of this step. The main intermediate facts are \texttt{r}, \texttt{h\ensuremath{\varphi}m}, \texttt{hjoint}, \texttt{hint}, \texttt{h\allowbreak{}D}, \texttt{h\allowbreak{}Lap} in that order. Named proof applications include \texttt{Measurable.\allowbreak{}aestrongly\allowbreak{}Measurable}, \texttt{abs\_\allowbreak{}integral\_\allowbreak{}le\_\allowbreak{}integral\_\allowbreak{}abs.\allowbreak{}trans}; the adjacent proof shows their hypotheses and order of use.
\end{leanbox}
\section{\texorpdfstring{Fresh Gaussian innovations with retained initial data}{Fresh Gaussian innovations with retained initial data}}
\label{proof-module-AverageLMC-LaplacianIndependentSteps-lean}


The weak consistency proof needs the actual joint law of the current recursive state and its next innovation, including a passive initial weight. These results obtain that law from independence of initialization and an iid noise path.



For iid innovations of law \(\eta\), independent of initialization, the actual recursion has the fresh joint law


\[X_{k+1}=F(X_k,W_k),\qquad
\operatorname{Law}(X_k,W_k)=\operatorname{Law}(X_k)\otimes\eta.\]


Copying a passive initial coordinate through each update retains it exactly:


\[\operatorname{Chain}_{\widetilde F}((a,x),w)_n
=\bigl(a,\operatorname{Chain}_F(x,w)_n\bigr),\qquad
\widetilde F((a,x),z)=(a,F(x,z)).\]

Earlier results used here: \cite[Section~3.33]{LangevinLean2026ULASC}.
\subsection{\texorpdfstring{Identify the current-state and next-noise law}{Identify the current-state and next-noise law}}
\label{report-step-AverageLMC-LaplacianIndependentSteps-lean-1}


The iid family has the infinite product path law. Independence makes its joint law with the initialization a product. Push this identity through the measurable recursion and use freshness of the canonical product chain to obtain the product of the actual current-state law and the innovation law.

\begin{leanbox}{Lean: \texttt{noiseDrivenChain\_\allowbreak{}joint\_\allowbreak{}fresh\_\allowbreak{}of\_\allowbreak{}iid}}
\leansource{\texttt{noiseDrivenChain\_\allowbreak{}joint\_\allowbreak{}fresh\_\allowbreak{}of\_\allowbreak{}iid} \quad AverageLMC/\allowbreak{}LaplacianIndependentSteps.lean, lines 25–55}
\begin{leancode}
\leanline{0}{theorem\ noiseDrivenChain\_joint\_fresh\_of\_iid}
\leanline{4}{(P\ :\ Measure\ Ω)\ [IsProbabilityMeasure\ P]}
\leanline{4}{(F\ :\ H\ →\ Z\ →\ H)\ (hF\ :\ Measurable\ (Function.uncurry\ F))}
\leanline{4}{(X\ :\ Ω\ →\ H)\ (W\ :\ ℕ\ →\ Ω\ →\ Z)\ (hX\ :\ AEMeasurable\ X\ P)}
\leanline{4}{(hW\ :\ ∀\ n,\ AEMeasurable\ (W\ n)\ P)}
\leanline{4}{(hind\ :\ IndepFun\ X\ (fun\ a\ n\ ↦\ W\ n\ a)\ P)\ (hiid\ :\ iIndepFun\ W\ P)}
\leanline{4}{(η\ :\ Measure\ Z)\ [IsProbabilityMeasure\ η]\ (hlaw\ :\ ∀\ n,\ P.map\ (W\ n)\ =\ η)\ (n\ :\ ℕ)\ :}
\leanline{4}{P.map\ (fun\ a\ ↦\ (ULA.noiseDrivenChain\ F\ (X\ a)\ (fun\ k\ ↦\ W\ k\ a)\ n,\ W\ n\ a))\ =}
\leanline{6}{(P.map\ (fun\ a\ ↦\ ULA.noiseDrivenChain\ F\ (X\ a)\ (fun\ k\ ↦\ W\ k\ a)\ n)).prod\ η}
\end{leancode}
\leanlabel{How the Lean proof follows the argument}
\texttt{noise\allowbreak{}Driven\allowbreak{}Chain\_\allowbreak{}joint\_\allowbreak{}fresh\_\allowbreak{}of\_\allowbreak{}iid} gives the formal counterpart of this step. The main intermediate facts are \texttt{Q}, \texttt{hpath}, \texttt{hpair}, \texttt{hmap}, \texttt{hprod}, \texttt{hjoint} in that order. Named proof applications include \texttt{hp}, \texttt{hjoint}; the adjacent proof shows their hypotheses and order of use.
\end{leanbox}


Earlier results used here: \cite[Section~3.20.1]{LangevinLean2026ULASC}; \cite[Section~3.20.4]{LangevinLean2026ULASC}; \cite[Section~3.20.8]{LangevinLean2026ULASC}; \cite[Section~3.33.7]{LangevinLean2026ULASC}.
\subsection{\texorpdfstring{Retain the initial weight exactly}{Retain the initial weight exactly}}
\label{report-step-AverageLMC-LaplacianIndependentSteps-lean-2}


Run the recursion on a pair whose first coordinate is copied at each step. Induction proves that this coordinate remains the original value while the second coordinate is precisely the ordinary noise-driven chain. This allows a bounded initial weight to be included in the freshness argument.

\begin{leanbox}{Lean: \texttt{noiseDrivenChain\_\allowbreak{}passive}}
\leansource{\texttt{noiseDrivenChain\_\allowbreak{}passive} \quad AverageLMC/\allowbreak{}LaplacianIndependentSteps.lean, lines 58–64}
\begin{leancode}
\leanline{0}{theorem\ noiseDrivenChain\_passive\ \{A\ :\ Type*\}\ (F\ :\ H\ →\ Z\ →\ H)}
\leanline{4}{(a\ :\ A)\ (x\ :\ H)\ (w\ :\ ℕ\ →\ Z)\ (n\ :\ ℕ)\ :}
\leanline{4}{ULA.noiseDrivenChain\ (fun\ p\ :\ A\ ×\ H\ ↦\ fun\ z\ ↦\ (p.1,\ F\ p.2\ z))\ (a,\ x)\ w\ n\ =}
\leanline{6}{(a,\ ULA.noiseDrivenChain\ F\ x\ w\ n)}
\end{leancode}
\leanlabel{How the Lean proof follows the argument}
\texttt{noise\allowbreak{}Driven\allowbreak{}Chain\_\allowbreak{}passive} proceeds by \texttt{induction n with}. The base cases establish the initial claims and the inductive branches use the induction hypotheses. Its remaining equalities are discharged by the simplifications displayed in the proof.
\end{leanbox}


Earlier results used here: \cite[Section~3.20.1]{LangevinLean2026ULASC}; \cite[Section~3.20.3]{LangevinLean2026ULASC}.
\section{\texorpdfstring{Weighted weak consistency of the Euler chain}{Weighted weak consistency of the Euler chain}}
\label{proof-module-AverageLMC-LaplacianWeakChain-lean}


This section proves a weak consistency estimate for the Euler (unadjusted Langevin) chain that is the discrete input of the Dynkin identity used in the proof of Theorem 2. Let \(H\) be a finite-dimensional real inner product space (in the application \(H=\mathbb R^p\)), \(\gamma\) the standard Gaussian measure on \(H\), \(U:H\to\mathbb R\) a potential with Lipschitz gradient and


\[X_{k+1}=X_k-h\nabla U(X_k)+\sqrt{2h}\,W_k\qquad(k\ge0),\]


where \(W_0,W_1,\dots\) are independent with law \(\gamma\) and independent of \((w,X_0)\) for a bounded random weight \(w\). For a test function \(\varphi\in C^2\) with \(\varphi\) and \(\mathcal L\varphi=\Delta\varphi-\langle\nabla U,\nabla\varphi\rangle\) bounded, write \(T(F,X)=\mathbb E[w\,F(X)]\). Assuming a uniform Gaussian Taylor estimate for \(\varphi\) with constants \(\varepsilon,K\) (proved elsewhere for compactly supported \(\varphi\)), we show


\[\Bigl|T(\varphi,X_n)-T(\varphi,X_0)-\sum_{k=0}^{n-1}h\,T(\mathcal L\varphi,X_k)\Bigr|\le A\,(nh)\,(\varepsilon+K\sqrt h)\,(1+R^2)\]


whenever \(|w|\le A\) almost surely, \(0\le h\le1\), and \(\|\nabla U(X_k)\|_{L_2}\le R\) for \(k\le n\). The weight \(w\) carries arbitrary bounded information about the initial state; it is needed because the Dynkin identity is applied to test functions multiplied by a function of the initial coupling. The proof combines the independence of each new Gaussian increment from the past (including the weight), Fubini's theorem, and a telescoping sum.



Earlier results used here: \cite[Section~3.42]{LangevinLean2026ULASC}.
\subsection{\texorpdfstring{One Euler step with a fresh Gaussian increment and a bounded weight}{One Euler step with a fresh Gaussian increment and a bounded weight}}
\label{report-step-AverageLMC-LaplacianWeakChain-lean-1}


Throughout, \((\Omega,\mathcal F,P)\) is a probability space and \(\gamma\) is the standard Gaussian measure on \(H\).




\textbf{Claim.} (1) Let \(w,F:\Omega\to\mathbb R\) be measurable, \(A\ge0\), \(C\in\mathbb R\), with \(|w|\le A\) almost surely and \(|F|\le C\) everywhere. Then \(wF\) is \(P\)-integrable.




(2) Let \(\varphi:H\to\mathbb R\) be continuous with \(|\varphi|\le C\), \(b:H\to H\) measurable, \(x,z:\Omega\to H\) and \(w:\Omega\to\mathbb R\) measurable with \(|w|\le A\) almost surely (\(A\ge0\)), and assume that the joint law of \(((w,x),z)\) is the product of the law of \((w,x)\) and \(\gamma\) (that is, \(z\sim\gamma\) is independent of \((w,x)\)). Then for every \(h\in\mathbb R\) (with \(\sqrt{2h}=0\) for \(h<0\)),


\[\mathbb E\bigl[w\,\varphi\bigl(x+hb(x)+\sqrt{2h}\,z\bigr)\bigr]=\mathbb E\Bigl[w\int_H\varphi\bigl(x+hb(x)+\sqrt{2h}\,y\bigr)\,\gamma(dy)\Bigr].\]


(3) In the setting of (2), let in addition \(\varphi\in C^2\), \(b(x)\in L_2(P)\), and \(w\,(D\varphi(x)b(x)+\Delta\varphi(x))\) integrable. Let \(\varepsilon>0\), \(K\ge0\), and assume the uniform Gaussian Taylor estimate


\[\Bigl|\int\varphi(y+hv+\sqrt{2h}\,u)\,\gamma(du)-\varphi(y)-h\bigl(D\varphi(y)v+\Delta\varphi(y)\bigr)\Bigr|\le h(\varepsilon+K\sqrt h)(1+\|v\|^2)\qquad(1)\]


for all \(y,v\in H\) and \(0\le h\le1\). Then for \(0\le h\le1\),


\[\begin{gathered}
\left|\begin{aligned}
&\mathbb E\bigl[w\varphi(x+hb(x)+\sqrt{2h}z)\bigr]-\mathbb E[w\varphi(x)]\\
&\quad-h\,\mathbb E\bigl[w\bigl(D\varphi(x)b(x)+\Delta\varphi(x)\bigr)\bigr]
\end{aligned}\right|\\
\le A\,h\,(\varepsilon+K\sqrt h)\bigl(1+\mathbb E\|b(x)\|^2\bigr).
\end{gathered}\]


\textbf{Proof.} (1) \(wF\) is measurable and \(|wF|\le AC\) almost surely; a constant is integrable for a probability measure, so \(wF\) is integrable by domination.




(2) Let \(G((s,u),y)=s\,\varphi(u+hb(u)+\sqrt{2h}y)\) on \((\mathbb R\times H)\times H\), a measurable function. By the change-of-variables formula and the product form of the law,


\[\mathbb E\bigl[G((w,x),z)\bigr]=\int G\,d\bigl(\mathcal L(w,x)\otimes\gamma\bigr)=\int_{\Omega\times H}w(a)\,\varphi\bigl(x(a)+hb(x(a))+\sqrt{2h}\,y\bigr)\,(P\otimes\gamma)(da,dy),\]


because \(\mathcal L(w,x)\otimes\gamma\) is the image of \(P\otimes\gamma\) under \((a,y)\mapsto((w(a),x(a)),y)\). The last integrand is integrable by (1), applied on \(\Omega\times H\) with \(P\otimes\gamma\) (the bound \(|w(a)|\le A\) holds for \(P\otimes\gamma\)-almost every \((a,y)\), the first marginal being \(P\)). Fubini's theorem integrates first in \(y\), and the factor \(w(a)\) comes out of the inner integral.




(3) Multiply (1) at \(y=x(a)\), \(v=b(x(a))\) by \(|w(a)|\le A\) and integrate; this gives (as proved in the section on compact \(C^2\) Taylor estimates)


\[\begin{gathered}
\left|\mathbb E\left[w\left(\begin{aligned}
&\int\varphi(x+hb(x)+\sqrt{2h}u)\gamma(du)-\varphi(x)\\
&\quad-h\bigl(D\varphi(x)b(x)+\Delta\varphi(x)\bigr)
\end{aligned}\right)\right]\right|\\
\le Ah(\varepsilon+K\sqrt h)\bigl(1+\mathbb E\|b(x)\|^2\bigr).
\end{gathered}\]


The three functions \(w\int\varphi(\cdot)\,d\gamma\), \(w\varphi(x)\) and \(w(D\varphi(x)b(x)+\Delta\varphi(x))\) are integrable: the first because its Gaussian integral is the inner integral of an integrable function on \(\Omega\times H\) (Fubini), the second by (1), the third by assumption. So the expectation splits into three expectations, and by (2) the first equals \(\mathbb E[w\varphi(x+hb(x)+\sqrt{2h}z)]\). \(\square\)

\begin{leanbox}{Lean: \texttt{integrable\_\allowbreak{}bounded\_\allowbreak{}weight\_\allowbreak{}observation}, \texttt{fresh\_\allowbreak{}weighted\_\allowbreak{}gaussian\_\allowbreak{}integral}, \texttt{fresh\_\allowbreak{}weighted\_\allowbreak{}weak\_\allowbreak{}step\_\allowbreak{}of\_\allowbreak{}bound}}
\leansource{\texttt{integrable\_\allowbreak{}bounded\_\allowbreak{}weight\_\allowbreak{}observation} \quad AverageLMC/\allowbreak{}LaplacianWeakChain.lean, lines 19–27}
\begin{leancode}
\leanline{0}{theorem\ integrable\_bounded\_weight\_observation}
\leanline{4}{(P\ :\ Measure\ Ω)\ [IsProbabilityMeasure\ P]\ (w\ F\ :\ Ω\ →\ ℝ)}
\leanline{4}{(hw\ :\ Measurable\ w)\ (hF\ :\ Measurable\ F)\ (A\ C\ :\ ℝ)\ (hA\ :\ 0\ ≤\ A)}
\leanline{4}{(hwA\ :\ ∀ᵐ\ a\ ∂P,\ |w\ a|\ ≤\ A)\ (hFC\ :\ ∀\ a,\ |F\ a|\ ≤\ C)\ :}
\leanline{4}{Integrable\ (fun\ a\ ↦\ w\ a\ *\ F\ a)\ P}
\end{leancode}
\leansource{\texttt{fresh\_\allowbreak{}weighted\_\allowbreak{}gaussian\_\allowbreak{}integral} \quad AverageLMC/\allowbreak{}LaplacianWeakChain.lean, lines 32–72}
\begin{leancode}
\leanline{0}{theorem\ fresh\_weighted\_gaussian\_integral}
\leanline{4}{(P\ :\ Measure\ Ω)\ [IsProbabilityMeasure\ P]}
\leanline{4}{(φ\ :\ H\ →\ ℝ)\ (hφ\ :\ Continuous\ φ)\ (C\ :\ ℝ)\ (hφC\ :\ ∀\ x,\ |φ\ x|\ ≤\ C)}
\leanline{4}{(b\ :\ H\ →\ H)\ (hb\ :\ Measurable\ b)\ (x\ z\ :\ Ω\ →\ H)\ (w\ :\ Ω\ →\ ℝ)}
\leanline{4}{(hx\ :\ Measurable\ x)\ (hz\ :\ Measurable\ z)\ (hw\ :\ Measurable\ w)}
\leanline{4}{(A\ :\ ℝ)\ (hA\ :\ 0\ ≤\ A)\ (hwA\ :\ ∀ᵐ\ a\ ∂P,\ |w\ a|\ ≤\ A)}
\leanline{4}{(hfresh\ :\ P.map\ (fun\ a\ ↦\ ((w\ a,\ x\ a),\ z\ a))\ =}
\leanline{6}{(P.map\ (fun\ a\ ↦\ (w\ a,\ x\ a))).prod\ (stdGaussian\ H))\ (h\ :\ ℝ)\ :}
\leanline{4}{(∫\ a,\ w\ a\ *\ φ\ (x\ a\ +\ h\ •\ b\ (x\ a)\ +\ Real.sqrt\ (2\ *\ h)\ •\ z\ a)\ ∂P)\ =}
\leanline{6}{∫\ a,\ w\ a\ *\ (∫\ y\ :\ H,\ φ\ (x\ a\ +\ h\ •\ b\ (x\ a)\ +\ Real.sqrt\ (2\ *\ h)\ •\ y)}
\leanline{8}{∂stdGaussian\ H)\ ∂P}
\end{leancode}
\leansource{\texttt{fresh\_\allowbreak{}weighted\_\allowbreak{}weak\_\allowbreak{}step\_\allowbreak{}of\_\allowbreak{}bound} \quad AverageLMC/\allowbreak{}LaplacianWeakChain.lean, lines 77–128}
\begin{leancode}
\leanline{0}{theorem\ fresh\_weighted\_weak\_step\_of\_bound}
\leanline{4}{(P\ :\ Measure\ Ω)\ [IsProbabilityMeasure\ P]}
\leanline{4}{(φ\ :\ H\ →\ ℝ)\ (hφ\ :\ ContDiff\ ℝ\ 2\ φ)\ (C\ :\ ℝ)\ (hφC\ :\ ∀\ x,\ |φ\ x|\ ≤\ C)}
\leanline{4}{(b\ :\ H\ →\ H)\ (hb\ :\ Measurable\ b)\ (x\ z\ :\ Ω\ →\ H)\ (w\ :\ Ω\ →\ ℝ)}
\leanline{4}{(hx\ :\ Measurable\ x)\ (hz\ :\ Measurable\ z)\ (hw\ :\ Measurable\ w)}
\leanline{4}{(hb₂\ :\ MemLp\ (fun\ a\ ↦\ b\ (x\ a))\ 2\ P)}
\leanline{4}{(hgen\ :\ Integrable\ (fun\ a\ ↦\ w\ a\ *}
\leanline{6}{(fderiv\ ℝ\ φ\ (x\ a)\ (b\ (x\ a))\ +\ Laplacian.laplacian\ φ\ (x\ a)))\ P)}
\leanline{4}{(A\ :\ ℝ)\ (hA\ :\ 0\ ≤\ A)\ (hwA\ :\ ∀ᵐ\ a\ ∂P,\ |w\ a|\ ≤\ A)}
\leanline{4}{(hfresh\ :\ P.map\ (fun\ a\ ↦\ ((w\ a,\ x\ a),\ z\ a))\ =}
\leanline{6}{(P.map\ (fun\ a\ ↦\ (w\ a,\ x\ a))).prod\ (stdGaussian\ H))}
\leanline{4}{(ε\ :\ ℝ)\ (hε\ :\ 0\ <\ ε)\ (K\ :\ ℝ)\ (hK\ :\ 0\ ≤\ K)}
\leanline{4}{(hbound\ :\ ∀\ (x\ v\ :\ H)\ (h\ :\ ℝ),\ 0\ ≤\ h\ →\ h\ ≤\ 1\ →}
\leanline{6}{|(∫\ z\ :\ H,\ φ\ (x\ +\ h\ •\ v\ +\ Real.sqrt\ (2\ *\ h)\ •\ z)\ ∂stdGaussian\ H)\ -}
\leanline{8}{φ\ x\ -\ h\ *\ (fderiv\ ℝ\ φ\ x\ v\ +\ Laplacian.laplacian\ φ\ x)|\ ≤}
\leanline{10}{h\ *\ (ε\ +\ K\ *\ Real.sqrt\ h)\ *\ (1\ +\ ‖v‖\ \textasciicircum{}\ 2))}
\leanline{4}{(h\ :\ ℝ)\ (hh\ :\ 0\ ≤\ h)\ (hh1\ :\ h\ ≤\ 1)\ :}
\leanline{4}{|(∫\ a,\ w\ a\ *\ φ\ (x\ a\ +\ h\ •\ b\ (x\ a)\ +\ Real.sqrt\ (2\ *\ h)\ •\ z\ a)\ ∂P)\ -}
\leanline{6}{(∫\ a,\ w\ a\ *\ φ\ (x\ a)\ ∂P)\ -}
\leanline{6}{h\ *\ (∫\ a,\ w\ a\ *\ (fderiv\ ℝ\ φ\ (x\ a)\ (b\ (x\ a))\ +\ Laplacian.laplacian\ φ\ (x\ a))\ ∂P)|\ ≤}
\leanline{6}{A\ *\ h\ *\ (ε\ +\ K\ *\ Real.sqrt\ h)\ *\ (1\ +\ ∫\ a,\ ‖b\ (x\ a)‖\ \textasciicircum{}\ 2\ ∂P)}
\end{leancode}
\leanlabel{How the Lean proof follows the argument}
\begin{enumerate}


\item \texttt{integrable\_\allowbreak{}bounded\_\allowbreak{}weight\_\allowbreak{}observation} is (1): domination of \(wF\) by the integrable constant \texttt{A * C} (\texttt{integrable\_\allowbreak{}const} and \texttt{Integrable.\allowbreak{}mono'}), with the pointwise bound from \texttt{abs\_\allowbreak{}mul} and \texttt{mul\_\allowbreak{}le\_\allowbreak{}mul}.

\item \texttt{fresh\_\allowbreak{}weighted\_\allowbreak{}gaussian\_\allowbreak{}integral} is (2). The freshness hypothesis is \texttt{hfresh : P.\allowbreak{}map (fun a \ensuremath{\mapsto} ((w a, x a), z a)) = (P.\allowbreak{}map (fun a \ensuremath{\mapsto} (w a, x a))).\allowbreak{}prod (std\allowbreak{}Gaussian H)}, and \texttt{std\allowbreak{}Gaussian H} is \(\gamma\). \texttt{F} is the function \(G\); \texttt{hmap} is the image identity from \texttt{Measure.\allowbreak{}map\_\allowbreak{}prod\_\allowbreak{}map}; \texttt{heq} is the displayed chain, a \texttt{calc} with two applications of \texttt{integral\_\allowbreak{}map} and the rewrite by \texttt{hfresh}. \texttt{hi} is the integrability on the product (part (1) together with \texttt{measure\allowbreak{}Preserving\_\allowbreak{}fst} to transfer the almost-sure bound), \texttt{integral\_\allowbreak{}prod} is Fubini's theorem, and \texttt{integral\_\allowbreak{}const\_\allowbreak{}mul} extracts \(w(a)\).

\item \texttt{fresh\_\allowbreak{}weighted\_\allowbreak{}weak\_\allowbreak{}step\_\allowbreak{}of\_\allowbreak{}bound} is (3); the hypothesis \texttt{hbound} is (1), with \texttt{fderiv \ensuremath{\mathbb R} \ensuremath{\varphi} x v} for \(D\varphi(x)v\) and \texttt{Laplacian.\allowbreak{}laplacian \ensuremath{\varphi} x} for \(\Delta\varphi(x)\). \texttt{hweak} is the integrated estimate \texttt{gaussian\_\allowbreak{}weighted\_\allowbreak{}weak\_\allowbreak{}step\_\allowbreak{}of\_\allowbreak{}bound}. \texttt{hprod}, \texttt{hmean} (integrability of the Gaussian average, from \texttt{integral\_\allowbreak{}prod\_\allowbreak{}left}) and \texttt{hbase} provide the integrability, \texttt{heq} regroups the integrand, two applications of \texttt{integral\_\allowbreak{}sub} and one of \texttt{integral\_\allowbreak{}const\_\allowbreak{}mul} split the expectation, and the final rewrite uses part (2).

\end{enumerate}
\end{leanbox}
\subsection{\texorpdfstring{Freshness along the chain and the one-step residual}{Freshness along the chain and the one-step residual}}
\label{report-step-AverageLMC-LaplacianWeakChain-lean-2}


\textbf{Setting.} Let \(U:H\to\mathbb R\) have an \(L\)-Lipschitz gradient, \(h\in\mathbb R\), \(X_0\) and \(W_0,W_1,\dots\) square-integrable random vectors in \(H\), and \(w\) a measurable real random variable, such that \((w,X_0)\) is independent of the sequence \((W_k)_{k\ge0}\), the \(W_k\) are mutually independent, and each \(W_k\) has law \(\gamma\). The Euler chain is \(X_{k+1}=X_k-h\nabla U(X_k)+\sqrt{2h}\,W_k\). For a function \(F:H\to\mathbb R\) write \(T(F,X)=\mathbb E[w\,F(X)]\), and let \(\mathcal L\varphi=\Delta\varphi-\langle\nabla U,\nabla\varphi\rangle\).




\textbf{Claim.} (1) For every \(k\), the joint law of \(((w,X_k),W_k)\) is the product of the law of \((w,X_k)\) and \(\gamma\).




(2) Let \(\varphi\in C^2\) with \(|\varphi|\le C\) and \(|\mathcal L\varphi|\le D\), let \(|w|\le A\) almost surely with \(A\ge0\), and let \(\varphi\) satisfy (1) of Step 1 with \(\varepsilon>0\), \(K\ge0\). Then for \(0\le h\le1\) and every \(k\),


\[\bigl|T(\varphi,X_{k+1})-T(\varphi,X_k)-h\,T(\mathcal L\varphi,X_k)\bigr|\le A\,h\,(\varepsilon+K\sqrt h)\bigl(1+\|\nabla U(X_k)\|_{L_2}^2\bigr).\]


\textbf{Proof.} (1) Consider the augmented chain \(Z_k=(w,X_k)\) with values in \(\mathbb R\times H\). It satisfies \(Z_{k+1}=\Phi(Z_k,W_k)\) with \(\Phi((s,u),y)=(s,\,u-h\nabla U(u)+\sqrt{2h}y)\), a measurable map, since the first coordinate is carried along unchanged. For a chain driven by independent noise with common law \(\gamma\) and an initial state independent of the noise, \(Z_k\) is a measurable function of \(Z_0,W_0,\dots,W_{k-1}\), which are independent of \(W_k\); hence the joint law of \((Z_k,W_k)\) is the product of the law of \(Z_k\) and \(\gamma\) (this is proved in the section on independent Euler steps). Since \(Z_k=(w,X_k)\), this is the claim.




(2) Apply (3) of Step 1 with \(x=X_k\), \(z=W_k\) and \(b=-\nabla U\). Its hypotheses hold: \(-\nabla U\) is continuous, hence measurable; the freshness of \(W_k\) is part (1); \(b(X_k)=-\nabla U(X_k)\) is square integrable because a Lipschitz map sends square-integrable vectors to square-integrable vectors (\(\|\nabla U(u)\|\le\|\nabla U(0)\|+L\|u\|\)). Since \(D\varphi(u)v=\langle\nabla\varphi(u),v\rangle\),


\[D\varphi(u)\bigl(-\nabla U(u)\bigr)+\Delta\varphi(u)=\Delta\varphi(u)-\langle\nabla U(u),\nabla\varphi(u)\rangle=\mathcal L\varphi(u),\]


and \(w\,\mathcal L\varphi(X_k)\) is integrable by (1) of Step 1, because \(\mathcal L\varphi\) is continuous and bounded by \(D\). Finally \(X_{k+1}=X_k+h(-\nabla U(X_k))+\sqrt{2h}W_k\) almost surely, so \(T(\varphi,X_{k+1})\) is the first expectation in (3) of Step 1, and \(\mathbb E\|b(X_k)\|^2=\|\nabla U(X_k)\|_{L_2}^2\). \(\square\)

\begin{leanbox}{Lean: \texttt{hilbert\_\allowbreak{}lpChain\_\allowbreak{}weighted\_\allowbreak{}fresh}, \texttt{hilbert\_\allowbreak{}lpChain\_\allowbreak{}weak\_\allowbreak{}step\_\allowbreak{}of\_\allowbreak{}bound}}
\leansource{\texttt{hilbert\_\allowbreak{}lpChain\_\allowbreak{}weighted\_\allowbreak{}fresh} \quad AverageLMC/\allowbreak{}LaplacianWeakChain.lean, lines 146–178}
\begin{leancode}
\leanline{0}{theorem\ hilbert\_lpChain\_weighted\_fresh}
\leanline{4}{(P\ :\ Measure\ Ω)\ [IsProbabilityMeasure\ P]}
\leanline{4}{(U\ :\ H\ →\ ℝ)\ (h\ :\ ℝ)\ (L\ :\ ℝ≥0)\ (hLip\ :\ LipschitzWith\ L\ (gradient\ U))}
\leanline{4}{(initial\ :\ Lp\ H\ 2\ P)\ (W\ :\ ℕ\ →\ Lp\ H\ 2\ P)\ (w\ :\ Ω\ →\ ℝ)\ (hw\ :\ Measurable\ w)}
\leanline{4}{(hind\ :\ IndepFun\ (fun\ a\ ↦\ (w\ a,\ initial\ a))\ (fun\ a\ k\ ↦\ W\ k\ a)\ P)}
\leanline{4}{(hiid\ :\ iIndepFun\ (fun\ k\ a\ ↦\ W\ k\ a)\ P)}
\leanline{4}{(hlaw\ :\ ∀\ k,\ P.map\ (W\ k)\ =\ stdGaussian\ H)\ (k\ :\ ℕ)\ :}
\leanline{4}{P.map\ (fun\ a\ ↦\ ((w\ a,\ ULA.Hilbert.lpChain\ P\ U\ h\ hLip\ initial\ W\ k\ a),\ W\ k\ a))\ =}
\leanline{6}{(P.map\ (fun\ a\ ↦\ (w\ a,\ ULA.Hilbert.lpChain\ P\ U\ h\ hLip\ initial\ W\ k\ a))).prod}
\leanline{8}{(stdGaussian\ H)}
\end{leancode}
\leansource{\texttt{hilbert\_\allowbreak{}lpChain\_\allowbreak{}weak\_\allowbreak{}step\_\allowbreak{}of\_\allowbreak{}bound} \quad AverageLMC/\allowbreak{}LaplacianWeakChain.lean, lines 182–239}
\begin{leancode}
\leanline{0}{theorem\ hilbert\_lpChain\_weak\_step\_of\_bound}
\leanline{4}{(P\ :\ Measure\ Ω)\ [IsProbabilityMeasure\ P]}
\leanline{4}{(U\ :\ H\ →\ ℝ)\ (h\ :\ ℝ)\ (L\ :\ ℝ≥0)\ (hLip\ :\ LipschitzWith\ L\ (gradient\ U))}
\leanline{4}{(initial\ :\ Lp\ H\ 2\ P)\ (W\ :\ ℕ\ →\ Lp\ H\ 2\ P)\ (w\ :\ Ω\ →\ ℝ)\ (hw\ :\ Measurable\ w)}
\leanline{4}{(hind\ :\ IndepFun\ (fun\ a\ ↦\ (w\ a,\ initial\ a))\ (fun\ a\ k\ ↦\ W\ k\ a)\ P)}
\leanline{4}{(hiid\ :\ iIndepFun\ (fun\ k\ a\ ↦\ W\ k\ a)\ P)}
\leanline{4}{(hlaw\ :\ ∀\ k,\ P.map\ (W\ k)\ =\ stdGaussian\ H)}
\leanline{4}{(φ\ :\ H\ →\ ℝ)\ (hφ\ :\ ContDiff\ ℝ\ 2\ φ)}
\leanline{4}{(C\ D\ :\ ℝ)\ (hφC\ :\ ∀\ x,\ |φ\ x|\ ≤\ C)\ (hGD\ :\ ∀\ x,\ |ULA.langevinGenerator\ U\ φ\ x|\ ≤\ D)}
\leanline{4}{(A\ :\ ℝ)\ (hA\ :\ 0\ ≤\ A)\ (hwA\ :\ ∀ᵐ\ a\ ∂P,\ |w\ a|\ ≤\ A)}
\leanline{4}{(ε\ :\ ℝ)\ (hε\ :\ 0\ <\ ε)\ (K\ :\ ℝ)\ (hK\ :\ 0\ ≤\ K)}
\leanline{4}{(hbound\ :\ ∀\ (x\ v\ :\ H)\ (h\ :\ ℝ),\ 0\ ≤\ h\ →\ h\ ≤\ 1\ →}
\leanline{6}{|(∫\ z\ :\ H,\ φ\ (x\ +\ h\ •\ v\ +\ Real.sqrt\ (2\ *\ h)\ •\ z)\ ∂stdGaussian\ H)\ -}
\leanline{8}{φ\ x\ -\ h\ *\ (fderiv\ ℝ\ φ\ x\ v\ +\ Laplacian.laplacian\ φ\ x)|\ ≤}
\leanline{10}{h\ *\ (ε\ +\ K\ *\ Real.sqrt\ h)\ *\ (1\ +\ ‖v‖\ \textasciicircum{}\ 2))}
\leanline{4}{(hh\ :\ 0\ ≤\ h)\ (hh1\ :\ h\ ≤\ 1)\ (k\ :\ ℕ)\ :}
\leanline{4}{|lpWeightedTest\ P\ w\ φ\ (ULA.Hilbert.lpChain\ P\ U\ h\ hLip\ initial\ W\ (k\ +\ 1))\ -}
\leanline{6}{lpWeightedTest\ P\ w\ φ\ (ULA.Hilbert.lpChain\ P\ U\ h\ hLip\ initial\ W\ k)\ -}
\leanline{6}{h\ *\ lpWeightedTest\ P\ w\ (ULA.langevinGenerator\ U\ φ)\ (ULA.Hilbert.lpChain\ P\ U\ h\ hLip\ initial\ W\ k)|\ ≤}
\leanline{6}{A\ *\ h\ *\ (ε\ +\ K\ *\ Real.sqrt\ h)\ *}
\leanline{8}{(1\ +\ ‖ULA.lpDrift\ P\ (gradient\ U)\ hLip\ (ULA.Hilbert.lpChain\ P\ U\ h\ hLip\ initial\ W\ k)‖\ \textasciicircum{}\ 2)}
\end{leancode}
\leanlabel{How the Lean proof follows the argument}
The chain is \texttt{ULA.\allowbreak{}Hilbert.\allowbreak{}lp\allowbreak{}Chain P U h h\allowbreak{}Lip initial W k}, a sequence in the Hilbert space \texttt{Lp H 2 P} of square-integrable classes, defined by recursion with the update \texttt{ULA.\allowbreak{}lp\allowbreak{}Ula\allowbreak{}Update}; \texttt{ULA.\allowbreak{}hilbert\_\allowbreak{}lp\allowbreak{}Chain\_\allowbreak{}coe} states that it agrees almost surely with the pointwise recursion \texttt{ULA.\allowbreak{}noise\allowbreak{}Driven\allowbreak{}Chain (ULA.\allowbreak{}Hilbert.\allowbreak{}pointwise\allowbreak{}Update U h) (initial \ensuremath{\omega}) (fun k \ensuremath{\mapsto} W k \ensuremath{\omega}) n}. Independence is \texttt{hind : Indep\allowbreak{}Fun (fun a \ensuremath{\mapsto} (w a, initial a)) (fun a k \ensuremath{\mapsto} W k a) P} and \texttt{hiid : i\allowbreak{}Indep\allowbreak{}Fun (fun k a \ensuremath{\mapsto} W k a) P}, and \texttt{hlaw} gives the Gaussian laws.



\begin{enumerate}


\item \texttt{hilbert\_\allowbreak{}lp\allowbreak{}Chain\_\allowbreak{}weighted\_\allowbreak{}fresh} is (1). \texttt{Fp} is the augmented update \(\Phi\), measurable by \texttt{ULA.\allowbreak{}Hilbert.\allowbreak{}pointwise\allowbreak{}Update\_\allowbreak{}measurable}; \texttt{hj} is \texttt{noise\allowbreak{}Driven\allowbreak{}Chain\_\allowbreak{}joint\_\allowbreak{}fresh\_\allowbreak{}of\_\allowbreak{}iid} for the augmented chain, and \texttt{noise\allowbreak{}Driven\allowbreak{}Chain\_\allowbreak{}passive} rewrites the augmented chain as \((w,X_k)\). \texttt{hm\ensuremath{{}_{1}}} and \texttt{hm\ensuremath{{}_{2}}} replace the pointwise chain by the \(L^2\) chain inside both laws, using \texttt{Measure.\allowbreak{}map\_\allowbreak{}congr} and the almost-sure identity \texttt{hc}.

\item \texttt{hilbert\_\allowbreak{}lp\allowbreak{}Chain\_\allowbreak{}weak\_\allowbreak{}step\_\allowbreak{}of\_\allowbreak{}bound} is (2), with \texttt{lp\allowbreak{}Weighted\allowbreak{}Test P w \ensuremath{\varphi} X} for \(T(\varphi,X)\) and \texttt{ULA.\allowbreak{}langevin\allowbreak{}Generator U \ensuremath{\varphi}} for \(\mathcal L\varphi\). \texttt{h\allowbreak{}G} is the generator identity, from \texttt{ULA.\allowbreak{}langevin\allowbreak{}Generator\_\allowbreak{}eq\_\allowbreak{}laplacian\_\allowbreak{}sub\_\allowbreak{}inner} and \texttt{inner\_\allowbreak{}gradient\_\allowbreak{}left}; \texttt{h\allowbreak{}Gc} is continuity of the generator; \texttt{hgen} the integrability of \(w\mathcal L\varphi(X_k)\). \texttt{hweak} is Step 1 (3) applied with freshness from part (1) and square integrability from \texttt{ULA.\allowbreak{}lp\allowbreak{}Drift\_\allowbreak{}mem\allowbreak{}Lp}. \texttt{hnext} identifies \(T(\varphi,X_{k+1})\) through \texttt{ULA.\allowbreak{}lp\allowbreak{}Ula\allowbreak{}Update\_\allowbreak{}coe}, and \texttt{hnorm} computes the squared norm of the \(L^2\) drift \texttt{ULA.\allowbreak{}lp\allowbreak{}Drift P (gradient U) h\allowbreak{}Lip X} as \(\mathbb E\|\nabla U(X_k)\|^2\) via \texttt{L2.\allowbreak{}inner\_\allowbreak{}def}.

\end{enumerate}
\end{leanbox}


Earlier results used here: \cite[Section~3.13.11]{LangevinLean2026ULASC}; \cite[Section~3.13.12]{LangevinLean2026ULASC}; \cite[Section~3.13.2]{LangevinLean2026ULASC}; \cite[Section~3.13.3]{LangevinLean2026ULASC}; \cite[Section~3.13.4]{LangevinLean2026ULASC}; \cite[Section~3.19.1]{LangevinLean2026ULASC}; \cite[Section~3.20.1]{LangevinLean2026ULASC}; \cite[Section~3.23.1]{LangevinLean2026ULASC}; \cite[Section~3.23.2]{LangevinLean2026ULASC}; \cite[Section~3.42.1]{LangevinLean2026ULASC}; \cite[Section~3.46.14]{LangevinLean2026ULASC}; \cite[Section~3.46.2]{LangevinLean2026ULASC}; \cite[Section~3.46.8]{LangevinLean2026ULASC}.
\subsection{\texorpdfstring{Summing the residuals over a time horizon}{Summing the residuals over a time horizon}}
\label{report-step-AverageLMC-LaplacianWeakChain-lean-3}


\textbf{Claim.} (1) Let \((a_k)\) and \((c_k)\) be real sequences, \(h\in\mathbb R\), \(E\ge0\) and \(n\in\mathbb N\). If \(|a_{k+1}-a_k-hc_k|\le E\) for every \(k<n\), then


\[\Bigl|a_n-a_0-\sum_{k=0}^{n-1}h\,c_k\Bigr|\le nE.\]


(2) In the setting of Step 2 (2), let \(R\ge0\) with \(\|\nabla U(X_k)\|_{L_2}\le R\) for every \(k\le n\). Then


\[\Bigl|T(\varphi,X_n)-T(\varphi,X_0)-\sum_{k=0}^{n-1}h\,T(\mathcal L\varphi,X_k)\Bigr|\le A\,(nh)\,(\varepsilon+K\sqrt h)\,(1+R^2).\]


\textbf{Proof.} (1) The sum telescopes: \(\sum_{k=0}^{n-1}(a_{k+1}-a_k-hc_k)=a_n-a_0-\sum_{k=0}^{n-1}hc_k\). By the triangle inequality for finite sums,


\[\Bigl|a_n-a_0-\sum_{k<n}hc_k\Bigr|\le\sum_{k<n}|a_{k+1}-a_k-hc_k|\le\sum_{k<n}E=nE.\]


(2) For \(k<n\) we have \(\|\nabla U(X_k)\|_{L_2}^2\le R^2\), so by Step 2 (2) each local residual is at most \(E=Ah(\varepsilon+K\sqrt h)(1+R^2)\), which is nonnegative. Part (1) with \(a_k=T(\varphi,X_k)\) and \(c_k=T(\mathcal L\varphi,X_k)\) gives the bound \(nE=A(nh)(\varepsilon+K\sqrt h)(1+R^2)\). \(\square\)




In the Dynkin identity, \(nh\) is the fixed time horizon \(t\), and the right-hand side tends to \(At\,\varepsilon(1+R^2)\) as \(h\to0\); since \(\varepsilon>0\) is arbitrary, the Euler sums approximate \(\int_0^tT(\mathcal L\varphi,\cdot)\,ds\).

\begin{leanbox}{Lean: \texttt{sum\_\allowbreak{}weighted\_\allowbreak{}weak\_\allowbreak{}residual}, \texttt{hilbert\_\allowbreak{}lpChain\_\allowbreak{}weak\_\allowbreak{}sum\_\allowbreak{}of\_\allowbreak{}bound}}
\leansource{\texttt{sum\_\allowbreak{}weighted\_\allowbreak{}weak\_\allowbreak{}residual} \quad AverageLMC/\allowbreak{}LaplacianWeakChain.lean, lines 132–142}
\begin{leancode}
\leanline{0}{theorem\ sum\_weighted\_weak\_residual\ (a\ c\ :\ ℕ\ →\ ℝ)\ (h\ E\ :\ ℝ)\ (n\ :\ ℕ)}
\leanline{4}{(\_hE\ :\ 0\ ≤\ E)\ (hlocal\ :\ ∀\ k\ <\ n,\ |a\ (k\ +\ 1)\ -\ a\ k\ -\ h\ *\ c\ k|\ ≤\ E)\ :}
\leanline{4}{|a\ n\ -\ a\ 0\ -\ ∑\ k\ ∈\ Finset.range\ n,\ h\ *\ c\ k|\ ≤\ (n\ :\ ℝ)\ *\ E}
\end{leancode}
\leansource{\texttt{hilbert\_\allowbreak{}lpChain\_\allowbreak{}weak\_\allowbreak{}sum\_\allowbreak{}of\_\allowbreak{}bound} \quad AverageLMC/\allowbreak{}LaplacianWeakChain.lean, lines 244–284}
\begin{leancode}
\leanline{0}{theorem\ hilbert\_lpChain\_weak\_sum\_of\_bound}
\leanline{4}{(P\ :\ Measure\ Ω)\ [IsProbabilityMeasure\ P]}
\leanline{4}{(U\ :\ H\ →\ ℝ)\ (h\ :\ ℝ)\ (L\ :\ ℝ≥0)\ (hLip\ :\ LipschitzWith\ L\ (gradient\ U))}
\leanline{4}{(initial\ :\ Lp\ H\ 2\ P)\ (W\ :\ ℕ\ →\ Lp\ H\ 2\ P)\ (w\ :\ Ω\ →\ ℝ)\ (hw\ :\ Measurable\ w)}
\leanline{4}{(hind\ :\ IndepFun\ (fun\ a\ ↦\ (w\ a,\ initial\ a))\ (fun\ a\ k\ ↦\ W\ k\ a)\ P)}
\leanline{4}{(hiid\ :\ iIndepFun\ (fun\ k\ a\ ↦\ W\ k\ a)\ P)}
\leanline{4}{(hlaw\ :\ ∀\ k,\ P.map\ (W\ k)\ =\ stdGaussian\ H)}
\leanline{4}{(φ\ :\ H\ →\ ℝ)\ (hφ\ :\ ContDiff\ ℝ\ 2\ φ)}
\leanline{4}{(C\ D\ :\ ℝ)\ (hφC\ :\ ∀\ x,\ |φ\ x|\ ≤\ C)\ (hGD\ :\ ∀\ x,\ |ULA.langevinGenerator\ U\ φ\ x|\ ≤\ D)}
\leanline{4}{(A\ :\ ℝ)\ (hA\ :\ 0\ ≤\ A)\ (hwA\ :\ ∀ᵐ\ a\ ∂P,\ |w\ a|\ ≤\ A)}
\leanline{4}{(ε\ :\ ℝ)\ (hε\ :\ 0\ <\ ε)\ (K\ :\ ℝ)\ (hK\ :\ 0\ ≤\ K)}
\leanline{4}{(hbound\ :\ ∀\ (x\ v\ :\ H)\ (h\ :\ ℝ),\ 0\ ≤\ h\ →\ h\ ≤\ 1\ →}
\leanline{6}{|(∫\ z\ :\ H,\ φ\ (x\ +\ h\ •\ v\ +\ Real.sqrt\ (2\ *\ h)\ •\ z)\ ∂stdGaussian\ H)\ -}
\leanline{8}{φ\ x\ -\ h\ *\ (fderiv\ ℝ\ φ\ x\ v\ +\ Laplacian.laplacian\ φ\ x)|\ ≤}
\leanline{10}{h\ *\ (ε\ +\ K\ *\ Real.sqrt\ h)\ *\ (1\ +\ ‖v‖\ \textasciicircum{}\ 2))}
\leanline{4}{(hh\ :\ 0\ ≤\ h)\ (hh1\ :\ h\ ≤\ 1)\ (n\ :\ ℕ)\ (R\ :\ ℝ)\ (hR\ :\ 0\ ≤\ R)}
\leanline{4}{(hboundR\ :\ ∀\ k\ ≤\ n,}
\leanline{6}{‖ULA.lpDrift\ P\ (gradient\ U)\ hLip\ (ULA.Hilbert.lpChain\ P\ U\ h\ hLip\ initial\ W\ k)‖\ ≤\ R)\ :}
\leanline{4}{|lpWeightedTest\ P\ w\ φ\ (ULA.Hilbert.lpChain\ P\ U\ h\ hLip\ initial\ W\ n)\ -}
\leanline{6}{lpWeightedTest\ P\ w\ φ\ initial\ -\ ∑\ k\ ∈\ Finset.range\ n,}
\leanline{8}{h\ *\ lpWeightedTest\ P\ w\ (ULA.langevinGenerator\ U\ φ)}
\leanline{10}{(ULA.Hilbert.lpChain\ P\ U\ h\ hLip\ initial\ W\ k)|\ ≤}
\leanline{6}{A\ *\ ((n\ :\ ℝ)\ *\ h)\ *\ (ε\ +\ K\ *\ Real.sqrt\ h)\ *\ (1\ +\ R\ \textasciicircum{}\ 2)}
\end{leancode}
\leanlabel{How the Lean proof follows the argument}
\begin{enumerate}


\item \texttt{sum\_\allowbreak{}weighted\_\allowbreak{}weak\_\allowbreak{}residual} has the hypothesis \texttt{hlocal : \ensuremath{\forall} k < n, |a (k + 1) -\allowbreak{} a k -\allowbreak{} h * c k| \ensuremath{\le} E} and an unused sign condition \texttt{\_\allowbreak{}h\allowbreak{}E}. \texttt{ht} is the telescoping identity (\texttt{Finset.\allowbreak{}sum\_\allowbreak{}sub\_\allowbreak{}distrib} and \texttt{Finset.\allowbreak{}sum\_\allowbreak{}range\_\allowbreak{}sub}), and the \texttt{calc} applies \texttt{Finset.\allowbreak{}abs\_\allowbreak{}sum\_\allowbreak{}le\_\allowbreak{}sum\_\allowbreak{}abs} and \texttt{Finset.\allowbreak{}sum\_\allowbreak{}le\_\allowbreak{}sum}.

\item \texttt{hilbert\_\allowbreak{}lp\allowbreak{}Chain\_\allowbreak{}weak\_\allowbreak{}sum\_\allowbreak{}of\_\allowbreak{}bound} adds \texttt{h\allowbreak{}R : 0 \ensuremath{\le} R} and \texttt{hbound\allowbreak{}R}, the bound on the \(L^2\) drift norms for \(k\le n\). \texttt{hlocal} combines \texttt{hilbert\_\allowbreak{}lp\allowbreak{}Chain\_\allowbreak{}weak\_\allowbreak{}step\_\allowbreak{}of\_\allowbreak{}bound} with the squared bound \texttt{hs} (\texttt{sq\_\allowbreak{}le\_\allowbreak{}sq\ensuremath{{}_{0}}}) and \texttt{mul\_\allowbreak{}le\_\allowbreak{}mul\_\allowbreak{}of\_\allowbreak{}nonneg\_\allowbreak{}left}; part (1) is applied as \texttt{hb}, and the final \texttt{simpa} rearranges \(n\cdot Ah(\cdots)\) as \(A(nh)(\cdots)\).

\end{enumerate}
\end{leanbox}


Earlier results used here: \cite[Section~3.13.2]{LangevinLean2026ULASC}; \cite[Section~3.19.1]{LangevinLean2026ULASC}; \cite[Section~3.46.2]{LangevinLean2026ULASC}.
\section{\texorpdfstring{Weighted Dynkin identity for the constructed Brownian Langevin solution}{Weighted Dynkin identity for the constructed Brownian Langevin solution}}
\label{proof-module-AverageLMC-LaplacianBrownianDynkin-lean}


This module completes the stochastic bridge for compact \(C^{2}\) tests. All formulas are derived from the actual Brownian innovations, the continuous \(L^{2}\) integral solution, and finite Gaussian Euler chains.



The actual Brownian Euler chains have an eventual uniform score bound on every time partition:


\[\exists R\ge0\ \exists N\ \forall n\ge N\ \forall k\le n+1:
\quad\|\nabla f(E_k^{(n)})\|_{L^2(P)}\le R.\]


Actual Gaussian weak consistency and the mesh limit prove, for bounded initial weights independent of the innovations and \(\phi\in C_c^2\),


\[\mathbb E[w\phi(Y_T)]-\mathbb E[w\phi(Y_0)]
=\int_0^T\mathbb E[w\mathcal L_f\phi(Y_t)]\,dt.\]

Earlier results used here: \cite[Section~3.42]{LangevinLean2026ULASC}.
\subsection{\texorpdfstring{Bound all actual Brownian Euler scores}{Bound all actual Brownian Euler scores}}
\label{report-step-AverageLMC-LaplacianBrownianDynkin-lean-1}


Use the deterministic uniform grid bound in \(L^{2}\), with the lifted gradient as drift and the actual Brownian noise path. The established equality between additive matching-noise Euler states and the Euler chain in the Hilbert space of square-integrable random vectors identifies the objects used in the probabilistic weak-step proof.

\begin{leanbox}{Lean: \texttt{brownian\_\allowbreak{}euler\_\allowbreak{}score\_\allowbreak{}grid\_\allowbreak{}eventually\_\allowbreak{}bounded}}
\leansource{\texttt{brownian\_\allowbreak{}euler\_\allowbreak{}score\_\allowbreak{}grid\_\allowbreak{}eventually\_\allowbreak{}bounded} \quad AverageLMC/\allowbreak{}LaplacianBrownianDynkin.lean, lines 28–48}
\begin{leancode}
\leanline{0}{theorem\ brownian\_euler\_score\_grid\_eventually\_bounded}
\leanline{4}{(f\ :\ EuclideanSpace\ ℝ\ ι\ →\ ℝ)\ (L\ :\ ℝ≥0)\ (hLip\ :\ LipschitzWith\ L\ (gradient\ f))}
\leanline{4}{(initial\ :\ Lp\ (EuclideanSpace\ ℝ\ ι)\ 2\ P)}
\leanline{4}{(B\ :\ ι\ →\ ℝ≥0\ →\ Ω\ →\ ℝ)\ (hB\ :\ ∀\ i,\ IsPreBrownianReal\ (B\ i)\ P)}
\leanline{4}{(Y\ :\ ℝ\ →\ Lp\ (EuclideanSpace\ ℝ\ ι)\ 2\ P)\ (hY\ :\ Continuous\ Y)}
\leanline{4}{(T\ :\ ℝ)\ (hT\ :\ 0\ <\ T)}
\leanline{4}{(hYe\ :\ ∀\ t\ ∈\ Icc\ 0\ T,\ Y\ t\ -\ initial\ =}
\leanline{6}{-(∫\ s\ in\ 0..t,\ ULA.lpDrift\ P\ (gradient\ f)\ hLip\ (Y\ s))\ +}
\leanline{8}{ULA.brownianLangevinNoise\ P\ B\ hB\ 0\ t)\ :}
\leanline{4}{∃\ R\ :\ ℝ,\ 0\ ≤\ R\ ∧\ ∃\ N\ :\ ℕ,\ ∀\ n\ ≥\ N,\ ∀\ k\ ≤\ n\ +\ 1,}
\leanline{6}{‖ULA.lpDrift\ P\ (gradient\ f)\ hLip}
\leanline{8}{(ULA.Hilbert.lpChain\ P\ f\ (T\ /\ ((n\ +\ 1\ :\ ℕ)\ :\ ℝ))\ hLip\ initial}
\leanline{10}{(ULA.Hilbert.brownianEulerNoise\ P\ B\ hB\ (T\ /\ ((n\ +\ 1\ :\ ℕ)\ :\ ℝ)))\ k)‖\ ≤\ R}
\end{leancode}
\leanlabel{How the Lean proof follows the argument}
\texttt{brownian\_\allowbreak{}euler\_\allowbreak{}score\_\allowbreak{}grid\_\allowbreak{}eventually\_\allowbreak{}bounded} gives the formal counterpart of this step. Named proof applications include \texttt{h\allowbreak{}N}; the adjacent proof shows their hypotheses and order of use.
\end{leanbox}


Earlier results used here: \cite[Section~3.13.2]{LangevinLean2026ULASC}; \cite[Section~3.13.5]{LangevinLean2026ULASC}; \cite[Section~3.19.1]{LangevinLean2026ULASC}; \cite[Section~3.25.1]{LangevinLean2026ULASC}; \cite[Section~3.25.3]{LangevinLean2026ULASC}; \cite[Section~3.29.7]{LangevinLean2026ULASC}; \cite[Section~3.42.5]{LangevinLean2026ULASC}.
\subsection{\texorpdfstring{Derive weighted Dynkin without an assumed formula}{Derive weighted Dynkin without an assumed formula}}
\label{report-step-AverageLMC-LaplacianBrownianDynkin-lean-2}


For positive \(T\), actual Brownian coordinate increments give independent standard Gaussian vector innovations on every refining partition. Keep the bounded initial weight paired with the initialization, so each next innovation is fresh. Apply compact \(C^{2}\) Gaussian weak consistency, sum its residual using the uniform score moment, and pass to the limit with bounded uniformly continuous test and generator observations. The zero-time case follows directly from the integral equation. The conclusion is \(\mathbb E[w\varphi(Y_T)]-\mathbb E[w\varphi(Y_0)]=\int_0^T\mathbb E[w\mathcal L\varphi(Y_t)]\,dt\).

\begin{leanbox}{Lean: \texttt{lp\_\allowbreak{}langevin\_\allowbreak{}weighted\_\allowbreak{}dynkin\_\allowbreak{}compact}}
\leansource{\texttt{lp\_\allowbreak{}langevin\_\allowbreak{}weighted\_\allowbreak{}dynkin\_\allowbreak{}compact} \quad AverageLMC/\allowbreak{}LaplacianBrownianDynkin.lean, lines 54–110}
\begin{leancode}
\leanline{0}{theorem\ lp\_langevin\_weighted\_dynkin\_compact}
\leanline{4}{(f\ :\ EuclideanSpace\ ℝ\ ι\ →\ ℝ)\ (L\ :\ ℝ≥0)\ (hLip\ :\ LipschitzWith\ L\ (gradient\ f))}
\leanline{4}{(initial\ :\ Lp\ (EuclideanSpace\ ℝ\ ι)\ 2\ P)}
\leanline{4}{(B\ :\ ι\ →\ ℝ≥0\ →\ Ω\ →\ ℝ)\ (hB\ :\ ∀\ i,\ IsPreBrownianReal\ (B\ i)\ P)}
\leanline{4}{(w\ :\ Ω\ →\ ℝ)\ (hw\ :\ Measurable\ w)\ (A\ :\ ℝ)\ (hA\ :\ 0\ ≤\ A)}
\leanline{4}{(hwA\ :\ ∀ᵐ\ a\ ∂P,\ |w\ a|\ ≤\ A)}
\leanline{4}{(hscalar\ :\ ∀\ r\ :\ ℝ,\ 0\ <\ r\ →\ iIndepFun\ (fun\ q\ :\ ℕ\ ×\ ι\ ↦\ fun\ a\ ↦}
\leanline{6}{ULA.Hilbert.brownianEulerNoise\ P\ B\ hB\ r\ q.1\ a\ q.2)\ P)}
\leanline{4}{(hind\ :\ ∀\ r\ :\ ℝ,\ 0\ <\ r\ →\ IndepFun\ (fun\ a\ ↦\ (w\ a,\ initial\ a))}
\leanline{6}{(fun\ a\ k\ ↦\ ULA.Hilbert.brownianEulerNoise\ P\ B\ hB\ r\ k\ a)\ P)}
\leanline{4}{(Y\ :\ ℝ\ →\ Lp\ (EuclideanSpace\ ℝ\ ι)\ 2\ P)\ (hY\ :\ Continuous\ Y)}
\leanline{4}{(T\ :\ ℝ)\ (hT\ :\ 0\ ≤\ T)}
\leanline{4}{(hYe\ :\ ∀\ t\ ∈\ Icc\ 0\ T,\ Y\ t\ -\ initial\ =}
\leanline{6}{-(∫\ s\ in\ 0..t,\ ULA.lpDrift\ P\ (gradient\ f)\ hLip\ (Y\ s))\ +}
\leanline{8}{ULA.brownianLangevinNoise\ P\ B\ hB\ 0\ t)}
\leanline{4}{(φ\ :\ EuclideanSpace\ ℝ\ ι\ →\ ℝ)\ (hφ\ :\ ContDiff\ ℝ\ 2\ φ)\ (hc\ :\ HasCompactSupport\ φ)\ :}
\leanline{4}{lpWeightedTest\ P\ w\ φ\ (Y\ T)\ -\ lpWeightedTest\ P\ w\ φ\ initial\ =}
\leanline{6}{∫\ t\ in\ 0..T,\ lpWeightedTest\ P\ w\ (ULA.langevinGenerator\ f\ φ)\ (Y\ t)}
\end{leancode}
\leanlabel{How the Lean proof follows the argument}
\texttt{lp\_\allowbreak{}langevin\_\allowbreak{}weighted\_\allowbreak{}dynkin\_\allowbreak{}compact} separates the cases \texttt{h\allowbreak{}Tzero : T = 0} and its negation before combining them. The main intermediate facts are \texttt{h\allowbreak{}Tpos}, \texttt{htest}, \texttt{hgen}, \texttt{hmesh}, \texttt{h}, \texttt{hh} in that order. Named proof applications include \texttt{weighted\_\allowbreak{}dynkin\_\allowbreak{}of\_\allowbreak{}weak\_\allowbreak{}euler\_\allowbreak{}consistency}; the adjacent proof shows their hypotheses and order of use.
\end{leanbox}


Earlier results used here: \cite[Section~3.13.2]{LangevinLean2026ULASC}; \cite[Section~3.13.5]{LangevinLean2026ULASC}; \cite[Section~3.25.1]{LangevinLean2026ULASC}; \cite[Section~3.25.3]{LangevinLean2026ULASC}; \cite[Section~3.29.7]{LangevinLean2026ULASC}; \cite[Section~3.42.5]{LangevinLean2026ULASC}; \cite[Section~3.42.6]{LangevinLean2026ULASC}; \cite[Section~3.46.2]{LangevinLean2026ULASC}.
\section{\texorpdfstring{Bounded radial weights approximate every square-integrable initial vector}{Bounded radial weights approximate every square-integrable initial vector}}
\label{proof-module-AverageLMC-LaplacianWeightClipping-lean}


The initial coupling error in the Laplacian proof is only square integrable. This module constructs bounded radial clippings whose \(L^{2}\) norms never exceed the original norm. Their pairings converge against every fixed \(L^{2}\) vector by dominated convergence, allowing the later compact-test argument to use bounded scalar weights without adding boundedness to the main theorem hypotheses.



The measurable radial clipping has \(c_n(0)=0\); for \(x\ne0\) it uses radius \(n+1\):


\[c_n(x)=\min\!\left\{1,\frac{n+1}{\|x\|}\right\}x,
\qquad\|c_n(x)\|\le\min\{\|x\|,n+1\}.\]


Clipping never increases the actual \(L^2\) norm and converges against every square-integrable vector:


\[\|c_n(\Psi)\|_{L^2(P)}\le\|\Psi\|_{L^2(P)},\qquad
\langle c_n(\Psi),D\rangle_{L^2(P)}
\longrightarrow\langle\Psi,D\rangle_{L^2(P)}.\]\subsection{\texorpdfstring{Define radial clipping and prove both norm bounds}{Define radial clipping and prove both norm bounds}}
\label{report-step-AverageLMC-LaplacianWeightClipping-lean-1}


Clip a vector \(x\) by multiplying it by \(\min(1,(n+1)/\|x\|)\). The coefficient lies between zero and one, so the clipped norm is at most \(\|x\|\). It is also at most \(n+1\): at a nonzero vector cancel the positive norm against the radial quotient, and handle the zero vector directly. These two bounds serve different purposes: the radius makes every coordinate a bounded test weight, while domination by the original vector supplies the integrable majorant and preserves the sharp final \(L^{2}\) norm.

\begin{leanbox}{Lean: \texttt{clipWeight}, \texttt{norm\_\allowbreak{}clipWeight\_\allowbreak{}le}, \texttt{norm\_\allowbreak{}clipWeight\_\allowbreak{}le\_\allowbreak{}radius}}
\leansource{\texttt{clipWeight} \quad AverageLMC/\allowbreak{}LaplacianWeightClipping.lean, lines 28–29}
\begin{leancode}
\leanline{0}{def\ clipWeight\ (n\ :\ ℕ)\ (x\ :\ H)\ :\ H\ :=}
\leanline{2}{min\ 1\ (((n\ :\ ℝ)\ +\ 1)\ /\ ‖x‖)\ •\ x}
\end{leancode}
\leansource{\texttt{norm\_\allowbreak{}clipWeight\_\allowbreak{}le} \quad AverageLMC/\allowbreak{}LaplacianWeightClipping.lean, lines 31–35}
\begin{leancode}
\leanline{0}{theorem\ norm\_clipWeight\_le\ (n\ :\ ℕ)\ (x\ :\ H)\ :\ ‖clipWeight\ n\ x‖\ ≤\ ‖x‖}
\end{leancode}
\leansource{\texttt{norm\_\allowbreak{}clipWeight\_\allowbreak{}le\_\allowbreak{}radius} \quad AverageLMC/\allowbreak{}LaplacianWeightClipping.lean, lines 37–47}
\begin{leancode}
\leanline{0}{theorem\ norm\_clipWeight\_le\_radius\ (n\ :\ ℕ)\ (x\ :\ H)\ :\ ‖clipWeight\ n\ x‖\ ≤\ (n\ :\ ℝ)\ +\ 1}
\end{leancode}
\leanlabel{How the Lean proof follows the argument}
\texttt{clip\allowbreak{}Weight} names the mathematical object defined above. Its defining expression appears in the adjacent statement; subsequent identities may unfold this definition directly.




\texttt{norm\_\allowbreak{}clip\allowbreak{}Weight\_\allowbreak{}le} gives the formal counterpart of this step. The main intermediate facts are \texttt{hc} in that order. Named proof applications include \texttt{mul\_\allowbreak{}le\_\allowbreak{}mul\_\allowbreak{}of\_\allowbreak{}nonneg\_\allowbreak{}right}; the adjacent proof shows their hypotheses and order of use.




\texttt{norm\_\allowbreak{}clip\allowbreak{}Weight\_\allowbreak{}le\_\allowbreak{}radius} separates the cases \texttt{hx : x = 0} and its negation before combining them. The main intermediate facts are \texttt{hn}, \texttt{hc} in that order. Its remaining equalities are discharged by the simplifications displayed in the proof.
\end{leanbox}
\subsection{\texorpdfstring{Obtain eventual equality and measurability}{Obtain eventual equality and measurability}}
\label{report-step-AverageLMC-LaplacianWeightClipping-lean-2}


If \(\|x\|\) is at most \(n+1\), the minimum coefficient is one and clipping fixes \(x\). For a fixed vector choose a natural number above its norm. Every larger index then fixes that vector exactly, giving eventual pointwise equality with no approximation-rate premise. Norm, real division, minimum and scalar multiplication are measurable, so the clipping map is measurable even with the real quotient convention at zero.

\begin{leanbox}{Lean: \texttt{clipWeight\_\allowbreak{}eq\_\allowbreak{}of\_\allowbreak{}norm\_\allowbreak{}le}, \texttt{clipWeight\_\allowbreak{}eventually\_\allowbreak{}eq}, \texttt{measurable\_\allowbreak{}clipWeight}}
\leansource{\texttt{clipWeight\_\allowbreak{}eq\_\allowbreak{}of\_\allowbreak{}norm\_\allowbreak{}le} \quad AverageLMC/\allowbreak{}LaplacianWeightClipping.lean, lines 49–55}
\begin{leancode}
\leanline{0}{theorem\ clipWeight\_eq\_of\_norm\_le\ (n\ :\ ℕ)\ (x\ :\ H)\ (hx\ :\ ‖x‖\ ≤\ (n\ :\ ℝ)\ +\ 1)\ :}
\leanline{4}{clipWeight\ n\ x\ =\ x}
\end{leancode}
\leansource{\texttt{clipWeight\_\allowbreak{}eventually\_\allowbreak{}eq} \quad AverageLMC/\allowbreak{}LaplacianWeightClipping.lean, lines 57–62}
\begin{leancode}
\leanline{0}{theorem\ clipWeight\_eventually\_eq\ (x\ :\ H)\ :\ ∀ᶠ\ n\ :\ ℕ\ in\ atTop,\ clipWeight\ n\ x\ =\ x}
\end{leancode}
\leansource{\texttt{measurable\_\allowbreak{}clipWeight} \quad AverageLMC/\allowbreak{}LaplacianWeightClipping.lean, lines 64–67}
\begin{leancode}
\leanline{0}{theorem\ measurable\_clipWeight\ [MeasurableSpace\ H]\ [BorelSpace\ H]\ (n\ :\ ℕ)\ :}
\leanline{4}{Measurable\ (@clipWeight\ H\ \_\ \_\ n)}
\end{leancode}
\leanlabel{How the Lean proof follows the argument}
\texttt{clip\allowbreak{}Weight\_\allowbreak{}eq\_\allowbreak{}of\_\allowbreak{}norm\_\allowbreak{}le} separates the cases \texttt{hx0 : x = 0} and its negation before combining them. The main intermediate facts are \texttt{hn}, \texttt{hc} in that order. Its remaining equalities are discharged by the simplifications displayed in the proof.




\texttt{clip\allowbreak{}Weight\_\allowbreak{}eventually\_\allowbreak{}eq} gives the formal counterpart of this step. The main intermediate facts are \texttt{hcast} in that order. Named proof applications include \texttt{clip\allowbreak{}Weight\_\allowbreak{}eq\_\allowbreak{}of\_\allowbreak{}norm\_\allowbreak{}le}; the adjacent proof shows their hypotheses and order of use.




\texttt{measurable\_\allowbreak{}clip\allowbreak{}Weight} gives the formal counterpart of this step.
\end{leanbox}
\subsection{\texorpdfstring{Lift the clippings into the genuine \(L^{2}\) space}{Lift the clippings into the genuine L² space}}
\label{report-step-AverageLMC-LaplacianWeightClipping-lean-3}


Compose the measurable clipping with a square-integrable representative. The pointwise norm domination by the original representative proves square integrability. The resulting clipped weight has an almost-everywhere representative equal to this composition. Monotonicity of the square-integrable norm under the same pointwise domination gives a clipped norm at most the original norm, uniformly over the clipping threshold.

\begin{leanbox}{Lean: \texttt{clippedWeight\_\allowbreak{}memLp}, \texttt{lpClippedWeight}, \texttt{lpClippedWeight\_\allowbreak{}coe} and 1 more}
\leansource{\texttt{clippedWeight\_\allowbreak{}memLp} \quad AverageLMC/\allowbreak{}LaplacianWeightClipping.lean, lines 78–82}
\begin{leancode}
\leanline{0}{theorem\ clippedWeight\_memLp\ (X\ :\ Lp\ H\ 2\ P)\ (n\ :\ ℕ)\ :}
\leanline{4}{MemLp\ (fun\ a\ ↦\ clipWeight\ n\ (X\ a))\ 2\ P}
\end{leancode}
\leansource{\texttt{lpClippedWeight} \quad AverageLMC/\allowbreak{}LaplacianWeightClipping.lean, lines 84–85}
\begin{leancode}
\leanline{0}{def\ lpClippedWeight\ (X\ :\ Lp\ H\ 2\ P)\ (n\ :\ ℕ)\ :\ Lp\ H\ 2\ P\ :=}
\leanline{2}{(clippedWeight\_memLp\ P\ X\ n).toLp\ (fun\ a\ ↦\ clipWeight\ n\ (X\ a))}
\end{leancode}
\leansource{\texttt{lpClippedWeight\_\allowbreak{}coe} \quad AverageLMC/\allowbreak{}LaplacianWeightClipping.lean, lines 87–89}
\begin{leancode}
\leanline{0}{theorem\ lpClippedWeight\_coe\ (X\ :\ Lp\ H\ 2\ P)\ (n\ :\ ℕ)\ :}
\leanline{4}{lpClippedWeight\ P\ X\ n\ =ᵐ[P]\ fun\ a\ ↦\ clipWeight\ n\ (X\ a)}
\end{leancode}
\leansource{\texttt{lpClippedWeight\_\allowbreak{}norm\_\allowbreak{}le} \quad AverageLMC/\allowbreak{}LaplacianWeightClipping.lean, lines 91–97}
\begin{leancode}
\leanline{0}{theorem\ lpClippedWeight\_norm\_le\ (X\ :\ Lp\ H\ 2\ P)\ (n\ :\ ℕ)\ :\ ‖lpClippedWeight\ P\ X\ n‖\ ≤\ ‖X‖}
\end{leancode}
\leanlabel{How the Lean proof follows the argument}
\textbf{Formal correspondence.} The \texttt{Mem\allowbreak{}Lp} proof at exponent two defines \texttt{lp\allowbreak{}Clipped\allowbreak{}Weight} through \texttt{to\allowbreak{}Lp}, with its representative identity proved explicitly.



\texttt{clipped\allowbreak{}Weight\_\allowbreak{}mem\allowbreak{}Lp} gives the formal counterpart of this step.




\texttt{lp\allowbreak{}Clipped\allowbreak{}Weight} names the mathematical object defined above. Its defining expression appears in the adjacent statement; subsequent identities may unfold this definition directly.




\texttt{lp\allowbreak{}Clipped\allowbreak{}Weight\_\allowbreak{}coe} gives the formal counterpart of this step.




\texttt{lp\allowbreak{}Clipped\allowbreak{}Weight\_\allowbreak{}norm\_\allowbreak{}le} gives the formal counterpart of this step. The main intermediate facts are \texttt{hbound} in that order. Named proof applications include \texttt{Lp.\allowbreak{}norm\_\allowbreak{}le\_\allowbreak{}mul\_\allowbreak{}norm\_\allowbreak{}of\_\allowbreak{}ae\_\allowbreak{}le\_\allowbreak{}mul}, \texttt{norm\_\allowbreak{}clip\allowbreak{}Weight\_\allowbreak{}le}, \texttt{hbound}; the adjacent proof shows their hypotheses and order of use.
\end{leanbox}
\subsection{\texorpdfstring{Pass every fixed \(L^{2}\) pairing through the clipping limit}{Pass every fixed L² pairing through the clipping limit}}
\label{report-step-AverageLMC-LaplacianWeightClipping-lean-4}


For \(X\) and \(D\) in \(L^{2}\), Holder gives integrability of \(\|X\|\) times \(\|D\|\). Cauchy-Schwarz bounds the absolute pointwise inner product of the clipping of \(X\) with \(D\) by this common majorant. Eventual pointwise equality gives pointwise convergence of the integrands, and dominated convergence gives convergence of their actual integrals. Rewriting the Hilbert-space \(L^{2}\) inner product as that integral and using the almost-everywhere clipping representative proves convergence of the exact \(L^{2}\) pairings.

\begin{leanbox}{Lean: \texttt{lpClippedWeight\_\allowbreak{}inner\_\allowbreak{}tendsto}}
\leansource{\texttt{lpClippedWeight\_\allowbreak{}inner\_\allowbreak{}tendsto} \quad AverageLMC/\allowbreak{}LaplacianWeightClipping.lean, lines 101–129}
\begin{leancode}
\leanline{0}{theorem\ lpClippedWeight\_inner\_tendsto\ (X\ D\ :\ Lp\ H\ 2\ P)\ :}
\leanline{4}{Tendsto\ (fun\ n\ ↦\ inner\ ℝ\ (lpClippedWeight\ P\ X\ n)\ D)\ atTop\ ({\leanfallback 𝓝}\ (inner\ ℝ\ X\ D))}
\end{leancode}
\leanlabel{How the Lean proof follows the argument}
\texttt{lp\allowbreak{}Clipped\allowbreak{}Weight\_\allowbreak{}inner\_\allowbreak{}tendsto} gives the formal counterpart of this step. The main intermediate facts are \texttt{hbound}, \texttt{hmeas}, \texttt{hdom}, \texttt{hpoint}, \texttt{hd}, \texttt{hi} in that order. Named proof applications include \texttt{ae\_\allowbreak{}of\_\allowbreak{}all}, \texttt{integral\_\allowbreak{}congr\_\allowbreak{}ae}, \texttt{hd}; the adjacent proof shows their hypotheses and order of use.
\end{leanbox}
\section{\texorpdfstring{Pass proved compact Dynkin identities to the sharp coupling pairing}{Pass proved compact Dynkin identities to the sharp coupling pairing}}
\label{proof-module-AverageLMC-LaplacianPairingLimit-lean}


These are analytic limit principles used after the stochastic identities have been proved. The first theorem assumes explicit identities for compact approximations and a uniform-in-time generator norm bound, and passes their endpoint and generator errors to zero. The second removes boundedness of the initial vector by the independently proved clipping limit. Neither principle claims a sampling conclusion without its displayed premises.



Suppose \(X_n\to X\), \(Y_n\to Y\), \(\epsilon_n\to0\), and the proved compact identities satisfy


\[\langle\Psi,X_n-Y_n\rangle
=\int_0^T\langle\Psi,G_n(t)\rangle\,dt,\qquad
\sup_{0\le t\le T}\|G_n(t)\|\le R+\epsilon_n.\]


The integral norm bound and passage to the limit yield, also after removing boundedness of the weight by radial clipping,


\[\langle\Psi,X-Y\rangle\le T\|\Psi\|R.\]\subsection{\texorpdfstring{Bound each compact approximation before taking its limit}{Bound each compact approximation before taking its limit}}
\label{report-step-AverageLMC-LaplacianPairingLimit-lean-1}


Suppose the pairing of a fixed vector \(\Psi\) with the difference of two endpoint approximations equals the time integral of its pairing with a generator approximation. If every generator norm on [0,\(T\)] is at most \(R+\operatorname{err}(n)\), Cauchy-Schwarz bounds the integrand by \(\|\Psi\|\) times that constant. The interval-integral norm bound then gives an endpoint pairing at most \(T\) \(\|\Psi\|\)\((R+err(n))\). Continuity of the inner product passes the convergent endpoint sequences to their limits. Since \(\operatorname{err}(n)\) tends to zero, closedness of the real order gives the exact bound \(T\) \(\|\Psi\|\) \(R\). This requires no continuity of the eventual unbounded generator path.

\begin{leanbox}{Lean: \texttt{pairing\_\allowbreak{}bound\_\allowbreak{}of\_\allowbreak{}compact\_\allowbreak{}dynkin\_\allowbreak{}approximations}}
\leansource{\texttt{pairing\_\allowbreak{}bound\_\allowbreak{}of\_\allowbreak{}compact\_\allowbreak{}dynkin\_\allowbreak{}approximations} \quad AverageLMC/\allowbreak{}LaplacianPairingLimit.lean, lines 27–54}
\begin{leancode}
\leanline{0}{theorem\ pairing\_bound\_of\_compact\_dynkin\_approximations}
\leanline{4}{(Ψ\ X\ Y\ :\ H)\ (Xn\ Yn\ :\ ℕ\ →\ H)\ (Gn\ :\ ℕ\ →\ ℝ\ →\ H)}
\leanline{4}{(T\ R\ :\ ℝ)\ (hT\ :\ 0\ ≤\ T)\ (hR\ :\ 0\ ≤\ R)\ (err\ :\ ℕ\ →\ ℝ)}
\leanline{4}{(herr\ :\ ∀\ n,\ 0\ ≤\ err\ n)\ (herrlim\ :\ Tendsto\ err\ atTop\ ({\leanfallback 𝓝}\ 0))}
\leanline{4}{(hX\ :\ Tendsto\ Xn\ atTop\ ({\leanfallback 𝓝}\ X))\ (hY\ :\ Tendsto\ Yn\ atTop\ ({\leanfallback 𝓝}\ Y))}
\leanline{4}{(hgen\ :\ ∀\ n\ t,\ t\ ∈\ Icc\ 0\ T\ →\ ‖Gn\ n\ t‖\ ≤\ R\ +\ err\ n)}
\leanline{4}{(hidentity\ :\ ∀\ n,\ inner\ ℝ\ Ψ\ (Xn\ n\ -\ Yn\ n)\ =}
\leanline{6}{∫\ t\ in\ 0..T,\ inner\ ℝ\ Ψ\ (Gn\ n\ t))\ :}
\leanline{4}{inner\ ℝ\ Ψ\ (X\ -\ Y)\ ≤\ T\ *\ ‖Ψ‖\ *\ R}
\end{leancode}
\leanlabel{How the Lean proof follows the argument}
\texttt{pairing\_\allowbreak{}bound\_\allowbreak{}of\_\allowbreak{}compact\_\allowbreak{}dynkin\_\allowbreak{}approximations} gives the formal counterpart of this step. The main intermediate facts are \texttt{hbound}, \texttt{hleft}, \texttt{hright} in that order. Named proof applications include \texttt{le\_\allowbreak{}abs\_\allowbreak{}self}, \texttt{hab.\allowbreak{}trans}, \texttt{le\_\allowbreak{}of\_\allowbreak{}tendsto\_\allowbreak{}of\_\allowbreak{}tendsto}; the adjacent proof shows their hypotheses and order of use.
\end{leanbox}
\subsection{\texorpdfstring{Remove the bounded-weight restriction without increasing the norm}{Remove the bounded-weight restriction without increasing the norm}}
\label{report-step-AverageLMC-LaplacianPairingLimit-lean-2}


Apply the assumed inequality to every radial clipping of the \(L^{2}\) vector \(\Psi\). Each clipped norm is at most \(\|\Psi\|\), so every right-hand side is bounded by the same \(T\) \(\|\Psi\|\) \(R\). The prior clipping theorem makes the left-hand pairings converge to the pairing of \(\Psi\) itself with the fixed \(L^{2}\) increment. Passing the real inequality to the limit gives the desired estimate for arbitrary square-integrable initial vectors, with no coordinate or dimension factor.

\begin{leanbox}{Lean: \texttt{pairing\_\allowbreak{}bound\_\allowbreak{}of\_\allowbreak{}clipped\_\allowbreak{}weights}}
\leansource{\texttt{pairing\_\allowbreak{}bound\_\allowbreak{}of\_\allowbreak{}clipped\_\allowbreak{}weights} \quad AverageLMC/\allowbreak{}LaplacianPairingLimit.lean, lines 68–76}
\begin{leancode}
\leanline{0}{theorem\ pairing\_bound\_of\_clipped\_weights}
\leanline{4}{(Ψ\ D\ :\ Lp\ H\ 2\ P)\ (T\ R\ :\ ℝ)\ (hT\ :\ 0\ ≤\ T)\ (hR\ :\ 0\ ≤\ R)}
\leanline{4}{(hclip\ :\ ∀\ n,\ inner\ ℝ\ (lpClippedWeight\ P\ Ψ\ n)\ D\ ≤}
\leanline{6}{T\ *\ ‖lpClippedWeight\ P\ Ψ\ n‖\ *\ R)\ :}
\leanline{4}{inner\ ℝ\ Ψ\ D\ ≤\ T\ *\ ‖Ψ‖\ *\ R}
\end{leancode}
\leanlabel{How the Lean proof follows the argument}
\texttt{pairing\_\allowbreak{}bound\_\allowbreak{}of\_\allowbreak{}clipped\_\allowbreak{}weights} gives the formal counterpart of this step. Named proof applications include \texttt{le\_\allowbreak{}of\_\allowbreak{}tendsto}, \texttt{Eventually.\allowbreak{}of\_\allowbreak{}forall}, \texttt{mul\_\allowbreak{}le\_\allowbreak{}mul\_\allowbreak{}of\_\allowbreak{}nonneg\_\allowbreak{}right}; the adjacent proof shows their hypotheses and order of use.
\end{leanbox}
\section{\texorpdfstring{Transfer Gibbs observations to the actual stationary process}{Transfer Gibbs observations to the actual stationary process}}
\label{proof-module-AverageLMC-LaplacianStationaryTest-lean}


Stationarity is an equality of actual marginal laws. This module transfers genuine \(L^{2}\) moments and approximation errors from the Gibbs measure to the constructed process, and converts finite coordinate Dynkin identities into vector inner products.

\[\mathcal L(Y)=\pi\quad\Longrightarrow\quad\|F(Y)-G(Y)\|_{\mathbb L_2}^2=\int\|F-G\|^2\,d\pi,\quad \langle\Psi,F(Y)\rangle_{\mathbb L_2}=\sum_j\mathbb E[\Psi_jF_j(Y)].\]\subsection{\texorpdfstring{Represent every square-integrable stationary observation}{Represent every square-integrable stationary observation}}
\label{report-step-AverageLMC-LaplacianStationaryTest-lean-1}


Rewrite the observation's Gibbs moment with the actual marginal map and pull it back to the original probability space. The resulting square-integrability estimate defines a genuine observation in the space of square-integrable random vectors modulo almost-everywhere equality. Its almost-everywhere representative is exactly the original vector function evaluated at the actual process state.

\begin{leanbox}{Lean: \texttt{stationary\_\allowbreak{}observation\_\allowbreak{}memLp}, \texttt{stationaryObservation}, \texttt{stationaryObservation\_\allowbreak{}coe}}
\leansource{\texttt{stationary\_\allowbreak{}observation\_\allowbreak{}memLp} \quad AverageLMC/\allowbreak{}LaplacianStationaryTest.lean, lines 22–26}
\begin{leancode}
\leanline{0}{theorem\ stationary\_observation\_memLp}
\leanline{4}{(P\ :\ Measure\ Ω)\ (π\ :\ Measure\ H)\ (F\ :\ H\ →\ E)\ (hF\ :\ MemLp\ F\ 2\ π)}
\leanline{4}{(Y\ :\ Lp\ H\ 2\ P)\ (hlaw\ :\ P.map\ Y\ =\ π)\ :\ MemLp\ (fun\ a\ ↦\ F\ (Y\ a))\ 2\ P}
\end{leancode}
\leansource{\texttt{stationaryObservation} \quad AverageLMC/\allowbreak{}LaplacianStationaryTest.lean, lines 29–32}
\begin{leancode}
\leanline{0}{def\ stationaryObservation}
\leanline{4}{(P\ :\ Measure\ Ω)\ (π\ :\ Measure\ H)\ (F\ :\ H\ →\ E)\ (hF\ :\ MemLp\ F\ 2\ π)}
\leanline{4}{(Y\ :\ Lp\ H\ 2\ P)\ (hlaw\ :\ P.map\ Y\ =\ π)\ :\ Lp\ E\ 2\ P\ :=}
\leanline{2}{(stationary\_observation\_memLp\ P\ π\ F\ hF\ Y\ hlaw).toLp\ (fun\ a\ ↦\ F\ (Y\ a))}
\end{leancode}
\leansource{\texttt{stationaryObservation\_\allowbreak{}coe} \quad AverageLMC/\allowbreak{}LaplacianStationaryTest.lean, lines 36–40}
\begin{leancode}
\leanline{0}{theorem\ stationaryObservation\_coe}
\leanline{4}{(P\ :\ Measure\ Ω)\ (π\ :\ Measure\ H)\ (F\ :\ H\ →\ E)\ (hF\ :\ MemLp\ F\ 2\ π)}
\leanline{4}{(Y\ :\ Lp\ H\ 2\ P)\ (hlaw\ :\ P.map\ Y\ =\ π)\ :}
\leanline{4}{stationaryObservation\ P\ π\ F\ hF\ Y\ hlaw\ =ᵐ[P]\ fun\ a\ ↦\ F\ (Y\ a)}
\end{leancode}
\leanlabel{How the Lean proof follows the argument}
\texttt{stationary\_\allowbreak{}observation\_\allowbreak{}mem\allowbreak{}Lp} gives the formal counterpart of this step. Named proof applications include \texttt{h\allowbreak{}F.\allowbreak{}comp\_\allowbreak{}of\_\allowbreak{}map}; the adjacent proof shows their hypotheses and order of use.




\texttt{stationary\allowbreak{}Observation} names the mathematical object defined above. Its defining expression appears in the adjacent statement; subsequent identities may unfold this definition directly.




\texttt{stationary\allowbreak{}Observation\_\allowbreak{}coe} gives the formal counterpart of this step.
\end{leanbox}
\subsection{\texorpdfstring{Preserve the exact squared approximation error}{Preserve the exact squared approximation error}}
\label{report-step-AverageLMC-LaplacianStationaryTest-lean-2}


The \(L^{2}\) inner-product formula identifies the squared norm with the integral of the squared representative norm. The proved representative identity and the stationary law map turn this into the actual Gibbs moment. Applying the same argument to differences preserves the complete squared approximation error, with no change of constants or dependence on time.

\begin{leanbox}{Lean: \texttt{stationaryObservation\_\allowbreak{}norm\_\allowbreak{}sq}, \texttt{stationaryObservation\_\allowbreak{}sub\_\allowbreak{}norm\_\allowbreak{}sq}}
\leansource{\texttt{stationaryObservation\_\allowbreak{}norm\_\allowbreak{}sq} \quad AverageLMC/\allowbreak{}LaplacianStationaryTest.lean, lines 45–58}
\begin{leancode}
\leanline{0}{theorem\ stationaryObservation\_norm\_sq}
\leanline{4}{(P\ :\ Measure\ Ω)\ (π\ :\ Measure\ H)\ (F\ :\ H\ →\ E)\ (hFc\ :\ Continuous\ F)\ (hF\ :\ MemLp\ F\ 2\ π)}
\leanline{4}{(Y\ :\ Lp\ H\ 2\ P)\ (hlaw\ :\ P.map\ Y\ =\ π)\ :}
\leanline{4}{‖stationaryObservation\ P\ π\ F\ hF\ Y\ hlaw‖\ \textasciicircum{}\ 2\ =\ ∫\ x,\ ‖F\ x‖\ \textasciicircum{}\ 2\ ∂π}
\end{leancode}
\leansource{\texttt{stationaryObservation\_\allowbreak{}sub\_\allowbreak{}norm\_\allowbreak{}sq} \quad AverageLMC/\allowbreak{}LaplacianStationaryTest.lean, lines 64–85}
\begin{leancode}
\leanline{0}{theorem\ stationaryObservation\_sub\_norm\_sq}
\leanline{4}{(P\ :\ Measure\ Ω)\ (π\ :\ Measure\ H)\ (F\ G\ :\ H\ →\ E)}
\leanline{4}{(hFc\ :\ Continuous\ F)\ (hGc\ :\ Continuous\ G)\ (hF\ :\ MemLp\ F\ 2\ π)\ (hG\ :\ MemLp\ G\ 2\ π)}
\leanline{4}{(Y\ :\ Lp\ H\ 2\ P)\ (hlaw\ :\ P.map\ Y\ =\ π)\ :}
\leanline{4}{‖stationaryObservation\ P\ π\ F\ hF\ Y\ hlaw\ -\ stationaryObservation\ P\ π\ G\ hG\ Y\ hlaw‖\ \textasciicircum{}\ 2\ =}
\leanline{6}{∫\ x,\ ‖F\ x\ -\ G\ x‖\ \textasciicircum{}\ 2\ ∂π}
\end{leancode}
\leanlabel{How the Lean proof follows the argument}
\texttt{stationary\allowbreak{}Observation\_\allowbreak{}norm\_\allowbreak{}sq} gives the formal counterpart of this step. Named proof applications include \texttt{integral\_\allowbreak{}congr\_\allowbreak{}ae}; the adjacent proof shows their hypotheses and order of use.




\texttt{stationary\allowbreak{}Observation\_\allowbreak{}sub\_\allowbreak{}norm\_\allowbreak{}sq} gives the formal counterpart of this step. Named proof applications include \texttt{integral\_\allowbreak{}congr\_\allowbreak{}ae}; the adjacent proof shows their hypotheses and order of use.
\end{leanbox}
\subsection{\texorpdfstring{Sum scalar weighted identities into vector duality}{Sum scalar weighted identities into vector duality}}
\label{report-step-AverageLMC-LaplacianStationaryTest-lean-3}


Each coordinate product is integrable by domination with the product of the two vector norms and the \(L^{2}\) Cauchy inequality. Expand the Euclidean inner product as the finite coordinate sum, then interchange this finite sum and the genuine integral. The result is exactly the sum of scalar weighted expectations used by the compact-test Dynkin theorem; no separate coordinatewise bound introduces a dimension factor.

\begin{leanbox}{Lean: \texttt{stationaryObservation\_\allowbreak{}inner\_\allowbreak{}eq\_\allowbreak{}sum}}
\leansource{\texttt{stationaryObservation\_\allowbreak{}inner\_\allowbreak{}eq\_\allowbreak{}sum} \quad AverageLMC/\allowbreak{}LaplacianStationaryTest.lean, lines 92–120}
\begin{leancode}
\leanline{0}{theorem\ stationaryObservation\_inner\_eq\_sum}
\leanline{4}{(π\ :\ Measure\ (State\ p))\ (F\ :\ State\ p\ →\ State\ p)\ (hFc\ :\ Continuous\ F)}
\leanline{4}{(hF\ :\ MemLp\ F\ 2\ π)\ (Y\ :\ Lp\ (State\ p)\ 2\ P)\ (hlaw\ :\ P.map\ Y\ =\ π)}
\leanline{4}{(Ψ\ :\ Lp\ (State\ p)\ 2\ P)\ :}
\leanline{4}{inner\ ℝ\ Ψ\ (stationaryObservation\ P\ π\ F\ hF\ Y\ hlaw)\ =}
\leanline{6}{∑\ j\ :\ Fin\ p,\ lpWeightedTest\ P\ (fun\ a\ ↦\ Ψ\ a\ j)\ (fun\ x\ ↦\ F\ x\ j)\ Y}
\end{leancode}
\leanlabel{How the Lean proof follows the argument}
\texttt{stationary\allowbreak{}Observation\_\allowbreak{}inner\_\allowbreak{}eq\_\allowbreak{}sum} gives the formal counterpart of this step. The main intermediate facts are \texttt{h\allowbreak{}Obs}, \texttt{hdom}, \texttt{hi} in that order. Named proof applications include \texttt{hdom.\allowbreak{}mono'}, \texttt{mul\_\allowbreak{}le\_\allowbreak{}mul}, \texttt{integral\_\allowbreak{}congr\_\allowbreak{}ae}; the adjacent proof shows their hypotheses and order of use.
\end{leanbox}
\section{\texorpdfstring{Product and coordinate identities for the true Langevin generator}{Product and coordinate identities for the true Langevin generator}}
\label{proof-module-AverageLMC-GeneratorProduct-lean}


The cutoff approximation for coordinate scores uses the exact generator product rule. These identities are deterministic calculus statements about the actual gradient, derivative and generator; no probabilistic cancellation is postulated.



For \(\mathcal L_f\psi=\Delta\psi-\langle\nabla f,\nabla\psi\rangle\), the compact cutoff construction uses the genuine product rule and the coordinate gradient identity:


\[\mathcal L_f(\chi\psi)=\chi\mathcal L_f\psi+\psi\mathcal L_f\chi+2\langle\nabla\chi,\nabla\psi\rangle.\]
\[\mathcal L_f(\partial_jf)=(G_f)_j,\qquad G_f=\nabla\Delta f-D(\nabla f)[\nabla f].\]

Earlier results used here: \cite[Section~3.46]{LangevinLean2026ULASC}.
\subsection{\texorpdfstring{Differentiate the product once and twice}{Differentiate the product once and twice}}
\label{report-step-AverageLMC-GeneratorProduct-lean-1}


The Fréchet product rule, evaluated on a fixed direction, gives \(D_{v}(\chi \psi )=\chi D_{v}\psi +\psi D_{v}\chi\). The corresponding addition rule allows this expression to be differentiated once more. For \(C^{2}\) factors their first directional derivatives are differentiable, so another product rule gives \(D_{v}^{2}(\chi \psi )=\chi D_{v}^{2}\psi +\psi D_{v}^{2}\chi +2D_{v}\chi D_{v}\psi\). The cross coefficient is exactly two.

\begin{leanbox}{Lean: \texttt{directionalDerivative\_\allowbreak{}mul}, \texttt{directionalDerivative\_\allowbreak{}add}, \texttt{directionalDerivative\_\allowbreak{}two\_\allowbreak{}mul}}
\leansource{\texttt{directionalDerivative\_\allowbreak{}mul} \quad AverageLMC/\allowbreak{}GeneratorProduct.lean, lines 12–18}
\begin{leancode}
\leanline{0}{theorem\ directionalDerivative\_mul\ (χ\ ψ\ :\ E\ →\ ℝ)\ (hχ\ :\ Differentiable\ ℝ\ χ)}
\leanline{4}{(hψ\ :\ Differentiable\ ℝ\ ψ)\ (v\ x\ :\ E)\ :}
\leanline{4}{directionalDerivative\ (fun\ y\ ↦\ χ\ y\ *\ ψ\ y)\ v\ x\ =}
\leanline{6}{χ\ x\ *\ directionalDerivative\ ψ\ v\ x\ +\ ψ\ x\ *\ directionalDerivative\ χ\ v\ x}
\end{leancode}
\leansource{\texttt{directionalDerivative\_\allowbreak{}add} \quad AverageLMC/\allowbreak{}GeneratorProduct.lean, lines 20–26}
\begin{leancode}
\leanline{0}{theorem\ directionalDerivative\_add\ (χ\ ψ\ :\ E\ →\ ℝ)\ (hχ\ :\ Differentiable\ ℝ\ χ)}
\leanline{4}{(hψ\ :\ Differentiable\ ℝ\ ψ)\ (v\ x\ :\ E)\ :}
\leanline{4}{directionalDerivative\ (fun\ y\ ↦\ χ\ y\ +\ ψ\ y)\ v\ x\ =}
\leanline{6}{directionalDerivative\ χ\ v\ x\ +\ directionalDerivative\ ψ\ v\ x}
\end{leancode}
\leansource{\texttt{directionalDerivative\_\allowbreak{}two\_\allowbreak{}mul} \quad AverageLMC/\allowbreak{}GeneratorProduct.lean, lines 29–45}
\begin{leancode}
\leanline{0}{theorem\ directionalDerivative\_two\_mul\ (χ\ ψ\ :\ E\ →\ ℝ)\ (hχ\ :\ ContDiff\ ℝ\ 2\ χ)}
\leanline{4}{(hψ\ :\ ContDiff\ ℝ\ 2\ ψ)\ (v\ x\ :\ E)\ :}
\leanline{4}{directionalDerivative\ (directionalDerivative\ (fun\ y\ ↦\ χ\ y\ *\ ψ\ y)\ v)\ v\ x\ =}
\leanline{6}{χ\ x\ *\ directionalDerivative\ (directionalDerivative\ ψ\ v)\ v\ x\ +}
\leanline{6}{ψ\ x\ *\ directionalDerivative\ (directionalDerivative\ χ\ v)\ v\ x\ +}
\leanline{6}{2\ *\ directionalDerivative\ χ\ v\ x\ *\ directionalDerivative\ ψ\ v\ x}
\end{leancode}
\leanlabel{How the Lean proof follows the argument}
\texttt{directional\allowbreak{}Derivative\_\allowbreak{}mul} gives the formal counterpart of this step. Its remaining equalities are discharged by the simplifications displayed in the proof.




\texttt{directional\allowbreak{}Derivative\_\allowbreak{}add} gives the formal counterpart of this step.




\texttt{directional\allowbreak{}Derivative\_\allowbreak{}two\_\allowbreak{}mul} gives the formal counterpart of this step. The main intermediate facts are \texttt{h\ensuremath{\chi}d}, \texttt{h\ensuremath{\psi}d}, \texttt{h\allowbreak{}D\ensuremath{\chi}}, \texttt{h\allowbreak{}D\ensuremath{\psi}}, \texttt{heq} in that order. The final step is the algebraic normalization or ordered arithmetic displayed in the proof.
\end{leanbox}
\subsection{\texorpdfstring{Sum the directional generator product rules}{Sum the directional generator product rules}}
\label{report-step-AverageLMC-GeneratorProduct-lean-2}


Subtract the first-order gradient drift from the second directional derivative product identity. Algebra separates the two factors’ directional generators and retains the cross term \(2D_{v}\chi D_{v}\psi\). Summing over the standard orthonormal basis produces the full Langevin generator product rule. The sum of the products of the directional derivatives equals the inner product of the two actual gradients by the orthonormal-basis inner-product identity.

\begin{leanbox}{Lean: \texttt{directionalGenerator\_\allowbreak{}mul}, \texttt{langevinGenerator\_\allowbreak{}mul}}
\leansource{\texttt{directionalGenerator\_\allowbreak{}mul} \quad AverageLMC/\allowbreak{}GeneratorProduct.lean, lines 49–69}
\begin{leancode}
\leanline{0}{theorem\ directionalGenerator\_mul\ (f\ χ\ ψ\ :\ E\ →\ ℝ)\ (hχ\ :\ ContDiff\ ℝ\ 2\ χ)}
\leanline{4}{(hψ\ :\ ContDiff\ ℝ\ 2\ ψ)\ (v\ x\ :\ E)\ :}
\leanline{4}{ULA.directionalGenerator\ f\ (fun\ y\ ↦\ χ\ y\ *\ ψ\ y)\ v\ x\ =}
\leanline{6}{χ\ x\ *\ ULA.directionalGenerator\ f\ ψ\ v\ x\ +}
\leanline{6}{ψ\ x\ *\ ULA.directionalGenerator\ f\ χ\ v\ x\ +}
\leanline{6}{2\ *\ directionalDerivative\ χ\ v\ x\ *\ directionalDerivative\ ψ\ v\ x}
\end{leancode}
\leansource{\texttt{langevinGenerator\_\allowbreak{}mul} \quad AverageLMC/\allowbreak{}GeneratorProduct.lean, lines 72–85}
\begin{leancode}
\leanline{0}{theorem\ langevinGenerator\_mul\ (f\ χ\ ψ\ :\ E\ →\ ℝ)\ (hχ\ :\ ContDiff\ ℝ\ 2\ χ)}
\leanline{4}{(hψ\ :\ ContDiff\ ℝ\ 2\ ψ)\ (x\ :\ E)\ :}
\leanline{4}{ULA.langevinGenerator\ f\ (fun\ y\ ↦\ χ\ y\ *\ ψ\ y)\ x\ =}
\leanline{6}{χ\ x\ *\ ULA.langevinGenerator\ f\ ψ\ x\ +\ ψ\ x\ *\ ULA.langevinGenerator\ f\ χ\ x\ +}
\leanline{6}{2\ *\ inner\ ℝ\ (gradient\ χ\ x)\ (gradient\ ψ\ x)}
\end{leancode}
\leanlabel{How the Lean proof follows the argument}
\texttt{directional\allowbreak{}Generator\_\allowbreak{}mul} gives the formal counterpart of this step. The final step is the algebraic normalization or ordered arithmetic displayed in the proof.




\texttt{langevin\allowbreak{}Generator\_\allowbreak{}mul} gives the formal counterpart of this step. Its remaining equalities are discharged by the simplifications displayed in the proof.
\end{leanbox}


Earlier results used here: \cite[Section~3.46.1]{LangevinLean2026ULASC}; \cite[Section~3.46.2]{LangevinLean2026ULASC}.
\subsection{\texorpdfstring{Identify the coordinate score generator}{Identify the coordinate score generator}}
\label{report-step-AverageLMC-GeneratorProduct-lean-3}


A coordinate score is the continuous linear projection of the \(C^{2}\) vector field \(\nabla f\). Projection commutes with its vector Laplacian, and the proved \(C^{3}\) commutation identity turns this into the corresponding coordinate of \(\nabla \Delta f\). The Fréchet chain rule identifies its derivative along \(\nabla f\) with the same coordinate of \(D(\nabla f)\nabla f\). Subtracting the drift term therefore gives exactly that coordinate of the actual vector gradient generator.

\begin{leanbox}{Lean: \texttt{langevinGenerator\_\allowbreak{}coordinate\_\allowbreak{}gradient}}
\leansource{\texttt{langevinGenerator\_\allowbreak{}coordinate\_\allowbreak{}gradient} \quad AverageLMC/\allowbreak{}GeneratorProduct.lean, lines 89–108}
\begin{leancode}
\leanline{0}{theorem\ langevinGenerator\_coordinate\_gradient\ \{p\ :\ ℕ\}}
\leanline{4}{(f\ :\ State\ p\ →\ ℝ)\ (hf\ :\ ContDiff\ ℝ\ 3\ f)\ (j\ :\ Fin\ p)\ (x\ :\ State\ p)\ :}
\leanline{4}{ULA.langevinGenerator\ f\ (fun\ y\ ↦\ (gradient\ f\ y)\ j)\ x\ =\ (gradientGenerator\ f\ x)\ j}
\end{leancode}
\leanlabel{How the Lean proof follows the argument}
\texttt{langevin\allowbreak{}Generator\_\allowbreak{}coordinate\_\allowbreak{}gradient} gives the formal counterpart of this step. The main intermediate facts are \texttt{hg}, \texttt{h\ensuremath{\psi}}, \texttt{h\allowbreak{}Lap}, \texttt{hd} in that order. Its remaining equalities are discharged by the simplifications displayed in the proof.
\end{leanbox}


Earlier results used here: \cite[Section~3.46.14]{LangevinLean2026ULASC}; \cite[Section~3.46.2]{LangevinLean2026ULASC}.
\section{\texorpdfstring{Compact \(C^{2}\) approximation of the actual Gibbs score and generator}{Compact C² approximation of the actual Gibbs score and generator}}
\label{proof-module-AverageLMC-GibbsScoreCutoff-lean}


Smooth radial cutoffs turn an unbounded coordinate of the potential gradient into compact \(C^{2}\) tests. Both the cut-off scores and their true Langevin generators converge in Gibbs \(L^{2}\). The proof derives all score moments from their sub-Gaussian MGFs, controls the product-generator cross terms using exactly the Laplacian-gradient assumption, and applies dominated convergence. This justifies extending compact-test Dynkin identities to the actual score under the paper \(C^{3}\) potential assumption.



Write the genuine scalar generator and vector score generator as

\[g=\nabla f,\qquad \mathcal L\phi=\Delta\phi-\langle g,\nabla\phi\rangle,\qquad G=\nabla\Delta f-Dg\,g.\]

There are compact \(C^{2}\) coordinate tests, with \(\Phi_n=(\phi_{n,j})_{j=1}^p\) and \(\mathcal L\Phi_n=(\mathcal L\phi_{n,j})_{j=1}^p\), for which every displayed vector belongs to Gibbs \(L^{2}\) and

\[\phi_{n,j}\in C_c^2(\mathbb R^p),\qquad \int\!\|\Phi_n-g\|^2\,d\pi\longrightarrow0,\qquad \int\!\|\mathcal L\Phi_n-G\|^2\,d\pi\longrightarrow0.\]\subsection{\texorpdfstring{Build an expanding smooth radial cutoff family}{Build an expanding smooth radial cutoff family}}
\label{report-step-AverageLMC-GibbsScoreCutoff-lean-1}


The \(n\)th bump equals one on the ball of radius \(n+1\) and has support inside radius \(2(n+1)\). It is smooth, compactly supported, and takes values in the interval from zero to one. It is exactly one fixed unit-scale bump composed with multiplication by \(1/(n+1)\). At every fixed spatial point, all sufficiently large cutoffs are identically one in a neighborhood of that point.

\begin{leanbox}{Lean: \texttt{radialCutoffBump}, \texttt{radialCutoff}, \texttt{radialCutoff\_\allowbreak{}contDiff} and 4 more}
\leansource{\texttt{radialCutoffBump} \quad AverageLMC/\allowbreak{}GibbsScoreCutoff.lean, lines 18–22}
\begin{leancode}
\leanline{0}{def\ radialCutoffBump\ \{p\ :\ ℕ\}\ (n\ :\ ℕ)\ :\ ContDiffBump\ (0\ :\ State\ p)\ where}
\leanline{2}{rIn\ :=\ (n\ :\ ℝ)\ +\ 1}
\leanline{2}{rOut\ :=\ 2\ *\ ((n\ :\ ℝ)\ +\ 1)}
\leanline{2}{rIn\_pos\ :=\ by\ positivity}
\leanline{2}{rIn\_lt\_rOut\ :=\ by\ linarith\ [Nat.cast\_nonneg\ (α\ :=\ ℝ)\ n]}
\end{leancode}
\leansource{\texttt{radialCutoff} \quad AverageLMC/\allowbreak{}GibbsScoreCutoff.lean, lines 26–26}
\begin{leancode}
\leanline{0}{def\ radialCutoff\ \{p\ :\ ℕ\}\ (n\ :\ ℕ)\ :\ State\ p\ →\ ℝ\ :=\ radialCutoffBump\ n}
\end{leancode}
\leansource{\texttt{radialCutoff\_\allowbreak{}contDiff} \quad AverageLMC/\allowbreak{}GibbsScoreCutoff.lean, lines 28–29}
\begin{leancode}
\leanline{0}{theorem\ radialCutoff\_contDiff\ \{p\ :\ ℕ\}\ (n\ :\ ℕ)\ :\ ContDiff\ ℝ\ ∞\ (@radialCutoff\ p\ n)}
\end{leancode}
\leansource{\texttt{radialCutoff\_\allowbreak{}compactSupport} \quad AverageLMC/\allowbreak{}GibbsScoreCutoff.lean, lines 31–32}
\begin{leancode}
\leanline{0}{theorem\ radialCutoff\_compactSupport\ \{p\ :\ ℕ\}\ (n\ :\ ℕ)\ :}
\leanline{4}{HasCompactSupport\ (@radialCutoff\ p\ n)}
\end{leancode}
\leansource{\texttt{radialCutoff\_\allowbreak{}mem\_\allowbreak{}Icc} \quad AverageLMC/\allowbreak{}GibbsScoreCutoff.lean, lines 34–36}
\begin{leancode}
\leanline{0}{theorem\ radialCutoff\_mem\_Icc\ \{p\ :\ ℕ\}\ (n\ :\ ℕ)\ (x\ :\ State\ p)\ :}
\leanline{4}{radialCutoff\ n\ x\ ∈\ Icc\ (0\ :\ ℝ)\ 1}
\end{leancode}
\leansource{\texttt{radialCutoff\_\allowbreak{}eq\_\allowbreak{}scaled} \quad AverageLMC/\allowbreak{}GibbsScoreCutoff.lean, lines 38–41}
\begin{leancode}
\leanline{0}{theorem\ radialCutoff\_eq\_scaled\ \{p\ :\ ℕ\}\ (n\ :\ ℕ)\ (x\ :\ State\ p)\ :}
\leanline{4}{radialCutoff\ n\ x\ =\ radialCutoff\ 0\ ((((n\ :\ ℝ)\ +\ 1)⁻¹)\ •\ x)}
\end{leancode}
\leansource{\texttt{radialCutoff\_\allowbreak{}eventuallyEq\_\allowbreak{}one} \quad AverageLMC/\allowbreak{}GibbsScoreCutoff.lean, lines 45–53}
\begin{leancode}
\leanline{0}{theorem\ radialCutoff\_eventuallyEq\_one\ \{p\ :\ ℕ\}\ (x\ :\ State\ p)\ :}
\leanline{4}{∀ᶠ\ n\ :\ ℕ\ in\ atTop,\ radialCutoff\ n\ =ᶠ[{\leanfallback 𝓝}\ x]\ 1}
\end{leancode}
\leanlabel{How the Lean proof follows the argument}
\texttt{radial\allowbreak{}Cutoff\allowbreak{}Bump} names the mathematical object defined above. Its defining expression appears in the adjacent statement; subsequent identities may unfold this definition directly.




\texttt{radial\allowbreak{}Cutoff} names the mathematical object defined above. Its defining expression appears in the adjacent statement; subsequent identities may unfold this definition directly.




\texttt{radial\allowbreak{}Cutoff\_\allowbreak{}cont\allowbreak{}Diff} gives the formal counterpart of this step.




\texttt{radial\allowbreak{}Cutoff\_\allowbreak{}compact\allowbreak{}Support} gives the formal counterpart of this step.




\texttt{radial\allowbreak{}Cutoff\_\allowbreak{}mem\_\allowbreak{}Icc} gives the formal counterpart of this step.




\texttt{radial\allowbreak{}Cutoff\_\allowbreak{}eq\_\allowbreak{}scaled} gives the formal counterpart of this step. The main intermediate facts are \texttt{hr} in that order. Its remaining equalities are discharged by the simplifications displayed in the proof.




\texttt{radial\allowbreak{}Cutoff\_\allowbreak{}eventually\allowbreak{}Eq\_\allowbreak{}one} gives the formal counterpart of this step. The main intermediate facts are \texttt{hcast} in that order. The final step is the algebraic normalization or ordered arithmetic displayed in the proof.
\end{leanbox}
\subsection{\texorpdfstring{Bound the first and second cutoff derivatives uniformly}{Bound the first and second cutoff derivatives uniformly}}
\label{report-step-AverageLMC-GibbsScoreCutoff-lean-2}


The fixed compact smooth bump has bounded derivatives through order two. Rescaling by \(1/(n+1)\) multiplies its \(k\)th derivative by that scalar to the \(k\)th power, which is at most one. Hence a common constant bounds the first two derivative norms for every cutoff. Gradient duality converts the first derivative bound into a gradient bound. Expanding the Laplacian in the canonical orthonormal basis bounds its absolute value by \(p\) times the same constant.

\begin{leanbox}{Lean: \texttt{radialCutoff\_\allowbreak{}derivative\_\allowbreak{}bounds}, \texttt{radialCutoff\_\allowbreak{}gradient\_\allowbreak{}laplacian\_\allowbreak{}bounds}}
\leansource{\texttt{radialCutoff\_\allowbreak{}derivative\_\allowbreak{}bounds} \quad AverageLMC/\allowbreak{}GibbsScoreCutoff.lean, lines 57–75}
\begin{leancode}
\leanline{0}{theorem\ radialCutoff\_derivative\_bounds\ \{p\ :\ ℕ\}\ :}
\leanline{4}{∃\ C\ :\ ℝ,\ 0\ ≤\ C\ ∧\ ∀\ n\ k,\ k\ ≤\ 2\ →\ ∀\ x\ :\ State\ p,}
\leanline{6}{‖iteratedFDeriv\ ℝ\ k\ (radialCutoff\ n)\ x‖\ ≤\ C}
\end{leancode}
\leansource{\texttt{radialCutoff\_\allowbreak{}gradient\_\allowbreak{}laplacian\_\allowbreak{}bounds} \quad AverageLMC/\allowbreak{}GibbsScoreCutoff.lean, lines 77–105}
\begin{leancode}
\leanline{0}{theorem\ radialCutoff\_gradient\_laplacian\_bounds\ \{p\ :\ ℕ\}\ :}
\leanline{4}{∃\ C\ :\ ℝ,\ 0\ ≤\ C\ ∧\ ∀\ n\ (x\ :\ State\ p),}
\leanline{6}{‖gradient\ (radialCutoff\ n)\ x‖\ ≤\ C\ ∧}
\leanline{8}{|laplacian\ (radialCutoff\ n)\ x|\ ≤\ (p\ :\ ℝ)\ *\ C}
\end{leancode}
\leanlabel{How the Lean proof follows the argument}
\texttt{radial\allowbreak{}Cutoff\_\allowbreak{}derivative\_\allowbreak{}bounds} gives the formal counterpart of this step. The main intermediate facts are \texttt{hfun}, \texttt{ha}, \texttt{hp} in that order. Its remaining equalities are discharged by the simplifications displayed in the proof.




\texttt{radial\allowbreak{}Cutoff\_\allowbreak{}gradient\_\allowbreak{}laplacian\_\allowbreak{}bounds} gives the formal counterpart of this step. Named proof applications include \texttt{Finset.\allowbreak{}abs\_\allowbreak{}sum\_\allowbreak{}le\_\allowbreak{}sum\_\allowbreak{}abs}, \texttt{Finset.\allowbreak{}sum\_\allowbreak{}le\_\allowbreak{}sum}, \texttt{h.\allowbreak{}trans}; the adjacent proof shows their hypotheses and order of use.
\end{leanbox}
\subsection{\texorpdfstring{Derive every finite Gibbs score moment}{Derive every finite Gibbs score moment}}
\label{report-step-AverageLMC-GibbsScoreCutoff-lean-3}


Each true coordinate of the Gibbs score has the existing Gaussian MGF bound obtained by translation of the density. Sub-Gaussianity supplies all finite scalar moments. Finite-dimensional coordinate membership then proves \(L^{q}\) membership of the entire actual gradient for every finite exponent \(q\). In particular its fourth moment is finite, so one plus its squared norm belongs to Gibbs \(L^{2}\).

\begin{leanbox}{Lean: \texttt{memLp\_\allowbreak{}gradient\_\allowbreak{}gibbs\_\allowbreak{}all}}
\leansource{\texttt{memLp\_\allowbreak{}gradient\_\allowbreak{}gibbs\_\allowbreak{}all} \quad AverageLMC/\allowbreak{}GibbsScoreCutoff.lean, lines 109–116}
\begin{leancode}
\leanline{0}{theorem\ memLp\_gradient\_gibbs\_all\ \{p\ :\ ℕ\}}
\leanline{4}{(f\ :\ State\ p\ →\ ℝ)\ (M\ :\ ℝ≥0)\ (hf\ :\ Differentiable\ ℝ\ f)}
\leanline{4}{(hLip\ :\ LipschitzWith\ M\ (gradient\ f))\ (hZ\ :\ ULA.gibbsPartition\ f\ ≠\ ⊤)}
\leanline{4}{(q\ :\ ℝ≥0)\ :\ MemLp\ (gradient\ f)\ q\ (ULA.gibbsMeasure\ f)}
\end{leancode}
\leanlabel{How the Lean proof follows the argument}
\texttt{mem\allowbreak{}Lp\_\allowbreak{}gradient\_\allowbreak{}gibbs\_\allowbreak{}all} gives the formal counterpart of this step. The main intermediate facts are \texttt{hs} in that order. Named proof applications include \texttt{mem\allowbreak{}Lp\_\allowbreak{}pi\allowbreak{}Lp\_\allowbreak{}iff.\allowbreak{}mpr}, \texttt{hs.\allowbreak{}mem\allowbreak{}Lp}; the adjacent proof shows their hypotheses and order of use.
\end{leanbox}


Earlier results used here: \cite[Section~3.26.2]{LangevinLean2026ULASC}; \cite[Section~3.26.3]{LangevinLean2026ULASC}; \cite[Section~3.27.3]{LangevinLean2026ULASC}.
\subsection{\texorpdfstring{Preserve compact \(C^{2}\) regularity and local generator equality}{Preserve compact C² regularity and local generator equality}}
\label{report-step-AverageLMC-GibbsScoreCutoff-lean-4}


Multiply a \(C^{2}\) test by the smooth radial cutoff. The product is \(C^{2}\) and has compact support. Because the cutoff is eventually locally one at a fixed point, the products eventually agree with the original test on a neighborhood. Their Frechet derivatives and Laplacians therefore agree there as well. The genuine Langevin generators consequently become pointwise equal. This uses locality of the actual derivatives rather than an assumed approximation identity.

\begin{leanbox}{Lean: \texttt{radialCutoff\_\allowbreak{}mul\_\allowbreak{}contDiff}, \texttt{radialCutoff\_\allowbreak{}mul\_\allowbreak{}compactSupport}, \texttt{radialCutoff\_\allowbreak{}generator\_\allowbreak{}eventually\_\allowbreak{}eq}}
\leansource{\texttt{radialCutoff\_\allowbreak{}mul\_\allowbreak{}contDiff} \quad AverageLMC/\allowbreak{}GibbsScoreCutoff.lean, lines 119–122}
\begin{leancode}
\leanline{0}{theorem\ radialCutoff\_mul\_contDiff\ \{p\ :\ ℕ\}\ (ψ\ :\ State\ p\ →\ ℝ)}
\leanline{4}{(hψ\ :\ ContDiff\ ℝ\ 2\ ψ)\ (n\ :\ ℕ)\ :}
\leanline{4}{ContDiff\ ℝ\ 2\ (fun\ x\ ↦\ radialCutoff\ n\ x\ *\ ψ\ x)}
\end{leancode}
\leansource{\texttt{radialCutoff\_\allowbreak{}mul\_\allowbreak{}compactSupport} \quad AverageLMC/\allowbreak{}GibbsScoreCutoff.lean, lines 124–126}
\begin{leancode}
\leanline{0}{theorem\ radialCutoff\_mul\_compactSupport\ \{p\ :\ ℕ\}\ (ψ\ :\ State\ p\ →\ ℝ)\ (n\ :\ ℕ)\ :}
\leanline{4}{HasCompactSupport\ (fun\ x\ ↦\ radialCutoff\ n\ x\ *\ ψ\ x)}
\end{leancode}
\leansource{\texttt{radialCutoff\_\allowbreak{}generator\_\allowbreak{}eventually\_\allowbreak{}eq} \quad AverageLMC/\allowbreak{}GibbsScoreCutoff.lean, lines 130–148}
\begin{leancode}
\leanline{0}{theorem\ radialCutoff\_generator\_eventually\_eq\ \{p\ :\ ℕ\}\ (f\ ψ\ :\ State\ p\ →\ ℝ)}
\leanline{4}{(hψ\ :\ ContDiff\ ℝ\ 2\ ψ)\ (x\ :\ State\ p)\ :}
\leanline{4}{∀ᶠ\ n\ :\ ℕ\ in\ atTop,}
\leanline{6}{(radialCutoff\ n\ x\ *\ ψ\ x\ =\ ψ\ x)\ ∧}
\leanline{6}{ULA.langevinGenerator\ f\ (fun\ y\ ↦\ radialCutoff\ n\ y\ *\ ψ\ y)\ x\ =}
\leanline{8}{ULA.langevinGenerator\ f\ ψ\ x}
\end{leancode}
\leanlabel{How the Lean proof follows the argument}
\texttt{radial\allowbreak{}Cutoff\_\allowbreak{}mul\_\allowbreak{}cont\allowbreak{}Diff} gives the formal counterpart of this step. Its remaining equalities are discharged by the simplifications displayed in the proof.




\texttt{radial\allowbreak{}Cutoff\_\allowbreak{}mul\_\allowbreak{}compact\allowbreak{}Support} gives the formal counterpart of this step.




\texttt{radial\allowbreak{}Cutoff\_\allowbreak{}generator\_\allowbreak{}eventually\_\allowbreak{}eq} gives the formal counterpart of this step. The main intermediate facts are \texttt{heq}, \texttt{h\allowbreak{}Lap}, \texttt{hgrad} in that order. Its remaining equalities are discharged by the simplifications displayed in the proof.
\end{leanbox}


Earlier results used here: \cite[Section~3.46.14]{LangevinLean2026ULASC}; \cite[Section~3.46.2]{LangevinLean2026ULASC}.
\subsection{\texorpdfstring{Prove dominated \(L^{2}\) convergence for the eventual equalities}{Prove dominated L² convergence for the eventual equalities}}
\label{report-step-AverageLMC-GibbsScoreCutoff-lean-5}


If every approximation and its limit are bounded in absolute value by a common \(L^{2}\) function \(B\), their squared differences are bounded by 4 \(B^2\). Eventual pointwise equality makes those differences tend to zero. The ordinary dominated convergence theorem then proves convergence of their actual squared integrals to zero. The module also proves continuity of the actual scalar Langevin generator for \(C^{2}\) tests and the bound \(M\) on the gradient norm of every score coordinate.

\begin{leanbox}{Lean: \texttt{integral\_\allowbreak{}sq\_\allowbreak{}sub\_\allowbreak{}tendsto\_\allowbreak{}zero\_\allowbreak{}of\_\allowbreak{}eventually\_\allowbreak{}eq}, \texttt{continuous\_\allowbreak{}langevinGenerator\_\allowbreak{}test}, \texttt{gradient\_\allowbreak{}score\_\allowbreak{}coordinate\_\allowbreak{}norm\_\allowbreak{}le}}
\leansource{\texttt{integral\_\allowbreak{}sq\_\allowbreak{}sub\_\allowbreak{}tendsto\_\allowbreak{}zero\_\allowbreak{}of\_\allowbreak{}eventually\_\allowbreak{}eq} \quad AverageLMC/\allowbreak{}GibbsScoreCutoff.lean, lines 153–178}
\begin{leancode}
\leanline{0}{theorem\ integral\_sq\_sub\_tendsto\_zero\_of\_eventually\_eq\ \{Ω\ :\ Type*\}\ [MeasurableSpace\ Ω]}
\leanline{4}{(μ\ :\ Measure\ Ω)\ (F\ :\ ℕ\ →\ Ω\ →\ ℝ)\ (G\ B\ :\ Ω\ →\ ℝ)}
\leanline{4}{(hF\ :\ ∀\ n,\ AEStronglyMeasurable\ (F\ n)\ μ)\ (hG\ :\ AEStronglyMeasurable\ G\ μ)}
\leanline{4}{(hB\ :\ Integrable\ (fun\ a\ ↦\ B\ a\ \textasciicircum{}\ 2)\ μ)}
\leanline{4}{(hbound\ :\ ∀\ n,\ ∀ᵐ\ a\ ∂μ,\ |F\ n\ a|\ ≤\ B\ a)}
\leanline{4}{(hGbound\ :\ ∀ᵐ\ a\ ∂μ,\ |G\ a|\ ≤\ B\ a)}
\leanline{4}{(heq\ :\ ∀ᵐ\ a\ ∂μ,\ ∀ᶠ\ n\ :\ ℕ\ in\ atTop,\ F\ n\ a\ =\ G\ a)\ :}
\leanline{4}{Tendsto\ (fun\ n\ ↦\ ∫\ a,\ (F\ n\ a\ -\ G\ a)\ \textasciicircum{}\ 2\ ∂μ)\ atTop\ ({\leanfallback 𝓝}\ 0)}
\end{leancode}
\leansource{\texttt{continuous\_\allowbreak{}langevinGenerator\_\allowbreak{}test} \quad AverageLMC/\allowbreak{}GibbsScoreCutoff.lean, lines 180–191}
\begin{leancode}
\leanline{0}{theorem\ continuous\_langevinGenerator\_test\ \{p\ :\ ℕ\}\ (f\ ψ\ :\ State\ p\ →\ ℝ)\ (M\ :\ ℝ≥0)}
\leanline{4}{(hLip\ :\ LipschitzWith\ M\ (gradient\ f))\ (hψ\ :\ ContDiff\ ℝ\ 2\ ψ)\ :}
\leanline{4}{Continuous\ (ULA.langevinGenerator\ f\ ψ)}
\end{leancode}
\leansource{\texttt{gradient\_\allowbreak{}score\_\allowbreak{}coordinate\_\allowbreak{}norm\_\allowbreak{}le} \quad AverageLMC/\allowbreak{}GibbsScoreCutoff.lean, lines 193–204}
\begin{leancode}
\leanline{0}{theorem\ gradient\_score\_coordinate\_norm\_le\ \{p\ :\ ℕ\}\ (f\ :\ State\ p\ →\ ℝ)\ (M\ :\ ℝ≥0)}
\leanline{4}{(hLip\ :\ LipschitzWith\ M\ (gradient\ f))\ (j\ :\ Fin\ p)\ (x\ :\ State\ p)\ :}
\leanline{4}{‖gradient\ (fun\ y\ ↦\ (gradient\ f\ y)\ j)\ x‖\ ≤\ (M\ :\ ℝ)}
\end{leancode}
\leanlabel{How the Lean proof follows the argument}
\texttt{integral\_\allowbreak{}sq\_\allowbreak{}sub\_\allowbreak{}tendsto\_\allowbreak{}zero\_\allowbreak{}of\_\allowbreak{}eventually\_\allowbreak{}eq} gives the formal counterpart of this step. The main intermediate facts are \texttt{hd} in that order. Named proof applications include \texttt{hs}, \texttt{tendsto\_\allowbreak{}const\_\allowbreak{}nhds.\allowbreak{}congr'}; the adjacent proof shows their hypotheses and order of use.




\texttt{continuous\_\allowbreak{}langevin\allowbreak{}Generator\_\allowbreak{}test} gives the formal counterpart of this step. The main intermediate facts are \texttt{heq}, \texttt{hgrad} in that order.




\texttt{gradient\_\allowbreak{}score\_\allowbreak{}coordinate\_\allowbreak{}norm\_\allowbreak{}le} gives the formal counterpart of this step. The main intermediate facts are \texttt{h\ensuremath{\psi}\allowbreak{}Lip} in that order. Named proof applications include \texttt{Lipschitz\allowbreak{}With.\allowbreak{}of\_\allowbreak{}dist\_\allowbreak{}le\_\allowbreak{}mul}, \texttt{norm\_\allowbreak{}fderiv\_\allowbreak{}le\_\allowbreak{}of\_\allowbreak{}lipschitz}; the adjacent proof shows their hypotheses and order of use.
\end{leanbox}


Earlier results used here: \cite[Section~3.46.14]{LangevinLean2026ULASC}; \cite[Section~3.46.2]{LangevinLean2026ULASC}.
\subsection{\texorpdfstring{Control the cutoff generator without a bounded third derivative}{Control the cutoff generator without a bounded third derivative}}
\label{report-step-AverageLMC-GibbsScoreCutoff-lean-6}


The generator product rule gives the cutoff times the original score generator, the score coordinate times the cutoff generator, and twice the inner product of their gradients. The true score-coordinate generator is the corresponding coordinate of \(\nabla\Delta f-D(\nabla f)\nabla f\). The first term is bounded by the printed Laplacian-gradient condition and the second by \(M\) times the score norm. Uniform cutoff gradient and Laplacian bounds control the remaining terms by a common constant times one plus the squared score norm. This is a genuine Gibbs \(L^{2}\) majorant by the already derived fourth score moment.

\begin{leanbox}{Lean: \texttt{scoreCutoff\_\allowbreak{}generator\_\allowbreak{}bound}}
\leansource{\texttt{scoreCutoff\_\allowbreak{}generator\_\allowbreak{}bound} \quad AverageLMC/\allowbreak{}GibbsScoreCutoff.lean, lines 208–268}
\begin{leancode}
\leanline{0}{theorem\ scoreCutoff\_generator\_bound\ \{p\ :\ ℕ\}\ (f\ :\ State\ p\ →\ ℝ)\ (M\ :\ ℝ≥0)}
\leanline{4}{(hLip\ :\ LipschitzWith\ M\ (gradient\ f))\ (Mdelta\ :\ ℝ)\ (hMdelta\ :\ 0\ ≤\ Mdelta)}
\leanline{4}{(hδ\ :\ LaplacianSmooth\ f\ Mdelta)\ (j\ :\ Fin\ p)\ :}
\leanline{4}{∃\ B\ :\ ℝ,\ 0\ ≤\ B\ ∧\ ∀\ n\ x,}
\leanline{6}{|ULA.langevinGenerator\ f\ (fun\ y\ ↦\ radialCutoff\ n\ y\ *\ (gradient\ f\ y)\ j)\ x|\ ≤}
\leanline{8}{B\ *\ (1\ +\ ‖gradient\ f\ x‖\ \textasciicircum{}\ 2)}
\end{leancode}
\leanlabel{How the Lean proof follows the argument}
\texttt{score\allowbreak{}Cutoff\_\allowbreak{}generator\_\allowbreak{}bound} gives the formal counterpart of this step. The main intermediate facts are \texttt{D}, \texttt{B}, \texttt{h\allowbreak{}D}, \texttt{h\allowbreak{}B}, \texttt{s}, \texttt{hs} in that order. Named proof applications include \texttt{Pi\allowbreak{}Lp.\allowbreak{}norm\_\allowbreak{}apply\_\allowbreak{}le}, \texttt{hraw.\allowbreak{}trans}; the adjacent proof shows their hypotheses and order of use.
\end{leanbox}


Earlier results used here: \cite[Section~3.46.14]{LangevinLean2026ULASC}; \cite[Section~3.46.2]{LangevinLean2026ULASC}.
\subsection{\texorpdfstring{Produce compact approximations for the real score Dynkin formula}{Produce compact approximations for the real score Dynkin formula}}
\label{report-step-AverageLMC-GibbsScoreCutoff-lean-7}


For each coordinate \(j\), choose the compact \(C^{2}\) tests equal to the cutoff times the \(j\)th actual gradient coordinate. The score cutoffs are dominated by the score norm. Their true generators are dominated by the common \(L^{2}\) majorant just constructed. Both sequences eventually agree pointwise with their targets. Applying the squared-integral dominated convergence principle twice yields convergence in Gibbs \(L^{2}\) of the tests to the coordinate score and of their generators to the corresponding coordinate of the actual gradient generator. Every normalization and moment assumption needed here follows from finite Gibbs partition, Lipschitz gradient, and the source \(C^{3}\) Laplacian-gradient condition.

\begin{leanbox}{Lean: \texttt{score\_\allowbreak{}coordinate\_\allowbreak{}compact\_\allowbreak{}approximation}}
\leansource{\texttt{score\_\allowbreak{}coordinate\_\allowbreak{}compact\_\allowbreak{}approximation} \quad AverageLMC/\allowbreak{}GibbsScoreCutoff.lean, lines 273–341}
\begin{leancode}
\leanline{0}{theorem\ score\_coordinate\_compact\_approximation\ \{p\ :\ ℕ\}}
\leanline{4}{(f\ :\ State\ p\ →\ ℝ)\ (M\ :\ ℝ≥0)\ (hLip\ :\ LipschitzWith\ M\ (gradient\ f))}
\leanline{4}{(hZ\ :\ ULA.gibbsPartition\ f\ ≠\ ⊤)\ (Mdelta\ :\ ℝ)\ (hMdelta\ :\ 0\ ≤\ Mdelta)}
\leanline{4}{(hδ\ :\ LaplacianSmooth\ f\ Mdelta)\ (j\ :\ Fin\ p)\ :}
\leanline{4}{∃\ φ\ :\ ℕ\ →\ State\ p\ →\ ℝ,}
\leanline{6}{(∀\ n,\ ContDiff\ ℝ\ 2\ (φ\ n)\ ∧\ HasCompactSupport\ (φ\ n))\ ∧}
\leanline{6}{Tendsto\ (fun\ n\ ↦\ ∫\ x,\ (φ\ n\ x\ -\ (gradient\ f\ x)\ j)\ \textasciicircum{}\ 2\ ∂ULA.gibbsMeasure\ f)}
\leanline{8}{atTop\ ({\leanfallback 𝓝}\ 0)\ ∧}
\leanline{6}{Tendsto\ (fun\ n\ ↦\ ∫\ x,\ (ULA.langevinGenerator\ f\ (φ\ n)\ x\ -}
\leanline{8}{(gradientGenerator\ f\ x)\ j)\ \textasciicircum{}\ 2\ ∂ULA.gibbsMeasure\ f)\ atTop\ ({\leanfallback 𝓝}\ 0)}
\end{leancode}
\leanlabel{How the Lean proof follows the argument}
\texttt{score\_\allowbreak{}coordinate\_\allowbreak{}compact\_\allowbreak{}approximation} gives the formal counterpart of this step. The main intermediate facts are \texttt{hf}, \texttt{\ensuremath{\psi}}, \texttt{\ensuremath{\varphi}}, \texttt{h\ensuremath{\psi}}, \texttt{h\ensuremath{\varphi}} in that order. Named proof applications include \texttt{Pi\allowbreak{}Lp.\allowbreak{}norm\_\allowbreak{}apply\_\allowbreak{}le}, \texttt{mem\allowbreak{}Lp\_\allowbreak{}gradient\_\allowbreak{}gibbs\_\allowbreak{}all}, \texttt{hb}; the adjacent proof shows their hypotheses and order of use.
\end{leanbox}


Earlier results used here: \cite[Section~3.26.2]{LangevinLean2026ULASC}; \cite[Section~3.26.3]{LangevinLean2026ULASC}; \cite[Section~3.26.9]{LangevinLean2026ULASC}; \cite[Section~3.46.2]{LangevinLean2026ULASC}; \cite[Section~3.11.1]{LangevinLean2026IPLA}.
\subsection{\texorpdfstring{Assemble one vector approximation with explicit \(L^{2}\) membership}{Assemble one vector approximation with explicit L² membership}}
\label{report-step-AverageLMC-GibbsScoreCutoff-lean-8}


Choose the proved scalar approximation for each of the finitely many coordinates and place them in the Euclidean coordinate vector. Each compact test and its compact continuous generator is bounded, hence belongs to Gibbs \(L^{2}\). The target generator belongs to Gibbs \(L^{2}\) because its norm is at most \(M\) times the true score norm plus the fixed Laplacian-gradient bound. Expanding the squared Euclidean error as a finite sum of squared coordinate errors identifies its integral with the sum of the scalar squared integrals. Finite-sum limit continuity proves convergence of both vector errors to zero. The packaged result includes \(L^{2}\) membership of every vector approximant and both actual targets, so later Cauchy-Schwarz arguments do not assume an unproved moment.

\begin{leanbox}{Lean: \texttt{integral\_\allowbreak{}norm\_\allowbreak{}sq\_\allowbreak{}toLp\_\allowbreak{}sub\_\allowbreak{}tendsto\_\allowbreak{}zero}, \texttt{gradientGenerator\_\allowbreak{}gibbs\_\allowbreak{}memLp}, \texttt{score\_\allowbreak{}vector\_\allowbreak{}compact\_\allowbreak{}approximation}}
\leansource{\texttt{integral\_\allowbreak{}norm\_\allowbreak{}sq\_\allowbreak{}toLp\_\allowbreak{}sub\_\allowbreak{}tendsto\_\allowbreak{}zero} \quad AverageLMC/\allowbreak{}GibbsScoreCutoff.lean, lines 343–359}
\begin{leancode}
\leanline{0}{theorem\ integral\_norm\_sq\_toLp\_sub\_tendsto\_zero\ \{Ω\ :\ Type*\}\ [MeasurableSpace\ Ω]}
\leanline{4}{\{p\ :\ ℕ\}\ (μ\ :\ Measure\ Ω)\ (F\ :\ ℕ\ →\ Fin\ p\ →\ Ω\ →\ ℝ)\ (G\ :\ Ω\ →\ State\ p)}
\leanline{4}{(hF\ :\ ∀\ n,\ MemLp\ (fun\ a\ ↦\ WithLp.toLp\ 2\ (fun\ j\ ↦\ F\ n\ j\ a))\ 2\ μ)}
\leanline{4}{(hG\ :\ MemLp\ G\ 2\ μ)}
\leanline{4}{(hconv\ :\ ∀\ j,\ Tendsto\ (fun\ n\ ↦\ ∫\ a,\ (F\ n\ j\ a\ -\ (G\ a)\ j)\ \textasciicircum{}\ 2\ ∂μ)\ atTop\ ({\leanfallback 𝓝}\ 0))\ :}
\leanline{4}{Tendsto\ (fun\ n\ ↦\ ∫\ a,\ ‖WithLp.toLp\ 2\ (fun\ j\ ↦\ F\ n\ j\ a)\ -\ G\ a‖\ \textasciicircum{}\ 2\ ∂μ)}
\leanline{6}{atTop\ ({\leanfallback 𝓝}\ 0)}
\end{leancode}
\leansource{\texttt{gradientGenerator\_\allowbreak{}gibbs\_\allowbreak{}memLp} \quad AverageLMC/\allowbreak{}GibbsScoreCutoff.lean, lines 361–388}
\begin{leancode}
\leanline{0}{theorem\ gradientGenerator\_gibbs\_memLp\ \{p\ :\ ℕ\}}
\leanline{4}{(f\ :\ State\ p\ →\ ℝ)\ (M\ :\ ℝ≥0)\ (hLip\ :\ LipschitzWith\ M\ (gradient\ f))}
\leanline{4}{(hZ\ :\ ULA.gibbsPartition\ f\ ≠\ ⊤)\ (Mdelta\ :\ ℝ)\ (hMdelta\ :\ 0\ ≤\ Mdelta)}
\leanline{4}{(hδ\ :\ LaplacianSmooth\ f\ Mdelta)\ :}
\leanline{4}{MemLp\ (gradientGenerator\ f)\ 2\ (ULA.gibbsMeasure\ f)}
\end{leancode}
\leansource{\texttt{score\_\allowbreak{}vector\_\allowbreak{}compact\_\allowbreak{}approximation} \quad AverageLMC/\allowbreak{}GibbsScoreCutoff.lean, lines 393–442}
\begin{leancode}
\leanline{0}{theorem\ score\_vector\_compact\_approximation\ \{p\ :\ ℕ\}}
\leanline{4}{(f\ :\ State\ p\ →\ ℝ)\ (M\ :\ ℝ≥0)\ (hLip\ :\ LipschitzWith\ M\ (gradient\ f))}
\leanline{4}{(hZ\ :\ ULA.gibbsPartition\ f\ ≠\ ⊤)\ (Mdelta\ :\ ℝ)\ (hMdelta\ :\ 0\ ≤\ Mdelta)}
\leanline{4}{(hδ\ :\ LaplacianSmooth\ f\ Mdelta)\ :}
\leanline{4}{∃\ φ\ :\ ℕ\ →\ Fin\ p\ →\ State\ p\ →\ ℝ,}
\leanline{6}{(∀\ n\ j,\ ContDiff\ ℝ\ 2\ (φ\ n\ j)\ ∧\ HasCompactSupport\ (φ\ n\ j))\ ∧}
\leanline{6}{(∀\ n,\ MemLp\ (fun\ x\ ↦\ WithLp.toLp\ 2\ (fun\ j\ ↦\ φ\ n\ j\ x))\ 2\ (ULA.gibbsMeasure\ f))\ ∧}
\leanline{6}{(∀\ n,\ MemLp\ (fun\ x\ ↦\ WithLp.toLp\ 2\ (fun\ j\ ↦}
\leanline{8}{ULA.langevinGenerator\ f\ (φ\ n\ j)\ x))\ 2\ (ULA.gibbsMeasure\ f))\ ∧}
\leanline{6}{MemLp\ (gradient\ f)\ 2\ (ULA.gibbsMeasure\ f)\ ∧}
\leanline{6}{MemLp\ (gradientGenerator\ f)\ 2\ (ULA.gibbsMeasure\ f)\ ∧}
\leanline{6}{Tendsto\ (fun\ n\ ↦\ ∫\ x,\ ‖WithLp.toLp\ 2\ (fun\ j\ ↦\ φ\ n\ j\ x)\ -\ gradient\ f\ x‖\ \textasciicircum{}\ 2}
\leanline{8}{∂ULA.gibbsMeasure\ f)\ atTop\ ({\leanfallback 𝓝}\ 0)\ ∧}
\leanline{6}{Tendsto\ (fun\ n\ ↦\ ∫\ x,\ ‖WithLp.toLp\ 2\ (fun\ j\ ↦\ ULA.langevinGenerator\ f\ (φ\ n\ j)\ x)\ -}
\leanline{8}{gradientGenerator\ f\ x‖\ \textasciicircum{}\ 2\ ∂ULA.gibbsMeasure\ f)\ atTop\ ({\leanfallback 𝓝}\ 0)}
\end{leancode}
\leanlabel{How the Lean proof follows the argument}
\texttt{integral\_\allowbreak{}norm\_\allowbreak{}sq\_\allowbreak{}to\allowbreak{}Lp\_\allowbreak{}sub\_\allowbreak{}tendsto\_\allowbreak{}zero} gives the formal counterpart of this step. The main intermediate facts are \texttt{heq} in that order. Named proof applications include \texttt{tendsto\_\allowbreak{}finset\allowbreak{}Sum}; the adjacent proof shows their hypotheses and order of use.




\texttt{gradient\allowbreak{}Generator\_\allowbreak{}gibbs\_\allowbreak{}mem\allowbreak{}Lp} gives the formal counterpart of this step. The main intermediate facts are \texttt{hf}, \texttt{hs}, \texttt{hb0}, \texttt{hb} in that order. Named proof applications include \texttt{ENNReal.\allowbreak{}inv\_\allowbreak{}mul\_\allowbreak{}cancel}, \texttt{hb.\allowbreak{}of\_\allowbreak{}le\_\allowbreak{}mul}, \texttt{ae\_\allowbreak{}of\_\allowbreak{}all}; the adjacent proof shows their hypotheses and order of use.




\texttt{score\_\allowbreak{}vector\_\allowbreak{}compact\_\allowbreak{}approximation} gives the formal counterpart of this step. The main intermediate facts are \texttt{hf}, \texttt{Φ}, \texttt{h\allowbreak{}Φ}, \texttt{h\allowbreak{}GΦ}, \texttt{hscore2}, \texttt{hgen2} in that order. Named proof applications include \texttt{ENNReal.\allowbreak{}inv\_\allowbreak{}mul\_\allowbreak{}cancel}, \texttt{mem\allowbreak{}Lp\_\allowbreak{}pi\allowbreak{}Lp\_\allowbreak{}iff.\allowbreak{}mpr}, \texttt{Mem\allowbreak{}Lp.\allowbreak{}of\_\allowbreak{}bound}; the adjacent proof shows their hypotheses and order of use.
\end{leanbox}


Earlier results used here: \cite[Section~3.26.2]{LangevinLean2026ULASC}; \cite[Section~3.26.3]{LangevinLean2026ULASC}; \cite[Section~3.26.9]{LangevinLean2026ULASC}; \cite[Section~3.46.2]{LangevinLean2026ULASC}.
\section{\texorpdfstring{Derive the actual stationary score-increment pairing under the printed \(C^{3}\) condition}{Derive the actual stationary score-increment pairing under the printed C³ condition}}
\label{proof-module-AverageLMC-LaplacianPairing-lean}


This module supplies the stochastic pairing estimate needed in the Laplacian one-step proof. It starts from the actual Brownian Langevin integral equation, Gaussian-innovation independence, and Gibbs stationarity. Compact \(C^{2}\) score approximations satisfy the proved weighted Dynkin identity. Their genuine generator errors have a stationary \(L^{2}\) norm that is uniform in time. Taking the compact approximation limit and then the initial-weight clipping limit yields the sharp coefficient \(M_\infty\sqrt{pM_{\mathrm{av}}}+M_\Delta\sqrt p\). Only the source \(C^{3}\) potential and its bound on \(\nabla\Delta f\) are used.



Stationarity transfers every genuine Gibbs squared approximation error to the actual observation:


\[\|F_n(Y_t)-\nabla f(Y_t)\|_{L^2(P)}^2
=\int\|F_n(x)-\nabla f(x)\|^2\,\pi(dx)\longrightarrow0.\]


Applying the compact Dynkin identities, the generator approximation, and radial clipping proves the sharp cross-term bound used in Theorem 2:


\[\left\langle\Psi,\nabla f(Y_T)-\nabla f(Y_0)\right\rangle_{L^2(P)}
\le T\|\Psi\|_{L^2(P)}
\left(M_\infty\sqrt{M_{\mathrm{av}}p}+M_\Delta\sqrt p\right).\]\subsection{\texorpdfstring{Transfer Gibbs squared-error convergence through an actual stationary law}{Transfer Gibbs squared-error convergence through an actual stationary law}}
\label{report-step-AverageLMC-LaplacianPairing-lean-1}


If a random \(L^{2}\) state has law \(\pi\), the squared norm of the difference of its two canonical observations is exactly the integral of the squared pointwise error under \(\pi\). Taking square roots identifies the \(L^{2}\) distance. Continuity of the square root at zero therefore turns convergence of the genuine squared-error integrals into convergence of the actual \(L^{2}\) observations. This theorem uses the displayed map-law equality and the square-integrability witnesses, rather than assuming convergence after composition.

\begin{leanbox}{Lean: \texttt{stationaryObservation\_\allowbreak{}tendsto\_\allowbreak{}of\_\allowbreak{}sq\_\allowbreak{}error}}
\leansource{\texttt{stationaryObservation\_\allowbreak{}tendsto\_\allowbreak{}of\_\allowbreak{}sq\_\allowbreak{}error} \quad AverageLMC/\allowbreak{}LaplacianPairing.lean, lines 29–44}
\begin{leancode}
\leanline{0}{theorem\ stationaryObservation\_tendsto\_of\_sq\_error\ \{p\ :\ ℕ\}}
\leanline{4}{(π\ :\ Measure\ (State\ p))\ (F\ :\ ℕ\ →\ State\ p\ →\ State\ p)\ (G\ :\ State\ p\ →\ State\ p)}
\leanline{4}{(hFc\ :\ ∀\ n,\ Continuous\ (F\ n))\ (hGc\ :\ Continuous\ G)}
\leanline{4}{(hF\ :\ ∀\ n,\ MemLp\ (F\ n)\ 2\ π)\ (hG\ :\ MemLp\ G\ 2\ π)}
\leanline{4}{(Y\ :\ Lp\ (State\ p)\ 2\ P)\ (hlaw\ :\ P.map\ Y\ =\ π)}
\leanline{4}{(herr\ :\ Tendsto\ (fun\ n\ ↦\ ∫\ x,\ ‖F\ n\ x\ -\ G\ x‖\ \textasciicircum{}\ 2\ ∂π)\ atTop\ ({\leanfallback 𝓝}\ 0))\ :}
\leanline{4}{Tendsto\ (fun\ n\ ↦\ stationaryObservation\ P\ π\ (F\ n)\ (hF\ n)\ Y\ hlaw)\ atTop}
\leanline{6}{({\leanfallback 𝓝}\ (stationaryObservation\ P\ π\ G\ hG\ Y\ hlaw))}
\end{leancode}
\leanlabel{How the Lean proof follows the argument}
\texttt{stationary\allowbreak{}Observation\_\allowbreak{}tendsto\_\allowbreak{}of\_\allowbreak{}sq\_\allowbreak{}error} gives the formal counterpart of this step. The main intermediate facts are \texttt{he} in that order. Its remaining equalities are discharged by the simplifications displayed in the proof.
\end{leanbox}
\subsection{\texorpdfstring{Construct and pass the actual compact Dynkin pairings to the score increment}{Construct and pass the actual compact Dynkin pairings to the score increment}}
\label{report-step-AverageLMC-LaplacianPairing-lean-2}


Strong convexity normalizes the Gibbs measure and gives its finite partition. The integral equation at time zero identifies the initial state with \(Y(0)\), so stationarity supplies its Gibbs law as well as that of the final state. The proved radial score-cutoff construction provides compact \(C^{2}\) scalar tests in every coordinate, with explicit \(L^{2}\) membership of their vector tests and actual Langevin generators. Their squared Gibbs errors tend to zero against the actual gradient and gradient generator. The observation convergence theorem yields both endpoint \(L^{2}\) limits.



The gradient generator is the actual vector \(\nabla\Delta f-D(\nabla f)\nabla f\). Its stationary \(L^{2}\) norm is at most \(R=M_\infty\sqrt{pM_{\mathrm{av}}}+M_\Delta\sqrt p\), by the exact average-smoothness Gibbs score bound and the printed Laplacian-gradient condition. At every time the observation of an approximating generator differs from this target by the same Gibbs \(L^{2}\) error \(\operatorname{err}(n)\), independent of time. The triangle inequality therefore bounds every approximating generator norm by \(R+\operatorname{err}(n)\), and \(\operatorname{err}(n)\) tends to zero. The cutoff construction needs only \(C^{2}\) regularity of the score, which follows from \(C^{3}\) regularity of \(f\).



Radially clip the initial vector \(\Psi\). The clipped vector has bounded coordinates and remains independent of future Gaussian innovations jointly with the initial state, because measurable maps preserve independence. Apply the actual compact-test weighted Dynkin theorem to every scalar coordinate weight and its compact \(C^{2}\) score approximation. Compact generators are bounded and uniformly continuous, so these weighted generator paths are interval integrable. Sum the coordinate identities, exchange only the finite sum with the time integral, and identify the sum with the exact Euclidean \(L^{2}\) pairing.



For a fixed clipped weight, apply the compact Dynkin approximation limit theorem to the endpoint \(L^{2}\) limits and the uniform generator bound \(R+\operatorname{err}(n)\). This gives the pairing of the clipped weight with the actual score increment at most \(T\) times its \(L^{2}\) norm times \(R\). Finally apply the clipping limit theorem: the pairings converge, and the clipped norms never exceed the original \(L^{2}\) norm. The resulting stationary gradient pairing bound applies to every square-integrable independent initial vector. Its stochastic hypotheses concern the actual Brownian integral equation and independence; no desired Wasserstein recurrence or score-increment bound is supplied as a premise.

\begin{leanbox}{Lean: \texttt{lp\_\allowbreak{}stationary\_\allowbreak{}gradient\_\allowbreak{}pairing\_\allowbreak{}bound}}
\leansource{\texttt{lp\_\allowbreak{}stationary\_\allowbreak{}gradient\_\allowbreak{}pairing\_\allowbreak{}bound} \quad AverageLMC/\allowbreak{}LaplacianPairing.lean, lines 50–210}
\begin{leancode}
\leanline{0}{theorem\ lp\_stationary\_gradient\_pairing\_bound\ \{p\ :\ ℕ\}\ (hp\ :\ 0\ <\ p)}
\leanline{4}{(f\ :\ State\ p\ →\ ℝ)\ (m\ :\ ℝ)\ (L\ :\ ℝ≥0)\ (hm\ :\ 0\ <\ m)}
\leanline{4}{(hsc\ :\ StrongConvexOn\ univ\ m\ f)\ (hLip\ :\ LipschitzWith\ L\ (gradient\ f))}
\leanline{4}{(Mdelta\ :\ ℝ)\ (hMdelta\ :\ 0\ ≤\ Mdelta)\ (hδ\ :\ LaplacianSmooth\ f\ Mdelta)}
\leanline{4}{(initial\ :\ Lp\ (State\ p)\ 2\ P)}
\leanline{4}{(B\ :\ Fin\ p\ →\ ℝ≥0\ →\ Ω\ →\ ℝ)\ (hB\ :\ ∀\ i,\ IsPreBrownianReal\ (B\ i)\ P)}
\leanline{4}{(hscalar\ :\ ∀\ r\ :\ ℝ,\ 0\ <\ r\ →\ iIndepFun\ (fun\ q\ :\ ℕ\ ×\ Fin\ p\ ↦\ fun\ a\ ↦}
\leanline{6}{ULA.Hilbert.brownianEulerNoise\ P\ B\ hB\ r\ q.1\ a\ q.2)\ P)}
\leanline{4}{(Ψ\ :\ Lp\ (State\ p)\ 2\ P)}
\leanline{4}{(hind\ :\ ∀\ r\ :\ ℝ,\ 0\ <\ r\ →\ IndepFun\ (fun\ a\ ↦\ (Ψ\ a,\ initial\ a))}
\leanline{6}{(fun\ a\ k\ ↦\ ULA.Hilbert.brownianEulerNoise\ P\ B\ hB\ r\ k\ a)\ P)}
\leanline{4}{(Y\ :\ ℝ\ →\ Lp\ (State\ p)\ 2\ P)\ (hY\ :\ Continuous\ Y)}
\leanline{4}{(T\ :\ ℝ)\ (hT\ :\ 0\ ≤\ T)}
\leanline{4}{(hYe\ :\ ∀\ t\ ∈\ Icc\ 0\ T,\ Y\ t\ -\ initial\ =}
\leanline{6}{-(∫\ s\ in\ 0..t,\ ULA.lpDrift\ P\ (gradient\ f)\ hLip\ (Y\ s))\ +}
\leanline{8}{ULA.brownianLangevinNoise\ P\ B\ hB\ 0\ t)}
\leanline{4}{(hstationary\ :\ ∀\ t\ ∈\ Icc\ 0\ T,\ P.map\ (Y\ t)\ =\ ULA.gibbsMeasure\ f)\ :}
\leanline{4}{inner\ ℝ\ Ψ\ (ULA.lpDrift\ P\ (gradient\ f)\ hLip\ (Y\ T)\ -}
\leanline{6}{ULA.lpDrift\ P\ (gradient\ f)\ hLip\ initial)\ ≤}
\leanline{8}{T\ *\ ‖Ψ‖\ *\ ((L\ :\ ℝ)\ *\ Real.sqrt\ (averageSmoothness\ f\ *\ (p\ :\ ℝ))\ +}
\leanline{10}{Mdelta\ *\ Real.sqrt\ (p\ :\ ℝ))}
\end{leancode}
\leanlabel{How the Lean proof follows the argument}
\textbf{Formal correspondence.} The resulting estimate is \texttt{lp\_\allowbreak{}stationary\_\allowbreak{}gradient\_\allowbreak{}pairing\_\allowbreak{}bound}.



\texttt{lp\_\allowbreak{}stationary\_\allowbreak{}gradient\_\allowbreak{}pairing\_\allowbreak{}bound} gives the formal counterpart of this step. The main intermediate facts are \texttt{hf}, \texttt{h\allowbreak{}Z}, \texttt{\ensuremath{\pi}}, \texttt{h\allowbreak{}Y0}, \texttt{hinitlaw}, \texttt{h\allowbreak{}Tlaw} in that order. Named proof applications include \texttt{interval\allowbreak{}Integral.\allowbreak{}integral\_\allowbreak{}congr\_\allowbreak{}ae\_\allowbreak{}restrict}, \texttt{\ensuremath{\langle}ht.\allowbreak{}1.\allowbreak{}le}, \texttt{pairing\_\allowbreak{}bound\_\allowbreak{}of\_\allowbreak{}compact\_\allowbreak{}dynkin\_\allowbreak{}approximations}; the adjacent proof shows their hypotheses and order of use.
\end{leanbox}


Earlier results used here: \cite[Section~3.13.2]{LangevinLean2026ULASC}; \cite[Section~3.13.3]{LangevinLean2026ULASC}; \cite[Section~3.25.1]{LangevinLean2026ULASC}; \cite[Section~3.25.3]{LangevinLean2026ULASC}; \cite[Section~3.26.10]{LangevinLean2026ULASC}; \cite[Section~3.26.13]{LangevinLean2026ULASC}; \cite[Section~3.26.3]{LangevinLean2026ULASC}; \cite[Section~3.29.7]{LangevinLean2026ULASC}; \cite[Section~3.46.2]{LangevinLean2026ULASC}.
\section{\texorpdfstring{The canonical stationary coupling for the third-order LMC estimate}{The canonical stationary coupling for the third-order LMC estimate}}
\label{proof-module-AverageLMC-LaplacianCanonical-lean}


The third-order coupling uses the same actual independent Gaussian update and constructed Gibbs-stationary Langevin diffusion. Every dual pairing premise is obtained from the weighted Dynkin theorem on the explicit Brownian extension.



For a finite-cost initial coupling \(\gamma\), set \(C_\gamma^2=\int\|x-y\|^2\,d\gamma\) and \(q=1-mh\). Weighted Dynkin, the Gibbs score moment and the actual Gaussian coupling establish


\[R_{\rm gen}=L\sqrt{M_{\rm av}p}+M_\Delta\sqrt p.\]
\[W_2^2(\nu^+,\pi)\le q^2C_\gamma^2+h^2qC_\gamma\bigl(L\sqrt{M_{\rm av}p}+M_\Delta\sqrt p\bigr)+\left(\frac32Lh\sqrt h\sqrt p\right)^2.\]

Earlier results used here: \cite[Section~3.45]{LangevinLean2026ULASC}.
\subsection{\texorpdfstring{Construct the initial-pair independence and derive the one-step transport estimate}{Construct the initial-pair independence and derive the one-step transport estimate}}
\label{report-step-AverageLMC-LaplacianCanonical-lean-1}


Extend each finite-cost coupling by the canonical Brownian probability space and construct the continuous integral solution started at its Gibbs coordinate. Stationarity, exact Brownian moments and standard Gaussian innovations follow from their existing constructions. The deterministic gradient-step difference \(\Psi\) and the initial reference coordinate are a measurable function of the coupled initial pair. The product measure therefore makes this pair independent of every full Brownian Euler innovation sequence; almost-everywhere \(L^{2}\) representatives preserve that independence. Apply the stationary gradient pairing theorem on each restricted interval [0,\(t\)] to obtain the cross-term bound with the true constant \(L\sqrt{M_{\mathrm{av}}p}+M_\Delta\sqrt p\). The local one-step energy theorem supplies the contracted error, exact cross coefficient and total-defect square. The endpoint Euler random vector is proved to have the ordinary Gaussian LMC law, while the reference endpoint has the Gibbs law. Their actual \(L^{2}\) coupling bounds Wasserstein distance. The exact initial-coupling cost identity then gives the claimed quadratic-plus-linear estimate for every finite initial coupling.

\begin{leanbox}{Lean: \texttt{laplacian\_\allowbreak{}coupling\_\allowbreak{}step}}
\leansource{\texttt{laplacian\_\allowbreak{}coupling\_\allowbreak{}step} \quad AverageLMC/\allowbreak{}LaplacianCanonical.lean, lines 18–183}
\begin{leancode}
\leanline{0}{theorem\ laplacian\_coupling\_step\ \{p\ :\ ℕ\}\ (hp\ :\ 0\ <\ p)}
\leanline{4}{(f\ :\ State\ p\ →\ ℝ)\ (m\ h\ :\ ℝ)\ (L\ :\ ℝ≥0)\ (hm\ :\ 0\ <\ m)\ (hh\ :\ 0\ <\ h)}
\leanline{4}{(hf\ :\ Differentiable\ ℝ\ f)\ (hsc\ :\ StrongConvexOn\ univ\ m\ f)}
\leanline{4}{(hLip\ :\ LipschitzWith\ L\ (gradient\ f))\ (hsmall\ :\ h\ *\ (L\ :\ ℝ)\ ≤\ 1)}
\leanline{4}{(Mdelta\ :\ ℝ)\ (hMdelta\ :\ 0\ ≤\ Mdelta)\ (hδ\ :\ LaplacianSmooth\ f\ Mdelta)}
\leanline{4}{(μ\ :\ Measure\ (State\ p))\ [IsProbabilityMeasure\ μ]}
\leanline{4}{(c\ :\ Coupling\ μ\ (gibbsMeasure\ f))\ (hc\ :\ couplingCost\ c.joint\ ≠\ ⊤)\ :}
\leanline{4}{wasserstein2Real\ (ulaLawStep\ f\ h\ μ)\ (gibbsMeasure\ f)\ \textasciicircum{}\ 2\ ≤}
\leanline{6}{(1\ -\ m\ *\ h)\ \textasciicircum{}\ 2\ *\ (couplingCost\ c.joint).toReal\ \textasciicircum{}\ 2\ +}
\leanline{6}{h\ \textasciicircum{}\ 2\ *\ (1\ -\ m\ *\ h)\ *\ (couplingCost\ c.joint).toReal\ *}
\leanline{8}{((L\ :\ ℝ)\ *\ Real.sqrt\ (averageSmoothness\ f\ *\ p)\ +\ Mdelta\ *\ Real.sqrt\ p)\ +}
\leanline{6}{((3\ /\ 2\ :\ ℝ)\ *\ (L\ :\ ℝ)\ *\ h\ *\ Real.sqrt\ h\ *\ Real.sqrt\ p)\ \textasciicircum{}\ 2}
\end{leancode}
\leanlabel{How the Lean proof follows the argument}
\texttt{laplacian\_\allowbreak{}coupling\_\allowbreak{}step} gives the formal counterpart of this step. The main intermediate facts are \texttt{Q}, \texttt{P}, \texttt{h\ensuremath{\pi}}, \texttt{X}, \texttt{Z}, \texttt{W} in that order. Named proof applications include \texttt{hi}, \texttt{hi'.\allowbreak{}comp}, \texttt{brownian\allowbreak{}Euler\allowbreak{}Noise\_\allowbreak{}endpoint}; the adjacent proof shows their hypotheses and order of use.
\end{leanbox}


Earlier results used here: \cite[Section~3.13.11]{LangevinLean2026ULASC}; \cite[Section~3.13.12]{LangevinLean2026ULASC}; \cite[Section~3.13.2]{LangevinLean2026ULASC}; \cite[Section~3.13.3]{LangevinLean2026ULASC}; \cite[Section~3.17.1]{LangevinLean2026ULASC}; \cite[Section~3.17.12]{LangevinLean2026ULASC}; \cite[Section~3.17.2]{LangevinLean2026ULASC}; \cite[Section~3.17.4]{LangevinLean2026ULASC}; \cite[Section~3.22.17]{LangevinLean2026ULASC}; \cite[Section~3.22.18]{LangevinLean2026ULASC}; \cite[Section~3.22.19]{LangevinLean2026ULASC}; \cite[Section~3.22.2]{LangevinLean2026ULASC}; \cite[Section~3.22.20]{LangevinLean2026ULASC}; \cite[Section~3.22.21]{LangevinLean2026ULASC}; \cite[Section~3.22.22]{LangevinLean2026ULASC}; \cite[Section~3.22.23]{LangevinLean2026ULASC}; \cite[Section~3.22.3]{LangevinLean2026ULASC}; \cite[Section~3.22.7]{LangevinLean2026ULASC}; \cite[Section~3.22.8]{LangevinLean2026ULASC}; \cite[Section~3.23.1]{LangevinLean2026ULASC}; \cite[Section~3.23.2]{LangevinLean2026ULASC}; \cite[Section~3.23.6]{LangevinLean2026ULASC}; \cite[Section~3.25.1]{LangevinLean2026ULASC}; \cite[Section~3.25.4]{LangevinLean2026ULASC}; \cite[Section~3.25.6]{LangevinLean2026ULASC}; \cite[Section~3.26.3]{LangevinLean2026ULASC}; \cite[Section~3.28.2]{LangevinLean2026ULASC}; \cite[Section~3.29.7]{LangevinLean2026ULASC}; \cite[Section~3.29.8]{LangevinLean2026ULASC}; \cite[Section~3.32.12]{LangevinLean2026ULASC}; \cite[Section~3.32.13]{LangevinLean2026ULASC}; \cite[Section~3.32.15]{LangevinLean2026ULASC}; \cite[Section~3.32.16]{LangevinLean2026ULASC}; \cite[Section~3.32.17]{LangevinLean2026ULASC}; \cite[Section~3.32.19]{LangevinLean2026ULASC}; \cite[Section~3.32.4]{LangevinLean2026ULASC}; \cite[Section~3.32.5]{LangevinLean2026ULASC}; \cite[Section~3.36.7]{LangevinLean2026ULASC}; \cite[Section~3.42.6]{LangevinLean2026ULASC}; \cite[Section~3.42.7]{LangevinLean2026ULASC}; \cite[Section~3.8.11]{LangevinLean2026ULASC}; \cite[Section~3.8.16]{LangevinLean2026ULASC}; \cite[Section~3.8.17]{LangevinLean2026ULASC}; \cite[Section~3.8.19]{LangevinLean2026ULASC}; \cite[Section~3.8.20]{LangevinLean2026ULASC}; \cite[Section~3.8.21]{LangevinLean2026ULASC}; \cite[Section~3.8.22]{LangevinLean2026ULASC}.
\section{\texorpdfstring{The exact third-order LMC theorem for the actual Gaussian algorithm}{The exact third-order LMC theorem for the actual Gaussian algorithm}}
\label{proof-module-AverageLMC-LaplacianTheorem-lean}


Laplacian theorem (source th:3) follows from the actual one-step coupling and the solved quadratic recurrence, with precisely three continuous derivatives and the paper’s Laplacian-gradient condition. No Hessian Lipschitz or fourth-derivative condition is added.



The Laplacian theorem (source th:3, rendered Theorem 2) assumes \(f\in C^3\), \(m\)-strong convexity, \(L\)-Lipschitz gradient, \(0<h<1/L\), and exactly the trace-type third-order condition


\[\|\nabla\Delta f(x)\|^2\le M_\Delta^2p\qquad(x\in\mathbb R^p).\]


The actual constant-step Gaussian LMC marginals then satisfy, also for arbitrary initial laws in extended transport,


\[W_2(\nu_k,\pi)\le(1-mh)^kW_2(\nu_0,\pi)+\bigl(2L\sqrt{M_{\rm av}}+M_\Delta\bigr)\frac{h\sqrt p}{m}.\]\subsection{\texorpdfstring{Pass the quadratic-plus-linear estimate to the transport infimum}{Pass the quadratic-plus-linear estimate to the transport infimum}}
\label{report-step-AverageLMC-LaplacianTheorem-lean-1}


For each positive \(\varepsilon\), the definition of Wasserstein distance supplies a finite initial coupling with cost at most \(W_{2}+\varepsilon\). Nonnegative quadratic and linear coefficients allow both powers of that cost to be replaced by \(W_{2}+\varepsilon\). The resulting polynomial upper bound is continuous as \(\varepsilon\) decreases to zero. Passing to this limit proves the exact quadratic-plus-linear estimate in the actual Wasserstein distance without assuming that any optimal coupling exists.

\begin{leanbox}{Lean: \texttt{quadratic\_\allowbreak{}linear\_\allowbreak{}bound\_\allowbreak{}of\_\allowbreak{}finite\_\allowbreak{}couplings}}
\leansource{\texttt{quadratic\_\allowbreak{}linear\_\allowbreak{}bound\_\allowbreak{}of\_\allowbreak{}finite\_\allowbreak{}couplings} \quad AverageLMC/\allowbreak{}LaplacianTheorem.lean, lines 12–41}
\begin{leancode}
\leanline{0}{theorem\ quadratic\_linear\_bound\_of\_finite\_couplings\ \{E\ :\ Type*\}}
\leanline{4}{[NormedAddCommGroup\ E]\ [MeasurableSpace\ E]\ (μ\ π\ :\ Measure\ E)\ (r\ a\ b\ c\ :\ ℝ)}
\leanline{4}{(ha\ :\ 0\ ≤\ a)\ (hb\ :\ 0\ ≤\ b)\ (hfinite\ :\ wasserstein2\ μ\ π\ ≠\ ⊤)}
\leanline{4}{(hfamily\ :\ ∀\ γ\ :\ Coupling\ μ\ π,\ couplingCost\ γ.joint\ ≠\ ⊤\ →}
\leanline{6}{r\ ≤\ a\ *\ (couplingCost\ γ.joint).toReal\ \textasciicircum{}\ 2\ +\ b\ *\ (couplingCost\ γ.joint).toReal\ +\ c)\ :}
\leanline{4}{r\ ≤\ a\ *\ wasserstein2Real\ μ\ π\ \textasciicircum{}\ 2\ +\ b\ *\ wasserstein2Real\ μ\ π\ +\ c}
\end{leancode}
\leanlabel{How the Lean proof follows the argument}
\texttt{quadratic\_\allowbreak{}linear\_\allowbreak{}bound\_\allowbreak{}of\_\allowbreak{}finite\_\allowbreak{}couplings} first obtains \texttt{hw}, nonnegativity of the real transport distance. For each positive \texttt{\ensuremath{\varepsilon}}, \texttt{heps} selects a finite-cost plan whose cost lies below that distance plus \texttt{\ensuremath{\varepsilon}} and applies the assumed planwise polynomial inequality. Nonnegativity of the quadratic and linear coefficients permits this enlargement. \texttt{hlim} proves convergence of that polynomial as \texttt{\ensuremath{\varepsilon}} tends to zero from the right, and \texttt{le\_\allowbreak{}of\_\allowbreak{}tendsto\_\allowbreak{}of\_\allowbreak{}tendsto} passes the bound to the limit.
\end{leanbox}


Earlier results used here: \cite[Section~3.17.1]{LangevinLean2026ULASC}; \cite[Section~3.17.2]{LangevinLean2026ULASC}; \cite[Section~3.17.3]{LangevinLean2026ULASC}; \cite[Section~3.17.4]{LangevinLean2026ULASC}; \cite[Section~3.1.1]{LangevinLean2026IPLA}.
\subsection{\texorpdfstring{Instantiate the error recurrence with actual marginal laws}{Instantiate the error recurrence with actual marginal laws}}
\label{report-step-AverageLMC-LaplacianTheorem-lean-2}


Finite initial transport propagates along the actual constant-step Gaussian LMC laws. At every step, use the canonical coupling estimate and the quadratic-plus-linear infimum theorem to obtain the scalar recurrence for their actual Wasserstein errors. Strong convexity and coordinate smoothness give \(m\le M_{\mathrm{av}}\le L\). Split \(\sqrt{M_{\mathrm{av}}p}\) into \(\sqrt{M_{\mathrm{av}}}\sqrt p\) and apply the proved recurrence with \(A=\sqrt{M_{\mathrm{av}}}\), \(\delta =M_\Delta\) and \(S=\sqrt p\). Its invariant floor is exactly \((2L\sqrt{M_{\mathrm{av}}}+M_\Delta )h\sqrt p/m\), while its transient is \((1-mh)^{k}\) times the true initial distance. The constant-schedule law identity converts the recursive marginals to iteration of the actual Gaussian transition.

\begin{leanbox}{Lean: \texttt{lmc\_\allowbreak{}laplacian}}
\leansource{\texttt{lmc\_\allowbreak{}laplacian} \quad AverageLMC/\allowbreak{}LaplacianTheorem.lean, lines 45–84}
\begin{leancode}
\leanline{0}{theorem\ lmc\_laplacian\ \{p\ :\ ℕ\}\ (hp\ :\ 0\ <\ p)\ (f\ :\ State\ p\ →\ ℝ)}
\leanline{4}{(m\ h\ :\ ℝ)\ (L\ :\ ℝ≥0)\ (hm\ :\ 0\ <\ m)\ (hh\ :\ 0\ <\ h)}
\leanline{4}{(hsc\ :\ StrongConvexOn\ univ\ m\ f)\ (hLip\ :\ LipschitzWith\ L\ (gradient\ f))}
\leanline{4}{(hsmall\ :\ h\ *\ (L\ :\ ℝ)\ <\ 1)\ (Mdelta\ :\ ℝ)\ (hMdelta\ :\ 0\ ≤\ Mdelta)}
\leanline{4}{(hδ\ :\ LaplacianSmooth\ f\ Mdelta)\ (μ\ :\ Measure\ (State\ p))\ [IsProbabilityMeasure\ μ]}
\leanline{4}{(hfinite\ :\ wasserstein2\ μ\ (gibbsMeasure\ f)\ ≠\ ⊤)\ (k\ :\ ℕ)\ :}
\leanline{4}{wasserstein2Real\ ((ulaLawStep\ f\ h)\textasciicircum{}[k]\ μ)\ (gibbsMeasure\ f)\ ≤}
\leanline{6}{(1\ -\ m\ *\ h)\ \textasciicircum{}\ k\ *\ wasserstein2Real\ μ\ (gibbsMeasure\ f)\ +}
\leanline{6}{(2\ *\ (L\ :\ ℝ)\ *\ Real.sqrt\ (averageSmoothness\ f)\ +\ Mdelta)\ *\ h\ *\ Real.sqrt\ p\ /\ m}
\end{leancode}
\leanlabel{How the Lean proof follows the argument}
\texttt{lp\_\allowbreak{}stationary\_\allowbreak{}gradient\_\allowbreak{}pairing\_\allowbreak{}bound} bounds the cross term with the actual stationary diffusion and the stated gradient-of-Laplacian assumption. \texttt{lp\_\allowbreak{}laplacian\_\allowbreak{}one\_\allowbreak{}step\_\allowbreak{}bound} combines that pairing with the local integral remainder and gradient-step contraction. \texttt{laplacian\_\allowbreak{}recurrence\_\allowbreak{}bound\_\allowbreak{}of\_\allowbreak{}global\_\allowbreak{}defect} iterates the resulting quadratic-plus-linear scalar estimate. The final theorem \texttt{lmc\_\allowbreak{}laplacian} obtains the required finite-distance recurrence for the actual Gaussian laws and simplifies its constant using the exact Gibbs score moment bound.
\end{leanbox}


Earlier results used here: \cite[Section~3.17.3]{LangevinLean2026ULASC}; \cite[Section~3.17.4]{LangevinLean2026ULASC}; \cite[Section~3.26.3]{LangevinLean2026ULASC}; \cite[Section~3.36.7]{LangevinLean2026ULASC}.
\subsection{\texorpdfstring{Include arbitrary initial probability laws exactly}{Include arbitrary initial probability laws exactly}}
\label{report-step-AverageLMC-LaplacianTheorem-lean-3}


The manuscript uses the strict step range \(hL<1\). Together with \(m\le L\) this makes \(1-mh\) strictly positive, so its \(k\)th power is positive for every \(k\). An infinite initial extended Wasserstein distance therefore gives an infinite right side, proving that case immediately. For finite initial distance, the actual-law theorem and derived finite output transport convert back to the extended nonnegative inequality. Thus the final statement imposes no independent initial second-moment assumption.

\begin{leanbox}{Lean: \texttt{lmc\_\allowbreak{}laplacian\_\allowbreak{}extended}}
\leansource{\texttt{lmc\_\allowbreak{}laplacian\_\allowbreak{}extended} \quad AverageLMC/\allowbreak{}LaplacianTheorem.lean, lines 88–111}
\begin{leancode}
\leanline{0}{theorem\ lmc\_laplacian\_extended\ \{p\ :\ ℕ\}\ (hp\ :\ 0\ <\ p)\ (f\ :\ State\ p\ →\ ℝ)}
\leanline{4}{(m\ h\ :\ ℝ)\ (L\ :\ ℝ≥0)\ (hm\ :\ 0\ <\ m)\ (hh\ :\ 0\ <\ h)}
\leanline{4}{(hsc\ :\ StrongConvexOn\ univ\ m\ f)\ (hLip\ :\ LipschitzWith\ L\ (gradient\ f))}
\leanline{4}{(hsmall\ :\ h\ *\ (L\ :\ ℝ)\ <\ 1)\ (Mdelta\ :\ ℝ)\ (hMdelta\ :\ 0\ ≤\ Mdelta)}
\leanline{4}{(hδ\ :\ LaplacianSmooth\ f\ Mdelta)\ (μ\ :\ Measure\ (State\ p))\ [IsProbabilityMeasure\ μ]\ (k\ :\ ℕ)\ :}
\leanline{4}{wasserstein2\ ((ulaLawStep\ f\ h)\textasciicircum{}[k]\ μ)\ (gibbsMeasure\ f)\ ≤}
\leanline{6}{ENNReal.ofReal\ ((1\ -\ m\ *\ h)\ \textasciicircum{}\ k)\ *\ wasserstein2\ μ\ (gibbsMeasure\ f)\ +}
\leanline{6}{ENNReal.ofReal\ ((2\ *\ (L\ :\ ℝ)\ *\ Real.sqrt\ (averageSmoothness\ f)\ +\ Mdelta)\ *\ h\ *\ Real.sqrt\ p\ /\ m)}
\end{leancode}
\leanlabel{How the Lean proof follows the argument}
\texttt{hq} proves strict positivity of the contraction base from \texttt{hsmall} and \(m\le L\). The proof splits \texttt{hi}, the infinite-initial-distance case. That branch simplifies because multiplication by the positive geometric factor preserves infinity. In the finite branch \texttt{ho} derives finite output distance, \texttt{hb} lifts \texttt{lmc\_\allowbreak{}laplacian} through \texttt{ENNReal.\allowbreak{}of\allowbreak{}Real}, and \texttt{hconst} and \texttt{hprod} justify splitting the nonnegative sum and product. The final rewrites recover both actual extended distances from their real values.
\end{leanbox}


Earlier results used here: \cite[Section~3.17.3]{LangevinLean2026ULASC}; \cite[Section~3.17.4]{LangevinLean2026ULASC}; \cite[Section~3.26.3]{LangevinLean2026ULASC}; \cite[Section~3.36.7]{LangevinLean2026ULASC}.
\section{\texorpdfstring{The genuine finite-sum control-variate estimator}{The genuine finite-sum control-variate estimator}}
\label{proof-module-AverageLMC-SGLDEstimator-lean}


This module constructs the finite-sum potential, the actual uniform component-index probability, and the fixed-point control-variate estimator in the SGLD update. It derives unbiasedness, the exact centered variance identity, the sharp RMS component bound, and all product-space second moments. The anchor is justified by existence and uniqueness of a true minimizer, including its zero gradient.

\subsection{\texorpdfstring{Construct the potential and the uniform index probability}{Construct the potential and the uniform index probability}}
\label{report-step-AverageLMC-SGLDEstimator-lean-1}


The total potential is the finite sum of the component functions. Differentiation commutes with this finite sum, and the gradient identification follows by the inner-product duality map. The uniform index law is counting measure multiplied by 1/\allowbreak{}\(n\). For \(n\) positive its mass is one. It is a finite measure even at \(n=0\), where the index space is empty. Its Bochner integral is exactly the finite vector sum divided by \(n\); this formula is used for all following index moments.

\begin{leanbox}{Lean: \texttt{componentPotential}, \texttt{uniformIndexLaw}, \texttt{uniformIndexLaw\_\allowbreak{}probability} and 3 more}
\leansource{\texttt{componentPotential} \quad AverageLMC/\allowbreak{}SGLDEstimator.lean, lines 15–16}
\begin{leancode}
\leanline{0}{def\ componentPotential\ \{p\ n\ :\ ℕ\}\ (f\ :\ Fin\ n\ →\ State\ p\ →\ ℝ)\ (x\ :\ State\ p)\ :\ ℝ\ :=}
\leanline{2}{∑\ i,\ f\ i\ x}
\end{leancode}
\leansource{\texttt{uniformIndexLaw} \quad AverageLMC/\allowbreak{}SGLDEstimator.lean, lines 19–20}
\begin{leancode}
\leanline{0}{def\ uniformIndexLaw\ (n\ :\ ℕ)\ :\ Measure\ (Fin\ n)\ :=}
\leanline{2}{(n\ :\ ℝ≥0∞)⁻¹\ •\ Measure.count}
\end{leancode}
\leansource{\texttt{uniformIndexLaw\_\allowbreak{}probability} \quad AverageLMC/\allowbreak{}SGLDEstimator.lean, lines 22–25}
\begin{leancode}
\leanline{0}{theorem\ uniformIndexLaw\_probability\ (n\ :\ ℕ)\ (hn\ :\ 0\ <\ n)\ :}
\leanline{4}{IsProbabilityMeasure\ (uniformIndexLaw\ n)}
\end{leancode}
\leansource{\texttt{uniformIndexLaw\_\allowbreak{}isFiniteMeasure} \quad AverageLMC/\allowbreak{}SGLDEstimator.lean, lines 27–34}
\begin{leancode}
\leanline{0}{instance\ uniformIndexLaw\_isFiniteMeasure\ (n\ :\ ℕ)\ :\ IsFiniteMeasure\ (uniformIndexLaw\ n)\ :=\ by}
\leanline{2}{by\_cases\ hn\ :\ 0\ <\ n}
\leanline{2}{·\ have\ :=\ uniformIndexLaw\_probability\ n\ hn}
\leanline{4}{infer\_instance}
\leanline{2}{·\ have\ hn0\ :\ n\ =\ 0\ :=\ Nat.eq\_zero\_of\_not\_pos\ hn}
\leanline{4}{subst\ n}
\leanline{4}{constructor}
\leanline{4}{simp\ [uniformIndexLaw]}
\end{leancode}
\leansource{\texttt{integral\_\allowbreak{}uniformIndexLaw} \quad AverageLMC/\allowbreak{}SGLDEstimator.lean, lines 36–40}
\begin{leancode}
\leanline{0}{theorem\ integral\_uniformIndexLaw\ \{E\ :\ Type*\}\ [NormedAddCommGroup\ E]\ [NormedSpace\ ℝ\ E]}
\leanline{4}{[CompleteSpace\ E]}
\leanline{4}{(n\ :\ ℕ)\ (g\ :\ Fin\ n\ →\ E)\ :}
\leanline{4}{(∫\ i,\ g\ i\ ∂uniformIndexLaw\ n)\ =\ (n\ :\ ℝ)⁻¹\ •\ ∑\ i,\ g\ i}
\end{leancode}
\leansource{\texttt{gradient\_\allowbreak{}componentPotential} \quad AverageLMC/\allowbreak{}SGLDEstimator.lean, lines 42–47}
\begin{leancode}
\leanline{0}{theorem\ gradient\_componentPotential\ \{p\ n\ :\ ℕ\}\ (f\ :\ Fin\ n\ →\ State\ p\ →\ ℝ)}
\leanline{4}{(hf\ :\ ∀\ i,\ Differentiable\ ℝ\ (f\ i))\ (x\ :\ State\ p)\ :}
\leanline{4}{gradient\ (componentPotential\ f)\ x\ =\ ∑\ i,\ gradient\ (f\ i)\ x}
\end{leancode}
\leanlabel{How the Lean proof follows the argument}
\texttt{component\allowbreak{}Potential} names the mathematical object defined above. Its defining expression appears in the adjacent statement; subsequent identities may unfold this definition directly.




\texttt{uniform\allowbreak{}Index\allowbreak{}Law} names the mathematical object defined above. Its defining expression appears in the adjacent statement; subsequent identities may unfold this definition directly.




\texttt{uniform\allowbreak{}Index\allowbreak{}Law\_\allowbreak{}probability} gives the formal counterpart of this step. Its remaining equalities are discharged by the simplifications displayed in the proof.




\texttt{uniform\allowbreak{}Index\allowbreak{}Law\_\allowbreak{}is\allowbreak{}Finite\allowbreak{}Measure} names the mathematical object defined above. Its defining expression appears in the adjacent statement; subsequent identities may unfold this definition directly.




\texttt{integral\_\allowbreak{}uniform\allowbreak{}Index\allowbreak{}Law} gives the formal counterpart of this step. Its remaining equalities are discharged by the simplifications displayed in the proof.




\texttt{gradient\_\allowbreak{}component\allowbreak{}Potential} gives the formal counterpart of this step. Named proof applications include \texttt{fderiv\_\allowbreak{}fun\_\allowbreak{}sum}; the adjacent proof shows their hypotheses and order of use.
\end{leanbox}
\subsection{\texorpdfstring{Use the actual fixed-point estimator and prove unbiasedness}{Use the actual fixed-point estimator and prove unbiasedness}}
\label{report-step-AverageLMC-SGLDEstimator-lean-2}


For an anchor \(x_\star\), define the estimator \(n(\nabla f_i(x)-\nabla f_i(x_\star))\). It vanishes at the anchor for every index. Averaging over the uniform index cancels \(n\) and sums the component gradients. If the true total gradient at the anchor is zero, this average is exactly \(\nabla f(x)\). Thus the paper estimator is unbiased for the gradient of the same potential defining its Gibbs target.

\begin{leanbox}{Lean: \texttt{controlVariateEstimator}, \texttt{controlVariateEstimator\_\allowbreak{}at\_\allowbreak{}anchor}, \texttt{controlVariateEstimator\_\allowbreak{}unbiased}}
\leansource{\texttt{controlVariateEstimator} \quad AverageLMC/\allowbreak{}SGLDEstimator.lean, lines 50–52}
\begin{leancode}
\leanline{0}{def\ controlVariateEstimator\ \{p\ n\ :\ ℕ\}\ (f\ :\ Fin\ n\ →\ State\ p\ →\ ℝ)}
\leanline{4}{(xstar\ x\ :\ State\ p)\ (i\ :\ Fin\ n)\ :\ State\ p\ :=}
\leanline{2}{(n\ :\ ℝ)\ •\ (gradient\ (f\ i)\ x\ -\ gradient\ (f\ i)\ xstar)}
\end{leancode}
\leansource{\texttt{controlVariateEstimator\_\allowbreak{}at\_\allowbreak{}anchor} \quad AverageLMC/\allowbreak{}SGLDEstimator.lean, lines 54–56}
\begin{leancode}
\leanline{0}{theorem\ controlVariateEstimator\_at\_anchor\ \{p\ n\ :\ ℕ\}\ (f\ :\ Fin\ n\ →\ State\ p\ →\ ℝ)}
\leanline{4}{(xstar\ :\ State\ p)\ (i\ :\ Fin\ n)\ :\ controlVariateEstimator\ f\ xstar\ xstar\ i\ =\ 0}
\end{leancode}
\leansource{\texttt{controlVariateEstimator\_\allowbreak{}unbiased} \quad AverageLMC/\allowbreak{}SGLDEstimator.lean, lines 58–68}
\begin{leancode}
\leanline{0}{theorem\ controlVariateEstimator\_unbiased\ \{p\ n\ :\ ℕ\}\ (hn\ :\ 0\ <\ n)}
\leanline{4}{(f\ :\ Fin\ n\ →\ State\ p\ →\ ℝ)\ (hf\ :\ ∀\ i,\ Differentiable\ ℝ\ (f\ i))}
\leanline{4}{(xstar\ :\ State\ p)\ (hstar\ :\ gradient\ (componentPotential\ f)\ xstar\ =\ 0)}
\leanline{4}{(x\ :\ State\ p)\ :}
\leanline{4}{(∫\ i,\ controlVariateEstimator\ f\ xstar\ x\ i\ ∂uniformIndexLaw\ n)\ =}
\leanline{6}{gradient\ (componentPotential\ f)\ x}
\end{leancode}
\leanlabel{How the Lean proof follows the argument}
\texttt{control\allowbreak{}Variate\allowbreak{}Estimator} names the mathematical object defined above. Its defining expression appears in the adjacent statement; subsequent identities may unfold this definition directly.




\texttt{control\allowbreak{}Variate\allowbreak{}Estimator\_\allowbreak{}at\_\allowbreak{}anchor} gives the formal counterpart of this step. Its remaining equalities are discharged by the simplifications displayed in the proof.




\texttt{control\allowbreak{}Variate\allowbreak{}Estimator\_\allowbreak{}unbiased} gives the formal counterpart of this step.
\end{leanbox}
\subsection{\texorpdfstring{Control the full gradient by the total component smoothness}{Control the full gradient by the total component smoothness}}
\label{report-step-AverageLMC-SGLDEstimator-lean-3}


Squaring the defining square root gives \(M_{\rm rms}\) squared equal to the average squared exact component-gradient seminorms. The triangle inequality bounds the gradient Lipschitz constant of the sum by the sum of component constants. Finite Cauchy-Schwarz bounds the sum of exact nonnegative seminorms by \(n\) \(M_{\rm rms}\). Replacing each supplied component bound by its proved exact seminorm gives a Lipschitz gradient bound for the full potential with exactly \(n\) \(M_{\rm rms}\). Comparing a strong-convexity lower model with the coordinate upper model then proves \(m\) at most \(n\) \(M_{\rm rms}\). These comparisons justify the dimensionless step-size conditions in the SGLD theorem.

\begin{leanbox}{Lean: \texttt{rmsComponentSmoothness\_\allowbreak{}sq}, \texttt{componentPotential\_\allowbreak{}gradient\_\allowbreak{}lipschitz}, \texttt{sum\_\allowbreak{}globalSmoothness\_\allowbreak{}le\_\allowbreak{}n\_\allowbreak{}mul\_\allowbreak{}rms} and 2 more}
\leansource{\texttt{rmsComponentSmoothness\_\allowbreak{}sq} \quad AverageLMC/\allowbreak{}SGLDEstimator.lean, lines 71–74}
\begin{leancode}
\leanline{0}{theorem\ rmsComponentSmoothness\_sq\ \{p\ n\ :\ ℕ\}\ (f\ :\ Fin\ n\ →\ State\ p\ →\ ℝ)\ :}
\leanline{4}{rmsComponentSmoothness\ f\ \textasciicircum{}\ 2\ =\ (∑\ i,\ globalSmoothness\ (f\ i)\ \textasciicircum{}\ 2)\ /\ (n\ :\ ℝ)}
\end{leancode}
\leansource{\texttt{componentPotential\_\allowbreak{}gradient\_\allowbreak{}lipschitz} \quad AverageLMC/\allowbreak{}SGLDEstimator.lean, lines 76–87}
\begin{leancode}
\leanline{0}{theorem\ componentPotential\_gradient\_lipschitz\ \{p\ n\ :\ ℕ\}\ (f\ :\ Fin\ n\ →\ State\ p\ →\ ℝ)}
\leanline{4}{(hf\ :\ ∀\ i,\ Differentiable\ ℝ\ (f\ i))\ (L\ :\ Fin\ n\ →\ ℝ≥0)}
\leanline{4}{(hLip\ :\ ∀\ i,\ LipschitzWith\ (L\ i)\ (gradient\ (f\ i)))\ :}
\leanline{4}{LipschitzWith\ (∑\ i,\ L\ i)\ (gradient\ (componentPotential\ f))}
\end{leancode}
\leansource{\texttt{sum\_\allowbreak{}globalSmoothness\_\allowbreak{}le\_\allowbreak{}n\_\allowbreak{}mul\_\allowbreak{}rms} \quad AverageLMC/\allowbreak{}SGLDEstimator.lean, lines 89–105}
\begin{leancode}
\leanline{0}{theorem\ sum\_globalSmoothness\_le\_n\_mul\_rms\ \{p\ n\ :\ ℕ\}\ (hp\ :\ 0\ <\ p)\ (hn\ :\ 0\ <\ n)}
\leanline{4}{(f\ :\ Fin\ n\ →\ State\ p\ →\ ℝ)\ (L\ :\ Fin\ n\ →\ ℝ≥0)}
\leanline{4}{(hLip\ :\ ∀\ i,\ LipschitzWith\ (L\ i)\ (gradient\ (f\ i)))\ :}
\leanline{4}{(∑\ i,\ globalSmoothness\ (f\ i))\ ≤\ (n\ :\ ℝ)\ *\ rmsComponentSmoothness\ f}
\end{leancode}
\leansource{\texttt{componentPotential\_\allowbreak{}gradient\_\allowbreak{}lipschitz\_\allowbreak{}rms} \quad AverageLMC/\allowbreak{}SGLDEstimator.lean, lines 108–122}
\begin{leancode}
\leanline{0}{theorem\ componentPotential\_gradient\_lipschitz\_rms\ \{p\ n\ :\ ℕ\}\ (hp\ :\ 0\ <\ p)\ (hn\ :\ 0\ <\ n)}
\leanline{4}{(f\ :\ Fin\ n\ →\ State\ p\ →\ ℝ)\ (hf\ :\ ∀\ i,\ Differentiable\ ℝ\ (f\ i))}
\leanline{4}{(L\ :\ Fin\ n\ →\ ℝ≥0)\ (hLip\ :\ ∀\ i,\ LipschitzWith\ (L\ i)\ (gradient\ (f\ i)))\ :}
\leanline{4}{LipschitzWith\ ⟨(n\ :\ ℝ)\ *\ rmsComponentSmoothness\ f,}
\leanline{6}{mul\_nonneg\ (Nat.cast\_nonneg\ \_)\ (Real.sqrt\_nonneg\ \_)⟩}
\leanline{6}{(gradient\ (componentPotential\ f))}
\end{leancode}
\leansource{\texttt{strongConvex\_\allowbreak{}le\_\allowbreak{}n\_\allowbreak{}mul\_\allowbreak{}rms} \quad AverageLMC/\allowbreak{}SGLDEstimator.lean, lines 124–136}
\begin{leancode}
\leanline{0}{theorem\ strongConvex\_le\_n\_mul\_rms\ \{p\ n\ :\ ℕ\}\ (hp\ :\ 0\ <\ p)\ (hn\ :\ 0\ <\ n)}
\leanline{4}{(f\ :\ Fin\ n\ →\ State\ p\ →\ ℝ)\ (hf\ :\ ∀\ i,\ Differentiable\ ℝ\ (f\ i))}
\leanline{4}{(L\ :\ Fin\ n\ →\ ℝ≥0)\ (hLip\ :\ ∀\ i,\ LipschitzWith\ (L\ i)\ (gradient\ (f\ i)))}
\leanline{4}{(m\ :\ ℝ)\ (hsc\ :\ StrongConvexOn\ Set.univ\ m\ (componentPotential\ f))\ :}
\leanline{4}{m\ ≤\ (n\ :\ ℝ)\ *\ rmsComponentSmoothness\ f}
\end{leancode}
\leanlabel{How the Lean proof follows the argument}
\texttt{rms\allowbreak{}Component\allowbreak{}Smoothness\_\allowbreak{}sq} gives the formal counterpart of this step. Named proof applications include \texttt{Real.\allowbreak{}sq\_\allowbreak{}sqrt}, \texttt{div\_\allowbreak{}nonneg}; the adjacent proof shows their hypotheses and order of use.




\texttt{component\allowbreak{}Potential\_\allowbreak{}gradient\_\allowbreak{}lipschitz} gives the formal counterpart of this step. Named proof applications include \texttt{Lipschitz\allowbreak{}With.\allowbreak{}of\_\allowbreak{}dist\_\allowbreak{}le\_\allowbreak{}mul}; the adjacent proof shows their hypotheses and order of use.




\texttt{sum\_\allowbreak{}global\allowbreak{}Smoothness\_\allowbreak{}le\_\allowbreak{}n\_\allowbreak{}mul\_\allowbreak{}rms} gives the formal counterpart of this step. The main intermediate facts are \texttt{hc}, \texttt{heq} in that order. Named proof applications include \texttt{hc}; the adjacent proof shows their hypotheses and order of use.




\texttt{component\allowbreak{}Potential\_\allowbreak{}gradient\_\allowbreak{}lipschitz\_\allowbreak{}rms} gives the formal counterpart of this step. The main intermediate facts are \texttt{C}, \texttt{hc} in that order. Named proof applications include \texttt{NNReal.\allowbreak{}coe\_\allowbreak{}le\_\allowbreak{}coe.\allowbreak{}mp}, \texttt{sum\_\allowbreak{}global\allowbreak{}Smoothness\_\allowbreak{}le\_\allowbreak{}n\_\allowbreak{}mul\_\allowbreak{}rms}; the adjacent proof shows their hypotheses and order of use.




\texttt{strong\allowbreak{}Convex\_\allowbreak{}le\_\allowbreak{}n\_\allowbreak{}mul\_\allowbreak{}rms} gives the formal counterpart of this step. The main intermediate facts are \texttt{hd}, \texttt{C}, \texttt{hc} in that order.
\end{leanbox}
\subsection{\texorpdfstring{Compute and bound the estimator second moment}{Compute and bound the estimator second moment}}
\label{report-step-AverageLMC-SGLDEstimator-lean-4}


The uniform average of the squared estimator norm equals \(n\) times the sum of squared component-gradient differences. Each difference is bounded by the component exact seminorm times the distance to the anchor. Square these nonnegative inequalities, sum, and substitute the exact RMS definition. This gives

\[\mathbb E_I\|G_I(x)\|^2\le n^2M_{\rm rms}^2\|x-x_*\|^2.\]

The integral is finite because the index space and measure are finite.

\begin{leanbox}{Lean: \texttt{controlVariateEstimator\_\allowbreak{}second\_\allowbreak{}moment}, \texttt{controlVariateEstimator\_\allowbreak{}second\_\allowbreak{}moment\_\allowbreak{}le}}
\leansource{\texttt{controlVariateEstimator\_\allowbreak{}second\_\allowbreak{}moment} \quad AverageLMC/\allowbreak{}SGLDEstimator.lean, lines 139–148}
\begin{leancode}
\leanline{0}{theorem\ controlVariateEstimator\_second\_moment\ \{p\ n\ :\ ℕ\}\ (hn\ :\ 0\ <\ n)}
\leanline{4}{(f\ :\ Fin\ n\ →\ State\ p\ →\ ℝ)\ (xstar\ x\ :\ State\ p)\ :}
\leanline{4}{(∫\ i,\ ‖controlVariateEstimator\ f\ xstar\ x\ i‖\ \textasciicircum{}\ 2\ ∂uniformIndexLaw\ n)\ =}
\leanline{6}{(n\ :\ ℝ)\ *\ ∑\ i,\ ‖gradient\ (f\ i)\ x\ -\ gradient\ (f\ i)\ xstar‖\ \textasciicircum{}\ 2}
\end{leancode}
\leansource{\texttt{controlVariateEstimator\_\allowbreak{}second\_\allowbreak{}moment\_\allowbreak{}le} \quad AverageLMC/\allowbreak{}SGLDEstimator.lean, lines 152–167}
\begin{leancode}
\leanline{0}{theorem\ controlVariateEstimator\_second\_moment\_le\ \{p\ n\ :\ ℕ\}\ (hp\ :\ 0\ <\ p)\ (hn\ :\ 0\ <\ n)}
\leanline{4}{(f\ :\ Fin\ n\ →\ State\ p\ →\ ℝ)\ (L\ :\ Fin\ n\ →\ ℝ≥0)}
\leanline{4}{(hLip\ :\ ∀\ i,\ LipschitzWith\ (L\ i)\ (gradient\ (f\ i)))\ (xstar\ x\ :\ State\ p)\ :}
\leanline{4}{(∫\ i,\ ‖controlVariateEstimator\ f\ xstar\ x\ i‖\ \textasciicircum{}\ 2\ ∂uniformIndexLaw\ n)\ ≤}
\leanline{6}{(n\ :\ ℝ)\ \textasciicircum{}\ 2\ *\ rmsComponentSmoothness\ f\ \textasciicircum{}\ 2\ *\ ‖x\ -\ xstar‖\ \textasciicircum{}\ 2}
\end{leancode}
\leanlabel{How the Lean proof follows the argument}
\texttt{control\allowbreak{}Variate\allowbreak{}Estimator\_\allowbreak{}second\_\allowbreak{}moment} gives the formal counterpart of this step. Its remaining equalities are discharged by the simplifications displayed in the proof.




\texttt{control\allowbreak{}Variate\allowbreak{}Estimator\_\allowbreak{}second\_\allowbreak{}moment\_\allowbreak{}le} gives the formal counterpart of this step. The main intermediate facts are \texttt{hi}, \texttt{hs} in that order. Its remaining equalities are discharged by the simplifications displayed in the proof.
\end{leanbox}
\subsection{\texorpdfstring{Center the error and retain the negative gradient energy}{Center the error and retain the negative gradient energy}}
\label{report-step-AverageLMC-SGLDEstimator-lean-5}


Define the estimator error as \(G_i(x)-\nabla f(x)\). Unbiasedness proves that its uniform-index expectation is zero. Expanding its squared norm, the cross expectation is the inner product of \(\nabla f(x)\) with the estimator expectation, hence the squared full-gradient norm. The exact variance identity is therefore

\[\mathbb E_I\|G_I(x)-\nabla f(x)\|^2=\mathbb E_I\|G_I(x)\|^2-\|\nabla f(x)\|^2.\]

Dropping the nonnegative final score energy gives the same \(n^2\) RMS distance bound for the centered error. Keeping it is essential for the larger-step one-step SGLD estimates.

\begin{leanbox}{Lean: \texttt{controlVariateError}, \texttt{controlVariateError\_\allowbreak{}centered}, \texttt{controlVariateError\_\allowbreak{}second\_\allowbreak{}moment} and 1 more}
\leansource{\texttt{controlVariateError} \quad AverageLMC/\allowbreak{}SGLDEstimator.lean, lines 170–172}
\begin{leancode}
\leanline{0}{def\ controlVariateError\ \{p\ n\ :\ ℕ\}\ (f\ :\ Fin\ n\ →\ State\ p\ →\ ℝ)}
\leanline{4}{(xstar\ x\ :\ State\ p)\ (i\ :\ Fin\ n)\ :\ State\ p\ :=}
\leanline{2}{controlVariateEstimator\ f\ xstar\ x\ i\ -\ gradient\ (componentPotential\ f)\ x}
\end{leancode}
\leansource{\texttt{controlVariateError\_\allowbreak{}centered} \quad AverageLMC/\allowbreak{}SGLDEstimator.lean, lines 174–183}
\begin{leancode}
\leanline{0}{theorem\ controlVariateError\_centered\ \{p\ n\ :\ ℕ\}\ (hn\ :\ 0\ <\ n)}
\leanline{4}{(f\ :\ Fin\ n\ →\ State\ p\ →\ ℝ)\ (hf\ :\ ∀\ i,\ Differentiable\ ℝ\ (f\ i))}
\leanline{4}{(xstar\ :\ State\ p)\ (hstar\ :\ gradient\ (componentPotential\ f)\ xstar\ =\ 0)}
\leanline{4}{(x\ :\ State\ p)\ :\ (∫\ i,\ controlVariateError\ f\ xstar\ x\ i\ ∂uniformIndexLaw\ n)\ =\ 0}
\end{leancode}
\leansource{\texttt{controlVariateError\_\allowbreak{}second\_\allowbreak{}moment} \quad AverageLMC/\allowbreak{}SGLDEstimator.lean, lines 187–211}
\begin{leancode}
\leanline{0}{theorem\ controlVariateError\_second\_moment\ \{p\ n\ :\ ℕ\}\ (hn\ :\ 0\ <\ n)}
\leanline{4}{(f\ :\ Fin\ n\ →\ State\ p\ →\ ℝ)\ (hf\ :\ ∀\ i,\ Differentiable\ ℝ\ (f\ i))}
\leanline{4}{(xstar\ :\ State\ p)\ (hstar\ :\ gradient\ (componentPotential\ f)\ xstar\ =\ 0)}
\leanline{4}{(x\ :\ State\ p)\ :}
\leanline{4}{(∫\ i,\ ‖controlVariateError\ f\ xstar\ x\ i‖\ \textasciicircum{}\ 2\ ∂uniformIndexLaw\ n)\ =}
\leanline{6}{(∫\ i,\ ‖controlVariateEstimator\ f\ xstar\ x\ i‖\ \textasciicircum{}\ 2\ ∂uniformIndexLaw\ n)\ -}
\leanline{8}{‖gradient\ (componentPotential\ f)\ x‖\ \textasciicircum{}\ 2}
\end{leancode}
\leansource{\texttt{controlVariateError\_\allowbreak{}second\_\allowbreak{}moment\_\allowbreak{}le} \quad AverageLMC/\allowbreak{}SGLDEstimator.lean, lines 213–243}
\begin{leancode}
\leanline{0}{theorem\ controlVariateError\_second\_moment\_le\ \{p\ n\ :\ ℕ\}\ (hp\ :\ 0\ <\ p)\ (hn\ :\ 0\ <\ n)}
\leanline{4}{(f\ :\ Fin\ n\ →\ State\ p\ →\ ℝ)\ (hf\ :\ ∀\ i,\ Differentiable\ ℝ\ (f\ i))}
\leanline{4}{(L\ :\ Fin\ n\ →\ ℝ≥0)\ (hLip\ :\ ∀\ i,\ LipschitzWith\ (L\ i)\ (gradient\ (f\ i)))}
\leanline{4}{(xstar\ :\ State\ p)\ (hstar\ :\ gradient\ (componentPotential\ f)\ xstar\ =\ 0)}
\leanline{4}{(x\ :\ State\ p)\ :}
\leanline{4}{(∫\ i,\ ‖controlVariateError\ f\ xstar\ x\ i‖\ \textasciicircum{}\ 2\ ∂uniformIndexLaw\ n)\ ≤}
\leanline{6}{(n\ :\ ℝ)\ \textasciicircum{}\ 2\ *\ rmsComponentSmoothness\ f\ \textasciicircum{}\ 2\ *\ ‖x\ -\ xstar‖\ \textasciicircum{}\ 2}
\end{leancode}
\leanlabel{How the Lean proof follows the argument}
\texttt{control\allowbreak{}Variate\allowbreak{}Error} names the mathematical object defined above. Its defining expression appears in the adjacent statement; subsequent identities may unfold this definition directly.




\texttt{control\allowbreak{}Variate\allowbreak{}Error\_\allowbreak{}centered} gives the formal counterpart of this step.




\texttt{control\allowbreak{}Variate\allowbreak{}Error\_\allowbreak{}second\_\allowbreak{}moment} gives the formal counterpart of this step. The main intermediate facts are \texttt{hmean}, \texttt{hcross} in that order. The final step is the algebraic normalization or ordered arithmetic displayed in the proof.




\texttt{control\allowbreak{}Variate\allowbreak{}Error\_\allowbreak{}second\_\allowbreak{}moment\_\allowbreak{}le} gives the formal counterpart of this step. The main intermediate facts are \texttt{hmean}, \texttt{hcross}, \texttt{hvar} in that order. The final step is the algebraic normalization or ordered arithmetic displayed in the proof.
\end{leanbox}
\subsection{\texorpdfstring{Supply the measurable product-space error with derived \(L^{2}\) moments}{Supply the measurable product-space error with derived L² moments}}
\label{report-step-AverageLMC-SGLDEstimator-lean-6}


For a random state \(X\) with finite second moment, each component gradient evaluated at \(X\) is \(L^{2}\) by Lipschitz composition. A finite independent index can select these \(L^{2}\) functions: express the selector as the finite sum of component functions times the indicators of the measurable index events. This proves product-space \(L^{2}\) membership of the actual estimator and centered error. Measurability follows by checking each fixed finite-index slice. Fubini then integrates the pointwise variance bound over \(X\), producing

\[\mathbb E\|G_I(X)-\nabla f(X)\|^2\le n^2M_{\rm rms}^2\mathbb E\|X-x_*\|^2.\]

These results supply the genuine random perturbation used by the transport comparison; no error variance is assumed.

\begin{leanbox}{Lean: \texttt{memLp\_\allowbreak{}finite\_\allowbreak{}index\_\allowbreak{}selector}, \texttt{controlVariateEstimator\_\allowbreak{}product\_\allowbreak{}memLp}, \texttt{controlVariateError\_\allowbreak{}product\_\allowbreak{}memLp} and 2 more}
\leansource{\texttt{memLp\_\allowbreak{}finite\_\allowbreak{}index\_\allowbreak{}selector} \quad AverageLMC/\allowbreak{}SGLDEstimator.lean, lines 246–257}
\begin{leancode}
\leanline{0}{theorem\ memLp\_finite\_index\_selector\ \{Ω\ E\ :\ Type*\}\ [MeasurableSpace\ Ω]}
\leanline{4}{[NormedAddCommGroup\ E]\ (P\ :\ Measure\ Ω)\ (n\ :\ ℕ)}
\leanline{4}{(g\ :\ Fin\ n\ →\ Ω\ →\ E)\ (hg\ :\ ∀\ i,\ MemLp\ (g\ i)\ 2\ P)\ :}
\leanline{4}{MemLp\ (fun\ q\ :\ Ω\ ×\ Fin\ n\ ↦\ g\ q.2\ q.1)\ 2\ (P.prod\ (uniformIndexLaw\ n))}
\end{leancode}
\leansource{\texttt{controlVariateEstimator\_\allowbreak{}product\_\allowbreak{}memLp} \quad AverageLMC/\allowbreak{}SGLDEstimator.lean, lines 259–274}
\begin{leancode}
\leanline{0}{theorem\ controlVariateEstimator\_product\_memLp\ \{Ω\ :\ Type*\}\ [MeasurableSpace\ Ω]}
\leanline{4}{\{p\ n\ :\ ℕ\}\ (P\ :\ Measure\ Ω)\ [IsProbabilityMeasure\ P]}
\leanline{4}{(f\ :\ Fin\ n\ →\ State\ p\ →\ ℝ)\ (L\ :\ Fin\ n\ →\ ℝ≥0)}
\leanline{4}{(hLip\ :\ ∀\ i,\ LipschitzWith\ (L\ i)\ (gradient\ (f\ i)))}
\leanline{4}{(xstar\ :\ State\ p)\ (X\ :\ Ω\ →\ State\ p)\ (hX\ :\ MemLp\ X\ 2\ P)\ :}
\leanline{4}{MemLp\ (fun\ q\ :\ Ω\ ×\ Fin\ n\ ↦\ controlVariateEstimator\ f\ xstar\ (X\ q.1)\ q.2)}
\leanline{6}{2\ (P.prod\ (uniformIndexLaw\ n))}
\end{leancode}
\leansource{\texttt{controlVariateError\_\allowbreak{}product\_\allowbreak{}memLp} \quad AverageLMC/\allowbreak{}SGLDEstimator.lean, lines 276–292}
\begin{leancode}
\leanline{0}{theorem\ controlVariateError\_product\_memLp\ \{Ω\ :\ Type*\}\ [MeasurableSpace\ Ω]}
\leanline{4}{\{p\ n\ :\ ℕ\}\ (P\ :\ Measure\ Ω)\ [IsProbabilityMeasure\ P]}
\leanline{4}{(f\ :\ Fin\ n\ →\ State\ p\ →\ ℝ)\ (hf\ :\ ∀\ i,\ Differentiable\ ℝ\ (f\ i))}
\leanline{4}{(L\ :\ Fin\ n\ →\ ℝ≥0)\ (hLip\ :\ ∀\ i,\ LipschitzWith\ (L\ i)\ (gradient\ (f\ i)))}
\leanline{4}{(xstar\ :\ State\ p)\ (X\ :\ Ω\ →\ State\ p)\ (hX\ :\ MemLp\ X\ 2\ P)\ :}
\leanline{4}{MemLp\ (fun\ q\ :\ Ω\ ×\ Fin\ n\ ↦\ controlVariateError\ f\ xstar\ (X\ q.1)\ q.2)}
\leanline{6}{2\ (P.prod\ (uniformIndexLaw\ n))}
\end{leancode}
\leansource{\texttt{controlVariateError\_\allowbreak{}product\_\allowbreak{}measurable} \quad AverageLMC/\allowbreak{}SGLDEstimator.lean, lines 294–303}
\begin{leancode}
\leanline{0}{theorem\ controlVariateError\_product\_measurable\ \{Ω\ :\ Type*\}\ [MeasurableSpace\ Ω]}
\leanline{4}{\{p\ n\ :\ ℕ\}\ (f\ :\ Fin\ n\ →\ State\ p\ →\ ℝ)\ (hf\ :\ ∀\ i,\ Differentiable\ ℝ\ (f\ i))}
\leanline{4}{(L\ :\ Fin\ n\ →\ ℝ≥0)\ (hLip\ :\ ∀\ i,\ LipschitzWith\ (L\ i)\ (gradient\ (f\ i)))}
\leanline{4}{(xstar\ :\ State\ p)\ (X\ :\ Ω\ →\ State\ p)\ (hX\ :\ Measurable\ X)\ :}
\leanline{4}{Measurable\ (fun\ q\ :\ Ω\ ×\ Fin\ n\ ↦\ controlVariateError\ f\ xstar\ (X\ q.1)\ q.2)}
\end{leancode}
\leansource{\texttt{controlVariateError\_\allowbreak{}product\_\allowbreak{}second\_\allowbreak{}moment\_\allowbreak{}le} \quad AverageLMC/\allowbreak{}SGLDEstimator.lean, lines 305–325}
\begin{leancode}
\leanline{0}{theorem\ controlVariateError\_product\_second\_moment\_le\ \{Ω\ :\ Type*\}\ [MeasurableSpace\ Ω]}
\leanline{4}{\{p\ n\ :\ ℕ\}\ (hp\ :\ 0\ <\ p)\ (hn\ :\ 0\ <\ n)\ (P\ :\ Measure\ Ω)\ [IsProbabilityMeasure\ P]}
\leanline{4}{(f\ :\ Fin\ n\ →\ State\ p\ →\ ℝ)\ (hf\ :\ ∀\ i,\ Differentiable\ ℝ\ (f\ i))}
\leanline{4}{(L\ :\ Fin\ n\ →\ ℝ≥0)\ (hLip\ :\ ∀\ i,\ LipschitzWith\ (L\ i)\ (gradient\ (f\ i)))}
\leanline{4}{(xstar\ :\ State\ p)\ (hstar\ :\ gradient\ (componentPotential\ f)\ xstar\ =\ 0)}
\leanline{4}{(X\ :\ Ω\ →\ State\ p)\ (hX\ :\ MemLp\ X\ 2\ P)\ :}
\leanline{4}{(∫\ q\ :\ Ω\ ×\ Fin\ n,\ ‖controlVariateError\ f\ xstar\ (X\ q.1)\ q.2‖\ \textasciicircum{}\ 2}
\leanline{6}{∂P.prod\ (uniformIndexLaw\ n))\ ≤}
\leanline{6}{(n\ :\ ℝ)\ \textasciicircum{}\ 2\ *\ rmsComponentSmoothness\ f\ \textasciicircum{}\ 2\ *}
\leanline{8}{∫\ ω,\ ‖X\ ω\ -\ xstar‖\ \textasciicircum{}\ 2\ ∂P}
\end{leancode}
\leanlabel{How the Lean proof follows the argument}
\texttt{mem\allowbreak{}Lp\_\allowbreak{}finite\_\allowbreak{}index\_\allowbreak{}selector} gives the formal counterpart of this step. The main intermediate facts are \texttt{hset}, \texttt{hsum} in that order. Its remaining equalities are discharged by the simplifications displayed in the proof.




\texttt{control\allowbreak{}Variate\allowbreak{}Estimator\_\allowbreak{}product\_\allowbreak{}mem\allowbreak{}Lp} gives the formal counterpart of this step. The main intermediate facts are \texttt{hg} in that order. Named proof applications include \texttt{mem\allowbreak{}Lp\_\allowbreak{}finite\_\allowbreak{}index\_\allowbreak{}selector}, \texttt{hg.\allowbreak{}ae\_\allowbreak{}eq}; the adjacent proof shows their hypotheses and order of use.




\texttt{control\allowbreak{}Variate\allowbreak{}Error\_\allowbreak{}product\_\allowbreak{}mem\allowbreak{}Lp} gives the formal counterpart of this step. The main intermediate facts are \texttt{hg}, \texttt{hsum} in that order. Named proof applications include \texttt{h.\allowbreak{}ae\_\allowbreak{}eq}, \texttt{mem\allowbreak{}Lp\_\allowbreak{}finset\allowbreak{}Sum}; the adjacent proof shows their hypotheses and order of use.




\texttt{control\allowbreak{}Variate\allowbreak{}Error\_\allowbreak{}product\_\allowbreak{}measurable} gives the formal counterpart of this step. Named proof applications include \texttt{measurable\_\allowbreak{}from\_\allowbreak{}prod\_\allowbreak{}countable\_\allowbreak{}left}; the adjacent proof shows their hypotheses and order of use.




\texttt{control\allowbreak{}Variate\allowbreak{}Error\_\allowbreak{}product\_\allowbreak{}second\_\allowbreak{}moment\_\allowbreak{}le} gives the formal counterpart of this step. The main intermediate facts are \texttt{hi}, \texttt{hdist} in that order.
\end{leanbox}


Earlier results used here: \cite[Section~3.13.4]{LangevinLean2026ULASC}.
\subsection{\texorpdfstring{Construct the unique zero-score minimizer}{Construct the unique zero-score minimizer}}
\label{report-step-AverageLMC-SGLDEstimator-lean-7}


Positive strong convexity gives a quadratic lower bound, so the sublevel set \(f(x)\) at most \(f(0)\) is bounded. Continuity and the extreme-value theorem yield a global minimizer. Fermat theorem makes its actual Frechet derivative and gradient zero. If two points had zero gradients, strong gradient monotonicity would force \(m\) times their squared distance to be at most zero. Since \(m\) is positive, the points coincide. The resulting unique minimizer provides the anchor required by the finite-sum algorithm.

\begin{leanbox}{Lean: \texttt{exists\_\allowbreak{}unique\_\allowbreak{}minimizer}}
\leansource{\texttt{exists\_\allowbreak{}unique\_\allowbreak{}minimizer} \quad AverageLMC/\allowbreak{}SGLDEstimator.lean, lines 329–355}
\begin{leancode}
\leanline{0}{theorem\ exists\_unique\_minimizer\ \{p\ :\ ℕ\}\ (f\ :\ State\ p\ →\ ℝ)\ (m\ :\ ℝ)\ (hm\ :\ 0\ <\ m)}
\leanline{4}{(hf\ :\ Differentiable\ ℝ\ f)\ (hsc\ :\ StrongConvexOn\ Set.univ\ m\ f)\ :}
\leanline{4}{∃!\ xstar,\ (∀\ x,\ f\ xstar\ ≤\ f\ x)\ ∧\ gradient\ f\ xstar\ =\ 0}
\end{leancode}
\leanlabel{How the Lean proof follows the argument}
\texttt{exists\_\allowbreak{}unique\_\allowbreak{}minimizer} gives the formal counterpart of this step. The main intermediate facts are \texttt{hb}, \texttt{hz}, \texttt{hmono}, \texttt{hsq} in that order. Named proof applications include \texttt{is\allowbreak{}Bounded\_\allowbreak{}iff\_\allowbreak{}forall\_\allowbreak{}norm\_\allowbreak{}le.\allowbreak{}mpr}, \texttt{sub\_\allowbreak{}eq\_\allowbreak{}zero.\allowbreak{}mp}; the adjacent proof shows their hypotheses and order of use.
\end{leanbox}


Earlier results used here: \cite[Section~3.2.3]{LangevinLean2026ULASC}; \cite[Section~3.26.5]{LangevinLean2026ULASC}.
\section{\texorpdfstring{Transport cost of a conditionally centred perturbation}{Transport cost of a conditionally centred perturbation}}
\label{proof-module-AverageLMC-CenteredPerturbation-lean}


This section proves the variance-addition inequality for quadratic transport that underlies the stochastic-gradient bound. Let \(s\) be a random input with law \(P\) on a measurable space \(S\), let \(i\) be an independent random index with law \(Q\) on a measurable space \(I\), let \(F(s)\) be an output in a Hilbert space \(E\) (for instance \(\mathbb R^p\)), and let \(D(s,i)\) be a perturbation whose mean over the index vanishes for every fixed input. Then, for every probability measure \(\pi\) on \(E\) with finite second moment,


\[\mathsf W_2^2\bigl(\mathcal L(F(s)+D(s,i)),\pi\bigr)\le\mathsf W_2^2\bigl(\mathcal L(F(s)),\pi\bigr)+\mathbb E\|D(s,i)\|^2,\]


where \(\mathcal L\) denotes the law under \(P\otimes Q\). For SGLD, \(s=(\vartheta,\xi)\), \(F(s)=\vartheta-h\nabla f(\vartheta)+\sqrt{2h}\,\xi\) is the full-gradient Langevin output, \(i=I\) is the uniform data index and \(D=-h(g_I(\vartheta)-\nabla f(\vartheta))\) is the stochastic-gradient error. The proof uses an arbitrary finite-cost coupling of \(\mathcal L(F(s))\) with \(\pi\), lifts it to the input by the construction of the previous section, adjoins the independent index, and shows that the cross term of the squared distance vanishes.



Earlier results used here: \cite[Section~3.21]{LangevinLean2026ULASC}.
\subsection{\texorpdfstring{The squared \(L_2\) norm as an integral}{The squared L₂ norm as an integral}}
\label{report-step-AverageLMC-CenteredPerturbation-lean-1}


\textbf{Claim.} Let \((\Omega,\mathcal F,P)\) be a measure space, \(E\) a real inner product space and \(f:\Omega\to E\) square integrable, \(\int\|f\|^2\,dP<\infty\) (with \(f\) almost everywhere strongly measurable). Let \(\|f\|_{L_2(P)}=\bigl(\int\|f\|^2\,dP\bigr)^{1/2}\), a finite nonnegative number. Then


\[\|f\|_{L_2(P)}^2=\int_\Omega\|f(\omega)\|^2\,P(d\omega).\]


\textbf{Proof.} Let \([f]\) denote the class of \(f\) in the Hilbert space \(L_2(P;E)\), whose inner product is \(\langle[g],[k]\rangle=\int\langle g(\omega),k(\omega)\rangle\,P(d\omega)\) and whose norm is \(\|f\|_{L_2(P)}\). Then


\[\|f\|_{L_2(P)}^2=\langle[f],[f]\rangle=\int_\Omega\langle f(\omega),f(\omega)\rangle\,P(d\omega)=\int_\Omega\|f(\omega)\|^2\,P(d\omega),\]


where the representative of \([f]\) used in the integral agrees with \(f\) almost everywhere, which does not change the integral. \(\square\)




The finiteness assumption matters because the \(L_2\) norm is defined with values in \([0,\infty]\): for a function that is not square integrable, the real number associated with \(\infty\) would be \(0\) and the identity would fail.

\begin{leanbox}{Lean: \texttt{eLpNorm\_\allowbreak{}toReal\_\allowbreak{}sq\_\allowbreak{}eq\_\allowbreak{}integral}}
\leansource{\texttt{eLpNorm\_\allowbreak{}toReal\_\allowbreak{}sq\_\allowbreak{}eq\_\allowbreak{}integral} \quad AverageLMC/\allowbreak{}CenteredPerturbation.lean, lines 29–38}
\begin{leancode}
\leanline{0}{theorem\ eLpNorm\_toReal\_sq\_eq\_integral\ (P\ :\ Measure\ Ω)\ (f\ :\ Ω\ →\ E)\ (hf\ :\ MemLp\ f\ 2\ P)\ :}
\leanline{4}{(eLpNorm\ f\ 2\ P).toReal\ \textasciicircum{}\ 2\ =\ ∫\ ω,\ ‖f\ ω‖\ \textasciicircum{}\ 2\ ∂P}
\end{leancode}
\leanlabel{How the Lean proof follows the argument}
\texttt{e\allowbreak{}Lp\allowbreak{}Norm\_\allowbreak{}to\allowbreak{}Real\_\allowbreak{}sq\_\allowbreak{}eq\_\allowbreak{}integral} takes \texttt{hf : Mem\allowbreak{}Lp f 2 P}, the statement that \(f\) is almost everywhere strongly measurable with finite \(L^2\) seminorm. \texttt{e\allowbreak{}Lp\allowbreak{}Norm f 2 P} is the seminorm with values in \texttt{\ensuremath{\mathbb R}\ensuremath{\ge}0\ensuremath{\infty}} and \texttt{.\allowbreak{}to\allowbreak{}Real} converts it to a real number. The \texttt{calc} block follows the proof: \texttt{Lp.\allowbreak{}norm\_\allowbreak{}to\allowbreak{}Lp} identifies the real seminorm with the norm of \texttt{hf.\allowbreak{}to\allowbreak{}Lp f}, the class of \(f\) in \(L^2\); \texttt{real\_\allowbreak{}inner\_\allowbreak{}self\_\allowbreak{}eq\_\allowbreak{}norm\_\allowbreak{}sq} and \texttt{L2.\allowbreak{}inner\_\allowbreak{}def} write the squared norm as the integral of the pointwise inner product; and \texttt{integral\_\allowbreak{}congr\_\allowbreak{}ae} with \texttt{hf.\allowbreak{}coe\allowbreak{}Fn\_\allowbreak{}to\allowbreak{}Lp} (the representative agrees with \(f\) almost everywhere) replaces the representative by \(f\).
\end{leanbox}
\subsection{\texorpdfstring{Adding a centred perturbation adds its second moment}{Adding a centred perturbation adds its second moment}}
\label{report-step-AverageLMC-CenteredPerturbation-lean-2}


\textbf{Claim.} Let \(S\) and \(I\) be measurable spaces with probability measures \(P\) and \(Q\). Let \(E\) be a real Hilbert space that is standard Borel and second countable with its Borel sets (for instance \(\mathbb R^p\)). Let \(F:S\to E\) be measurable and square integrable for \(P\), and let \(D:S\times I\to E\) be jointly measurable and square integrable for \(P\otimes Q\), such that for every \(s\in S\) the function \(D(s,\cdot)\) is \(Q\)-integrable with


\[\int_ID(s,i)\,Q(di)=0.\qquad(1)\]


Let \(\pi\) be a probability measure on \(E\) with \(\int\|x\|^2\pi(dx)<\infty\), let \(\mu=P\circ F^{-1}\) be the law of \(F\), and let \(\nu\) be the law of \((s,i)\mapsto F(s)+D(s,i)\) under \(P\otimes Q\). Then \(\mathsf W_2(\mu,\pi)\) and \(\mathsf W_2(\nu,\pi)\) are finite and


\[\mathsf W_2(\nu,\pi)^2\le\mathsf W_2(\mu,\pi)^2+\int_{S\times I}\|D(s,i)\|^2\,(P\otimes Q)(ds,di).\]


\textbf{Proof.} \emph{Finiteness.} The identity map has finite second moment under \(\mu\), because \(\int\|x\|^2\mu(dx)=\int\|F\|^2dP<\infty\), and under \(\nu\), because \(F(s)+D(s,i)\) is a sum of two square-integrable functions on \(S\times I\). For two probability measures with finite second moments the product coupling has finite cost, since \(\|x-y\|^2\le2\|x\|^2+2\|y\|^2\); hence \(\mathsf W_2(\mu,\pi)<\infty\) and \(\mathsf W_2(\nu,\pi)<\infty\).




\emph{Reduction.} By the first result of the previous section (with \(a=1\) and \(b=\int\|D\|^2d(P\otimes Q)\)), it suffices to show that \(\mathsf W_2(\nu,\pi)^2\le C(\gamma)^2+\int\|D\|^2d(P\otimes Q)\) for every coupling \(\gamma\) of \(\mu\) and \(\pi\) with finite cost \(C(\gamma)\).




\emph{Construction.} Fix such a \(\gamma\) and let \(R\) be the lifted coupling of the previous section: a probability measure on \(S\times E\) with first marginal \(P\), second marginal \(\pi\), under which \((F(s),y)\) has law \(\gamma\). Work on \((S\times E)\times I\) with the product probability \(R\otimes Q\) and the random variables


\[A=F(s)-y,\qquad B=D(s,i),\qquad X=F(s)+D(s,i),\qquad Y=y.\]


Since the first marginal of \(R\) is \(P\), the image of \(R\otimes Q\) under \(((s,y),i)\mapsto(s,i)\) is \(P\otimes Q\). Consequently \(X\) has law \(\nu\), and \(Y\) has law \(\pi\). Both \(A\) and \(B\) are square integrable: \(\int\|A\|^2d(R\otimes Q)=\int\|F(s)-y\|^2R(ds,dy)=\int\|x-y\|^2\gamma(dx,dy)=C(\gamma)^2<\infty\), where the first equality integrates out the probability measure \(Q\) and the second is the change of variables under \((s,y)\mapsto(F(s),y)\); and \(\int\|B\|^2d(R\otimes Q)=\int\|D\|^2d(P\otimes Q)\) by the change of variables under \(((s,y),i)\mapsto(s,i)\). We record these as


\[\int\|A\|^2\,d(R\otimes Q)=C(\gamma)^2,\qquad\int\|B\|^2\,d(R\otimes Q)=\int\|D\|^2\,d(P\otimes Q).\qquad(2)\]


\emph{The cross term vanishes.} By the Cauchy–Schwarz inequality \(|\langle A,B\rangle|\le\|A\|\,\|B\|\), and \(\|A\|\,\|B\|\) is integrable as a product of two square-integrable functions (Hölder's inequality); so \(\langle A,B\rangle\) is integrable. By Fubini's theorem,


\[\int\langle A,B\rangle\,d(R\otimes Q)=\int_{S\times E}\Bigl(\int_I\bigl\langle F(s)-y,\,D(s,i)\bigr\rangle\,Q(di)\Bigr)R(ds,dy).\]


For fixed \((s,y)\), the map \(v\mapsto\langle F(s)-y,v\rangle\) is a continuous linear functional on \(E\), and continuous linear functionals commute with Bochner integrals of integrable functions. Using the integrability of \(D(s,\cdot)\) and (1), the inner integral equals \(\langle F(s)-y,\int_ID(s,i)Q(di)\rangle=\langle F(s)-y,0\rangle=0\). Hence \(\int\langle A,B\rangle\,d(R\otimes Q)=0\).




\emph{Energy identity.} Pointwise \(\|A+B\|^2=\|A\|^2+2\langle A,B\rangle+\|B\|^2\), and the three terms are integrable. Integrating and using the vanishing of the cross term and (2),


\[\int\|A+B\|^2\,d(R\otimes Q)=C(\gamma)^2+\int\|D\|^2\,d(P\otimes Q).\qquad(3)\]


\emph{Conclusion.} The pair \((X,Y)\) defines a coupling of \(\nu\) and \(\pi\), and \(X-Y=A+B\). Therefore \(\mathsf W_2(\nu,\pi)\le\bigl(\int\|A+B\|^2d(R\otimes Q)\bigr)^{1/2}\). Squaring this inequality between nonnegative numbers, using the identity of Step 1 for the square-integrable function \(A+B\), and inserting (3) gives the required bound. \(\square\)




Two hypotheses deserve comment. Condition (1) must hold for every input \(s\), not only on average; it is what makes the cross term vanish after integrating out the index. The finiteness of the second moment of \(\pi\) is used only to make both distances finite, so that the bound over couplings passes to the distance.

\begin{leanbox}{Lean: \texttt{wasserstein2\_\allowbreak{}sq\_\allowbreak{}centered\_\allowbreak{}perturbation}}
\leansource{\texttt{wasserstein2\_\allowbreak{}sq\_\allowbreak{}centered\_\allowbreak{}perturbation} \quad AverageLMC/\allowbreak{}CenteredPerturbation.lean, lines 42–146}
\begin{leancode}
\leanline{0}{theorem\ wasserstein2\_sq\_centered\_perturbation}
\leanline{4}{(P\ :\ Measure\ S)\ [IsProbabilityMeasure\ P]\ (Q\ :\ Measure\ I)\ [IsProbabilityMeasure\ Q]}
\leanline{4}{(F\ :\ S\ →\ E)\ (hF\ :\ Measurable\ F)\ (hF2\ :\ MemLp\ F\ 2\ P)}
\leanline{4}{(D\ :\ S\ →\ I\ →\ E)\ (hD\ :\ Measurable\ (Function.uncurry\ D))}
\leanline{4}{(hD2\ :\ MemLp\ (Function.uncurry\ D)\ 2\ (P.prod\ Q))}
\leanline{4}{(hDi\ :\ ∀\ s,\ Integrable\ (D\ s)\ Q)\ (hmean\ :\ ∀\ s,\ ∫\ i,\ D\ s\ i\ ∂Q\ =\ 0)}
\leanline{4}{(π\ :\ Measure\ E)\ [IsProbabilityMeasure\ π]\ (hπ\ :\ MemLp\ (fun\ x\ :\ E\ ↦\ x)\ 2\ π)\ :}
\leanline{4}{ULA.wasserstein2Real\ ((P.prod\ Q).map\ (fun\ q\ ↦\ F\ q.1\ +\ D\ q.1\ q.2))\ π\ \textasciicircum{}\ 2\ ≤}
\leanline{6}{ULA.wasserstein2Real\ (P.map\ F)\ π\ \textasciicircum{}\ 2\ +}
\leanline{8}{∫\ q\ :\ S\ ×\ I,\ ‖D\ q.1\ q.2‖\ \textasciicircum{}\ 2\ ∂P.prod\ Q}
\end{leancode}
\leanlabel{How the Lean proof follows the argument}
\texttt{wasserstein2\_\allowbreak{}sq\_\allowbreak{}centered\_\allowbreak{}perturbation} has the hypotheses \texttt{h\allowbreak{}F : Measurable F}, \texttt{h\allowbreak{}F2 : Mem\allowbreak{}Lp F 2 P}, \texttt{h\allowbreak{}D : Measurable (Function.\allowbreak{}uncurry D)} (joint measurability), \texttt{h\allowbreak{}D2 : Mem\allowbreak{}Lp (Function.\allowbreak{}uncurry D) 2 (P.\allowbreak{}prod Q)}, \texttt{h\allowbreak{}Di : \ensuremath{\forall} s, Integrable (D s) Q}, \texttt{hmean : \ensuremath{\forall} s, \ensuremath{\int} i, D s i \ensuremath{\partial}Q = 0}, which is (1), and \texttt{h\ensuremath{\pi} : Mem\allowbreak{}Lp (fun x : E \ensuremath{\mapsto} x) 2 \ensuremath{\pi}}, the finite second moment of \(\pi\). The law \(\nu\) is the pushforward \texttt{(P.\allowbreak{}prod Q).\allowbreak{}map (fun q \ensuremath{\mapsto} F q.\allowbreak{}1 + D q.\allowbreak{}1 q.\allowbreak{}2)}, and distances are \texttt{ULA.\allowbreak{}wasserstein2\allowbreak{}Real}.



\begin{enumerate}


\item \emph{Finiteness.} \texttt{h\allowbreak{}Fm}, \texttt{hplus} and \texttt{h\ensuremath{\nu}} give the second moments of \(\mu\) and \(\nu\) (through \texttt{mem\allowbreak{}Lp\_\allowbreak{}map\_\allowbreak{}measure\_\allowbreak{}iff}), and \texttt{hfinite}, \texttt{houtput} are the finiteness of both distances from \texttt{ULA.\allowbreak{}wasserstein2\_\allowbreak{}ne\_\allowbreak{}top\_\allowbreak{}of\_\allowbreak{}mem\allowbreak{}Lp}.

\item \emph{Reduction.} The \texttt{suffices} line inserts the factor \(1\), and \texttt{apply wasserstein2\_\allowbreak{}sq\_\allowbreak{}bound\_\allowbreak{}of\_\allowbreak{}finite\_\allowbreak{}couplings} reduces the goal to a fixed coupling \texttt{c} with finite cost \texttt{hc}.

\item \emph{Construction.} \texttt{R := lifted\allowbreak{}Output\allowbreak{}Coupling P F h\allowbreak{}F c} with its marginal identities \texttt{h\allowbreak{}Rfst}, \texttt{h\allowbreak{}Rout}, \texttt{h\allowbreak{}Rsnd}; \texttt{h\allowbreak{}RI} is the identity of the image of \texttt{R.\allowbreak{}prod Q} under \(((s,y),i)\mapsto(s,i)\) with \texttt{P.\allowbreak{}prod Q}, from \texttt{Measure.\allowbreak{}map\_\allowbreak{}prod\_\allowbreak{}map}. \texttt{A} and \texttt{B} are the variables of the proof; \texttt{h\allowbreak{}A} and \texttt{h\allowbreak{}B} are their square integrability, transported from \texttt{h\allowbreak{}Ac} (finite cost of the coupling) and from \texttt{h\allowbreak{}D2} by \texttt{comp\_\allowbreak{}of\_\allowbreak{}map}.

\item \emph{Cross term.} \texttt{hcross} is the integrability of \(\langle A,B\rangle\), from the integrability of \(\|A\|\|B\|\) (\texttt{integrable\_\allowbreak{}mul}) and \texttt{abs\_\allowbreak{}real\_\allowbreak{}inner\_\allowbreak{}le\_\allowbreak{}norm}. In \texttt{hzero}, \texttt{integral\_\allowbreak{}prod} is Fubini's theorem and \texttt{hz} evaluates the inner integral: \texttt{(inner\allowbreak{}SL \ensuremath{\mathbb R} (F r.\allowbreak{}1 -\allowbreak{} r.\allowbreak{}2)).\allowbreak{}integral\_\allowbreak{}comp\_\allowbreak{}comm (h\allowbreak{}Di r.\allowbreak{}1)} moves the continuous linear functional \(\langle F(s)-y,\cdot\rangle\) through the integral, and \texttt{hmean} with \texttt{inner\_\allowbreak{}zero\_\allowbreak{}right} gives \(0\).

\item \emph{Energy.} \texttt{henergy} is (3): \texttt{norm\_\allowbreak{}add\_\allowbreak{}sq\_\allowbreak{}real} expands the square and \texttt{integral\_\allowbreak{}add}, \texttt{integral\_\allowbreak{}const\_\allowbreak{}mul} and \texttt{hzero} evaluate the integral. \texttt{h\allowbreak{}Acost} and \texttt{h\allowbreak{}Bcost} are the two identities (2), using \texttt{integral\_\allowbreak{}fun\_\allowbreak{}fst} (integration of a function of the first factor against a product with the probability measure \(Q\)), \texttt{integral\_\allowbreak{}map} with \texttt{h\allowbreak{}Rout} or \texttt{h\allowbreak{}RI}, and Step 1.

\item \emph{Conclusion.} \texttt{h\allowbreak{}Xlaw} and \texttt{h\allowbreak{}Ylaw} identify the laws of \(X\) and \(Y\); \texttt{hw} is the bound of \(\mathsf W_2\) by the \(L^2\) distance of two random variables, \texttt{ULA.\allowbreak{}wasserstein2\_\allowbreak{}le\_\allowbreak{}random\_\allowbreak{}cost\_\allowbreak{}real}; \texttt{hs} squares it with \texttt{sq\_\allowbreak{}le\_\allowbreak{}sq\ensuremath{{}_{0}}}, Step 1 turns the \(L^2\) norm into an integral, \texttt{heq} rewrites \(X-Y=A+B\), and the final \texttt{simpa} inserts \texttt{henergy}, \texttt{h\allowbreak{}Acost} and \texttt{h\allowbreak{}Bcost}.

\end{enumerate}
\end{leanbox}


Earlier results used here: \cite[Section~3.17.11]{LangevinLean2026ULASC}; \cite[Section~3.17.2]{LangevinLean2026ULASC}; \cite[Section~3.17.4]{LangevinLean2026ULASC}; \cite[Section~3.21.2]{LangevinLean2026ULASC}.
\section{\texorpdfstring{The SGLD transition and its variance decomposition}{The SGLD transition and its variance decomposition}}
\label{proof-module-AverageLMC-SGLDLaw-lean}


This section defines the stochastic-gradient Langevin dynamics with fixed-point control variates and proves the variance decomposition on which its analysis rests. The potential is a finite sum \(f=\sum_{i=1}^nf_i\) of functions \(f_i:\mathbb R^p\to\mathbb R\), and \(\theta_*\in\mathbb R^p\) is a point with \(\nabla f(\theta_*)=0\) (the minimizer of \(f\) in the application). The control-variate estimator of \(\nabla f(x)\) and its error are


\[g_i(x)=n\bigl(\nabla f_i(x)-\nabla f_i(\theta_*)\bigr),\qquad e_i(x)=g_i(x)-\nabla f(x)\qquad(i=1,\dots,n).\]


One SGLD step from \(\vartheta\sim\nu\) is \(\vartheta^+=\vartheta-h\,g_I(\vartheta)+\sqrt{2h}\,\xi\), where \(\xi\sim\mathcal N(0,I_p)\) and the index \(I\), uniform on \(\{1,\dots,n\}\), are independent of each other and of \(\vartheta\); \(\nu^+\) denotes the law of \(\vartheta^+\), and \(\nu_k\) the law after \(k\) steps with fresh innovations. Writing \(\nu^+_{\rm LMC}\) for the law of the full-gradient step \(\vartheta-h\nabla f(\vartheta)+\sqrt{2h}\,\xi\), the main result is


\[\mathsf W_2^2(\nu^+,\pi)\le\mathsf W_2^2(\nu^+_{\rm LMC},\pi)+h^2\,\mathbb E\|e_I(\vartheta)\|^2,\qquad\mathbb E\|e_I(\vartheta)\|^2=\mathbb E\|g_I(\vartheta)\|^2-\mathbb E\|\nabla f(\vartheta)\|^2,\]


for every probability measure \(\pi\) with finite second moment, together with \(\mathbb E\|g_I(\vartheta)\|^2=n\sum_{i=1}^n\mathbb E\|\nabla f_i(\vartheta)-\nabla f_i(\theta_*)\|^2\). The SGLD step is thus the LMC step plus a perturbation that is centred given \((\vartheta,\xi)\), and the variance-addition inequality of the previous sections applies. The section also checks that all laws involved are probability measures with finite second moments.



Earlier results used here: \cite[Section~3.10]{LangevinLean2026ULALSI}.
\subsection{\texorpdfstring{The update, the transition and the chain of laws}{The update, the transition and the chain of laws}}
\label{report-step-AverageLMC-SGLDLaw-lean-1}


\textbf{Definitions.} Let \(n,p\in\mathbb N\), \(f_1,\dots,f_n:\mathbb R^p\to\mathbb R\), \(f=\sum_if_i\), \(\theta_*\in\mathbb R^p\) and \(h\in\mathbb R\) (with \(\sqrt{2h}\) read as \(0\) if \(h<0\); only \(h\ge0\) is of interest). The SGLD update map is


\[\Phi(x,z,i)=x-h\,g_i(x)+\sqrt{2h}\,z\qquad(x,z\in\mathbb R^p,\ i\in\{1,\dots,n\}).\]


For a measure \(\nu\) on \(\mathbb R^p\), the SGLD transition gives the image measure \(\nu^+=(\nu\otimes\gamma_p\otimes U_n)\circ\Phi^{-1}\), where \(\gamma_p=\mathcal N(0,I_p)\) is the standard Gaussian measure and \(U_n\) the uniform probability on \(\{1,\dots,n\}\), \(U_n(\{i\})=1/n\). The chain of laws is \(\nu_0=\nu\), \(\nu_{k+1}=(\nu_k)^+\); each step uses a new independent pair \((\xi,I)\).




\textbf{Claim.} (1) For all \(x,z,i\),


\[\Phi(x,z,i)=\bigl(x-h\nabla f(x)+\sqrt{2h}\,z\bigr)+(-h)\,e_i(x).\qquad(1)\]


(2) If every \(f_i\) is differentiable with Lipschitz gradient and \(\nabla f\) is Lipschitz, then \(\Phi\) is measurable on \(\mathbb R^p\times\mathbb R^p\times\{1,\dots,n\}\).




(3) If \(n\ge1\) and \(\nu\) is a probability measure, then \(\nu^+\) and every \(\nu_k\) are probability measures.




\textbf{Proof.} (1) By definition \(g_i(x)=e_i(x)+\nabla f(x)\), so \(-h\,g_i(x)=-h\nabla f(x)+(-h)e_i(x)\); insert this in \(\Phi\).




(2) By (1), \(\Phi\) is the sum of \((x,z)\mapsto x-h\nabla f(x)+\sqrt{2h}z\), which is continuous because \(\nabla f\) is Lipschitz, and \((x,i)\mapsto-h\,e_i(x)\). For each fixed \(i\) the map \(x\mapsto e_i(x)=n(\nabla f_i(x)-\nabla f_i(\theta_*))-\nabla f(x)\) is continuous, and a function on \(\mathbb R^p\times\{1,\dots,n\}\) whose sections are measurable is measurable, since \(\{1,\dots,n\}\) is finite with all subsets measurable. Sums of measurable functions are measurable.




(3) For \(n\ge1\), \(U_n\) is a probability measure (it gives mass \(1/n\) to each of the \(n\) points). The product \(\nu\otimes\gamma_p\otimes U_n\) of probability measures is a probability measure, and so is its image under \(\Phi\). Induction on \(k\) gives the statement for \(\nu_k\): \(\nu_0=\nu\), and if \(\nu_k\) is a probability measure, so is \(\nu_{k+1}=(\nu_k)^+\). \(\square\)




Identity (1) displays the SGLD step as the LMC step of Theorem 1 plus the perturbation \(-h\,e_I(\vartheta)\); everything that follows is built on it.

\begin{leanbox}{Lean: \texttt{sgldProposal}, \texttt{sgldLawStep}, \texttt{sgldLaws} and 5 more}
\leansource{\texttt{sgldProposal} \quad AverageLMC/\allowbreak{}SGLDLaw.lean, lines 14–16}
\begin{leancode}
\leanline{0}{def\ sgldProposal\ \{p\ n\ :\ ℕ\}\ (f\ :\ Fin\ n\ →\ State\ p\ →\ ℝ)\ (anchor\ :\ State\ p)}
\leanline{4}{(h\ :\ ℝ)\ (x\ z\ :\ State\ p)\ (i\ :\ Fin\ n)\ :\ State\ p\ :=}
\leanline{2}{x\ -\ h\ •\ controlVariateEstimator\ f\ anchor\ x\ i\ +\ Real.sqrt\ (2\ *\ h)\ •\ z}
\end{leancode}
\leansource{\texttt{sgldLawStep} \quad AverageLMC/\allowbreak{}SGLDLaw.lean, lines 20–23}
\begin{leancode}
\leanline{0}{def\ sgldLawStep\ \{p\ n\ :\ ℕ\}\ (f\ :\ Fin\ n\ →\ State\ p\ →\ ℝ)\ (anchor\ :\ State\ p)}
\leanline{4}{(h\ :\ ℝ)\ (μ\ :\ Measure\ (State\ p))\ :\ Measure\ (State\ p)\ :=}
\leanline{2}{((μ.prod\ (stdGaussian\ (State\ p))).prod\ (uniformIndexLaw\ n)).map}
\leanline{4}{(fun\ q\ ↦\ sgldProposal\ f\ anchor\ h\ q.1.1\ q.1.2\ q.2)}
\end{leancode}
\leansource{\texttt{sgldLaws} \quad AverageLMC/\allowbreak{}SGLDLaw.lean, lines 27–30}
\begin{leancode}
\leanline{0}{def\ sgldLaws\ \{p\ n\ :\ ℕ\}\ (f\ :\ Fin\ n\ →\ State\ p\ →\ ℝ)\ (anchor\ :\ State\ p)}
\leanline{4}{(h\ :\ ℝ)\ (μ\ :\ Measure\ (State\ p))\ :\ ℕ\ →\ Measure\ (State\ p)}
\leanline{2}{|\ 0\ =>\ μ}
\leanline{2}{|\ k\ +\ 1\ =>\ sgldLawStep\ f\ anchor\ h\ (sgldLaws\ f\ anchor\ h\ μ\ k)}
\end{leancode}
\leansource{\texttt{sgldProposal\_\allowbreak{}eq\_\allowbreak{}full\_\allowbreak{}gradient} \quad AverageLMC/\allowbreak{}SGLDLaw.lean, lines 32–37}
\begin{leancode}
\leanline{0}{theorem\ sgldProposal\_eq\_full\_gradient\ \{p\ n\ :\ ℕ\}\ (f\ :\ Fin\ n\ →\ State\ p\ →\ ℝ)}
\leanline{4}{(anchor\ x\ z\ :\ State\ p)\ (h\ :\ ℝ)\ (i\ :\ Fin\ n)\ :}
\leanline{4}{sgldProposal\ f\ anchor\ h\ x\ z\ i\ =\ ULA.malaProposal\ (componentPotential\ f)\ h\ x\ z\ +}
\leanline{6}{(-h)\ •\ controlVariateError\ f\ anchor\ x\ i}
\end{leancode}
\leansource{\texttt{sgldProposal\_\allowbreak{}measurable} \quad AverageLMC/\allowbreak{}SGLDLaw.lean, lines 39–48}
\begin{leancode}
\leanline{0}{theorem\ sgldProposal\_measurable\ \{p\ n\ :\ ℕ\}\ (f\ :\ Fin\ n\ →\ State\ p\ →\ ℝ)}
\leanline{4}{(hf\ :\ ∀\ i,\ Differentiable\ ℝ\ (f\ i))\ (L\ :\ Fin\ n\ →\ ℝ≥0)}
\leanline{4}{(hLip\ :\ ∀\ i,\ LipschitzWith\ (L\ i)\ (gradient\ (f\ i)))}
\leanline{4}{(M\ :\ ℝ≥0)\ (hfull\ :\ LipschitzWith\ M\ (gradient\ (componentPotential\ f)))}
\leanline{4}{(anchor\ :\ State\ p)\ (h\ :\ ℝ)\ :}
\leanline{4}{Measurable\ (fun\ q\ :\ (State\ p\ ×\ State\ p)\ ×\ Fin\ n\ ↦}
\leanline{6}{sgldProposal\ f\ anchor\ h\ q.1.1\ q.1.2\ q.2)}
\end{leancode}
\leansource{\texttt{sgldLawStep\_\allowbreak{}isProbability} \quad AverageLMC/\allowbreak{}SGLDLaw.lean, lines 50–58}
\begin{leancode}
\leanline{0}{theorem\ sgldLawStep\_isProbability\ \{p\ n\ :\ ℕ\}\ (hn\ :\ 0\ <\ n)}
\leanline{4}{(f\ :\ Fin\ n\ →\ State\ p\ →\ ℝ)\ (\_hf\ :\ ∀\ i,\ Differentiable\ ℝ\ (f\ i))}
\leanline{4}{(L\ :\ Fin\ n\ →\ ℝ≥0)\ (\_hLip\ :\ ∀\ i,\ LipschitzWith\ (L\ i)\ (gradient\ (f\ i)))}
\leanline{4}{(M\ :\ ℝ≥0)\ (\_hfull\ :\ LipschitzWith\ M\ (gradient\ (componentPotential\ f)))}
\leanline{4}{(anchor\ :\ State\ p)\ (h\ :\ ℝ)\ (μ\ :\ Measure\ (State\ p))\ [IsProbabilityMeasure\ μ]\ :}
\leanline{4}{IsProbabilityMeasure\ (sgldLawStep\ f\ anchor\ h\ μ)}
\end{leancode}
\leansource{\texttt{sgldLawStep\_\allowbreak{}probability} \quad AverageLMC/\allowbreak{}SGLDLaw.lean, lines 60–66}
\begin{leancode}
\leanline{0}{instance\ sgldLawStep\_probability\ \{p\ n\ :\ ℕ\}\ [NeZero\ n]}
\leanline{4}{(f\ :\ Fin\ n\ →\ State\ p\ →\ ℝ)\ (anchor\ :\ State\ p)\ (h\ :\ ℝ)}
\leanline{4}{(μ\ :\ Measure\ (State\ p))\ [IsProbabilityMeasure\ μ]\ :}
\leanline{4}{IsProbabilityMeasure\ (sgldLawStep\ f\ anchor\ h\ μ)\ :=\ by}
\leanline{2}{have\ :=\ uniformIndexLaw\_probability\ n\ (Nat.pos\_of\_ne\_zero\ (NeZero.ne\ n))}
\leanline{2}{unfold\ sgldLawStep}
\leanline{2}{infer\_instance}
\end{leancode}
\leansource{\texttt{sgldLaws\_\allowbreak{}probability} \quad AverageLMC/\allowbreak{}SGLDLaw.lean, lines 68–74}
\begin{leancode}
\leanline{0}{instance\ sgldLaws\_probability\ \{p\ n\ :\ ℕ\}\ [NeZero\ n]}
\leanline{4}{(f\ :\ Fin\ n\ →\ State\ p\ →\ ℝ)\ (anchor\ :\ State\ p)\ (h\ :\ ℝ)}
\leanline{4}{(μ\ :\ Measure\ (State\ p))\ [IsProbabilityMeasure\ μ]\ (k\ :\ ℕ)\ :}
\leanline{4}{IsProbabilityMeasure\ (sgldLaws\ f\ anchor\ h\ μ\ k)\ :=\ by}
\leanline{2}{induction\ k\ with}
\leanline{2}{|\ zero\ =>\ exact\ inferInstanceAs\ (IsProbabilityMeasure\ μ)}
\leanline{2}{|\ succ\ k\ ih\ =>\ exact\ sgldLawStep\_probability\ f\ anchor\ h\ \_}
\end{leancode}
\leanlabel{How the Lean proof follows the argument}
The components are a family \texttt{f : Fin n \ensuremath{\to} State p \ensuremath{\to} \ensuremath{\mathbb R}}; their sum \(f\) is \texttt{component\allowbreak{}Potential f}, the anchor \(\theta_*\) is \texttt{anchor}, \(g_i\) is \texttt{control\allowbreak{}Variate\allowbreak{}Estimator f anchor x i} and \(e_i\) is \texttt{control\allowbreak{}Variate\allowbreak{}Error f anchor x i} (all defined in the estimator section). The uniform law \(U_n\) is \texttt{uniform\allowbreak{}Index\allowbreak{}Law n}, the measure \texttt{(n : \ensuremath{\mathbb R}\ensuremath{\ge}0\ensuremath{\infty})\ensuremath{{}^{-1}} \ensuremath{\bullet} Measure.\allowbreak{}count} on \texttt{Fin n}, and \(\gamma_p\) is Mathlib's \texttt{std\allowbreak{}Gaussian (State p)}.



\begin{enumerate}


\item \texttt{sgld\allowbreak{}Proposal f anchor h x z i} is \(\Phi(x,z,i)\), with \texttt{Real.\allowbreak{}sqrt} (which returns \(0\) on negative numbers). \texttt{sgld\allowbreak{}Law\allowbreak{}Step f anchor h \ensuremath{\mu}} is the image \texttt{((\ensuremath{\mu}.\allowbreak{}prod (std\allowbreak{}Gaussian (State p))).\allowbreak{}prod (uniform\allowbreak{}Index\allowbreak{}Law n)).\allowbreak{}map …} of the product measure under \(\Phi\), and \texttt{sgld\allowbreak{}Laws} is defined by recursion on \texttt{k}, with \texttt{0 => \ensuremath{\mu}} and \texttt{k + 1 => sgld\allowbreak{}Law\allowbreak{}Step …}.

\item \texttt{sgld\allowbreak{}Proposal\_\allowbreak{}eq\_\allowbreak{}full\_\allowbreak{}gradient} is (1): the full-gradient part is \texttt{ULA.\allowbreak{}mala\allowbreak{}Proposal (component\allowbreak{}Potential f) h x z}, the map \(x-h\nabla f(x)+\sqrt{2h}z\) of the unadjusted Langevin library. After unfolding, \texttt{module} proves the vector identity.

\item \texttt{sgld\allowbreak{}Proposal\_\allowbreak{}measurable} rewrites with (1) and adds the measurability of \texttt{ULA.\allowbreak{}mala\allowbreak{}Proposal\_\allowbreak{}measurable} (composed with the projection) to that of the error, \texttt{control\allowbreak{}Variate\allowbreak{}Error\_\allowbreak{}product\_\allowbreak{}measurable}, multiplied by \(-h\).

\item \texttt{sgld\allowbreak{}Law\allowbreak{}Step\_\allowbreak{}is\allowbreak{}Probability} (with \texttt{hn : 0 < n}) and the instance \texttt{sgld\allowbreak{}Law\allowbreak{}Step\_\allowbreak{}probability} (with \texttt{[Ne\allowbreak{}Zero n]}, the same condition in typeclass form) use \texttt{uniform\allowbreak{}Index\allowbreak{}Law\_\allowbreak{}probability} and \texttt{infer\_\allowbreak{}instance}: products and images of probability measures are probability measures. In the Mathlib version used, the image of a probability measure under any map is a probability measure, so the unused hypotheses \texttt{\_\allowbreak{}hf}, \texttt{\_\allowbreak{}h\allowbreak{}Lip}, \texttt{\_\allowbreak{}hfull} of the first theorem are not needed. \texttt{sgld\allowbreak{}Laws\_\allowbreak{}probability} is the induction on \texttt{k}.

\end{enumerate}
\end{leanbox}


Earlier results used here: \cite[Section~3.34.1]{LangevinLean2026ULASC}; \cite[Section~3.36.5]{LangevinLean2026ULASC}.
\subsection{\texorpdfstring{Second moments are preserved}{Second moments are preserved}}
\label{report-step-AverageLMC-SGLDLaw-lean-2}


\textbf{Claim.} Let \(\nu\) be a probability measure on \(\mathbb R^p\) with finite second moment, \(\int\|x\|^2\nu(dx)<\infty\).



\begin{enumerate}


\item If \(F:\mathbb R^p\to\mathbb R\) has an \(M\)-Lipschitz gradient for some \(M\ge0\), then \(\nabla F\) is square integrable for \(\nu\).

\item Under the same hypothesis on \(F\) and for every \(h\in\mathbb R\), the map \((x,z)\mapsto x-h\nabla F(x)+\sqrt{2h}\,z\) is square integrable for \(\nu\otimes\gamma_p\).

\item If every \(f_i\) is differentiable with Lipschitz gradient and \(\nabla f\) is Lipschitz, then \(\nu^+\) has finite second moment.

\end{enumerate}



\textbf{Proof.} (1) The Lipschitz bound and the triangle inequality give \(\|\nabla F(x)\|\le\|\nabla F(0)\|+M\|x\|\), hence \(\|\nabla F(x)\|^2\le2\|\nabla F(0)\|^2+2M^2\|x\|^2\) by \((u+v)^2\le2u^2+2v^2\). The right-hand side is \(\nu\)-integrable because \(\nu\) is a probability measure with finite second moment. (Measurability holds because \(\nabla F\) is continuous.)




(2) Under \(\nu\otimes\gamma_p\), the coordinate \(x\) is square integrable by assumption, \(\nabla F(x)\) is square integrable by (1), and \(z\) is square integrable because \(\int\|z\|^2\gamma_p(dz)=p<\infty\). Square-integrable functions form a vector space, so the linear combination \(x-h\nabla F(x)+\sqrt{2h}z\) is square integrable.




(3) By the change-of-variables formula, \(\int\|y\|^2\nu^+(dy)=\int\|\Phi\|^2\,d(\nu\otimes\gamma_p\otimes U_n)\), so it suffices that \(\Phi\) is square integrable. By (1) of Step 1, \(\Phi\) is the full-gradient map, square integrable by (2) applied to \(F=f\), plus \(-h\,e_i(x)\). For each \(i\), \(e_i(x)=n(\nabla f_i(x)-\nabla f_i(\theta_*))-\sum_k\nabla f_k(x)\) is square integrable in \(x\) by (1) applied to the \(f_k\) (constants are square integrable for a probability measure), and since there are finitely many indices, \(\|e_i(x)\|^2\le\sum_{k=1}^n\|e_k(x)\|^2\) shows that \((x,z,i)\mapsto e_i(x)\) is square integrable for \(\nu\otimes\gamma_p\otimes U_n\). \(\square\)




These facts are what make the transport distances in the next step finite; no moment assumption beyond that of \(\nu\) is needed.

\begin{leanbox}{Lean: \texttt{full\_\allowbreak{}gradient\_\allowbreak{}proposal\_\allowbreak{}memLp}, \texttt{gradient\_\allowbreak{}memLp}, \texttt{sgldLawStep\_\allowbreak{}memLp}}
\leansource{\texttt{full\_\allowbreak{}gradient\_\allowbreak{}proposal\_\allowbreak{}memLp} \quad AverageLMC/\allowbreak{}SGLDLaw.lean, lines 78–93}
\begin{leancode}
\leanline{0}{theorem\ full\_gradient\_proposal\_memLp\ \{p\ :\ ℕ\}\ (F\ :\ State\ p\ →\ ℝ)\ (M\ :\ ℝ≥0)}
\leanline{4}{(hLip\ :\ LipschitzWith\ M\ (gradient\ F))\ (h\ :\ ℝ)}
\leanline{4}{(μ\ :\ Measure\ (State\ p))\ [IsProbabilityMeasure\ μ]}
\leanline{4}{(hμ\ :\ MemLp\ (fun\ x\ :\ State\ p\ ↦\ x)\ 2\ μ)\ :}
\leanline{4}{MemLp\ (fun\ q\ :\ State\ p\ ×\ State\ p\ ↦\ ULA.malaProposal\ F\ h\ q.1\ q.2)}
\leanline{6}{2\ (μ.prod\ (stdGaussian\ (State\ p)))}
\end{leancode}
\leansource{\texttt{gradient\_\allowbreak{}memLp} \quad AverageLMC/\allowbreak{}SGLDLaw.lean, lines 144–151}
\begin{leancode}
\leanline{0}{theorem\ gradient\_memLp\ \{p\ :\ ℕ\}\ (F\ :\ State\ p\ →\ ℝ)\ (M\ :\ ℝ≥0)}
\leanline{4}{(hLip\ :\ LipschitzWith\ M\ (gradient\ F))}
\leanline{4}{(μ\ :\ Measure\ (State\ p))\ [IsProbabilityMeasure\ μ]}
\leanline{4}{(hμ\ :\ MemLp\ (fun\ x\ :\ State\ p\ ↦\ x)\ 2\ μ)\ :\ MemLp\ (gradient\ F)\ 2\ μ}
\end{leancode}
\leansource{\texttt{sgldLawStep\_\allowbreak{}memLp} \quad AverageLMC/\allowbreak{}SGLDLaw.lean, lines 154–168}
\begin{leancode}
\leanline{0}{theorem\ sgldLawStep\_memLp\ \{p\ n\ :\ ℕ\}}
\leanline{4}{(f\ :\ Fin\ n\ →\ State\ p\ →\ ℝ)\ (hf\ :\ ∀\ i,\ Differentiable\ ℝ\ (f\ i))}
\leanline{4}{(L\ :\ Fin\ n\ →\ ℝ≥0)\ (hLip\ :\ ∀\ i,\ LipschitzWith\ (L\ i)\ (gradient\ (f\ i)))}
\leanline{4}{(M\ :\ ℝ≥0)\ (hfull\ :\ LipschitzWith\ M\ (gradient\ (componentPotential\ f)))}
\leanline{4}{(anchor\ :\ State\ p)\ (h\ :\ ℝ)\ (μ\ :\ Measure\ (State\ p))\ [IsProbabilityMeasure\ μ]}
\leanline{4}{(hμ\ :\ MemLp\ (fun\ x\ :\ State\ p\ ↦\ x)\ 2\ μ)\ :}
\leanline{4}{MemLp\ (fun\ x\ :\ State\ p\ ↦\ x)\ 2\ (sgldLawStep\ f\ anchor\ h\ μ)}
\end{leancode}
\leanlabel{How the Lean proof follows the argument}
Square integrability is \texttt{Mem\allowbreak{}Lp g 2 \ensuremath{\mu}}, and the finite second moment of a law \texttt{\ensuremath{\mu}} is \texttt{Mem\allowbreak{}Lp (fun x : State p \ensuremath{\mapsto} x) 2 \ensuremath{\mu}}.



\begin{enumerate}


\item \texttt{gradient\_\allowbreak{}mem\allowbreak{}Lp} is claim (1). It applies \texttt{ULA.\allowbreak{}lp\allowbreak{}Drift\_\allowbreak{}mem\allowbreak{}Lp}, which states that a Lipschitz vector field maps an \(L^2\) random variable to an \(L^2\) random variable, to the class \texttt{h\ensuremath{\mu}.\allowbreak{}to\allowbreak{}Lp \_\allowbreak{}} of the identity, and passes back to the function itself with \texttt{ae\_\allowbreak{}eq} and \texttt{coe\allowbreak{}Fn\_\allowbreak{}to\allowbreak{}Lp}.

\item \texttt{full\_\allowbreak{}gradient\_\allowbreak{}proposal\_\allowbreak{}mem\allowbreak{}Lp} is claim (2): \texttt{h\allowbreak{}X} is the square integrability of the first coordinate under the product (\texttt{comp\_\allowbreak{}fst}), \texttt{hg} and \texttt{hg'} that of \(\nabla F(x)\) by the same argument, and \texttt{h\allowbreak{}Z} that of the Gaussian coordinate, from \texttt{Is\allowbreak{}Gaussian.\allowbreak{}mem\allowbreak{}Lp\_\allowbreak{}id} composed with the second projection (\texttt{comp\_\allowbreak{}snd}). The conclusion combines them with \texttt{sub}, \texttt{add} and \texttt{const\_\allowbreak{}smul}.

\item \texttt{sgld\allowbreak{}Law\allowbreak{}Step\_\allowbreak{}mem\allowbreak{}Lp} is claim (3). \texttt{mem\allowbreak{}Lp\_\allowbreak{}map\_\allowbreak{}measure\_\allowbreak{}iff}, which needs the measurability \texttt{sgld\allowbreak{}Proposal\_\allowbreak{}measurable}, transfers the question to the update map; \texttt{hd} is the square integrability of the error (\texttt{control\allowbreak{}Variate\allowbreak{}Error\_\allowbreak{}product\_\allowbreak{}mem\allowbreak{}Lp}, scaled by \(-h\)) and \texttt{hg} that of the full-gradient map, and the final \texttt{simpa} uses identity (1).

\end{enumerate}
\end{leanbox}


Earlier results used here: \cite[Section~3.13.4]{LangevinLean2026ULASC}; \cite[Section~3.34.1]{LangevinLean2026ULASC}.
\subsection{\texorpdfstring{Variance addition for one SGLD step}{Variance addition for one SGLD step}}
\label{report-step-AverageLMC-SGLDLaw-lean-3}


\textbf{Claim.} Let \(n\ge1\), let every \(f_i\) be differentiable with Lipschitz gradient, let \(\nabla f\) be \(M\)-Lipschitz, and let \(\nabla f(\theta_*)=0\). Let \(h\in\mathbb R\), and let \(\nu\) and \(\pi\) be probability measures on \(\mathbb R^p\) with finite second moments. Let \(\nu^+_{\rm LMC}\) be the law of \(x-h\nabla f(x)+\sqrt{2h}\,z\) under \(\nu\otimes\gamma_p\). Then


\[\mathsf W_2(\nu^+,\pi)^2\le\mathsf W_2(\nu^+_{\rm LMC},\pi)^2+h^2\int\|e_i(x)\|^2\,(\nu\otimes\gamma_p\otimes U_n)(dx,dz,di),\qquad(2)\]


and


\[\int\|e_i(x)\|^2\,d(\nu\otimes\gamma_p\otimes U_n)=\int\|g_i(x)\|^2\,(\nu\otimes U_n)(dx,di)-\int\|\nabla f(x)\|^2\,\nu(dx).\qquad(3)\]


\textbf{Proof.} \emph{Pointwise facts for fixed \(x\).} Since each \(f_i\) is differentiable, \(\nabla f=\sum_i\nabla f_i\). Averaging over the index and using \(\nabla f(\theta_*)=0\),


\[\frac1n\sum_{i=1}^ng_i(x)=\sum_{i=1}^n\bigl(\nabla f_i(x)-\nabla f_i(\theta_*)\bigr)=\nabla f(x)-\nabla f(\theta_*)=\nabla f(x),\]


so the estimator is unbiased and \(\frac1n\sum_ie_i(x)=0\). Expanding \(\|g_i(x)-\nabla f(x)\|^2\) and averaging,


\[\begin{aligned}
\frac1n\sum_i\|e_i(x)\|^2
&=\frac1n\sum_i\|g_i(x)\|^2\\
&\quad-2\Bigl\langle\frac1n\sum_ig_i(x),\nabla f(x)\Bigr\rangle
  +\|\nabla f(x)\|^2\\
&=\frac1n\sum_i\|g_i(x)\|^2-\|\nabla f(x)\|^2.
\qquad(4)
\end{aligned}\]


\emph{Proof of (2).} Apply the centred-perturbation inequality of the earlier section with \(S=\mathbb R^p\times\mathbb R^p\), \(P=\nu\otimes\gamma_p\), index space \(\{1,\dots,n\}\) with \(Q=U_n\), \(F(x,z)=x-h\nabla f(x)+\sqrt{2h}z\) and \(D((x,z),i)=-h\,e_i(x)\). Its hypotheses hold: \(F\) is continuous, hence measurable, and square integrable by Step 2; \(D\) is jointly measurable and square integrable by Steps 1 and 2; for fixed \((x,z)\), \(D((x,z),\cdot)\) is integrable (the index set is finite) and its mean \(-h\cdot\frac1n\sum_ie_i(x)\) is \(0\); and \(\pi\) is a probability measure with finite second moment. By (1) of Step 1, the law of \(F+D\) under \(P\otimes Q\) is \(\nu^+\), and the law of \(F\) under \(P\) is \(\nu^+_{\rm LMC}\). Finally \(\|D\|^2=h^2\|e_i(x)\|^2\), and the constant \(h^2\) comes out of the integral. This gives (2).




\emph{Proof of (3).} The integrand \(\|e_i(x)\|^2\) is integrable by Step 2. By Fubini's theorem we may integrate first over the index; by (4) this gives the function \(x\mapsto\int\|g_i(x)\|^2U_n(di)-\|\nabla f(x)\|^2\), which does not depend on \(z\). Integrating it against the probability measure \(\gamma_p\) leaves it unchanged, and integrating over \(\nu\), splitting the difference (both terms are integrable by Step 2) and using Fubini's theorem once more for \(\int\!\int\|g_i(x)\|^2U_n(di)\,\nu(dx)\) gives (3). \(\square\)




The hypothesis \(\nabla f(\theta_*)=0\) is used only for the unbiasedness of the estimator, and hence for the vanishing of the cross term. The inequality is between real numbers; Step 2 guarantees that all distances are finite.

\begin{leanbox}{Lean: \texttt{sgld\_\allowbreak{}variance\_\allowbreak{}addition}, \texttt{sgld\_\allowbreak{}error\_\allowbreak{}variance\_\allowbreak{}identity}}
\leansource{\texttt{sgld\_\allowbreak{}variance\_\allowbreak{}addition} \quad AverageLMC/\allowbreak{}SGLDLaw.lean, lines 98–141}
\begin{leancode}
\leanline{0}{theorem\ sgld\_variance\_addition\ \{p\ n\ :\ ℕ\}\ (hn\ :\ 0\ <\ n)}
\leanline{4}{(f\ :\ Fin\ n\ →\ State\ p\ →\ ℝ)\ (hf\ :\ ∀\ i,\ Differentiable\ ℝ\ (f\ i))}
\leanline{4}{(L\ :\ Fin\ n\ →\ ℝ≥0)\ (hLip\ :\ ∀\ i,\ LipschitzWith\ (L\ i)\ (gradient\ (f\ i)))}
\leanline{4}{(M\ :\ ℝ≥0)\ (hfull\ :\ LipschitzWith\ M\ (gradient\ (componentPotential\ f)))}
\leanline{4}{(anchor\ :\ State\ p)\ (hstar\ :\ gradient\ (componentPotential\ f)\ anchor\ =\ 0)}
\leanline{4}{(h\ :\ ℝ)\ (μ\ π\ :\ Measure\ (State\ p))\ [IsProbabilityMeasure\ μ]\ [IsProbabilityMeasure\ π]}
\leanline{4}{(hμ\ :\ MemLp\ (fun\ x\ :\ State\ p\ ↦\ x)\ 2\ μ)\ (hπ\ :\ MemLp\ (fun\ x\ :\ State\ p\ ↦\ x)\ 2\ π)\ :}
\leanline{4}{ULA.wasserstein2Real\ (sgldLawStep\ f\ anchor\ h\ μ)\ π\ \textasciicircum{}\ 2\ ≤}
\leanline{6}{ULA.wasserstein2Real\ (ULA.ulaLawStep\ (componentPotential\ f)\ h\ μ)\ π\ \textasciicircum{}\ 2\ +}
\leanline{8}{h\ \textasciicircum{}\ 2\ *\ ∫\ q\ :\ (State\ p\ ×\ State\ p)\ ×\ Fin\ n,}
\leanline{10}{‖controlVariateError\ f\ anchor\ q.1.1\ q.2‖\ \textasciicircum{}\ 2}
\leanline{12}{∂(μ.prod\ (stdGaussian\ (State\ p))).prod\ (uniformIndexLaw\ n)}
\end{leancode}
\leansource{\texttt{sgld\_\allowbreak{}error\_\allowbreak{}variance\_\allowbreak{}identity} \quad AverageLMC/\allowbreak{}SGLDLaw.lean, lines 172–200}
\begin{leancode}
\leanline{0}{theorem\ sgld\_error\_variance\_identity\ \{p\ n\ :\ ℕ\}\ (hn\ :\ 0\ <\ n)}
\leanline{4}{(f\ :\ Fin\ n\ →\ State\ p\ →\ ℝ)\ (hf\ :\ ∀\ i,\ Differentiable\ ℝ\ (f\ i))}
\leanline{4}{(L\ :\ Fin\ n\ →\ ℝ≥0)\ (hLip\ :\ ∀\ i,\ LipschitzWith\ (L\ i)\ (gradient\ (f\ i)))}
\leanline{4}{(M\ :\ ℝ≥0)\ (hfull\ :\ LipschitzWith\ M\ (gradient\ (componentPotential\ f)))}
\leanline{4}{(anchor\ :\ State\ p)\ (hstar\ :\ gradient\ (componentPotential\ f)\ anchor\ =\ 0)}
\leanline{4}{(μ\ :\ Measure\ (State\ p))\ [IsProbabilityMeasure\ μ]}
\leanline{4}{(hμ\ :\ MemLp\ (fun\ x\ :\ State\ p\ ↦\ x)\ 2\ μ)\ :}
\leanline{4}{(∫\ q\ :\ (State\ p\ ×\ State\ p)\ ×\ Fin\ n,}
\leanline{6}{‖controlVariateError\ f\ anchor\ q.1.1\ q.2‖\ \textasciicircum{}\ 2}
\leanline{8}{∂(μ.prod\ (stdGaussian\ (State\ p))).prod\ (uniformIndexLaw\ n))\ =}
\leanline{4}{(∫\ q\ :\ State\ p\ ×\ Fin\ n,\ ‖controlVariateEstimator\ f\ anchor\ q.1\ q.2‖\ \textasciicircum{}\ 2}
\leanline{6}{∂μ.prod\ (uniformIndexLaw\ n))\ -\ ∫\ x,\ ‖gradient\ (componentPotential\ f)\ x‖\ \textasciicircum{}\ 2\ ∂μ}
\end{leancode}
\leanlabel{How the Lean proof follows the argument}
Both theorems take \texttt{hn : 0 < n}, \texttt{hf : \ensuremath{\forall} i, Differentiable \ensuremath{\mathbb R} (f i)}, Lipschitz constants \texttt{L : Fin n \ensuremath{\to} \ensuremath{\mathbb R}\ensuremath{\ge}0} with \texttt{h\allowbreak{}Lip : \ensuremath{\forall} i, Lipschitz\allowbreak{}With (L i) (gradient (f i))}, \texttt{hfull : Lipschitz\allowbreak{}With M (gradient (component\allowbreak{}Potential f))} and \texttt{hstar : gradient (component\allowbreak{}Potential f) anchor = 0}.



\begin{enumerate}


\item In \texttt{sgld\_\allowbreak{}variance\_\allowbreak{}addition}, \texttt{P} is the product \(\nu\otimes\gamma_p\), \texttt{D} the perturbation, \texttt{h\allowbreak{}D} and \texttt{h\allowbreak{}D2} its measurability and square integrability, and \texttt{hmean} the vanishing mean: \texttt{integral\_\allowbreak{}smul} extracts \(-h\) and \texttt{control\allowbreak{}Variate\allowbreak{}Error\_\allowbreak{}centered} (the identity \(\frac1n\sum_ie_i(x)=0\), proved in the estimator section) gives \(0\). \texttt{hw} is \texttt{wasserstein2\_\allowbreak{}sq\_\allowbreak{}centered\_\allowbreak{}perturbation} with the full-gradient map, its measurability \texttt{ULA.\allowbreak{}mala\allowbreak{}Proposal\_\allowbreak{}measurable} and square integrability \texttt{full\_\allowbreak{}gradient\_\allowbreak{}proposal\_\allowbreak{}mem\allowbreak{}Lp}; the integrability over the finite index set is \texttt{Integrable.\allowbreak{}of\_\allowbreak{}finite}. \texttt{hmap} identifies the law of \(F+D\) with \texttt{sgld\allowbreak{}Law\allowbreak{}Step f anchor h \ensuremath{\mu}} through \texttt{sgld\allowbreak{}Proposal\_\allowbreak{}eq\_\allowbreak{}full\_\allowbreak{}gradient}, and \texttt{hvariance} pulls out \(h^2\) with \texttt{norm\_\allowbreak{}smul}, \texttt{mul\_\allowbreak{}pow}, \texttt{sq\_\allowbreak{}abs}, \texttt{neg\_\allowbreak{}sq} and \texttt{integral\_\allowbreak{}const\_\allowbreak{}mul}. The law \(\nu^+_{\rm LMC}\) is \texttt{ULA.\allowbreak{}ula\allowbreak{}Law\allowbreak{}Step (component\allowbreak{}Potential f) h \ensuremath{\mu}}.

\item In \texttt{sgld\_\allowbreak{}error\_\allowbreak{}variance\_\allowbreak{}identity}, \texttt{he}, \texttt{hs} and \texttt{hg} are the integrability of \(\|e_i(x)\|^2\), \(\|g_i(x)\|^2\) and \(\|\nabla f(x)\|^2\). \texttt{rw [integral\_\allowbreak{}prod \_\allowbreak{} he]} is Fubini's theorem; \texttt{control\allowbreak{}Variate\allowbreak{}Error\_\allowbreak{}second\_\allowbreak{}moment} is the pointwise identity (4); \texttt{hi} is the integrability of the resulting function of \(x\); a second \texttt{integral\_\allowbreak{}prod} and \texttt{integral\_\allowbreak{}const} with \texttt{prob\allowbreak{}Real\_\allowbreak{}univ} integrate out the Gaussian coordinate; and \texttt{integral\_\allowbreak{}sub} with a last \texttt{integral\_\allowbreak{}prod} gives (3).

\end{enumerate}
\end{leanbox}


Earlier results used here: \cite[Section~3.17.4]{LangevinLean2026ULASC}; \cite[Section~3.34.1]{LangevinLean2026ULASC}; \cite[Section~3.36.5]{LangevinLean2026ULASC}; \cite[Section~3.36.7]{LangevinLean2026ULASC}.
\subsection{\texorpdfstring{The second moment of the estimator}{The second moment of the estimator}}
\label{report-step-AverageLMC-SGLDLaw-lean-4}


\textbf{Claim.} Let \(n\ge1\), let every \(\nabla f_i\) be Lipschitz, let \(\theta_*\in\mathbb R^p\), and let \(\nu\) be a probability measure on \(\mathbb R^p\) with finite second moment. Then


\[\int\|g_i(x)\|^2\,(\nu\otimes U_n)(dx,di)=n\sum_{i=1}^n\int\|\nabla f_i(x)-\nabla f_i(\theta_*)\|^2\,\nu(dx).\]


\textbf{Proof.} For fixed \(x\),


\[\int\|g_i(x)\|^2\,U_n(di)=\frac1n\sum_{i=1}^n n^2\|\nabla f_i(x)-\nabla f_i(\theta_*)\|^2=n\sum_{i=1}^n\|\nabla f_i(x)-\nabla f_i(\theta_*)\|^2.\]


The function \((x,i)\mapsto\|g_i(x)\|^2\) is integrable for \(\nu\otimes U_n\) (Step 2), so by Fubini's theorem its integral equals the \(\nu\)-integral of the last display. Each \(x\mapsto\|\nabla f_i(x)-\nabla f_i(\theta_*)\|^2\) is \(\nu\)-integrable by (1) of Step 2, so the integral of the finite sum is the sum of the integrals, and the constant \(n\) comes out of the integral. \(\square\)




With (3) of Step 3, \(h^2\mathbb E\|e_I(\vartheta)\|^2=nh^2\sum_i\|\nabla f_i(\vartheta)-\nabla f_i(\theta_*)\|_{L_2}^2-h^2\|\nabla f(\vartheta)\|_{L_2}^2\); this is the last term of the first claim of Proposition 3 of the paper, together with the extra \(-h^2\|\nabla f(\vartheta)\|_{L_2}^2\) that changes the coefficient \(h^2(hM_\infty-1)\) of Theorem 1 into \(h^2(hM_\infty-2)\).

\begin{leanbox}{Lean: \texttt{integral\_\allowbreak{}controlVariateEstimator\_\allowbreak{}sq}}
\leansource{\texttt{integral\_\allowbreak{}controlVariateEstimator\_\allowbreak{}sq} \quad AverageLMC/\allowbreak{}SGLDLaw.lean, lines 204–220}
\begin{leancode}
\leanline{0}{theorem\ integral\_controlVariateEstimator\_sq\ \{p\ n\ :\ ℕ\}\ (hn\ :\ 0\ <\ n)}
\leanline{4}{(f\ :\ Fin\ n\ →\ State\ p\ →\ ℝ)\ (L\ :\ Fin\ n\ →\ ℝ≥0)}
\leanline{4}{(hLip\ :\ ∀\ i,\ LipschitzWith\ (L\ i)\ (gradient\ (f\ i)))}
\leanline{4}{(anchor\ :\ State\ p)\ (μ\ :\ Measure\ (State\ p))\ [IsProbabilityMeasure\ μ]}
\leanline{4}{(hμ\ :\ MemLp\ (fun\ x\ :\ State\ p\ ↦\ x)\ 2\ μ)\ :}
\leanline{4}{(∫\ q\ :\ State\ p\ ×\ Fin\ n,\ ‖controlVariateEstimator\ f\ anchor\ q.1\ q.2‖\ \textasciicircum{}\ 2}
\leanline{6}{∂μ.prod\ (uniformIndexLaw\ n))\ =}
\leanline{6}{(n\ :\ ℝ)\ *\ ∑\ i,\ ∫\ x,\ ‖gradient\ (f\ i)\ x\ -\ gradient\ (f\ i)\ anchor‖\ \textasciicircum{}\ 2\ ∂μ}
\end{leancode}
\leanlabel{How the Lean proof follows the argument}
\texttt{integral\_\allowbreak{}control\allowbreak{}Variate\allowbreak{}Estimator\_\allowbreak{}sq} takes \texttt{hn : 0 < n}, the Lipschitz hypotheses \texttt{h\allowbreak{}Lip} and the second moment \texttt{h\ensuremath{\mu}}; no differentiability is needed. \texttt{he} is the integrability of \(\|g_i(x)\|^2\) (from \texttt{control\allowbreak{}Variate\allowbreak{}Estimator\_\allowbreak{}product\_\allowbreak{}mem\allowbreak{}Lp} applied to the identity), and \texttt{hi i} the integrability of each \(\|\nabla f_i(x)-\nabla f_i(\theta_*)\|^2\) (from \texttt{gradient\_\allowbreak{}mem\allowbreak{}Lp} and \texttt{mem\allowbreak{}Lp\_\allowbreak{}const}). \texttt{integral\_\allowbreak{}prod} is Fubini's theorem, \texttt{control\allowbreak{}Variate\allowbreak{}Estimator\_\allowbreak{}second\_\allowbreak{}moment} is the pointwise identity, and \texttt{integral\_\allowbreak{}const\_\allowbreak{}mul} with \texttt{integral\_\allowbreak{}finset\allowbreak{}Sum} exchanges the integral with the constant and the finite sum.
\end{leanbox}
\section{\texorpdfstring{Exact scaling of coordinate smoothness constants}{Exact scaling of coordinate smoothness constants}}
\label{proof-module-AverageLMC-SmoothnessScaling-lean}


This module proves that the literal supremum constants commute with nonnegative rescaling of the potential. In particular the barred average-coordinate constant of the SGLD average potential \(f/n\) is exactly the full-potential constant divided by \(n\). The proof does not replace either quantity by an arbitrary upper bound.



For a differentiable potential with Lipschitz gradient, nonnegative multiplication commutes with the literal coordinate Taylor suprema:

\[c\ge0\quad\Longrightarrow\quad M_j(cf)=cM_j(f),\qquad M_{\rm av}(cf)=cM_{\rm av}(f).\]

In particular the SGLD barred constant is exactly the constant of the average potential:

\[n\ge1\quad\Longrightarrow\quad \bar M_{\rm av}=M_{\rm av}(f/n)=\frac{M_{\rm av}(f)}{n}.\]\subsection{\texorpdfstring{Rescale the actual gradient and coordinate Taylor quotients}{Rescale the actual gradient and coordinate Taylor quotients}}
\label{report-step-AverageLMC-SmoothnessScaling-lean-1}


The Frechet derivative of \(c\) times \(f\) is \(c\) times the derivative of \(f\). The Riesz representation defining the gradient commutes with this real scalar, so the true gradient is \(c\) times the original gradient. Substitution in the exact coordinate Taylor quotient factors out \(c\). Consequently the entire quotient set of the rescaled potential is the image of the original set under multiplication by \(c\).

\begin{leanbox}{Lean: \texttt{gradient\_\allowbreak{}const\_\allowbreak{}mul\_\allowbreak{}potential}, \texttt{coordinateTaylorQuotient\_\allowbreak{}const\_\allowbreak{}mul}, \texttt{coordinateTaylorQuotients\_\allowbreak{}const\_\allowbreak{}mul}}
\leansource{\texttt{gradient\_\allowbreak{}const\_\allowbreak{}mul\_\allowbreak{}potential} \quad AverageLMC/\allowbreak{}SmoothnessScaling.lean, lines 18–21}
\begin{leancode}
\leanline{0}{theorem\ gradient\_const\_mul\_potential\ \{p\ :\ ℕ\}\ (f\ :\ State\ p\ →\ ℝ)}
\leanline{4}{(hf\ :\ Differentiable\ ℝ\ f)\ (c\ :\ ℝ)\ (x\ :\ State\ p)\ :}
\leanline{4}{gradient\ (fun\ y\ ↦\ c\ *\ f\ y)\ x\ =\ c\ •\ gradient\ f\ x}
\end{leancode}
\leansource{\texttt{coordinateTaylorQuotient\_\allowbreak{}const\_\allowbreak{}mul} \quad AverageLMC/\allowbreak{}SmoothnessScaling.lean, lines 23–30}
\begin{leancode}
\leanline{0}{theorem\ coordinateTaylorQuotient\_const\_mul\ \{p\ :\ ℕ\}\ (f\ :\ State\ p\ →\ ℝ)}
\leanline{4}{(hf\ :\ Differentiable\ ℝ\ f)\ (c\ :\ ℝ)\ (j\ :\ Fin\ p)\ (x\ :\ State\ p)\ (a\ :\ ℝ)\ :}
\leanline{4}{coordinateTaylorQuotient\ (fun\ y\ ↦\ c\ *\ f\ y)\ j\ x\ a\ =}
\leanline{6}{c\ *\ coordinateTaylorQuotient\ f\ j\ x\ a}
\end{leancode}
\leansource{\texttt{coordinateTaylorQuotients\_\allowbreak{}const\_\allowbreak{}mul} \quad AverageLMC/\allowbreak{}SmoothnessScaling.lean, lines 32–42}
\begin{leancode}
\leanline{0}{theorem\ coordinateTaylorQuotients\_const\_mul\ \{p\ :\ ℕ\}\ (f\ :\ State\ p\ →\ ℝ)}
\leanline{4}{(hf\ :\ Differentiable\ ℝ\ f)\ (c\ :\ ℝ)\ (j\ :\ Fin\ p)\ :}
\leanline{4}{coordinateTaylorQuotients\ (fun\ y\ ↦\ c\ *\ f\ y)\ j\ =}
\leanline{6}{(fun\ q\ ↦\ c\ *\ q)\ ''\ coordinateTaylorQuotients\ f\ j}
\end{leancode}
\leanlabel{How the Lean proof follows the argument}
\texttt{gradient\_\allowbreak{}const\_\allowbreak{}mul\_\allowbreak{}potential} gives the formal counterpart of this step. Its remaining equalities are discharged by the simplifications displayed in the proof.




\texttt{coordinate\allowbreak{}Taylor\allowbreak{}Quotient\_\allowbreak{}const\_\allowbreak{}mul} gives the formal counterpart of this step. Its remaining equalities are discharged by the simplifications displayed in the proof.




\texttt{coordinate\allowbreak{}Taylor\allowbreak{}Quotients\_\allowbreak{}const\_\allowbreak{}mul} gives the formal counterpart of this step. Named proof applications include \texttt{\ensuremath{\langle}coordinate\allowbreak{}Taylor\allowbreak{}Quotient}, \texttt{\ensuremath{\langle}x}; the adjacent proof shows their hypotheses and order of use.
\end{leanbox}
\subsection{\texorpdfstring{Transport the least upper bound through nonnegative multiplication}{Transport the least upper bound through nonnegative multiplication}}
\label{report-step-AverageLMC-SmoothnessScaling-lean-2}


The Lipschitz-gradient assumption supplies an upper bound on the original Taylor-quotient set, and that set is nonempty. Its supremum therefore is a genuine least upper bound. Nonnegative multiplication transports that least upper bound to the image set, including scalar zero. This proves exact homogeneity of every coordinate supremum. A finite sum and division by the dimension prove the same identity for \(M_{\rm av}\).

\begin{leanbox}{Lean: \texttt{coordinateSmoothness\_\allowbreak{}const\_\allowbreak{}mul}, \texttt{averageSmoothness\_\allowbreak{}const\_\allowbreak{}mul}}
\leansource{\texttt{coordinateSmoothness\_\allowbreak{}const\_\allowbreak{}mul} \quad AverageLMC/\allowbreak{}SmoothnessScaling.lean, lines 45–53}
\begin{leancode}
\leanline{0}{theorem\ coordinateSmoothness\_const\_mul\ \{p\ :\ ℕ\}\ (f\ :\ State\ p\ →\ ℝ)}
\leanline{4}{(hf\ :\ Differentiable\ ℝ\ f)\ (M\ :\ ℝ≥0)\ (hLip\ :\ LipschitzWith\ M\ (gradient\ f))}
\leanline{4}{(c\ :\ ℝ)\ (hc\ :\ 0\ ≤\ c)\ (j\ :\ Fin\ p)\ :}
\leanline{4}{coordinateSmoothness\ (fun\ y\ ↦\ c\ *\ f\ y)\ j\ =\ c\ *\ coordinateSmoothness\ f\ j}
\end{leancode}
\leansource{\texttt{averageSmoothness\_\allowbreak{}const\_\allowbreak{}mul} \quad AverageLMC/\allowbreak{}SmoothnessScaling.lean, lines 56–63}
\begin{leancode}
\leanline{0}{theorem\ averageSmoothness\_const\_mul\ \{p\ :\ ℕ\}\ (f\ :\ State\ p\ →\ ℝ)}
\leanline{4}{(hf\ :\ Differentiable\ ℝ\ f)\ (M\ :\ ℝ≥0)\ (hLip\ :\ LipschitzWith\ M\ (gradient\ f))}
\leanline{4}{(c\ :\ ℝ)\ (hc\ :\ 0\ ≤\ c)\ :}
\leanline{4}{averageSmoothness\ (fun\ y\ ↦\ c\ *\ f\ y)\ =\ c\ *\ averageSmoothness\ f}
\end{leancode}
\leanlabel{How the Lean proof follows the argument}
\texttt{coordinate\allowbreak{}Smoothness\_\allowbreak{}const\_\allowbreak{}mul} gives the formal counterpart of this step.




\texttt{average\allowbreak{}Smoothness\_\allowbreak{}const\_\allowbreak{}mul} gives the formal counterpart of this step. The final step is the algebraic normalization or ordered arithmetic displayed in the proof.
\end{leanbox}
\subsection{\texorpdfstring{Identify the paper's barred average-potential constant}{Identify the paper's barred average-potential constant}}
\label{report-step-AverageLMC-SmoothnessScaling-lean-3}


Choose the scalar to be the reciprocal of the finite component count. It is nonnegative, and real division is multiplication by this reciprocal. The proved homogeneity yields \(M_{\mathrm{av}}(x\mapsto f(x)/n)=M_{\mathrm{av}}(f)/n\). Thus the source-facing SGLD theorem can state its discretization term with the literal constant of the average potential.

\begin{leanbox}{Lean: \texttt{averageSmoothness\_\allowbreak{}div\_\allowbreak{}nat}}
\leansource{\texttt{averageSmoothness\_\allowbreak{}div\_\allowbreak{}nat} \quad AverageLMC/\allowbreak{}SmoothnessScaling.lean, lines 66–71}
\begin{leancode}
\leanline{0}{theorem\ averageSmoothness\_div\_nat\ \{p\ :\ ℕ\}\ (f\ :\ State\ p\ →\ ℝ)}
\leanline{4}{(hf\ :\ Differentiable\ ℝ\ f)\ (M\ :\ ℝ≥0)\ (hLip\ :\ LipschitzWith\ M\ (gradient\ f))}
\leanline{4}{(n\ :\ ℕ)\ :}
\leanline{4}{averageSmoothness\ (fun\ x\ ↦\ f\ x\ /\ (n\ :\ ℝ))\ =\ averageSmoothness\ f\ /\ n}
\end{leancode}
\leanlabel{How the Lean proof follows the argument}
\texttt{average\allowbreak{}Smoothness\_\allowbreak{}div\_\allowbreak{}nat} gives the formal counterpart of this step. Its remaining equalities are discharged by the simplifications displayed in the proof.
\end{leanbox}
\section{\texorpdfstring{The scalar SGLD recurrence}{The scalar SGLD recurrence}}
\label{proof-module-AverageLMC-SGLDRecurrence-lean}


This section carries out the scalar part of the constant-step SGLD theorem. The letters denote real numbers: \(m>0\) is the strong convexity constant of the full potential \(f=\sum_if_i\), \(R=nM_{\rm rms}\), \(A\ge0\) is the average coordinate smoothness of \(f\), \(h\ge0\) the step and \(p\ge0\) the dimension. Under the step condition \(3R^2h\le m\) of the paper (equivalent to \(h\le\bar m/(3nM_{\rm rms}^2)\) with \(m=n\bar m\)) and \(m^2\le R^2\), we prove


\[\frac{1-mh}{1+mh}+2R^2h^2\le\frac1{1+mh},\qquad\Bigl(\frac{A+m}{1+mh}+\frac{2R^2}m\Bigr)h^2p\le\frac{(A+2m+2R^2/m)\,h^2p}{1+mh},\]


and deduce that every nonnegative sequence satisfying the one-step SGLD recurrence obeys


\[r_k\le(1+mh)^{-k}r_0+\frac{(A+2m+2R^2/m)\,hp}{m}.\]


With \(m=n\bar m\), \(R=nM_{\rm rms}\) and \(A=n\bar M_{\rm av}\), the last term is \((\bar M_{\rm av}+2\bar m+2M_{\rm rms}^2/\bar m)\,hp/\bar m\), the persistent error of the SGLD theorem. The recurrence itself is proved for the SGLD laws in a later section.

\subsection{\texorpdfstring{Comparison of the contraction and forcing coefficients}{Comparison of the contraction and forcing coefficients}}
\label{report-step-AverageLMC-SGLDRecurrence-lean-1}


\textbf{Claim.} Let \(m>0\), \(h\ge0\) and \(R\in\mathbb R\) with \(3R^2h\le m\).



\begin{enumerate}


\item If \(m^2\le R^2\), then \(mh\le\frac13\).

\item If \(m^2\le R^2\), then \(\displaystyle\frac{1-mh}{1+mh}+2R^2h^2\le\frac1{1+mh}\).

\item For all \(A\in\mathbb R\) and \(p\ge0\),
\[\Bigl(\frac{A+m}{1+mh}+\frac{2R^2}m\Bigr)h^2p\le\frac{(A+2m+2R^2/m)\,h^2p}{1+mh}.\]

\end{enumerate}



\textbf{Proof.} Write \(t=mh\ge0\); then \(1+t>0\).




(1) Multiplying \(m^2\le R^2\) by \(h\ge0\) gives \(m^2h\le R^2h\), so \(3m^2h\le3R^2h\le m\), that is, \(m(3mh-1)\le0\). Since \(m>0\), this forces \(3mh-1\le0\).




(2) Multiplying by \(1+t>0\), the claim is equivalent to \((1-t)+2R^2h^2(1+t)\le1\), that is, to \(2R^2h^2(1+t)\le t\). Multiplying \(3R^2h\le m\) by \(h\ge0\) gives \(3R^2h^2\le t\). By (1), \(1+t\le\frac43\), and \(2R^2h^2\ge0\), so


\[2R^2h^2(1+t)\le\frac43\cdot2R^2h^2=\frac89\cdot3R^2h^2\le\frac89\,t\le t,\]


where the last inequality uses \(t\ge0\).




(3) It suffices to prove the inequality between the coefficients of \(h^2p\) and then multiply by \(h^2p\ge0\). Multiplying by \(1+t>0\), the coefficient inequality is equivalent to


\[(A+m)+\frac{2R^2}m(1+t)\le A+2m+\frac{2R^2}m,\qquad\text{that is,}\qquad\frac{2R^2}m\,t\le m.\]


Since \(\frac{2R^2}m\,t=2R^2h\), and \(2R^2h\le3R^2h\le m\) because \(R^2h\ge0\), the inequality holds. \(\square\)




Part (2) shows that the SGLD contraction factor, which is the LMC factor \(\frac{1-mh}{1+mh}\) increased by the stochastic-gradient contribution \(2R^2h^2\), is still at most \((1+mh)^{-1}<1\). Part (1) is used for this purpose; in the application \(m^2\le R^2\) is \(\bar m\le M_{\rm rms}\).

\begin{leanbox}{Lean: \texttt{sgld\_\allowbreak{}dimensionless\_\allowbreak{}step}, \texttt{sgld\_\allowbreak{}contraction\_\allowbreak{}le}, \texttt{sgld\_\allowbreak{}bias\_\allowbreak{}le}}
\leansource{\texttt{sgld\_\allowbreak{}dimensionless\_\allowbreak{}step} \quad AverageLMC/\allowbreak{}SGLDRecurrence.lean, lines 14–21}
\begin{leancode}
\leanline{0}{theorem\ sgld\_dimensionless\_step\ \{m\ R\ h\ :\ ℝ\}\ (hm\ :\ 0\ <\ m)\ (hh\ :\ 0\ ≤\ h)}
\leanline{4}{(hR\ :\ m\ \textasciicircum{}\ 2\ ≤\ R\ \textasciicircum{}\ 2)\ (hstep\ :\ 3\ *\ R\ \textasciicircum{}\ 2\ *\ h\ ≤\ m)\ :\ m\ *\ h\ ≤\ 1\ /\ 3}
\end{leancode}
\leansource{\texttt{sgld\_\allowbreak{}contraction\_\allowbreak{}le} \quad AverageLMC/\allowbreak{}SGLDRecurrence.lean, lines 24–34}
\begin{leancode}
\leanline{0}{theorem\ sgld\_contraction\_le\ \{m\ R\ h\ :\ ℝ\}\ (hm\ :\ 0\ <\ m)\ (hh\ :\ 0\ ≤\ h)}
\leanline{4}{(hR\ :\ m\ \textasciicircum{}\ 2\ ≤\ R\ \textasciicircum{}\ 2)\ (hstep\ :\ 3\ *\ R\ \textasciicircum{}\ 2\ *\ h\ ≤\ m)\ :}
\leanline{4}{(1\ -\ m\ *\ h)\ /\ (1\ +\ m\ *\ h)\ +\ 2\ *\ R\ \textasciicircum{}\ 2\ *\ h\ \textasciicircum{}\ 2\ ≤\ 1\ /\ (1\ +\ m\ *\ h)}
\end{leancode}
\leansource{\texttt{sgld\_\allowbreak{}bias\_\allowbreak{}le} \quad AverageLMC/\allowbreak{}SGLDRecurrence.lean, lines 37–48}
\begin{leancode}
\leanline{0}{theorem\ sgld\_bias\_le\ \{m\ R\ A\ h\ p\ :\ ℝ\}\ (hm\ :\ 0\ <\ m)\ (hh\ :\ 0\ ≤\ h)\ (hp\ :\ 0\ ≤\ p)}
\leanline{4}{(hstep\ :\ 3\ *\ R\ \textasciicircum{}\ 2\ *\ h\ ≤\ m)\ :}
\leanline{4}{((A\ +\ m)\ /\ (1\ +\ m\ *\ h)\ +\ 2\ *\ R\ \textasciicircum{}\ 2\ /\ m)\ *\ h\ \textasciicircum{}\ 2\ *\ p\ ≤}
\leanline{6}{(A\ +\ 2\ *\ m\ +\ 2\ *\ R\ \textasciicircum{}\ 2\ /\ m)\ *\ h\ \textasciicircum{}\ 2\ *\ p\ /\ (1\ +\ m\ *\ h)}
\end{leancode}
\leanlabel{How the Lean proof follows the argument}
\begin{enumerate}


\item \texttt{sgld\_\allowbreak{}dimensionless\_\allowbreak{}step} has the hypotheses \texttt{hm : 0 < m}, \texttt{hh : 0 \ensuremath{\le} h}, \texttt{h\allowbreak{}R : m \textasciicircum{} 2 \ensuremath{\le} R \textasciicircum{} 2} and \texttt{hstep : 3 * R \textasciicircum{} 2 * h \ensuremath{\le} m}. \texttt{hr} multiplies \texttt{h\allowbreak{}R} by \(h\), \texttt{hmul} is \(m(3mh-1)\le0\) (\texttt{nlinarith}), and the conclusion is proved by contradiction (\texttt{by\_\allowbreak{}contra}): if \(mh>\frac13\), then \(m(3mh-1)>0\) by \texttt{mul\_\allowbreak{}pos}.

\item \texttt{sgld\_\allowbreak{}contraction\_\allowbreak{}le} uses part (1) as \texttt{ht}, clears the denominator with \texttt{le\_\allowbreak{}div\_\allowbreak{}iff\ensuremath{{}_{0}} hd} and \texttt{field\_\allowbreak{}simp}, and lets \texttt{nlinarith} combine \texttt{h1} (\(3R^2h^2\le mh\)) and \texttt{h2} (\(2R^2h^2(1+mh)\le2R^2h^2\cdot\frac43\)).

\item \texttt{sgld\_\allowbreak{}bias\_\allowbreak{}le} proves the coefficient inequality \texttt{hc} by the same denominator clearing and \texttt{nlinarith}, multiplies it by \(h^2p\ge0\) to obtain \texttt{hs}, and matches the two sides with \texttt{convert hs using 1 <;> ring}.

\end{enumerate}
\end{leanbox}
\subsection{\texorpdfstring{Iteration of the SGLD recurrence}{Iteration of the SGLD recurrence}}
\label{report-step-AverageLMC-SGLDRecurrence-lean-2}


\textbf{Claim.} Let \(m>0\), \(h\ge0\), \(p\ge0\), \(A\ge0\) and \(R\in\mathbb R\) with \(m^2\le R^2\) and \(3R^2h\le m\). Let \((r_k)_{k\ge0}\) be nonnegative real numbers such that, for every \(k\),


\[r_{k+1}\le\Bigl(\frac{1-mh}{1+mh}+2R^2h^2\Bigr)r_k+\Bigl(\frac{A+m}{1+mh}+\frac{2R^2}m\Bigr)h^2p.\qquad(1)\]


Then for every \(k\in\mathbb N\),


\[r_k\le(1+mh)^{-k}r_0+\frac{(A+2m+2R^2/m)\,hp}m.\]


\textbf{Proof.} Put \(q=(1+mh)^{-1}\) and \(c=(A+2m+2R^2/m)hp/m\). Then \(q\ge0\), and \(c\ge0\) because \(A,m,h,p\ge0\) and \(R^2/m\ge0\). We verify the hypothesis of the finite-horizon affine iteration of the scalar-recurrence section, \(r_{j+1}\le q\,r_j+(1-q)c\) for every \(j\). Multiplying part (2) of Step 1 by \(r_j\ge0\) bounds the first term of (1) by \(q\,r_j\). Part (3) of Step 1 bounds the second term by \((A+2m+2R^2/m)h^2p/(1+mh)\), and this equals \((1-q)c\):


\[(1-q)\,c=\frac{mh}{1+mh}\cdot\frac{(A+2m+2R^2/m)\,hp}m=\frac{(A+2m+2R^2/m)\,h^2p}{1+mh}.\]


Hence \(r_{j+1}\le q\,r_j+(1-q)c\) for every \(j\), and the affine iteration gives \(r_k\le q^kr_0+c\), which is the claim since \(q^k=(1+mh)^{-k}\). \(\square\)




The nonnegativity of the \(r_k\) is used to multiply the contraction inequality; in the application \(r_k=\mathsf W_2^2(\nu_k,\pi)\).

\begin{leanbox}{Lean: \texttt{sgld\_\allowbreak{}recurrence\_\allowbreak{}bound}}
\leansource{\texttt{sgld\_\allowbreak{}recurrence\_\allowbreak{}bound} \quad AverageLMC/\allowbreak{}SGLDRecurrence.lean, lines 52–80}
\begin{leancode}
\leanline{0}{theorem\ sgld\_recurrence\_bound\ (r\ :\ ℕ\ →\ ℝ)\ (m\ R\ A\ h\ p\ :\ ℝ)}
\leanline{4}{(hm\ :\ 0\ <\ m)\ (hh\ :\ 0\ ≤\ h)\ (hp\ :\ 0\ ≤\ p)\ (hA\ :\ 0\ ≤\ A)}
\leanline{4}{(hR\ :\ m\ \textasciicircum{}\ 2\ ≤\ R\ \textasciicircum{}\ 2)\ (hstep\ :\ 3\ *\ R\ \textasciicircum{}\ 2\ *\ h\ ≤\ m)}
\leanline{4}{(hr\ :\ ∀\ n,\ 0\ ≤\ r\ n)}
\leanline{4}{(hrec\ :\ ∀\ n,\ r\ (n\ +\ 1)\ ≤}
\leanline{6}{((1\ -\ m\ *\ h)\ /\ (1\ +\ m\ *\ h)\ +\ 2\ *\ R\ \textasciicircum{}\ 2\ *\ h\ \textasciicircum{}\ 2)\ *\ r\ n\ +}
\leanline{6}{((A\ +\ m)\ /\ (1\ +\ m\ *\ h)\ +\ 2\ *\ R\ \textasciicircum{}\ 2\ /\ m)\ *\ h\ \textasciicircum{}\ 2\ *\ p)}
\leanline{4}{(k\ :\ ℕ)\ :}
\leanline{4}{r\ k\ ≤\ (1\ +\ m\ *\ h)⁻¹\ \textasciicircum{}\ k\ *\ r\ 0\ +}
\leanline{6}{(A\ +\ 2\ *\ m\ +\ 2\ *\ R\ \textasciicircum{}\ 2\ /\ m)\ *\ h\ *\ p\ /\ m}
\end{leancode}
\leanlabel{How the Lean proof follows the argument}
\texttt{sgld\_\allowbreak{}recurrence\_\allowbreak{}bound} takes the sign conditions \texttt{hm}, \texttt{hh}, \texttt{hp}, \texttt{h\allowbreak{}A}, the conditions \texttt{h\allowbreak{}R} and \texttt{hstep}, the nonnegativity \texttt{hr : \ensuremath{\forall} n, 0 \ensuremath{\le} r n} and the recurrence \texttt{hrec}, which is (1). With \texttt{q := (1 + m * h)\ensuremath{{}^{-1}}} and \texttt{c} as above (\texttt{hq}, \texttt{hc} their signs), \texttt{apply affine\_\allowbreak{}floor\_\allowbreak{}up\_\allowbreak{}to r q c hq hc k} reduces the goal to the one-step affine inequality. In it, \texttt{hcontract} multiplies \texttt{sgld\_\allowbreak{}contraction\_\allowbreak{}le} by \(r_j\ge0\), \texttt{hbias} is \texttt{sgld\_\allowbreak{}bias\_\allowbreak{}le}, \texttt{hfix} is the identity for \((1-q)c\) (by \texttt{field\_\allowbreak{}simp} and \texttt{ring}), and \texttt{linarith} combines them with \texttt{hrec j}.
\end{leanbox}
\section{\texorpdfstring{Sharp initialization radius from actual transport}{Sharp initialization radius from actual transport}}
\label{proof-module-AverageLMC-InitialMoment-lean}


The centered stochastic-gradient proof needs the initial squared distance from the true minimizer. This theorem derives it from finite initial transport and the sharp Gibbs concentration radius, rather than assuming a separate initialization moment bound.



For any true minimizer \(x_\star\), strong convexity gives the sharp Gibbs radius, and the coupling argument bounds the initial radius by the actual transport distance:


\[\int\|x-x_\star\|^2\,d\pi(x)\le\frac pm.\]
\[\int\|x-x_\star\|^2\,d\nu(x)\le2W_2^2(\nu,\pi)+\frac{2p}{m}.\]

Earlier results used here: \cite[Section~3.5]{LangevinLean2026Proximal}; \cite[Section~3.32]{LangevinLean2026ULASC}.
\subsection{\texorpdfstring{Integrate the two-term radius inequality and take the infimum}{Integrate the two-term radius inequality and take the infimum}}
\label{report-step-AverageLMC-InitialMoment-lean-1}


Choose any finite-cost coupling of \(\mu\) and the actual Gibbs law. Both marginals have finite second moments: the Gibbs moment follows from strong convexity, and the \(\mu\) moment follows from finite transport. At each pair (\(x\),\(y\)), write \(x-anchor=(x-y)+(y-anchor)\); the squared norm is at most twice the squared coupling displacement plus twice the squared target radius. Integrate, use the two exact marginal pushforwards, and identify the coupling displacement integral with the square of its \(L^{2}\) coupling cost. The true minimizer and strong convexity give the sharp target radius at most \(p/m\). Thus every finite coupling bounds the initial radius by twice its squared cost plus \(2p/m\). Passing to the actual transport infimum proves the asserted inequality without assuming an optimal coupling or a Hessian regularity condition.

\begin{leanbox}{Lean: \texttt{initial\_\allowbreak{}radius\_\allowbreak{}le}}
\leansource{\texttt{initial\_\allowbreak{}radius\_\allowbreak{}le} \quad AverageLMC/\allowbreak{}InitialMoment.lean, lines 16–76}
\begin{leancode}
\leanline{0}{theorem\ initial\_radius\_le\ \{p\ :\ ℕ\}\ (f\ :\ State\ p\ →\ ℝ)\ (m\ :\ ℝ)\ (hm\ :\ 0\ <\ m)}
\leanline{4}{(hf\ :\ Differentiable\ ℝ\ f)\ (hsc\ :\ StrongConvexOn\ univ\ m\ f)}
\leanline{4}{(anchor\ :\ State\ p)\ (hanchor\ :\ ∀\ x,\ f\ anchor\ ≤\ f\ x)}
\leanline{4}{(μ\ :\ Measure\ (State\ p))\ [IsProbabilityMeasure\ μ]}
\leanline{4}{(hfinite\ :\ wasserstein2\ μ\ (gibbsMeasure\ f)\ ≠\ ⊤)\ :}
\leanline{4}{(∫\ x,\ ‖x\ -\ anchor‖\ \textasciicircum{}\ 2\ ∂μ)\ ≤}
\leanline{6}{2\ *\ wasserstein2Real\ μ\ (gibbsMeasure\ f)\ \textasciicircum{}\ 2\ +\ 2\ *\ (p\ :\ ℝ)\ /\ m}
\end{leancode}
\leanlabel{How the Lean proof follows the argument}
\texttt{initial\_\allowbreak{}radius\_\allowbreak{}le} gives the formal counterpart of this step. The main intermediate facts are \texttt{h\ensuremath{\pi}}, \texttt{h\ensuremath{\mu}}, \texttt{hradius}, \texttt{hfst}, \texttt{hsnd}, \texttt{hdiff} in that order. Named proof applications include \texttt{real\_\allowbreak{}bound\_\allowbreak{}of\_\allowbreak{}finite\_\allowbreak{}couplings}, \texttt{h.\allowbreak{}comp\_\allowbreak{}of\_\allowbreak{}map}; the adjacent proof shows their hypotheses and order of use.
\end{leanbox}


Earlier results used here: \cite[Section~3.17.2]{LangevinLean2026ULASC}; \cite[Section~3.17.3]{LangevinLean2026ULASC}; \cite[Section~3.17.4]{LangevinLean2026ULASC}; \cite[Section~3.21.1]{LangevinLean2026ULASC}; \cite[Section~3.26.15]{LangevinLean2026ULASC}; \cite[Section~3.26.3]{LangevinLean2026ULASC}; \cite[Section~3.11.1]{LangevinLean2026IPLA}; \cite[Section~3.15.1]{LangevinLean2026Proximal}; \cite[Section~3.5.2]{LangevinLean2026Proximal}.
\section{\texorpdfstring{Convergence of control-variate SGLD}{Convergence of control-variate SGLD}}
\label{proof-module-AverageLMC-SGLDTheorem-lean}


This section proves the convergence of control-variate SGLD, Theorem 3 of the paper, together with the intermediate claims of Proposition 3. The potential is \(f=\sum_{i=1}^nf_i\) on \(\mathbb R^p\), each \(f_i\) differentiable with Lipschitz gradient, \(f\) is \(m\)-strongly convex with \(m>0\), \(\pi\) is the Gibbs measure of \(f\), \(\theta_*\) is a global minimizer of \(f\), and \(R=nM_{\rm rms}\), where \(M_{\rm rms}^2=\frac1n\sum_i\|\nabla f_i\|_{\rm Lip}^2\). With the notation of the SGLD transition section, the chain of laws is \(\nu_0=\nu\), \(\nu_{k+1}=(\nu_k)^+\).




The steps are: the one-step bound (first claim of Proposition 3), obtained by adding the variance decomposition of SGLD to the one-step bound of Theorem 1; the bound \(\mathbb E\|g_I(\vartheta)\|^2\le R^2\,\mathbb E\|\vartheta-\theta_*\|^2\) and the moment bound \(\mathbb E\|\vartheta-\theta_*\|^2\le2\mathsf W_2^2(\nu,\pi)+2p/m\), which give the recurrence of the second claim; and the iteration of this recurrence under the step condition \(3R^2h\le m\). In the paper's units \(m=n\bar m\) and \(\bar M_{\rm av}=M_{\rm av}(f/n)\), the result is


\[\mathsf W_2^2(\nu_k,\pi)\le(1+n\bar mh)^{-k}\,\mathsf W_2^2(\nu_0,\pi)+\Bigl(\bar M_{\rm av}+2\bar m+\frac{2M_{\rm rms}^2}{\bar m}\Bigr)\frac{hp}{\bar m}\qquad\Bigl(0\le h\le\frac{\bar m}{3nM_{\rm rms}^2}\Bigr),\]


valid for every initial probability measure, in \([0,\infty]\).

\subsection{\texorpdfstring{The one-step bound with the estimator energy}{The one-step bound with the estimator energy}}
\label{report-step-AverageLMC-SGLDTheorem-lean-1}


\textbf{Claim.} Let \(p,n\ge1\), let each \(f_i\) be differentiable with Lipschitz gradient, let \(f=\sum_if_i\) be \(m\)-strongly convex with \(m>0\) and have an \(M\)-Lipschitz gradient (\(M\ge0\)), let \(h\ge0\), and let \(\theta_*\) satisfy \(\nabla f(\theta_*)=0\). Let \(\pi\) be the Gibbs measure of \(f\) and \(\nu\) a probability measure with \(\mathsf W_2(\nu,\pi)<\infty\). Then, with \(G=\int\|\nabla f\|^2d\nu\) and \(M_{\rm av}\) the average coordinate smoothness of \(f\),


\[\mathsf W_2^2(\nu^+,\pi)\le\frac{1-mh}{1+mh}\,\mathsf W_2^2(\nu,\pi)+\frac{(M_{\rm av}+m)h^2p}{1+mh}+\frac{h^2(hM-2)}{1+mh}\,G+h^2\int\|g_i(x)\|^2\,(\nu\otimes U_n)(dx,di).\]


\textbf{Proof.} \emph{Preliminaries.} \(f\) is differentiable as a finite sum of differentiable functions. The Gibbs measure of a differentiable strongly convex potential is a probability measure with finite second moment. Because \(\mathsf W_2(\nu,\pi)<\infty\), there is a coupling \(\gamma\) of \(\nu\) and \(\pi\) with \(\int\|x-y\|^2d\gamma<\infty\); from \(\|x\|^2\le2\|x-y\|^2+2\|y\|^2\) it follows that \(\nu\) has a finite second moment.




\emph{Two inequalities.} The variance decomposition of SGLD (inequalities (2) and (3) of its section) gives


\[\mathsf W_2^2(\nu^+,\pi)\le\mathsf W_2^2(\nu^+_{\rm LMC},\pi)+h^2\int\|g_i\|^2\,d(\nu\otimes U_n)-h^2G.\qquad(1)\]


The one-step bound of Theorem 1 for the full-gradient step gives


\[\mathsf W_2^2(\nu^+_{\rm LMC},\pi)\le\frac{1-mh}{1+mh}\,\mathsf W_2^2(\nu,\pi)+\frac{(M_{\rm av}+m)h^2p}{1+mh}+\frac{h^2(hM-1)}{1+mh}\,G.\qquad(2)\]


\emph{The coefficient of \(G\).} Inserting (2) into (1), \(G\) appears with the coefficient \(\frac{h^2(hM-1)}{1+mh}-h^2\). Multiplying by \(1+mh>0\),


\[h^2(hM-1)-h^2(1+mh)=h^2(hM-2)-mh^3\le h^2(hM-2),\]


because \(mh^3\ge0\). Hence \(\frac{h^2(hM-1)}{1+mh}-h^2\le\frac{h^2(hM-2)}{1+mh}\), and multiplying by \(G\ge0\) and combining (1) and (2) gives the claim. \(\square\)




The last term is \(nh^2\sum_i\|\nabla f_i(\vartheta)-\nabla f_i(\theta_*)\|_{L_2}^2\) by the estimator identity of the SGLD transition section, so this is the first claim of Proposition 3. No bound on the estimator variance is used yet.

\begin{leanbox}{Lean: \texttt{sgld\_\allowbreak{}one\_\allowbreak{}step}}
\leansource{\texttt{sgld\_\allowbreak{}one\_\allowbreak{}step} \quad AverageLMC/\allowbreak{}SGLDTheorem.lean, lines 19–53}
\begin{leancode}
\leanline{0}{theorem\ sgld\_one\_step\ \{p\ n\ :\ ℕ\}\ (hp\ :\ 0\ <\ p)\ (hn\ :\ 0\ <\ n)}
\leanline{4}{(f\ :\ Fin\ n\ →\ State\ p\ →\ ℝ)\ (hf\ :\ ∀\ i,\ Differentiable\ ℝ\ (f\ i))}
\leanline{4}{(L\ :\ Fin\ n\ →\ ℝ≥0)\ (hLip\ :\ ∀\ i,\ LipschitzWith\ (L\ i)\ (gradient\ (f\ i)))}
\leanline{4}{(m\ h\ :\ ℝ)\ (M\ :\ ℝ≥0)\ (hm\ :\ 0\ <\ m)\ (hh\ :\ 0\ ≤\ h)}
\leanline{4}{(hsc\ :\ StrongConvexOn\ univ\ m\ (componentPotential\ f))}
\leanline{4}{(hfull\ :\ LipschitzWith\ M\ (gradient\ (componentPotential\ f)))}
\leanline{4}{(anchor\ :\ State\ p)\ (hstar\ :\ gradient\ (componentPotential\ f)\ anchor\ =\ 0)}
\leanline{4}{(μ\ :\ Measure\ (State\ p))\ [IsProbabilityMeasure\ μ]}
\leanline{4}{(hfinite\ :\ wasserstein2\ μ\ (gibbsMeasure\ (componentPotential\ f))\ ≠\ ⊤)\ :}
\leanline{4}{wasserstein2Real\ (sgldLawStep\ f\ anchor\ h\ μ)\ (gibbsMeasure\ (componentPotential\ f))\ \textasciicircum{}\ 2\ ≤}
\leanline{6}{(1\ -\ m\ *\ h)\ /\ (1\ +\ m\ *\ h)\ *\ wasserstein2Real\ μ\ (gibbsMeasure\ (componentPotential\ f))\ \textasciicircum{}\ 2\ +}
\leanline{6}{(averageSmoothness\ (componentPotential\ f)\ +\ m)\ *\ h\ \textasciicircum{}\ 2\ *\ p\ /\ (1\ +\ m\ *\ h)\ +}
\leanline{6}{h\ \textasciicircum{}\ 2\ *\ (h\ *\ (M\ :\ ℝ)\ -\ 2)\ /\ (1\ +\ m\ *\ h)\ *}
\leanline{8}{(∫\ x,\ ‖gradient\ (componentPotential\ f)\ x‖\ \textasciicircum{}\ 2\ ∂μ)\ +}
\leanline{6}{h\ \textasciicircum{}\ 2\ *\ ∫\ q\ :\ State\ p\ ×\ Fin\ n,\ ‖controlVariateEstimator\ f\ anchor\ q.1\ q.2‖\ \textasciicircum{}\ 2}
\leanline{8}{∂μ.prod\ (uniformIndexLaw\ n)}
\end{leancode}
\leanlabel{How the Lean proof follows the argument}
The Gibbs measure is \texttt{gibbs\allowbreak{}Measure (component\allowbreak{}Potential f)}, defined in the unadjusted Langevin library as the normalized measure with density \(e^{-f}\) with respect to Lebesgue measure \texttt{volume}. In \texttt{sgld\_\allowbreak{}one\_\allowbreak{}step}, \texttt{hdiff} is the differentiability of the sum (\texttt{Differentiable.\allowbreak{}fun\_\allowbreak{}sum}); \texttt{is\allowbreak{}Probability\allowbreak{}Measure\_\allowbreak{}gibbs\allowbreak{}Measure} and \texttt{mem\allowbreak{}Lp\_\allowbreak{}id\_\allowbreak{}gibbs\allowbreak{}Measure} give that \(\pi\) is a probability measure with finite second moment, and \texttt{mem\allowbreak{}Lp\_\allowbreak{}id\_\allowbreak{}of\_\allowbreak{}wasserstein2\_\allowbreak{}ne\_\allowbreak{}top} transfers the finite second moment to \(\mu=\nu\) through the finite distance. \texttt{hvar} is \texttt{sgld\_\allowbreak{}variance\_\allowbreak{}addition}, rewritten by \texttt{sgld\_\allowbreak{}error\_\allowbreak{}variance\_\allowbreak{}identity}, which is (1); \texttt{hlmc} is \texttt{lmc\_\allowbreak{}one\_\allowbreak{}step}, which is (2). \texttt{h\allowbreak{}G} is \(G\ge0\) (\texttt{integral\_\allowbreak{}nonneg}), \texttt{hc} is the coefficient inequality (cleared of the denominator by \texttt{le\_\allowbreak{}div\_\allowbreak{}iff\ensuremath{{}_{0}} hd} and \texttt{field\_\allowbreak{}simp}, then \texttt{nlinarith}), \texttt{hg} multiplies it by \(G\), and the final \texttt{nlinarith} adds everything.
\end{leanbox}


Earlier results used here: \cite[Section~3.17.3]{LangevinLean2026ULASC}; \cite[Section~3.17.4]{LangevinLean2026ULASC}; \cite[Section~3.21.1]{LangevinLean2026ULASC}; \cite[Section~3.26.13]{LangevinLean2026ULASC}; \cite[Section~3.26.15]{LangevinLean2026ULASC}; \cite[Section~3.26.3]{LangevinLean2026ULASC}.
\subsection{\texorpdfstring{The recurrence in terms of the root-mean-square constant; finiteness along the chain}{The recurrence in terms of the root-mean-square constant; finiteness along the chain}}
\label{report-step-AverageLMC-SGLDTheorem-lean-2}


\textbf{Claim.} (1) In addition to the hypotheses of Step 1, assume \(hM\le2\) and that \(\theta_*\) is a global minimizer of \(f\), \(f(\theta_*)\le f(x)\) for all \(x\). Then, with \(R=nM_{\rm rms}\),


\[\mathsf W_2^2(\nu^+,\pi)\le\Bigl(\frac{1-mh}{1+mh}+2R^2h^2\Bigr)\mathsf W_2^2(\nu,\pi)+\Bigl(\frac{M_{\rm av}+m}{1+mh}+\frac{2R^2}m\Bigr)h^2p.\]


(2) Let \(n\ge1\), let each \(f_i\) be differentiable with Lipschitz gradient, let \(f\) be \(m\)-strongly convex (\(m>0\)) with Lipschitz gradient, and let \(\theta_*\in\mathbb R^p\) and \(h\in\mathbb R\) be arbitrary. If \(\nu\) is a probability measure with \(\mathsf W_2(\nu,\pi)<\infty\), then \(\mathsf W_2(\nu_k,\pi)<\infty\) for every \(k\).




\textbf{Proof.} (1) \emph{The gradient term is dropped.} Since \(hM\le2\), the factor \(h^2(hM-2)\) is \(\le0\); divided by \(1+mh>0\) and multiplied by \(G\ge0\), the third term of Step 1 is \(\le0\).




\emph{The estimator moment.} For each \(x\), by the definition of the gradient seminorm, \(\|\nabla f_i(x)-\nabla f_i(\theta_*)\|\le M_\infty(f_i)\|x-\theta_*\|\), so


\[\int\|g_i(x)\|^2U_n(di)=n\sum_{i=1}^n\|\nabla f_i(x)-\nabla f_i(\theta_*)\|^2\le n\sum_{i=1}^nM_\infty(f_i)^2\,\|x-\theta_*\|^2=R^2\|x-\theta_*\|^2,\]


since \(\sum_iM_\infty(f_i)^2=nM_{\rm rms}^2\). (This pointwise bound is proved in the estimator section.) Integrating over \(\nu\) with Fubini's theorem (the integrands are integrable because \(\nu\) has a finite second moment),


\[\int\|g_i(x)\|^2\,d(\nu\otimes U_n)\le R^2\int\|x-\theta_*\|^2\,\nu(dx).\]


\emph{The moment around the minimizer.} For a global minimizer \(\theta_*\) of the strongly convex \(f\), the section on the initial moment proves \(\int\|x-\theta_*\|^2\nu(dx)\le2\mathsf W_2^2(\nu,\pi)+2p/m\). Hence


\[h^2\int\|g_i\|^2\,d(\nu\otimes U_n)\le h^2R^2\Bigl(2\mathsf W_2^2(\nu,\pi)+\frac{2p}m\Bigr)=2R^2h^2\,\mathsf W_2^2(\nu,\pi)+\frac{2R^2}m\,h^2p.\]


Adding this to the remaining terms of Step 1 and grouping \(\frac{(M_{\rm av}+m)h^2p}{1+mh}+\frac{2R^2}mh^2p\) gives the claim.




(2) Induction on \(k\). For \(k=0\) this is the hypothesis. If \(\mathsf W_2(\nu_k,\pi)<\infty\), then \(\nu_k\) has a finite second moment (as in Step 1, using the finite second moment of \(\pi\)); by the moment propagation of the SGLD transition section, \(\nu_{k+1}\) has a finite second moment; and two probability measures with finite second moments are at finite distance, since the product coupling has finite cost. \(\square\)

\begin{leanbox}{Lean: \texttt{sgld\_\allowbreak{}one\_\allowbreak{}step\_\allowbreak{}rms}, \texttt{sgldLaws\_\allowbreak{}finite}}
\leansource{\texttt{sgld\_\allowbreak{}one\_\allowbreak{}step\_\allowbreak{}rms} \quad AverageLMC/\allowbreak{}SGLDTheorem.lean, lines 56–103}
\begin{leancode}
\leanline{0}{theorem\ sgld\_one\_step\_rms\ \{p\ n\ :\ ℕ\}\ (hp\ :\ 0\ <\ p)\ (hn\ :\ 0\ <\ n)}
\leanline{4}{(f\ :\ Fin\ n\ →\ State\ p\ →\ ℝ)\ (hf\ :\ ∀\ i,\ Differentiable\ ℝ\ (f\ i))}
\leanline{4}{(L\ :\ Fin\ n\ →\ ℝ≥0)\ (hLip\ :\ ∀\ i,\ LipschitzWith\ (L\ i)\ (gradient\ (f\ i)))}
\leanline{4}{(m\ h\ :\ ℝ)\ (M\ :\ ℝ≥0)\ (hm\ :\ 0\ <\ m)\ (hh\ :\ 0\ ≤\ h)}
\leanline{4}{(hsc\ :\ StrongConvexOn\ univ\ m\ (componentPotential\ f))}
\leanline{4}{(hfull\ :\ LipschitzWith\ M\ (gradient\ (componentPotential\ f)))}
\leanline{4}{(hstep\ :\ h\ *\ (M\ :\ ℝ)\ ≤\ 2)}
\leanline{4}{(anchor\ :\ State\ p)\ (hanchor\ :\ ∀\ x,\ componentPotential\ f\ anchor\ ≤\ componentPotential\ f\ x)}
\leanline{4}{(hstar\ :\ gradient\ (componentPotential\ f)\ anchor\ =\ 0)}
\leanline{4}{(μ\ :\ Measure\ (State\ p))\ [IsProbabilityMeasure\ μ]}
\leanline{4}{(hfinite\ :\ wasserstein2\ μ\ (gibbsMeasure\ (componentPotential\ f))\ ≠\ ⊤)\ :}
\leanline{4}{let\ R\ :=\ (n\ :\ ℝ)\ *\ rmsComponentSmoothness\ f}
\leanline{4}{wasserstein2Real\ (sgldLawStep\ f\ anchor\ h\ μ)\ (gibbsMeasure\ (componentPotential\ f))\ \textasciicircum{}\ 2\ ≤}
\leanline{6}{((1\ -\ m\ *\ h)\ /\ (1\ +\ m\ *\ h)\ +\ 2\ *\ R\ \textasciicircum{}\ 2\ *\ h\ \textasciicircum{}\ 2)\ *}
\leanline{8}{wasserstein2Real\ μ\ (gibbsMeasure\ (componentPotential\ f))\ \textasciicircum{}\ 2\ +}
\leanline{6}{((averageSmoothness\ (componentPotential\ f)\ +\ m)\ /\ (1\ +\ m\ *\ h)\ +\ 2\ *\ R\ \textasciicircum{}\ 2\ /\ m)\ *}
\leanline{8}{h\ \textasciicircum{}\ 2\ *\ p}
\end{leancode}
\leansource{\texttt{sgldLaws\_\allowbreak{}finite} \quad AverageLMC/\allowbreak{}SGLDTheorem.lean, lines 106–125}
\begin{leancode}
\leanline{0}{theorem\ sgldLaws\_finite\ \{p\ n\ :\ ℕ\}\ (hn\ :\ 0\ <\ n)}
\leanline{4}{(f\ :\ Fin\ n\ →\ State\ p\ →\ ℝ)\ (hf\ :\ ∀\ i,\ Differentiable\ ℝ\ (f\ i))}
\leanline{4}{(L\ :\ Fin\ n\ →\ ℝ≥0)\ (hLip\ :\ ∀\ i,\ LipschitzWith\ (L\ i)\ (gradient\ (f\ i)))}
\leanline{4}{(m\ :\ ℝ)\ (M\ :\ ℝ≥0)\ (hm\ :\ 0\ <\ m)}
\leanline{4}{(hsc\ :\ StrongConvexOn\ univ\ m\ (componentPotential\ f))}
\leanline{4}{(hfull\ :\ LipschitzWith\ M\ (gradient\ (componentPotential\ f)))}
\leanline{4}{(anchor\ :\ State\ p)\ (h\ :\ ℝ)\ (μ\ :\ Measure\ (State\ p))\ [IsProbabilityMeasure\ μ]}
\leanline{4}{(hfinite\ :\ wasserstein2\ μ\ (gibbsMeasure\ (componentPotential\ f))\ ≠\ ⊤)\ (k\ :\ ℕ)\ :}
\leanline{4}{wasserstein2\ (sgldLaws\ f\ anchor\ h\ μ\ k)\ (gibbsMeasure\ (componentPotential\ f))\ ≠\ ⊤}
\end{leancode}
\leanlabel{How the Lean proof follows the argument}
\begin{enumerate}


\item \texttt{sgld\_\allowbreak{}one\_\allowbreak{}step\_\allowbreak{}rms} adds \texttt{hstep : h * (M : \ensuremath{\mathbb R}) \ensuremath{\le} 2} and \texttt{hanchor : \ensuremath{\forall} x, component\allowbreak{}Potential f anchor \ensuremath{\le} component\allowbreak{}Potential f x}; its conclusion begins with \texttt{let R := (n : \ensuremath{\mathbb R}) * rms\allowbreak{}Component\allowbreak{}Smoothness f}, a local abbreviation for \(R\). \texttt{hs} is Step 1 and \texttt{hneg} shows that its gradient term is \(\le0\) (\texttt{mul\_\allowbreak{}nonpos\_\allowbreak{}of\_\allowbreak{}nonpos\_\allowbreak{}of\_\allowbreak{}nonneg}, \texttt{div\_\allowbreak{}nonpos\_\allowbreak{}of\_\allowbreak{}nonpos\_\allowbreak{}of\_\allowbreak{}nonneg}, \texttt{integral\_\allowbreak{}nonneg}). \texttt{hmoment} is the integrated estimator bound: Fubini's theorem \texttt{integral\_\allowbreak{}prod}, the pointwise bound \texttt{control\allowbreak{}Variate\allowbreak{}Estimator\_\allowbreak{}second\_\allowbreak{}moment\_\allowbreak{}le} and \texttt{integral\_\allowbreak{}mono}. \texttt{hradius} is \texttt{initial\_\allowbreak{}radius\_\allowbreak{}le}, and \texttt{hmoment'}, \texttt{hscaled} chain the two bounds and multiply by \(h^2\). After \texttt{dsimp only} unfolds \texttt{R} and divisions are written as multiplications by inverses, \texttt{nlinarith} combines the inequalities.

\item \texttt{sgld\allowbreak{}Laws\_\allowbreak{}finite} is the induction \texttt{induction k with}; the step uses \texttt{mem\allowbreak{}Lp\_\allowbreak{}id\_\allowbreak{}of\_\allowbreak{}wasserstein2\_\allowbreak{}ne\_\allowbreak{}top}, \texttt{sgld\allowbreak{}Law\allowbreak{}Step\_\allowbreak{}mem\allowbreak{}Lp} and \texttt{wasserstein2\_\allowbreak{}ne\_\allowbreak{}top\_\allowbreak{}of\_\allowbreak{}mem\allowbreak{}Lp} in that order. The line \texttt{have : Ne\allowbreak{}Zero n := \ensuremath{\langle}hn.\allowbreak{}ne'\ensuremath{\rangle}} provides the typeclass form of \(n\ne0\) used by the probability instances of the SGLD laws.

\end{enumerate}
\end{leanbox}


Earlier results used here: \cite[Section~3.17.3]{LangevinLean2026ULASC}; \cite[Section~3.17.4]{LangevinLean2026ULASC}; \cite[Section~3.21.1]{LangevinLean2026ULASC}; \cite[Section~3.21.2]{LangevinLean2026ULASC}; \cite[Section~3.26.13]{LangevinLean2026ULASC}; \cite[Section~3.26.15]{LangevinLean2026ULASC}; \cite[Section~3.26.3]{LangevinLean2026ULASC}.
\subsection{\texorpdfstring{The constant-step bound}{The constant-step bound}}
\label{report-step-AverageLMC-SGLDTheorem-lean-3}


\textbf{Claim.} (1) Let \(p,n\ge1\), let each \(f_i\) be differentiable with Lipschitz gradient, let \(m>0\), \(h\ge0\), let \(f\) be \(m\)-strongly convex, let \(R=nM_{\rm rms}\) satisfy \(3R^2h\le m\), let \(\theta_*\) be a global minimizer of \(f\), and let \(\nu\) be a probability measure with \(\mathsf W_2(\nu,\pi)<\infty\). Then for every \(k\in\mathbb N\),


\[\mathsf W_2^2(\nu_k,\pi)\le(1+mh)^{-k}\,\mathsf W_2^2(\nu,\pi)+\Bigl(M_{\rm av}+2m+\frac{2R^2}m\Bigr)\frac{hp}m.\]


(2) Let \(\bar m>0\), let \(f\) be \(n\bar m\)-strongly convex and let \(0\le h\le\bar m/(3nM_{\rm rms}^2)\); keep the other hypotheses of (1). Then for every \(k\),


\[\mathsf W_2^2(\nu_k,\pi)\le(1+n\bar mh)^{-k}\,\mathsf W_2^2(\nu,\pi)+\Bigl(\frac{M_{\rm av}}n+2\bar m+\frac{2M_{\rm rms}^2}{\bar m}\Bigr)\frac{hp}{\bar m}.\]


\textbf{Proof.} (1) We collect the facts needed for Step 2 and for the scalar SGLD recurrence.



\begin{itemize}


\item \(\nabla f\) is \(R\)-Lipschitz: \(\|\nabla f(x)-\nabla f(y)\|\le\sum_iM_\infty(f_i)\|x-y\|\le\sqrt{n\sum_iM_\infty(f_i)^2}\,\|x-y\|=R\|x-y\|\), by the triangle inequality and the Cauchy–Schwarz inequality in \(\mathbb R^n\) (proved in the estimator section).

\item \(m\le R\) (also proved there, from \(m\le M_\infty(f)\le R\)); hence \(R>0\) and \(m^2\le R^2\).

\item \(hR\le\frac13\): indeed \(R(3Rh-1)=3R^2h-R\le m-R\le0\), and \(R>0\) forces \(3Rh\le1\). In particular the step condition \(hR\le2\) of Step 2 holds with \(M=R\).

\item \(\nabla f(\theta_*)=0\): a global minimizer is a local minimizer, and the derivative of a differentiable function vanishes at a local minimizer.

\item \(M_{\rm av}\ge m>0\), by the comparison of constants in the first section.

\end{itemize}



Let \(r_j=\mathsf W_2^2(\nu_j,\pi)\), finite by Step 2 (2). Applying Step 2 (1) to \(\nu_j\) (with \(M=R\)) gives, for every \(j\),


\[r_{j+1}\le\Bigl(\frac{1-mh}{1+mh}+2R^2h^2\Bigr)r_j+\Bigl(\frac{M_{\rm av}+m}{1+mh}+\frac{2R^2}m\Bigr)h^2p.\]


The \(r_j\) are nonnegative, \(A=M_{\rm av}\ge0\), \(p\ge0\), \(m^2\le R^2\) and \(3R^2h\le m\), so the iteration of the scalar SGLD recurrence applies and gives the claim.




(2) Put \(m=n\bar m>0\). Since \(n\bar m\le nM_{\rm rms}\) by (1), \(M_{\rm rms}\ge\bar m>0\), so \(3nM_{\rm rms}^2>0\) and the step condition is equivalent to \(3nM_{\rm rms}^2h\le\bar m\); multiplying by \(n\) gives \(3R^2h\le m\). Part (1) applies, and


\[\Bigl(M_{\rm av}+2n\bar m+\frac{2n^2M_{\rm rms}^2}{n\bar m}\Bigr)\frac{hp}{n\bar m}=\Bigl(\frac{M_{\rm av}}n+2\bar m+\frac{2M_{\rm rms}^2}{\bar m}\Bigr)\frac{hp}{\bar m}.\qquad\square\]\begin{leanbox}{Lean: \texttt{sgld\_\allowbreak{}constant\_\allowbreak{}step}, \texttt{sgld\_\allowbreak{}constant\_\allowbreak{}step\_\allowbreak{}average}}
\leansource{\texttt{sgld\_\allowbreak{}constant\_\allowbreak{}step} \quad AverageLMC/\allowbreak{}SGLDTheorem.lean, lines 130–171}
\begin{leancode}
\leanline{0}{theorem\ sgld\_constant\_step\ \{p\ n\ :\ ℕ\}\ (hp\ :\ 0\ <\ p)\ (hn\ :\ 0\ <\ n)}
\leanline{4}{(f\ :\ Fin\ n\ →\ State\ p\ →\ ℝ)\ (hf\ :\ ∀\ i,\ Differentiable\ ℝ\ (f\ i))}
\leanline{4}{(L\ :\ Fin\ n\ →\ ℝ≥0)\ (hLip\ :\ ∀\ i,\ LipschitzWith\ (L\ i)\ (gradient\ (f\ i)))}
\leanline{4}{(m\ h\ :\ ℝ)\ (hm\ :\ 0\ <\ m)\ (hh\ :\ 0\ ≤\ h)}
\leanline{4}{(hsc\ :\ StrongConvexOn\ univ\ m\ (componentPotential\ f))}
\leanline{4}{(hstep\ :\ 3\ *\ ((n\ :\ ℝ)\ *\ rmsComponentSmoothness\ f)\ \textasciicircum{}\ 2\ *\ h\ ≤\ m)}
\leanline{4}{(anchor\ :\ State\ p)\ (hanchor\ :\ ∀\ x,\ componentPotential\ f\ anchor\ ≤\ componentPotential\ f\ x)}
\leanline{4}{(μ\ :\ Measure\ (State\ p))\ [IsProbabilityMeasure\ μ]}
\leanline{4}{(hfinite\ :\ wasserstein2\ μ\ (gibbsMeasure\ (componentPotential\ f))\ ≠\ ⊤)\ (k\ :\ ℕ)\ :}
\leanline{4}{let\ R\ :=\ (n\ :\ ℝ)\ *\ rmsComponentSmoothness\ f}
\leanline{4}{wasserstein2Real\ (sgldLaws\ f\ anchor\ h\ μ\ k)\ (gibbsMeasure\ (componentPotential\ f))\ \textasciicircum{}\ 2\ ≤}
\leanline{6}{(1\ +\ m\ *\ h)⁻¹\ \textasciicircum{}\ k\ *\ wasserstein2Real\ μ\ (gibbsMeasure\ (componentPotential\ f))\ \textasciicircum{}\ 2\ +}
\leanline{6}{(averageSmoothness\ (componentPotential\ f)\ +\ 2\ *\ m\ +\ 2\ *\ R\ \textasciicircum{}\ 2\ /\ m)\ *\ h\ *\ p\ /\ m}
\end{leancode}
\leansource{\texttt{sgld\_\allowbreak{}constant\_\allowbreak{}step\_\allowbreak{}average} \quad AverageLMC/\allowbreak{}SGLDTheorem.lean, lines 175–201}
\begin{leancode}
\leanline{0}{theorem\ sgld\_constant\_step\_average\ \{p\ n\ :\ ℕ\}\ (hp\ :\ 0\ <\ p)\ (hn\ :\ 0\ <\ n)}
\leanline{4}{(f\ :\ Fin\ n\ →\ State\ p\ →\ ℝ)\ (hf\ :\ ∀\ i,\ Differentiable\ ℝ\ (f\ i))}
\leanline{4}{(L\ :\ Fin\ n\ →\ ℝ≥0)\ (hLip\ :\ ∀\ i,\ LipschitzWith\ (L\ i)\ (gradient\ (f\ i)))}
\leanline{4}{(mbar\ h\ :\ ℝ)\ (hmbar\ :\ 0\ <\ mbar)\ (hh\ :\ 0\ ≤\ h)}
\leanline{4}{(hsc\ :\ StrongConvexOn\ univ\ ((n\ :\ ℝ)\ *\ mbar)\ (componentPotential\ f))}
\leanline{4}{(hstep\ :\ h\ ≤\ mbar\ /\ (3\ *\ n\ *\ rmsComponentSmoothness\ f\ \textasciicircum{}\ 2))}
\leanline{4}{(anchor\ :\ State\ p)\ (hanchor\ :\ ∀\ x,\ componentPotential\ f\ anchor\ ≤\ componentPotential\ f\ x)}
\leanline{4}{(μ\ :\ Measure\ (State\ p))\ [IsProbabilityMeasure\ μ]}
\leanline{4}{(hfinite\ :\ wasserstein2\ μ\ (gibbsMeasure\ (componentPotential\ f))\ ≠\ ⊤)\ (k\ :\ ℕ)\ :}
\leanline{4}{wasserstein2Real\ (sgldLaws\ f\ anchor\ h\ μ\ k)\ (gibbsMeasure\ (componentPotential\ f))\ \textasciicircum{}\ 2\ ≤}
\leanline{6}{(1\ +\ (n\ :\ ℝ)\ *\ mbar\ *\ h)⁻¹\ \textasciicircum{}\ k\ *\ wasserstein2Real\ μ\ (gibbsMeasure\ (componentPotential\ f))\ \textasciicircum{}\ 2\ +}
\leanline{6}{(averageSmoothness\ (componentPotential\ f)\ /\ n\ +\ 2\ *\ mbar\ +}
\leanline{8}{2\ *\ rmsComponentSmoothness\ f\ \textasciicircum{}\ 2\ /\ mbar)\ *\ h\ *\ p\ /\ mbar}
\end{leancode}
\leanlabel{How the Lean proof follows the argument}
\begin{enumerate}


\item In \texttt{sgld\_\allowbreak{}constant\_\allowbreak{}step}, \texttt{M} is \(R\) packaged as a nonnegative real, and \texttt{hfull} is its Lipschitz property \texttt{component\allowbreak{}Potential\_\allowbreak{}gradient\_\allowbreak{}lipschitz\_\allowbreak{}rms}. \texttt{hm\allowbreak{}R} is \texttt{strong\allowbreak{}Convex\_\allowbreak{}le\_\allowbreak{}n\_\allowbreak{}mul\_\allowbreak{}rms}, \texttt{h\allowbreak{}R} squares it with \texttt{sq\_\allowbreak{}le\_\allowbreak{}sq\ensuremath{{}_{0}}}, and \texttt{h\allowbreak{}Rh} is \(hR\le\frac13\), proved by contradiction from \texttt{hmul}. \texttt{hstar} is \(\nabla f(\theta_*)=0\): \texttt{hlocal} turns the global minimum into \texttt{Is\allowbreak{}Local\allowbreak{}Min}, and \texttt{Is\allowbreak{}Local\allowbreak{}Min.\allowbreak{}fderiv\_\allowbreak{}eq\_\allowbreak{}zero} (used as \texttt{hlocal.\allowbreak{}fderiv\_\allowbreak{}eq\_\allowbreak{}zero}) is Fermat's theorem. \texttt{h\allowbreak{}A} is \(M_{\rm av}\ge0\) from \texttt{strong\allowbreak{}Convex\_\allowbreak{}le\_\allowbreak{}average\allowbreak{}Smoothness}. \texttt{hrec j} is \texttt{sgld\_\allowbreak{}one\_\allowbreak{}step\_\allowbreak{}rms} applied to the law at step \texttt{j}, with finiteness from \texttt{sgld\allowbreak{}Laws\_\allowbreak{}finite}, and the final line is \texttt{sgld\_\allowbreak{}recurrence\_\allowbreak{}bound}.

\item \texttt{sgld\_\allowbreak{}constant\_\allowbreak{}step\_\allowbreak{}average} has \texttt{hsc : Strong\allowbreak{}Convex\allowbreak{}On univ ((n : \ensuremath{\mathbb R}) * mbar) (component\allowbreak{}Potential f)} and \texttt{hstep : h \ensuremath{\le} mbar /\allowbreak{} (3 * n * rms\allowbreak{}Component\allowbreak{}Smoothness f \textasciicircum{} 2)}. \texttt{h\allowbreak{}R} is \(M_{\rm rms}>0\), \texttt{hsmall} clears the denominator (\texttt{le\_\allowbreak{}div\_\allowbreak{}iff\ensuremath{{}_{0}} hd}), \texttt{hnorm} is \(3R^2h\le n\bar m\), \texttt{hb} is part (1), and \texttt{convert hb using 1} followed by \texttt{field\_\allowbreak{}simp} identifies the constants.

\end{enumerate}
\end{leanbox}


Earlier results used here: \cite[Section~3.17.3]{LangevinLean2026ULASC}; \cite[Section~3.17.4]{LangevinLean2026ULASC}; \cite[Section~3.26.3]{LangevinLean2026ULASC}.
\subsection{\texorpdfstring{Arbitrary initial distributions and the paper's constant}{Arbitrary initial distributions and the paper's constant}}
\label{report-step-AverageLMC-SGLDTheorem-lean-4}


\textbf{Claim.} Under the hypotheses of Step 3 (2), but for an arbitrary initial probability measure \(\nu\) (no finiteness of \(\mathsf W_2(\nu,\pi)\) assumed), the inequality


\[\mathsf W_2^2(\nu_k,\pi)\le(1+n\bar mh)^{-k}\,\mathsf W_2^2(\nu,\pi)+\Bigl(\bar M_{\rm av}+2\bar m+\frac{2M_{\rm rms}^2}{\bar m}\Bigr)\frac{hp}{\bar m}\]


holds in \([0,\infty]\) for every \(k\), first with \(\bar M_{\rm av}\) replaced by \(M_{\rm av}(f)/n\), and then with \(\bar M_{\rm av}=M_{\rm av}(f/n)\), the average coordinate smoothness of the average potential \(f/n\). This is Theorem 3 of the paper.




\textbf{Proof.} Put \(a=(1+n\bar mh)^{-k}\) and \(b=(M_{\rm av}(f)/n+2\bar m+2M_{\rm rms}^2/\bar m)hp/\bar m\). Then \(a>0\), and \(b\ge0\) because \(M_{\rm av}(f)\ge n\bar m>0\) (comparison of constants for the \(n\bar m\)-strongly convex \(f\), whose gradient is Lipschitz by Step 3), \(\bar m>0\), \(h\ge0\) and \(p\ge0\). The extension result for affine transport bounds applies with these \(a,b\): if \(\mathsf W_2(\nu,\pi)<\infty\), then \(\mathsf W_2(\nu_k,\pi)<\infty\) by Step 2 (2), and the real inequality is Step 3 (2). This proves the first form.




For the second form, the coordinate constants scale linearly: for \(c>0\), the Taylor quotients of \(cf\) are \(c\) times those of \(f\), so \(M_j(cf)=cM_j(f)\) and \(M_{\rm av}(cf)=cM_{\rm av}(f)\) (proved in the section on smoothness scaling, which needs \(f\) differentiable with Lipschitz gradient). With \(c=1/n\), \(M_{\rm av}(f/n)=M_{\rm av}(f)/n\), and the first form becomes the second. \(\square\)




The statement holds for every initial probability measure: if \(\mathsf W_2(\nu,\pi)=\infty\), the right-hand side is \(\infty\).

\begin{leanbox}{Lean: \texttt{sgld\_\allowbreak{}constant\_\allowbreak{}step\_\allowbreak{}extended}, \texttt{sgld\_\allowbreak{}constant\_\allowbreak{}step\_\allowbreak{}paper}}
\leansource{\texttt{sgld\_\allowbreak{}constant\_\allowbreak{}step\_\allowbreak{}extended} \quad AverageLMC/\allowbreak{}SGLDTheorem.lean, lines 206–235}
\begin{leancode}
\leanline{0}{theorem\ sgld\_constant\_step\_extended\ \{p\ n\ :\ ℕ\}\ (hp\ :\ 0\ <\ p)\ (hn\ :\ 0\ <\ n)}
\leanline{4}{(f\ :\ Fin\ n\ →\ State\ p\ →\ ℝ)\ (hf\ :\ ∀\ i,\ Differentiable\ ℝ\ (f\ i))}
\leanline{4}{(L\ :\ Fin\ n\ →\ ℝ≥0)\ (hLip\ :\ ∀\ i,\ LipschitzWith\ (L\ i)\ (gradient\ (f\ i)))}
\leanline{4}{(mbar\ h\ :\ ℝ)\ (hmbar\ :\ 0\ <\ mbar)\ (hh\ :\ 0\ ≤\ h)}
\leanline{4}{(hsc\ :\ StrongConvexOn\ univ\ ((n\ :\ ℝ)\ *\ mbar)\ (componentPotential\ f))}
\leanline{4}{(hstep\ :\ h\ ≤\ mbar\ /\ (3\ *\ n\ *\ rmsComponentSmoothness\ f\ \textasciicircum{}\ 2))}
\leanline{4}{(anchor\ :\ State\ p)\ (hanchor\ :\ ∀\ x,\ componentPotential\ f\ anchor\ ≤\ componentPotential\ f\ x)}
\leanline{4}{(μ\ :\ Measure\ (State\ p))\ [IsProbabilityMeasure\ μ]\ (k\ :\ ℕ)\ :}
\leanline{4}{wasserstein2\ (sgldLaws\ f\ anchor\ h\ μ\ k)\ (gibbsMeasure\ (componentPotential\ f))\ \textasciicircum{}\ 2\ ≤}
\leanline{6}{ENNReal.ofReal\ ((1\ +\ (n\ :\ ℝ)\ *\ mbar\ *\ h)⁻¹\ \textasciicircum{}\ k)\ *}
\leanline{8}{wasserstein2\ μ\ (gibbsMeasure\ (componentPotential\ f))\ \textasciicircum{}\ 2\ +}
\leanline{6}{ENNReal.ofReal\ ((averageSmoothness\ (componentPotential\ f)\ /\ n\ +\ 2\ *\ mbar\ +}
\leanline{8}{2\ *\ rmsComponentSmoothness\ f\ \textasciicircum{}\ 2\ /\ mbar)\ *\ h\ *\ p\ /\ mbar)}
\end{leancode}
\leansource{\texttt{sgld\_\allowbreak{}constant\_\allowbreak{}step\_\allowbreak{}paper} \quad AverageLMC/\allowbreak{}SGLDTheorem.lean, lines 239–256}
\begin{leancode}
\leanline{0}{theorem\ sgld\_constant\_step\_paper\ \{p\ n\ :\ ℕ\}\ (hp\ :\ 0\ <\ p)\ (hn\ :\ 0\ <\ n)}
\leanline{4}{(f\ :\ Fin\ n\ →\ State\ p\ →\ ℝ)\ (hf\ :\ ∀\ i,\ Differentiable\ ℝ\ (f\ i))}
\leanline{4}{(L\ :\ Fin\ n\ →\ ℝ≥0)\ (hLip\ :\ ∀\ i,\ LipschitzWith\ (L\ i)\ (gradient\ (f\ i)))}
\leanline{4}{(mbar\ h\ :\ ℝ)\ (hmbar\ :\ 0\ <\ mbar)\ (hh\ :\ 0\ ≤\ h)}
\leanline{4}{(hsc\ :\ StrongConvexOn\ univ\ ((n\ :\ ℝ)\ *\ mbar)\ (componentPotential\ f))}
\leanline{4}{(hstep\ :\ h\ ≤\ mbar\ /\ (3\ *\ n\ *\ rmsComponentSmoothness\ f\ \textasciicircum{}\ 2))}
\leanline{4}{(anchor\ :\ State\ p)\ (hanchor\ :\ ∀\ x,\ componentPotential\ f\ anchor\ ≤\ componentPotential\ f\ x)}
\leanline{4}{(μ\ :\ Measure\ (State\ p))\ [IsProbabilityMeasure\ μ]\ (k\ :\ ℕ)\ :}
\leanline{4}{wasserstein2\ (sgldLaws\ f\ anchor\ h\ μ\ k)\ (gibbsMeasure\ (componentPotential\ f))\ \textasciicircum{}\ 2\ ≤}
\leanline{6}{ENNReal.ofReal\ ((1\ +\ (n\ :\ ℝ)\ *\ mbar\ *\ h)⁻¹\ \textasciicircum{}\ k)\ *}
\leanline{8}{wasserstein2\ μ\ (gibbsMeasure\ (componentPotential\ f))\ \textasciicircum{}\ 2\ +}
\leanline{6}{ENNReal.ofReal\ ((averageSmoothness\ (fun\ x\ ↦\ componentPotential\ f\ x\ /\ n)\ +}
\leanline{8}{2\ *\ mbar\ +\ 2\ *\ rmsComponentSmoothness\ f\ \textasciicircum{}\ 2\ /\ mbar)\ *\ h\ *\ p\ /\ mbar)}
\end{leancode}
\leanlabel{How the Lean proof follows the argument}
\begin{enumerate}


\item \texttt{sgld\_\allowbreak{}constant\_\allowbreak{}step\_\allowbreak{}extended} states the conclusion in \texttt{\ensuremath{\mathbb R}\ensuremath{\ge}0\ensuremath{\infty}}, with the real coefficients embedded by \texttt{ENNReal.\allowbreak{}of\allowbreak{}Real}. After establishing \texttt{h\allowbreak{}A}, it applies \texttt{squared\_\allowbreak{}transport\_\allowbreak{}extended} with \(a=(1+n\bar mh)^{-k}\); the two sign conditions are discharged by \texttt{positivity} (which uses \texttt{h\allowbreak{}A} for \(b\ge0\)), and the two remaining goals are \texttt{sgld\allowbreak{}Laws\_\allowbreak{}finite} and \texttt{sgld\_\allowbreak{}constant\_\allowbreak{}step\_\allowbreak{}average}.

\item \texttt{sgld\_\allowbreak{}constant\_\allowbreak{}step\_\allowbreak{}paper} writes the constant as \texttt{average\allowbreak{}Smoothness (fun x \ensuremath{\mapsto} component\allowbreak{}Potential f x /\allowbreak{} n)}, the literal \(\bar M_{\rm av}\) of the average potential. It rewrites this constant with \texttt{average\allowbreak{}Smoothness\_\allowbreak{}div\_\allowbreak{}nat} (using the Lipschitz constant \(\sum_iL_i\) of \(\nabla f\) from \texttt{component\allowbreak{}Potential\_\allowbreak{}gradient\_\allowbreak{}lipschitz}) and concludes with \texttt{sgld\_\allowbreak{}constant\_\allowbreak{}step\_\allowbreak{}extended}.

\end{enumerate}
\end{leanbox}


Earlier results used here: \cite[Section~3.17.3]{LangevinLean2026ULASC}; \cite[Section~3.26.3]{LangevinLean2026ULASC}.
\section{\texorpdfstring{Main average-smoothness sampling theorems}{Main average-smoothness sampling theorems}}
\label{proof-module-AverageLMC-lean}


This import bundle exports the Gaussian LMC, decreasing-step LMC, Laplacian LMC and control-variate SGLD bounds. Their coordinate-average smoothness constants, analytic estimates and distributional recurrences are proved. The source inventory distinguishes the main sampling results from broader auxiliary coverage.

\section{\texorpdfstring{Audit the main sampling proofs and analytic dependencies}{Audit the main sampling proofs and analytic dependencies}}
\label{proof-module-AverageLMC-Check-lean}


The audit brings together the sampling convergence theorems and the analytic estimates used in their proofs.


\begin{thebibliography}{ULA-LSI}
\bibitem[ULA-SC]{LangevinLean2026ULASC}
LangevinLean Project. \emph{Unadjusted Langevin under Strong Convexity}. LL-01, 2026. Draft report, content revision d50c0f64815b. \url{https://odakyildiz.github.io/LangevinLean/docs/pdf/UnadjustedLangevin/FullProof.pdf}.
\bibitem[ULA-LSI]{LangevinLean2026ULALSI}
LangevinLean Project. \emph{Unadjusted Langevin under LSI}. LL-02, 2026. Draft report, content revision 7046bc2a17e4. \url{https://odakyildiz.github.io/LangevinLean/docs/pdf/UnadjustedLangevin/LSI/FullProof.pdf}.
\bibitem[IPLA-SC]{LangevinLean2026IPLA}
LangevinLean Project. \emph{Interacting Particle Langevin under Strong Convexity}. LL-03, 2026. Draft report, content revision bfbb88fd1d79. \url{https://odakyildiz.github.io/LangevinLean/docs/pdf/IPLA/FullProof.pdf}.
\bibitem[Proximal]{LangevinLean2026Proximal}
LangevinLean Project. \emph{Proximal particle Langevin}. LL-05, 2026. Draft report, content revision f201d59ac57b. \url{https://odakyildiz.github.io/LangevinLean/docs/pdf/ProximalLangevin/FullProof.pdf}.
\end{thebibliography}

\end{document}
